% #21 | Archive/SP_SauLuoi_2024-2025/SP_HKI/1-TN-DS-TLN-THPT-ChuyenHungVuong-PhuTho-HKI-NH24-25.tex | mã gốc 1D1V5-6
\begin{ex}%[1D1V5-6]%[Dự án đề kiểm tra Toán khối 11 học kì 1 NH24-25-Dot 5-Thái Bảo ]%[THPT Chuyên Hùng Vương - Phú Thọ]
	Một quả bóng được ném xiên một góc $\alpha\left(0^{\circ} < \alpha < 90^{\circ}\right)$ từ mặt đất với tốc độ $v_0$ (m/s). Khoảng cách theo phương ngang từ vị trí ban đầu của quả bóng đến vị trí bóng chạm đất được tính theo công thức $d=\dfrac{v_0^2\sin 2\alpha}{10}$. Nếu tốc độ ban đầu của quả bóng là $10$ m/s thì tồn tại hai góc ném $\alpha_1$, $\alpha_2$ $\left(\alpha_1 > \alpha_2\right)$ để khoảng cách $d$ là $5$ m. Hiệu số đo hai góc ném $\alpha_1-\alpha_2$ bằng bao nhiêu độ?
	\shortans{$60$}
	\loigiai{
		Theo giả thuyết bài toán ta có 
		$$\dfrac{10^2\cdot \sin 2\alpha}{10}=5\Leftrightarrow\sin 2\alpha=\dfrac{1}{2}\Leftrightarrow\heva{&2\alpha =30^\circ + k360^\circ\\&2\alpha=150^\circ+k360^\circ}\Leftrightarrow\heva{&\alpha =15^\circ + k180^\circ\\&\alpha=75^\circ+k180^\circ}\,(k\in\mathbb{Z}).$$
		Vì $0^{\circ} < \alpha < 90^{\circ}$ và $\alpha_1 > \alpha_2$ nên $\alpha_1=75^\circ$, $\alpha_2=15^\circ$.\\
		Suy ra $\alpha_1-\alpha_2=75^\circ-15^\circ=60^\circ$.
	}
\end{ex}

% #38 | Archive/SP_SauLuoi_2024-2025/SP_HKI/1-TN-TL-THPT-MacDinhChi-TPHCM-HKI-NH24-25.tex | mã gốc 1D1H5-3
\begin{ex}%[1D1H5-3]%[Dự án đề kiểm tra Toán 11 HKI NH24-25 - Đoàn Minh Tâm]%[THPT Mạc Đĩnh Chi]
	Giải phương trình $ 3\tan \left(x-\dfrac{\pi}{6}\right)-\sqrt 3=0$.
	\loigiai{
			Ta có
			\begin{eqnarray*}
				&3 \tan \left(x-\dfrac{\pi}{6}\right)-\sqrt{3}=0 &\Leftrightarrow \tan \left(x-\dfrac{\pi}{6}\right)=\dfrac{\sqrt{3}}{3}=\tan \dfrac{\pi}{6}\\
				& &\Leftrightarrow x-\dfrac{\pi}{6}=\dfrac{\pi}{6}+k \pi\\
				& &\Leftrightarrow x=\dfrac{\pi}{3}+k \pi \quad (k \in \mathbb{Z}).
			\end{eqnarray*}
	}
\end{ex}

% #46 | Archive/SP_SauLuoi_2024-2025/SP_HKI/1-TL-TN-DS-THPT-Hoang-Viet-Daklak-HKI-NH24-25.tex | mã gốc 1D1V5-6
\begin{bt}%[1D1V5-6]%[Dự án đề kiểm tra Toán 11 HKI NH24-25-Dot5-Thanh Phong]%[TH,THCS & THPT Hoàng Việt - Đăklăk]
	Nhiệt độ ngoài trời ở một thành phố vào các thời điểm khác nhau trong ngày có thể
	được mô phỏng bởi công thức $h(t)=29+3\sin\left[\dfrac{\pi}{12}\left(t-9\right)\right]$ với $h$ tính bằng $^\circ$C và $t$ là thời gian trong ngày tính bằng giờ. Nhiệt độ trong ngày là thấp nhất là bao nhiêu?
\loigiai{
	Ta có $-1\le\sin\left[\dfrac{\pi}{12}\left(t-9\right)\right]\leq 1\Leftrightarrow 26\leq 29+3\sin\left[\dfrac{\pi}{12}\left(t-9\right)\right]\leq 32$,\, $\forall t\ge 0$.\\
	Vậy nhiệt độ thấp nhất trong ngày là $26^\circ$C, xảy ra khi
	$$\sin\left[\dfrac{\pi}{12}\left(t-9\right)\right]=-1\Leftrightarrow\dfrac{\pi}{12}\left(t-9\right)=-\dfrac{\pi}{2}+k2\pi\Leftrightarrow t=3+24k,\ k\in\mathbb{Z}.$$
	Do $t$ là thời gian trong ngày tính bằng giờ nên $0\leq t\le 24\Rightarrow k=0$ .\\
	Suy ra $t=3$ .\\
	Vậy lúc 3 giờ thì nhiệt độ trong ngày là thấp nhất.}
\end{bt}

% #60 | Archive/SP_SauLuoi_2024-2025/SP_HKI/1-TN-DS-TLN-ChuyenLeQuyDon-NinhThuan-HKI-NH24-25.tex | mã gốc 1D1H5-3
\begin{ex}%[1D1H5-3]%[1-TN-DS-TLN-ChuyenLeQuyDon-NinhThuan-HKI-NH24-25]%[Thầy Hải Toán]
	Phương trình $\cos 2x=\dfrac{1}{2}$ có bao nhiêu nghiệm thuộc khoảng $\left(0; \dfrac{3\pi}{2}\right)$?
	\loigiai{
		Ta có \allowdisplaybreaks
		\begin{eqnarray*}
			&&\cos 2x=\dfrac{1}{2}\\
			&\Leftrightarrow& \cos{2x}=\cos{\dfrac{\pi}{3}}\\
			&\Leftrightarrow& \hoac{& 2x=\dfrac{\pi}{3}+k2\pi \\ & 2x=-\dfrac{\pi}{3}+\ell2\pi}\\
			&\Leftrightarrow& \hoac{& x=\dfrac{\pi}{6}+k\pi \\ & x=-\dfrac{\pi}{6}+\ell\pi}\; k,\ell\in \mathbb{Z}.
		\end{eqnarray*}
		Xét $0 < \dfrac{\pi}{6}+k\pi < \dfrac{3\pi}{2}\Leftrightarrow -\dfrac{1}{6}< k< \dfrac{4}{3}\xrightarrow{\quad k\ \in\ \mathbb{Z}\quad } k\in\left\{0;1\right\}$.\\
		$0<-\dfrac{\pi}{6}+\ell\pi<\dfrac{3\pi}{4}\Leftrightarrow \dfrac{1}{6}<\ell<\dfrac{5}{3}\xrightarrow{\quad  \ell\ \in \ \mathbb{Z}\quad} \ell\in\left\{1\right\}$.\\
		Với $k=0\Rightarrow x=\dfrac{\pi}{6}$; $k=1\Rightarrow x=\dfrac{7\pi}{6}$.\\
		Với $\ell=1\Rightarrow x=\dfrac{5\pi}{6}$.\\
		Vậy phương trình có $3$ nghiệm thuộc khoảng $\left(0;\dfrac{3\pi}{2}\right)$.
	}
\end{ex}

% #75 | Archive/SP_SauLuoi_2024-2025/SP_HKI/1-TN-DS-TLN-THPT-NguyenThaiBinh-HCM-HKI-NH24-25.tex | mã gốc 1D1H4-5
\begin{ex}%[1D1H4-5]%[Dự án đề kiểm tra Toán 11 HKI NH24-25-THPT Nguyễn Thái Bình - Tp HCM]%[Trần Ngọc Thành]
Hằng ngày mực nước của con kênh lên xuống theo thủy triều. Độ sâu $h$ (mét) của mực nước trong kênh tính theo thời gian $t$ (giờ) được cho bởi công thức $h(t)=3\cos\left(\dfrac{\pi t}{6}+\dfrac{\pi}{3}\right)+10$. Hỏi có bao nhiêu thời điểm trong ngày độ sâu của mực nước trong kênh là nhỏ nhất.
\shortans[4]{2}
\loigiai{
	Vì $\cos\left(\dfrac{\pi t}{6}+\dfrac{\pi}{3}\right) \ge -1, \forall t$ nên $h(t)=3\cos\left(\dfrac{\pi t}{6}+\dfrac{\pi}{3}\right)+10 \ge 7$ (m).\\
	Thời điểm trong ngày độ sâu của mực nước trong kênh là nhỏ nhất ứng với 
	$$\cos\left(\dfrac{\pi t}{6}+\dfrac{\pi}{3}\right) =-1\Leftrightarrow \dfrac{\pi t}{6}+\dfrac{\pi}{3}=\pi+k2\pi\Leftrightarrow t=4+12k,\,(k\in\mathbb{Z}).$$
	Vì $0<t\le 24$ nên $k\in {0,1}$.\\
	Như vậy, có $2$ thời điểm trong ngày độ sâu của mực nước trong kênh là nhỏ nhất.
}
\end{ex}

% #88 | Archive/SP_SauLuoi_2024-2025/SP_HKI/1-TN-DS-TLN-THPT-ChuyenLeHongPhong-TpHCM-HKI-NH24-25.tex | mã gốc 1D1V5-5
\begin{bt}%[1D1V5-5][Dự án đề kiểm tra môn Toán khối 11 HKI NH24-25 - Đợt 5 - Hieu Hieu Minh Minh]%[THPT Chuyên Lê Hồng Phong - Tp HCM]
	Phương trình $\cos{2x} - 3\cos{x} + 2 = 0$ có bao nhiêu nghiệm trên khoảng $\left(\dfrac{\pi}{2}; 4\pi\right)$?
	\loigiai{
	\begin{eqnarray*}
		&&\cos {2x}-3 \cos x + 2 = 0 \\
		&\Leftrightarrow& 2 \cos^2 x-1-3 \cos x+2=0 \\
		&\Leftrightarrow& 2 \cos ^2 x-3 \cos x+1=0\\
		&\Leftrightarrow& \hoac{&\cos x = 1 \\&\cos x = \dfrac{1}{2}}\\
		&\Leftrightarrow& \hoac{&x=k2\pi \\&x= \pm \dfrac{\pi}{3}+k2\pi (k \in \mathbb{Z})}\\
	\end{eqnarray*}
	Do $x \in (\dfrac{\pi}{2}; 4\pi)$ nên $x\in \left\{ 2\pi; \dfrac{7\pi}{3}; \dfrac{5\pi}{3}; \dfrac{11\pi}{3}\right\}$.\\
	Vậy phương trình đã cho có $4$ nghiệm thuộc $\left(\dfrac{\pi}{2}; 4\pi\right)$.
	}
\end{bt}

% #93 | Archive/SP_SauLuoi_2024-2025/SP_GHKI/1-TN-DS-TLN-THPT-TayThanh-HCM-GHKI-NH24-25.tex | mã gốc 1D1V5-6
\begin{ex}%[1D1V5-6]%[Dự án đề kiểm tra Toán 11 GHKI NH24-25- Lê Hữu Kiệt]%[THPT Tây Thạnh- Tp HCM]
	Giả sử nhiệt độ bên trong một ngôi nhà sau $t$ (giờ) kể từ $12$ giờ trưa, gọi là $T(t)$ được tính bởi $T(t)=5\cos\left(\dfrac{\pi}{2}-\dfrac{\pi t}{6}\right)+25$ ($^\circ$C) với $0\leq t \leq 24$.
	\choiceTF
		{$T(0)=29$}
		{\True $T(12)=T(6)$}
		{Nhiệt độ cao nhất trong ngôi nhà là $36^\circ$C}
		{\True Nhiệt độ trong ngôi nhà lúc $5$ giờ chiều là $27{,}5^\circ$C}
	\loigiai{
	\begin{itemchoice}
		\itemch \textbf{Sai.}\\
		Ta có $T(0)=25$.
		\itemch \textbf{Đúng.}\\
		Ta có $T(12)=25$ và $T(6)=25$. Vậy $T(12)=T(6)$.
		\itemch \textbf{Sai.}\\
		Do $-1\leq \cos\left(\dfrac{\pi}{2}-\dfrac{\pi t}{6}\right) \leq 1$ nên $20\leq 5\cos\left(\dfrac{\pi}{2}-\dfrac{\pi t}{6}\right)+25 \leq 30$.\\\
		Suy ra nhiệt độ cao nhất trong ngôi nhà là $30^\circ$C khi $\cos\left(\dfrac{\pi}{2}-\dfrac{\pi t}{6}\right)=1$ hay $t=3$, $t=15$.
		\itemch \textbf{Đúng.}\\
		Lúc $5$ giờ chiều, tức sau $12$ giờ trưa là $5$ giờ, suy ra $t=5$.\\
		Ta có $T(5)=27{,}5^\circ$C.
	\end{itemchoice}
	}
\end{ex}

% #95 | Archive/SP_SauLuoi_2024-2025/SP_GHKI/1-TN-DS-TLN-THPT-TayThanh-HCM-GHKI-NH24-25.tex | mã gốc 1D1C4-8
\begin{ex}%[1D1C4-8]%[GHKI-THPT Tây Thạnh -2425 ]%[Hiếu Phan]
    \immini{Giả sử một con tàu vũ trụ được phóng lên từ mũi Canaveral ở Mĩ. Nó chuyển động theo một quỹ đạo được mô tả trên một bản đồ phẳng (quanh đường xích đạo) của mặt đất như hình mô phỏng bên dưới. Điểm $M$ mô tả cho con tàu, đường thẳng $\Delta$ mô tả cho đường xích đạo. Khoảng cách $h$ (km) từ $M$ đến $\Delta$ được tính theo công thức $h=|d|$, trong đó $d=4000\cos\left[\dfrac{\pi}{45}(t-10)\right]$, với $t$ (phút) là thời gian trôi qua kể từ khi con tàu đi vào quỹ đạo, $d>0$ nếu $M$ ở phía trên $\Delta$, $d<0$ nếu $M$ ở phía dưới $\Delta$. Hãy tìm thời điểm sớm nhất sau khi con tàu đi vào quỹ đạo để có $d$ lớn nhất.}{\begin{tikzpicture}[line width=2pt,font=\footnotesize, line join=round, line cap=round, >=stealth]
            \fill[blue!20](-5,-5)rectangle(5,2.6);
            \fill[teal!50!green](-4.65,0)..controls++(45:0.5)and++(160:0.75)..(-3.9,0)--++(-60:0.5)--++(-91:0.5)--++(-45:0.75)--++(-90:0.15)--++(-25:0.85)--++(-90:0.2)--++(-95:0.35)--++(0:0.1)--++(-100:0.2)--++(-50:0.35)--++(-170:0.1)--++(-92:1.35)--++(-5:0.2)--++(30:0.1)--++(88:0.45)--++(70:0.1)--++(60:0.3)--++(5:0.1)--++(50:0.2)--++(62:0.2)--++(50:0.2)--++(88:0.3)--++(40:0.2)--++(88:0.1)--++(130:0.2)--++(160:0.25)--++(134:0.15)--++(85:0.1)--++(170:0.2)--++(150:0.2)--++(175:0.2)--++(155:0.1)--++(160:0.1)--++(145:0.2)--++(150:0.1)--++(90:0.2)--++(-135:0.2)--++(140:0.2)--++(90:0.2)--++(40:0.1)--++(0:0.1)--++(30:0.2)--++(-88:0.2)--++(75:0.3)--++(20:0.2)--++(92:0.2)--++(35:0.2)--++(90:0.1)--++(0:0.15)--++(70:0.2)--++(140:0.2)--++(30:0.1)--++(-20:0.1)--++(40:0.2)--++(145:0.3)--++(90:0.2)--++(-140:0.2)--++(140:0.3)--++(-98:0.55)--++(135:0.2)--++(155:0.2)--++(110:0.3)--++(45:0.4)--++(140:0.5)--++(40:0.1)--++(-35:0.4)--++(90:0.1)--++(-30:0.3)--++(-95:0.2)--++(0:0.2)--++(90:0.2)--++(120:0.2)--++(170:0.3)--++(70:0.55)--++(90:0.1)--++(25:0.2)--++(-90:0.1)--++(-135:0.1)--++(-50:0.2)--++(0:0.2)--++(-75:0.3)--++(-90:0.1)--++(-80:0.2)--++(-70:0.45)--++(-10:0.15)--++(70:0.3)--++(60:0.2)--++(0:0.2)--++(60:0.3)--++(160:0.1)--++(70:0.1)--++(-20:0.1)--++(82:0.5)--++(20:0.15)--++(100:0.1)--++(-170:0.2)--++(65:0.2)--++(-170:0.2)--++(170:0.2)--++(160:0.15)--++(-145:0.3)--++(-175:0.4)--++(90:0.1)--++(-165:0.75)--++(-90:0.2)--++(-40:0.2)--++(-160:0.2)--++(-170:0.5)--++(-90:0.5)--++(130:0.2)--++(-182:0.5)--++(-150:0.2)--++(160:0.5)--++(-175:0.2)--++(-140:0.2)--++(-90:0.2)--++(-10:0.1)--++(-90:0.2)--++(-70:0.2)--cycle;
            \fill[teal!60!green](0,-0.8)--++(178:0.2)--++(-170:0.2)--++(-130:0.2)--++(-110:0.3)--++(-115:0.2)--++(-86:0.4)--++(-20:0.2)--++(-88:0.1)--++(-30:0.2)--++(40:0.2)--++(-30:0.3)--++(-92:0.1)--++(-80:0.4)--++(-92:0.3)--++(-45:0.2)--++(-85:0.4)--++(25:0.3)--++(60:0.4)--++(100:0.2)--++(30:0.2)--++(90:0.4)--++(60:0.6)--++(135:0.1)--++(180:0.2)--++(130:0.1)--++(115:0.6)--++(-40:0.5)--++(-90:0.1)--++(20:0.4)--++(120:0.2)--++(-140:0.1)--++(130:0.2)--++(0:0.3)--++(-40:0.4)--++(-70:0.4)--++(55:0.65)--++(-80:0.3)--++(10:0.2)--++(-80:0.2)--++(0:0.1)--++(95:0.1)--++(-40:0.15)--++(20:0.2)--++(115:0.3)--++(20:0.3)--++(80:0.2)--++(0:0.1)--++(110:0.4)--++(90:0.2)--++(-25:0.1)--++(-80:0.25)--++(0:0.15)--++(80:0.3)--++(0:0.1)--++(70:0.5)--++(130:0.2)--++(65:0.3)--++(10:0.3)--++(80:0.1)--++(0:0.3)--++(-90:0.1)--++(25:0.4)--++(90:0.1)--++(30:0.1)--++(-20:0.2)--++(90:0.2)--++(170:0.4)--++(180:0.3)--++(160:0.2)--++(165:0.4)--++(-160:0.3)--++(150:0.2)--++(165:0.2)--++(180:0.3)--++(50:0.2)--++(140:0.3)--++(-150:0.4)--++(-150:0.3)--++(-130:0.5)--++(-170:0.75)--++(-130:0.3)--++(130:0.2)--++(40:0.2)--++(150:0.3)--++(190:0.3)--++(-110:0.4)--++(180:0.2)--++(-90:0.2)--++(0:0.2)--++(-90:0.1)--++(40:0.2)--++(90:0.2)--++(60:0.2)--++(-90:0.3)--++(-20:0.2)--++(-140:0.4)--++(180:0.3)--++(-130:0.3)--++(-92:0.3)--++(180:0.2)--++(-90:0.2)--++(0:0.2)--++(45:0.3)--++(-45:0.3)--++(45:0.1)--++(130:0.2)--++(30:0.2)--++(-50:0.3)--++(-30:0.3)--++(-90:0.2)--++(175:0.5)--cycle;
            \fill[teal!60!green](3,-3)--++(80:0.2)--++(0:0.1)--++(85:0.2)--++(30:0.2)--++(20:0.2)--++(-45:0.2)--++(60:0.15)--++(-40:0.3)--++(-60:0.3)--++(-90:0.1)--++(-30:0.1)--++(-100:0.4)--++(-145:0.1)--++(-90:0.2)--++(130:0.3)--++(170:0.2)--++(90:0.2)--++(160:0.2)--++(-130:0.3)--++(180:0.2)--++(90:0.2)--cycle;
            \draw[blue](-2.5,-1.3)..controls++(50:3)and++(-140:3)..(5,-3)coordinate[pos=0.2](M);
            \fill[red](M)circle(3pt)node[above]{\textbf{M}};
            \path(M)node[below=2mm,scale=1.25,red]{\textbf{Canaveral}};
            \draw[dashed,] (-5,-2.2)--(5,-2.2);
            \path(3.5,-2.3)node[above=0mm,scale=1.25]{\textbf{Xích đạo}};
    \end{tikzpicture}}
    \shortans[]{$10$}
    \loigiai
    {\begin{itemize}
            \item Trường hợp 1: Xét $M$ ở phía trên $\Delta$, khi đó thời điểm tàu đi vào quỹ đạo có $d=4000$ là lớn nhất với điều kiện 
            $$\cos\left[\dfrac{\pi}{45}(t-10)\right]=1\Leftrightarrow \dfrac{\pi}{45}(t-10)=k2\pi\Leftrightarrow t=10+90k, k\in\mathbb{Z}.$$
            Do $t>0$ nên thời điểm sớm nhất trong trường hợp này là $t=10$ (phút).
            \item Trường hợp 2: Xét $M$ ở phía dưới $\Delta$, khi đó thời điểm tàu đi vào quỹ đạo khi $d=-4000$ thì $h$ là lớn nhất với điều kiện 
            $$\cos\left[\dfrac{\pi}{45}(t-10)\right]=-1\Leftrightarrow \dfrac{\pi}{45}(t-10)=\pi+k2\pi\Leftrightarrow t=55+90k, k\in\mathbb{Z}.$$
            Do $t>0$ nên thời điểm sớm nhất trong trường hợp này là $t=55$ (phút).
        \end{itemize}
        Vậy sau $10$ phút kể từ khi con tàu đi vào quỹ đạo thì có $d$ lớn nhất.
    }
\end{ex}

% #96 | Archive/SP_SauLuoi_2024-2025/SP_GHKI/1-TN-DS-TLN-THPT-TayThanh-HCM-GHKI-NH24-25.tex | mã gốc 1D1H5-6
\begin{ex}%[1D1H5-6]%[GHKI-THPT Tây Thạnh -2425 ]%[Hiếu Phan]
    \immini{Một người dẫn em gái của mình đến công viên chơi xích đu. Lực đẩy $F$ (N) theo phương ngang  mà người đó dùng để đẩy em gái trong trò chơi này (ở vị trí $A$ trên hình vẽ) được xác định bởi công thức $F=mg\cdot\tan\theta$, trong đó $m$ (kg) là khối lượng của em gái, $g$ là gia tốc trọng trường và $\theta$ là góc tạo bởi xích đu khi bắt đầu được đẩy với  phương thẳng đứng $(0^\circ<\theta<90^\circ)$ được minh hoạ như hình bên. Xác định góc $\theta$ (đơn vị độ) khi $F=400\sqrt{3}$ N, $m=40$ kg và $g=10$ m/s$^2$.
        \shortans[]{$60$}
    }{
        \begin{tikzpicture}[scale=.8,font=\footnotesize,line join = round, line cap = round, >= stealth]
            \def \a{2.5}
            %\draw[opacity=0.3] (-5,0) grid (5,5);
%            \path (0,0) node{\hypersetup{hidelinks}\href{Twhb8wr}{ }};
            \draw[line width =5pt,blue] (-\a,0)--(\a,0); 
            \draw[line width =2pt,violet] (-2.5,0)--(0,5)--(2.5,0); 
            \draw[line width =1pt,violet] (-3,2)--(0,5); 
            \draw[fill] (-3,2)  circle (3pt);
            \draw (-3,2.4) node{$A$};
            \draw[->,line width =1pt,violet] (-3,2)--(-2,1.7);
            \draw (0,2.5)--(0,5);
            \clip (-3,2)--(0,5)--(0,0);
            \draw [<->] (0,5) circle (2cm);
            \draw (-0.5,3.5) node[]{$\theta$}; 
            \foreach \p/\g in {} \draw[fill] (\p) circle(.5pt)node [shift={(\g:.3)}] {$\p$};
    \end{tikzpicture}
}
    \loigiai
    {
        Ta có $F=mg\cdot\tan\theta\Rightarrow \tan\theta=\dfrac{F}{mg}=\dfrac{400\sqrt{3}}{40\cdot 10}=\sqrt{3}$.\\
        Suy ra $\theta =60^\circ$ do $0^\circ<\theta<90^\circ$.
    }
\end{ex}

% #103 | Archive/SP_SauLuoi_2024-2025/SP_GHKI/1-TN-DS-TLN-THPT-NguyenBinhKhiem-HaNoi-GHKI-NH23-24.tex | mã gốc 1D1H5-6
\begin{ex}%[1D1H5-6]%[Dự án đề kiểm tra Toán khối 11 HK1 NH23-24-Dot-15-Huỳnh Thanh Chí]%[Nguyễn Bĩnh Khiêm Hà Nội]
	Thời gian tập thể dục của ông A vào ngày thứ $x$ trong một tháng được tính bằng công thức $t(x)=\left|\tan \left(\dfrac{\pi}{3} x\right)\right|$ giờ. Hỏi trong một tháng $30$ ngày) có bao nhiêu ngày ông A không tập thể dục?
	
	\shortans[]{$10$}
	\loigiai{
		Để xác định số ngày ông A không tập thể dục, tức là tìm số ngày mà $t(x) = 0$, ta có
		\allowdisplaybreaks
		\begin{eqnarray*}
			&& t(x)= 0\\
			&\Leftrightarrow& \left|\tan \left(\dfrac{\pi}{3} x\right)\right| =0\\
			&\Leftrightarrow& \tan \left(\dfrac{\pi}{3} x\right)=0\\
			&\Leftrightarrow& \dfrac{\pi}{3}x=k\pi\\
			&\Leftrightarrow& x=3k\,\left(k\in\mathbb{Z}\right).
		\end{eqnarray*}
		Vì $x$ là số nguyên và $x\in [1;30]$, nên
		\[
		1 \leq 3k \leq 30 \Rightarrow \dfrac{1}{3} \leq k \leq 10
		\]
		Vì $k$ là số nguyên, nên $k\in\left\{1;2;\ldots;9;10\right\}$.\\
		Vậy, có tất cả $10$ ngày trong tháng mà ông A không tập thể dục.
	}
\end{ex}

% #111 | Archive/SP_SauLuoi_2024-2025/SP_GHKI/1-TN-DS-TLN-chuyenLeQuyDon-NinhThuan-GHKI-NH24-25..tex | mã gốc 1D1V4-8
\begin{ex}%[1D1V4-8]%[Dự án đề kiểm tra Toán 10 - 11 GHKI NH24-25- Thái Bảo]%[THPT Chuyên Bình Thuận]
	\immini{
		Một vật dao động xung quanh vị trí cân bằng theo phương trình $x=1{,}5\cos \left(\dfrac{t\pi}{4} \right)$ (tham khảo hình vẽ); trong đó $t$ là thời gian được tính bằng giây và quãng đường $h=|x|$ được tính bằng mét là khoảng cách theo phương ngang của chất điểm đối với vị trí cân bằng. Khi đó
	}{
		\definecolor{darkpastelgreen}{rgb}{0.01, 0.75, 0.24}
		\definecolor{almond}{rgb}{0.94, 0.87, 0.8}
		\definecolor{arsenic}{rgb}{0.23, 0.27, 0.29}
		\definecolor{eggplant}{rgb}{0.38, 0.25, 0.32}
		\definecolor{crimson}{rgb}{0.86, 0.08, 0.24}
		\begin{tikzpicture}[line join=round, line cap=round,scale=.5,transform shape]
			\clip (-2.5,-2.5) rectangle (4.5,4.2);
			\tikzset{q_h/.pic={
					% Tre
					\draw[fill=darkpastelgreen] (-.05,4)--(.9,-1.8)
					..controls +(30:.05) and +(-30:.05) ..(1,-1.75)--(.05,4);
					\draw[fill=darkpastelgreen] (-.1,4)--(.6,-1.4)
					..controls +(30:.05) and +(-30:.05) ..(.5,-1.4)--(-.2,4)
					;
					%----------------------
					\def\M{ 
						(-.2,1.76)
						..controls +(-80:.1) and +(80:.1) ..(-.18,1.6)
						..controls +(-40:.1) and +(30:.06) ..(-.18,1.5)--(-.17,1.46)--(-.22,1.44)--(-.19,1.43)--(-.2,1.43)
						..controls +(-80:.1) and +(-12:.05) ..(-.3,1.35)--(-.35,1.28)--(-.55,1.35)
						..controls +(60:.1) and +(-90:.05) ..(-.45,1.5)
						..controls +(100:.1) and +(140:.05) ..(-.38,1.55)
						..controls +(70:.1) and +(-140:.05) ..cycle
						;
					}
					
					\draw[black]\M;
					\fill[almond] \M;
					\draw (-.2,1.76)
					..controls +(-140:.1) and +(50:.05) ..(-.3,1.6)
					(-.2,1.76)
					..controls +(-120:.1) and +(40:.05) ..(-.25,1.6)
					(-.28,1.578)
					..controls +(-20:.02) and +(-150:.02) ..(-.22,1.58)
					(-.25,1.568)--(-.23,1.567);
					%------------------Áo
					\def\A3{ 
						(-.1,.9)
						..controls +(-30:.05) and +(110:0) ..(.18,1)
						--(.19,.95)
						..controls +(-140:.05) and +(30:.05) ..(-.1,.7)--cycle
						;
					}
					\fill[arsenic] \A3;
					\draw[black]\A3;
					\def\A2{ 
						(-.35,1.28)
						..controls +(-30:.05) and +(110:.05) ..(-.25,1.18)
						..controls +(-70:.05) and +(90:.02) ..(-.05,.9)
						..controls +(-90:.02) and +(50:.05) ..(-.05,.4)
						..controls +(-130:.02) and +(50:.05) ..(.6,-.8)
						..controls +(-130:.02) and +(50:.05) ..(.9,-1.8)
						..controls +(-130:.2) and +(30:.3) ..(-.5,-2.3)
						..controls +(130:.2) and +(-90:.1) ..(-.2,-1.5)
						..controls +(130:.2) and +(-140:.1) ..cycle
						;
					}
					\fill[arsenic] \A2;
					\draw[black]\A2;
					\def\A{ 
						(-.55,1.35)
						..controls +(-150:.2) and +(120:.1) ..(-.45,0)
						..controls +(-60:.1) and +(50:.1) ..(-.3,-1.6)
						..controls +(-10:.1) and +(-120:.1) ..(-.15,-1.5)
						..controls +(60:.1) and +(-120:.1) ..(0,-1.45)
						..controls +(120:.3) and +(-50:.1) ..(-.35,.77)
						..controls +(-50:.1) and +(130:.1) ..(-.15,.48)
						..controls +(-50:.1) and +(-130:.1) ..(.43,.93)
						..controls +(-130:.05) and +(-40:.05) ..(.3,1.05)
						..controls +(-130:.05) and +(40:.05) ..(-.05,.75)
						..controls +(-140:.05) and +(-40:.05) ..(-.3,1.15)
						..controls +(100:.05) and +(-40:.05) ..(-.35,1.28)--cycle
						;
					}
					\fill[eggplant] \A;
					\draw[black]\A;
					%-------------------------------
					%nón
					\def\N{ 
						(-.14,1.95)
						..controls +(-80:.1) and +(80:.1) ..(-.17,1.7)
						..controls +(110:.02) and +(-40:.02) ..(-.2,1.76)
						..controls +(-110:.1) and +(50:.4) ..(-.48,1.45)--(-.45,1.35)..controls +(-150:.05) and +(20:.05) ..(-.7,1.32)
						..controls +(50:.05) and +(-120:.05) ..(-.65,1.45)
						..controls +(160:.02) and +(-50:.02) ..(-.7,1.48)
						..controls +(160:.02) and +(-50:.02) ..(-.8,1.6)--(-.75,1.63)
						..controls +(-20:.06) and +(130:.04) ..(-.6,1.47)
						..controls +(130:.2) and +(-90:.1) ..(-.7,1.78)
						..controls +(90:.1) and +(170:.1) ..(-.4,1.9)
						..controls +(20:.1) and +(160:.1) ..cycle
						;
					}
					\fill[arsenic] \N;
					\draw[black]\N;
					\draw (-.14,1.95)
					..controls +(-100:.05) and +(40:.06) ..(-.2,1.76)
					(-.5,1.55)
					..controls +(30:.1) and +(-100:.06) ..(-.3,1.76)
					(-.6,1.7)
					..controls +(-55:.1) and +(120:.06) ..(-.55,1.5)
					;
					
					\def\T1{ 
						(.3,1.03)
						..controls +(60:.05) and +(-120:.05) ..(.35,1.15)
						..controls +(60:.05) and +(-140:.05) ..(.45,1.26)
						..controls +(30:.05) and +(60:.04) ..(.5,1.22)
						..controls +(30:.03) and +(50:.03) ..(.54,1.18)
						..controls +(30:.03) and +(50:.03) ..(.56,1.12)
						..controls +(30:.03) and +(50:.03) ..(.58,1.07)
						..controls +(160:.03) and +(40:.03) ..(.49,1)
						..controls +(-140:.03) and +(40:.03) ..(.41,.92)--cycle
						;
					}
					\fill[almond] \T1;
					\draw[black]\T1;
					
					\fill[crimson](-.27,.4)
					..controls +(-80:.6) and +(-80:.1) ..(-.5,-1.9)--(-.45,-1.85)
					..controls +(90:.1) and +(-80:1) ..(-.28,.48)--cycle;
					
					\fill[crimson](-.27,.4)
					..controls +(-80:.6) and +(-80:.1) ..(-.2,-1.9)--(-.15,-1.85)
					..controls +(90:.1) and +(-80:1) ..(-.28,.48)--cycle;
					
			}}
			\tikzset{no/.pic={
					\draw[crimson] (1.8,-1)
					..controls +(150:.4) and +(-120:.5) ..(1.7,-1)
					..controls +(30:.4) and +(-40:.3) ..(1.75,-1)
					..controls +(-30:.2) and +(40:.2) ..(1.7,-1.4)
					(1.75,-1)
					..controls +(-50:.1) and +(20:.1) ..(1.7,-1.2)
					
					;
			}}
			\draw[thick] (4,-1.25) arc(-60:-120:6cm);
			\draw[dashed,thick] (4,-1.25)--(1.3,4)--(-2,-1.25) (1.3,4)--(1.3,-2);
			\node at (1.6,.5) [rotate=90]{ Vị trí cân bằng};
			\draw[thick,<->] (.35,-2)--(1.3,-2);
			\node at (.85,-2) [above]{ $h$};
			\fill (4,-1.25) circle (0.05);
			\fill (-2,-1.25) circle (0.05);
			\fill (1.3,4) circle (0.05);
			\path
			(0,0)pic[scale=1,rotate=-20]{q_h}(-.88,1.48)pic[scale=1, scale=.7,rotate=-20]{no};
			
		\end{tikzpicture}
	}
	\choiceTF
	{\True Vật ở xa vị trí cân bằng nhất nghĩa là $h=1{,}5$ m}
	{Trong $10$ giây đầu tiên, có hai thời điểm vật ở xa vị trí cân bằng nhất}
	{\True Khi vật ở vị trí cân bằng thì $\cos \left(\dfrac{t\pi}{4} \right)=0$}
	{Trong khoảng từ $0$ đến $20$ giây thì vật đi qua vị trí cân bằng $4$ lần}
	\loigiai{
		\begin{itemchoice}
			\itemch \textbf{Đúng}.\\
			Ta có $h=|x|= \left|1{,}5 \cos \left(\dfrac{t\pi}{4}\right)\right|\leq 1{,}5$.\\
			Vật ở xa vị trí cân bằng nhất nghĩa là $h=1{,}5$ m.
			\itemch \textbf{Sai}.\\
			Khi đó, $\cos \left( \dfrac{t\pi}{4} \right) = \pm 1 \Leftrightarrow \hoac{& \dfrac{t\pi}{4} = k2\pi\\ & \dfrac{t\pi}{4} = \pi + k2\pi} \Leftrightarrow \hoac{& t=8k\\& t=4+8k}\quad (k \in \mathbb{Z})$.\\
			Vậy trong 10 giây đầu tiên thì vật ở xa vị trí cân bằng nhất tại các thời điểm $t=0$, $t=4$, $t=8$ (giây).
			\itemch \textbf{Đúng}.\\
			Khi vật ở vị trí cân bằng thì  
			\begin{eqnarray*}
			& & x = 0\\
			&\Leftrightarrow & 1{,}5 \cos \left(\dfrac{t\pi}{4}\right)=0\\
			&\Leftrightarrow & \cos \left( \dfrac{t\pi}{4} \right)=0\\
			&\Leftrightarrow & \dfrac{t\pi}{4} = \dfrac{\pi}{2}+k\pi \Rightarrow t=2+4k\quad (k \in \mathbb{Z}).
			\end{eqnarray*}			
			\itemch \textbf{Sai}.\\
			Trong khoảng từ $0$ đến $20$ giây thì vật ở vị trí cân bằng tại các thời điểm $t=2$, $t=6$, $t=10$, $t=14$, $t=18$ (giây); tức là có $5$ lần vật qua vị trí cân bằng.
		\end{itemchoice}
	}
\end{ex}

% #112 | Archive/SP_SauLuoi_2024-2025/SP_GHKI/1-TN-DS-TLN-chuyenLeQuyDon-NinhThuan-GHKI-NH24-25..tex | mã gốc 1D1H1-6
\begin{ex}%[1D1H1-6]%[1-TN-DS-TLN-THPT-ChuyenBinhThuan-BinhThuan-GHKI-NH24-25]%[Quan Ón]
	Một vòng quay Mặt Trời trong trò chơi cảm giác mạnh tại một khu du lịch, quay mỗi vòng mất $15$ phút. Tại vị trí quan sát, bạn Hoa thấy vòng quay chuyển động theo chiều kim đồng hồ. Khi vòng quay chuyển động được $10$ phút, bán kính của vòng quay quét một góc lượng giác có số đo bằng $-\dfrac{a\pi}{b}$ (rad) trong đó $\dfrac{a}{b}$ là phân số tối giản và $a$, $b > 0$). Tính $a + b$.
	\shortans{$7$}
	\loigiai{
		Do vòng quay Mặt Trời quay mỗi vòng khoảng $15$ phút và chuyển động theo chiều kim đồng hồ nên sau $15$ phút, bán kính của vòng quay quét một góc lượng giác có số đo bằng $-2\pi$ (rad).\\
		Do đó, sau $10$ phút, bán kính của vòng quay quét một góc lượng giác có số đo bằng 
		$$\dfrac{-2\pi}{15}\cdot 10 = -\dfrac{4\pi}{3} \textrm{ (rad).}$$
		Vậy $a = 4$ và $b = 3$, suy ra $a + b = 4 + 3 = 7$.
	}
\end{ex}

% #116 | Archive/SP_SauLuoi_2024-2025/SP_GHKI/1-TN-DS-TLN-chuyenLeQuyDon-NinhThuan-GHKI-NH24-25..tex | mã gốc 1D1V4-6
\begin{ex}%[1D1V4-6]%[Dự án đề kiểm tra Toán 11 GKI NH24-25-Trung Kiên]%[THPT Chuyên Bình Thuận - Bình Thuận]
	Hằng ngày mực nước tại một bến cảng lên xuống theo thủy triều. Độ sâu $h$ (m) của mực nước theo thời gian $t$ (giờ) trong một ngày được cho bởi công thức:
	$$h(t)=11+2\sin\left(\dfrac{\pi}{12}t\right),$$
	với $0\le t\le 24$. Mực nước thấp nhất tại cảng là bao nhiêu?
	\shortans{9}
	\loigiai{
		Với mọi $t$, ta có $-1\le \sin\left(\dfrac{\pi}{12}t\right)\le 1$. Suy ra 
		$$9\le 11+2\sin\left(\dfrac{\pi}{12}t\right) \le 13\Leftrightarrow 9\le h(t)\le 13.$$
		Vậy mực nước thấp nhất tại cảng là $9$ m.
	}
\end{ex}

% #117 | Archive/SP_SauLuoi_2024-2025/SP_GHKI/1-TN-DS-TLN-chuyenLeQuyDon-NinhThuan-GHKI-NH24-25..tex | mã gốc 1D1V3-2
\begin{ex}%[1D1V3-2]%[Dự án đề kiểm tra Toán 11 GKI NH24-25-Trung Kiên]%[THPT Chuyên Bình Thuận - Bình Thuận]
	Cho sơ đồ cấu trúc kim loại của một cây cầu ({\it tham khảo hình vẽ sau}).
	\begin{center}
		\begin{tikzpicture}[line join=round,line cap=round,>=stealth,font=\footnotesize,scale=0.6]
		\path
			(0,0) coordinate (O)
			(7.5,0) coordinate (A)
			(7.5,4) coordinate (Y)
			(1.66,-3.11) coordinate (X)
			(13.5,-2.5) coordinate (M)
			(13.5,1.5) coordinate (N)
		;
		\draw (O)--(Y)--(A)--(O)--(X)--(A)--(M)--(N)--(Y);
		\draw[ultra thick] (O)--(Y)--(N);
		\fill[gray!50] (O)--(Y)--(N)--(M)--(A)--(X)--(O);
		
		\foreach \x in {0.5,1,1.5} \draw (0,\x)--(7.5,\x+4)--(13.5,\x+1.5);
		\foreach \x/\i in {0/0,1/1,2/2,3/3,4/4,5/5}\draw (1.5*\x,0.8*\i)--(1.5*\x,1.5+0.8*\i);
		\foreach \x/\i in {6/3.375,7/2.75,8/2.125,9/1.5}\draw (1.5*\x,\i)--(1.5*\x,1.5+\i);
		
		\foreach \x/\y in {O/180,A/-90,X/-90}{
			\draw ($(\x)+(\y:4mm)$) node{$\x$};\fill[black] (\x) circle(1pt);}
		\draw 
			(O)--(A) node[sloped,pos=0.5,above]{$15$ m}
			(O)--(X) node[sloped,pos=0.5,below]{$9$ m}
			(A)--(Y) node[sloped,pos=0.5,below]{$8$ m}
			(Y) node[shift={(-60:4mm)}]{$Y$};
			;
		
		\draw pic[draw,angle radius=0.3cm]{right angle=A--X--O};
		\draw pic[draw,angle radius=0.3cm]{right angle=O--A--Y};
		\draw pic[draw,angle radius=0.3cm,angle eccentricity=1.5, "$\alpha$"]{angle=X--O--A};
		\draw pic[draw,angle radius=0.5cm,angle eccentricity=1.4, "$\beta$"]{angle=A--O--Y};
		\end{tikzpicture}
	\end{center}
	Biết $OX=9$ m, $OA=15$ m, $AY=8$ m, $\sin\widehat{XOY}=\dfrac{a}{b}$. Tính $b-a$.
	\shortans{$1$}
	\loigiai{
		Đặt $\alpha=\widehat{XOA}$ và $\beta=\widehat{AOY}$.\\
		Trong tam giác vuông $XOA$ vuông tại $X$, $AX=\sqrt{OA^2-OX^2}=12$ và
		$$\sin\alpha=\dfrac{AX}{OA}=\dfrac{12}{15}=\dfrac{4}{5},\quad \cos\alpha=\dfrac{OX}{OA}=\dfrac{9}{15}=\dfrac{3}{5}.$$
		Trong tam giác vuông $AOY$ vuông tại $A$, $OY=\sqrt{OA^2+AY^2}=17$ và
		$$\sin\beta=\dfrac{AY}{OY}=\dfrac{8}{17}, \quad \cos\beta=\dfrac{OA}{OY}=\dfrac{15}{17}.$$
		Ta có
		\begin{align*}
		\sin\widehat{XOY}
		&=\sin\left(\alpha+\beta\right)\\
		&=\sin\alpha\cos\beta+\sin\beta\cos\alpha\\
		&=\dfrac{4}{5}\cdot \dfrac{15}{17}+\dfrac{8}{17}\cdot \dfrac{3}{5}\\
		&=\dfrac{84}{85}.
		\end{align*}
		Vậy $b-a=85-84=1$.
	}
\end{ex}

% #118 | Archive/SP_SauLuoi_2024-2025/SP_GHKI/1-TN-DS-TL-THPT-TranPhu-TpHCM-GHKI-NH24-25.tex | mã gốc 1D1V5-3
\begin{ex}%[1D1V5-3]%[Dự án đề kiểm tra Toán khối 11 HK1 NH23-24-Đợt 3-Khắc Thiên]%[THPT Trần Phú - TP HCM]
	\choiceTF
	{\True Phương trình: $2\sin\left(x-\dfrac{\pi}{6}\right)=1$ có tập nghiệm là $\left\lbrace \dfrac{\pi}{3}+k2\pi; \pi+k2\pi| k\in\mathbb{Z}\right\rbrace$}
	{Số lượng hươu sau $t$ năm kể từ năm $2024$ (nghĩa là năm $2025$ là năm thứ nhất) được xác định bởi công thức: $N(t)=30+20\sin\left(\dfrac{\pi t}{10} \right)$ (nghìn con hươu). Năm $2030$ là năm đầu tiên số lượng hươu đạt $50$ nghìn con}
	{\True Nghiệm dương nhỏ nhất của phương trình: $\sqrt{3}\cot x-1=0$ là $\dfrac{\pi}{3}$}
	{\True Phương trình $\tan\left(2x-\dfrac{\pi}{6}\right)-\sqrt{3}=0$ có đúng sáu nghiệm thuộc $\left(0;3\pi\right)$}
	\loigiai{
		\begin{itemchoice}
			\itemch \textbf{Đúng}. Ta có 
			\begin{eqnarray*}
				&\Rightarrow&2\sin\left(x-\dfrac{\pi}{6}\right)=1\\
				&\Rightarrow&\sin\left(x-\dfrac{\pi}{6}\right)=\dfrac{1}{2}\\
				&\Rightarrow&\sin\left(x-\dfrac{\pi}{6}\right)=\sin\left(\dfrac{\pi}{6}\right)\\
				&\Rightarrow&\hoac{&x-\dfrac{\pi}{6}=\dfrac{\pi}{6}+k2\pi\\&x-\dfrac{\pi}{6}=\dfrac{5\pi}{6}+k2\pi}\left(k\in\mathbb{Z} \right) \\
				&\Rightarrow&\hoac{&x=\dfrac{\pi}{3}+k2\pi\\&x=\pi+k2\pi.}\left(k\in \mathbb{Z} \right) 
			\end{eqnarray*}
			Vậy tập nghiệm của phương trình là $\left\lbrace \dfrac{\pi}{3}+k2\pi; \pi+k2\pi| k\in\mathbb{Z}\right\rbrace$.
			\itemch \textbf{Sai}. Số lượng hươu năm $2030$ là $N(6)=30+20\sin\left(\dfrac{6\pi}{10}\right) \approx 49$ (nghìn con).
			\itemch \textbf{Đúng}. Ta có \begin{eqnarray*}
				&&\sqrt{3}\cot x-1=0\\
				&\Rightarrow&\cot x=\dfrac{1}{\sqrt{3}}\\
				&\Rightarrow&\tan x=\sqrt{3}\\
				&\Rightarrow& \tan x= \tan \left(\dfrac{\pi}{3} \right)\\
				&\Rightarrow& x=\dfrac{\pi}{3}+k\pi~(k\in \mathbb{Z})
			\end{eqnarray*}
			Chọn $k=0$ ta có $x=\dfrac{\pi}{3}$ là nghiệm dương nhỏ nhất của phương trình đã cho.
			\itemch \textbf{Đúng}. Ta có \begin{eqnarray*}
				&&\tan \left(2x-\dfrac{\pi}{6} \right) =\sqrt{3}\\
				&\Rightarrow& \tan \left(2x-\dfrac{\pi}{6} \right)= \tan \left(\dfrac{\pi}{3} \right)\\
				&\Rightarrow& 2x-\dfrac{\pi}{6} =\dfrac{\pi}{3}+k\pi\\
				&\Rightarrow& x=\dfrac{\pi}{4}+\dfrac{k\pi}{2}~(k\in\mathbb{Z}).
			\end{eqnarray*}
			Vì $x\in (0;3\pi)$ nên $0<\dfrac{\pi}{4}+\dfrac{k\pi}{2}<3\pi$ hay $-\dfrac{1}{2}<k<\dfrac{11}{2}$.\\
			Lại có $k\in\mathbb{Z}$ nên $k\in \left\lbrace 0;1;2;3;4;5 \right\rbrace$.\\
			Vậy phương trình đã cho có đúng sáu nghiệm thuộc $(0;3\pi)$.  
		\end{itemchoice}
	}
\end{ex}

% #122 | Archive/SP_SauLuoi_2024-2025/SP_GHKI/1-TN-DS-TL-THPT-TranPhu-TpHCM-GHKI-NH24-25.tex | mã gốc 1D1H6-8
\begin{ex}%[1D1H6-8]%[THPT Trần Phú - TP HCM - 2024 - đợt 3]%[TheHung Nguyen]
Để rèn luyện sức khỏe bạn Bình tập thể dục vào mỗi buổi sáng  và kết hợp chạy bộ trong một khoảng thời gian nhất định của mỗi buổi tập. Số phút chạy bộ của bạn Bình trong ngày thứ $t$ của một tháng được cho bởi hàm số sau: $p(t)=45+30 \cos \left(\dfrac{\pi t}{6}\right)$, trong đó $p(t)$ tính bằng phút $t\in \mathbb{Z}$ và $1 \leq t \leq 31$. Bằng cách giải phương trình lượng giác, tính số ngày bạn Bình chạy bộ được $60$ phút.
\loigiai{Ta có 
	\begin{eqnarray*}\allowdisplaybreaks
		& & 45+30 \cos \left(\dfrac{\pi t}{6}\right)=60\\
		&\Leftrightarrow & \cos \left(\dfrac{\pi t}{6}\right)=\dfrac{1}{2}\\
		&\Leftrightarrow & \cos \left(\dfrac{\pi t}{6}\right)=\cos \dfrac{\pi}{3}\\
		&\Leftrightarrow & \hoac{&\dfrac{\pi t}{6}=\dfrac{\pi}{3}+k2\pi\\& \dfrac{\pi t}{6}=-\dfrac{\pi}{3}+k2\pi}\\
		&\Leftrightarrow &\hoac{&t=2+12k\\&t=-2+12k.}
	\end{eqnarray*}
	Ta có  $1 \leq t \leq 31$, $k$, $t\in \mathbb{Z}$ và $t=2+12k$. Khi đó
	 $$1 \leq 2+12k \leq 31 \Leftrightarrow \leq -\dfrac{1}{12}\leq k \leq \dfrac{29}{12}\Rightarrow k\in \left\{0;1;2\right\}.$$ 
	Ta có  $1 \leq t \leq 31$, $k$, $t\in \mathbb{Z}$ và $t=-2+12k$. Khi đó	
		 $$1 \leq -2+12k \leq 31 \Leftrightarrow \leq \dfrac{3}{12}\leq k \leq \dfrac{33}{12}\Rightarrow k\in \left\{1;2\right\}.$$ 	
	Vậy có $5$ ngày bạn Bình chạy bộ được $60$ phút.
}
\end{ex}

% #125 | Archive/SP_SauLuoi_2024-2025/SP_GHKI/1-TL-TN-DS-NguyenThiMinhKhai-HCM-GHKI-NH24-25.tex | mã gốc 1D1V5-6
\begin{ex}%[1D1V5-6]%[Dự án đề kiểm tra Toán 11 GHKI NH24-25-Dot3-Hieu Hieu Minh Minh]%[THPT Nguyễn Thị Minh Khai - Tp HCM]
	Huyết áp là áp lực máu cần thiết tác động lên thành động mạch nhằm đưa máu đi nuôi dưỡng các mô trong cơ thể. Nhờ lực co bóp của tim và sức cản của động mạch mà huyết áp được tạo ra. Giả sử huyết áp của một người thay đổi theo thời gian được cho bởi công thức: $p(t)=120+15 \cos 150 \pi t$, trong đó $p(t)$ là huyết áp tính theo đơn vị mmHg (milimét thuỷ ngân) và thời gian $t$ tính theo đơn vị phút. Huyết áp cao nhất được gọi là huyết áp tâm thu và huyết áp thấp nhất được gọi là huyết áp tâm trương.
	\choiceTF[t]
	{\True $-1 \leq \cos 150 \pi t \leq 1, \forall t \geq 0$}
	{$p(2)=120 \mathrm{mmHg}$}
	{\True Huyết áp tâm thu của người đó là $135$ mmHg và huyết áp tâm trương của người đó là $105$ mmHg}
	{Trong thời gian từ $0$ giây đến $30$ giây, huyết áp của một người có $75$ lần đạt $120$ mmHg}
	\loigiai{
		\begin{itemchoice}
			\itemch \textbf{Đúng}. Vì $-1 \leq \cos 150 \pi t \leq 1, \forall t \geq 0$.
			\itemch \textbf{Sai}. \\
			Thay $t=2$ vào biểu thức $p(t)=120+15 \cos 150 \pi t$, \\
			ta được $p(2)=120+15 \cos (150 \cdot \pi \cdot 2)=135$ mmHg.
			\itemch \textbf{Đúng}. \\
			Ta có
			\begin{eqnarray*}
				&&-1 \leq \cos 150 \pi t \leq 1 \\&\Leftrightarrow& -15 \leq 15 \cos 150 \pi t \leq 15\\ &\Leftrightarrow& 105 \leq 120 + 15 \cos 150 \pi t \leq 135
			\end{eqnarray*}
			Vậy huyết áp tâm thu của người đó là $135$ mmHg, huyết áp tâm trương là $105$ mmHg.
			\itemch \textbf{Sai}. Vì huyết áp của một người đạt $120$ mmHg nên \begin{eqnarray*}
				&& 120 + 15 \cos 150 \pi t = 120 \\&\Leftrightarrow& \cos 150 \pi t =0 \\&\Leftrightarrow & 150 \pi t=\dfrac{\pi}{2}+k \pi\\
				&\Leftrightarrow& t=\dfrac{1}{300}+\dfrac{1}{150}k.
			\end{eqnarray*}
			Xét thời gian từ $0$ giây đến $30$ giây, ta có
			\begin{eqnarray*}
				&&0\leq t \leq 30 \\
				&\Leftrightarrow& 0 \leq \dfrac{1}{300}+\dfrac{1}{150}k \leq 30\\ &\Leftrightarrow& -\dfrac{1}{300} \leq \dfrac{1}{150}k \leq \dfrac{8999}{300}\\ &\Leftrightarrow& -\dfrac{1}{4} \leq k \leq \dfrac{8999}{2}.	
			\end{eqnarray*}
			Vì $k \in \mathbb{Z}$ nên $k \in \{0;1;2; \ldots ;4499\}$.\\
			Vậy có $4450$ lần huyết áp đạt $120$ mhHg.
		\end{itemchoice}
	}
\end{ex}

% #127 | Archive/SP_SauLuoi_2024-2025/SP_GHKI/1-TL-TN-DS-NguyenThiMinhKhai-HCM-GHKI-NH24-25.tex | mã gốc 1D1N5-4
\begin{bt}%[1D1N5-4]%[1D1H5-5]
	Giải các phương trình sau:
	\begin{enumerate}
		\item $\tan \left(2x-55^{\circ}\right)=1$.
		\item $1-2\sin^2x=\cos (x-12)$.
		\item $\sin x-\sqrt{3}\cos x=-2\cos \left(3x+\dfrac{\pi}{2}\right)$.
	\end{enumerate}
\loigiai{
	\begin{enumerate}
		\item $\tan \left(2x-55^{\circ}\right)=1\Leftrightarrow 2x-55^{\circ}=45^{\circ}+k\cdot 180^{\circ}\Leftrightarrow x=50^{\circ}+k\cdot 90^{\circ}$, $k\in\mathbb{Z}$.
		\item Ta có \allowdisplaybreaks
		\begin{eqnarray*}
		& &1-2\sin^2x=\cos (x-12)\\
		&\Leftrightarrow& \cos 2x=\cos (x-12)\\
		&\Leftrightarrow&\hoac{&2x=x-12+k\cdot 2\pi\\&2x=12-x+k\cdot 2\pi}\\
		&\Leftrightarrow&\hoac{&x=-12+k\cdot 2\pi\\&x=4+k\cdot\dfrac{2\pi}{3}}, k\in\mathbb{Z}.
		\end{eqnarray*}
		\item  Ta có \allowdisplaybreaks
		\begin{eqnarray*}
			& &\sin x-\sqrt{3}\cos x=-2\cos \left(3x+\dfrac{\pi}{2}\right)\\
			&\Leftrightarrow& \dfrac{1}{2}\cdot\sin x-\dfrac{\sqrt{3}}{2}\cdot\cos x=\sin 3x\\
			&\Leftrightarrow& \sin x\cdot\cos\dfrac{\pi}{3}-\sin\dfrac{\pi}{3}\cdot\cos x=\sin 3x\\
			&\Leftrightarrow& \sin\left(x-\dfrac{\pi}{3}\right)=\sin 3x\\ 
			&\Leftrightarrow&\hoac{&x-\dfrac{\pi}{3}=3x+k\cdot 2\pi\\&x-\dfrac{\pi}{3}=\pi-3x+k\cdot 2\pi}\\
			&\Leftrightarrow&\hoac{&x=-\dfrac{\pi}{6}-k\cdot \pi\\&x=\dfrac{\pi}{3}+k\cdot\dfrac{\pi}{2}}, k\in\mathbb{Z}.
		\end{eqnarray*}
	\end{enumerate}
}
\end{bt}

% #133 | Archive/SP_SauLuoi_2024-2025/SP_GHKI/1-TN-DS-TLN-THPT-NguyenAnNinh-HCM-GHKI-NH24-25.tex | mã gốc 1D1H4-8
\begin{ex}%[1D1H4-8]%%[Dự án đề kiểm tra Toán 11 GHKI NH24-25- Nguyễn Hữu Đức]%[THPT-NguyenAnNinh - Tp HCM]
Độ sâu $D(t)$m của nước ở một cảng biển sau $t$ giờ kể từ lúc nữa đêm được tính bởi công thức $D(t) = 4 \cos \left( \frac{\pi t}{6} \right) + 6 \, (m)$; $(0\le t\le 24)$. Độ sâu của nước ở cảng này sau $14$h là bao nhiêu mét?
\shortans{$8$}
\loigiai{
Dựa vào công thức đã cho
\[ D(t) = 4 \cos \left( \frac{\pi t}{6} \right) + 6 \, (m). \]
Tính độ sâu của nước tại thời điểm $t = 14$ giờ.\\
Thay $t = 14$ vào công thức
\[
D(14) = 4 \cos \left( \frac{\pi \cdot 14}{6} \right) + 6=8.
\]
Độ sâu của nước ở cảng này lúc $14$ giờ là $8$ mét.

}
\end{ex}

% #135 | Archive/SP_SauLuoi_2024-2025/SP_GHKI/1-TN-DS-TLN-THPT-NguyenAnNinh-HCM-GHKI-NH24-25.tex | mã gốc 1D1V2-5
\begin{ex}%[1D1V2-5]%[Dự án đề kiểm tra Toán 11 GHKI NH24-25- Nguyễn Hữu Đức]%[THPT-NguyenAnNinh - Tp HCM]
	Một xe đạp đang di chuyển, van $M$ của bánh xe quay quanh trục $O$ theo chiều kim đồng hồ với tốc độ góc không đổi $10\,\text{rad/s}$. Khoảng cách từ van $M$ đến mặt đất là đoạn $MH$. Ban đầu, van $M$ nằm ở vị trí $A$ (xem hình). Hỏi sau 2 phút khoảng cách từ van $M$ đến mặt đất là bao nhiêu cm? Biết bán kính bánh xe $R = 35\, \text{cm}$. Giả sử độ dày của lốp xe không đáng kể. Làm tròn kết quả đến hàng phần mười.
	
	\shortans{$2$}
	\loigiai{
		\begin{center}
			\begin{tikzpicture}
					% Vẽ mặt đất
					\draw[thick] (-4,-2) -- (4,-2) node[anchor=north] {Mặt đất};
					
					% Vẽ trục x và trục y
					\draw[->] (-3,0) -- (3,0) node[anchor=north west] {$x$};
					\draw[->] (0,-2) -- (0,3) node[anchor=east] {$y$};
					
					% Vẽ bánh xe
					\draw[thick] (0,0) circle (2);
					
					% Vẽ bán kính OM và đoạn thẳng AH
					\draw[thick] (1.3,1.5) -- (1.3,-2) node[below] {H};
					\draw[thick] (0,0) -- (2,0) node[anchor=north west] {A};
					\draw[thick] (0,0) -- (1.3,1.5) node[above right] {M};
					
					% Đặt các điểm O và các kí hiệu
					\node at (0,0) [below left] {O};
					\node at (0.4,0.182) {$\alpha$};
					
				\end{tikzpicture}
		\end{center}
		Ta có
		\begin{itemize}
			\item Tốc độ góc: $\omega = 10 \, \text{rad/s}$.
			\item Thời gian: $t = 2 \, \text{phút} = 120 \, \text{giây}$.
			\item Bán kính bánh xe: $R = 35 \, \text{cm}$.
		\end{itemize}
	Tính góc quay sau $2$ phút
	\[
	\theta = \omega \cdot t = 10 \cdot 120 = 1200 \, \text{rad}.
	\]
%	Xác định vị trí van $M$ sau $2$ phút.\\
	Số vòng quay $M$ sau $2$ phút.
	\[
	\text{Số vòng quay} = \frac{1200 \, \text{rad}}{2\pi \, \text{rad/vòng}} = \frac{1200}{2\pi} \approx 191 \, \text{vòng}.
	\]
	Sau mỗi vòng quay, van $M$ sẽ trở lại vị trí ban đầu, vì vậy sau $191$ vòng quay, van $M$ trở lại vị trí ban đầu.\\
	Khi $M$ ở vị trí $A$, khoảng cách từ van $M$ đến mặt đất là
	\[
	MH = R = 35 \, \text{cm}.
	\]
	}
\end{ex}

% #1064 | Archive/2024/TOAN_11_XE/TOAN_11_XE/data4/11C1/BTTN1.tex | mã gốc 1D1H1-3
\begin{ex}%[1D1H1-3]%[1D1V1-3]%[Trương Đăng Khoa - TeX Đúng Sai Toán 11]
	Xét tính đúng - sai của các mệnh đề sau.
	\choiceTF
	{\True Cho góc hình học $\widehat{u O v}=75^{\circ}$. Công thức số đo tổng quát của góc lượng giác $(O u, O v)$ là $75^{\circ}+k 360^{\circ}, k \in \mathbb{Z}$}
	{Cho một góc lượng giác $(O x, O u)$ có số đo $50^{\circ}$ và một góc lượng giác $(O x, O v)$ có số đo $-125^{\circ}$. Số đo của góc lượng giác $(O u, O v)$ bằng $175^{\circ}+k 360^{\circ}$, $k \in \mathbb{Z}$}
	{Trên đồng hồ hình dưới đây, kim phút đang chỉ đúng số $2$. Đến lúc $13$ giờ $30$ phút thì kim phút đã quay được một góc lượng giác có số đo bằng $480^{\circ}$\\
		\centerline{\begin{tikzpicture}[scale=1, font=\footnotesize, line join=round, line cap=round,>=stealth]
				\coordinate (O) at (0,0);
				\draw (O) circle (3 cm);
				%\fill[gray!50] (O) circle ( 1 cm);
				%Trục Oxy.
				\coordinate (3) at ($(O) + (0:2.5)$);
				\coordinate (2) at ($(O) + (30:2.5)$);
				\coordinate (1) at ($(O) + (60:2.5)$);
				\coordinate (12) at ($(O) + (90:2.5)$);
				\coordinate (11) at ($(O) + (120:2.5)$);
				\coordinate (10) at ($(O) + (150:2.5)$);
				\coordinate (9) at ($(O) + (180:2.5)$);
				\coordinate (8) at ($(O) + (210:2.5)$);
				\coordinate (7) at ($(O) + (240:2.5)$);
				\coordinate (6) at ($(O) + (270:2.5)$);
				\coordinate (5) at ($(O) + (300:2.5)$);
				\coordinate (4) at ($(O) + (330:2.5)$);
				\draw[line width=3pt, ->] (O)--(0.1,1.3);
				\draw[line width=3pt, ->]  (O)--(2);
				\foreach \i/\g in {3/0, 2/60,1/60,12/90,11/120,10/120,9/180,8/-120,7/-120,6/-90,5/-45,4/-45}{\draw[fill=black](\i) circle (2pt) ($(\i)+(\g:3mm)$) node[scale=1]{$\i$};}
				\draw (0,-3.7) node {$12$ giờ $30$ phút};
		\end{tikzpicture}}
	}
	{\True Trong hình dưới, cánh quạt được chia thành $3$ phần bằng nhau  $\widehat{A O D}=28^{\circ}$.  Công thức số đo tồng quát của góc lượng giác $(O x, O C)$ bằng $-148^{\circ}+k \cdot 360^{\circ}$\\
		\centerline{\begin{tikzpicture}[scale=1, font=\footnotesize, line join=round, line cap=round,>=stealth]
				\coordinate (O) at (0,0);
				\draw[dashed] (O) circle (2.5 cm);
				\fill[gray!50] (O) circle ( 1 cm);
				%Trục Oxy.
				\draw[->] (-3.5,0)--(3.5,0) node[below]{$x$};
				\draw[->] (0,-3.5)--(0,3.5) node[left]{$y$};
				\coordinate (A) at ($(O) + (0:2.5)$);
				\coordinate (D) at ($(O) + (-30:2.5)$);
				\coordinate (C) at ($(O) + (210:2.5)$);
				\coordinate (B) at ($(O) + (90:2.5)$);
				\draw[line width=3pt] (O)--(C) (O)--(B) (O)--(D);
				\foreach \i/\g in {A/45,B/45,C/-120}{\draw[fill=white](\i) circle (1.5pt) ($(\i)+(\g:3mm)$) node[scale=1]{$\i$};}
		\end{tikzpicture}}
	}
	\loigiai{
		\begin{itemchoice}
			\itemch Đúng.
			\itemch  Sai. $ \text{sđ}( {Ou},  {Ov})= \text{sđ}( {Ox},  {Ov})- {s} d( {Ox},  {Ou})+ {k} 360^{\circ}=-175^{\circ}+k 360^{\circ}$, $k \in \mathbb{Z}$.
			\itemch  Sai. Kim phút đã quay được một góc lượng giác có số đo bằng \break  $-\left(360^{\circ}+120^{\circ}\right)=-480^{\circ}$.
			\itemch  Đúng. Cánh quạt được chia thành $3$ phần bằng nhau, vậy mỗi phần sẽ có số đo là $\dfrac{360^{\circ}}{ 3}=120^{\circ}$.\\
			Số đo góc lượng giác $(O D, O C)=-120^{\circ}$.\\
			Theo hệ thức Chasles ta có
			$$
			(O x, O C)=(O x, O D)+(O D, O C)=-28^{\circ}+\left(-120^{\circ}\right)+k 360^{\circ}=-148^{\circ}+k 360^{\circ} .
			$$
		\end{itemchoice}
	}
\end{ex}

% #1100 | Archive/2024/TOAN_11_XE/TOAN_11_XE/data4/11C1/KNTT-1.4.tex | mã gốc 1D1K3-8
\begin{bt}%[1D1K3-8]
	Số giờ có ánh sáng mặt trời của một thành phố $A$ ở vĩ độ $40^\circ$ Bắc trong ngày thứ $t$ của một năm không nhuận được cho bởi hàm số
	\[
	d(t)=3\sin \left[\dfrac{\pi}{182}(t-80)\right] +12 \text{ với } t\in \mathbb{Z} \text{ và } 0<t\leq 365.
	\]
	\textit{(Nguồn: Đại số và Giải tích 11 Nâng cao, NXBGD Việt Nam, 2020}
	\begin{listEX}[1]
		\item Thành phố $A$ có đúng $12$ giờ có ánh sáng mặt trời vào ngày nào trong năm?
		\item Vào ngày nào trong năm thành phố $A$ có đúng $9$ giờ có ánh sáng mặt trời?
		\item Vào ngày nào trong năm thì thành phố $A$ có đúng $15$ giờ có ánh sáng mặt trời?
	\end{listEX}
	\loigiai{
		\begin{listEX}[1]
			\item Thành phố $A$ có đúng $12$ giờ có ánh sáng mặt trời vào ngày nào trong năm?\\
			Ta giải phương trình 
			\begin{eqnarray*}
				d(t)=12 &\Leftrightarrow& 3\sin\left[\dfrac{\pi}{182}(t-80)\right] +12=12 \\
				&\Leftrightarrow& \sin\left[\dfrac{\pi}{182}(t-80)\right]=0\\
				&\Leftrightarrow& \dfrac{\pi}{182}(t-80) =k \pi \\
				&\Leftrightarrow& t-80 = 182k\\
				&\Leftrightarrow& t= 182k+80 (k\in \mathbb{Z}).	
			\end{eqnarray*}
			Ta lại có  \[0<182k+80 \le 365 \Leftrightarrow -\dfrac{80}{182}<k \le \dfrac{285}{182} \Leftrightarrow \hoac{&k=0\\ &k=1.}\]
			Vậy thành phố $A$ có đúng $12$ giờ có ánh sáng mặt trời vào ngày thứ $80$ (ứng với $k=0$) và ngày thứ $262$ (ứng với $k=1$) trong năm.
			\item Vào ngày nào trong năm thành phố $A$ có đúng $9$ giờ có ánh sáng mặt trời?\\
			Ta giải phương trình 
			\begin{eqnarray*}
				d(t)=9 &\Leftrightarrow& 3\sin\left[\dfrac{\pi}{182}(t-80)\right] +12=9 \\
				&\Leftrightarrow& \sin\left[\dfrac{\pi}{182}(t-80)\right]=-1\\
				&\Leftrightarrow& \dfrac{\pi}{182}(t-80) =-\dfrac{\pi}{2}+k2 \pi \\
				&\Leftrightarrow& t-80 = -91+364k\\
				&\Leftrightarrow& t= 364k-11 (k\in \mathbb{Z}).	
			\end{eqnarray*}
			Ta lại có  \[0<364k-11 \le 365 \Leftrightarrow \dfrac{11}{364}<k \le \dfrac{376}{364} \Leftrightarrow k=1.\]
			Vậy thành phố $A$ có đúng $9$ giờ có ánh sáng mặt trời vào ngày thứ $353$ trong năm.
			\item Vào ngày nào trong năm thì thành phố $A$ có đúng $15$ giờ có ánh sáng mặt trời?
			Ta giải phương trình 
			\begin{eqnarray*}
				d(t)=15 &\Leftrightarrow& 3\sin\left[\dfrac{\pi}{182}(t-80)\right] +12=15 \\
				&\Leftrightarrow& \sin\left[\dfrac{\pi}{182}(t-80)\right]=1\\
				&\Leftrightarrow& \dfrac{\pi}{182}(t-80) =\dfrac{\pi}{2}+k2 \pi \\
				&\Leftrightarrow& t-80 = 91+364k\\
				&\Leftrightarrow& t= 364k+171 (k\in \mathbb{Z}).	
			\end{eqnarray*}
			Ta lại có  \[0<364k+171 \le 365 \Leftrightarrow -\dfrac{171}{364}<k \le \dfrac{196}{364} \Leftrightarrow k=0.\]
			Vậy thành phố $A$ có đúng $15$ giờ có ánh sáng mặt trời vào ngày thứ $171$ trong năm.
		\end{listEX}
	}
\end{bt}

% #1573 | Archive/2024/Update-SP-BG10-11-3/SP-3/BG11-3/data_1C5/Bai-2-Dang-4-Tinhxacsuatbangcongthuccong-Quan-On.tex | mã gốc 1D1H2-3
\begin{ex}[Mức độ 2]%[BG Toán 11 CT2025, Quan Ón]%[1D1H2-3]
	Một hộp có $6$ bi xanh, $4$ bi đỏ, $5$ bi vàng. Lấy ngẫu nhiên từ hộp $1$ viên bi. Tính xác suất lấy được một viên bi màu đỏ hoặc màu vàng (kết quả được ghi dưới dạng số thập phân).
	\shortans{$0{,}6$}
	\loigiai{
		Số phần tử của không gian mẫu $n(\Omega ) = 6 + 4 + 5 = 15$.\\
		Gọi $A$ là biến cố \lq\lq Lấy được một viên bi đỏ\rq\rq. Khi đó $n(A) = 4\Rightarrow \mathrm{P}(A) = \dfrac{4}{15}$.\\
		Gọi $B$ là biến cố \lq\lq Lấy được một viên bi vàng\rq\rq. Khi đó $n(B) = 5\Rightarrow \mathrm{P}(B) = \dfrac{5}{15} = \dfrac{1}{3}$.\\
		Biến cố $A\cup B$ là \lq\lq Lấy được một viên bi màu đỏ hoặc màu vàng\rq\rq.\\
		Vì $A$ và $B$ là các biến cố xung khắc nên
		$$\mathrm{P}(A\cup B) = \mathrm{P}(A) + \mathrm{P}(B) = \dfrac{4}{15} + \dfrac{1}{3} = \dfrac{3}{5} = 0{,}6.$$
	}
\end{ex}

% #1574 | Archive/2024/Update-SP-BG10-11-3/SP-3/BG11-3/data_1C5/Bai-2-Dang-4-Tinhxacsuatbangcongthuccong-Quan-On.tex | mã gốc 1D1H2-3
\begin{ex}[Mức độ 2]%[BG Toán 11 CT2025, Quan Ón]%[1D1H2-3]
	Một lớp có $35$ học sinh trong đó có $26$ bạn thích học Toán, $14$ người thích học cả Toán và Văn, $30$ người thích ít nhất một trong hai môn Toán hoặc Văn. Chọn ngẫu nhiên một bạn trong lớp. Xác suất để học sinh đó thích học môn Văn bằng bao nhiêu? (kết quả làm tròn đến hàng phần trăm).
	\shortans{$0{,}51$}
	\loigiai{
		Xét các biến cố $A\colon$\lq\lq Học sinh đó thích học Toán\rq\rq; $B\colon$\lq\lq Học sinh đó thích học Văn\rq\rq.\\
		Theo đề bài ta có $\mathrm{P}(A) = \dfrac{26}{35}$; $\mathrm{P}(AB) = \dfrac{14}{35}$; $\mathrm{P}(A\cup B) = \dfrac{30}{35}$.\\
		Áp dụng công thức cộng xác suất, ta có $\mathrm{P}(A\cup B) = \mathrm{P}(A) + \mathrm{P}(B) - \mathrm{P}(AB)$.\\
		Suy ra $\mathrm{P}(B) = \mathrm{P}(A\cup B) + \mathrm{P}(AB) - \mathrm{P}(A) = \dfrac{30}{35} + \dfrac{14}{35} - \dfrac{26}{35} = \dfrac{18}{35} \approx 0{,}51$.
	}
\end{ex}

% #1575 | Archive/2024/Update-SP-BG10-11-3/SP-3/BG11-3/data_1C5/Bai-2-Dang-4-Tinhxacsuatbangcongthuccong-Quan-On.tex | mã gốc 1D1H2-3
\begin{ex}[Mức độ 2]%[BG Toán 11 CT2025, Quan Ón]%[1D1H2-3]
	Một câu lạc bộ ngoại ngữ có $40$ học sinh, trong đó có $25$ học sinh giỏi tiếng Anh, $20$ học sinh giỏi tiếng Pháp và $15$ học sinh giỏi cả hai ngoại ngữ tiếng Anh và tiếng Pháp. Học sinh nào giỏi ít nhất một trong hai thứ tiếng sẽ được cử đi phiên dịch cho một Hội nghị quốc tế. Chọn ngẫu nhiên một trong những học sinh của câu lạc bộ. Tính xác suất để học sinh đó có thể tham gia Hội nghị (kết quả được ghi dưới dạng số thập phân).
	\shortans{$0{,}75$}
	\loigiai{ 
		Gọi $X$ là biến cố: \lq \lq Học sinh chọn ra có thể tham gia Hội nghị\rq\rq.\\
		$A$ là biến cố: \lq \lq Học sinh chọn ra giỏi tiếng Anh\rq\rq.\\
		$B$ là biến cố: \lq \lq Học sinh chọn ra giỏi tiếng Pháp\rq\rq.\\
		Khi đó $X = A\cup B$. \\
		Theo công thức cộng xác suất, ta có $\mathrm{P}(A\cup B) = \mathrm{P}(A) + \mathrm{P}(B) - \mathrm{P}(A\cap B) = \dfrac{25}{40} + \dfrac{20}{40} - \dfrac{15}{40} =\dfrac{3}{4} = 0{,}75$.
		}
\end{ex}

% #1576 | Archive/2024/Update-SP-BG10-11-3/SP-3/BG11-3/data_1C5/Bai-2-Dang-4-Tinhxacsuatbangcongthuccong-Quan-On.tex | mã gốc 1D1V2-3
\begin{ex}[Mức độ 3]%[BG Toán 11 CT2025, Quan Ón]%[1D1V2-3]
	Hộp $I$ chứa $4$ viên bi trắng, $5$ viên bi đỏ và $6$ viên bi xanh. Hộp $II$ có $7$ viên bi trắng, $6$ viên bi đỏ và $5$ viên bi xanh. Lấy ngẫu nhiên từ mỗi hộp $1$ viên bi. Xác suất để hai viên bi được lấy ra có cùng màu là $\dfrac{a}{b}$ với $\dfrac{a}{b}$ là phân số tối giản và $a, b\in\mathbb{Z}$. Tính $T = a + b$.
	\shortans{$179$}
	\loigiai{
		Gọi $A$ là biến cố: \lq\lq Hai viên bi được lấy ra có cùng màu\rq\rq.\\
		Xét các biến cố sau
		\begin{itemize}
			\item $A_1\colon$ \lq\lq Hai viên bi lấy ra màu trắng\rq\rq. Ta được $\mathrm{P}(A_1)=\dfrac{4}{15}\cdot \dfrac{7}{18}$.
			\item  $A_2\colon$ \lq\lq Hai viên bi lấy ra màu đỏ\rq\rq. Ta được $\mathrm{P}(A_2)=\dfrac{5}{15}\cdot \dfrac{6}{18}$.
			\item $A_3\colon$ \lq\lq Hai viên bi lấy ra màu xanh\rq\rq. Ta được $\mathrm{P}(A_3)=\dfrac{6}{15}\cdot \dfrac{5}{18}$.
		\end{itemize}
		Lúc đó $A=A_1\cup A_2\cup A_3$ và $A_1$, $A_2$, $A_3$ là các biến cố xung khắc nên
		\begin{eqnarray*}
			\mathrm{P}(A) = \mathrm{P}(A_1) + \mathrm{P}(A_2) + \mathrm{P}(A_3) = \dfrac{44}{135}.
		\end{eqnarray*}
		Do đó $\heva{&a = 44\\&b = 135.}$\\
		Vậy $T = a + b = 44 + 135 = 179$.
	}
\end{ex}

% #1577 | Archive/2024/Update-SP-BG10-11-3/SP-3/BG11-3/data_1C5/Bai-2-Dang-4-Tinhxacsuatbangcongthuccong-Quan-On.tex | mã gốc 1D1C2-3
\begin{ex}[Mức độ 4]%[BG Toán 11 CT2025, Quan Ón]%[1D1C2-3]
	Cho đa giác đều $(H)$ có $20$ đỉnh. Lấy tùy ý $3$ đỉnh của $(H)$. Gọi $A$ là biến cố \lq\lq $3$ đỉnh lấy được tạo thành một tam giác vuông sao cho không có cạnh nào là cạnh của $(H)$\rq\rq; $B$ là biến cố \lq\lq $3$ đỉnh lấy được tạo thành một tam giác có đúng một cạnh là cạnh của đa giác $(H)$\rq\rq. Khi đó giá trị của $\mathrm{P}(A\cup B)$ là $\dfrac{a}{b}$ với $\dfrac{a}{b}$ là phân số tối giản và. Tính $T=2a+3b$.
	\shortans{$103$}
	\loigiai{
		Ta có $A$ và $B$ là hai biến cố xung khắc nên $\mathrm{P}(A\cup B)=\mathrm{P}(A)+\mathrm{P}(B)$.\\
		Số phần tử của không gian mẫu $n(\Omega)=\mathrm{C}_{20}^3$.\\
		Gọi $A$ là biến cố \lq\lq $3$ đỉnh lấy được tạo thành một tam giác vuông sao cho không có cạnh nào là cạnh của $(H)$\rq\rq.\\
		Để ba đỉnh được chọn tạo thành tam giác vuông sao cho không có cạnh nào là cạnh của $(H)$ thì hai trong ba đỉnh được chọn là hai đầu mút của một đường kính của đường tròn ngoại tiếp đa giác đều $(H)$ và đỉnh còn lại là một trong số các đỉnh còn lại của đa giác không kề với hai đỉnh đã chọn.\\
		Đa giác có $20$ đỉnh nên có $10$ đường kính.\\
		Số đỉnh còn lại của đa giác mà không kề với hai đỉnh đã chọn là $20-2-4=14$.\\
		Suy ra $n(A)=\mathrm{C}_{10}^1\cdot \mathrm{C}_{14}^1\Rightarrow \mathrm{P}(A) = \dfrac{n(A)}{n(\Omega)}=\dfrac{\mathrm{C}_{10}^1\cdot \mathrm{C}_{14}^1}{\mathrm{C}_{20}^3}=\dfrac{7}{57}$.\\
		Gọi $B$ là biến cố: \lq\lq $3$ đỉnh lấy được tạo thành một tam giác có đúng một cạnh là cạnh của đa giác $(H)$\rq\rq. Giả sử đa giác đã cho có các cạnh $A_1A_2A_3\ldots A_{20}$.\\
		Với cạnh $A_1A_2$ ta có các tam giác là: $A_1A_2A_4$, $A_1A_2A_5 ,\ldots,A_1A_2A_{19}$ có $16$ tam giác có một cạnh là cạnh của đa giác.\\
		Tương tự với $19$ cạnh còn lại, mỗi cạnh đều lập được $16$ tam giác có một cạnh là cạnh của đa giác.\\
		Suy ra $n(B)=16\cdot 20=320 \Rightarrow \mathrm{P}(B) = \dfrac{n(B)}{n(\Omega)} = \dfrac{320}{C_{20}^3}=\dfrac{16}{57}$.\\
		Do đó $\mathrm{P}(A\cup B) = \mathrm{P}(A) + \mathrm{P}(B) = \dfrac{23}{57}\Rightarrow\heva{&a=23\\&b=57}\Rightarrow T = 103$.
	}
\end{ex}

% #1627 | Archive/2024/SP-K10-K11-LyVanHoang_TVN-023/11. C1-LuongGiac-lvhoang_TVN-023/dataC1K11/Chuong1-Bai4-Dang3-PhuongTrinhCBDungRadian.tex | mã gốc 1D1N5-3
\begin{vd}[Mức N, CTST]%[BG Toán 11 CT2025, Quan Ón]%[1D1N5-3]
	Giải các phương trình
	\begin{listEX}[5]
		\item $\sin x=\dfrac{1}{2}$.
		\item $\cos x=-\dfrac{1}{2}$.
		\item $\tan x=\sqrt{3}$.
		\item $\cot x=-\dfrac{\sqrt{3}}{3}$.
		\item $\sin x = -\dfrac{5}{2}$.
	\end{listEX}
	\loigiai{
		\begin{enumerate}
			\item $\sin x = \dfrac{1}{2} \Leftrightarrow \sin x = \sin\dfrac{\pi}{6} \Leftrightarrow \hoac{&x = \dfrac{\pi}{6} + k2\pi\\&x = \pi - \dfrac{\pi}{6} + k2\pi} \Leftrightarrow \hoac{&x = \dfrac{\pi}{6} + k2\pi\\&x = \dfrac{5\pi}{6} + k2\pi}$ $(k\in\mathbb{Z})$.
			\item $\cos x = -\dfrac{1}{2} \Leftrightarrow \cos x = \cos\dfrac{2\pi}{3} \Leftrightarrow \hoac{&x = \dfrac{2\pi}{3} + k2\pi\\&x = -\dfrac{2\pi}{3} + k2\pi}$ $(k\in\mathbb{Z})$.
			\item $\tan x = \sqrt{3} \Leftrightarrow \tan x = \tan\dfrac{\pi}{3} \Leftrightarrow x = \dfrac{\pi}{3} + k\pi$, $(k\in\mathbb{Z})$.
			\item $\cot x = -\dfrac{\sqrt{3}}{3} \Leftrightarrow \cot x = \cot\left(-\dfrac{\pi}{3}\right) \Leftrightarrow x = -\dfrac{\pi}{3} + k\pi$, $(k\in\mathbb{Z})$.
			\item Vì $\left|-\dfrac{5}{2} \right| = \dfrac{5}{2} > 1$, do đó phương trình $\sin x = -\dfrac{5}{2}$ vô nghiệm.
		\end{enumerate}
	}
\end{vd}

% #1629 | Archive/2024/SP-K10-K11-LyVanHoang_TVN-023/11. C1-LuongGiac-lvhoang_TVN-023/dataC1K11/Chuong1-Bai4-Dang3-PhuongTrinhCBDungRadian.tex | mã gốc 1D1H5-3
\begin{bt}[Mức H, CTST]%[BG Toán 11 CT2025, Quan Ón]%[1D1H5-3] 
	Giải các phương trình
	\begin{listEX}[3]
		\item $\sin x=-\dfrac{3}{2}$;
		\item $\sin x=\dfrac{\sqrt{3}}{2}$;
		\item $\cos x=-3$;
		\item $\cos\left(x+\dfrac{\pi}{12}\right)=\cos\dfrac{3\pi}{12}$;
		\item $\tan2x=\tan\dfrac{\pi}{11}$;
		\item $\tan x=0$;
		\item $\cot3x=\cot\dfrac{\pi}{7}$;
		\item $\cot x=1$;
		\item $\sin2x=\sin3x$;
		\item $\cos3x=\sin x$.
	\end{listEX}
	\loigiai{
		\begin{enumerate}
			\item $\sin x=-\dfrac{3}{2}$ vô nghiệm vì $\left|-\dfrac{3}{2} \right| = \dfrac{3}{2} > 1$.
			\item $\sin x=\dfrac{\sqrt{3}}{2}\Leftrightarrow\sin x=\sin\dfrac{\pi}{3}\Leftrightarrow\hoac{&x=\dfrac{\pi}{3}+k2\pi\\&x=\dfrac{2\pi}{3}+k2\pi}\,(k\in\mathbb{Z})$.
			\item $\cos x=-3$ vô nghiệm vì $|-3| = 3 > 1$.
			\item $\cos\left(x+\dfrac{\pi}{12}\right)=\cos\dfrac{3\pi}{12}\Leftrightarrow\hoac{&x+\dfrac{\pi}{12}=\dfrac{3\pi}{12}+k2\pi\\&x+\dfrac{\pi}{12}=-\dfrac{3\pi}{12}+k2\pi}\Leftrightarrow\hoac{&x=\dfrac{\pi}{6}+k2\pi\\&x=-\dfrac{\pi}{3}+k2\pi}\,(k\in\mathbb{Z})$.
			\item $\tan2x=\tan\dfrac{\pi}{11}\Leftrightarrow2x=\dfrac{\pi}{11}+k\pi\Leftrightarrow x=\dfrac{\pi}{22}+\dfrac{k\pi}{2},\,(k\in\mathbb{Z})$.
			\item $\tan x=0 \Leftrightarrow x=k\pi,\,(k\in\mathbb{Z})$.
			\item $\cot3x=\cot\dfrac{\pi}{7}\Leftrightarrow 3x=\dfrac{\pi}{7}+k\pi\Leftrightarrow x=\dfrac{\pi}{21}+\dfrac{k\pi}{3},\,(k\in\mathbb{Z})$.
			\item $\cot x=1\Leftrightarrow\cot x=\cot\dfrac{\pi}{4}\Leftrightarrow x=\dfrac{\pi}{4}+k\pi,\,(k\in\mathbb{Z})$.
			\item $\sin2x=\sin3x\Leftrightarrow\hoac{&2x=3x+k2\pi\\&2x=\pi-3x+k2\pi}\Leftrightarrow\hoac{&x=-k2\pi\\&x=\dfrac{\pi}{5}+\dfrac{k2\pi}{5}}\,(k\in\mathbb{Z})$.
			\item $\cos3x=\sin x\Leftrightarrow \cos3x=\cos\left(\dfrac{\pi}{2}-x\right)\Leftrightarrow\hoac{&3x=\dfrac{\pi}{2}-x+k2\pi\\&3x=-\dfrac{\pi}{2}+x+k2\pi}\Leftrightarrow\hoac{&x=\dfrac{\pi}{8}+\dfrac{k\pi}{2}\\&x=-\dfrac{\pi}{4}+k\pi}\,(k\in\mathbb{Z})$.
	\end{enumerate}}
\end{bt}

% #1630 | Archive/2024/SP-K10-K11-LyVanHoang_TVN-023/11. C1-LuongGiac-lvhoang_TVN-023/dataC1K11/Chuong1-Bai4-Dang3-PhuongTrinhCBDungRadian.tex | mã gốc 1D1H5-3
\begin{bt}[Mức H, Cánh Diều]%[BG Toán 11 CT2025, Quan Ón]%[1D1H5-3]
	Giải các phương trình
	\begin{multicols}{3}
		\begin{enumerate}
			\item $\sin x=-\dfrac{1}{2}$;
			\item $\sin x=\dfrac{\sqrt{2}}{2}$;
			\item $\cos x=\dfrac{\sqrt{3}}{2}$;
			\item $\cos x=-\dfrac{\sqrt{2}}{2}$;
			\item $\cos x=-\dfrac{1}{2}$;
			\item $\cos3x=\cos\left(x+\dfrac{\pi}{3}\right)$;
			\item $\tan x=\dfrac{1}{\sqrt{3}}$;
			\item $\tan x=-1$;
			\item $\cot2x=-\sqrt{3}$.
		\end{enumerate}
	\end{multicols}
	\loigiai{
		\begin{enumerate}
			\item Do $\sin\left(-\dfrac{\pi}{6}\right)=-\dfrac{1}{2}$ nên $$\sin x=\sin\left(-\dfrac{\pi}{6}\right)\Leftrightarrow\hoac{&x=-\dfrac{\pi}{6}+k2\pi\\&x=\pi-\left(-\dfrac{\pi}{6}\right)+k2\pi}\Leftrightarrow\hoac{&x=-\dfrac{\pi}{6}+k2\pi\\&x=\dfrac{7\pi}{6}+k2\pi}\,(k\in\mathbb{Z}).$$
			\item Do $\sin\dfrac{\pi}{4}=\dfrac{\sqrt{2}}{2}$ nên $$\sin x=\sin\dfrac{\pi}{4}\Leftrightarrow\hoac{&x=\dfrac{\pi}{4}+k2\pi\\&x=\pi-\dfrac{\pi}{4}+k2\pi}\Leftrightarrow\hoac{&x=\dfrac{\pi}{4}+k2\pi\\&x=\dfrac{3\pi}{4}+k2\pi}\,(k\in\mathbb{Z}).$$
			\item $\cos x=\dfrac{\sqrt{3}}{2}\Leftrightarrow\cos x=\cos\dfrac{\pi}{6}\Leftrightarrow x=\pm\dfrac{\pi}{6}+k2\pi,\,(k\in\mathbb{Z})$.
			\item $\cos x=-\dfrac{\sqrt{2}}{2}\Leftrightarrow \cos x=\cos\dfrac{3\pi}{4}\Leftrightarrow x=\pm\dfrac{3\pi}{4}+k2\pi,\,(k\in\mathbb{Z})$.
			\item $\cos x=-\dfrac{1}{2}\Leftrightarrow\cos x=\cos\dfrac{2\pi}{3}\Leftrightarrow x=\pm\dfrac{2\pi}{3}+k2\pi,\,(k\in\mathbb{Z})$.
			\item $\cos3x=\cos\left(x+\dfrac{\pi}{3}\right)\Leftrightarrow\hoac{&3x=x+\dfrac{\pi}{3}+k2\pi\\&3x=-x-\dfrac{\pi}{3}+k2\pi}\Leftrightarrow\hoac{&x=\dfrac{\pi}{6}+k\pi\\&x=-\dfrac{\pi}{12}+\dfrac{k\pi}{2}},\,(k\in\mathbb{Z})$.
			\item $\tan x=\dfrac{1}{\sqrt{3}}\Leftrightarrow\tan x=\tan\dfrac{\pi}{6}\Leftrightarrow x=\dfrac{\pi}{6}+k\pi,\,(k\in\mathbb{Z})$.
			\item $\tan x=-1\Leftrightarrow\tan x=\tan\left(-\dfrac{\pi}{4}\right)\Leftrightarrow x=-\dfrac{\pi}{4}+k\pi,\,(k\in\mathbb{Z})$.
			\item $\cot2x=-\sqrt{3}\Leftrightarrow\cot2x=\cot\left(-\dfrac{\pi}{6}\right)\Leftrightarrow x=-\dfrac{\pi}{6}+k\pi,\,(k\in\mathbb{Z})$.
	\end{enumerate}}
\end{bt}

% #1631 | Archive/2024/SP-K10-K11-LyVanHoang_TVN-023/11. C1-LuongGiac-lvhoang_TVN-023/dataC1K11/Chuong1-Bai4-Dang3-PhuongTrinhCBDungRadian.tex | mã gốc 1D1H5-3
\begin{bt}[Mức H, KNTT]%[BG Toán 11 CT2025, Quan Ón]%[1D1H5-3] 
	Giải các phương trình
	\begin{multicols}{3}
		\begin{enumerate}
			\item $\sin x=-\dfrac{\sqrt{3}}{2}$;
			\item $\sin x=\dfrac{1}{3}$;
			\item $\sin x=\dfrac{\sqrt{2}}{2}$;
			\item $\sin3x=-\sin5x$;
			\item $\cos x=-\dfrac{\sqrt{3}}{2}$;
			\item $\cos x=0{,}1$;
			\item $2\cos x=-\sqrt{2}$;
			\item $\cos3x-\sin5x=0$;
			\item $\tan x=-\sqrt{3}$;
			\item $\tan x=2$;
			\item $\sqrt{3}\tan2x=-1$;
			\item $\tan3x+\tan5x=0$;
			\item $\cot x=-\sqrt{3}$;
			\item $\cot x=5$;
			\item $\cot x=1$;
			\item $\sqrt{3}\cot x+1=0$.
		\end{enumerate}
	\end{multicols}
	\loigiai{
		\begin{enumerate}
			\item $\sin x=-\dfrac{\sqrt{3}}{2}\Leftrightarrow \sin x=\sin\left(-\dfrac{\pi}{3}\right)\Leftrightarrow \hoac{&x=-\dfrac{\pi}{3}+k2\pi\\&x=\pi-\left(-\dfrac{\pi}{3}\right)+k2\pi}\Leftrightarrow\hoac{&x=-\dfrac{\pi}{3}+k2\pi\\&x=\dfrac{4\pi}{3}+k2\pi}\,(k\in\mathbb{Z})$.
			\item Gọi $\alpha\in\left[-\dfrac{\pi}{2};\dfrac{\pi}{2}\right]$ là góc thỏa mãn $\sin\alpha=\dfrac{1}{3}$.\\Khi đó ta có $\sin x=\dfrac{1}{3}\Leftrightarrow\sin x=\sin\alpha\Leftrightarrow \hoac{&x=\alpha+k2\pi\\&x=\pi-\alpha+k2\pi}\,(k\in\mathbb{Z})$.
			\item $\sin x=\dfrac{\sqrt{2}}{2}\Leftrightarrow\sin x=\sin\dfrac{\pi}{4}\Leftrightarrow\hoac{&x=\dfrac{\pi}{4}+k2\pi\\&x=\pi-\dfrac{\pi}{4}+k2\pi}\Leftrightarrow\hoac{&x=\dfrac{\pi}{4}+k2\pi\\&x=\dfrac{5\pi}{4}+k2\pi}\,(k\in\mathbb{Z})$.
			\item $\sin3x=-\sin5x\Leftrightarrow\sin3x=\sin(-5x)\Leftrightarrow\hoac{&3x=-5x+k2\pi\\&3x=\pi+5x+k2\pi}\Leftrightarrow\hoac{&x=\dfrac{k\pi}{4}\\&x=-\dfrac{\pi}{2}-k\pi}\,(k\in\mathbb{Z})$.
			\item $\cos x=-\dfrac{\sqrt{3}}{2}\Leftrightarrow\cos x=\cos\dfrac{5\pi}{6}\Leftrightarrow x=\pm\dfrac{5\pi}{6}+k2\pi,\,(k\in\mathbb{Z})$.
			\item Gọi $\alpha\in\left[0;\pi\right]$ là góc thỏa mãn $\cos\alpha=0{,}1$. Khi đó ta có $\cos x=\cos\alpha\Leftrightarrow x=\pm\alpha+k2\pi,\,(k\in\mathbb{Z})$.
			\item $2\cos x=-\sqrt{2}\Leftrightarrow\cos x=-\dfrac{\sqrt{2}}{2}\Leftrightarrow\cos x=\cos\dfrac{3\pi}{4}\Leftrightarrow x=\pm\dfrac{3\pi}{4}+k2\pi,\,(k\in\mathbb{Z})$.
			\item $\cos3x-\sin5x=0\Leftrightarrow\cos3x=\sin5x\Leftrightarrow\cos3x=\cos\left(\dfrac{\pi}{2}-5x\right)\Leftrightarrow\hoac{&3x=\dfrac{\pi}{2}-5x+k2\pi\\&3x=-\dfrac{\pi}{2}+5x+k2\pi}\\\Leftrightarrow\hoac{&x=\dfrac{\pi}{16}+\dfrac{k\pi}{4}\\&x=\dfrac{\pi}{4}-k\pi}\,(k\in\mathbb{Z})$.
			\item $\tan x=-\sqrt{3}\Leftrightarrow\tan x=\tan\left(-\dfrac{\pi}{3}\right)\Leftrightarrow x=-\dfrac{\pi}{3}+k\pi,\,(k\in\mathbb{Z})$.
			\item Gọi $\alpha\in\left(-\dfrac{\pi}{2};\dfrac{\pi}{2}\right)$ sao cho $\tan\alpha=2$.\\Khi đó $\tan x=2\Leftrightarrow\tan x=\tan\alpha\Leftrightarrow x=\alpha+k\pi,\,(k\in\mathbb{Z})$.
			\item $\sqrt{3}\tan2x=-1\Leftrightarrow\tan2x=-\dfrac{\sqrt{3}}{3}\Leftrightarrow\tan2x=\tan\left(-\dfrac{\pi}{6}\right)\Leftrightarrow2x=-\dfrac{\pi}{6}+k\pi\\\Leftrightarrow x=-\dfrac{\pi}{12}+\dfrac{k\pi}{2},\,(k\in\mathbb{Z})$.
			\item Điều kiện $\heva{&3x\ne\dfrac{\pi}{2}+k\pi\\&5x\ne\dfrac{\pi}{2}+k\pi}\Leftrightarrow\heva{&x\ne\dfrac{\pi}{6}+\dfrac{k\pi}{3}\\&x\ne\dfrac{\pi}{5}+\dfrac{k\pi}{5}}\,(k\in\mathbb{Z})$.\\ $\tan3x+\tan5x=0\Leftrightarrow\tan3x=\tan(-5x)\Leftrightarrow3x=-5x+k\pi\Leftrightarrow x=\dfrac{k\pi}{8},\,(k\in\mathbb{Z})$.
			\item $\cot x=-\sqrt{3}\Leftrightarrow\cot x=\cot\left(-\dfrac{\pi}{6}\right)\Leftrightarrow x=-\dfrac{\pi}{6}+k\pi,\,(k\in\mathbb{Z})$.
			\item Gọi $\alpha\in\left(0;\pi\right)$ sao cho $\cot\alpha=5$.\\Khi đó $\cot x=5\Leftrightarrow\cot x=\cot\alpha\Leftrightarrow x=\alpha+k\pi,\,(k\in\mathbb{Z})$.
			\item $\cot x=1\Leftrightarrow\cot x=\cot\dfrac{\pi}{4}\Leftrightarrow x=\dfrac{\pi}{4}+k\pi,\,(k\in\mathbb{Z})$.
			\item $\sqrt{3}\cot x+1=0\Leftrightarrow\cot x=\dfrac{-1}{\sqrt{3}}\Leftrightarrow\cot x=\cot\left(-\dfrac{\pi}{3}\right)\Leftrightarrow x=-\dfrac{\pi}{3}+k\pi,\,(k\in\mathbb{Z})$.
	\end{enumerate}}
\end{bt}

% #1634 | Archive/2024/SP-K10-K11-LyVanHoang_TVN-023/11. C1-LuongGiac-lvhoang_TVN-023/dataC1K11/Chuong1-Bai4-Dang3-PhuongTrinhCBDungRadian.tex | mã gốc 1D1H5-3
\begin{bt}[Mức H, KNTT]%[BG Toán 11 CT2025, Quan Ón]%[1D1H5-3]
	Giải các phương trình sau
	\begin{listEX}[3]
		\item $\sin x=\dfrac{\sqrt{3}}{2}$;
		\item $2\cos x=-\sqrt{2}$;
		\item $\cot(2x-1)=\cot\dfrac{\pi}{5}$.
	\end{listEX}
	\loigiai{
		\begin{enumerate}
			\item $\sin x=\dfrac{\sqrt{3}}{2}\Leftrightarrow\sin x=\sin\dfrac{\pi}{3}\Leftrightarrow \hoac{&x=\dfrac{\pi}{3}+k2\pi\\&x=\pi-\dfrac{\pi}{3}+k2\pi}\Leftrightarrow\hoac{&x=\dfrac{\pi}{3}+k2\pi\\&x=\dfrac{2\pi}{3}+k2\pi}\,(k\in\mathbb{Z})$.
			\item $2\cos x=-\sqrt{2}\Leftrightarrow\cos x=-\dfrac{\sqrt{2}}{2}\Leftrightarrow\cos x=\cos\dfrac{3\pi}{4}\Leftrightarrow x=\pm\dfrac{3\pi}{4}+k2\pi,\,(k\in\mathbb{Z})$.
			\item $\cot(2x-1)=\cot\dfrac{\pi}{5}\Leftrightarrow2x-1=\dfrac{\pi}{5}+k\pi\Leftrightarrow x=\dfrac{\pi}{10}+\dfrac{1}{2}+\dfrac{k\pi}{2},\,(k\in\mathbb{Z})$.
		\end{enumerate}
	}
\end{bt}

% #1635 | Archive/2024/SP-K10-K11-LyVanHoang_TVN-023/11. C1-LuongGiac-lvhoang_TVN-023/dataC1K11/Chuong1-Bai4-Dang3-PhuongTrinhCBDungRadian.tex | mã gốc 1D1H5-3
\begin{bt}[Mức H, KNTT]%[BG Toán 11 CT2025, Quan Ón]%[1D1H5-3]
	Giải các phương trình sau
	\begin{listEX}[2]
		\item $\sin2x+\cos4x=0$;
		\item $\cos3x=-\cos7x$.
	\end{listEX}
	\loigiai{
		\begin{enumerate}
			\item $\sin2x+\cos4x=0\Leftrightarrow \cos4x=\sin(-2x)\Leftrightarrow\cos4x=\cos\left(\dfrac{\pi}{2}+2x\right)\Leftrightarrow\hoac{&4x=\dfrac{\pi}{2}+2x+k2\pi\\&4x=-\dfrac{\pi}{2}-2x+k2\pi}\\\Leftrightarrow\hoac{&x=\dfrac{\pi}{4}+k\pi\\&x=-\dfrac{\pi}{12}+\dfrac{k\pi}{3}}\,(k\in\mathbb{Z})$.
			\item $\cos3x=-\cos7x\Leftrightarrow\cos3x=\cos(\pi-7x)\Leftrightarrow\hoac{&3x=\pi-7x+k2\pi\\&3x=-\pi+7x+k2\pi}\Leftrightarrow \hoac{&x=\dfrac{\pi}{10}+\dfrac{k\pi}{5}\\&x=\dfrac{\pi}{4}-\dfrac{k\pi}{2}}\,(k\in\mathbb{Z})$.
		\end{enumerate}
	}
\end{bt}

% #1636 | Archive/2024/SP-K10-K11-LyVanHoang_TVN-023/11. C1-LuongGiac-lvhoang_TVN-023/dataC1K11/Chuong1-Bai4-Dang3-PhuongTrinhCBDungRadian.tex | mã gốc 1D1H5-3
\begin{bt}[Mức H, CTST]%[BG Toán 11 CT2025, Quan Ón]%[1D1H5-3]
	Giải các phương trình lượng giác sau:
	\begin{listEX}[3]
		\item $\sin 2 x=\dfrac{1}{2}$;
		\item $\sin \left(x-\dfrac{\pi}{7}\right)=\sin \dfrac{2 \pi}{7}$;
		\item $\sin 4 x-\cos \left(x+\dfrac{\pi}{6}\right)=0$;
		\item $\cos \left(x+\dfrac{\pi}{3}\right)=\dfrac{\sqrt{3}}{2}$;
		\item $\cos 4 x=\cos \dfrac{5 \pi}{12}$;
		\item $\cos ^2 x=1$;
		\item $\tan \left(2 x+\dfrac{\pi}{4}\right)=0$;
		\item $\cot \left(\dfrac{1}{2} x+\dfrac{\pi}{4}\right)=-1$;
		\item $\cot 3 x=-\dfrac{\sqrt{3}}{3}$.
	\end{listEX}
	\loigiai{
		\begin{enumerate}
			\item $\sin2x=\dfrac{1}{2}\Leftrightarrow\sin2x=\sin\dfrac{\pi}{6}\Leftrightarrow\hoac{&2x=\dfrac{\pi}{6}+k2\pi\\&2x=\dfrac{5\pi}{6}+k2\pi}\Leftrightarrow\hoac{&x=\dfrac{\pi}{12}+k\pi\\&x=\dfrac{5\pi}{12}+k\pi}\,(k\in\mathbb{Z})$.
			\item $\sin\left(x-\dfrac{\pi}{7}\right)=\sin \dfrac{2\pi}{7}\Leftrightarrow\hoac{&x-\dfrac{\pi}{7}=\dfrac{2\pi}{7}+k2\pi\\&x-\dfrac{\pi}{7}=\dfrac{5\pi}{7}+k2\pi}\Leftrightarrow\hoac{&x=\dfrac{3\pi}{7}+k2\pi\\&x=\dfrac{6\pi}{7}+k2\pi}\,(k\in\mathbb{Z})$.
			\item Ta có
			\allowdisplaybreaks
			\begin{eqnarray*}
				\sin4x-\cos \left(x+\dfrac{\pi}{6}\right)=0 \Leftrightarrow \sin4x=\sin\left(\dfrac{\pi}{3}-x\right)
				&\Leftrightarrow&
				\hoac{&4x=\dfrac{\pi}{3}-x+k2\pi\\&4x=\dfrac{2\pi}{3}+x+k2\pi}\\
				&\Leftrightarrow& \hoac{&x=\dfrac{\pi}{15}+\dfrac{k2\pi}{5}\\&x=\dfrac{2\pi}{9}+\dfrac{k2\pi}{3}}\,\, (k\in\mathbb{Z}).
			\end{eqnarray*}
			\item Ta có
			\allowdisplaybreaks
			\begin{eqnarray*}
				\cos \left(x+\dfrac{\pi}{3}\right) = \dfrac{\sqrt{3}}{2} \Leftrightarrow \cos\left(x+\dfrac{\pi}{3}\right)=\cos\dfrac{\pi}{6} &\Leftrightarrow& \hoac{&x+\dfrac{\pi}{3}=\dfrac{\pi}{6}+k2\pi\\&x+\dfrac{\pi}{3}=-\dfrac{\pi}{6}+k2\pi}\\
				&\Leftrightarrow& \hoac{&x=-\dfrac{\pi}{6}+k2\pi\\&x=-\dfrac{\pi}{2}+k2\pi}\,\,(k\in\mathbb{Z}).
			\end{eqnarray*}
			\item $\cos 4 x=\cos \dfrac{5 \pi}{12}\Leftrightarrow4x=\pm\dfrac{5\pi}{12}+k2\pi\Leftrightarrow x=\pm\dfrac{5\pi}{48}+\dfrac{k\pi}{2},\,(k\in\mathbb{Z})$.
			\item $\cos^2x=1\Leftrightarrow\sin^2x=0\Leftrightarrow\sin x=0\Leftrightarrow x=k\pi,\,(k\in\mathbb{Z})$.
			\item $\tan \left(2 x+\dfrac{\pi}{4}\right)=0\Leftrightarrow2x+\dfrac{\pi}{4}=k\pi\Leftrightarrow x=-\dfrac{\pi}{8}+\dfrac{k\pi}{2},\,(k\in\mathbb{Z})$.
			\item Ta có
			\allowdisplaybreaks
			\begin{eqnarray*}
				\cot \left(\dfrac{1}{2} x+\dfrac{\pi}{4}\right)=-1 \Leftrightarrow \cot\left(\dfrac{1}{2}x+\dfrac{\pi}{4}\right)=\cot\left(-\dfrac{\pi}{4}\right) &\Leftrightarrow& \dfrac{1}{2}x+\dfrac{\pi}{4}=-\dfrac{\pi}{4}+k\pi\\
				&\Leftrightarrow& x=(2k-1)\pi,\,\,(k\in\mathbb{Z}).
			\end{eqnarray*}
			\item $\cot3x=-\dfrac{\sqrt{3}}{3}\Leftrightarrow\cot3x=\cot\left(-\dfrac{\pi}{3}\right)\Leftrightarrow3x=-\dfrac{\pi}{3}+k\pi\Leftrightarrow x=-\dfrac{\pi}{9}+\dfrac{k\pi}{3},\,(k\in\mathbb{Z})$.
		\end{enumerate}
	}
\end{bt}

% #1637 | Archive/2024/SP-K10-K11-LyVanHoang_TVN-023/11. C1-LuongGiac-lvhoang_TVN-023/dataC1K11/Chuong1-Bai4-Dang3-PhuongTrinhCBDungRadian.tex | mã gốc 1D1H5-3
\begin{bt}[Mức H, CTST]%[BG Toán 11 CT2025, Quan Ón]%[1D1H5-3]
	Tại các giá trị nào của $x$ thì đồ thị hàm số $y=\cos x$ và $y=\sin x$ giao nhau?
	\loigiai{
		Phương trình hoành độ giao điểm của hai đồ thị hàm số $y=\cos x$ và $y=\sin x$ là $$\cos x=\sin x\Leftrightarrow\cos x=\cos\left(\dfrac{\pi}{2}-x\right)\Leftrightarrow x=\dfrac{\pi}{2}-x+k2\pi\Leftrightarrow x=\dfrac{\pi}{4}+k\pi,\,(k\in\mathbb{Z}).$$
		Tại các giá trị $x=\dfrac{\pi}{4}+k\pi,\,(k\in\mathbb{Z})$ thì đồ thị hàm số $y=\cos x$ và $y=\sin x$ giao nhau.
	}
\end{bt}

% #1649 | Archive/2024/SP-K10-K11-LyVanHoang_TVN-023/11. C1-LuongGiac-lvhoang_TVN-023/dataC1K11/Chuong1-Bai4-Dang2-Dieukienconghiem.tex | mã gốc 1D1H5-2
\begin{vd}[Mức độ 1]%[Bài Giảng 10-11 - Phan Anh]%[1D1H5-2]
	Tìm tất các các giá trị thực của tham số $m$ để phương trình $\sin x=m^2$ có nghiệm.
	\loigiai{
		Với mọi $x\in\mathbb{R}$, ta luôn có $-1\le\sin x\le 1$.\\
		Do đó, phương trình $\sin x=m^2$ có nghiệm khi và chỉ khi $m^2\le 1\Leftrightarrow-1\le m\le1$.}
\end{vd}

% #1652 | Archive/2024/SP-K10-K11-LyVanHoang_TVN-023/11. C1-LuongGiac-lvhoang_TVN-023/dataC1K11/Chuong1-Bai4-Dang2-Dieukienconghiem.tex | mã gốc 1D1H5-2
\begin{bt}[Mức độ 2]%[Bài Giảng 10-11 - Phan Anh]%[1D1H5-2]
	Xác định tất cả các giá trị của tham số $m$ để
	\begin{enumEX}{2}
		\item $\sin^2x-\cos^2x=m+1$ có nghiệm;
		\item $2\sin x+1=m$ vô nghiệm.
	\end{enumEX}
	\loigiai{
		\begin{enumerate}
			\item $\sin^2x-\cos^2x=m+1\Leftrightarrow\cos2x=-m-1$ có nghiệm khi và chỉ khi $-1\leq-m-1\leq1\Leftrightarrow-2\leq m\leq0$;
			\item $2\sin x+1=m\Leftrightarrow\sin x=\dfrac{m-1}{2}$ vô nghiệm khi và chỉ khi
			\[\left|\dfrac{m-1}{2}\right|>1\Leftrightarrow|m-1|>2\Leftrightarrow\hoac{&m-1>2\\&m-1<-2}\\\Leftrightarrow\hoac{&m>3\\&m<-1.}\]
		\end{enumerate}
	}
\end{bt}

% #1653 | Archive/2024/SP-K10-K11-LyVanHoang_TVN-023/11. C1-LuongGiac-lvhoang_TVN-023/dataC1K11/Chuong1-Bai4-Dang2-Dieukienconghiem.tex | mã gốc 1D1H5-2
\begin{bt}[Mức độ 2]%[Bài Giảng 10-11 - Phan Anh]%[1D1H5-2]
	Tìm tất các các giá trị thực của tham số $m$ để phương trình $\cos x-m=0$ vô nghiệm.
	\loigiai{
		Phương trình $\cos x-m=0\Leftrightarrow\cos x=m.$\\
		Phương trình $\cos x=m$ vô nghiệm $\Leftrightarrow\left| m\right|>1\Leftrightarrow\hoac{& m<-1\\& m>1.}$}
\end{bt}

% #1654 | Archive/2024/SP-K10-K11-LyVanHoang_TVN-023/11. C1-LuongGiac-lvhoang_TVN-023/dataC1K11/Chuong1-Bai4-Dang2-Dieukienconghiem.tex | mã gốc 1D1H5-2
\begin{bt}[Mức độ 2]%[Bài Giảng 10-11 - Phan Anh]%[1D1H5-2]
	Có bao nhiêu giá trị nguyên của tham số $m$ để phương trình $\sin x=m+1$ có nghiệm?
	\loigiai{
		Phương trình $\sin x=m+1$ có nghiệm khi và chỉ khi 
		\[\left| m+1\right|\le 1\Leftrightarrow-1\le m+1\le 1\Leftrightarrow-2\le m\le 0.\]
		Do $m\in\mathbb{Z}$ nên $m\in\left\{-2;-1;0\right\}$.}
\end{bt}

% #1662 | Archive/2024/SP-K10-K11-LyVanHoang_TVN-023/11. C1-LuongGiac-lvhoang_TVN-023/dataC1K11/Chuong1-Bai4-Dang4-PhuongTrinhCBDungDo.tex | mã gốc 1D1N5-4
\begin{vd}[Mức N, Cánh Diều]%[BG Toán 11 CT2025, Quan Ón]%[1D1N5-4]
	Giải các phương trình lượng giác sau
	\begin{multicols}{2}
		\begin{enumerate}
			\item $\sin x=\sin 55^\circ$;
			\item $\cos x=\cos(-87^\circ)$;
			\item $\tan x=\tan 67^\circ$;
			\item $\cot x=\cot (-83^\circ)$.
		\end{enumerate}
	\end{multicols}
	\loigiai{
		\begin{enumerate}
			\item Ta có
			\[
			\sin x=\sin 55^\circ \Leftrightarrow \hoac{&x=55^\circ+k360^\circ\\&x=180^\circ-55^\circ+k360^\circ} \Leftrightarrow
			\hoac{
				&x=55^\circ+k360^\circ\\
				&x=125^\circ+k360^\circ
			}\quad (k\in \mathbb{Z}).
			\]
			\item Ta có 
			\[
			\cos x=\cos (-87^\circ)
			\Leftrightarrow \hoac{
				&x=-87^\circ+k360^\circ\\
				&x=-(-87^\circ)+k360^\circ
			} \Leftrightarrow
			\hoac{
				&x=-87^\circ+k360^\circ\\
				&x=87^\circ+k360^\circ\\
			} \quad (k\in \mathbb{Z}).
			\]
			\item Ta có $\tan x=\tan 67^\circ \Leftrightarrow x=67^\circ +k180^\circ$, $k\in \mathbb{Z}$.
			\item Ta có $\cot x=\cot (-83^\circ) \Leftrightarrow x=-83^\circ+k180^\circ$, $(k\in \mathbb{Z})$.
		\end{enumerate}
	}
\end{vd}

% #1665 | Archive/2024/SP-K10-K11-LyVanHoang_TVN-023/11. C1-LuongGiac-lvhoang_TVN-023/dataC1K11/Chuong1-Bai4-Dang4-PhuongTrinhCBDungDo.tex | mã gốc 1D1H5-3
\begin{bt}[Mức H, CTST]%[BG Toán 11 CT2025, Quan Ón]%[1D1H5-3]
	Giải các phương trình sau
	\begin{listEX}[2]
		\item $\cos x=\cos15^\circ$;
		\item $\tan x=\tan 55^{\circ}$;
		\item $\tan(30^\circ-3x)=\tan75^\circ$;
		\item $\cot(3x+30^\circ)=\cot75^\circ$;
		\item $\sin(x+30^\circ)=\sin(x+60^\circ)$;
		\item $\cos 2x=\cos\left(x + 60^{\circ}\right)$.
	\end{listEX}
	\loigiai{
		\begin{enumerate}
			\item $\cos x=\cos15^\circ\Leftrightarrow x=\pm15^\circ+k360^\circ,\,k\in\mathbb{Z}$.
			\item $\tan x=\tan 55^{\circ}\Leftrightarrow x=55^\circ+k180^\circ,\,k\in\mathbb{Z}$.
			\item $\tan(30^\circ-3x)=\tan75^\circ\Leftrightarrow30^\circ-3x=75^\circ+k180^\circ\Leftrightarrow x=-15^\circ-k120^\circ,\,k\in\mathbb{Z}$.
			\item $\cot(3x+30^\circ)=\cot75^\circ\Leftrightarrow3x+30^\circ=75^\circ+k180^\circ\Leftrightarrow x=15^\circ+k60^\circ,\,k\in\mathbb{Z}$.
			\item $\sin(x+30^\circ)=\sin(x+60^\circ)|\Leftrightarrow\hoac{&x+30^\circ=x+60^\circ+k360^\circ\\&x+30^\circ=120^\circ-x+k360^\circ}\Leftrightarrow x=75^\circ+k180^\circ,\,k\in\mathbb{Z}$.
			\item $\cos 2x=\cos\left(x + 60^{\circ}\right)
			\Leftrightarrow \hoac{&2x=x + 60^{\circ} + k 360^{\circ}\\&2x= - \left(x + 60^{\circ}\right) + k 360^{\circ}}
			\Leftrightarrow \hoac{&x=60^{\circ} + k 360^{\circ}\\&x= - 20^{\circ} + k 120^{\circ}},\,k\in\mathbb{Z}$.
		\end{enumerate}
	}
\end{bt}

% #1666 | Archive/2024/SP-K10-K11-LyVanHoang_TVN-023/11. C1-LuongGiac-lvhoang_TVN-023/dataC1K11/Chuong1-Bai4-Dang4-PhuongTrinhCBDungDo.tex | mã gốc 1D1H5-3
\begin{bt}[Mức H, KNTT]%[BG Toán 11 CT2025, Quan Ón]%[1D1H5-3]
	Giải các phương trình sau
	\begin{listEX}[2]
		\item $\sin 2x  = \sin (60^{\circ} - 3x)$;
		\item $\cos 2 x=\cos \left(45^{\circ}-x\right)$.
	\end{listEX}
	\loigiai{
		\begin{enumerate}
			\item $\sin2x=\sin(60^{\circ}-3x)\Leftrightarrow\hoac{&2x=60^\circ-3x+k360^\circ\\&2x=3x+120^\circ+k360^\circ}\Leftrightarrow\hoac{&x=12^\circ+k72^\circ\\&x=-120^\circ-k360^\circ},\,(k \in \mathbb{Z})$.
			\item $\cos2x=\cos\left(45^{\circ}-x\right)  \Leftrightarrow\hoac{&2 x=45^{\circ}-x+k 360^{\circ}\\&2x=-\left(45^{\circ}-x\right)+k 360^{\circ}} \Leftrightarrow\hoac{&3x=45^{\circ}+k360^{\circ}\\&x=-45^{\circ}+k360^{\circ}}\\\Leftrightarrow \hoac{&x=15^{\circ}+k 120^{\circ} \\&x=-45^{\circ}+k 360^{\circ}},\,(k \in \mathbb{Z})$.
		\end{enumerate}
	}
	
\end{bt}

% #1676 | Archive/2024/SP-K10-K11-LyVanHoang_TVN-023/11. C1-LuongGiac-lvhoang_TVN-023/dataC1K11/Bai-3-Dang-5-gtlngtnncuahamsoluonggiac.tex | mã gốc 1D1H4-6
\begin{vd}%[Dự án BGK10-11,2024]%[Nguyễn Văn Nay]%[1D1H4-6]
	[Mức độ 2]
	Tìm giá trị lớn nhất và giá trị nhỏ nhất của hàm số $ y=5-3 \cos 4 x $.
	\loigiai
	{
		Tập xác định $ \mathscr{D}=\mathbb{R}$.\\
		Với mọi $x \in \mathbb{R} $, ta có
		\begin{eqnarray*}
			&&-1 \leq \cos 4 x \leq 1 \\
			&\Leftrightarrow& 3 \geq-3 \cos 4 x \geq-3\\
			&\Leftrightarrow& 5+3 \geq 5-3 \cos 4 x \geq 5-3 \\
			&\Leftrightarrow& 2 \leq y \leq 8. 
		\end{eqnarray*}
	\begin{itemize}
		\item $ \max y=8 $ khi $ \cos 4 x=-1 \Leftrightarrow x=\dfrac{\pi}{4}+\dfrac{k\pi}{2}$, $k\in\mathbb{Z}$.
		\item $ \min y=2 $ khi $ \cos 4 x=1 \Leftrightarrow x=\dfrac{k\pi}{2}$, $k\in\mathbb{Z}$.
	\end{itemize}
	}
\end{vd}

% #1677 | Archive/2024/SP-K10-K11-LyVanHoang_TVN-023/11. C1-LuongGiac-lvhoang_TVN-023/dataC1K11/Bai-3-Dang-5-gtlngtnncuahamsoluonggiac.tex | mã gốc 1D1H4-6
\begin{vd}%[Dự án BGK10-11,2024]%[Nguyễn Văn Nay]%[1D1H4-6]
	[Mức độ 2]
	Tìm giá trị lớn nhất và giá trị nhỏ nhất của hàm số $ y=\dfrac{1+4 \cos^2 x}{3} $.
	\loigiai
	{
		Tập xác định: $ \mathscr{D}=\mathbb{R}$. $ \forall x \in \mathbb{R} $, ta có 
		\begin{eqnarray*}
			&& 0\leq \cos^2 x \leq 1\\
			&\Leftrightarrow& 0\leq 4\cos^2 x \leq 4\\
			&\Leftrightarrow& 1\leq 1+4\cos^2 x \leq 5\\
			&\Leftrightarrow& \dfrac{1}{3}\leq \dfrac{1+4\cos^2 x}{3} \leq \dfrac{5}{3}.
		\end{eqnarray*}
	\begin{itemize}
		\item $ \max y=\dfrac{5}{3} $ khi $ \cos x=0 \Leftrightarrow x=\dfrac{\pi}{2}+k\pi$, $k\in \mathbb{Z}$.
		\item $ \min y=\dfrac{1}{3} $ khi $ \cos  x=\pm1 \Leftrightarrow x=k\pi$, $k\in \mathbb{Z}$.
	\end{itemize}
	}
\end{vd}

% #1678 | Archive/2024/SP-K10-K11-LyVanHoang_TVN-023/11. C1-LuongGiac-lvhoang_TVN-023/dataC1K11/Bai-3-Dang-5-gtlngtnncuahamsoluonggiac.tex | mã gốc 1D1H4-6
\begin{vd}%[Dự án BGK10-11,2024]%[Nguyễn Văn Nay]%[1D1H4-6]
	[Mức độ 2]
	Tìm giá trị lớn nhất và giá trị nhỏ nhất của hàm số $ y=\cos x+\cos \left(x+\dfrac{\pi}{3}\right) $.
	\loigiai
	{
		Áp dụng công thức \fbox{$ \cos a+\cos b=2 \cos \dfrac{a+b}{2} \cos \dfrac{a-b}{2} $}, ta có
		\begin{equation*}
			y=\cos x+\cos \left (x+\dfrac{\pi}{3}\right )=\sqrt{3}\cos \left (x+\dfrac{\pi}{6}\right ).
		\end{equation*}
		Tập xác định: $ \mathscr{D}=\mathbb{R}$. $ \forall x \in \mathbb{R} $, ta có
		\begin{eqnarray*}
			&&-1 \leq \cos \left (x+\dfrac{\pi}{6}\right ) \leq 1\\
			&\Leftrightarrow&  -\sqrt{3} \leq \sqrt{3}\cos \left (x+\dfrac{\pi}{6}\right ) \leq \sqrt{3}.
		\end{eqnarray*}
	\begin{itemize}
		\item $ \max y=\sqrt{3} $ khi $ \cos \left (x+\dfrac{\pi}{6}\right ) =1 \Leftrightarrow x=-\dfrac{\pi}{6}+k2\pi$, $k\in \mathbb{Z}$.
		\item $ \min y=-\sqrt{3} $ khi $ \cos \left (x+\dfrac{\pi}{6}\right )=-1 \Leftrightarrow x=\dfrac{5\pi}{6}+k2\pi$, $k\in \mathbb{Z}$.
	\end{itemize}
	}
\end{vd}

% #1681 | Archive/2024/SP-K10-K11-LyVanHoang_TVN-023/11. C1-LuongGiac-lvhoang_TVN-023/dataC1K11/Bai-3-Dang-5-gtlngtnncuahamsoluonggiac.tex | mã gốc 1D1H4-6
\begin{bt}%[Dự án BGK10-11,2024]%[Nguyễn Văn Nay]%[1D1H4-6]
	[Mức độ 2]
	Tìm giá trị lớn nhất và giá trị nhỏ nhất của hàm số $ y=1-\dfrac{1}{2} \sin^2 2 x $.
	\loigiai
	{
		Tập xác định: $ \mathscr{D}=\mathbb{R}$. $ \forall x \in \mathbb{R} $, ta có 
		\begin{eqnarray*}
			&& 0 \leq  \sin^2 2 x \leq 1\\
			&\Leftrightarrow& -\dfrac{1}{2} \leq  -\dfrac{1}{2}\cdot\sin^2 2 x \leq 0\\
			&\Leftrightarrow& \dfrac{1}{2} \leq  1-\dfrac{1}{2}\cdot\sin^2 2 x \leq 1.
		\end{eqnarray*}
	\begin{itemize}
		\item $ \max y=1 $ khi $ \sin  2 x=0 \Leftrightarrow x=\dfrac{k\pi}{2}$, $k\in \mathbb{Z}$.
		\item $ \min y=\dfrac{1}{2} $ khi $ \sin  2 x=\pm1 \Leftrightarrow x=\dfrac{\pi}{4}+\dfrac{k\pi}{2}$, $k\in \mathbb{Z}$.
	\end{itemize}
	}
\end{bt}

% #1682 | Archive/2024/SP-K10-K11-LyVanHoang_TVN-023/11. C1-LuongGiac-lvhoang_TVN-023/dataC1K11/Bai-3-Dang-5-gtlngtnncuahamsoluonggiac.tex | mã gốc 1D1H4-6
\begin{bt}%[Dự án BGK10-11,2024]%[Nguyễn Văn Nay]%[1D1H4-6]
	[Mức độ 2]
	Tìm giá trị lớn nhất và giá trị nhỏ nhất của hàm số $ y=3-2|\sin 2 x| $.
	\loigiai
	{
		Tập xác định: $ \mathscr{D}=\mathbb{R}$. $ \forall x \in \mathbb{R} $, ta có 
		\begin{eqnarray*}
			&&0\leq	|\sin 2 x|\leq 1\\
			&\Leftrightarrow&  -2\leq	-2|\sin 2 x|\leq 0\\
			&\Leftrightarrow&  1\leq	3-2|\sin 2 x|\leq 3.
		\end{eqnarray*}
	\begin{itemize}
		\item $ \max y=3 $ khi $ \sin 2 x=0 \Leftrightarrow x=\dfrac{k\pi}{2}$, $k\in \mathbb{Z}$.
		\item $ \min y=1$ khi $ \sin 2 x=\pm1 \Leftrightarrow x=\dfrac{\pm\pi}{4}+k\pi$, $k\in \mathbb{Z}$.
	\end{itemize}
	}
\end{bt}

% #1683 | Archive/2024/SP-K10-K11-LyVanHoang_TVN-023/11. C1-LuongGiac-lvhoang_TVN-023/dataC1K11/Bai-3-Dang-5-gtlngtnncuahamsoluonggiac.tex | mã gốc 1D1V4-6
\begin{bt}%[Dự án BGK10-11,2024]%[Nguyễn Văn Nay]%[1D1V4-6]
	[Mức độ 3]
	Tìm giá trị lớn nhất và giá trị nhỏ nhất của hàm số $ y=\sin x+\sin \left(x+{\dfrac{2\pi}{3}}\right)$.
	\loigiai
	{
		Áp dụng công thức \fbox{$ \sin a+\sin b=2 \sin \dfrac{a+b}{2} \cos \dfrac{a-b}{2} $}, ta có
		\begin{equation*}
			y=\sin x + \sin \left (x+\dfrac{2\pi}{3}\right )=-\sin \left (x+\dfrac{\pi}{3}\right ).
		\end{equation*}
		Tập xác định $ \mathscr{D}=\mathbb{R}$. $ \forall x \in \mathbb{R} $, ta có 
		\begin{equation*}
			-1 \leq -\sin \left (x+\dfrac{\pi}{3}\right ) \leq 1.
		\end{equation*}
	\begin{itemize}
		\item $ \max y=1 $ khi $ \sin \left (x+\dfrac{\pi}{3}\right )=-1 \Leftrightarrow x=-\dfrac{5\pi}{6}+k2\pi$, $k\in \mathbb{Z}$.
		\item $ \min y=-1 $ khi $ \sin \left (x+\dfrac{\pi}{3}\right )=1 \Leftrightarrow x=\dfrac{\pi}{6}+k2\pi$, $k\in \mathbb{Z}$.
	\end{itemize}
	}
\end{bt}

% #1684 | Archive/2024/SP-K10-K11-LyVanHoang_TVN-023/11. C1-LuongGiac-lvhoang_TVN-023/dataC1K11/Bai-3-Dang-5-gtlngtnncuahamsoluonggiac.tex | mã gốc 1D1V4-6
\begin{bt}%[Dự án BGK10-11,2024]%[Nguyễn Văn Nay]%[1D1V4-6]
	[Mức độ 3]
	Tìm giá trị lớn nhất và giá trị nhỏ nhất của hàm số $y=\sqrt{3} \sin x-\cos x+5$.
	\loigiai
	{
		Tập xác định $\mathscr{D}=\mathbb{R}$.\\
		Biến đổi $y=\sqrt{3} \sin x-\cos x+5=2\left(\dfrac{\sqrt{3}}{2}\cdot\sin x-\dfrac{1}{2}\cdot\cos x\right)+5=2\sin\left(x-\dfrac{\pi}{6}\right)+5$.\\
		Với mọi $x\in \mathbb{R}$ ta có
		\allowdisplaybreaks
		\begin{eqnarray*}
			& & -1\leq \sin\left(x-\dfrac{\pi}{6}\right)\leq 1\\
			&\Leftrightarrow& -2\leq 2\sin\left(x-\dfrac{\pi}{6}\right)\leq 2\\
			&\Leftrightarrow&3\leq  2\sin\left(x-\dfrac{\pi}{6}\right)+5\leq 7.
		\end{eqnarray*}
		Vậy $\max\limits_{x \in \mathbb{R}} y=7$ khi $x=\dfrac{2\pi}{3}+k2\pi$, $k\in \mathbb{Z}$ và $\min\limits_{x \in \mathbb{R}} y=3$ khi $x=-\dfrac{\pi}{3}+k2\pi$, $k\in \mathbb{Z}$.
	}
\end{bt}

% #1685 | Archive/2024/SP-K10-K11-LyVanHoang_TVN-023/11. C1-LuongGiac-lvhoang_TVN-023/dataC1K11/Bai-3-Dang-5-gtlngtnncuahamsoluonggiac.tex | mã gốc 1D1V4-6
\begin{bt}%[Dự án BGK10-11,2024]%[Nguyễn Văn Nay]%[1D1V4-6]
	[Mức độ 3]
	Tìm giá trị lớn nhất và giá trị nhỏ nhất của hàm số  $y=4 \sin^2 x-4 \sin x+3$.
	\loigiai
	{
		Tập xác định $\mathscr{D}=\mathbb{R}$.\\
		Biến đổi $y=4 \sin^2 x-4 \sin x+3=\left(2\sin x+1\right)^2+2$.\\
		Với mọi $x\in \mathbb{R}$ ta có
		\allowdisplaybreaks
		\begin{eqnarray*}
			& & -1\leq \sin x\leq 1\\
			&\Leftrightarrow& -1\leq 2\sin x+1\leq 3\\
			&\Leftrightarrow&0\leq \left( 2\sin x+1\right)^2\leq 9\\
			&\Leftrightarrow&2\leq \left( 2\sin x+1\right)^2+2\leq 11.
		\end{eqnarray*}
		Vậy $\max\limits_{x \in \mathbb{R}} y=11$ khi $x=\dfrac{\pi}{2}+k2\pi$, $k\in \mathbb{Z}$ và $\min\limits_{x \in \mathbb{R}} y=2$ khi $x=-\dfrac{\pi}{6}+k2\pi$ hoặc $x=\dfrac{7\pi}{6}+k2\pi$, $k\in \mathbb{Z}$.
	}
\end{bt}

% #1686 | Archive/2024/SP-K10-K11-LyVanHoang_TVN-023/11. C1-LuongGiac-lvhoang_TVN-023/dataC1K11/Bai-3-Dang-5-gtlngtnncuahamsoluonggiac.tex | mã gốc 1D1V4-6
\begin{bt}%[Dự án BGK10-11,2024]%[Nguyễn Văn Nay]%[1D1V4-6]
	[Mức độ 3]
	Tìm giá trị lớn nhất và giá trị nhỏ nhất của hàm số  $y=\cos^2 x-2 \cos x-4$.
	\loigiai
	{
		Tập xác định $\mathscr{D}=\mathbb{R}$.\\
		Biến đổi $y=\cos^2 x-2 \cos x-4=\left(\cos x-1\right)^2-5$.\\
		Với mọi $x\in \mathbb{R}$ ta có
		\allowdisplaybreaks
		\begin{eqnarray*}
			& & -1\leq \cos x\leq 1\\
			&\Leftrightarrow& -2\leq \cos x-1\leq 0\\
			&\Rightarrow&0\leq \left( \cos x-1\right)^2\leq 4\\
			&\Leftrightarrow&-5\leq \left( \cos x-1\right)^2-5\leq -1.
		\end{eqnarray*}
		Vậy $\max\limits_{x \in \mathbb{R}} y=-1$ khi $x=\pi+k2\pi$, $k\in \mathbb{Z}$ và $\min\limits_{x \in \mathbb{R}} y=-5$ khi $x=k2\pi$, $k\in \mathbb{Z}$.
	}
\end{bt}

% #1687 | Archive/2024/SP-K10-K11-LyVanHoang_TVN-023/11. C1-LuongGiac-lvhoang_TVN-023/dataC1K11/Bai-3-Dang-5-gtlngtnncuahamsoluonggiac.tex | mã gốc 1D1V4-6
\begin{bt}%[Dự án BGK10-11,2024]%[Nguyễn Văn Nay]%[1D1V4-6]
	[Mức độ 3]
	Tìm giá trị lớn nhất và giá trị nhỏ nhất của hàm số $y=\cos^4 x-2 \sin^2 x+1$.
	\loigiai
	{
		Tập xác định $\mathscr{D}=\mathbb{R}$.\\
		Biến đổi $y=\cos^4 x-2 \sin^2 x+1=\cos^4 x-2 \left(1-\cos^2 x\right)+1=\cos^4 x+2 \cos^2 x-1=\left(\cos^2 x+1\right)^2-2 $.\\
		Với mọi $x\in \mathbb{R}$ ta có
		\allowdisplaybreaks
		\begin{eqnarray*}
			& & 0\leq \cos^2 x\leq 1\\
			&\Leftrightarrow& 1\leq \cos^2 x+1\leq 2\\
			&\Leftrightarrow&1\leq \left( \cos^2 x+1\right)^2\leq 4\\
			&\Leftrightarrow&-1\leq \left( \cos^2 x+1\right)^2-2\leq 2.
		\end{eqnarray*}
		Vậy $\max\limits_{x \in \mathbb{R}} y=2$ khi $x=k\pi$, $k\in \mathbb{Z}$ và $\min\limits_{x \in \mathbb{R}} y=-1$ khi $x=\dfrac{\pi}{2}+k\pi$, $k\in \mathbb{Z}$.
	}
\end{bt}

% #1688 | Archive/2024/SP-K10-K11-LyVanHoang_TVN-023/11. C1-LuongGiac-lvhoang_TVN-023/dataC1K11/Bai-3-Dang-5-gtlngtnncuahamsoluonggiac.tex | mã gốc 1D1C4-6
\begin{bt}%[Dự án BGK10-11,2024]%[Nguyễn Văn Nay]%[1D1C4-6]
	[Mức độ 4]
	Tìm giá trị lớn nhất và giá trị nhỏ nhất của hàm số $y=\sin x+\cos x+2 \sin x \cos x-1$.
	\loigiai
	{
		Tập xác định $\mathscr{D}=\mathbb{R}$.\\
		Đặt $t=\sin x+ \cos x=\sqrt{2}\sin \left(x+\dfrac{\pi}{4}\right)$, $t\in \left[-\sqrt{2};\sqrt{2}\right]$.\\
		Ta có $t^2=\left(\sin x+ \cos x\right)^2=1+2\sin x\cos x\Rightarrow 2\sin x\cos x =t^2-1$.\\
		Hàm số trở thành $y=g(t)=t^2+t-2$. \\
		Bảng biến thiên của hàm số $y=g(t)$ trên đoạn $ \left[-\sqrt{2};\sqrt{2}\right]$
		\begin{center}
			\begin{tikzpicture}
				\tkzTabInit[nocadre=true,lgt=1,espcl=3,deltacl=1]%nocadre,
				{$t$/1, $g(t)$ /2}
				{$-\sqrt{2}$ , $-\dfrac{1}{2}$ ,  $\sqrt{2}$}
				\tkzTabVar {+/$-\sqrt{2}$,-/$-\dfrac{9}{4}$ ,+/ $\sqrt{2}$}
			\end{tikzpicture}
		\end{center}
		Vậy $\max\limits_{x \in \mathbb{R}} y=\sqrt{2}$ và $\min\limits_{x \in \mathbb{R}} y=-\dfrac{9}{4}$.
	}
\end{bt}

% #1689 | Archive/2024/SP-K10-K11-LyVanHoang_TVN-023/11. C1-LuongGiac-lvhoang_TVN-023/dataC1K11/Bai-3-Dang-5-gtlngtnncuahamsoluonggiac.tex | mã gốc 1D1C4-6
\begin{bt}%[Dự án BGK10-11,2024]%[Nguyễn Văn Nay]%[1D1C4-6]
	[Mức độ 4]
	Tìm giá trị lớn nhất và giá trị nhỏ nhất của hàm số $y=\sin \left(2 x+\dfrac{\pi}{4}\right)+\dfrac{1}{2}$ trên $\left[-\dfrac{\pi}{4} ; \dfrac{\pi}{4}\right]$.
	\loigiai
	{\immini
		{Ta có $x\in \left[ -\dfrac{\pi}{4} ; \dfrac{\pi}{4}\right]\Rightarrow 2x+\dfrac{\pi}{4}\in \left[ -\dfrac{\pi}{4}; \dfrac{3\pi}{4}\right]$.\\
			Do đó: 
			\allowdisplaybreaks
			\begin{eqnarray*}
				& & -\dfrac{\sqrt{2}}{2}\leq \sin \left(2 x+\dfrac{\pi}{4}\right)\leq 1\\
				&\Rightarrow & \dfrac{1-\sqrt{2}}{2}\leq \sin \left(2 x+\dfrac{\pi}{4}\right)+\dfrac{1}{2}\leq \dfrac{3}{2}.
			\end{eqnarray*}
			Vậy $\max\limits_{x \in\left[-\tfrac{\pi}{4} ; \tfrac{\pi}{4}\right]} y=\dfrac{3}{2}$ và $\min\limits_{x \in \left[-\tfrac{\pi}{4} ; \tfrac{\pi}{4}\right]} y=\dfrac{1-\sqrt{2}}{2}$.
		}
		{
			\begin{tikzpicture}[line join=round, line cap=round,>=stealth,thick,scale=1.4]
				\tikzset{label style/.style={font=\footnotesize}}
				\draw[->] (-1.2,0)--(1.4,0) node[below] {$x$};
				\draw[->] (0,-1.2)--(0,1.3) node[left] {$y$};
				\draw (0,0) circle (1);
				\draw (0,0) node [below left] {$O$};
				\path
				(0,0) coordinate (O)
				($(O) + (-45:1.1)$) coordinate (A)
				($(O) + (135:1.1)$) coordinate (B)
				;
				\draw[red,<->] (A) arc (-45:135:1.1 cm);
				\draw (O) -- (A) node [right] {\scriptsize $-\frac{\pi}{4}$};
				\draw (O) -- (B) node [left] {\scriptsize $\frac{3\pi}{4}$};
				\draw[dashed] (-45:1) -- (0,-0.7) node [left] {\scriptsize $-\frac{\sqrt{2}}{2}$};
				\draw[dashed] (135:1) -- (0,0.7) node [right] {\scriptsize $\frac{\sqrt{2}}{2}$};
			\end{tikzpicture}
		}
	}
\end{bt}

% #1696 | Archive/2024/SP-K10-K11-LyVanHoang_TVN-023/11. C1-LuongGiac-lvhoang_TVN-023/dataC1K11/Bai-3-Dang-5-gtlngtnncuahamsoluonggiac.tex | mã gốc 1D1V4-6
\begin{ex}%[Dự án BGK10-11,2024]%[Nguyễn Văn Nay]%[1D1V4-6]
	[Mức độ 3]
	Số giờ có ánh sáng của một thành phố $X$ ở vĩ độ $40^\circ$ Bắc trong ngày thứ $t$ của một năm không nhuận được cho bởi hàm số $d(t)=3\sin\left[\dfrac{\pi}{182}(t-80)\right]+12$, $t\in\mathbb{Z}$ và $0<t\le 365$. Vào ngày thứ bao nhiêu trong năm thì thành phố $X$ có nhiều giờ ánh sáng nhất?
	\shortans[]{$171$}
	\loigiai
	{
		Ta có $d(t)=3\sin\left[\dfrac{\pi}{182}(t-80)\right]+12\le 3+12=15$.\\
		Nên thành phố $X$ có nhiều giờ ánh sáng nhất thì $d(t)=15$.\\
		Đẳng thức xảy ra khi $\sin\left[\dfrac{\pi}{182}(t-80)\right]=1\Leftrightarrow \dfrac{\pi}{182}(t-80)=\dfrac{\pi}{2}+k2\pi$, $k\in\mathbb{Z}$
		\[\Leftrightarrow \dfrac{t-80}{182}=\dfrac{4k+1}{2}\Leftrightarrow t=91(4k+1)+80=364k+171.\]
		Do $0<t\le 365$ nên $k=0$, khi đó với $t=171$ thì thành phố $X$ có nhiều giờ ánh sáng nhất.
	}
\end{ex}

% #1697 | Archive/2024/SP-K10-K11-LyVanHoang_TVN-023/11. C1-LuongGiac-lvhoang_TVN-023/dataC1K11/Bai-3-Dang-5-gtlngtnncuahamsoluonggiac.tex | mã gốc 1D1V4-6
\begin{ex}%[Dự án BGK10-11,2024]%[Nguyễn Văn Nay]%[1D1V4-6]
	[Mức độ 3]
	Hằng ngày, mực nước của con kênh lên xuống theo thủy triều. Độ cao $h$ (m) của mực nước trong kênh tính theo thời gian $t$ (h) được cho bởi công thức $h=3\cos\left(\dfrac{\pi t}{6}+\dfrac{\pi}{3}\right)+12$. Sớm nhất vào lúc mấy giờ thì mực nước của kênh là cao nhất?
	\shortans[]{$10$}
	\loigiai
	{
		Ta có $h=3\cos\left(\dfrac{\pi t}{6}+\dfrac{\pi}{3}\right)+12\le 3+12=15$ nên mực nước của kênh cao nhất là $h=15$ (m).\\
		Đẳng thức xảy ra khi $\cos\left(\dfrac{\pi t}{6}+\dfrac{\pi}{3}\right)=1\Leftrightarrow \dfrac{\pi t}{6}+\dfrac{\pi}{3}=k2\pi\Leftrightarrow t=12k-2$, $k\in\mathbb{Z}$.\\
		Do thời gian sớm nhất và dương nên $k=1$, khi đó $t=10$ (h).
	}
\end{ex}

% #1700 | Archive/2024/SP-K10-K11-LyVanHoang_TVN-023/11. C1-LuongGiac-lvhoang_TVN-023/dataC1K11/Chuong1-Bai4-Dang6-Toanthucte-lienmon.tex | mã gốc 1D1V5-6
\begin{vd}[KNTT]%[1D1V5-6]%[Anh Duy]
	Khi mặt trăng quay quanh trái đất, mặt đối diện với trái đất thường chỉ được mặt trời chiếu sáng một phần. Các pha của Mặt Trăng mô tả mức độ phần bề mặt của nó được Mặt Trời chiếu sáng. Khi các góc giữa Mặt Trời, Trái Đất và Mặt Trăng là $\alpha$ ($0^\circ\leq\alpha\leq360^\circ$) thì tỉ lệ $F$ của phần Mặt Trăng được chiếu sáng cho bởi công thức $$F=\dfrac{1}{2}(1-\cos\alpha).$$
	Xác định góc $\alpha$ tương ứng với các pha sau của Mặt Trăng:
	\immini{\begin{enumerate}[a)]
			\item $F=0$ (trăng mới);
			\item $F=0{,}25$ (trăng lưỡi liềm);
			\item $F=0{,}5$ (trăng bán nguyệt đầu tháng hoặc trăng bán nguyệt cuối tháng);
			\item $F=1$ (trăng tròn).
	\end{enumerate}}{\begin{tikzpicture}[scale=1, font=\footnotesize, line join=round, line cap=round, >=stealth]
			\coordinate[label=below:$\text{Trái đất}$] (A) at (0,0);
			\coordinate[label=above:$\text{Mặt trời}$] (B) at (8,0);
			\draw[dashed] (0,0) circle[radius=1cm]; 
			\coordinate [label=above right:$\text{Mặt trăng}$](C) at ($(0,0)+(40:1cm)$); 
			\draw (C)--(A)--(B);
			\draw pic[draw,blue,angle radius=3mm,angle eccentricity=1.5,"$\alpha$"] {angle = B--A--C}; % Vẽ góc BAC với tên nhãn \beta. Nếu vẽ góc vuông thì thay angle = thành right angle =
			
			\foreach \diem in {A,B,C}	\fill (\diem)circle(3pt);
			%\newcommand{\gocv}[4][black]{\draw[#1] ($(#3)!5pt!(#2)$)--($(#3)!2!($($(#3)!5pt!(#2)$)!.5!($(#3)!5pt!(#4)$)$)$)--($(#3)!5pt!(#4)$);}
			%\gocv{S}{A}{D}
		\end{tikzpicture}
	}
	\loigiai{
		\begin{enumerate}[a)]
			\item $F=0\Leftrightarrow \dfrac{1}{2}(1-\cos\alpha)=0\Leftrightarrow\cos\alpha=1\Leftrightarrow\alpha=k360^\circ\Rightarrow\alpha=0^\circ,360^\circ$;
			\item $F=0{,}25\Leftrightarrow\dfrac{1}{2}(1-\cos\alpha)=0{,}25\Leftrightarrow\cos\alpha=\dfrac{1}{2}\Leftrightarrow\alpha=\pm60^\circ+k360^\circ\Rightarrow\alpha=60^\circ,300^\circ$;
			\item $F=0{,}5\Leftrightarrow\dfrac{1}{2}(1-\cos\alpha)=0{,}5\Leftrightarrow\cos\alpha=0\Leftrightarrow\alpha=90^\circ+k180^\circ\Rightarrow\alpha=90^\circ,270^\circ$;
			\item $F=1\Leftrightarrow\dfrac{1}{2}(1-\cos\alpha)=1\Leftrightarrow\cos\alpha=-1\Leftrightarrow\alpha=180^\circ+k360^\circ\Rightarrow\alpha=180^\circ$.
		\end{enumerate}
	}
\end{vd}

% #1702 | Archive/2024/SP-K10-K11-LyVanHoang_TVN-023/11. C1-LuongGiac-lvhoang_TVN-023/dataC1K11/Chuong1-Bai4-Dang6-Toanthucte-lienmon.tex | mã gốc 1D1H5-6
\begin{vd}%[1D1H5-6]%[Anh Duy]
	Giả sử một vật dao động điều hoà xung quanh vị trí cân bằng theo phương trình
	$$
	x=2 \cos \left(5 t-\dfrac{\pi}{6}\right).
	$$
	Ở đây, thời gian $t$ tính bằng giây và quãng đường $x$ tính bằng centimét. 
	\choiceTF[t]
	{\True Phương trình dao động là hàm số có chu kì $T=\dfrac{2\pi}{5}$}
	{Vật qua vị trí cân bằng khi $x=\pm 1$}
	{\True Vật qua vị trí cân bằng khi $t=\dfrac{2\pi}{15}+\dfrac{k\pi}{5} \,\left(k\in \mathbb{Z}\right)$}
	{Trong khoảng thời gian từ $ 0 $ đến $ 6 $ giây, vật đi qua vị trí cân bằng $8$ lần}
	\loigiai{
		\begin{itemchoice}
			\itemch Đúng. Chu kì $T=\dfrac{2\pi}{5}$.
			\itemch Sai. Vật qua vị trí cân bằng khi $x=0$.
			\itemch	Đúng. Vật đi qua vị trí cân bằng khi và chỉ khi \[ x=0 \Leftrightarrow 2 \cos \left(5 t-\dfrac{\pi}{6}\right)=0 \Leftrightarrow \cos \left(5 t-\dfrac{\pi}{6}\right)=0\\
			\Leftrightarrow 5 t-\dfrac{\pi}{6}=\dfrac{\pi}{2}+k\pi \Leftrightarrow t=\dfrac{2\pi}{15}+\dfrac{k\pi}{5} \left(k\in \mathbb{Z}\right).\]
			\itemch Sai. Vì $ 0\leq t \leq 6 $ nên $ 0\leq\dfrac{2\pi}{15}+\dfrac{k\pi}{5}\leq 6 \Leftrightarrow -0{,}6\leq k \leq 8{,}8$ mà $ k \in \mathbb{Z} $. Do đó $ k\in \{0;1;2;3;4;5;6;7;8\} $.\\
			Vậy vật đi qua vị trí cân bằng $ 9 $ lần. 
		\end{itemchoice}
		
	}
\end{vd}

% #1703 | Archive/2024/SP-K10-K11-LyVanHoang_TVN-023/11. C1-LuongGiac-lvhoang_TVN-023/dataC1K11/Chuong1-Bai4-Dang6-Toanthucte-lienmon.tex | mã gốc 1D1H5-6
\begin{vd}%[1D1H5-6]%[Anh Duy]
	Độ sâu $h$ (m) của mực nước biến tại thời điểm $t$ (giờ) sau khi thủy triều lên lần đầu tiên trong ngày được tính xấp xỉ bởi công thức $h(t)=0{,}8\cos 0{,}5t+4$. Tính độ sâu của nước biển tại thời điểm $t=2$ là bao nhiêu mét? (Làm tròn đến hàng phần trăm).
	\loigiai{
		Độ sâu mực nước biển tại thời điểm $t=2$ giây là $0{,}8\cos 0{,}5\cdot 2+4=4{,}43$ (m).
	}
\end{vd}

% #1705 | Archive/2024/SP-K10-K11-LyVanHoang_TVN-023/11. C1-LuongGiac-lvhoang_TVN-023/dataC1K11/Chuong1-Bai4-Dang6-Toanthucte-lienmon.tex | mã gốc 1D1V5-6
\begin{ex}%[1D1V5-6]%[Anh Duy]
	Số giờ có ánh sáng mặt trời của một thành phố $A$ ở vĩ độ $40^\circ$ Bắc trong ngày thứ $t$ của một năm không nhuận được cho bởi hàm số
	\[
	d(t)=3\sin \left[\dfrac{\pi}{182}(t-80)\right] +12 \text{ với } t\in \mathbb{Z} \text{ và } 0<t\leq 365.
	\]
	\begin{listEX}[1]
		\item Thành phố $A$ có đúng $12$ giờ có ánh sáng mặt trời vào ngày nào trong năm?
		\item Vào ngày nào trong năm thành phố $A$ có đúng $9$ giờ có ánh sáng mặt trời?
	\end{listEX}
	\loigiai{
		\begin{listEX}[1]
			\item Thành phố $A$ có đúng $12$ giờ có ánh sáng mặt trời vào ngày nào trong năm?\\
			Ta giải phương trình 
			\begin{eqnarray*}
				d(t)=12 &\Leftrightarrow& 3\sin\left[\dfrac{\pi}{182}(t-80)\right] +12=12 \\
				&\Leftrightarrow& \sin\left[\dfrac{\pi}{182}(t-80)\right]=0\\
				&\Leftrightarrow& \dfrac{\pi}{182}(t-80) =k \pi \\
				&\Leftrightarrow& t-80 = 182k\\
				&\Leftrightarrow& t= 182k+80 (k\in \mathbb{Z}).	
			\end{eqnarray*}
			Ta lại có  \[0<182k+80 \le 365 \Leftrightarrow -\dfrac{80}{182}<k \le \dfrac{285}{182} \Leftrightarrow \hoac{&k=0\\ &k=1.}\]
			Vậy thành phố $A$ có đúng $12$ giờ có ánh sáng mặt trời vào ngày thứ $80$ (ứng với $k=0$) và ngày thứ $262$ (ứng với $k=1$) trong năm.
			\item Vào ngày nào trong năm thành phố $A$ có đúng $9$ giờ có ánh sáng mặt trời?\\
			Ta giải phương trình 
			\begin{eqnarray*}
				d(t)=9 &\Leftrightarrow& 3\sin\left[\dfrac{\pi}{182}(t-80)\right] +12=9 \\
				&\Leftrightarrow& \sin\left[\dfrac{\pi}{182}(t-80)\right]=-1\\
				&\Leftrightarrow& \dfrac{\pi}{182}(t-80) =-\dfrac{\pi}{2}+k2 \pi \\
				&\Leftrightarrow& t-80 = -91+364k\\
				&\Leftrightarrow& t= 364k-11 (k\in \mathbb{Z}).	
			\end{eqnarray*}
			Ta lại có  \[0<364k-11 \le 365 \Leftrightarrow \dfrac{11}{364}<k \le \dfrac{376}{364} \Leftrightarrow k=1.\]
			Vậy thành phố $A$ có đúng $9$ giờ có ánh sáng mặt trời vào ngày thứ $353$ trong năm.
		\end{listEX}
	}
\end{ex}

% #1706 | Archive/2024/SP-K10-K11-LyVanHoang_TVN-023/11. C1-LuongGiac-lvhoang_TVN-023/dataC1K11/Chuong1-Bai4-Dang6-Toanthucte-lienmon.tex | mã gốc 1D1V5-6
\begin{ex}%[1D1V5-6]%[Anh Duy]
	Hội Lim (tỉnh Bắc Ninh) được tổ chức vào mùa xuân thường có trò chơi đánh đu. 
	\immini{
		Khi người chơi đu nhún đều, cây đu sẽ đưa người chơi đu dao động quanh vị trí cân bằng (\textit{Hình 39}). Nghiên cứu trò chơi này, người ta thấy khoảng cách $h$ (m) từ vị trí người chơi đu đến vị trí cân bằng được biểu diễn qua thời gian $t\,(s)$ (với $t\geq 0$) vởi hệ thức $h=|d|$ với $d=3\cos \left[\dfrac{\pi}{3}(2t-1)\right]$, trong đó ta quy ước $d>0$ khi vị trí cân bằng ở phía sau lưng người chơi đu và $d<0$ trong trường hợp ngược lại (\textit{Nguồn: Đại số và giải tích 11 Nâng cao, NXBGD Việt Nam, 2022)}. Vào thời điểm $t$ nào thì khoảng cách $h$ là $3$ m; $0$ m?
	}{
	\begin{tikzpicture}[line join=round, line cap=round,scale=1,transform shape]
		\definecolor{darkpastelgreen}{rgb}{0.01, 0.75, 0.24}
		\definecolor{almond}{rgb}{0.94, 0.87, 0.8}
		\definecolor{arsenic}{rgb}{0.23, 0.27, 0.29}
		\definecolor{eggplant}{rgb}{0.38, 0.25, 0.32}
		\definecolor{crimson}{rgb}{0.86, 0.08, 0.24}
		\clip (-2.5,-2.5) rectangle (4.5,4.2);
		\tikzset{q_h/.pic={
				% Tre
				\draw[fill=darkpastelgreen] (-.05,4)--(.9,-1.9)
				..controls +(30:.05) and +(-30:.05) ..(1,-1.75)--(.05,4);
				\draw[fill=darkpastelgreen] (-.1,4)--(.6,-1.4)
				..controls +(30:.05) and +(-30:.05) ..(.5,-1.4)--(-.2,4)
				;
				%----------------------
				\def\M{ 
					(-.2,1.76)
					..controls +(-80:.1) and +(80:.1) ..(-.18,1.6)
					..controls +(-40:.1) and +(30:.06) ..(-.18,1.5)--(-.17,1.46)--(-.22,1.44)--(-.19,1.43)--(-.2,1.43)
					..controls +(-80:.1) and +(-12:.05) ..(-.3,1.35)--(-.35,1.28)--(-.55,1.35)
					..controls +(60:.1) and +(-90:.05) ..(-.45,1.5)
					..controls +(100:.1) and +(140:.05) ..(-.38,1.55)
					..controls +(70:.1) and +(-140:.05) ..cycle
					;
				}
				
				\draw[black]\M;
				\fill[almond] \M;
				\draw (-.2,1.76)
				..controls +(-140:.1) and +(50:.05) ..(-.3,1.6)
				(-.2,1.76)
				..controls +(-120:.1) and +(40:.05) ..(-.25,1.6)
				(-.28,1.578)
				..controls +(-20:.02) and +(-150:.02) ..(-.22,1.58)
				(-.25,1.568)--(-.23,1.567);
				%------------------Áo
				\def\A3{ 
					(-.1,.9)
					..controls +(-30:.05) and +(110:0) ..(.18,1)
					--(.19,.95)
					..controls +(-140:.05) and +(30:.05) ..(-.1,.7)--cycle
					;
				}
				\fill[arsenic] \A3;
				\draw[black]\A3;
				\def\A2{ 
					(-.35,1.28)
					..controls +(-30:.05) and +(110:.05) ..(-.25,1.18)
					..controls +(-70:.05) and +(90:.02) ..(-.05,.9)
					..controls +(-90:.02) and +(50:.05) ..(-.05,.4)
					..controls +(-130:.02) and +(50:.05) ..(.6,-.9)
					..controls +(-130:.02) and +(50:.05) ..(.9,-1.9)
					..controls +(-130:.2) and +(30:.3) ..(-.5,-2.3)
					..controls +(130:.2) and +(-90:.1) ..(-.2,-1.5)
					..controls +(130:.2) and +(-140:.1) ..cycle
					;
				}
				\fill[arsenic] \A2;
				\draw[black]\A2;
				\def\A{ 
					(-.55,1.35)
					..controls +(-150:.2) and +(120:.1) ..(-.45,0)
					..controls +(-60:.1) and +(50:.1) ..(-.3,-1.6)
					..controls +(-10:.1) and +(-120:.1) ..(-.15,-1.5)
					..controls +(60:.1) and +(-120:.1) ..(0,-1.45)
					..controls +(120:.3) and +(-50:.1) ..(-.35,.77)
					..controls +(-50:.1) and +(130:.1) ..(-.15,.48)
					..controls +(-50:.1) and +(-130:.1) ..(.43,.93)
					..controls +(-130:.05) and +(-40:.05) ..(.3,1.05)
					..controls +(-130:.05) and +(40:.05) ..(-.05,.75)
					..controls +(-140:.05) and +(-40:.05) ..(-.3,1.15)
					..controls +(100:.05) and +(-40:.05) ..(-.35,1.28)--cycle
					;
				}
				\fill[eggplant] \A;
				\draw[black]\A;
				%-------------------------------
				%nón
				\def\N{ 
					(-.14,1.95)
					..controls +(-80:.1) and +(80:.1) ..(-.17,1.7)
					..controls +(110:.02) and +(-40:.02) ..(-.2,1.76)
					..controls +(-110:.1) and +(50:.4) ..(-.48,1.45)--(-.45,1.35)
					..controls +(-150:.05) and +(20:.05) ..(-.7,1.32)
					..controls +(50:.05) and +(-120:.05) ..(-.65,1.45)
					..controls +(160:.02) and +(-50:.02) ..(-.7,1.48)
					..controls +(160:.02) and +(-50:.02) ..(-.8,1.6)--(-.75,1.63)
					..controls +(-20:.06) and +(130:.04) ..(-.6,1.47)
					..controls +(130:.2) and +(-90:.1) ..(-.7,1.78)
					..controls +(90:.1) and +(170:.1) ..(-.4,1.9)
					..controls +(20:.1) and +(160:.1) ..cycle
					;
				}
				\fill[arsenic] \N;
				\draw[black]\N;
				\draw (-.14,1.95)
				..controls +(-100:.05) and +(40:.06) ..(-.2,1.76)
				(-.5,1.55)
				..controls +(30:.1) and +(-100:.06) ..(-.3,1.76)
				(-.6,1.7)
				..controls +(-55:.1) and +(120:.06) ..(-.55,1.5)
				;
				
				\def\T1{ 
					(.3,1.03)
					..controls +(60:.05) and +(-120:.05) ..(.35,1.15)
					..controls +(60:.05) and +(-140:.05) ..(.45,1.26)
					..controls +(30:.05) and +(60:.04) ..(.5,1.22)
					..controls +(30:.03) and +(50:.03) ..(.54,1.18)
					..controls +(30:.03) and +(50:.03) ..(.56,1.12)
					..controls +(30:.03) and +(50:.03) ..(.58,1.07)
					..controls +(160:.03) and +(40:.03) ..(.49,1)
					..controls +(-140:.03) and +(40:.03) ..(.41,.92)--cycle
					;
				}
				\fill[almond] \T1;
				\draw[black]\T1;
				
				\fill[crimson](-.27,.4)
				..controls +(-80:.6) and +(-80:.1) ..(-.5,-1.9)--(-.45,-1.85)
				..controls +(90:.1) and +(-80:1) ..(-.28,.48)--cycle;
				
				\fill[crimson](-.27,.4)
				..controls +(-80:.6) and +(-80:.1) ..(-.2,-1.9)--(-.15,-1.85)
				..controls +(90:.1) and +(-80:1) ..(-.28,.48)--cycle;
				
		}}
		\tikzset{no/.pic={
				\draw[crimson] (1.8,-1)
				..controls +(150:.4) and +(-120:.5) ..(1.7,-1)
				..controls +(30:.4) and +(-40:.3) ..(1.75,-1)
				..controls +(-30:.2) and +(40:.2) ..(1.7,-1.4)
				(1.75,-1)
				..controls +(-50:.1) and +(20:.1) ..(1.7,-1.2)
				
				;
		}}
		\draw[thick] (4,-1.25) arc(-60:-120:6cm);
		\draw[dashed,thick] (4,-1.25)--(1.3,4)--(-2,-1.25) (1.3,4)--(1.3,-2);
		\node at (1.5,.5) [rotate=90]{\tiny Vị trí cân bằng};
		\draw[thick,<->] (.35,-2)--(1.3,-2);
		\node at (.85,-2) [above]{\tiny $h$};
		\fill (4,-1.25) circle (0.05);
		\fill (-2,-1.25) circle (0.05);
		\fill (1.3,4) circle (0.05);
		\path
		(0,0)pic[scale=1,rotate=-20]{q_h}(-.88,1.48)pic[scale=1, scale=.7,rotate=-20]{no};
		
		
	\end{tikzpicture}
	}
	\loigiai{
		Khoảng cách $h$ là $3$ m khi 
		\[3\cos \left[\dfrac{\pi}{3}(2t-1)\right] = - 3 \Leftrightarrow  \cos \left[\dfrac{\pi}{3}(2t-1)\right] = -1 \Leftrightarrow \dfrac{\pi}{3}(2t-1) =-\pi +k2\pi \Leftrightarrow t=-1+3k,\, k\in \mathbb{Z}.\]
		Vậy vào thời điểm $t=-1+3k,\, k\in \mathbb{Z}$ nào thì khoảng cách $h$ là $3$ m.\\
		Khoảng cách $h$ là $0$ m khi 
		\[3\cos \left[\dfrac{\pi}{3}(2t-1)\right] = 0 \Leftrightarrow  \cos \left[\dfrac{\pi}{3}(2t-1)\right] = 0 \Leftrightarrow \dfrac{\pi}{3}(2t-1) =\dfrac{\pi}{2} +k\pi \Leftrightarrow t=\dfrac{5}{4}+\dfrac{3}{2}k,\, k\in \mathbb{Z}.\]
		Vậy vào thời điểm $t=\dfrac{5}{4}+\dfrac{3}{2}k,\, k\in \mathbb{Z}$ nào thì khoảng cách $h$ là $0$ m.
	}
\end{ex}

% #1708 | Archive/2024/SP-K10-K11-LyVanHoang_TVN-023/11. C1-LuongGiac-lvhoang_TVN-023/dataC1K11/Chuong1-Bai4-Dang6-Toanthucte-lienmon.tex | mã gốc 1D1H5-6
\begin{ex}%[1D1H5-6]%[Anh Duy]
	Một cây cầu có dạng cung $A B$ của đồ thị hàm số $y=4{,}2 \cdot \cos \dfrac{x}{8}$ và được mô tả trong hệ trục toạ độ với đơn vị trục là mét như ở hình sau. Một sà lan chở khối hàng hoá được xếp thành hình hộp chữ nhật với độ cao $3 \mathrm{~m}$ so với mực nước sông sao cho sà lan có thể đi qua được gầm cầu.
	\choiceTF
	{\True Giá trị lớn nhất của hàm số là $4{,}2$}
	{Hàm số đã cho là hàm số lẻ}
	{\True $4{,}2 \cdot \cos \dfrac{x}{8}=3 \Leftrightarrow \dfrac{x}{8} \approx \pm 0{,}775$}
	{\True Chiều rộng của khối hàng hoá đó phải nhỏ hơn $12{,}5 \mathrm{~m}$}
	\begin{center}
		\begin{tikzpicture}[>=stealth,x=.75cm,y=.75cm,scale=0.4]
			%		\draw[color=gray,dash pattern=on 1pt off 1pt,xstep=1.0cm,ystep=1.0cm] (-5.2,-5.2) grid (5.2,5.2);
			\draw[->] (-13.5,0) -- (14.5,0)node[below]{\scriptsize $x$};
			\draw[->] (0,-1) -- (0,5) node[left] {\scriptsize $y$};
			\draw (0,0)node[below left]{\scriptsize $O$};
			\draw (-12.6,0) node[below]{$-4\pi$};
			\draw (12.6,0) node[below]{$4\pi$};
			\clip (-13,-1)rectangle(13,5);
			\draw[thick,samples=150,smooth,domain=-12.6:12.6] plot(\x,{4.2*cos(\x/8 r)});
		\end{tikzpicture}
	\end{center}
	\loigiai{
		\begin{itemchoice}
			\itemch Đúng. Giá trị lớn nhất của hàm số là $4{,}2$.
			\itemch Sai. Tập xác định $\mathscr{D}=\mathbb{R}$. Khi đó $\forall x\in \mathscr{D}$ ta có $-x\in \mathscr{D}$ và 
			$$y(-x)=4{,}2\cos \dfrac{-x}{8}=4{,}2\cos \dfrac{x}{8}=y(x).$$
			nên hàm số đã cho là hàm số chẵn.
			\itemch Đúng. Với mỗi điểm $M$ nằm trên mặt cầu, khoảng cách từ điểm $M$ đến mặt nước tương ứng với giá trị tung độ $y$ của điểm $M$.\\
			Xét phương trình: $4{,}2 \cdot \cos \dfrac{x}{8}=3 \Leftrightarrow \cos \dfrac{x}{8}=\dfrac{5}{7}$. Do $x \in[-4 \pi ; 4 \pi]$ nên $\dfrac{x}{8} \in\left[-\dfrac{\pi}{2};\dfrac{\pi}{2}\right]$.\\
			Khi đó, ta có $\cos \dfrac{x}{8}=\dfrac{5}{7} \Leftrightarrow \dfrac{x}{8} \approx \pm 0{,}775$ suy ra $\left|\dfrac{x}{8}\right|<0{,}78 \Leftrightarrow|x|<6{,}24$.
			\itemch Đúng. Do sà lan có thể đi qua được gầm cầu nên chiều rộng của khối hàng hoá là
			$$
			2|x|<12{,}48<12{,}5(\mathrm{~m}).
			$$
		\end{itemchoice}
	}
\end{ex}

% #1709 | Archive/2024/SP-K10-K11-LyVanHoang_TVN-023/11. C1-LuongGiac-lvhoang_TVN-023/dataC1K11/Chuong1-Bai4-Dang6-Toanthucte-lienmon.tex | mã gốc 1D1H5-6
\begin{ex}[Cánh Diều]%[1D1H5-6]%[Anh Duy]
	Một vệ tinh nhân tạo bay quanh Trái Đất theo một quỹ đạo là đường elip (\textit{Hình 33}). Độ cao $h$ (km) của vệ tinh so với Trái Đất được xác định bởi công thức $h=550+450\cos \dfrac{\pi}{50}t$ (\textit{Nguồn: Đại số và Giải tích 11 Nâng cao, NXBGD Việt Nam, 2021}), trong đó $t$ là thời gian tính bằng phút kể từ lúc vệ tinh bay vào quỹ đạo. Tại thời điểm $t$ bằng bao nhiêu thì vệ tinh cách mặt đất $1\,000$ km; $250$ km; $100$ km?	
	\choiceTF[t]
	{Với $t=250$ thì độ cao của vệ tinh so với Trái Đất là $1000$ km}
	{\True Tại các thời điểm $t=\pm36{,}61+100k,\,(k\in\mathbb{Z})$ thì vệ tính cách mặt đất $250$ km}
	{\True Tại các thời điểm $t=50+100k, \,(k\in\mathbb{Z})$ thì vệ tính cách mặt đất $100$ km}
	{Trong khoảng thời gian $60$ phút đầu tiên thì có 2 lần vệ tính cách mặt đất $100$ km}
	\loigiai{
		\begin{itemchoice}
			\itemch Đúng. Với $t=20$ thì độ cao của vệ tinh so với Trái Đất là $550$ km.
			\itemch Đúng. Với $h=250$ km ta có $$250=550+450\cos\dfrac{\pi}{50}t\Leftrightarrow\cos\dfrac{\pi}{50}t=-\dfrac{2}{3}\Leftrightarrow\dfrac{\pi}{50}t \Leftrightarrow t=\pm36{,}61+100k,\,(k\in\mathbb{Z}).$$
			\itemch Đúng. Với $h=100$ km ta có $$100=550+450\cos\dfrac{\pi}{50}t\Leftrightarrow\cos\dfrac{\pi}{50}t=-1\Leftrightarrow\dfrac{\pi}{50}t=\pi+k2\pi\Leftrightarrow t=50+100k,\,(k\in\mathbb{Z}).$$
			\itemch Ta có $0<50+100k\le 60\Leftrightarrow -0{,}5<k\le 0{,}6\Rightarrow k=0$.\\
			Vậy trong khoảng thời gian $60$ phút đầu tiên thì có 1 lần vệ tính cách mặt đất $100$ km.
		\end{itemchoice}
	}
\end{ex}

% #1710 | Archive/2024/SP-K10-K11-LyVanHoang_TVN-023/11. C1-LuongGiac-lvhoang_TVN-023/dataC1K11/Chuong1-Bai4-Dang6-Toanthucte-lienmon.tex | mã gốc 1D1H5-6
\begin{ex}%[1D1H5-6]%[Anh Duy]
	Một quả đạn pháo được bắn ra khỏi nòng pháo với vận tốc ban đầu $v_0=500 ~\mathrm{m} / \mathrm{s}$ hợp với phương ngang một góc $\alpha$. Trong Vật lí, ta biết rằng, nếu bỏ qua sức cản của không khí và coi quả đạn pháo được bắn ra từ mặt đất thì quỹ đạo của quả đạn tuân theo phương trình $y=\dfrac{-g}{2 v_0^2 \cos ^2 \alpha} x^2+x \tan \alpha$, ở đó $g=9{,}8 \mathrm{~m} / \mathrm{s}^2$ là gia tốc trọng trường.
	\choiceTF[t]
	{Quả đạn chạm đất khi $\dfrac{-g}{2 v_0^2 \cos ^2 \alpha} x^2+x \tan \alpha=\pm 1$}
	{\True Tính theo góc bắn $\alpha$ tầm xa mà quả đạn đạt tới (tức là khoảng cách từ vị trí bắn đến điểm quả đạn chạm đất) là $x=\dfrac{2 v_0^2 \cos ^2 \alpha.\tan \alpha}{g}$}
	{Góc bắn $\alpha$ để quả đạn trúng mục tiêu cách vị trí đặt khẩu pháo $22000 \mathrm{~m}$ là $40^\circ$}
	{\True Góc bắn $\alpha=45^\circ$ thì quả đạn bay xa nhất}
	\loigiai{
		\begin{itemchoice}
			\itemch Quả đạn chạm đất $\Rightarrow y=0\Leftrightarrow \dfrac{-g}{2 v_0^2 \cos ^2 \alpha} x^2+x \tan \alpha=0$.
			\itemch Ta có \begin{center}
				$\dfrac{-g}{2 v_0^2 \cos ^2 \alpha} x^2+x \tan \alpha=0 \Leftrightarrow x\left(\dfrac{-g}{2 v_0^2 \cos ^2 \alpha} x+ \tan \alpha\right)=0\Rightarrow \hoac{&x=0\\&x=\dfrac{2 v_0^2 \cos ^2 \alpha.\tan \alpha}{g}}.$
			\end{center}
			Vậy tầm xa của quả đạn tính theo góc $\alpha$ là $x=\dfrac{2 v_0^2 \cos ^2 \alpha.\tan \alpha}{g}$.
			\itemch Sai. Để quả đạn trúng mục tiêu cách vị trí đặt khẩu pháo $22000 \mathrm{~m}$ thì
			\begin{eqnarray*}
				&&\dfrac{2 v_0^2 \cos ^2 \alpha.\tan \alpha}{g}=22000 \\
				&\Leftrightarrow &	\dfrac{2. 500^2. \cos ^2 \alpha.\tan \alpha}{9,8}=22000 \\
				&\Leftrightarrow &\cos ^2 \alpha.\tan \alpha=0,4312.\\
				&\Leftrightarrow &\dfrac{1}{1+\tan^2\alpha}.\tan \alpha=0,4312\\
				&\Leftrightarrow &0,4312\tan^2 \alpha-\tan\alpha+0,4312=0.\\
				&\Leftrightarrow &\hoac{&\tan\alpha=0{,}57\\&\tan\alpha=1{,}75}\\
				&\Leftrightarrow &\hoac{&\alpha\approx 30^\circ +k180^\circ \\&\alpha\approx 60^\circ +k180^\circ.}
			\end{eqnarray*} Do góc $0<\alpha<90^\circ$ nên ta chọn $\alpha\approx 30^\circ$ hoặc $\alpha\approx 60^\circ$.
			\itemch Tầm  xa của quả đạn  là $x=\dfrac{2 v_0^2 \cos ^2 \alpha.\tan \alpha}{g}$, do đó quả đạn bay xa nhất khi $ \cos ^2 \alpha.\tan \alpha$ lớn nhất.\\
			Ta có $\cos ^2 \alpha.\tan \alpha =\dfrac{\tan \alpha}{1+\tan^2\alpha}$.\\ Mặt khác $1+\tan^2\alpha \ge 2\tan\alpha$  nên ta có $\dfrac{\tan \alpha}{1+\tan^2\alpha} \le \dfrac{1}{2}$.\\  Dấu bằng xảy ra khi $\tan\alpha =1 \Leftrightarrow \alpha=45^\circ $ (do $0<\alpha<90^\circ$). \\
			Vậy khi góc bắn $\alpha=45^\circ$ thì quả đạn bay xa nhất.
		\end{itemchoice}
	}
\end{ex}

% #1713 | Archive/2024/SP-K10-K11-LyVanHoang_TVN-023/11. C1-LuongGiac-lvhoang_TVN-023/dataC1K11/Chuong1-Bai4-Dang6-Toanthucte-lienmon.tex | mã gốc 1D1H5-6
\begin{ex}%[1D1H5-6]%[Anh Duy]
	Số giờ có ánh sáng mặt trời của một thành phố $A$ ở vĩ độ $40^\circ$ Bắc trong ngày thứ $t$ của một năm không nhuận được cho bởi hàm số
	\[
	d(t)=3\sin \left[\dfrac{\pi}{182}(t-80)\right] +12 \text{ với } t\in \mathbb{Z} \text{ và } 0<t\leq 365.
	\]
	Vào ngày nào trong năm thì thành phố $A$ có đúng $15$ giờ có ánh sáng mặt trời?
	\shortans[]{$171$}
	\loigiai{
		Ta giải phương trình 
		\begin{eqnarray*}
			d(t)=15 &\Leftrightarrow& 3\sin\left[\dfrac{\pi}{182}(t-80)\right] +12=15 \\
			&\Leftrightarrow& \sin\left[\dfrac{\pi}{182}(t-80)\right]=1\\
			&\Leftrightarrow& \dfrac{\pi}{182}(t-80) =\dfrac{\pi}{2}+k2 \pi \\
			&\Leftrightarrow& t-80 = 91+364k\\
			&\Leftrightarrow& t= 364k+171 (k\in \mathbb{Z}).	
		\end{eqnarray*}
		Ta lại có  \[0<364k+171 \le 365 \Leftrightarrow -\dfrac{171}{364}<k \le \dfrac{196}{364} \Leftrightarrow k=0.\]
		Vậy thành phố $A$ có đúng $15$ giờ có ánh sáng mặt trời vào ngày thứ $171$ trong năm.}
\end{ex}

% #1714 | Archive/2024/SP-K10-K11-LyVanHoang_TVN-023/11. C1-LuongGiac-lvhoang_TVN-023/dataC1K11/Chuong1-Bai4-Dang6-Toanthucte-lienmon.tex | mã gốc 1D1H5-6
\begin{ex}%[1D1H5-6]%[Anh Duy]
	Huyết áp là áp lực cần thiết tác động lên thành của động mạch để đưa máu từ tim đến nuôi dưỡng các mô trong cơ thể. Huyết áp được tạo ra do lực co bóp của cơ tim và sức cản của thành động mạch. Mỗi lần tim đập, huyết áp của chúng ta tăng rồi giảm giữa các nhịp. Huyết áp tối đa và huyết áp tối thiểu gọi là huyết áp tâm thu và tâm trương, tương ứng. Chỉ số huyết áp của chúng ta được viết là tâm thu/tâm trương. Chỉ số huyết áp $120 / 80 $là bình thường. Giả sử một người nào đó có nhịp tim là $70$ lần trên phút và huyết áp của người đó được mô hình hoá bởi hàm số $P(t)=100+20\sin \left(\dfrac{7\pi}{3}t\right)$, trong đó $P(t)$ là huyết áp tính theo đơn vị mmHg (milimét thuỷ ngân) và thời gian tính theo giây. Trong khoảng từ $0$ đến $1$ giây, hãy xác định số lần huyết áp là $100$ mmHg.
	\shortans[]{$2$}
	\loigiai{
		Ta có $100+20\sin \left(\dfrac{7\pi}{3}t\right)=100\Leftrightarrow \sin \left(\dfrac{7\pi}{3}t\right)=0 \Leftrightarrow \dfrac{7\pi}{3}t =k\pi \Leftrightarrow t=\dfrac{3k}{7}$.\\
		Khi đó $0< \dfrac{3k}{7}\le 1\Rightarrow 0< k\le \dfrac{7}{3}\Rightarrow k\in \{1;2\}$.\\
		Vậy trong khoảng từ $0$ đến $1$ giây, có 2 lần lần huyết áp là $100$ mmHg.
	}
\end{ex}

% #1716 | Archive/2024/SP-K10-K11-LyVanHoang_TVN-023/11. C1-LuongGiac-lvhoang_TVN-023/dataC1K11/Bai-1-Dang-7-GTLG-GocLienQuanDacBiet.tex | mã gốc 1D1N2-4
\begin{vd}[Cánh Diều, MĐ1]%[1D1N2-4]
	Tính 
	\begin{listEX}[2]
		\item $\sin \dfrac{13\pi}{4}$;
		\item $\sin \dfrac{\pi}{10}-\cos \dfrac{2\pi}{5}$.
	\end{listEX}
	\loigiai{
		Ta có 
		\begin{enumerate}
			\item $\sin \dfrac{13\pi}{4}=\sin \left(3\pi+\dfrac{\pi}{4}\right)=\sin \left(\pi+\dfrac{\pi}{4}\right)=-\sin \dfrac{\pi}{4}=-\dfrac{\sqrt{2}}{2}$;
			\item $\sin \dfrac{\pi}{10}-\cos \dfrac{2\pi}{5}=\sin \dfrac{\pi}{10}-\sin \left(\dfrac{\pi}{2}-\dfrac{2\pi}{5}\right)=\sin \dfrac{\pi}{10}-\sin \dfrac{\pi}{10}=0$.
		\end{enumerate}
	}
\end{vd}

% #1717 | Archive/2024/SP-K10-K11-LyVanHoang_TVN-023/11. C1-LuongGiac-lvhoang_TVN-023/dataC1K11/Bai-1-Dang-7-GTLG-GocLienQuanDacBiet.tex | mã gốc 1D1H2-4
\begin{vd}[Cánh Diều,MĐ2]%[1D1H2-4]
	Tính 
	\begin{listEX}[2]
		\item $\cos^2 \dfrac{\pi}{8}+\cos^2 \dfrac{3\pi}{8}$.
		\item $\tan 1^\circ\cdot \tan 2^\circ\cdot \tan 45^\circ\cdot \tan 88^\circ\cdot \tan 89^\circ$.
	\end{listEX}
	\loigiai{
		\begin{enumerate}
			\item Ta có 
			\[\cos^2 \dfrac{\pi}{8}+\cos^2 \dfrac{3\pi}{8}
			=\cos^2 \dfrac{\pi}{8}+\cos^2 \left(\dfrac{\pi}{2}-\dfrac{\pi}{8}\right)
			=\cos^2 \dfrac{\pi}{8}+\sin^2 \dfrac{\pi}{8}=1.\]
			\item Ta có
			\begin{eqnarray*}
				& &\tan 1^\circ\cdot \tan 2^\circ\cdot \tan 45^\circ\cdot \tan 88^\circ\cdot \tan 89^\circ\\
				&=&\tan 1^\circ\cdot \tan 2^\circ\cdot \tan 45^\circ\cdot \tan (90^\circ-2^\circ)\cdot \tan (90^\circ-1^\circ)\\
				&=&\tan 1^\circ\cdot \tan 2^\circ\cdot \tan 45^\circ\cdot \cot 2^\circ\cdot \cot 1^\circ\\
				&=&\tan 1^\circ\cdot \cot 1^\circ\cdot \tan 2^\circ\cdot \cot 2^\circ\cdot \tan 45^\circ\\
				&=&\tan 45^\circ=1.
			\end{eqnarray*}
		\end{enumerate}
	}
\end{vd}

% #1718 | Archive/2024/SP-K10-K11-LyVanHoang_TVN-023/11. C1-LuongGiac-lvhoang_TVN-023/dataC1K11/Bai-1-Dang-7-GTLG-GocLienQuanDacBiet.tex | mã gốc 1D1V2-4
\begin{vd}[MĐ3]%[1D1V2-4]
	Tính
	\begin{enumerate}
		\item $A=\sin ^2 5^{\circ}+\sin ^2 10^{\circ}+\sin ^2 15^{\circ}+\ldots+\sin ^2 85^{\circ}$ ($17$ số hạng).
		\item $B=\cos 5^{\circ}+\cos 10^{\circ}+\cos 15^{\circ}+\ldots+\cos 175^{\circ}$ ($35$ số hạng).
	\end{enumerate}
	\loigiai{
		\begin{enumerate}
			\item Vận dụng $\sin \left(90^{\circ}-\alpha\right)=\cos \alpha$.\\
			Ta có $\sin 85^{\circ}=\cos 5^{\circ}$,  $\sin 80^{\circ}=\cos 10^{\circ}$, $\sin 75^{\circ}=\cos 15^{\circ}, \ldots, \sin 50^{\circ}=\cos 40^{\circ}$.
			\begin{eqnarray*}
				A & =&\left(\sin ^2 5^{\circ}+\cos ^2 5^{\circ}\right)+\left(\sin ^2 10^{\circ}+\cos ^2 10^{\circ}\right)+\ldots+\left(\sin ^2 40^{\circ}+\cos ^2 40^{\circ}\right)+\sin ^2 45^{\circ} \\
				& =&\underbrace{1+1+\ldots+1}_{\text {8 số 1 }}+\left(\dfrac{\sqrt{2}}{2}\right)^2=\dfrac{17}{2} .
			\end{eqnarray*}
			\item Vận dụng $\cos \left(180^{\circ}-\alpha\right)=-\cos \alpha$.\\
			Ta có $\cos 175^{\circ}=-\cos 5^{\circ}, \cos 170^{\circ}=-\cos 10^{\circ}$, $\cos 165^{\circ}=-\cos 15^{\circ}, \ldots, \cos 95^{\circ}=-\cos 85^{\circ}$.
			\begin{eqnarray*}
				B&=&\left(\cos 5^{\circ}+\cos 175^{\circ}\right)+\left(\cos 10^{\circ}+\cos 170^{\circ}\right)+\ldots+\left(\cos 85^{\circ}+\cos 95^{\circ}\right)+\cos 90^{\circ}\\
				&=&\underbrace{0+0+\ldots+0}_{17 \text { số } 0}+\cos 90^{\circ}=0.
			\end{eqnarray*}
		\end{enumerate} 
	}
\end{vd}

% #1719 | Archive/2024/SP-K10-K11-LyVanHoang_TVN-023/11. C1-LuongGiac-lvhoang_TVN-023/dataC1K11/Bai-1-Dang-7-GTLG-GocLienQuanDacBiet.tex | mã gốc 1D1H2-4
\begin{vd}[MĐ2]%[1D1H2-4]
	Cho $\sin\alpha=-\dfrac{3}{5}$ $\left(\dfrac{3\pi}{2}<\alpha<2\pi\right)$. Tính $\cos\alpha$; $\tan\left(\dfrac{5\pi}{2}+\alpha\right)$.
	\loigiai{
		Do $\alpha\in \left(\dfrac{3\pi}{2};2\pi\right)$ nên $\cos\alpha>0$.\\
		Suy ra $\cos\alpha=\sqrt{1-\sin^2 \alpha}=\sqrt{1-\left(-\dfrac{3}{5}\right)^2}=\dfrac{4}{5}$.\\
		Ta có $\tan\left(\dfrac{5\pi}{2}+\alpha\right)=\tan\left(\dfrac{\pi}{2}+\alpha\right)=-\cot\alpha=-\dfrac{\cos\alpha}{\sin\alpha}=\dfrac{4}{3}$.
	}
\end{vd}

% #1722 | Archive/2024/SP-K10-K11-LyVanHoang_TVN-023/11. C1-LuongGiac-lvhoang_TVN-023/dataC1K11/Bai-1-Dang-7-GTLG-GocLienQuanDacBiet.tex | mã gốc 1D1H2-4
\begin{bt}%[1D1H2-4]
	Tính giá trị của biểu thức $A=\cos ^2 \dfrac{\pi}{8}+\cos ^2 \dfrac{3 \pi}{8}+\cos ^2 \dfrac{5 \pi}{8}+\cos ^2 \dfrac{7 \pi}{8}$.
	\loigiai{
		Vì $\cos ^2 \dfrac{\pi}{8}=\sin^2\dfrac{7\pi}{8}$ và $\cos ^2 \dfrac{3\pi}{8}=\sin^2\dfrac{5\pi}{8}$ nên ta có
		$A=\sin ^2 \dfrac{7\pi}{8}+\cos ^2 \dfrac{7 \pi}{8}+\sin ^2 \dfrac{5 \pi}{8}+\cos ^2 \dfrac{5 \pi}{8}=2$.
		
	}
\end{bt}

% #1723 | Archive/2024/SP-K10-K11-LyVanHoang_TVN-023/11. C1-LuongGiac-lvhoang_TVN-023/dataC1K11/Bai-1-Dang-7-GTLG-GocLienQuanDacBiet.tex | mã gốc 1D1H2-4
\begin{bt}%[1D1H2-4]
	Tính giá trị biểu thức $A=\cos 10^\circ+\cos 20^\circ+\ldots +\cos 90^\circ$.
	\loigiai{
		Ta có $\cos 10^\circ=-\cos 80^\circ$,..., $\cos 40^\circ=-\cos 50^\circ$. \\
		Vậy $A=\left(\cos 10^\circ+\cos 80^\circ\right)+\ldots +\left(\cos 40^\circ+\cos 50^\circ\right)+\cos 90^\circ=0$.
	}	
\end{bt}

% #1724 | Archive/2024/SP-K10-K11-LyVanHoang_TVN-023/11. C1-LuongGiac-lvhoang_TVN-023/dataC1K11/Bai-1-Dang-7-GTLG-GocLienQuanDacBiet.tex | mã gốc 1D1H2-4
\begin{bt}%[1D1H2-4]
	Tính giá trị biểu thức $B=\cot 1^\circ \cdot \cot 2^\circ \ldots \cot 88^\circ \cdot \cot 89^\circ$.
	\loigiai{
		Ta có $\cot 89^\circ =\tan 1^\circ$,..., $\cot 46^\circ =\tan 44^\circ$. \\
		Vậy $B=\left(\tan 1^\circ\cdot \cot 1^\circ\right)\ldots \left(\tan 44^\circ\cdot \cot 44^\circ\right)\cdot \cot 45^\circ=1$.
	}
\end{bt}

% #1725 | Archive/2024/SP-K10-K11-LyVanHoang_TVN-023/11. C1-LuongGiac-lvhoang_TVN-023/dataC1K11/Bai-1-Dang-7-GTLG-GocLienQuanDacBiet.tex | mã gốc 1D1V2-4
\begin{bt}%[1D1V2-4]
	Cho $\sin x =\dfrac{1}{3}$ với $\dfrac{\pi}{2} < x < \pi $. Tính giá trị của $\cos x$, $\cos (5\pi -x)$, $\cot \left(\dfrac{3\pi}{2} +x\right)$.
	\loigiai{
		Do $\dfrac{\pi}{2} < x < \pi $ nên $\cos x =- \sqrt{1-\sin^2 x} =-\sqrt{1-\dfrac{1}{9}}=-\dfrac{2\sqrt{2}}{3}$.\\
		$\sin (5\pi -x) =\sin (\pi -x) =\sin x =\dfrac{1}{3}$.\\
		$\cot \left(\dfrac{3\pi}{2} +x\right) = \cot \left(\pi + \dfrac{\pi}{2} +x\right) =\cot \left(\dfrac{\pi}{2} +x\right) =-\tan x =-\dfrac{\sin x}{\cos x} =\dfrac{1}{2\sqrt{2}}$.
	}
\end{bt}

% #1726 | Archive/2024/SP-K10-K11-LyVanHoang_TVN-023/11. C1-LuongGiac-lvhoang_TVN-023/dataC1K11/Bai-1-Dang-7-GTLG-GocLienQuanDacBiet.tex | mã gốc 1D1V2-4
\begin{bt}%[1D1V2-4]
	Rút gọn của biểu thức $A = \sin(\pi + x)+ \cos \left(\dfrac{\pi}{2} - x\right) + \cot(2\pi-x)+ \tan \left(\dfrac{3\pi }{2} + x\right)$.
	\loigiai{
		Ta có
		\allowdisplaybreaks
		\begin{eqnarray*}
			A&=&\sin(\pi + x)+ \cos \left(\dfrac{\pi}{2} - x\right) + \cot(2\pi-x)+ \tan \left(\dfrac{3\pi }{2} + x\right)\\
			&=& - \sin x + \sin x + \cot(-x)+ \tan \left(-\dfrac{\pi}{2} + x\right)\\
			&=& - \cot x - \tan \left(\dfrac{\pi}{2} - x\right)\\
			&=&- \cot x - \cot x\\
			&=&-2\cot x.
		\end{eqnarray*}		
	}
\end{bt}

% #1727 | Archive/2024/SP-K10-K11-LyVanHoang_TVN-023/11. C1-LuongGiac-lvhoang_TVN-023/dataC1K11/Bai-1-Dang-7-GTLG-GocLienQuanDacBiet.tex | mã gốc 1D1V2-4
\begin{bt}%[1D1V2-4]
	Cho $\sin\alpha=m$. Tính 
	\begin{center}
		$P=\cos\left( \dfrac{\pi}{2}-\alpha \right) \sin(\pi-\alpha)-\sin\left( \dfrac{\pi}{2}-\alpha \right) \cos(\pi-\alpha)+\sin^2(\alpha+2018\pi).$
	\end{center}
	\loigiai{
		Ta có $P=\sin\alpha\cdot\sin\alpha-\cos\alpha\cdot(-\cos\alpha)+\sin^2\alpha=\sin^2\alpha+\cos^2\alpha+\sin^2\alpha=1+\sin^2\alpha=1+m^2$.
	}
\end{bt}

% #1745 | Archive/2024/SP-K10-K11-LyVanHoang_TVN-023/11. C1-LuongGiac-lvhoang_TVN-023/dataC1K11/Chuong1-Bai4-Dang1-Phuongtrinhtuongduong.tex | mã gốc 1D1H5-1
\begin{vd}[Mức độ 1]%[Bài Giảng 10-11 - Phan Anh]%[1D1H5-1]
	Hai phương trình $x-1=0$ và $\dfrac{x^2-1}{x+1}=0$ có tương đương không? Vì sao?
	\loigiai{
		Tập nghiệm của phương trình $x-1=0$ là $S_1=\{1\}$.\\
		Tập nghiệm của phương trình $\dfrac{x^2-1}{x+1}=0$ là $S_2=\{1\}$.\\
		Vì $S_1=S_2$ nên hai phương trình $x-1=0$ và $\dfrac{x^2-1}{x+1}=0$ tương đương.
	}	
\end{vd}

% #1748 | Archive/2024/SP-K10-K11-LyVanHoang_TVN-023/11. C1-LuongGiac-lvhoang_TVN-023/dataC1K11/Chuong1-Bai4-Dang1-Phuongtrinhtuongduong.tex | mã gốc 1D1H5-1
\begin{bt}[Mức độ 2]%[Bài Giảng 10-11 - Phan Anh]%[1D1H5-1]
	Xét sự tương đương của hai phương trình $\dfrac{x-1}{x+1}=0$ và $x^2-1=0$.
	\loigiai{
		Tập nghiệm của phương trình $\dfrac{x-1}{x+1}=0$ là $S_1=\{1\}$.\\
		Tập nghiệm của phương trình $x^2-1=$ là $S_2=\{\pm1\}$.\\
		Vì $S_1\ne S_2$ nên hai phương trình $\dfrac{x-1}{x+1}=0$ và $x^2-1=0$ không tương đương.
	}
\end{bt}

% #1749 | Archive/2024/SP-K10-K11-LyVanHoang_TVN-023/11. C1-LuongGiac-lvhoang_TVN-023/dataC1K11/Chuong1-Bai4-Dang1-Phuongtrinhtuongduong.tex | mã gốc 1D1V5-1
\begin{bt}[Mức độ 3]%[Bài Giảng 10-11 - Phan Anh]%[1D1V5-1]
	Xác định giá trị $m$ để hai phương trình $\dfrac{x^2-4x-4}{\sqrt{x-4}}=\sqrt{x-4}$ và $\dfrac{x^2-4x-m}{\sqrt{x-1}}=0$ tương đương.
	\loigiai{Xét phương trình $\dfrac{x^2-4x-4}{\sqrt{x-4}}=\sqrt{x-4}$.\hfill (1)\\
		Điều kiện $x>4$.\\
		$(1)\Leftrightarrow x^2-4x-4=x-4\Leftrightarrow\hoac{&x=0\ (\text{loại})\\ &x=5\ (\text{nhận})}\Rightarrow S_{1}=\{5\}$.\\
		Hai phương trình tương đương khi $x=5$ là nghiệm của phương trình $\dfrac{x^2-4x-m}{\sqrt{x-1}}=0\Rightarrow m=5$.\\Xét phương trình $\dfrac{x^2-4x-5}{\sqrt{x-1}}=0$.\hfill (2)\\
		Điều kiện $x>1$.\\
		$(2)\Leftrightarrow x^2-4x-5=0\Leftrightarrow\hoac{&x=-1\, (\text{loại})\\ &x=5\, (\text{nhận})}\Rightarrow S_{2}=\{5\}$.\\
		Vì $S_{1}=S_{2}$ nên hai phương trình đã cho tương đương.\\
		Vậy $m=5$ thỏa mãn yêu cầu bài toán.
	}
\end{bt}

% #1757 | Archive/2024/SP-K10-K11-LyVanHoang_TVN-023/11. C1-LuongGiac-lvhoang_TVN-023/dataC1K11/Bai1-Dang6.tex | mã gốc 1D1H2-2
\begin{vd}[MĐ2]%[Dự án BG K10-K11]%[Dao-V- Thuy]%[1D1H2-2][Cánh Diều]
	Cho góc lượng giác $\alpha$ sao cho $-\dfrac{\pi}{2}<\alpha<0$ và $\tan \alpha=-2$. Tính $\cos \alpha$, $\sin \alpha$.
	\loigiai{
		Do $\tan \alpha=-2$ nên $\dfrac{\sin \alpha}{\cos \alpha}=-2$, suy ra $\sin \alpha=-2\cos \alpha$.\\
		Vì $\cos^2\alpha+\sin^2\alpha=1$ nên $\cos^2\alpha+(-2\cos\alpha)^2=1$, suy ra $\cos^2 \alpha=\dfrac{1}{5}$.\\
		Do $-\dfrac{\pi}{2}<\alpha<0$ nên $\cos \alpha>0$.\\
		Từ đó ta có $\cos \alpha=\sqrt{\dfrac{1}{5}}=\dfrac{1}{\sqrt{5}}$, suy ra $\sin \alpha=-2\cdot \dfrac{1}{\sqrt{5}}=-\dfrac{2}{\sqrt{5}}$.
	}
\end{vd}

% #1758 | Archive/2024/SP-K10-K11-LyVanHoang_TVN-023/11. C1-LuongGiac-lvhoang_TVN-023/dataC1K11/Bai1-Dang6.tex | mã gốc 1D1H2-2
\begin{vd}[MĐ2]%[Dự án BG K10-K11]%[Dao-V- Thuy]%[1D1H2-2][Cánh Diều]
	Cho góc lượng giác $\alpha$ sao cho $\pi<\alpha<\dfrac{3\pi}{2}$ và $\sin \alpha=-\dfrac{4}{5}$. Tìm $\cos \alpha$.
	\loigiai{
		Vì $\cos^2\alpha+\sin^2\alpha=1$ nên $\cos^2\alpha=1-\sin^2\alpha=1-\left(-\dfrac{4}{5}\right)^2=\dfrac{9}{25}$.\\
		Do $\pi<\alpha<\dfrac{3\pi}{2}$ nên $\cos \alpha<0$.\\
		Từ đó ta có $\cos \alpha=-\sqrt{\dfrac{9}{25}}=-\dfrac{3}{4}$.
	}
\end{vd}

% #1759 | Archive/2024/SP-K10-K11-LyVanHoang_TVN-023/11. C1-LuongGiac-lvhoang_TVN-023/dataC1K11/Bai1-Dang6.tex | mã gốc 1D1H2-2
\begin{vd}[MĐ2]%[Dự án BG K10-K11]%[Dao-V- Thuy]%[1D1H2-2]
	Cho $\tan \alpha=2$, tính giá trị biểu thức $M=\cos^2\alpha-\sin^2\alpha$.
	\loigiai{
		Ta có $M=\dfrac{\cos^2\alpha-\sin^2\alpha}{\cos^2\alpha+\sin^2\alpha}$.\\
		Chia cả tử và mẫu cho $\cos^2\alpha$ ta được $M=\dfrac{1-\tan^2\alpha}{1+\tan^2\alpha}\Rightarrow M=\dfrac{1-4}{1+4}=-\dfrac{3}{5}$.
	}
\end{vd}

% #1760 | Archive/2024/SP-K10-K11-LyVanHoang_TVN-023/11. C1-LuongGiac-lvhoang_TVN-023/dataC1K11/Bai1-Dang6.tex | mã gốc 1D1V2-2
\begin{vd}[MĐ3]%[Dự án BG K10-K11]%[Dao-V- Thuy]%[1D1V2-2]
	Cho $\cot\alpha =3$. Tính giá trị biểu thức $M=\dfrac{2\sin \alpha-3\cos\alpha}{5\sin^3 \alpha+\cos^3 \alpha}$.
	\loigiai{
		Ta có
		\begin{align*}
			M&=\dfrac{2\sin \alpha-3\cos\alpha}{5\sin^3 \alpha+\cos^3 \alpha}\\
			&=\dfrac{2\left({\dfrac{1}{\sin^2 \alpha}}\right)-3\cot\alpha \left({\dfrac{1}{\sin^2 \alpha}}\right)}{5+\cot^3\alpha}\\
			&=\dfrac{-3\cot^3\alpha+2\cot^2\alpha-3\cot \alpha+2 }{5+\cot^3\alpha}\\
			&=-\dfrac{35}{16}.
		\end{align*}
	}
\end{vd}

% #1766 | Archive/2024/SP-K10-K11-LyVanHoang_TVN-023/11. C1-LuongGiac-lvhoang_TVN-023/dataC1K11/Bai1-Dang6.tex | mã gốc 1D1H2-2
\begin{bt}%[Dự án BG K10-K11]%[Dao-V- Thuy]%[1D1H2-2]
	Tính giá trị của biểu thức $P=\sin 30^{\circ}\cos60^{\circ}+\sin 60^{\circ}\cos 30^{\circ}$.
	\loigiai{
		Ta có $P=\sin 30^{\circ}\cos60^{\circ}+\sin 60^{\circ}\cos 30^{\circ}=\dfrac{1}{2}\cdot \dfrac{1}{2}+ \dfrac{\sqrt{3}}{2}\cdot \dfrac{\sqrt{3}}{2}=1$.
	}
\end{bt}

% #1767 | Archive/2024/SP-K10-K11-LyVanHoang_TVN-023/11. C1-LuongGiac-lvhoang_TVN-023/dataC1K11/Bai1-Dang6.tex | mã gốc 1D1H2-2
\begin{bt}%[Dự án BG K10-K11]%[Dao-V- Thuy]%[1D1H2-2]
	Cho $\sin\alpha=\dfrac{3}{5}$ với $\dfrac{\pi}{2}<\alpha<\pi$. Tính giá trị của $\cos\alpha$.
	\loigiai{
		Vì $\dfrac{\pi}{2}<\alpha<\pi$ nên $\cos\alpha<0$, suy ra $\cos\alpha=-\sqrt{1-\sin^2\alpha}=-\sqrt{1-\dfrac{9}{25}}=-\dfrac{4}{5}$.
	}
\end{bt}

% #1768 | Archive/2024/SP-K10-K11-LyVanHoang_TVN-023/11. C1-LuongGiac-lvhoang_TVN-023/dataC1K11/Bai1-Dang6.tex | mã gốc 1D1H2-2
\begin{bt}%[Dự án BG K10-K11]%[Dao-V- Thuy]%[1D1H2-2]
	Biết $\tan\alpha=2$ và $-\pi<\alpha<0$. Tính $\cos\alpha$.
	\loigiai{
		Ta có $\dfrac{1}{\cos^2\alpha}=1+\tan^2\alpha=5\Rightarrow \cos\alpha=\pm\dfrac{1}{\sqrt{5}}$.\\
		Mà $-\pi<\alpha<0$ và $\tan\alpha>0$ nên $\cos\alpha=-\dfrac{1}{\sqrt{5}}$.
	}
\end{bt}

% #1769 | Archive/2024/SP-K10-K11-LyVanHoang_TVN-023/11. C1-LuongGiac-lvhoang_TVN-023/dataC1K11/Bai1-Dang6.tex | mã gốc 1D1V2-2
\begin{bt}%[Dự án BG K10-K11]%[Dao-V- Thuy]%[1D1V2-2]
	Nếu $\sin x+\cos x=\dfrac{1}{2}$ và $0^\circ <x<180^\circ $ thì $\tan x=-\dfrac{p+\sqrt{q}}{3}$. Tìm cặp số nguyên $(p;q)$.
	\loigiai{
		Ta có $\sin x+\cos x=\dfrac{1}{2}\Rightarrow \sin x\cdot \cos x=\dfrac{-3}{8}$.\\
		$\Rightarrow \sin x$, $\cos x$ là nghiệm của phương trình $X^2-\dfrac{1}{2}X-\dfrac{3}{8}=0\Leftrightarrow\hoac{& X=\dfrac{1+\sqrt{7}}{4} \\ & X=\dfrac{1-\sqrt{7}}{4}.}$\\
		Vì $0^\circ <x<180^\circ $ nên $\sin x>0\Rightarrow \sin x=\dfrac{1+\sqrt{7}}{4}$; $\cos x=\dfrac{1-\sqrt{7}}{4}$
		\begin{center}
			$\Rightarrow \tan x=\dfrac{\sin x}{\cos x}=-\dfrac{4+\sqrt{7}}{3}\Rightarrow p=4$; $q=7$.
		\end{center}
	}
\end{bt}

% #1779 | Archive/2024/SP-K10-K11-LyVanHoang_TVN-023/11. C1-LuongGiac-lvhoang_TVN-023/dataC1K11/Bai2-Dang3-Congthucbiendoitichthanhtong.tex | mã gốc 1D1N3-4
\begin{vd}%[Huỳnh Đức Vũ, dự án bài giảng K10-11]%[1D1N3-4]
	Biết $\sin (a+b)=1$, $\sin (a-b)=\dfrac{1}{2}$. Tính $\sin a\cos b$.
	\loigiai{
		Ta có $\sin a\cos b=\dfrac{1}{2}\left[\sin (a-b)+\sin (a+b)\right]=\dfrac{1}{2}\left(\dfrac{1}{2}+1\right)=\dfrac{3}{4}$.
	}
\end{vd}

% #1780 | Archive/2024/SP-K10-K11-LyVanHoang_TVN-023/11. C1-LuongGiac-lvhoang_TVN-023/dataC1K11/Bai2-Dang3-Congthucbiendoitichthanhtong.tex | mã gốc 1D1H3-4
\begin{vd}%[Huỳnh Đức Vũ, dự án bài giảng K10-11]%[1D1H3-4] 
	Cho $\sin 2 x=-\dfrac{1}{3}$. Tính $A=\sin \left(x+\dfrac{\pi}{4}\right) \cos \left(x-\dfrac{\pi}{4}\right)$.
	\loigiai{
		$$
		\begin{aligned}
			A&=\sin \left(x+\frac{\pi}{4}\right) \cos \left(x-\dfrac{\pi}{4}\right)=\dfrac{1}{2}\left[\sin \left(x+\dfrac{\pi}{4}+x-\dfrac{\pi}{4}\right)+\sin \left(x+\dfrac{\pi}{4}-x+\dfrac{\pi}{4}\right)\right] \\
			&=\dfrac{1}{2}\left(\sin 2 x+\sin \dfrac{\pi}{2}\right)=\dfrac{1}{2}\left(-\dfrac{1}{3}+1\right)=\dfrac{1}{3} . 
		\end{aligned}
		$$
	}
\end{vd}

% #1781 | Archive/2024/SP-K10-K11-LyVanHoang_TVN-023/11. C1-LuongGiac-lvhoang_TVN-023/dataC1K11/Bai2-Dang3-Congthucbiendoitichthanhtong.tex | mã gốc 1D1H3-4
\begin{vd}%[Huỳnh Đức Vũ, dự án bài giảng K10-11]%[1D1H3-4] 
	Cho $\cos a=\dfrac{2}{3}$. Tính $B=\cos \dfrac{3 a}{2} \cos \dfrac{a}{2}$.
	\loigiai{
		Ta có $\cos2a=2\cos^2 a-1=-\dfrac{1}{9}$.\\
		Do đó
		$$B=\cos \dfrac{3 a}{2} \cos \dfrac{a}{2}=\dfrac{1}{2}\left[\cos \left(\dfrac{3 a}{2}+\dfrac{a}{2}\right)+\cos \left(\dfrac{3 a}{2}-\dfrac{a}{2}\right)\right]=\dfrac{1}{2}\left(\cos 2a+\cos a \right)=\dfrac{5}{18}.
		$$
	}
\end{vd}

% #1782 | Archive/2024/SP-K10-K11-LyVanHoang_TVN-023/11. C1-LuongGiac-lvhoang_TVN-023/dataC1K11/Bai2-Dang3-Congthucbiendoitichthanhtong.tex | mã gốc 1D1H3-4
\begin{vd}%[Huỳnh Đức Vũ, dự án bài giảng K10-11]%[1D1H3-4]
	Biến đổi thành tổng các biểu thức sau
	\begin{listEX}[2]
		\item $B= \sin 5x\cos 3x$.
		\item $E=\sin (x+y) \cos (x-y)$.
		\item $F=\sin \left( x+30^{\circ}\right) \cos \left( x-30^{\circ}\right)$.
		\end{listEX}
	\loigiai{
		\begin{enumerate}
			\item $B= \sin 5x\cos 3x= \dfrac{1}{2}\left( \sin 2x + \sin 8x\right)$.
			\item $E=\sin (x+y) \cos (x-y)= \dfrac{1}{2}\left( \sin 2y+ \sin 2x\right)$.
			\item $F=\sin \left( x+30^{\circ}\right) \cos \left( x-30^{\circ}\right)= \dfrac{1}{2}\left( \sin 60^\circ + \sin 2x\right)$.
			\end{enumerate}
	}
\end{vd}

% #1783 | Archive/2024/SP-K10-K11-LyVanHoang_TVN-023/11. C1-LuongGiac-lvhoang_TVN-023/dataC1K11/Bai2-Dang3-Congthucbiendoitichthanhtong.tex | mã gốc 1D1V3-4
\begin{vd}%[Huỳnh Đức Vũ, dự án bài giảng K10-11]%[1D1V3-4]
	Biến đổi thành tổng các biểu thức sau
	\begin{listEX}[2]
		\item $G= 2\sin x\sin 2x\sin 3x$.
		\item $H= 8\cos x\sin 2x\sin 3x$.
		\end{listEX}
	\loigiai{
		\begin{enumerate}
			\item $G= 2\sin x\sin 2x\sin 3x= \left( \cos x- \cos 3x\right)\sin 3x= \dfrac{1}{2}\left( \sin 2x+ \sin 4x- \sin 6x\right)$.
			\item 
			\begin{eqnarray*}
			H &=& 8\cos x\sin 2x\sin 3x=(2\sin2x\cos x)\cdot4\sin3x=(\sin x+\sin 3x)\cdot4\sin3x \\
			&=&2(2\sin3x\sin x+2\sin^23x)=2(\cos 2x-\cos 4x+1-\cos6x).
			\end{eqnarray*}
		\end{enumerate}
	}
\end{vd}

% #1786 | Archive/2024/SP-K10-K11-LyVanHoang_TVN-023/11. C1-LuongGiac-lvhoang_TVN-023/dataC1K11/Bai2-Dang3-Congthucbiendoitichthanhtong.tex | mã gốc 1D1H3-4
\begin{bt}%[Huỳnh Đức Vũ, dự án bài giảng K10-11]%[1D1H3-4]
	Cho $\cos 2 x=\dfrac{1}{4}$. Tính 
	\begin{listEX}[2]
		\item $A=\cos \left(x+\dfrac{\pi}{6}\right) \cos \left(x-\dfrac{\pi}{6}\right)$.
		\item $B=\sin \left(x+\dfrac{\pi}{3}\right) \sin \left(x-\dfrac{\pi}{3}\right)$.  	
	\end{listEX}   
	\loigiai{
		\begin{listEX}
			\item $A=\dfrac{1}{2}\left(\cos 2 x+\cos \dfrac{\pi}{3}\right)=\dfrac{1}{2}\left(\dfrac{1}{4}+\dfrac{1}{2}\right)=\dfrac{3}{8}$.
			\item $B=-\dfrac{1}{2}\left(\cos 2 x-\cos \dfrac{2\pi}{3}\right)=-\dfrac{1}{2}\left(\dfrac{1}{4}+\dfrac{1}{2}\right)=-\dfrac{3}{8}$.			
		\end{listEX}
	}
\end{bt}

% #1787 | Archive/2024/SP-K10-K11-LyVanHoang_TVN-023/11. C1-LuongGiac-lvhoang_TVN-023/dataC1K11/Bai2-Dang3-Congthucbiendoitichthanhtong.tex | mã gốc 1D1V3-4
\begin{bt}%[Huỳnh Đức Vũ, dự án bài giảng K10-11]%[1D1V3-4]
	Rút gọn các biểu thức sau
	\begin{enumerate}
		\item $A= \cos 11x \cos 3x- \cos 17x \cos 9x$.
		\item $B= \sin 18x \cos 13x- \sin 9x \cos 4x$.
		\item $C= \sin x\sin 3x+ \sin 4x\sin 8x$.
	\end{enumerate}
	\loigiai{
		\begin{enumerate}
			\item Ta có\\
			$\begin{aligned}
				A= \cos 11x \cos 3x- \cos 17x \cos 9x&= \dfrac{1}{2}(\cos 8x+ \cos 14x- \cos 8x- \cos 26x)\\
				&= \dfrac{1}{2}(\cos 14x- \cos 26x)= \sin 20x\sin 6x.
			\end{aligned}$
			\item $B= \sin 18x \cos 13x- \sin 9x \cos 4x= \dfrac{1}{2}\left( \sin 5x+ \sin 31x- \sin 5x- \sin 13x\right)= \cos 22x\sin 9x$.
			\item $C= \sin x\sin 3x+ \sin 4x\sin 8x= \dfrac{1}{2}\left( \cos 2x- \cos 4x+ \cos 4x- \cos 12x\right)= \sin 7x \sin 5x$.
		\end{enumerate}
	}
\end{bt}

% #1788 | Archive/2024/SP-K10-K11-LyVanHoang_TVN-023/11. C1-LuongGiac-lvhoang_TVN-023/dataC1K11/Bai2-Dang3-Congthucbiendoitichthanhtong.tex | mã gốc 1D1H3-4
\begin{bt}%[Huỳnh Đức Vũ, dự án bài giảng K10-11]%[1D1H3-4]
	Rút gọn các biểu thức sau
	\begin{enumerate}
		\item $D=\sin 2x\sin 6x- \cos x\cos 3x$.
		\item $E=\cos 3x\cos 6x- \cos 4x\cos 7x$.
		\item $F= \sin x\sin \left( 60^\circ - x\right) \sin \left( 60^\circ + x\right)$.
	\end{enumerate}
	\loigiai{
		\begin{enumerate}
			\item $D=\sin 2x\sin 6x- \cos x\cos 3x= \dfrac{1}{2}\left( \cos 4x- \cos 8x- \cos 2x- \cos 4x\right)= - \cos 5x\cos 3x$.
			\item $E=\cos 3x\cos 6x- \cos 4x\cos 7x= \dfrac{1}{2}\left( \cos 3x+ \cos 9x- \cos 3x- \cos 11x\right)= \sin 10x \sin x$.
			\item Ta có\\ $\begin{aligned}
				F= \sin x\sin \left( 60^\circ - x\right) \sin \left( 60^\circ + x\right)&= \dfrac{1}{2}\sin x\left( \cos 2x- \cos 120^\circ\right)\\
				&= \dfrac{1}{2}\cdot \dfrac{1}{2}\left( \sin (-x)+ \sin 3x\right)+ \dfrac{1}{4}\sin x\\
				&= \dfrac{1}{4}\sin 3x.
			\end{aligned}$
		\end{enumerate}
	}
\end{bt}

% #1789 | Archive/2024/SP-K10-K11-LyVanHoang_TVN-023/11. C1-LuongGiac-lvhoang_TVN-023/dataC1K11/Bai2-Dang3-Congthucbiendoitichthanhtong.tex | mã gốc 1D1V3-4
\begin{bt}%[Huỳnh Đức Vũ, dự án bài giảng K10-11]]%[1D1V3-4]
	Biến đổi thành tổng các biểu thức sau
	\begin{enumerate}
		\item $I=4\sin \left( x+\dfrac{\pi}{6}\right) \sin \left( x-\dfrac{\pi}{6}\right) \cos 2x$.
		\item $J=4\cos (a-b)\cos (b-c)\cos (c-a)$.
	\end{enumerate}
	\loigiai{
		\begin{enumerate}
			\item 
			\begin{eqnarray*}
			I&=&2\cdot2\sin \left( x+\dfrac{\pi}{6}\right)\sin \left( x-\dfrac{\pi}{6}\right) \cos 2x=2\cdot \left(\cos\dfrac{\pi}{3}-\cos2x\right)\cos2x\\
			&=&2\cos2x-2\cos^22x=2\cos2x-1-\cos4x.
			\end{eqnarray*}
			\item
			\begin{eqnarray*}
			J&=&2\cdot \left[2\cos (a-b)\cos (b-c)\right]\cos (c-a)=2\left[\cos (a-2b+c)+\cos (a-c)\right]\cos (c-a)\\
			&=&2\cos (a-2b+c) \cos(c-a)+2\cos^2(c-a)=\cos(2c-2b)+\cos(2a-2b)+1+\cos(2c-2a).
			\end{eqnarray*}
		\end{enumerate}
	}
\end{bt}

% #1798 | Archive/2024/SP-K10-K11-LyVanHoang_TVN-023/11. C1-LuongGiac-lvhoang_TVN-023/dataC1K11/Bai1-Dang4.tex | mã gốc 1D1N1-5
\begin{vd}[MĐ1]%[1D1N1-5]
	Biểu diễn trên đường tròn lượng giác góc lượng giác có số đo là $865^{\circ}$.
	\loigiai{
		\immini{Ta có $865^{\circ}=145^{\circ}+2\cdot 360^{\circ}$. Vậy điểm biểu diễn góc lượng giác có số đo $865^{\circ}$ là điểm $M$ trên phần đường tròn lượng giác thuộc góc phần tư thứ II sao cho $\widehat{A O M}=145^{\circ}$.}
		{\begin{tikzpicture}[scale=0.9, font=\footnotesize,line join=round, line cap=round, >=stealth];
				\begin{scope}
					\draw[->] (-2.5,0) -- (2.5,0) node[below] {\scriptsize $x$};
					\draw[->] (0,-2.5) -- (0,2.5) node[left] {\scriptsize $y$};
					\draw  (0:0)--(145:1.5) ;
					\draw (-0.3,0) node[below] {$O$};
					\draw  (148:2.1) node[right]{$M$};
					\draw  (45:2.2) node[right]{I} (135:2.3) node[right]{II} (220:2.2) node[right]{III} (315:2) node[right]{IV};
					\draw (0:1.5) arc (0:360:1.5);
					\draw (45:0.8) node{$145^\circ$};
					\draw (1.7,0) node[above] {$A$};
					\draw[-stealth] (0:.4) arc (0:145:.4);
					\draw[fill=black]  (-1.5,0) circle (1.2pt) (1.5,0) circle (1.2pt)
					(0,1.5) circle (1.2pt) (0,-1.5) circle (1.2pt) (145:1.5) circle (1.2pt);
					\path (270:2.8) node{a)};
				\end{scope}
		\end{tikzpicture}}
	}
\end{vd}

% #1799 | Archive/2024/SP-K10-K11-LyVanHoang_TVN-023/11. C1-LuongGiac-lvhoang_TVN-023/dataC1K11/Bai1-Dang4.tex | mã gốc 1D1N1-5
\begin{vd}[MĐ1]%[1D1N1-5]
	Biểu diễn trên đường tròn lượng giác góc lượng giác có số đo là $\dfrac{-7 \pi}{3}$.
	\loigiai{
		\immini{Ta có $\dfrac{-7 \pi}{3}=\dfrac{-\pi}{3}+(-1) \cdot 2 \pi$. Vậy điểm biểu diễn góc lượng giác có số đo $\dfrac{-7 \pi}{3}$ là điểm $N$ trên phần đường tròn lượng giác thuộc góc phần tư thứ IV sao cho $\widehat{A O N}=\dfrac{\pi}{3}$.}
		{\begin{tikzpicture}[scale=0.9, font=\footnotesize,line join=round, line cap=round, >=stealth];
				\begin{scope}
					\draw[->] (-2.5,0) -- (2.5,0) node[below] {\scriptsize $x$};
					\draw[->] (0,-2.5) -- (0,2.5) node[left] {\scriptsize $y$};
					\draw  (0:0)--(-60:1.5) ;
					\draw (-0.3,0) node[below] {$O$};
					\draw  (-68:1.7) node[right]{$N$};
					\draw  (45:2.2) node[right]{I} (135:2.3) node[right]{II} (220:2.2) node[right]{III} (315:2) node[right]{IV};
					\draw (0:1.5) arc (0:360:1.5);
					\draw (-30:0.9) node{$-\dfrac{\pi}{3}$};
					\draw (1.7,0) node[above] {$A$};
					\draw[-stealth] (0:.4) arc (0:-60:.4);
					\draw[fill=black]  (-1.5,0) circle (1.2pt) (1.5,0) circle (1.2pt)
					(0,1.5) circle (1.2pt) (0,-1.5) circle (1.2pt) (-60:1.5) circle (1.2pt);
					\path (270:2.8) node{b)};
				\end{scope}
		\end{tikzpicture}}
	}
\end{vd}

% #1800 | Archive/2024/SP-K10-K11-LyVanHoang_TVN-023/11. C1-LuongGiac-lvhoang_TVN-023/dataC1K11/Bai1-Dang4.tex | mã gốc 1D1N1-5
\begin{vd}[MĐ1]%[1D1N1-5]
	Xác định điểm trên đường tròn lượng giác biểu diễn góc lượng giác có số đo bằng $\dfrac{13 \pi}{4}$.
	\loigiai{
		\immini{Điểm $M$ trên đường đường tròn lượng giác biểu diễn góc lượng giác có số đo bằng $\dfrac{13 \pi}{4}$ được xác định trong hình bên.}
		{\begin{tikzpicture}[scale=1, font=\footnotesize, line join=round, line cap=round, >=stealth]
				\draw[-stealth] (-2.5,0) -- (3,0)node[below] {$x$};
				\draw[-stealth] (0,-2.5) -- (0,2.5)node[right] {$y$};
				\draw[orange] (0,0) circle (2);
				\draw[red,-stealth,smooth,samples=100] plot[domain =0:3.25*pi]({.5*(1.02)^(\x) *cos(\x r)},{.5*(1.02)^(\x) *sin(\x r)});
				\draw[red,fill=black] (0:0) node[below right]{$O$}--(585:2) circle(1pt) node[left]{$M$};
				\path (130:1.2) node[below=-2pt]{$\frac{13\pi}{4}$};
				\path (0:2) node[below right]{$A$};
		\end{tikzpicture}}
	}
\end{vd}

% #1801 | Archive/2024/SP-K10-K11-LyVanHoang_TVN-023/11. C1-LuongGiac-lvhoang_TVN-023/dataC1K11/Bai1-Dang4.tex | mã gốc 1D1N1-5
\begin{vd}[MĐ1]%[1D1N1-5]
	Xác định các điểm $M$ và $N$ trên đường tròn lượng giác lần lượt biểu diễn các góc lượng giác có số đo bằng $-150^{\circ}$.
	\loigiai{
		\immini{Điểm $N$ trên đường tròn lượng giác biểu diễn góc lượng giác có số đo bằng $-150^{\circ}$ được xác định trong hình bên.}
		{\begin{tikzpicture}[scale=1, font=\footnotesize, line join=round, line cap=round, >=stealth]
				\draw[-stealth] (-2.5,0) -- (3,0)node[below] {$x$};
				\draw[-stealth] (0,-2.5) -- (0,2.5)node[right] {$y$};
				\draw[orange] (0,0) circle (2);
				\draw[-stealth,blue] (0:1) arc (0:-150:1);
				\draw[blue,fill=black] (0:0) node[below right]{$O$}--(-150:2) circle(1pt) node[left]{$N$};
				\path (-120:1.2) node[below=-2pt]{$-150^{\circ}$};
				\path (0:2) node[below right]{$A$};
		\end{tikzpicture}}
	}
\end{vd}

% #1802 | Archive/2024/SP-K10-K11-LyVanHoang_TVN-023/11. C1-LuongGiac-lvhoang_TVN-023/dataC1K11/Bai1-Dang4.tex | mã gốc 1D1H1-5
\begin{vd}[MĐ2]%[1D1H1-5]
	\immini{	Trên đường tròn lượng giác với điểm gốc là $A$. Điểm $M$ thuộc đường tròn sao cho cung $AM$ có số đo $60^\circ$ (tham khảo hình vẽ bên). Gọi $N$, $P$, $Q$ lần lượt là các điểm đối xứng với điểm $M$ qua  gốc tọa độ, trục $Oy$ và trục $Ox$, $A'$ là điểm đối xứng với $A$ qua gốc tọa độ $O$. Mỗi khẳng định dưới đây đúng hay sai?
		\choiceTF[t]
		{\True Số đo của cung lượng giác $AP$ bằng $120^\circ+ k360^\circ$, $k\in\mathbb{Z}$}
		{\True Số đo của cung lượng giác $AQ$ bằng $-30^\circ+ k360^\circ$, $k\in\mathbb{Z}$}
		{\True Số đo của cung lượng giác $AA'$ bằng $180^\circ+ k360^\circ$, $k\in\mathbb{Z}$}
		{  Số đo của cung lượng giác $AN$ bằng $120^\circ+ k360^\circ$, $k\in\mathbb{Z}$}
	}
	{\begin{tikzpicture}[scale=1.4, font=\footnotesize, line join=round, line cap=round, >=stealth]
			\draw[->] (-1.3,0)--(1.5,0) node[below]{$x$};
			\draw[->] (0,-1.5)--(0,1.5) node[right]{$y$};
			\path
			(1,0) coordinate (A)
			(0:0) coordinate (O)
			(60:1) coordinate (M)
			(120:1) coordinate (P)
			(240:1) coordinate (N)
			(-60:1) coordinate (Q)
			(-1,0) coordinate (A')
			;
			\draw 
			(O) circle (1)
			(O)--(M)			
			;
			\foreach \p/\g in {A/-60, O/-30, M/60}
			\draw[fill=black] (\p) circle (1pt) node[shift=(\g:3mm)] {$\p$};
	\end{tikzpicture}}
	
	\loigiai{
		\immini{
			\begin{itemchoice}
				\itemch Ta có $\widehat{AOP}=120^\circ$ nên số đo của cung lượng giác $AP$ bằng \break $120^\circ+ k360^\circ$, $k\in\mathbb{Z}$.
				\itemch Ta có $\widehat{AOQ}=30^\circ$ nên số đo của cung lượng giác $AQ$ bằng \break $-30^\circ+ k360^\circ$, $k\in\mathbb{Z}$.
				\itemch Ta có $\widehat{AOA'}=180^\circ$ nên số đo của cung lượng giác $AA'$ bằng \break $180^\circ+ k360^\circ$, $k\in\mathbb{Z}$.
				\itemch Ta có $\widehat{AON}=120^\circ$ nên số đo của cung lượng giác $AA'$ bằng \break $-120^\circ+ k360^\circ$, $k\in\mathbb{Z}$.
			\end{itemchoice}			
		}{\begin{tikzpicture}[scale=1.4, font=\footnotesize, line join=round, line cap=round, >=stealth]
				\draw[->] (-1.3,0)--(1.5,0) node[below]{$x$};
				\draw[->] (0,-1.5)--(0,1.5) node[right]{$y$};
				\path
				(1,0) coordinate (A)
				(0:0) coordinate (O)
				(60:1) coordinate (M)
				(120:1) coordinate (P)
				(240:1) coordinate (N)
				(-60:1) coordinate (Q)
				(-1,0) coordinate (A')
				;
				\draw 
				(O) circle (1)
				(O)--(M) (O)--(P)
				(O)--(N) (O)--(Q)				
				;
				\foreach \p/\g in {A/-60, O/-30, M/60, P/90,N/-90, Q/-60,A'/130}
				\draw[fill=black] (\p) circle (1pt) node[shift=(\g:3mm)] {$\p$};
		\end{tikzpicture}}
	}
\end{vd}

% #1804 | Archive/2024/SP-K10-K11-LyVanHoang_TVN-023/11. C1-LuongGiac-lvhoang_TVN-023/dataC1K11/Bai1-Dang4.tex | mã gốc 1D1H1-5
\begin{bt}%[1D1H1-5]
	Biểu diễn góc $630^\circ$ trên đường tròn lượng giác gốc $A$.
	\loigiai
	{
		\immini
		{
			Ta có $630^\circ=270^\circ+360^\circ$.\\
			Vậy điểm biểu diễn góc $630^\circ$ trên đường tròn lượng giác gốc $A$ là $J$ như hình bên
		}
		{
			\begin{tikzpicture}[line join=round, line cap=round, >=stealth,font=\footnotesize, scale=0.8]
				\def\R{2}
				\draw[->](-\R-1,0)--(\R+1,0) node[right] {$x$};
				\draw[->](0,-\R-1)--(0,\R+1) node[right] {$y$};
				\coordinate (O) at (0,0)[fill=black]circle(1pt) node[below left]{$O$};
				\draw (O) circle (\R);
				\pgfmathsetmacro\xH{\R*cos(630)}\pgfmathsetmacro\yH{\R*sin(630)}
				\draw[fill=black] (\xH,\yH)circle(1pt) node[below right]{$J$};
				\draw[fill=black] (2,0)circle(1pt) node[above right]{$A$};
				\draw (O)--(\xH,\yH);
			\end{tikzpicture}
		}
	}
\end{bt}

% #1805 | Archive/2024/SP-K10-K11-LyVanHoang_TVN-023/11. C1-LuongGiac-lvhoang_TVN-023/dataC1K11/Bai1-Dang4.tex | mã gốc 1D1H1-5
\begin{bt}%[1D1H1-5]
	Biểu diễn góc $\dfrac{2010\pi}{4}$ trên đường tròn lượng giác gốc $A$.
	\loigiai
	{
		\immini
		{
			Ta có $\dfrac{2010\pi}{4}=\dfrac{\pi}{2}+251\cdot2\pi$.\\
			Vậy điểm biểu diễn góc $\dfrac{2010\pi}{4}$ trên đường tròn lượng giác gốc $A$ là $B$ như hình bên
		}
		{
			\begin{tikzpicture}[line join=round, line cap=round, >=stealth,font=\footnotesize, scale=0.8]
				\def\R{2}
				\draw[->](-\R-1,0)--(\R+1,0) node[right] {$x$};
				\draw[->](0,-\R-1)--(0,\R+1) node[right] {$y$};
				\coordinate (O) at (0,0)[fill=black]circle(1pt) node[below left]{$O$};
				\draw (O) circle (\R);
				\pgfmathsetmacro\xB{\R*cos(90)}\pgfmathsetmacro\yB{\R*sin(90)}
				\draw[fill=black] (\xB,\yB)circle(1pt) node[above left]{$B$};
				\draw[fill=black] (\R,0)circle(1pt) node[above right]{$A$};
				\draw (O)--(\xB,\yB);
			\end{tikzpicture}
		}
	}
\end{bt}

% #1806 | Archive/2024/SP-K10-K11-LyVanHoang_TVN-023/11. C1-LuongGiac-lvhoang_TVN-023/dataC1K11/Bai1-Dang4.tex | mã gốc 1D1H1-5
\begin{bt}%[1D1H1-5]
	Trên đường tròn lượng giác, hãy biểu diễn góc lượng giác có số đo có dạng là $\dfrac{\pi}{2}+k \pi\,(k \in \mathbb{Z})$.
	\loigiai{
		\immini{Ta có $\dfrac{\pi}{2}+k \pi=\dfrac{\pi}{2}+\dfrac{k2\pi}{2}\,(k \in \mathbb{Z})$.\\
			Vậy điểm biểu diễn góc lượng giác $\dfrac{\pi}{2}+k \pi\,(k \in \mathbb{Z})$ là $2$ điểm $M_1$, $M_2$ như hình vẽ bên.}
		{\begin{tikzpicture}[scale=0.9, font=\footnotesize,line join=round, line cap=round, >=stealth];
				\begin{scope}
					\draw[->] (-2.5,0) -- (2.5,0) node[below] {\scriptsize $x$};
					\draw[->] (0,-2.5) -- (0,2.5) node[left] {\scriptsize $y$};
					\draw (-0.3,0) node[below] {$O$};
					\draw (0:1.5) arc (0:360:1.5);
					\draw (0.3,1.5) node[above] {$M_1$};
					\draw (0.3,-1.5) node[below] {$M_2$};
					\draw[fill=black]  (0,1.5) circle (1.2pt) (0,-1.5) circle (1.2pt);
				\end{scope}
		\end{tikzpicture}}
	}
\end{bt}

% #1807 | Archive/2024/SP-K10-K11-LyVanHoang_TVN-023/11. C1-LuongGiac-lvhoang_TVN-023/dataC1K11/Bai1-Dang4.tex | mã gốc 1D1H1-5
\begin{bt}%[1D1H1-5]
	Trên đường tròn lượng giác, hãy biểu diễn góc lượng giác có số đo có dạng là $k \dfrac{\pi}{4}\,(k \in \mathbb{Z})$
	\loigiai{
		\immini{Ta có $k \dfrac{\pi}{4}=\dfrac{k 2\pi}{8}\,(k \in \mathbb{Z})$.\\
			Vậy điểm biểu diễn góc lượng giác $k \dfrac{\pi}{4}\,(k \in \mathbb{Z})$ là $8$ điểm $M_1$, $M_2$, $M_3$, $M_4$, $M_5$, $M_6$, $M_7$ và $M_8$ như hình vẽ bên.}
		{\begin{tikzpicture}[scale=0.9, font=\footnotesize,line join=round, line cap=round, >=stealth];
				\begin{scope}
					\draw[->] (-2.5,0) -- (2.5,0) node[below] {\scriptsize $x$};
					\draw[->] (0,-2.5) -- (0,2.5) node[left] {\scriptsize $y$};
					\draw  (1.5,0.2)node[right] {$M_1$}   (45:1.5)node[above right] {$M_2$}   (90:1.5)node[above left] {$M_3$}   (135:1.5)node[above left] {$M_4$}   (180:1.5)node[above left] {$M_5$}   (225:1.5)node[left] {$M_6$}   (270:1.5)node[below left] {$M_7$}   (315:1.5)node[below right] {$M_8$};
					\draw (0,0) node[above left] {$O$};
					\draw (0:1.5) arc (0:360:1.5);
					\draw[fill=black]  (1.5,0) circle (1.2pt) (45:1.5) circle (1.2pt) (90:1.5) circle (1.2pt) (135:1.5) circle (1.2pt) (180:1.5) circle (1.2pt) (225:1.5) circle (1.2pt) (270:1.5) circle (1.2pt) (315:1.5) circle (1.2pt);
					\path (270:2.8) node{b)};
				\end{scope}
		\end{tikzpicture}}
	}
\end{bt}

% #1810 | Archive/2024/SP-K10-K11-LyVanHoang_TVN-023/11. C1-LuongGiac-lvhoang_TVN-023/dataC1K11/Bai1-Dang4.tex | mã gốc 1D1H1-5
\begin{ex}%[1D1H1-5]
	Cho hình vẽ sau:\\
	\begin{center}
		\begin{tikzpicture}[>=stealth,line join=round,line cap=round,font=\footnotesize,scale=.7]
			\def\aa{3cm}
			\draw (0,0)coordinate(O) circle[radius=3cm];
			\draw[->] (-4,0)--(4,0)node[below]{$x$};
			\draw[->] (0,-4)--(0,4)node[right]{$y$};
			\coordinate (M) at ($(0,0)+(60:3cm)$);
			\path(\aa,0)coordinate(A)(-\aa,0)coordinate(A')(0,\aa)coordinate(B)(0,-\aa)coordinate(B')($(M)!2!(O)$)coordinate(N)(O)--++(30:.9)coordinate(G);
			\draw(N)--(M);
			\draw pic[draw,blue,fill=green!50,opacity=.5,angle radius=5mm] {angle = A--O--M};
			\draw (G)node[rotate=10]{\scriptsize $\dfrac{\pi}{3}$};
			\foreach \diem/\goc in {O/-45,A/45,B/45,B'/-45,A'/135,N/-95,M/45}
			\fill[black](\diem) circle (1pt) ($(\diem)+(\goc:3mm)$) node{$\diem$};		
		\end{tikzpicture}
	\end{center}
	Xét tính đúng sai của các khẳng định sau
	\choiceTF[t]
	{Số đo góc lượng giác $(OM,OA)$ là $\text{sđ} (OM,OA)=\dfrac{\pi}{3}+k 2\pi(k\in\mathbb{Z})$}
	{\True $\text{sđ}(ON,OA)=\text{sđ}(ON,OM)-\text{sđ} (OA,OM)$}
	{Độ dài cung tròn $AM$ lớn là $l_{AM}=\dfrac{2\pi}{3}$}
	{\True Hai điểm $M,N$ biểu diễn các cung có số đo là $x=\dfrac{\pi}{3}+k\pi(k\in\mathbb{Z})$}
	\loigiai{
		Xét tính đúng-sai của các phát biểu sau:
		\begin{enumerate}
			\item Số đo góc lượng giác $(OM,OA)$ là  $\text{sđ}(OM,OA)=-\dfrac{\pi}{3}+k 2\pi(k\in\mathbb{Z})$ nên mệnh đề sai.
			\item Theo hệ thức Chales $\text{sđ}(ON,OA)+\text{sđ}(OA,OM)=\text{sđ}(ON,OM)$.
			$\Rightarrow \text{sđ} (ON,OA)=\text{sđ}(ON,OM)-\text{sđ}(OA,OM)$ nên mệnh đề đúng.
			\item Độ dài cung tròn $AM$ lớn là $\ell_{AM}=\dfrac{5\pi}{3}$ nên mệnh đề sai.
			\item Hai điểm $M,N$ biểu diễn các cung có số đo là $x=\dfrac{\pi}{3}+k\pi(k\in\mathbb{Z})$ nên mệnh đề đúng. 
		\end{enumerate}
	}
\end{ex}

% #1811 | Archive/2024/SP-K10-K11-LyVanHoang_TVN-023/11. C1-LuongGiac-lvhoang_TVN-023/dataC1K11/Bai1-Dang4.tex | mã gốc 1D1H1-5
\begin{ex}%[1D1H1-5]
	\immini{Trên đường tròn lượng giác điểm gốc $A$ như hình vẽ bên có bao nhiêu điểm biểu diễn góc lượng giác $\alpha=\dfrac{k\pi}{2}, k\in \mathbb{Z}$ trên đường tròn lượng giác.}
	{\begin{tikzpicture}[scale=1, font=\footnotesize, line join=round, line cap=round,>=stealth]
			\def \xmin{-2.0};
			\def \xmax{2.5};
			\def \ymin{-2.0};
			\def \ymax{2.3};
			\def \r{1.5};
			\draw[->] (\xmin, 0.) -- (\xmax,0.) node[anchor=north] {$x$};
			\draw[->] (0.,\ymin) -- (0.,\ymax) node[anchor=west] {$y$};
			\clip(\xmin+0.1,\ymin+0.1) rectangle (\xmax-0.1,\ymax-0.1);
			\path (0,0) coordinate (O)++(0:\r) coordinate (A) (O)++(45:\r) coordinate (M)  (O)++(135:\r) coordinate (N)  (O)++(225:\r) coordinate (P) (O)++(-45:\r) coordinate (Q) (O)++(90:\r) coordinate (B) (O)++(180:\r) coordinate (A') (O)++(-90:\r) coordinate (B');
			\draw (0,0) circle (\r);
			\foreach \x/\g  in {O/225,A/30,M/45,N/135,P/225,Q/-45,B/60,A'/150,B'/-60}{\fill[black] (\x) circle(1pt)+(\g:0.3) node{$\x$};}
	\end{tikzpicture}}
	
	\shortans{$4$}
	\loigiai{
		\immini{Cho $k=0 \Rightarrow \alpha=0$ ta có điểm biểu diễn $A$ như hình bên.\\
			Cho $k=1 \Rightarrow \alpha=\dfrac{\pi}{2}$ ta có điểm biểu diễn $B$ như hình bên.\\
			Cho $k=2 \Rightarrow \alpha=\pi$ ta có điểm biểu diễn $A'$ như hình bên.\\
			Cho $k=3 \Rightarrow \alpha=\dfrac{3 \pi}{2}$ ta có điểm biểu diễn $B'$ như hình bên.\\
		Vậy có $4$ điểm biểu diễn góc lượng giác $\alpha=\dfrac{k\pi}{2}$ trên đường tròn lượng giác.
		}
		{\begin{tikzpicture}[scale=1, font=\footnotesize, line join=round, line cap=round,>=stealth]
				\def \xmin{-2.0};
				\def \xmax{2.5};
				\def \ymin{-2.0};
				\def \ymax{2.3};
				\def \r{1.5};
				\draw[->] (\xmin, 0.) -- (\xmax,0.) node[anchor=north] {$x$};
				\draw[->] (0.,\ymin) -- (0.,\ymax) node[anchor=west] {$y$};
				\clip(\xmin+0.1,\ymin+0.1) rectangle (\xmax-0.1,\ymax-0.1);
				\path (0,0) coordinate (O)++(0:\r) coordinate (A) (O)++(45:\r) coordinate (M)  (O)++(135:\r) coordinate (N)  (O)++(225:\r) coordinate (P) (O)++(-45:\r) coordinate (Q) (O)++(90:\r) coordinate (B) (O)++(180:\r) coordinate (A') (O)++(-90:\r) coordinate (B');
				\draw (0,0) circle (\r);
				\foreach \x/\g  in {O/225,A/30,M/45,N/135,P/225,Q/-45,B/60,A'/150,B'/-60}{\fill[black] (\x) circle(1pt)+(\g:0.3) node{$\x$};}
		\end{tikzpicture}}
		}
\end{ex}

% #1812 | Archive/2024/SP-K10-K11-LyVanHoang_TVN-023/11. C1-LuongGiac-lvhoang_TVN-023/dataC1K11/Bai1-Dang4.tex | mã gốc 1D1H1-5
\begin{ex}%[1D1H1-5]
	\immini{Trên đường tròn lượng giác điểm gốc $A$ như hình vẽ bên các điểm $M$, $N$ là những điểm biểu diễn của các góc lượng giác có số đo là $\dfrac{4\pi}{a}+k\pi, k\in\mathbb{Z}$, $a$ là số nguyên dương. Tính $a^2+a$.}
	{\begin{tikzpicture}[declare function={r=1.5;},>=stealth]
			\path (0:r) coordinate (A)
			(90:r) coordinate (B)
			(180:r) coordinate (A')
			(270:r) coordinate (B')
			(60:r) coordinate (M)
			(240:r) coordinate (N)
			(0,0) coordinate(O);
			\path pic["\scriptsize$\dfrac{\pi}{3}$", angle eccentricity=1.5,draw,angle radius=19pt,fill=cyan!35]{angle= A--O--M};
			\draw[->] ({-1.25*r},0)--({1.25*r},0);
			\draw[->] (0,{-1.25*r})--(0,{1.25*r});
			\draw (N)--(M) (O) circle (r);	
			\foreach \t/\g in {O/140,A/35,B/60,A'/140,B'/310,M/60,N/240}{
				\draw(\t) circle (1pt) node[shift={(\g:7pt)},font=\scriptsize]{$ \t $};
			}
	\end{tikzpicture}}
	\shortans{$12$}
	\loigiai{
		Theo hình vẽ ta thấy hai điểm $ M$, $ N$ biểu diễn góc lượng giác có số đo $\dfrac{4\pi}{3}+k\pi, k\in\mathbb{Z}$.\\
		Suy ra $a=3$.\\
		Vậy $a^2+a=12$.	
	}
\end{ex}

% #1816 | Archive/2024/SP-K10-K11-LyVanHoang_TVN-023/11. C1-LuongGiac-lvhoang_TVN-023/dataC1K11/Bai-1-Dang-8-RutGonBieuThucLuongGiac.tex | mã gốc 1D1H2-3
\begin{vd}[MĐ2]%[1D1H2-3][CTST]
	Chứng minh đẳng thức lượng giác sau	
	$\tan \alpha + \cot \alpha=\dfrac{1}{\sin \alpha \cos \alpha}$.
	\loigiai{Ta có $VT=\tan \alpha+\cot \alpha =\dfrac{\sin \alpha}{\cos \alpha}+\dfrac{\cos \alpha}{\sin \alpha}=\dfrac{\sin^2\alpha+\cos^2\alpha}{\sin \alpha\cdot \cos \alpha}=\dfrac{1}{\sin \alpha\cdot \cos \alpha}=VP$.
	}
\end{vd}

% #1817 | Archive/2024/SP-K10-K11-LyVanHoang_TVN-023/11. C1-LuongGiac-lvhoang_TVN-023/dataC1K11/Bai-1-Dang-8-RutGonBieuThucLuongGiac.tex | mã gốc 1D1H2-3
\begin{vd}[MĐ2]%[1D1H2-3][CTST]
	Rút gọn các biểu thức sau
	\begin{enumerate}
		\item $\dfrac{1}{\tan \alpha+1}+\dfrac{1}{\cot \alpha+1}$;
		\item $\cos \left(\dfrac{\pi}{2}-\alpha\right)-\sin (\pi+\alpha)$;
		\item $\sin \left(\alpha-\dfrac{\pi}{2}\right)+\cos (-\alpha+6 \pi)-\tan (\alpha+\pi) \cot (3 \pi-\alpha)$.
	\end{enumerate}
	\loigiai{
		\begin{enumerate}
			\item Ta có $\dfrac{1}{\tan \alpha+1}+\dfrac{1}{\cot \alpha+1}=\dfrac{1}{\dfrac{\sin \alpha}{\cos \alpha}+1}+\dfrac{1}{\dfrac{\cos \alpha}{\sin \alpha}+1}=\dfrac{\cos \alpha}{\sin \alpha + \cos \alpha}+\dfrac{\sin \alpha}{\sin \alpha + \cos \alpha}=1$;
			\item Ta có $\cos \left(\dfrac{\pi}{2}-\alpha\right)-\sin (\pi+\alpha)=\sin \alpha + \sin \alpha =2\sin \alpha$;
			\item Ta có
			\begin{eqnarray*}
				\sin \left(\alpha-\dfrac{\pi}{2}\right)+\cos (-\alpha+6 \pi)-\tan (\alpha+\pi) \cot (3 \pi-\alpha)
				& = & -\sin \alpha + \cos (-\alpha) + \tan \alpha \cot \alpha\\
				& = & -\sin \alpha + \cos \alpha + 1\\
				&=& 1-\sin \alpha +\cos \alpha.
			\end{eqnarray*}
		\end{enumerate}
	}
\end{vd}

% #1818 | Archive/2024/SP-K10-K11-LyVanHoang_TVN-023/11. C1-LuongGiac-lvhoang_TVN-023/dataC1K11/Bai-1-Dang-8-RutGonBieuThucLuongGiac.tex | mã gốc 1D1V2-3
\begin{vd}[KNTT,MĐ3]%[1D1V2-3]
	Chứng minh các đẳng thức:
	\begin{enumerate}
		\item $\cos^4\alpha-\sin^4\alpha=2\cos^2\alpha-1$;
		\item $\dfrac{\cos^2\alpha+\tan^2\alpha-1}{\sin^2\alpha}=\tan^2\alpha$.
	\end{enumerate}
	\loigiai{
		\begin{enumerate}
			\item Ta có
			\begin{eqnarray*}
				\cos^4\alpha-\sin^4\alpha&=& \left(\cos^2\alpha-\sin^2\alpha\right)\left(\cos^2\alpha+\sin^2\alpha\right)\\
				&= & \cos^2\alpha-\left(1-\sin^2\alpha\right)\\
				&= & 2\cos^2\alpha-1.
			\end{eqnarray*}
			\item Ta có
			\begin{eqnarray*}
				\dfrac{\cos^2\alpha+\tan^2\alpha-1}{\sin^2\alpha}&=& \dfrac{\cos^2\alpha-1}{\sin^2\alpha}+\dfrac{\tan^2\alpha}{\sin^2\alpha}\\
				&= & \dfrac{-\sin^2\alpha}{\sin^2\alpha}+\dfrac{\sin^2\alpha}{\cos^2\alpha}\cdot\dfrac{1}{\sin^2\alpha}\\
				&= & \dfrac{1}{\cos^2\alpha}-1
				=\dfrac{1-\cos^2\alpha}{\cos^2\alpha}\\
				&=&\dfrac{\sin^2\alpha}{\cos^2\alpha}=\tan^2\alpha.
			\end{eqnarray*}
		\end{enumerate}	
	}
\end{vd}

% #1819 | Archive/2024/SP-K10-K11-LyVanHoang_TVN-023/11. C1-LuongGiac-lvhoang_TVN-023/dataC1K11/Bai-1-Dang-8-RutGonBieuThucLuongGiac.tex | mã gốc 1D1V2-3
\begin{vd}[SBT KNTT,MĐ3]%[1D1V2-3]
	Biết $\sin a+\cos a=m$. Hãy tính theo $m$
	\begin{enumerate}
		\item $\sin a\cos a$;
		\item $\sin^3a+\cos^3a$;
		\item $\sin^4a+\cos^4a$.
	\end{enumerate}
	\loigiai{
		\begin{enumerate}
			\item Ta có $( \sin a+\cos  a)^2=\sin ^2a+2\sin a\cos a+{\cos ^2}a=1+2\sin a\cos a\ (*)$.\\
			Mặt khác $\sin a+\cos a=m$ suy ra $m^2=1+2\sin a\cos a$ hay $\sin a \cos a =\dfrac{m^2-1}{2}$.
			\item $\sin^3a+\cos^3a=(\sin a+\cos a)^3-3\sin a\cos a\left(\sin^2a+\cos^2a\right)=m^3-\dfrac{3}{2}\left(m^2-1\right).$
			\item $\left(\sin^4a+\cos^4a\right)=\left(\sin^2a+\cos^2a\right)^2-2\sin^2a\cos^2a=1-2\left(\dfrac{m^2-1}{2}\right)^2=\dfrac{-m^4+2m^2+1}{2}$.
		\end{enumerate}
	}
\end{vd}

% #1824 | Archive/2024/SP-K10-K11-LyVanHoang_TVN-023/11. C1-LuongGiac-lvhoang_TVN-023/dataC1K11/Bai-1-Dang-8-RutGonBieuThucLuongGiac.tex | mã gốc 1D1V2-3
\begin{bt}%[1D1V2-3]
	Cho $\sin a+\cos a=m$. Tính giá trị của $\sin a\cos a$, từ đó tính giá trị của $\left|\sin a-\cos a\right|$ theo $m$.
	\loigiai{
		\begin{itemize}
			\item Ta có $\sin a+\cos a =m$. \\
			Suy ra $(\sin a+\cos a)^2=m^2\Leftrightarrow 1+2\sin a\cos a=m^2\Leftrightarrow \sin a\cos a=\dfrac{m^2-1}{2}$.
			\item Ta có 
			$$(\sin a-\cos a)^2=1-2\sin a\cos a=1-(m^2-1)=2-m^2.$$
			Vậy $|\sin a-\cos a|=\sqrt{2-m^2}$.
		\end{itemize}
	}
\end{bt}

% #1825 | Archive/2024/SP-K10-K11-LyVanHoang_TVN-023/11. C1-LuongGiac-lvhoang_TVN-023/dataC1K11/Bai-1-Dang-8-RutGonBieuThucLuongGiac.tex | mã gốc 1D1V2-3
\begin{bt}[SBT CTST]%[1D1V2-3]
	Chứng minh các đẳng thức lượng giác sau:
	\begin{enumerate}
		\item $\dfrac{\sin ^3 \alpha+\cos ^3 \alpha}{\sin \alpha+\cos \alpha}=1-\sin \alpha \cos \alpha$;
		\item $\sin ^4 x+\cos ^4 x=1-2 \sin ^2 x \cos ^2 x$,
		\item $\dfrac{1+\cot x}{1-\cot x}=\dfrac{\tan x+1}{\tan x-1}$;
		\item $\dfrac{\sin \alpha+\cos \alpha}{\sin ^3 \alpha}=\dfrac{1-\cot ^4 \alpha}{1-\cot \alpha}$.
	\end{enumerate}
	\loigiai{
		\begin{enumerate}
			\item $\begin{aligned}[t]\dfrac{\sin ^3 \alpha+\cos ^3 \alpha}{\sin \alpha+\cos \alpha} & =\dfrac{(\sin \alpha+\cos \alpha)\left(\sin ^2 \alpha-\sin \alpha \cos \alpha+\cos ^2 \alpha\right)}{\sin \alpha+\cos \alpha} \\
				& =\sin ^2 \alpha-\sin \alpha \cos \alpha+\cos ^2 \alpha \\
				& =1-\sin \alpha \cos \alpha .
			\end{aligned}$
			\item $\sin ^4 x+\cos ^4 x=\left(\sin ^2 x+\cos ^2 x\right)^2-2 \sin ^2 x \cos ^2 x=1-2 \sin ^2 x \cos ^2 x$;
			\item $\dfrac{1+\cot x}{1-\cot x}=\dfrac{1+\dfrac{1}{\tan x}}{1-\dfrac{1}{\tan x}}=\dfrac{\tan x+1}{\tan x-1}$;
			\item 
			$\dfrac{\sin \alpha+\cos \alpha}{\sin ^3 \alpha}  =\dfrac{1}{\sin ^2 \alpha}+\dfrac{\cos \alpha}{\sin \alpha} \cdot \dfrac{1}{\sin ^2 \alpha}=\left(1+\cot ^2 \alpha\right)+\cot \alpha\left(1+\cot ^2 \alpha\right) 
			=(1+\cot \alpha)\left(1+\cot ^2 \alpha\right) \\
			=\dfrac{\left(1-\cot ^2 \alpha\right)\left(1+\cot ^2 \alpha\right)}{1-\cot \alpha}=\dfrac{1-\cot ^4 \alpha}{1-\cot \alpha}
			$
		\end{enumerate}
	}
\end{bt}

% #1826 | Archive/2024/SP-K10-K11-LyVanHoang_TVN-023/11. C1-LuongGiac-lvhoang_TVN-023/dataC1K11/Bai-1-Dang-8-RutGonBieuThucLuongGiac.tex | mã gốc 1D1V2-3
\begin{bt}%[1D1V2-3]
	Tính giá trị biểu thức $A=\dfrac{\sin^2 a+2\sin a\cos a-2\cos^2 a}{2\sin^2 a-3\sin a\cos a+4\cos^2 a}$ khi $\cot a=-3$.
	\loigiai{
		Vì $\cot a=-3$ nên $\sin a\neq 0$. Chia cả tử và mẫu của $A$ cho $\sin a$, ta có 
		$$D=\dfrac{1+2\cot a-2\cot^2 a}{2-3\cot a+4\cot^2 a}=-\dfrac{23}{47}.$$
	}
\end{bt}

% #1827 | Archive/2024/SP-K10-K11-LyVanHoang_TVN-023/11. C1-LuongGiac-lvhoang_TVN-023/dataC1K11/Bai-1-Dang-8-RutGonBieuThucLuongGiac.tex | mã gốc 1D1V2-3
\begin{bt}%[1D1V2-3]
	Tính giá trị biểu thức $A=\dfrac{3\cos x+4\sin x +2}{-2\sin x+\cos x}$, với $\sin x=\dfrac{3}{5}$ và $\dfrac{\pi}{2}<x<\pi$.
	\loigiai{
		Ta có $\dfrac{\pi}{2}<x<\pi$ nên $\cos x<0$. 
		Vậy $\cos x=-\sqrt{1-\sin^2 x}=-\dfrac{4}{5}$. \\
		Vậy $A=\dfrac{-3\cdot \dfrac{4}{5}+4\cdot\dfrac{3}{5}+2}{-2\cdot \dfrac{3}{5}-\dfrac{4}{5}}=-1$.
	}
\end{bt}

% #1829 | Archive/2024/SP-K10-K11-LyVanHoang_TVN-023/11. C1-LuongGiac-lvhoang_TVN-023/dataC1K11/Bai-1-Dang-8-RutGonBieuThucLuongGiac.tex | mã gốc 1D1V2-3
\begin{bt}%[1D1V2-3]
	Cho $\tan\alpha=2$. Tính giá trị của biểu thức $C=\dfrac{\sin\alpha}{\sin^3\alpha+2\cos^3\alpha}$.
	\loigiai{
		Ta có $C=\dfrac{\sin\alpha}{\sin^3\alpha+2\cos^3\alpha}$ $=\dfrac{\tan\alpha\cdot\dfrac{1}{\cos^2\alpha}}{\tan^3\alpha+2}$ $=\dfrac{\tan\alpha\left(1+\tan^2\alpha\right)}{\tan^3\alpha+2}$ $=\dfrac{2\left(1+2^2\right)}{2^3+2}=1$.}
\end{bt}

% #1842 | Archive/2024/SP-K10-K11-LyVanHoang_TVN-023/11. C1-LuongGiac-lvhoang_TVN-023/dataC1K11/Bai-3-Dang-1-Timtapxacdinhcuahamso.tex | mã gốc 1D1N4-2
\begin{vd}[Chân trời sáng tạo, Mức độ 1]%[Dự án BG-K11-Lần-2]%[Ngô Quang Anh]%[1D1N4-2]
	Tìm tập xác định của hàm số $y=\tan 3x$.
	\loigiai{
		Điều kiện xác định của hàm số $y=\tan 3x$ ta có $3x\ne \dfrac{\pi}{2}+k\pi \Leftrightarrow x\neq\dfrac{\pi}{6}+k\dfrac{\pi}{3},\, k\in \mathbb{Z}$.\\
		Vậy tập xác định của hàm số $ \mathscr{D}=\mathbb{R}\setminus\left\{\dfrac{\pi}{6}+k\dfrac{\pi}{3} \,\middle|\, k\in \mathbb{Z} \right\}$.}
\end{vd}

% #1843 | Archive/2024/SP-K10-K11-LyVanHoang_TVN-023/11. C1-LuongGiac-lvhoang_TVN-023/dataC1K11/Bai-3-Dang-1-Timtapxacdinhcuahamso.tex | mã gốc 1D1N4-2
\begin{vd}[Kết nối tri thức với cuộc sống, Mức độ 1]%[Dự án BG-K11-Lần-2]%[Ngô Quang Anh]%[1D1N4-2]
	Tìm tập xác định của hàm số $y=\cot 2x$.
	\loigiai{
		Điều kiện xác định của hàm số $y=\cot 2x$ ta có $2x\ne k\pi \Leftrightarrow x\neq k\dfrac{\pi}{2},\, k\in \mathbb{Z}$.\\
		Vậy tập xác định của hàm số $ \mathscr{D}=\mathbb{R}\setminus\left\{k\dfrac{\pi}{2}\,\middle|\, k\in \mathbb{Z} \right\}$.}
\end{vd}

% #1846 | Archive/2024/SP-K10-K11-LyVanHoang_TVN-023/11. C1-LuongGiac-lvhoang_TVN-023/dataC1K11/Bai-3-Dang-1-Timtapxacdinhcuahamso.tex | mã gốc 1D1H4-2
\begin{bt}[Mức độ 2]%[Dự án BG-K11-Lần-2]%[Ngô Quang Anh]%[1D1H4-2]
	Tìm tập xác định $\mathscr{D}$ của hàm số $y=3\tan^2\left(\dfrac{x}{2}-\dfrac{\pi}{4}\right).$
	\loigiai{
		Hàm số xác định $\Leftrightarrow \cos\left(\dfrac{x}{2}-\dfrac{\pi}{4}\right)\ne 0\Leftrightarrow \dfrac{x}{2}-\dfrac{\pi}{4}\ne \dfrac{\pi}{2}+k\pi \Leftrightarrow x\ne \dfrac{3\pi}{2}+k2\pi, k\in \mathbb{Z}.$\\
		Vậy tập xác định $ \mathscr{D}=\mathbb{R}\setminus\left\{\dfrac{3\pi}{2}+k2\pi \,\middle|\, k\in \mathbb{Z}\right\}$.
		}
\end{bt}

% #1847 | Archive/2024/SP-K10-K11-LyVanHoang_TVN-023/11. C1-LuongGiac-lvhoang_TVN-023/dataC1K11/Bai-3-Dang-1-Timtapxacdinhcuahamso.tex | mã gốc 1D1H4-2
\begin{bt}[Mức độ 2]%[Dự án BG-K11-Lần-2]%[Ngô Quang Anh]%[1D1H4-2]
	Tìm tập xác định $\mathscr{D} $ của hàm số $y=\dfrac{3\tan x - 5}{1-\sin^2 x}$.
	\loigiai{
		Hàm số xác định khi và chỉ khi $$\heva{&\cos x \ne 0 \\& 1-\sin^2 x \ne 0} \Leftrightarrow \cos x \ne 0 \Leftrightarrow x \ne \dfrac{\pi}{2}+k\pi,k\in \mathbb{Z} .$$
		Vậy tập xác định của hàm số là $ \mathscr{D} =\mathbb{R} \setminus \left\{\dfrac{\pi}{2}+k\pi \,\middle|\, k\in \mathbb{Z} \right\}$.
	}
\end{bt}

% #1848 | Archive/2024/SP-K10-K11-LyVanHoang_TVN-023/11. C1-LuongGiac-lvhoang_TVN-023/dataC1K11/Bai-3-Dang-1-Timtapxacdinhcuahamso.tex | mã gốc 1D1H4-2
\begin{bt}[Mức độ 2]%[Dự án BG-K11-Lần-2]%[Ngô Quang Anh]%[1D1H4-2]
	Tìm tập xác định của hàm số $y=\sqrt{\dfrac{2+\sin x}{1-\cos x}}$.
	\loigiai{
		Hàm số xác định khi $\heva{&\dfrac{2+\sin x}{1-\cos x}\ge 0\\&1-\cos x\ne 0}\Leftrightarrow 1-\cos x\ne 0\Leftrightarrow \cos x\ne 1\Leftrightarrow x\ne k2\pi$, $k\in\mathbb{Z}$.\\
		Vậy tập xác định $\mathscr{D}=\mathbb{R}\setminus \left\{k2\pi \mid k\in\mathbb{Z}\right\}$.
	}
\end{bt}

% #1849 | Archive/2024/SP-K10-K11-LyVanHoang_TVN-023/11. C1-LuongGiac-lvhoang_TVN-023/dataC1K11/Bai-3-Dang-1-Timtapxacdinhcuahamso.tex | mã gốc 1D1V4-2
\begin{bt}[Mức độ 3]%[Dự án BG-K11-Lần-2]%[Ngô Quang Anh]%[1D1V4-2]
	Tìm tập xác định của hàm số $y=\cos 2x+\dfrac{1}{\tan x}+5$.
	\loigiai{
		Hàm số xác định khi $\heva{&\tan x\ne 0\\&\cos x\ne 0}\Leftrightarrow \heva{&x\ne k\pi\\&x\ne \dfrac{\pi}{2}+k\pi}\Leftrightarrow x\ne \dfrac{k\pi}{2}$ ($k\in \mathbb{Z}$).\\
		Vậy tập xác định của hàm số là $\mathscr{D}=\mathbb{R}\setminus \left\{ \dfrac{k\pi}{2} \mid k\in \mathbb{Z} \right\}$.
	}
\end{bt}

% #1850 | Archive/2024/SP-K10-K11-LyVanHoang_TVN-023/11. C1-LuongGiac-lvhoang_TVN-023/dataC1K11/Bai-3-Dang-1-Timtapxacdinhcuahamso.tex | mã gốc 1D1V4-2
\begin{bt}[Mức độ 3]%[Dự án BG-K11-Lần-2]%[Ngô Quang Anh]%[1D1V4-2]
	Tìm tập xác định của hàm số $y=\dfrac{\cot x}{\cos x-1}$.
	\loigiai{
		Hàm số xác định khi và chỉ khi $\heva{&\sin x\ne 0\\ &\cos x\ne 1}\Leftrightarrow\heva{&x\ne k\pi\\ &x\ne \ell2\pi}\ (k, \ell\in\mathbb{Z})\Leftrightarrow x\ne k\pi, k\in\mathbb{Z}$.\\
		Vậy, tập xác định của hàm số $y=\dfrac{\cot x}{\cos x-1}$ là $\mathbb{R}\setminus\left\{k\pi \mid k\in\mathbb{Z}\right\}$.
	}
\end{bt}

% #1860 | Archive/2024/SP-K10-K11-LyVanHoang_TVN-023/11. C1-LuongGiac-lvhoang_TVN-023/dataC1K11/Bai1-Dang5.tex | mã gốc 1D1N2-2
\begin{vd}[MĐ1]%[Dự án BG K10-K11]%[Dao-V- Thuy]%[1D1N2-2][Cánh Diều]
	Xét dấu các giá trị lượng giác của góc lượng giác $\alpha=-\dfrac{3\pi}{4}$.
	\loigiai{
		Do $-\pi<-\dfrac{3\pi}{4}<-\dfrac{\pi}{2}$ nên $\sin \left(-\dfrac{3\pi}{4}\right)<0$; $\cos \left(-\dfrac{3\pi}{4}\right)<0$; $\tan \left(-\dfrac{3\pi}{4}\right)>0$; $\cot \left(-\dfrac{3\pi}{4}\right)>0$.
	}
\end{vd}

% #1861 | Archive/2024/SP-K10-K11-LyVanHoang_TVN-023/11. C1-LuongGiac-lvhoang_TVN-023/dataC1K11/Bai1-Dang5.tex | mã gốc 1D1N2-2
\begin{vd}[MĐ1]%[Dự án BG K10-K11]%[Dao-V- Thuy]%[1D1N2-2][Cánh Diều]
	Xét dấu các giá trị lượng giác của góc lượng giác $\alpha=\dfrac{5\pi}{6}$.
	\loigiai{
		Do $\dfrac{\pi}{2}<\dfrac{5\pi}{6}<\pi$ nên $\sin \left(\dfrac{5\pi}{6}\right)>0$; $\cos \left(\dfrac{5\pi}{6}\right)<0$; $\tan \left(\dfrac{5\pi}{6}\right)<0$; $\cot \left(\dfrac{5\pi}{6}\right)<0$.
	}
\end{vd}

% #1862 | Archive/2024/SP-K10-K11-LyVanHoang_TVN-023/11. C1-LuongGiac-lvhoang_TVN-023/dataC1K11/Bai1-Dang5.tex | mã gốc 1D1H2-2
\begin{vd}[MĐ2]%[Dự án BG K10-K11]%[Dao-V- Thuy]%[1D1H2-2]
	Xác định dấu các biểu thức
	\begin{listEX}[2]
		\item $A=\sin 50^\circ \cdot \cos (-100^\circ)$.
		\item $B=\sin 195^\circ \cdot \tan \dfrac{20\pi}{7}$.
	\end{listEX}
	\loigiai{
		\begin{enumerate}
			\item $A=\sin 50^\circ \cdot \cos (-100^\circ)$.\\
			Ta có điểm cuối của cung $50^\circ$ thuộc góc phần tư thứ $\mathit{I}$ nên $\sin 50^\circ>0$.\\
			Điểm cuối của cung $-100^\circ$ thuộc góc phần tư thứ $\mathit{III}$ nên $\cos (-100^\circ)<0$.\\
			Do đó, $A<0$.
			\item $B=\sin 195^\circ \cdot \tan \dfrac{20\pi}{7}$.\\
			Ta có điểm cuối của cung $195^\circ$ thuộc góc phần tư thứ $\mathit{III}$ nên $\sin 195^\circ<0$.\\
			Điểm cuối của cung $\dfrac{20\pi}{7}=\dfrac{6\pi}{7}+2\pi$ thuộc góc phần tư thứ $\mathit{II}$ nên $\tan \dfrac{20\pi}{7}<0$.\\
			Do đó, $B>0$.
		\end{enumerate}
	}
\end{vd}

% #1863 | Archive/2024/SP-K10-K11-LyVanHoang_TVN-023/11. C1-LuongGiac-lvhoang_TVN-023/dataC1K11/Bai1-Dang5.tex | mã gốc 1D1H2-2
\begin{vd}[MĐ2]%[Dự án BG K10-K11]%[Dao-V- Thuy]%[1D1H2-2]
	Xác định dấu các biểu thức
	\begin{listEX}[2]
		\item $A=\cot \dfrac{2\pi}{5} \cdot \sin \left( -\dfrac{2\pi}{3}\right)$.
		\item $B=\cos \dfrac{4\pi}{5}\cdot\sin \dfrac{\pi}{3}\cdot\tan\dfrac{4\pi}{3}\cdot\cot\dfrac{9\pi}{5}$.
	\end{listEX}
	\loigiai{
		\begin{enumerate}
			\item $A=\cot \dfrac{2\pi}{5} \cdot \sin \left( -\dfrac{2\pi}{3}\right)$.\\
			Ta có điểm cuối của cung $\dfrac{2\pi}{5}$ thuộc góc phần tư thứ $\mathit{I}$ nên $\cot \dfrac{2\pi}{5}>0$.\\
			Điểm cuối của cung $-\dfrac{2\pi}{3}$ thuộc góc phần tư thứ $\mathit{III}$ nên $\sin \left( -\dfrac{2\pi}{3}\right)<0$.\\
			Do đó, $A<0$.
			\item $B=\cos \dfrac{4\pi}{5}\cdot\sin \dfrac{\pi}{3}\cdot\tan\dfrac{4\pi}{3}\cdot\cot\dfrac{9\pi}{5}$.\\
			Ta có điểm cuối của cung $\dfrac{4\pi}{5}$ thuộc góc phần tư thứ $\mathit{II}$ nên $\cos \dfrac{4\pi}{5}<0$.\\
			Điểm cuối của cung $\dfrac{\pi}{3}$ thuộc góc phần tư thứ $\mathit{I}$ nên $\sin \dfrac{\pi}{3}>0$.\\
			Điểm cuối của cung $\dfrac{4\pi}{3}$ thuộc góc phần tư thứ $\mathit{III}$ nên $\tan \dfrac{4\pi}{3}>0$.\\
			Điểm cuối của cung $\dfrac{9\pi}{5}=-\dfrac{\pi}{5}+2\pi$ thuộc góc phần tư thứ $\mathit{IV}$ nên $\cot \dfrac{9\pi}{5}<0$.\\
			Do đó, $B>0$.
		\end{enumerate}
	}
\end{vd}

% #1864 | Archive/2024/SP-K10-K11-LyVanHoang_TVN-023/11. C1-LuongGiac-lvhoang_TVN-023/dataC1K11/Bai1-Dang5.tex | mã gốc 1D1H2-2
\begin{vd}[MĐ2]%[Dự án BG K10-K11]%[Dao-V- Thuy]%[1D1H2-2]
	Cho $\pi<\alpha<\dfrac{3\pi}{2}$. Xét dấu các biểu thức sau:
	\begin{listEX}[2]
		\item $A=\cos\left(\alpha-\dfrac{\pi}{2}\right)$.
		\item $B=\tan\left(\dfrac{2025\pi}{2}-\alpha\right)$.
	\end{listEX}
	\loigiai{
		\begin{enumerate}
			\item $A=\cos\left(\alpha-\dfrac{\pi}{2}\right)=\cos\left(\dfrac{\pi}{2}-\alpha\right)=\sin \alpha <0$.
			\item $B=\tan\left(\dfrac{2025\pi}{2}-\alpha\right)=\tan\left(\dfrac{\pi}{2}-\alpha+1012\pi\right)=\tan\left(\dfrac{\pi}{2}-\alpha\right)=\cot \alpha>0$.
		\end{enumerate}
	}
\end{vd}

% #1867 | Archive/2024/SP-K10-K11-LyVanHoang_TVN-023/11. C1-LuongGiac-lvhoang_TVN-023/dataC1K11/Bai1-Dang5.tex | mã gốc 1D1N2-2
\begin{bt}%[Dự án BG K10-K11]%[Dao-V- Thuy]%[1D1N2-2]
	Xét dấu các giá trị lượng giác của góc $\alpha =\dfrac{2\pi}{5}$.
	\loigiai{
		Vì $0<\dfrac{2\pi}{5}<\dfrac{\pi}{2}$ nên $\sin \alpha >0$, $\cos \alpha >0$, $\tan\alpha >0$, $\cot\alpha >0$.	
	}
\end{bt}

% #1868 | Archive/2024/SP-K10-K11-LyVanHoang_TVN-023/11. C1-LuongGiac-lvhoang_TVN-023/dataC1K11/Bai1-Dang5.tex | mã gốc 1D1N2-2
\begin{bt}%[Dự án BG K10-K11]%[Dao-V- Thuy]%[1D1N2-2]
	Xét dấu các giá trị lượng giác của góc $\alpha =-\dfrac{\pi}{7}$.
	\loigiai{
		Vì $-\dfrac{\pi}{2}<-\dfrac{\pi}{7}<0$ nên $\sin \alpha <0$, $\cos \alpha >0$, $\tan\alpha <0$, $\cot\alpha <0$.	
	}
\end{bt}

% #1869 | Archive/2024/SP-K10-K11-LyVanHoang_TVN-023/11. C1-LuongGiac-lvhoang_TVN-023/dataC1K11/Bai1-Dang5.tex | mã gốc 1D1N2-2
\begin{bt}%[Dự án BG K10-K11]%[Dao-V- Thuy]%[1D1N2-2]
	Xét dấu các giá trị lượng giác của góc $\beta =-\dfrac{2\pi}{3}$.
	\loigiai{
		Vì $-\pi<-\dfrac{2\pi}{3}<-\dfrac{\pi}{2}$ nên $\sin \beta<0$, $\cos \beta<0$, $\tan\alpha >0$, $\cot\alpha >0$.	
	}
\end{bt}

% #1870 | Archive/2024/SP-K10-K11-LyVanHoang_TVN-023/11. C1-LuongGiac-lvhoang_TVN-023/dataC1K11/Bai1-Dang5.tex | mã gốc 1D1N2-2
\begin{bt}%[Dự án BG K10-K11]%[Dao-V- Thuy]%[1D1N2-2]
	Xét dấu các giá trị lượng giác của góc $\alpha =\dfrac{11\pi}{4}$.
	\loigiai{
		Ta có $\dfrac{11\pi}{4}=\dfrac{3\pi}{4}+2\pi$ nên các giá trị lượng giác của góc $\dfrac{11\pi}{4}$ bằng các giá trị lượng giác của góc $\dfrac{3\pi}{4}$. \\
		Có $\dfrac{\pi}{2}<\dfrac{3\pi}{4}<\pi$,
		suy ra $\sin \dfrac{3\pi}{4}>0$, $\cos \dfrac{3\pi}{4}<0$, $\tan \dfrac{3\pi}{4}<0$ và $\cot\dfrac{3\pi}{4}<0$. \\
		Vậy $\sin \dfrac{11\pi}{4}>0$, $\cos \dfrac{11\pi}{4}<0$, $\tan \dfrac{11\pi}{4}<0$ và $\cot \dfrac{11\pi}{4}<0$.
	}
\end{bt}

% #1871 | Archive/2024/SP-K10-K11-LyVanHoang_TVN-023/11. C1-LuongGiac-lvhoang_TVN-023/dataC1K11/Bai1-Dang5.tex | mã gốc 1D1H2-2
\begin{bt}%[Dự án BG K10-K11]%[Dao-V- Thuy]%[1D1H2-2]
	Xác định dấu của $\sin \alpha$, $\cos \alpha$, $\tan \alpha$, biết
	\begin{listEX}[3]
		\item $\dfrac{3\pi}{2}<\alpha<\dfrac{7\pi}{4}$.
		\item $3\pi<\alpha<\dfrac{10\pi}{3}$.
		\item $\dfrac{5\pi}{2}<\alpha<\dfrac{11\pi}{4}$.
	\end{listEX}
	\loigiai{
		\begin{enumerate}
			\item $\dfrac{3\pi}{2}<\alpha<\dfrac{7\pi}{4}$.\\
			Ta có: điểm cuối của cung $\alpha$ thuộc góc phần tư thứ $\mathit{IV}$ nên $\sin \alpha<0$, $\cos \alpha>0$, $\tan \alpha<0$.
			\item $3\pi<\alpha<\dfrac{10\pi}{3}$.\\
			Ta có: điểm cuối của cung $\alpha$ thuộc góc phần tư thứ $\mathit{III}$ nên $\sin \alpha<0$, $\cos \alpha<0$, $\tan \alpha>0$.	
			\item $\dfrac{5\pi}{2}<\alpha<\dfrac{11\pi}{4}$.\\
			Ta có: điểm cuối của cung $\alpha$ thuộc góc phần tư thứ $\mathit{II}$ nên $\sin \alpha>0$, $\cos \alpha<0$, $\tan \alpha<0$.\\	
		\end{enumerate}
	}
\end{bt}

% #1872 | Archive/2024/SP-K10-K11-LyVanHoang_TVN-023/11. C1-LuongGiac-lvhoang_TVN-023/dataC1K11/Bai1-Dang5.tex | mã gốc 1D1H2-2
\begin{bt}%[Dự án BG K10-K11]%[Dao-V- Thuy]%[1D1H2-2]
	Cho $90^\circ<\alpha<180^\circ$. Xét dấu các biểu thức sau:
	\begin{listEX}[2]
		\item $A=\sin (270^\circ-\alpha)$.
		\item $B=\cos (2\alpha+90^\circ)$.
	\end{listEX}
	\loigiai{
		\begin{enumerate}
			\item $A=\sin (270^\circ-\alpha)$.\\
			Vì $-180^\circ<-\alpha<-90^\circ$ nên $90^\circ<270^\circ-\alpha<180^\circ$.\\
			Do đó, điểm cuối của cung $270^\circ-\alpha$ thuộc góc phần tư thứ $\mathit{II}$ nên $A>0$.
			\item $B=\cos (2\alpha+90^\circ)$.\\
			Ta có $B=\cos (2\alpha+90^\circ)=\cos (90^\circ-(-2\alpha))=\sin(-2\alpha)=-\sin(2\alpha)$.\\
			Vì $180^\circ<2\alpha<360^\circ$ nên $\sin (2\alpha)<0$.\\
			Do đó, $B>0$.
		\end{enumerate}
	}
\end{bt}

% #1873 | Archive/2024/SP-K10-K11-LyVanHoang_TVN-023/11. C1-LuongGiac-lvhoang_TVN-023/dataC1K11/Bai1-Dang5.tex | mã gốc 1D1H2-2
\begin{bt}%[Dự án BG K10-K11]%[Dao-V- Thuy]%[1D1H2-2]
	Cho $0<\alpha<\dfrac{\pi}{2}$. Xét dấu các biểu thức sau
	\begin{listEX}[2]
		\item $A=\cos \left(\alpha+\dfrac{3\pi}{5}\right)$.
		\item $B=\cos \left(\alpha-\dfrac{\pi}{8}\right)$.
	\end{listEX}
	\loigiai{
		\begin{enumerate}
			\item $A=\cos (\alpha+\dfrac{3\pi}{5})$.\\
			Vì $0<\alpha<\dfrac{\pi}{2}$ nên $\dfrac{3\pi}{5}<\alpha+\dfrac{3\pi}{5}<\dfrac{11\pi}{10}$.\\
			Do đó, điểm cuối của cung $\alpha+\dfrac{3\pi}{5}$ thuộc góc phần tư thứ $\mathit{II}$ hoặc thứ $\mathit{III}$.\\
			Vậy $A<0$.
			\item $B=\cos (\alpha-\dfrac{\pi}{8})$.\\
			Vì $0<\alpha<\dfrac{\pi}{2}$ nên $-\dfrac{\pi}{8}<\alpha-\dfrac{\pi}{8}<\dfrac{3\pi}{8}$.\\
			Do đó, điểm cuối của cung $\alpha-\dfrac{\pi}{8}$ thuộc góc phần tư thứ $\mathit{IV}$ hoặc thứ $\mathit{I}$.\\
			Vậy $B>0$.
		\end{enumerate}
	}
\end{bt}

% #1874 | Archive/2024/SP-K10-K11-LyVanHoang_TVN-023/11. C1-LuongGiac-lvhoang_TVN-023/dataC1K11/Bai1-Dang5.tex | mã gốc 1D1H2-2
\begin{bt}%[Dự án BG K10-K11]%[Dao-V- Thuy]%[1D1H2-2]
	Cho $\pi<\alpha<\dfrac{3\pi}{2}$. Xét dấu các biểu thức sau:
	\begin{listEX}[2]
		\item $A=\sin\left(\alpha+\dfrac{\pi}{2}\right)$.
		\item $B=\sin\left(\alpha+\dfrac{1119\pi}{2}\right)$.
	\end{listEX}
	\loigiai{
		\begin{enumerate}
			\item $A=\sin\left(\alpha+\dfrac{\pi}{2}\right)=\sin\left(\dfrac{\pi}{2}-(-\alpha)\right)=\cos(-\alpha)=\cos \alpha<0$.
			\item $B=\sin\left(\alpha+\dfrac{1119\pi}{2}\right)=\sin\left(\alpha-\dfrac{\pi}{2}+280\cdot2\pi\right)\\=\sin\left(\alpha-\dfrac{\pi}{2}\right)=-\sin\left(\dfrac{\pi}{2}-\alpha\right)=-\cos \alpha>0$.
		\end{enumerate}
	}
\end{bt}

% #1875 | Archive/2024/SP-K10-K11-LyVanHoang_TVN-023/11. C1-LuongGiac-lvhoang_TVN-023/dataC1K11/Bai-3-Dang-4-Xettinhtuanhoancuahamso.tex | mã gốc 1D1N4-5
\begin{vd}%[Dự án bài giảng new 4in1, Nguyễn Tiến]%[1D1N4-5]
	[Chân trời sáng tạo - Mức độ 1]
	Xét tính tuần hoàn và tìm chu kì (nếu có) của hàm số $y=\sin x$.
	\loigiai{
		Ta đã biết, với mỗi số nguyên $k$, số $k2\pi$ thoả mãn
		$$\sin (x+k2\pi)=\sin x, \forall x.$$
		Ngược lại, nếu có số $T\ne0$ sao cho
		$$\sin (x+T)=\sin x, \, \forall x \in \mathbb{R}$$
		phải có dạng $T=k2\pi$, với $k$ là một số nguyên.\\
		Rõ ràng, trong các số dạng $k2\pi$ $(k \in \mathbb{Z})$, số dương nhỏ nhất là $2\pi$.\\
		Vậy hàm số $y=\sin x$ là hàm số tuần hoàn với chu kì số $T=2\pi$.
	}	
\end{vd}

% #1876 | Archive/2024/SP-K10-K11-LyVanHoang_TVN-023/11. C1-LuongGiac-lvhoang_TVN-023/dataC1K11/Bai-3-Dang-4-Xettinhtuanhoancuahamso.tex | mã gốc 1D1H4-5
\begin{vd}%[Dự án bài giảng new 4in1, Nguyễn Tiến]%[1D1H4-5]
	[Kết nối tri thức với cuộc sống - Mức độ 2]
	Xét tính tuần hoàn và tìm chu kì (nếu có) của hàm số $y=\sin 2x$.
	\loigiai{
		Đặt $f(x)=\sin 2x$.\\
		Hàm số có tập xác định là $\mathbb{R}$.\\
		Với mọi số thực $x$, ta có $f(x)=\sin 2x=\sin (2x+k2\pi)=\sin 2(x+k\pi)$, $(k\in \mathbb{Z})$.\\
		Lấy $T=k\pi$, $k\in \mathbb{Z}$, $k\neq 0$ thì ta có
		\begin{align*}
			& x-T \in \mathbb{R}, x+T \in \mathbb{R},\\
			& f(x+T)= \sin 2(x+T)=\sin 2(x+k\pi)=\sin 2x=f(x), \, \forall x\in \mathbb{R}.
		\end{align*}
		Suy ra hàm số $y=\sin 2 x$ là hàm số tuần hoàn.\\
		Mặt khác, nếu có số $T\neq 0$ sao cho $\sin 2(x+T)=\sin 2x$ thì $T$ phải có dạng $T=k\pi$, $k\in \mathbb{Z}$.\\
		Rõ ràng, trong các số có dạng $k\pi$, $k\in \mathbb{Z}$ thì số dương nhỏ nhất là $\pi$.\\
		Vậy hàm số $f(x)=\sin 2x$ tuần hoàn với chu kì $T=\pi$.
		\begin{luuy}
			Ta có thể áp dụng nhanh cách tìm chu kì $T$ của hàm số $y=\sin 2x$ như sau:
			$$T=\dfrac{2\pi}{\left|a\right|}=\pi.$$
		\end{luuy}
	}	
\end{vd}

% #1881 | Archive/2024/SP-K10-K11-LyVanHoang_TVN-023/11. C1-LuongGiac-lvhoang_TVN-023/dataC1K11/Bai-3-Dang-4-Xettinhtuanhoancuahamso.tex | mã gốc 1D1H4-5
\begin{bt}%[Dự án bài giảng new 4in1, Nguyễn Tiến]%[1D1H4-5]
	[Mức độ 2]
	Tìm chu kì của hàm số $y=\sin 2x$.
	\loigiai{
		Xét phương trình sau với $T>0$:
		$$\sin[2(x+T)]=\sin 2x \Leftrightarrow\hoac{& T=k\pi\\& T=\dfrac{\pi}{2}-2x+k\pi},\, (k\in\mathbb{Z}).$$
		Do đó, số $T>0$ nhỏ nhất thỏa mãn $\sin[2(x+T)]=\sin 2x$, $\forall x\in\mathbb{R}$ là $T=\pi$.
	}
\end{bt}

% #1882 | Archive/2024/SP-K10-K11-LyVanHoang_TVN-023/11. C1-LuongGiac-lvhoang_TVN-023/dataC1K11/Bai-3-Dang-4-Xettinhtuanhoancuahamso.tex | mã gốc 1D1H4-5
\begin{bt}%[Dự án bài giảng new 4in1, Nguyễn Tiến]%[1D1H4-5]
	[Mức độ 2]
	Tìm chu kì của hàm số $f(x)=\tan\dfrac{x}{4}+2\sin\dfrac{x}{2}$.
	\loigiai{
		Hàm số $\tan x$ tuần hoàn với chu kì $\pi$ nên hàm số $\tan\dfrac{x}{4}$ tuần hoàn với chu kì $4\pi$.\\
		Hàm số $\sin x$ tuần hoàn với chu kì $2\pi$ nên hàm số $\sin\dfrac{x}{2}$ tuần hoàn với chu kì $4\pi$.\\
		Do đó hàm số $f(x)=\tan\dfrac{x}{4}+2\sin\dfrac{x}{2}$ tuần hoàn với chu kì $4\pi$.
	}
\end{bt}

% #1883 | Archive/2024/SP-K10-K11-LyVanHoang_TVN-023/11. C1-LuongGiac-lvhoang_TVN-023/dataC1K11/Bai-3-Dang-4-Xettinhtuanhoancuahamso.tex | mã gốc 1D1H4-5
\begin{bt}%[Dự án bài giảng new 4in1, Nguyễn Tiến]%[1D1H4-5]
	[Mức độ 2]
	Xét tính tuần hoàn và tìm chu kì (nếu có) của hàm số $y=\tan 2x$.
	\loigiai{
		Đặt $f(x)=\tan 2x$.\\
		Hàm số có tập xác định là $\mathscr{D}=\left\lbrace x\in\mathbb{R} \,\middle|\, x\neq\dfrac{\pi}{4}+\dfrac{k\pi}{2},\, k\in\mathbb{Z}\right\rbrace$.\\
		Với mọi số thực $x$, ta có
		$$f(x)=\tan 2x=\tan (2x+k\pi)=\tan 2\left(x+\dfrac{k\pi}{2}\right), k\in \mathbb{Z}.$$
		Lấy $T=\dfrac{k\pi}{2}$, $k\in\mathbb{Z}$, $k\neq 0$ thì ta có
		\begin{align*}
			& x-T\in\mathscr{D},\, x+T\in\mathscr{D},\\
			& f(x+T)=\tan 2(x+T)=\tan 2\left(x+\dfrac{k\pi}{2}\right)=\tan 2x=f(x),\, \forall x\in\mathscr{D}.
		\end{align*}
		Suy ra hàm số $y=\tan 2x$ là hàm số tuần hoàn.\\
		Mặt khác, nếu có số $T\neq 0$ sao cho $\tan 2(x+T)=\sin 2x$ thì $T$ phải có dạng $T=\dfrac{k\pi}{2}$, $k\in \mathbb{Z}$.\\
		Rõ ràng, trong các số có dạng $\dfrac{k\pi}{2}$, $k\in\mathbb{Z}$ thì số dương nhỏ nhất là $\dfrac{\pi}{2}$.\\
		Vậy hàm số $f(x)=\tan 2x$ tuần hoàn với chu kì $T=\dfrac{\pi}{2}$.
	}
\end{bt}

% #1884 | Archive/2024/SP-K10-K11-LyVanHoang_TVN-023/11. C1-LuongGiac-lvhoang_TVN-023/dataC1K11/Bai-3-Dang-4-Xettinhtuanhoancuahamso.tex | mã gốc 1D1H4-5
\begin{bt}%[Dự án bài giảng new 4in1, Nguyễn Tiến]%[1D1H4-5]
	[Mức độ 2]
	Xét tính tuần hoàn và tìm chu kì (nếu có) của hàm số $y=\cot 3x$.
	\loigiai{
		Đặt $f(x)=\cot 3x$.\\
		Hàm số có tập xác định là $\mathscr{D}=\left\lbrace x\in\mathbb{R} \,\middle|\, x\neq\dfrac{k\pi}{3},\, k\in\mathbb{Z}\right\rbrace$.\\
		Với mọi số thực $x$, ta có
		$$f(x)=\cot 3x=\cot (3x+k\pi)=\cot 3\left(x+\dfrac{k\pi}{3}\right),\, k\in\mathbb{Z}.$$
		Lấy $T=\dfrac{k\pi}{3}$, $k\in\mathbb{Z}$, $k\neq 0$ thì ta có
		\begin{align*}
			& x-T\in\mathscr{D},\, x+T\in\mathscr{D},\\
			& f(x+T)=\cot 3(x+T)=\cot 3\left(x+\dfrac{k\pi}{3}\right)=\cot 3x=f(x),\, \forall x\in\mathscr{D}.
		\end{align*}
		Suy ra hàm số $y=\cot 3x$ là hàm số tuần hoàn.\\
		Mặt khác, nếu có số $T\neq 0$ sao cho $\cot 3(x+T)=\cot 3x$ thì $T$ phải có dạng $T=\dfrac{k\pi}{3}$, $k\in \mathbb{Z}$.\\
		Rõ ràng, trong các số có dạng $\dfrac{k\pi}{3}$, $k\in\mathbb{Z}$ thì số dương nhỏ nhất là $\dfrac{\pi}{3}$.\\
		Vậy hàm số $f(x)=\cot 3x$ tuần hoàn với chu kì $T=\dfrac{\pi}{3}$.
	}
\end{bt}

% #1885 | Archive/2024/SP-K10-K11-LyVanHoang_TVN-023/11. C1-LuongGiac-lvhoang_TVN-023/dataC1K11/Bai-3-Dang-4-Xettinhtuanhoancuahamso.tex | mã gốc 1D1V4-5
\begin{bt}%[Dự án bài giảng new 4in1, Nguyễn Tiến]%[1D1V4-5]
	[Mức độ 3]
	Xét tính tuần hoàn của hàm số $y=x+\sin x$.
	\loigiai{
		Tập xác định $\mathscr{D} = \mathbb{R}$.\\
		Giả sử tồn tại số thực $T\neq 0$ sao cho $f\left(x+T\right)=f(x)$, $\forall x\in\mathbb{R}$, ta có
		\allowdisplaybreaks
		\begin{eqnarray*}
			& & \left(x+T\right)+\sin\left(x+T\right)=x+\sin x,\, \forall x\in\mathscr{D}\\
			&\Leftrightarrow & T+\sin\left(x+T\right)=\sin x,\, \forall x\in\mathbb{R}.
		\end{eqnarray*}
		Cho $x=0$ và $x=\pi $, ta được
		$$\heva{
			& T+\sin x=\sin 0=0 \\ 
			& T+\sin \left(\pi+T\right)=\sin \pi =0}
		\Rightarrow 2T+\sin T+\sin \left(\pi+T\right)=0\Leftrightarrow T=0.$$
		Điều này mâu thuẫn với giả thiết $T\neq 0$.\\
		Vậy hàm số $y=x+\sin x$ không phải là hàm số tuần hoàn.	
	}
\end{bt}

% #1894 | Archive/2024/SP-K10-K11-LyVanHoang_TVN-023/11. C1-LuongGiac-lvhoang_TVN-023/dataC1K11/Bai1-Dang1.tex | mã gốc 1D1H1-2
\begin{vd}[Cánh Diều, Mức độ 2]%[1D1H1-2]
	Hoàn thành bảng chuyển đổi số đo độ và số đo radian của một số góc đặc biệt sau
	\begin{center}
		\begin{longtable}{|c|c|c|c|c|c|c|c|}
			\hline
			Độ & $30^\circ$ & ? & $60^\circ$ & ? & $120^\circ$ & ? & $180^\circ$\\
			\hline
			Radian & ? & $\dfrac{\pi}{4}$ & ? & $\dfrac{\pi}{2}$ & ? & $\dfrac{3\pi}{4}$ & ? \\
			\hline
		\end{longtable}
	\end{center}
	\loigiai{
		Ta có bảng chuyển đổi số đo độ và số đo radian của một số góc đặc biệt.
		\begin{center}
			\begin{longtable}{|c|c|c|c|c|c|c|c|}
				\hline
				Độ & $30^\circ$ & $45^\circ$ & $60^\circ$ & $90^\circ$ & $120^\circ$ & $135^\circ$ & $180^\circ$\\
				\hline
				Radian & $\dfrac{\pi}{6}$ & $\dfrac{\pi}{4}$ & $\dfrac{\pi}{3}$ & $\dfrac{\pi}{2}$ & $\dfrac{2\pi}{3}$ & $\dfrac{3\pi}{4}$ & $\pi$ \\
				\hline
			\end{longtable}
		\end{center}
	}
\end{vd}

% #1895 | Archive/2024/SP-K10-K11-LyVanHoang_TVN-023/11. C1-LuongGiac-lvhoang_TVN-023/dataC1K11/Bai1-Dang1.tex | mã gốc 1D1H1-2
\begin{vd}[Cánh Diều, Mức độ 2]%[1D1H1-2]
	Hãy hoàn thành bảng chuyển đổi số đo độ và số đo radian của một số góc sau.
	\begin{center}
		\begin{longtable}{|c|c|c|c|c|}
			\hline
			Độ & $18^\circ$ & ? & $72^\circ$ & ? \\
			\hline
			Radian & ? & $\dfrac{2\pi}{9}$ & ? & $\dfrac{5\pi}{6}$\\
			\hline
		\end{longtable}
	\end{center}
	\loigiai{
		Ta có bảng chuyển đổi số đo độ và số đo radian của một số góc.
		\begin{center}
			\begin{longtable}{|c|c|c|c|c|}
				\hline
				Độ & $18^\circ$ & $40^\circ$ & $72^\circ$ & $150^\circ$ \\
				\hline
				Radian & $\dfrac{\pi}{10}$ & $\dfrac{2\pi}{9}$ & $\dfrac{2\pi}{5}$ & $\dfrac{5\pi}{6}$\\
				\hline
			\end{longtable}
		\end{center}
	}
\end{vd}

% #1896 | Archive/2024/SP-K10-K11-LyVanHoang_TVN-023/11. C1-LuongGiac-lvhoang_TVN-023/dataC1K11/Bai1-Dang1.tex | mã gốc 1D1H1-2
\begin{vd}[Kết nối tri thức, Mức độ 2]%[1D1H1-2]
	\begin{enumerate}
		\item Đổi từ độ sang rađian các số đo sau: $45^{\circ} ; 150^{\circ}$.
		\item Đổi từ rađian sang độ các số đo sau: $\dfrac{\pi}{3} ; \dfrac{5 \pi}{4}$.
	\end{enumerate}
	\loigiai{
		a) Ta có:
		$$
		\begin{aligned}
			& 45^{\circ}=45 \cdot \frac{\pi}{180}=\frac{\pi}{4}, \\
			& 150^{\circ}=150 \cdot \frac{\pi}{180}=\frac{5 \pi}{6} .
		\end{aligned}
		$$
		b) Ta có:
		$$
		\begin{aligned}
			& \frac{\pi}{3}=\frac{\pi}{3} \cdot\left(\frac{180}{\pi}\right)^\circ=60^{\circ}, \\
			& \frac{5 \pi}{4}=\frac{5 \pi}{4} \cdot\left(\frac{180}{\pi}\right)^{\circ}=225^{\circ} .
		\end{aligned}
		$$
	}
\end{vd}

% #1897 | Archive/2024/SP-K10-K11-LyVanHoang_TVN-023/11. C1-LuongGiac-lvhoang_TVN-023/dataC1K11/Bai1-Dang1.tex | mã gốc 1D1H1-2
\begin{vd}[Chân trời sáng tạo, Mức độ 2]%[1D1H1-2]
	Đổi các số đo góc sau đây từ radian sang độ hoặc ngược lại
	\begin{listEX}[3]
		\item $-60^{\circ}$.
		\item $\dfrac{2 \pi}{5}\, \mathrm{rad}$.
		\item $3\, \mathrm{rad}$.
	\end{listEX}
	\loigiai{
		\begin{listEX}[1]
			\item $-60^{\circ}=-\dfrac{60 \pi}{180}\, \mathrm{rad}=-\dfrac{\pi}{3}\, \mathrm{rad}$.
			\item $\dfrac{2 \pi}{5}\, \mathrm{rad}=\left(\dfrac{2 \pi}{5}\cdot \dfrac{180}{\pi}\right)^{\circ}=72^{\circ}$.
			\item $3\, \mathrm{rad}=\left(3 \cdot \dfrac{180}{\pi}\right)^{\circ}=\left(\dfrac{540}{\pi}\right)^{\circ} \approx 171{,}89^{\circ}$.
		\end{listEX}
	}
\end{vd}

% #1900 | Archive/2024/SP-K10-K11-LyVanHoang_TVN-023/11. C1-LuongGiac-lvhoang_TVN-023/dataC1K11/Bai1-Dang1.tex | mã gốc 1D1H1-2
\begin{bt}%[1D1H1-2]
	Đổi số đo của các góc sau đây sang radian:
	\begin{listEX}[3]
		\item $38^{\circ}$;
		\item $-115^{\circ}$;
		\item $\left(\dfrac{3}{\pi}\right)^{\circ}$.
	\end{listEX}
	\loigiai{
		\begin{listEX}[1]
			\item $38^{\circ}=\dfrac{38\pi}{180}=\dfrac{19\pi}{90}$.
			\item $-115^{\circ}=\dfrac{-115\pi}{180}=-\dfrac{23\pi}{36}$;
			\item $\left(\dfrac{3}{\pi}\right)^{\circ}=\dfrac{3\pi
			}{\pi\cdot 180}=\dfrac{1}{60}$.
		\end{listEX}
		
	}
	
\end{bt}

% #1901 | Archive/2024/SP-K10-K11-LyVanHoang_TVN-023/11. C1-LuongGiac-lvhoang_TVN-023/dataC1K11/Bai1-Dang1.tex | mã gốc 1D1H1-2
\begin{bt}%[1D1H1-2]
	Đổi số đo của các góc sau đây sang độ:
	\begin{listEX}[3]
		\item $\dfrac{\pi}{12}$;
		\item $-5$;
		\item $\dfrac{13 \pi}{9}$.
	\end{listEX}
	\loigiai{
		\begin{listEX}[1]
			\item $\dfrac{\pi}{12}=\left(\dfrac{\pi\cdot 180}{12\cdot \pi}\right)^\circ=15^\circ$;
			\item $-5=\left(\dfrac{-5\cdot 180}{\pi}\right)^\circ=\left(\dfrac{-900}{\pi}\right)^\circ$;
			\item $\dfrac{13 \pi}{9}=\left(\dfrac{13 \pi\cdot 180}{9\cdot \pi}\right)=260^\circ$.
		\end{listEX}
		
	}
	
\end{bt}

% #1902 | Archive/2024/SP-K10-K11-LyVanHoang_TVN-023/11. C1-LuongGiac-lvhoang_TVN-023/dataC1K11/Bai1-Dang1.tex | mã gốc 1D1H1-2
\begin{bt}%[1D1H1-2]
	Đổi số đo các cung $\dfrac{13\pi}{3}$; $-\dfrac{3\pi}{4}$; $2$; $\dfrac{15\pi}{7}$; $\dfrac{\pi}{15}$ ra độ, phút, giây.
	\dapso{$780^\circ$; $-135^\circ$; $114^\circ 	35' 30''$; $385^\circ 42'51''$; $12^\circ$}
	\loigiai
	{
		Ta có
		\begin{listEX}[3]
			\item $\dfrac{13\pi}{3} = \left(\dfrac{13\cdot 180}{3}\right)^\circ = 780^\circ$.
			\item $-\dfrac{3\pi}{4} = -\left(\dfrac{3\cdot 180}{4}\right)^\circ = -135^\circ$.
			\item $2 = \left(\dfrac{2\cdot 180}{\pi}\right)^\circ \approx 114^\circ 	35' 30''$.
			\item $\dfrac{15\pi}{7} = \left(\dfrac{15\cdot 180}{7}\right)^\circ \approx 385^\circ 42'51''$.
			\item $\dfrac{\pi}{15} = \left(\dfrac{1\cdot 180}{15}\right)^\circ = 12^\circ$.
		\end{listEX}
	}
\end{bt}

% #1903 | Archive/2024/SP-K10-K11-LyVanHoang_TVN-023/11. C1-LuongGiac-lvhoang_TVN-023/dataC1K11/Bai1-Dang1.tex | mã gốc 1D1H1-2
\begin{bt}%[1D1H1-2]
	Đổi số đo các góc có số đo $270^\circ$; $240^\circ$; $540^\circ$; $-750^\circ$; $21^\circ$ ra đơn vị ra-đi-an (rad).
	\dapso{$\dfrac{3\pi}{2}$; $\dfrac{4\pi}{3}$; $3\pi$; $-\dfrac{25\pi}{6}$; $\dfrac{7\pi}{60}$}
	\loigiai
	{
		Ta có
		\begin{listEX}[3]
			\item $270^\circ = 270\cdot \dfrac{\pi}{180} = \dfrac{3\pi}{2}$.
			\item $240^\circ = 240\cdot \dfrac{\pi}{180} = \dfrac{4\pi}{3}$.
			\item $540^\circ = 540\cdot \dfrac{\pi}{180} = 3\pi$.
			\item $-750^\circ = -750\cdot \dfrac{\pi}{180} = -\dfrac{25\pi}{6}$.
			\item $21^\circ = 21\cdot \dfrac{\pi}{180} = \dfrac{7\pi}{60}$.
		\end{listEX}
	}
\end{bt}

% #1904 | Archive/2024/SP-K10-K11-LyVanHoang_TVN-023/11. C1-LuongGiac-lvhoang_TVN-023/dataC1K11/Bai1-Dang1.tex | mã gốc 1D1H1-2
\begin{bt}%[1D1H1-2]
	Đổi số đo của các cung tròn có số đo $12^\circ 36'$; $15^\circ 30'$; $27^\circ 38'49''$ ra số đo ra-đi-an (chính xác đến phần nghìn).
	\dapso{$0{,}220$; $0{,}271$; $0{,}500$}
	\loigiai
	{
		Ta có
		\begin{enumerate}[$\circ$]
			\item $12^\circ 36' = \left(\dfrac{63}{5}\right)^\circ = \dfrac{63}{5}\cdot \dfrac{\pi}{180} \approx 0{,}220$.
			\item $15^\circ 30' = \left(\dfrac{31}{2}\right)^\circ = \dfrac{31}{2}\cdot \dfrac{\pi}{180} \approx 0{,}271$.
			\item $27^\circ 38'49'' = \left(\dfrac{103129}{3600}\right)^\circ = \dfrac{103129}{3600}\cdot \dfrac{\pi}{180} \approx 0{,}500$.
		\end{enumerate}
	}
\end{bt}

% #1905 | Archive/2024/SP-K10-K11-LyVanHoang_TVN-023/11. C1-LuongGiac-lvhoang_TVN-023/dataC1K11/Bai1-Dang1.tex | mã gốc 1D1H1-2
\begin{bt}%[1D1H1-2]
	Điền các giá trị thích hợp vào ô trống.
	\begin{center}
		\renewcommand{\arraystretch}{2}
		\begin{tabular}{|>{\centering\arraybackslash}m{1cm}|>{\centering\arraybackslash}m{1cm}|>{\centering\arraybackslash}m{1cm}|>{\centering\arraybackslash}m{1cm}|>{\centering\arraybackslash}m{1cm}|>{\centering\arraybackslash}m{1cm}|>{\centering\arraybackslash}m{1cm}|>{\centering\arraybackslash}m{1cm}|>{\centering\arraybackslash}m{1cm}|>{\centering\arraybackslash}m{1cm}|}
			\hline
			$0^\circ$ & $15^\circ$ & $30^\circ$ & $45^\circ$ & $60^\circ$ & $75^\circ$ & $90^\circ$ & $120^\circ$ & $150^\circ$ & $180^\circ$\\
			\hline
			&&&&&&&&&\\
			\hline
		\end{tabular}
	\end{center}
	\begin{center}
		\renewcommand{\arraystretch}{2}
		\begin{tabular}{|>{\centering\arraybackslash}m{1cm}|>{\centering\arraybackslash}m{1cm}|>{\centering\arraybackslash}m{1cm}|>{\centering\arraybackslash}m{1cm}|>{\centering\arraybackslash}m{1cm}|>{\centering\arraybackslash}m{1cm}|>{\centering\arraybackslash}m{1cm}|>{\centering\arraybackslash}m{1cm}|>{\centering\arraybackslash}m{1cm}|>{\centering\arraybackslash}m{1cm}|}
			\hline
			$\dfrac{\pi}{16}$ & $\dfrac{\pi}{8}$ & $\dfrac{\pi}{6}$ & $\dfrac{5\pi}{12}$ & $\dfrac{\pi}{2}$ & $\dfrac{2\pi}{3}$ & $\dfrac{3\pi}{4}$ & $\dfrac{5\pi}{6}$ & $\pi$ & $\dfrac{4\pi}{3}$\\
			\hline
			&&&&&&&&&\\
			\hline
		\end{tabular}
	\end{center}
	\begin{center}
		\renewcommand{\arraystretch}{2}
		\begin{tabular}{|>{\centering\arraybackslash}m{2cm}|>{\centering\arraybackslash}m{2cm}|>{\centering\arraybackslash}m{2cm}|>{\centering\arraybackslash}m{2cm}|>{\centering\arraybackslash}m{2cm}|>{\centering\arraybackslash}m{2cm}|>{\centering\arraybackslash}m{2cm}|}
			\hline
			Ra-đi-an & $\dfrac{7\pi}{4}$ & $\dfrac{\pi}{12}$ & $-\dfrac{5\pi}{6}$ & & & \\
			\hline
			Độ & & & & $135^\circ$ & $-50^\circ$ & $810^\circ$\\
			\hline
		\end{tabular}
	\end{center}
	\loigiai
	{
		\begin{center}
			\renewcommand{\arraystretch}{2}
			\begin{tabular}{|>{\centering\arraybackslash}m{1cm}|>{\centering\arraybackslash}m{1cm}|>{\centering\arraybackslash}m{1cm}|>{\centering\arraybackslash}m{1cm}|>{\centering\arraybackslash}m{1cm}|>{\centering\arraybackslash}m{1cm}|>{\centering\arraybackslash}m{1cm}|>{\centering\arraybackslash}m{1cm}|>{\centering\arraybackslash}m{1cm}|>{\centering\arraybackslash}m{1cm}|}
				\hline
				$0^\circ$ & $15^\circ$ & $30^\circ$ & $45^\circ$ & $60^\circ$ & $75^\circ$ & $90^\circ$ & $120^\circ$ & $150^\circ$ & $180^\circ$\\
				\hline
				$0$ & $\dfrac{\pi}{12}$ & $\dfrac{\pi}{6}$ & $\dfrac{\pi}{4}$ & $\dfrac{\pi}{3}$ & $\dfrac{5\pi}{12}$ & $\dfrac{\pi}{2}$ & $\dfrac{2\pi}{3}$ & $\dfrac{5\pi}{6}$ & $\pi$\\
				\hline
			\end{tabular}
		\end{center}
		\begin{center}
			\renewcommand{\arraystretch}{2}
			\begin{tabular}{|>{\centering\arraybackslash}m{1cm}|>{\centering\arraybackslash}m{1cm}|>{\centering\arraybackslash}m{1cm}|>{\centering\arraybackslash}m{1cm}|>{\centering\arraybackslash}m{1cm}|>{\centering\arraybackslash}m{1cm}|>{\centering\arraybackslash}m{1cm}|>{\centering\arraybackslash}m{1cm}|>{\centering\arraybackslash}m{1cm}|>{\centering\arraybackslash}m{1cm}|}
				\hline
				$\dfrac{\pi}{16}$ & $\dfrac{\pi}{8}$ & $\dfrac{\pi}{6}$ & $\dfrac{5\pi}{12}$ & $\dfrac{\pi}{2}$ & $\dfrac{2\pi}{3}$ & $\dfrac{3\pi}{4}$ & $\dfrac{5\pi}{6}$ & $\pi$ & $\dfrac{4\pi}{3}$\\
				\hline
				$11{,}25^\circ$ & $22{,}5^\circ$ & $30^\circ$ & $75^\circ$ & $90^\circ$ & $120^\circ$ & $135^\circ$ & $150^\circ$ & $180^\circ$ & $240^\circ$\\
				\hline
			\end{tabular}
		\end{center}
		\begin{center}
			\renewcommand{\arraystretch}{2}
			\begin{tabular}{|>{\centering\arraybackslash}m{2cm}|>{\centering\arraybackslash}m{2cm}|>{\centering\arraybackslash}m{2cm}|>{\centering\arraybackslash}m{2cm}|>{\centering\arraybackslash}m{2cm}|>{\centering\arraybackslash}m{2cm}|>{\centering\arraybackslash}m{2cm}|}
				\hline
				Ra-đi-an & $\dfrac{7\pi}{4}$ & $\dfrac{\pi}{12}$ & $-\dfrac{5\pi}{6}$ & $\dfrac{3\pi}{4}$ & $-\dfrac{5\pi}{18}$ & $\dfrac{9\pi}{2}$\\
				\hline
				Độ & $315^\circ$ & $15^\circ$ & $-150^\circ$ & $135^\circ$ & $-50^\circ$ & $810^\circ$\\
				\hline
			\end{tabular}
		\end{center}
	}
\end{bt}

% #1921 | Archive/2024/SP-K10-K11-LyVanHoang_TVN-023/11. C1-LuongGiac-lvhoang_TVN-023/dataC1K11/Chuong1-Ontap-De01.tex | mã gốc 1D1K3-1
\begin{ex}%[Mức V]%[24-25- Bài Giảng K10-K11, Bùi Mạnh Tiến]%[1D1K3-1]
	Giả sử một cái xích đu dao động điều hoà xung quanh vị trí cân bằng theo phương trình $x=60\cos \left(5t-\dfrac{\pi}{6}\right)$. Trong đó thời gian $t$ tính bằng giây và quãng đường $S$ tính bằng cm. Hãy cho biết trong khoảng thời gian từ $0$ đến $6$ giây, xích đu đó đi qua vị trí cân bằng bao nhiêu lần? \shortans{$9$}
	\loigiai{
		Khi vật đi qua vị trí cân bằng thì 
		\begin{eqnarray*}
			x=0&\Leftrightarrow & 60\cos \left(5t-\dfrac{\pi}{6}\right)=0\\
			&\Leftrightarrow & \cos \left(5t-\dfrac{\pi}{6}\right)=0\\
			&\Leftrightarrow& 5t-\dfrac{\pi}{6}=\dfrac{\pi}{2}+k\pi\\
			&\Leftrightarrow & t=\dfrac{2\pi}{15}+\dfrac{k\pi}{5} \text{ } (k\in \mathbb{Z}).
		\end{eqnarray*}
		Trong khoảng thời gian từ $0$ đến $6$ giây thì
		\begin{eqnarray*}
			0\leq t\leq 6&\Leftrightarrow & 0< \dfrac{2\pi}{15}+\dfrac{k\pi}{5}< 6\\
			&\Leftrightarrow & -\dfrac{2}{3}< k< \left(\dfrac{6}{\pi}-\dfrac{2}{15}\right)\cdot 5\approx 8{,}88. 
		\end{eqnarray*}
		Vì $k\in \mathbb{Z}$ nên $k\in \left\{0;1;2;3;4;5;6;7;8\right\}$.\\
		Vậy trong khoảng từ $0$ đến $6$ giây thì xích đu đi qua vị trí cân bằng $9$ lần.
	}
\end{ex}

% #1922 | Archive/2024/SP-K10-K11-LyVanHoang_TVN-023/11. C1-LuongGiac-lvhoang_TVN-023/dataC1K11/Bai2-Dang1-Congthuccong.tex | mã gốc 1D1N3-2
\begin{vd}%[Cánh Diều]%[Đỗ Vũ Minh Thắng, 24-25 BG 10 11]%[1D1N3-2]
	Tính
	\begin{enumEX}{2}
		\item $\sin 75^{\circ}$.
		\item $\sin \dfrac{\pi}{12}$.
	\end{enumEX}
	\loigiai{
		\begin{listEX}[1]
			\item Áp dụng công thức cộng, ta có\\
			$\sin 75^{\circ}=\sin \left(30^{\circ}+45^{\circ}\right)=\sin 30^{\circ} \cos 45^{\circ}+\cos 30^{\circ} \sin 45^{\circ}=\dfrac{\sqrt{6}+\sqrt{2}}{4}$.
			\item Áp dụng công thức cộng, ta có\\
			$\sin \dfrac{\pi}{12}=\sin \left(\dfrac{\pi}{3}-\dfrac{\pi}{4}\right)=\sin \dfrac{\pi}{3} \cos \dfrac{\pi}{4}-\cos \dfrac{\pi}{3} \sin \dfrac{\pi}{4}=\dfrac{\sqrt{6}-\sqrt{2}}{4}$.
		\end{listEX}
	}
\end{vd}

% #1923 | Archive/2024/SP-K10-K11-LyVanHoang_TVN-023/11. C1-LuongGiac-lvhoang_TVN-023/dataC1K11/Bai2-Dang1-Congthuccong.tex | mã gốc 1D1N3-2
\begin{vd}%[Cánh Diều]%[Đỗ Vũ Minh Thắng, 24-25 BG 10 11]%[1D1N3-2]
	Tính
	\begin{enumEX}{2}
		\item $\cos \dfrac{5\pi}{12}$.
		\item $\cos 15^{\circ}$.
	\end{enumEX}
	\loigiai{
		\begin{listEX}[1]
			\item Áp dụng công thức cộng, ta có\\
			$\cos \dfrac{5 \pi}{12}=\cos \left(\dfrac{\pi}{6}+\dfrac{\pi}{4}\right)=\cos \dfrac{\pi}{6} \cos \dfrac{\pi}{4}-\sin \dfrac{\pi}{6} \sin \dfrac{\pi}{4}=\dfrac{\sqrt{6}-\sqrt{2}}{4}$.
			\item Áp dụng công thức cộng, ta có\\
			$\sin 15^{\circ}=\sin \left(60^{\circ}-45^{\circ}\right)=\sin 60^{\circ} \cos 45^{\circ}+\cos 60^{\circ} \sin 45^{\circ}=\dfrac{\sqrt{6}+\sqrt{2}}{4}$.
		\end{listEX}
	}
\end{vd}

% #1924 | Archive/2024/SP-K10-K11-LyVanHoang_TVN-023/11. C1-LuongGiac-lvhoang_TVN-023/dataC1K11/Bai2-Dang1-Congthuccong.tex | mã gốc 1D1N3-2
\begin{vd}%[Cánh Diều]%[Đỗ Vũ Minh Thắng, 24-25 BG 10 11]%[1D1N3-2]
	Tính
	\begin{enumEX}{2}
		\item $\tan \dfrac{7 \pi}{12}$.
		\item $\tan 165^{\circ}$.
	\end{enumEX}
	\loigiai{
		\begin{listEX}[1]
			\item Áp dụng công thức cộng đối với tang, ta có\\
			$\tan \dfrac{7 \pi}{12}=\tan \left(\dfrac{\pi}{4}+\dfrac{\pi}{3}\right)=\dfrac{\tan \frac{\pi}{4}+\tan \frac{\pi}{3}}{1-\tan \frac{\pi}{4} \tan \frac{\pi}{3}}=\dfrac{1+\sqrt{3}}{1-\sqrt{3}}=-2-\sqrt{3}$.
			\item Áp dụng công thức cộng đối với tang, ta có\\
			$\tan 15^{\circ}=\tan \left(45^{\circ}-30^{\circ}\right)=\dfrac{\tan 45^{\circ}-\tan 30^{\circ}}{1+\tan 45^{\circ} \tan 30^{\circ}}=\dfrac{1-\frac{\sqrt{3}}{3}}{1+\frac{\sqrt{3}}{3}}=2-\sqrt{3}$.
		\end{listEX}
	}
\end{vd}

% #1927 | Archive/2024/SP-K10-K11-LyVanHoang_TVN-023/11. C1-LuongGiac-lvhoang_TVN-023/dataC1K11/Bai2-Dang1-Congthuccong.tex | mã gốc 1D1N3-2
\begin{bt}%[Cánh Diều]%[Đỗ Vũ Minh Thắng, 24-25 BG 10 11]%[1D1N3-2]
	Tính
	\begin{multicols}{3}
		\begin{enumerate}
			\item $\sin \dfrac{\pi}{12}$;
			\item $\cos 15^\circ$;
			\item $\tan 165^\circ$.
		\end{enumerate}
	\end{multicols}	
\end{bt}

% #1928 | Archive/2024/SP-K10-K11-LyVanHoang_TVN-023/11. C1-LuongGiac-lvhoang_TVN-023/dataC1K11/Bai2-Dang1-Congthuccong.tex | mã gốc 1D1N3-2
\begin{bt}%[CTST]%[Đỗ Vũ Minh Thắng, 24-25 BG 10 11]%[1D1N3-2]
	Tính $\sin \dfrac{\pi}{12}$ và $\tan \dfrac{\pi}{12}$.
	\loigiai{
		\begin{itemize}
			\item Ta có $\sin \dfrac{\pi}{12} = \sin \left(\dfrac{\pi}{3} - \dfrac{\pi}{4}\right) = \sin \dfrac{\pi}{3} \cos \dfrac{\pi}{4} - \cos \dfrac{\pi}{3} \sin \dfrac{\pi}{4} = \dfrac{\sqrt{3}}{2} \cdot \dfrac{\sqrt{2}}{2} - \dfrac{1}{2} \cdot \dfrac{\sqrt{2}}{2} = \dfrac{\sqrt{6} - \sqrt{2}}{4}$.
			\item Ta có $\tan \dfrac{\pi}{12} = \tan \left(\dfrac{\pi}{3} - \dfrac{\pi}{4}\right) = \dfrac{\tan \dfrac{\pi}{3} - \tan \dfrac{\pi}{4} }{1 + \tan \dfrac{\pi}{3} \tan \dfrac{\pi}{4}} = \dfrac{\sqrt{3} -1}{1 + \sqrt{3}}$.
		\end{itemize}
	} 
\end{bt}

% #1929 | Archive/2024/SP-K10-K11-LyVanHoang_TVN-023/11. C1-LuongGiac-lvhoang_TVN-023/dataC1K11/Bai2-Dang1-Congthuccong.tex | mã gốc 1D1N3-2
\begin{bt}%[CTST]%[Đỗ Vũ Minh Thắng, 24-25 BG 10 11]%[1D1N3-2]
	Tính $\sin \dfrac{7\pi}{12}$ và $\tan \dfrac{\pi}{12}$.
	\loigiai{
		\begin{itemize}
			\item Ta có $\sin \dfrac{7\pi}{12} = \sin \left(\dfrac{\pi}{3} + \dfrac{\pi}{4}\right) = \sin \dfrac{\pi}{3} \cos \dfrac{\pi}{4} + \cos \dfrac{\pi}{3} \sin \dfrac{\pi}{4} = \dfrac{\sqrt{3}}{2} \cdot \dfrac{\sqrt{2}}{2} + \dfrac{1}{2} \cdot \dfrac{\sqrt{2}}{2} = \dfrac{\sqrt{6} + \sqrt{2}}{4}$.
			\item Ta có $\tan \dfrac{\pi}{12} = \tan \left(\dfrac{\pi}{3} - \dfrac{\pi}{4}\right) = \dfrac{\tan \dfrac{\pi}{3} - \tan \dfrac{\pi}{4} }{1 + \tan \dfrac{\pi}{3} \tan \dfrac{\pi}{4}} = \dfrac{\sqrt{3} -1}{1 + \sqrt{3}}$.
		\end{itemize}
	} 
\end{bt}

% #1930 | Archive/2024/SP-K10-K11-LyVanHoang_TVN-023/11. C1-LuongGiac-lvhoang_TVN-023/dataC1K11/Bai2-Dang1-Congthuccong.tex | mã gốc 1D1N3-2
\begin{bt}%[CTST]%[Đỗ Vũ Minh Thắng, 24-25 BG 10 11]%[1D1N3-2]
	Không dùng máy tính cầm tay, tính các giá trị lượng giác của các góc:
	\begin{listEX}[2] 
		\item[a)] $\dfrac{5 \pi}{12}$.
		\item[b)] $-555^{\circ}$.
	\end{listEX}
	\loigiai{
		\begin{itemize}
			\item[a)] Ta có $\sin \dfrac{5 \pi}{12}=\sin\left( \dfrac{3\pi}{12}+\dfrac{2\pi}{12} \right)=\sin\left(\dfrac{\pi}{4}+\dfrac{\pi}{6}\right)=\sin\dfrac{\pi}{4}\cos\dfrac{\pi}{6}+\cos\dfrac{\pi}{4} \sin\dfrac{\pi}{6}=\dfrac{\sqrt{6}+\sqrt{2}}{4}$.\\
			Ta có $\cos \dfrac{5 \pi}{12}=\cos\left(\dfrac{\pi}{4}+\dfrac{\pi}{6}\right)=\cos\dfrac{\pi}{4}\cos\dfrac{\pi}{6}-\sin\dfrac{\pi}{4} \sin\dfrac{\pi}{6}=\dfrac{\sqrt{6}-\sqrt{2}}{4}$.\\
			Suy ra $\tan\dfrac{5\pi}{12}=\dfrac{\sin\dfrac{5\pi}{12}}{\cos\dfrac{5\pi}{12}}=\dfrac{\sqrt 6+\sqrt 2}{\sqrt 6-\sqrt 2}\Rightarrow \cot\dfrac{5\pi}{12}=\dfrac{\sqrt 6-\sqrt 2}{\sqrt 6+\sqrt 2}$.
			\item[b)] Ta có $-555^{\circ}=-2\cdot360^\circ+165^\circ$. 
			Suy ra \begin{eqnarray*}
				\sin(-555^\circ)&=&\sin 165^\circ \\
				&=&\sin 15^\circ=\sin\left(45^\circ-30^\circ\right)\\
				&=&\sin 45^\circ\cos 30^\circ-\sin 30^\circ \cos 45^\circ\\
				&=&\dfrac{\sqrt{6}-\sqrt{2}}{4};
			\end{eqnarray*}
			\begin{eqnarray*}
				\cos(-555^\circ)&=&\cos 165^\circ \\
				&=&-\cos 15^\circ=-\cos\left(45^\circ-30^\circ\right)\\
				&=&-\left( \cos 45^\circ\cos 30^\circ+\sin 30^\circ \sin 45^\circ\right) \\
				&=&-\dfrac{\sqrt{6}+\sqrt{2}}{4}.
			\end{eqnarray*}
		\end{itemize}
		Suy ra $\tan (-555^\circ)=\dfrac{\sin (-555^\circ)}{\cos(-555^\circ)}=-\dfrac{\sqrt6-\sqrt2}{\sqrt6+\sqrt2}\Rightarrow \cot(-555^\circ)=-\dfrac{\sqrt6+\sqrt2}{\sqrt6-\sqrt2}$.
	} 
\end{bt}

% #1931 | Archive/2024/SP-K10-K11-LyVanHoang_TVN-023/11. C1-LuongGiac-lvhoang_TVN-023/dataC1K11/Bai2-Dang1-Congthuccong.tex | mã gốc 1D1V3-2
\begin{bt}%[CTST]%[Đỗ Vũ Minh Thắng, 24-25 BG 10 11]%[1D1V3-2]
	Trong Hình bên dưới, ba điểm $M$, $N$, $P$ nằm ở đầu các cánh quạt của tua-bin gió. Biết các cánh quạt dài $31$ m, độ cao của điểm $M$ so với mặt đất là $30$ m, góc giữa các cánh quạt là $\dfrac{2 \pi}{3}$ và số đo góc $(OA, OM)$ là $\alpha$.
	\begin{itemize}
		\item[a)] Tính $\sin \alpha$ và $\cos \alpha$.
		\item[b)] Tính $\sin$ của các góc lượng giác $(OA, ON)$ và $(OA, OP)$, từ đó tính chiều cao của các điểm $N$ và $P$ so với mặt đất (theo đơn vị mét). Làm tròn kết quả đến hàng phần trăm.
	\end{itemize}
	\begin{minipage}{0.5\textwidth}
		\centering
		\tikzset{
			ex_markstyle/.style={},
			ex_mark/.style  n args={1}{decoration={ markings, %
					mark= at position 0.5 with
					with{
						\ifnum#1=1
						\draw[ex_markstyle] (0pt,-2pt) -- (0pt,2pt);
						\fi
						\ifnum#1=2
						\draw[ex_markstyle] (-1pt,-2pt) -- (-1pt,2pt);
						\draw[ex_markstyle] (1pt,-2pt) -- (1pt,2pt);
						\fi
						\ifnum#1=3
						\draw[ex_markstyle] (-2pt,-2pt) -- (-2pt,2pt);
						\draw[ex_markstyle] (0pt,-2pt) -- (0pt,2pt);
						\draw[ex_markstyle] (2pt,-2pt) -- (2pt,2pt);
						\fi
						\ifnum#1=4
						\draw[ex_markstyle] (-1pt,-1pt) -- (1pt,1pt);
						\draw[ex_markstyle] (-1pt,1pt) -- (1pt,-1pt);
						\fi
				} },
				pic actions/.append code=\tikzset{postaction=decorate}},
		}
		\definecolor{deepskyblue}{rgb}{0.0, 0.75, 1.0}
		\definecolor{green(ryb)}{rgb}{0.4, 0.69, 0.2}
		\definecolor{anti-flashwhite}{rgb}{0.95, 0.95, 0.96}
		\definecolor{bittersweet}{rgb}{1.0, 0.44, 0.37}
		\begin{tikzpicture}[scale=1, font=\footnotesize, line join=round, line cap=round,>=stealth]
			%\draw[gray!50] (-3,-3) grid (4,4);
			
			\fill[bottom color=white,top color=deepskyblue!90, middle color=deepskyblue!90] (-3,0) rectangle (4,4);  
			
			\fill[green(ryb)] (-3,-3) rectangle (4,0);  
			
			\tikzset{quat_gio/.pic={
					\coordinate (O) at (0,0);
					\coordinate (N) at ($(O)+(160:1.5)$);
					\coordinate (P) at ($(O)+(40:1.5)$);
					\coordinate (M) at ($(O)+(-70:1.5)$);
					
					\draw[fill=anti-flashwhite] (P)..controls +(-155:.3) and +(65:.4) ..  (O)--cycle;				
					\draw[fill=anti-flashwhite] (N)..controls +(-40:.3) and +(-140:.2) ..  (O)--cycle;
					\draw[fill=anti-flashwhite] (O)..controls +(-30:.3) and +(100:.3) ..  (M)--cycle;
					\draw[fill=anti-flashwhite] (O)..controls +(-95:.3) and +(70:.3) ..  (-.15,-2.8)--(.15,-2.8)
					..controls +(110:.3) and +(-80:.3) ..  (O)
					;
			}}
	
\path (-1.7,1)pic[scale=.7,yscale=.85,yslant=-.05]{quat_gio}
(0,1.5)pic[scale=.8,yscale=1,yslant=-.08]{quat_gio}
(3,1.1)pic[scale=.6,yscale=.8,yslant=-.08]{quat_gio}
(2,1.35)pic[scale=1,yscale=1.2,yslant=-.1]{quat_gio};
		\end{tikzpicture}
		
	\end{minipage}%
	\hspace{0.5cm} % Khoảng cách giữa hai hình
	\begin{minipage}{0.5\textwidth}
		\centering
		\begin{tikzpicture}[scale=1.2, font=\footnotesize, line join=round, line cap=round,>=stealth]
			\def\a{1/16}  
			\def\xmin{-2.5} \def\xmax{4}
			\def\ymin{-3.1} \def\ymax{2} 
			\draw[->] (\xmin,0)--(\xmax,0) node [below]{$x$};
			\draw[->] (0,\ymin)--(0,\ymax) node [left]{$y$};
			\coordinate[label=above left:$O$] (O) at (0,0);
			\coordinate[label=above left:$N$] (N) at ($(O)+(160:1.5)$);
			\coordinate[label=above right:$P$] (P) at ($(O)+(40:1.5)$);
			\coordinate[label=below:$M$] (M) at ($(O)+(-70:1.5)$);
			\coordinate[label=above:$A$] (A) at ($(O)+(0:1.5)$);
			\draw pic["$\dfrac{2\pi}{3}$", draw=black,below, angle eccentricity=1.4, angle radius=.5cm]{angle=M--O--P}; %góc
			\draw pic["$\dfrac{2\pi}{3}$", draw=black,angle eccentricity=1.8, angle radius=.4cm]{angle=N--O--M}; %góc
			\draw pic["$\alpha$", draw=black, angle eccentricity=1.6, angle radius=.2cm]{angle=M--O--A}; 
			\draw[dashed] (M)--($(M)+(0:1.8)$) 
			(0,-2.8)--($(0,-2.8)+(0:3)$);
			\coordinate (E) at ($(M)+(0:1.2)$);
			\draw[<->](E)--($(E)+(-90:1.4)$); 
			\draw[<->](2.8,0)--(2.8,-2.8);
			\draw[fill=gray] (P)..controls +(-155:.3) and +(65:.4) ..  (O)--cycle;				\draw[fill=gray] (N)..controls +(-40:.3) and +(-140:.2) ..  (O)--cycle;
			\draw[fill=gray] (O)..controls +(-30:.3) and +(100:.3) ..  (M)--cycle;
			\draw[fill=gray] (O)..controls +(-95:.3) and +(70:.3) ..  (-.15,-2.8)--(.15,-2.8)
			..controls +(110:.3) and +(-80:.3) ..  (O)
			;
			\node at (2.8,-1.5) [right]{$60$ m};
			\node at (1.7,-2.1) [right]{$30$ m};
		\end{tikzpicture}
	\end{minipage}
	\loigiai{
		\begin{itemize}
			\item[a)] Ta có tọa độ điểm $M$ trong hệ trục tọa độ $Oxy$ là $M(x_M; y_M)$.\\
			Với $\sin \alpha = \dfrac{y_M}{OM} = -\dfrac{30}{31}$ và $\cos \alpha =\sqrt{1 - \sin^2 \alpha} = \dfrac{\sqrt{61}}{31}$.
			\item[b)] 		
			Ta có $(OA, OP)=(OA,OM)+(OM,OP)=\alpha+\dfrac{2\pi}{3}$.\\ Suy ra \begin{eqnarray*}
				\sin (OA,OP)=\sin\left(\alpha+\dfrac{2\pi}{3}\right)&=&\sin\alpha\cos\dfrac{2\pi}{3}+\sin\dfrac{2\pi}{3}\cos\alpha\\
				&=&-\dfrac{30}{31}\cdot \left(-\dfrac{1}{2}\right)+\dfrac{\sqrt{3}}{2}\cdot\dfrac{\sqrt{61}}{31}=\dfrac{30+\sqrt{183}}{62}.
			\end{eqnarray*}
			Do đó chiều cao của điểm $P$ so với mặt đất là $60+31\sin(OA,OP)=60+31\cdot \dfrac{30+\sqrt{183}}{62}\approx 81{,}76$ m.\\
			Tương tự, ta có $(OA, ON)=(OA,OM)+(OM,ON)=\alpha+\dfrac{4\pi}{3}$. \\
			Suy ra \begin{eqnarray*}
				\sin (OA,ON)=\sin\left(\alpha+\dfrac{4\pi}{3}\right)&=&\sin\alpha\cos\dfrac{4\pi}{3}+\sin\dfrac{4\pi}{3}\cos\alpha\\
				&=&-\dfrac{30}{31}\cdot \left(-\dfrac{1}{2}\right)+\left(-\dfrac{\sqrt{3}}{2}\right) \cdot\dfrac{\sqrt{61}}{31}=\dfrac{30-\sqrt{183}}{62}.
			\end{eqnarray*}
			Do đó chiều cao của điểm $N$ so với mặt đất là $60+31\sin(OA,ON)=60+31\cdot \dfrac{30-\sqrt{183}}{62}\approx 68{,}24$ m.
		\end{itemize}
	} 
\end{bt}

% #1937 | Archive/2024/SP-K10-K11-LyVanHoang_TVN-023/11. C1-LuongGiac-lvhoang_TVN-023/dataC1K11/Bai2-Dang1-Congthuccong.tex | mã gốc 1D1V3-2
\begin{ex}%[TeX hóa SGK CTST]%[Đỗ Vũ Minh Thắng, 24-25 BG 10 11]%[1D1V3-2]
	\immini{Trong hình bên, tam giác $ABC$ vuông tại $B$ và có hai cạnh góc vuông là $AB = 4$, $BC = 3$. Vẽ điểm $D$ nằm trên tia đối của tia $CB$ thoả mãn $\widehat{CAD} = 30^\circ$. Tính $\tan \widehat{BAD}$ (kết quả làm tròn đến hàng phần trăm).
		\shortans{2,34}
	}{\begin{tikzpicture}[line join = round, line cap = round, scale=.6] 
		\coordinate[label = left:$B$] (B) at (0,0); 
		\coordinate[label = below:$C$] (C) at ($(B)+(2,0)$); 
		\coordinate[label = above:$A$] (A) at ($(B)+(0,4)$); 
		\coordinate[label = right:$D$] (D) at ($(B)+(6,0)$); 
		\draw[fill = gray!50] ($(B)!5pt!(C)$)--($(B)!5pt!(C)+(B)!5pt!(A)-(B)$)--($(B)!5pt!(A)$)--(B)--cycle;
		\draw (B)--(D)--(A)--cycle (A)--(C); 
		\foreach \x in {B,C,A,D} \fill[black] (\x) circle (1.5pt);
		\path (A)--(B) node[midway, left]{4} (B)--(C) node[midway, below]{3};
		\draw (1,3.42) node [below] {\tiny$30^\circ$};	
		\end{tikzpicture} 
	}
	\loigiai{
		\immini{Trong $\triangle ABC$ vuông tại $B$ ta có $AC = \sqrt{AB^2 + BC^2} =5$.\\
			Mặt khác $\tan \widehat{BAC} = \dfrac{BC}{AB} = \dfrac{3}{4}$.\\
			Suy ra
			\allowdisplaybreaks
			\begin{eqnarray*}
				\tan \widehat{BAD}&=&\tan \left(\widehat{BAC} + \widehat{CAD}\right)\\
				&=&\dfrac{\tan \widehat{BAC} + \tan 30^\circ}{1 - \tan \widehat{BAC} \cdot  \tan 30^\circ}\\
				&=&\dfrac{\dfrac{3}{4} + \dfrac{\sqrt{3}}{3}}{1 - \dfrac{3}{4} \cdot  \dfrac{\sqrt{3}}{3}} = \dfrac{9 + 4\sqrt{3}}{12 - 3\sqrt{3}}\\
				&=& \dfrac{48 + 25\sqrt{3}}{39}\approx 2{,}34.
			\end{eqnarray*}
		}{\begin{tikzpicture}[line join = round, line cap = round, scale=.6] 
			\coordinate[label = left:$B$] (B) at (0,0); 
			\coordinate[label = below:$C$] (C) at ($(B)+(2,0)$); 
			\coordinate[label = above:$A$] (A) at ($(B)+(0,4)$); 
			\coordinate[label = right:$D$] (D) at ($(B)+(6,0)$); 
			\draw[fill = gray!50] ($(B)!5pt!(C)$)--($(B)!5pt!(C)+(B)!5pt!(A)-(B)$)--($(B)!5pt!(A)$)--(B)--cycle;
			\draw (B)--(D)--(A)--cycle (A)--(C); 
			\foreach \x in {B,C,A,D} \fill[black] (\x) circle (1.5pt);
			\path (A)--(B) node[midway, left]{4} (B)--(C) node[midway, below]{3};	
			\end{tikzpicture} 
		}
	} 
\end{ex}

% #1939 | Archive/2024/SP-K10-K11-LyVanHoang_TVN-023/11. C1-LuongGiac-lvhoang_TVN-023/dataC1K11/Bai2-Dang1-Congthuccong.tex | mã gốc 1D1V3-2
\begin{ex}%[Cánh diều]%[Đỗ Vũ Minh Thắng, 24-25 BG 10 11]%[1D1V3-2]
	Có hai chung cư cao tầng xây cạnh nhau với khoảng cách giữa chúng là $H K=20 \mathrm{~m}$. Để đảm bảo an ninh, trên nóc chung cư thứ hai người ta lắp camera ở vị trí $C$. Gọi $A, B$ lần lượt là vị trí thấp nhất, cao nhất trên chung cư thứ nhất mà camera có thể quan sát được (Hình dưới). Hãy tính số đo góc $A C B$ (phạm vi camera có thể quan sát được ở chung cư thứ nhất). Biết rằng chiều cao của chung cư thứ hai là $C K=32 \mathrm{~m}$, $A H=6 \mathrm{~m}, B H=24 \mathrm{~m}$ (làm tròn kết quả đến hàng phần mười theo đơn vị độ).
	\shortans{30,6}
	\begin{center}
		\begin{tikzpicture}[>=stealth,line join=round,line cap=round,font=\footnotesize,scale=.7]	
		\path
		(0,0) coordinate (H)node[below]{$H$}
		(-1,0) coordinate (H1)
		(0,1) coordinate (A)node[right]{$A$}
		(0,4) coordinate (B)node[above]{$B$}
		(-1,4) coordinate (B1)
		(3,0) coordinate (K)node[below]{$K$}
		(4,0) coordinate (K1)
		(3,5) coordinate (C)node[above]{$C$}
		(4,5) coordinate (C1)
		;
		\fill[green!45]
		(H)--(B)--(B1)--(H1)--cycle
		;
		\fill[blue!45]
		(C)--(K)--(K1)--(C1)--cycle
		;
		\draw[red]
		(B)--(C)
		(A)--(C)
		(H)--(K)
		;
		\draw[black]
		(H)--(B)--(B1)--(H1)--cycle
		(C)--(K)--(K1)--(C1)--cycle
		;
		\end{tikzpicture}
	\end{center}
	\loigiai{
		\immini{Ta có $CD=CK-BH=8$, $CI=CK-AH=26$, $BD=AI=HK=20$.\\
			$\tan\widehat{BCD}=\dfrac{BD}{CD}=\dfrac{20}{8}=\dfrac{5}{2}$, $\tan\widehat{ACI}=\dfrac{AI}{CI}=\dfrac{20}{26}=\dfrac{10}{13}$.\\
			$\tan\widehat{BCA}=\tan\left(\widehat{BCD}-\widehat{ACI}\right)=\dfrac{\tan\widehat{BCD}-\tan\widehat{ACI}}{1+\tan\widehat{BCD}.\tan\widehat{ACI}}=\dfrac{\frac{5}{2}-\frac{10}{13}}{1+\frac{5}{2}\cdot\frac{10}{13}}$\\
			$\tan\widehat{BCA}=\dfrac{45}{76}\Rightarrow \widehat{BCA}\approx 30,6^\circ $.
			
		}
		{\begin{tikzpicture}[>=stealth,line join=round,line cap=round,font=\footnotesize,scale=.7]	
			\path
			(0,0) coordinate (H)node[below]{$H$}
			(-1,0) coordinate (H1)
			(0,1) coordinate (A)node[below right]{$A$}
			(0,4) coordinate (B)node[above]{$B$}
			(-1,4) coordinate (B1)
			(3,0) coordinate (K)node[below]{$K$}
			(4,0) coordinate (K1)
			(3,5) coordinate (C)node[above]{$C$}
			(4,5) coordinate (C1)
			(3,4) coordinate (D)node[below left]{$D$}
			(3,1) coordinate (I)node[below left]{$I$}
			;
			\fill[green!45]
			(H)--(B)--(B1)--(H1)--cycle
			;
			\fill[blue!45]
			(C)--(K)--(K1)--(C1)--cycle
			;
			\draw[red]
			(B)--(C)
			(A)--(C)
			(H)--(K)
			(B)--(D)
			(A)--(I)
			;
			\draw[black]
			(H)--(B)--(B1)--(H1)--cycle
			(C)--(K)--(K1)--(C1)--cycle
			;
			\end{tikzpicture}
		}
	}	
\end{ex}

% #1940 | Archive/2024/SP-K10-K11-LyVanHoang_TVN-023/11. C1-LuongGiac-lvhoang_TVN-023/dataC1K11/Bai1-Dang2.tex | mã gốc 1D1N2-3
\begin{vd}[CD, MĐ1]%[1D1N2-3]
	Đọc tên góc lượng giác, tia đầu và tia cuối của góc lượng giác đó trong hình.
	\begin{center}
		\begin{tikzpicture}[scale=1, font=\footnotesize, line join=round, line cap=round, >=stealth]
			\begin{axis}[
				axis line style={draw=none},
				axis lines=middle,
				axis equal image,
				enlargelimits,
				xtick=\empty,
				ytick=\empty,
				data cs=polar,
				samples=200,
				thick,
				line cap=round,
				line join=round,
				>=stealth
				]
				\addplot [smooth, domain=0:390,->] {1+x/5000};
				\addplot [smooth, domain=420:450,->] {1.5+x/5000} node[above,midway]{$+$};
				\addplot [mark=none] (0.475,0) node [below left] {$O$};
				\addplot [mark=none] coordinates {(0,0) (390,3+390/5000)} node[below]{$y$};
				\addplot [mark=none] coordinates {(0,0) (0,3+0/5000)} node[below]{$x$};
				\addplot [mark=none,dashed] coordinates {(0,0) (420,3+420/5000)} node[right]{$m$};
			\end{axis}
		\end{tikzpicture}
	\end{center}
	\loigiai{
		Trong hình, góc lượng giác là $(Ox,Oy)$ với tia đầu $Ox$ và tia cuối $Oy$.
	}
\end{vd}

% #1941 | Archive/2024/SP-K10-K11-LyVanHoang_TVN-023/11. C1-LuongGiac-lvhoang_TVN-023/dataC1K11/Bai1-Dang2.tex | mã gốc 1D1N2-3
\begin{vd}[CD, MĐ1]%[1D1N2-3]
	Đọc tên góc lượng giác, tia đầu và tia cuối của góc lượng giác đó trong hình.
	\begin{center}
		\begin{tikzpicture}[scale=1, font=\footnotesize, line join=round, line cap=round, >=stealth]
			\begin{axis}[
				axis line style={draw=none},
				axis lines=middle,
				axis equal image,
				enlargelimits,
				xtick=\empty,
				ytick=\empty,
				data cs=polar,
				samples=200,
				thick,
				line cap=round,
				line join=round,
				>=stealth
				]
				\addplot [smooth, domain=0:-330,->] {1+x/5000};
				\addplot [smooth, domain=-240:-270,->] {1.5+x/5000} node[above,midway]{$-$};
				\addplot [mark=none] (0.475,0) node [below left] {$O$};
				\addplot [mark=none] coordinates {(0,0) (30,3+30/5000)} node[below]{$t$};
				\addplot [mark=none] coordinates {(0,0) (0,3+0/5000)} node[below]{$z$};
				\addplot [mark=none,dashed] coordinates {(0,0) (120,3+120/5000)} node[right]{$m$};
			\end{axis}
		\end{tikzpicture}
	\end{center}
	\loigiai{
		Trong hình, góc lượng giác là $(Oz,Ot)$ với tia đầu $Oz$ và tia cuối $Ot$.
	}
\end{vd}

% #1942 | Archive/2024/SP-K10-K11-LyVanHoang_TVN-023/11. C1-LuongGiac-lvhoang_TVN-023/dataC1K11/Bai1-Dang2.tex | mã gốc 1D1N2-3
\begin{vd}[KNTT, MĐ1]%[1D1N2-3]
	\immini{Cho góc hình học $uOv$ có số đo $60^\circ$. Xác định số đo của góc lượng giác ($Ou$, $Ov$) và ($Ov$, $Ou$).}
	{\begin{tikzpicture}[scale=0.77, font=\footnotesize, line join=round, line cap=round, >=stealth]
			\draw (0,0) -- (4,0)node[below] {$u$};
			\draw[red] (0,0) -- (60:4)node[below right] {$v$};
			\draw (0:.6) arc (0:60:.6);
			\draw[-stealth] (2.75,1) arc (0:60:1);
			\path (30:3.5) node[below=-2pt]{$+$};
			\path (30:1) node{$60^\circ$};
			\path (0:0) node[below left]{$O$};
	\end{tikzpicture}}
	\loigiai{
		Ta có:
		\begin{itemize}
			\item Các góc lượng giác tia đầu $Ou$, tia cuối $Ov$ có số đo là sđ$(Ou,Ov)=60^\circ+k360^\circ$ $(k\in\mathbb{Z})$.
			\item Các góc lượng giác tia đầu $Ov$, tia cuối $Ou$ có số đo là sđ$(Ov,Ou)=-60^\circ+k360^\circ$ $(k\in\mathbb{Z})$.
		\end{itemize}
	}
\end{vd}

% #1943 | Archive/2024/SP-K10-K11-LyVanHoang_TVN-023/11. C1-LuongGiac-lvhoang_TVN-023/dataC1K11/Bai1-Dang2.tex | mã gốc 1D1N2-3
\begin{vd}[KNTT, MĐ1]%[1D1N2-3]
	Cho một góc lượng giác $(Ox, Ou)$ có số đo $-270^{\circ}$ và một góc lượng giác $(Ox, Ov)$ có số đo $135^{\circ}$. Tính số đo của các góc lượng giác $(Ou, Ov)$.
	\loigiai{
		Số đo của các góc lượng giác tia đầu $Ou$, tia cuối $Ov$ là
		$$
		\begin{aligned}
			\text{sđ}(Ou, Ov) & =\text{sđ}(Ox, Ov)-\text{sđ}(Ox, Ou)+k 360^{\circ} \\
			& =135^{\circ}-\left(-270^{\circ}\right)+k 360^{\circ}=405^{\circ}+k 360^{\circ} \\
			& =45^{\circ}+(k+1) 360^{\circ}=45^{\circ}+m 360^{\circ}(m=k+1, m \in \mathbb{Z}).
		\end{aligned}
		$$
		Vậy các góc lượng giác $(Ou, Ov)$ có số đo là $45^{\circ}+m 360^{\circ}(m \in \mathbb{Z})$.
	}
\end{vd}

% #1944 | Archive/2024/SP-K10-K11-LyVanHoang_TVN-023/11. C1-LuongGiac-lvhoang_TVN-023/dataC1K11/Bai1-Dang2.tex | mã gốc 1D1N2-3
\begin{vd}[CTST, MĐ1]%[1D1N2-3]
	Xác định số đo của các góc lượng giác $(Oa,Ob)$ trong hình sau
	\begin{listEX}[4]
		\item
		\begin{tikzpicture}[>=stealth,line join=round,line cap=round,font=\footnotesize,scale=1]
			\def\r{2.5}
			\path
			(0,0) coordinate (O)
			(\r,0) coordinate (a)
			(0,\r) coordinate (b)
			;
			\pic[draw,angle radius=2mm,angle eccentricity=1.5] {right angle = a--O--b};
			\draw (a)--(O)--(b);
			\draw[->] (0:0.5*\r) arc (0:90:0.5*\r);
			\fill[black] (O) circle (1pt);
			\foreach \l/\g in {a/-90,b/180,O/-135}
			\draw[fill=black] (\l)+(\g:.3) node{$\l$};
		\end{tikzpicture}
		\item
		\begin{tikzpicture}[>=stealth,line join=round,line cap=round,font=\footnotesize,scale=1]
			\def\r{2.5}
			\path
			(0,0) coordinate (O)
			(\r,0) coordinate (a)
			(0,\r) coordinate (b)
			;
			\pic[draw,angle radius=2mm,angle eccentricity=1.5] {right angle = a--O--b};
			\draw (a)--(O)--(b);
			\draw[->,domain=0:450,variable=\t,samples=200,>=latex]
			plot ({0.6*(\t+2*450)*cos(\t)/(2*450)},
			{0.6*(\t+ 2*450)*sin(\t)/(2*450)})
			;
			\fill[black] (O) circle (1pt);
			\foreach \l/\g in {a/-90,b/180,O/-135}
			\draw[fill=black] (\l) +(\g:.3) node{$\l$};
		\end{tikzpicture}
		\item
		\begin{tikzpicture}[>=stealth,line join=round,line cap=round,font=\footnotesize,scale=1]
			\def\r{2.5}
			\path
			(0,0) coordinate (O)
			(\r,0) coordinate (a)
			(0,\r) coordinate (b)
			;
			\pic[draw,angle radius=2mm,angle eccentricity=1.5] {right angle = a--O--b};
			\draw (a)--(O)--(b);
			\draw[->,domain=0:810,variable=\t,samples=200,>=latex]
			plot ({0.6*(\t+2*810)*cos(\t)/(2*810)},
			{0.6*(\t+ 2*810)*sin(\t)/(2*810)})
			;
			\fill[black] (O) circle (1pt);
			\foreach \l/\g in {a/-90,b/180,O/-135}
			\draw[fill=black] (\l) +(\g:.3) node{$\l$};
		\end{tikzpicture}
		\item
		\begin{tikzpicture}[>=stealth,line join=round,line cap=round,font=\footnotesize,scale=1]
			\def\r{2.5}
			\path
			(0,0) coordinate (O)
			(\r,0) coordinate (a)
			(0,\r) coordinate (b)
			;
			\pic[draw,angle radius=2mm,angle eccentricity=1.5] {right angle = a--O--b};
			\draw (a)--(O)--(b);
			\draw[->] (0:0.3*\r) arc (0:-270:0.3*\r);
			\fill[black] (O) circle (1pt);
			\foreach \l/\g in {a/-90,b/180,O/-135}
			\draw[fill=black] (\l) +(\g:.3) node{$\l$};
		\end{tikzpicture}
	\end{listEX}
	\loigiai{
		\begin{listEX}[1]
			\item Số đo của góc lượng giác $(Oa,Ob)$ trong hình vẽ là $90^\circ$.
			\item Số đo của góc lượng giác $(Oa,Ob)$ trong hình vẽ là $90^\circ + 360^\circ = 450^\circ$.
			\item Số đo của góc lượng giác $(Oa,Ob)$ trong hình vẽ là $90^\circ + 2\cdot 360^\circ = 810^\circ$.
			\item Số đo của góc lượng giác $(Oa,Ob)$ trong hình vẽ là $\dfrac{3}{4}\cdot \left(-360^\circ\right)=-270^\circ$.
		\end{listEX}
	}
\end{vd}

% #1945 | Archive/2024/SP-K10-K11-LyVanHoang_TVN-023/11. C1-LuongGiac-lvhoang_TVN-023/dataC1K11/Bai1-Dang2.tex | mã gốc 1D1N2-3
\begin{bt}%[Tex hóa SGK CD-CT,T7/22, TVN-099]%[1D1N2-3]
	\immini{Viết công thức số đo tổng quát của các góc lượng giác ($OA, OM$) và $(OA, ON)$ trong hình bên.}
	{\begin{tikzpicture}[scale=0.8, font=\footnotesize,line join=round, line cap=round, >=stealth];
			\draw[->] (-2.5,0) -- (2.5,0) node[below] {\scriptsize $x$};
			\draw[->] (0,-2.5) -- (0,2.5) node[left] {\scriptsize $y$};
			\draw (1.7,0) node[above] {$A$};
			\draw  (0:0)--(120:1.5) ;
			\draw (-0.3,0) node[below] {$O$};
			\draw  (110:1.7) node[left]{$M$};
			\draw (0:1.5) arc (0:360:1.5);
			\draw (60:0.9) node{$120^\circ$};
			\draw[-stealth] (0:.4) arc (0:120:.4);
			\draw  (0:0)--(-75:1.5) ;
			\draw  (-75:1.9) node {$N$};
			\draw (0:1.5) arc (0:-75:1.5);
			\draw (-25:0.9) node{$-75^\circ$};
			\draw[-stealth] (0:.5) arc (0:-75:.5);
			
			\draw[fill=black]  (120:1.5) circle (1.2pt) (-75:1.5) circle (1.2pt);
			\path (270:2.8) node{Hình 14};
		\end{tikzpicture}
	}
	\loigiai{
		\begin{itemize}
			\item $(OA, OM)=120^\circ+k360^\circ \; (k \in \mathbb{Z})$.
			\item $(OA, ON)=-75^\circ+k360^\circ \; (k \in \mathbb{Z})$.
		\end{itemize}
	}
	
\end{bt}

% #1946 | Archive/2024/SP-K10-K11-LyVanHoang_TVN-023/11. C1-LuongGiac-lvhoang_TVN-023/dataC1K11/Bai1-Dang2.tex | mã gốc 1D1V1-3
\begin{bt}%[Tex hóa SGK CD-CT,T7/22, TVN-099]%[1D1V1-3]
	\immini{Trong hình vẽ bên, mâm bánh xe ô tô được chia thành năm phần bằng nhau. Viết công thức số đo tổng quát của góc lượng giác $(Ox, ON)$.}
	{
		\begin{tikzpicture}[scale=1.0, font=\footnotesize,line join=round, line cap=round, >=stealth];
			\draw[line width=5,color=gray] (0:0)--(45:1.3) ;
			\draw[line width=5,color=gray] (0:0)--(117:1.3) ;
			\draw[line width=5,color=gray] (0:0)--(189:1.3) ;
			\draw[line width=5,color=gray] (0:0)--(261:1.3) ;
			\draw[line width=5,color=gray] (0:0)--(333:1.3) ;
			\draw[line width=10] (0:1.5) arc (0:360:1.5);
			\draw[line width=4,color=cyan] (0:1.3) arc (0:360:1.3);
			\fill[ball color=gray] (0:0) circle (.23);
			\draw[->] (-2.5,0) -- (2.5,0) node[below] {\scriptsize $x$};
			\draw[->] (0,-2.5) -- (0,2.5) node[left] {\scriptsize $y$};
			\draw (1.7,0) node[above] {$A$};
			\draw (-0.3,0) node[below] {$O$};
			\draw  (45:2) node {$M$};
			\draw  (261:2) node {$N$};
			\draw[color=red]  (0:0)--(45:1.7)  (0:0)--(261:1.7) ;
			\draw (20:0.9) node{$45^\circ$};
			\draw[-stealth] (0:.5) arc (0:45:.5);
			\draw[fill=black]  (120:1.5) circle (1.2pt) (-75:1.5) circle (1.2pt) (45:1.7) circle (1.2pt) (261:1.7) circle (1.2pt);
			\path (270:2.8) node{Hình 15};
		\end{tikzpicture}
	}
	\loigiai{
		Ta có $\widehat{AON}=\left(\dfrac{360^\circ}{5} -45^\circ\right)+72^\circ=99^\circ$.\\
		Vậy $(Ox, ON)=-99^\circ+k360^\circ \; (k \in \mathbb{Z})$.
		
	}
	
\end{bt}

% #1947 | Archive/2024/SP-K10-K11-LyVanHoang_TVN-023/11. C1-LuongGiac-lvhoang_TVN-023/dataC1K11/Bai1-Dang2.tex | mã gốc 1D1N2-3
\begin{bt}%[TeX hóa CD 11]%[Nhật Thiện]%[1D1N2-3]
	Cho góc lượng giác gốc $O$ có tia đầu $Ou$, tia cuối $Ov$ và có số đo $60^\circ$. Cho góc lượng giác $(O'u',O'v')$ có tia đầu $O'u'\equiv Ou$, tia cuối $O'v'\equiv Ov$. Viết công thức biểu thị số đo góc lượng giác $(O'u',O'v')$.
	\loigiai{
		Ta có $(O'u',Ov')=(Ou,Ov)+k360^\circ=60^\circ+k360^\circ\ (k\in \mathbb{Z})$.
	}
\end{bt}

% #1948 | Archive/2024/SP-K10-K11-LyVanHoang_TVN-023/11. C1-LuongGiac-lvhoang_TVN-023/dataC1K11/Bai1-Dang2.tex | mã gốc 1D1N2-3
\begin{bt}%[TeX hóa CD 11]%[Nhật Thiện]%[1D1N2-3]
	Cho góc lượng giác gốc $O$ có tia đầu $Ou$, tia cuối $Ov$ và có số đo $-\dfrac{4\pi}{3}$. Cho góc lượng giác $(O'u',O'v')$ có tia đầu $O'u'\equiv Ou$, tia cuối $O'v'\equiv Ov$. Viết công thức biểu thị số đo góc lượng giác $(O'u',O'v')$.
	\loigiai{
		Ta có $(O'u',Ov')=(Ou,Ov)+k2\pi=-\dfrac{4\pi}{3}+k2\pi\ (k\in \mathbb{Z})$.
	}
\end{bt}

% #1949 | Archive/2024/SP-K10-K11-LyVanHoang_TVN-023/11. C1-LuongGiac-lvhoang_TVN-023/dataC1K11/Bai1-Dang2.tex | mã gốc 1D1N2-3
\begin{bt}%[TeX hóa CD 11]%[Nhật Thiện]%[1D1N2-3]
	Cho góc lượng giác $(Ou,Ov)$ có số đo là $\dfrac{3\pi}{4}$, góc lượng giác $(Ou,Ow)$ có số đo là $\dfrac{5\pi}{4}$. Tìm số đo của góc lượng giác $(Ov,Ow)$.
	\loigiai{
		Theo hệ thức Chasles, ta có
		\begin{eqnarray*}
			(Ov,Ow)&=&(Ou,Ow)-(Ou,Ov)+k2\pi\\
			&=&\dfrac{5\pi}{4}-\dfrac{3\pi}{4}+k2\pi\\
			&=&\dfrac{\pi}{2}+k2\pi\ (k\in \mathbb{Z}).
		\end{eqnarray*}
	}
\end{bt}

% #1959 | Archive/2024/SP-K10-K11-LyVanHoang_TVN-023/11. C1-LuongGiac-lvhoang_TVN-023/dataC1K11/Chuong1-Ontap-De02.tex | mã gốc 1D1H1-3
\begin{ex}%[Câu 13]%[1D1H1-3]%[24-5-Bài giảng K10-K11, Nguyễn Văn Hồng]
Trong hình vẽ bên, ta xem hình ảnh đường tròn trên một bánh lái tàu thuỷ tương ứng với một đường tròn lượng giác. Các mệnh đề sau \textbf{ĐÚNG} hay \textbf{SAI} ?
\immini{
\choiceTF[t]
{\True Công thức tổng quát biểu diễn góc lượng giác $(OA,OB)$ theo đơn vị radian là\break $(OA,OB)=\dfrac{\pi}{4}+k2\pi \, (k\in \mathbb{Z})$}
{\True Công thức tổng quát chỉ ra góc lượng giác tương ứng với hai điểm biểu diễn là $A$, $E$ theo đơn vị độ là $k180 ^{\circ},\, (k\in \mathbb{Z})$}
{Công thức tổng quát chỉ ra góc lượng giác tương ứng với bốn điểm biểu diễn là $A$, $C$, $E$, $G$ theo đơn vị radian là $k\dfrac{\pi}{3}\, (k\in \mathbb{Z})$}
{Công thức tổng quát biểu diễn góc lượng giác $(OA,OC)+(OC,OH)$ theo đơn vị radian là\break	$\dfrac{\pi}{4}+k2\pi,\, (k\in \mathbb{Z})$}
}{
\begin{tikzpicture}[font=\footnotesize]
\draw[->] (180:2.5)--(0:2.5) node[below] {$x$};
\draw[->] (-90:2.5)--(90:2.5) node[left] {$y$};
\draw (0:0) circle (1.5);
\fill (00:0) circle (0.05)
(-140:0.18) node[below] {$O$}
(45:1.5) circle (0.05) node[right] {$B$}
(-45:1.5) circle (0.05) node[right] {$H$}
(135:1.5) circle (0.05) node[left] {$D$}
(225:1.5) circle (0.05) node[left] {$F$}
(0:1.5) circle (0.05) node[above right] {$A$}
(180:1.5) circle (0.05) node[above left] {$E$}
(90:1.5) circle (0.05) node[above left] {$C$}
(-90:1.5) circle (0.05) node[below left] {$G$};
\draw (0:0)--(45:1.5)  (0:0)--(135:1.5)  (0:0)--(225:1.5) (0:0)--(-45:1.5) ;
\end{tikzpicture}}
\loigiai{
\begin{itemchoice}
\itemch Ta có $(OA,OB)=\dfrac{\pi}{4}+k2\pi \, (k\in \mathbb{Z})$.
\itemch Ta thấy hai điểm $A$, $E$ lần lượt biểu diễn cho các góc lượng giác $0 ^{\circ}$, $180 ^{\circ}$, $360 ^{\circ}$, $540 ^{\circ}, \ldots$ Tất cả các góc này theo thứ tự chênh lệch nhau $180 ^{\circ}$.\\
Vì vậy công thức tổng quát biểu diễn cho các góc lượng giác ấy là $k180 ^{\circ}\, (k\in \mathbb{Z})$.
\itemch 	Ta thấy $A$, $C $, $E$, $G$ lần lượt biểu diễn cho các góc lượng giác $0$ rad, $\dfrac{\pi}{2}$ rad, $\pi$ rad, $\dfrac{3\pi}{2}$ rad, $2\pi$ rad, $\dfrac{5\pi}{2}$ rad, $\ldots$ Tất cả các góc này theo thứ tự chênh lệch nhau $\dfrac{\pi}{2}$ rad.\\
Vì vậy công thức tổng quát biểu diễn cho các góc lượng giác ấy là $k\dfrac{\pi}{2}\,  (k\in \mathbb{Z})$.
\itemch Theo hệ thức Sa-lơ, ta có
\[(OA,OC)+(OC,OH)=(OA,OH)=-\dfrac{\pi}{4}+k2\pi, \, (k\in \mathbb{Z}).
\]
\end{itemchoice}
}
\end{ex}

% #1966 | Archive/2024/SP-K10-K11-LyVanHoang_TVN-023/11. C1-LuongGiac-lvhoang_TVN-023/dataC1K11/Chuong1-Ontap-De02.tex | mã gốc 1D1H3-7
\begin{ex}%[Câu 20]%[1D1H3-7]%[24-5-Bài giảng K10-K11, Nguyễn Văn Hồng]
Nhiệt độ ngoài trời ở một thành phố vào các thời điểm khác nhau trong ngày có
thể được mô phỏng bởi công thức $h(t)=31+3\sin \dfrac{\pi}{12}(t-9)$, với $h$ tính bằng độ $\mathrm{C}$ và $t$ là thời gian trong ngày tính bằng giờ $(0<t\leq 24)$.
Tính nhiệt độ ngoài trời (theo đơn vị độ) ở thành phố đó vào lúc $7$ giờ tối. \shortans{$32{,}5$}
\loigiai{
Nhiệt độ ngoài trời lúc $7$ giờ tối (ứng vơi $t=19$) là
\[h(19)=31+3 \sin \dfrac{\pi}{12}(19-9)=31+3 \sin \dfrac{5 \pi}{6}=32{,}5^{\circ} C.\]
}
\end{ex}

% #1967 | Archive/2024/SP-K10-K11-LyVanHoang_TVN-023/11. C1-LuongGiac-lvhoang_TVN-023/dataC1K11/Chuong1-Ontap-De02.tex | mã gốc 1D1V5-6
\begin{ex}%[Câu 21]%[1D1V5-6]%[24-5-Bài giảng K10-K11, Nguyễn Văn Hồng]
Độ sâu $h$(m) của mực nước ở một cảng biển vào thời điểm $t$ (giờ) sau khi thuỷ triều lên lần đầu tiên trong ngày được tính xấp xỉ bởi công thức $h(t)=0{,}8 \cdot \cos 0{,}5 t+5$.
\begin{flushright}
\textit{(Theo \href{https://noc.ac.uk/files/documents/}{https://noc.ac.uk/files/documents/}business/an-introduction-to-tidalmodelling.pdf)}
\end{flushright}
Một con tàu cần mực nước sâu tối thiểu $4{,}6$m để có thể di chuyển ra vào cảng an toàn. Hỏi có bao nhiêu thời điểm trong vòng $12$ tiếng sau khi thuỷ triều lên lần đầu tiên trong ngày tàu có thể hạ thuỷ.
\shortans{ $2$}

\loigiai{
Theo đề bài, tàu cần mực nước sâu tối thiểu $4{,}6$m để có thể di chuyển ra vào cảng an toàn.\\
Suy ra $h(t)=0{,}8 \cos (0{,}5 t)+5=4{,}6 \quad (*)$ có nghiệm thuộc $\left[0;12\right]$.\\
Ta có
\allowdisplaybreaks
\begin{eqnarray*}
(*) &\Leftrightarrow& \cos (0{,}5 t)=-\dfrac{1}{2} \\
& \Leftrightarrow& 0{,}5 t= \pm \dfrac{2 \pi}{3}+k 2 \pi \\
&\Leftrightarrow&\hoac{&t=\dfrac{4 \pi}{3}+k 4 \pi \\&t=-\dfrac{4 \pi}{3}+l 4 \pi}\quad (k, l \in \mathbb{Z}.)
\end{eqnarray*}
Do nghiệm thuộc $\left[0;12\right]$ nên $\hoac{
&0 \leq \dfrac { 4 \pi } { 3 } + k 4 \pi \leq 1 2 \\
&0 \leq - \dfrac { 4 \pi } { 3 } + l 4 \pi \leq 1 2} \Leftrightarrow \hoac{
&-\dfrac{1}{3} \leq k \leq \dfrac{1}{4 \pi}\left(12-\dfrac{4 \pi}{3}\right) \approx 0{,}62 \\
&\dfrac{1}{3} \leq l \leq \dfrac{1}{4 \pi}\left(12+\dfrac{4 \pi}{3}\right) \approx 1{,}29. }
\Leftrightarrow 
\hoac{
& k=0 \\
	&l=1.}$
\begin{itemize}
\item Với $k=0$, suy ra $t=\dfrac{4 \pi}{3} \approx 4{,}19$ (giờ).
\item Với $l=1$, suy ra $t=\dfrac{2 \pi}{3} \approx 2{,}09$ (giờ).
\end{itemize}
Vậy, có $2$ thời điểm trong vòng $12$ tiếng sau khi thủy triều lên lần đầu tiên trong ngày tàu có thể hạ thủy.
}
\end{ex}

% #1968 | Archive/2024/SP-K10-K11-LyVanHoang_TVN-023/11. C1-LuongGiac-lvhoang_TVN-023/dataC1K11/Chuong1-Ontap-De02.tex | mã gốc 1D1V5-6
\begin{ex}%[Câu 22]%[1D1V5-6]%[24-5-Bài giảng K10-K11, Nguyễn Văn Hồng]
Giả sử một vật dao động điều hoà xung quanh vị trí cân bằng theo phương trình $\break x=2\cos\left(2t+\dfrac{\pi}{4}\right)$. Ở đây, thời gian $t$ tính bằng giây và quãng đường $x$ tính bằng centimét. Hãy cho biết trong thời gian từ $0$ đến $20$ giây, vật đi qua vị trí cân bằng bao nhiêu lần?
\shortans{ $13$}
\loigiai{
Vị trí cân bằng của vật dao động điều hòa là vị trí vật đứng yên, khi đó $x=0$.\\
Xét phương trình $2\cos\left(2t+\dfrac{\pi}{4}\right)=0$ ta có
$$2\cos\left(2t+\dfrac{\pi}{4}\right)=0
\Leftrightarrow
2t+\dfrac{\pi}{4}=\dfrac{\pi}{2}+ k\pi \Leftrightarrow
t= \dfrac{\pi}{8}+ k\dfrac{\pi}{2}, k\in\mathbb{Z}.$$
Trong thời gian từ $0$ đến $20$ giây, tức là $0\le t\le 20$ hay $0\le\dfrac{\pi}{8}+k\dfrac{\pi}{2}\le 20$
$\Leftrightarrow-\dfrac{1}{4}\le k\le\dfrac{160-\pi}{4\pi}$.\\
Mà $k\in\mathbb{Z}$ nên $k\in\left\{{0;1;2;3;\ldots;12}\right\}$.\\
Vậy trong khoảng thời gian từ $0$ đến $20$ giây, vật đi qua vị trí cân bằng $13$ lần.}
\end{ex}

% #1969 | Archive/2024/SP-K10-K11-LyVanHoang_TVN-023/11. C1-LuongGiac-lvhoang_TVN-023/dataC1K11/Bai-3-Dang-3-Xettinhchanlecuahamso.tex | mã gốc 1D1N4-4
\begin{vd}%[Bài giảng K10-11, Nguyễn Đắc Giáp]%[1D1N4-4] 
[Mức độ 1]
	Xét tính chẵn, lẻ của hàm số $y=\sin 2x$.
	\loigiai{
		Tập xác định $ \mathscr{D}=\mathbb{R}$.\\
		Với mọi $x\in\mathscr{D} \Rightarrow -x \in \mathscr{D}$, ta có
		$$f(-x)=\sin (-2x)=-\sin 2x = -f(x).$$
		Vậy hàm số $y=\sin 2x$ là hàm số lẻ.
	}
\end{vd}

% #1970 | Archive/2024/SP-K10-K11-LyVanHoang_TVN-023/11. C1-LuongGiac-lvhoang_TVN-023/dataC1K11/Bai-3-Dang-3-Xettinhchanlecuahamso.tex | mã gốc 1D1N4-4
\begin{vd}%[Bài giảng K10-11, Nguyễn Đắc Giáp]%[1D1N4-4]
	[Mức độ 1]
	Xét tính chẵn lẻ của hàm số $y=\sin \dfrac{1}{x^2-1}$.
	\loigiai{
		Tập xác định $ \mathscr{D}=\mathbb{R}\setminus\{\pm 1\}$.\\
		Với mọi $x\in\mathscr{D} \Rightarrow -x \in \mathscr{D}$, ta có
		$$f(-x)=\sin\dfrac{1}{(-x)^2-1}=\sin\dfrac{1}{x^2-1}=f(x).$$
		Vậy hàm số $y=\sin \dfrac{1}{x^2-1}$ là hàm số chẵn.
	}
\end{vd}

% #1971 | Archive/2024/SP-K10-K11-LyVanHoang_TVN-023/11. C1-LuongGiac-lvhoang_TVN-023/dataC1K11/Bai-3-Dang-3-Xettinhchanlecuahamso.tex | mã gốc 1D1N4-4
\begin{vd}%[Bài giảng K10-11, Nguyễn Đắc Giáp]%[1D1N4-4]
	[Cánh diều - Mức độ 1]
	Xét tính chẵn lẻ của hàm số $f(x)=\sin x+\tan x$.
	\loigiai{
		Tập xác định của hàm số $f(x)$ là $\mathscr{D}=\mathbb{R} \setminus \left\{\dfrac{\pi}{2}+k \pi \,\middle|\, k \in \mathbb{Z}\right\}$.\\
		Với mọi $x \in \mathscr{D}$, ta có $-x \in \mathscr{D}$
		và $$f(-x)=\sin (-x)+\tan (-x)=-\sin x-\tan x =-(\sin x+\tan x)=-f(x).$$	
		Vậy hàm số $f(x)=\sin x+\tan x$ là hàm số lẻ.
	}
\end{vd}

% #1972 | Archive/2024/SP-K10-K11-LyVanHoang_TVN-023/11. C1-LuongGiac-lvhoang_TVN-023/dataC1K11/Bai-3-Dang-3-Xettinhchanlecuahamso.tex | mã gốc 1D1H4-4
\begin{vd}%[Bài giảng K10-11, Nguyễn Đắc Giáp]%[1D1H4-4] 
	[Kết nối tri thức với cuộc sống - Mức độ 2]
	Xét tính chẵn lẻ của hàm số $f(x)=x \sin x$.
	\loigiai{
		Tập xác định của hàm số là $\mathscr{D}=\mathbb{R}$.\\
		Do đó, nếu $x$ thuộc tập xác định $\mathscr{D}$ thì $-x$ cũng thuộc tập xác định $\mathscr{D}$.\\
		Ta có $f(-x)=(-x) \cdot \sin(-x)=x \sin x=f(x), \forall x \in \mathscr{D}$.\\
		Vậy $f(x)=x \sin x$ là hàm số chẵn.
	}
\end{vd}

% #1973 | Archive/2024/SP-K10-K11-LyVanHoang_TVN-023/11. C1-LuongGiac-lvhoang_TVN-023/dataC1K11/Bai-3-Dang-3-Xettinhchanlecuahamso.tex | mã gốc 1D1H4-4
\begin{vd}%[Bài giảng K10-11, Nguyễn Đắc Giáp]%[1D1H4-4]
	[Cùng khám phá - Mức độ 2]
	Xét tính chẵn lẻ của hàm số $y=f(x)=|\sin x| \cos x$.
	\loigiai{
		Hàm số $y=f(x)=|\sin x| \cos x$ có tập xác định là $\mathscr{D}=\mathbb{R}$.\\
		Ta có với mọi $x \in \mathscr{D}$ thì $-x \in \mathscr{D}$ và 
		$$f(-x)=|\sin (-x)| \cdot\cos (-x)=|-\sin x|\cdot \cos x=|\sin x| \cos x=f(x).$$
		Vậy hàm số $y=f(x)=|\sin x| \cos x$ là hàm số chẵn trên $\mathbb{R}$.
	}
\end{vd}

% #1974 | Archive/2024/SP-K10-K11-LyVanHoang_TVN-023/11. C1-LuongGiac-lvhoang_TVN-023/dataC1K11/Bai-3-Dang-3-Xettinhchanlecuahamso.tex | mã gốc 1D1H4-4
\begin{vd}%[Bài giảng K10-11, Nguyễn Đắc Giáp]%[1D1H4-4]
	[Mức độ 2]
	Xét tính chẵn lẻ của các hàm số sau.
	\begin{enumEX}{2}
		\item $y=-\sin x$;
		\item $y=\cos x-\sin x$;
		\item $y=\cos x+\sin^2x$;
		\item $y=\cos x\sin x$.
	\end{enumEX}
	\loigiai{
		Tất các các hàm số đều có tập xác định $ \mathscr{D}=\mathbb{R}$.\\
		Do đó $\forall x\in \mathscr{D}\Rightarrow -x\in \mathscr{D}$ và
		\begin{enumerate}
			\item Với $y=f(x)=-\sin x$. Ta có $$f(-x)=-\sin (-x)=\sin x=-(-\sin x)=-f(x).$$ Suy ra hàm số $y=-\sin x$ là hàm số lẻ.
			\item  Với $y=f(x)=\cos x-\sin x$. Ta có $$f(-x)=\cos (-x)-\sin (-x)=\cos x+\sin x.$$
			$\Rightarrow f(-x)\ne-f(x)$ và $f(-x)\ne f(x)$.\\
			Suy ra hàm số $y=\cos x-\sin x$ không là hàm số chẵn, không là hàm số lẻ.
			\item  Với $y=f(x)=\cos x+\sin^2x$. Ta có $$f(-x)=\cos (-x)+\sin^2(-x)=\cos (-x)+\left(-\sin x\right)^2=\cos x+\sin^2x.$$
			$\Rightarrow f(-x)=f(x)$.\\
			Suy ra hàm số $y=\cos x+\sin^2x$ là hàm số chẵn.
			\item  Với $y=f(x)=\cos x\cdot \sin x$. Ta có $$f(-x)=\cos (-x)\cdot \sin (-x)=-\cos x\sin x=-f(x).$$
			Suy ra hàm số $y=\cos x\sin x$ là hàm số lẻ.
		\end{enumerate}
	}
\end{vd}

% #1977 | Archive/2024/SP-K10-K11-LyVanHoang_TVN-023/11. C1-LuongGiac-lvhoang_TVN-023/dataC1K11/Bai-3-Dang-3-Xettinhchanlecuahamso.tex | mã gốc 1D1H4-4
\begin{bt}%[Bài giảng K10-11, Nguyễn Đắc Giáp]%[1D1H4-4]
  	[Mức độ 1]
 	Xét tính chẵn, lẻ của hàm số $y=\sin x^2$.
 	\loigiai{
 		Tập xác định $\mathscr{D}=\mathbb{R}$.\\
 		$\forall x \in \mathscr{D} \Rightarrow -x \in \mathscr{D}$.\\
 		Ta có $f(-x)=\sin (-x)^2=\sin x^2 =f(x)$.\\
 		Vậy hàm số $y=\sin x^2$ là hàm số chẵn.
 	}
 \end{bt}

% #1978 | Archive/2024/SP-K10-K11-LyVanHoang_TVN-023/11. C1-LuongGiac-lvhoang_TVN-023/dataC1K11/Bai-3-Dang-3-Xettinhchanlecuahamso.tex | mã gốc 1D1B1-3
\begin{bt}%[Bài giảng K10-11, Nguyễn Đắc Giáp]%[1D1B1-3]
 	[Mức độ 2]
 	Xét tính chẵn, lẻ của hàm số $y=\sin x + \cos x$.
 	\loigiai{
 		Tập xác định $ \mathscr{D}=\mathbb{R}$.\\
 		$\forall x \in \mathscr{D} \Rightarrow -x \in \mathscr{D}$.\\
 		Ta có $f(-x)=\sin (-x)+\cos (-x)=-\sin x +\cos x\ne \pm f(x)$.\\
 		Vậy hàm số $y=\sin x + \cos x$ là hàm số không chẵn, không lẻ.
 	}
 \end{bt}

% #1979 | Archive/2024/SP-K10-K11-LyVanHoang_TVN-023/11. C1-LuongGiac-lvhoang_TVN-023/dataC1K11/Bai-3-Dang-3-Xettinhchanlecuahamso.tex | mã gốc 1D1H4-4
\begin{bt}%[Bài giảng K10-11, Nguyễn Đắc Giáp]%[1D1H4-4]
 	[Kết nối tri thức với cuộc sống - Mức độ 2]
 	Xét tính chẵn lẻ của các hàm số sau:
 	\begin{enumEX}{2}
 		\item $y=\sin 2x+\tan 2x$;
 		\item  $y=\cos x+\sin ^2x$;
 		\item $y=\sin x \cos 2x$;
 		\item $y=\sin x+\cos x$.
 	\end{enumEX}
 	\loigiai{
 		\begin{enumerate}
 			\item Tập xác định của hàm số $y=f(x)=\sin 2x+\tan 2x$ là $\mathscr{D}=\mathbb{R} \setminus \left\{\left.\dfrac{\pi}{4}+k \dfrac{\pi}{2} \right\rvert\, k \in \mathbb{Z}\right\}$.\\			
 			Do đó, nếu $x$ thuộc tập xác định $\mathscr{D}$ thì $-x$ cũng thuộc tập xác định $\mathscr{D}$ và
 			\begin{eqnarray*}
 				f(-x)&=&\sin (-2x)+\tan (-2x)\\
 				&=&-\sin 2x-\tan 2x\\
 				&=&-(\sin 2x+\tan 2x)=-f(x).
 			\end{eqnarray*}
 			Vậy $y=\sin 2x+\tan 2x$ là hàm số lẻ.
 			\item Tập xác định của hàm số $y=f(x)=\cos x+\sin ^2x$ là $\mathscr{D}=\mathbb{R}$.\\
 			Do đó, nếu $x$ thuộc tập xác định $\mathscr{D}$ thì $-x$ cũng thuộc tập xác định $\mathscr{D}$ và
 			\begin{eqnarray*}
 				f(-x)&=& \cos (-x)+\sin^2(-x) \\
 				&=& \cos x+(-\sin x)^2 \\
 				&=&\cos x+\sin ^2x=f(x). 
 			\end{eqnarray*}
 			Vậy $y=\cos x+\sin ^2x$ là hàm số chẵn.
 			\item Tập xác định của hàm số $y=f(x)=\sin x \cos 2x$ là $\mathscr{D}=\mathbb{R}$.\\
 			Do đó, nếu $x$ thuộc tập xác định $\mathscr{D}$ thì $-x$ cũng thuộc tập xác định $\mathscr{D}$ và
 			\begin{eqnarray*}
 				f(-x)&=& \sin (-x) \cdot \cos (-2x)\\
 				&=&  -\sin x \cdot \cos 2x=-f(x).
 			\end{eqnarray*}
 			Vậy $y=\sin x \cos 2x$ là hàm số lẻ.
 			\item Tập xác định của hàm số $y=f(x)=\sin x+\cos x$ là $\mathscr{D}=\mathbb{R}$.\\
 			Do đó, nếu $x$ thuộc tập xác định $\mathscr{D}$ thì $-x$ cũng thuộc tập xác định $\mathscr{D}$ và
 			$$f(-x)=\sin (-x)+\cos (-x)=-\sin x+\cos x \neq\pm f(x).$$
 			Vậy $y=\sin x+\cos x$ là hàm số không chẵn, không lẻ.
 		\end{enumerate}
 	}
 \end{bt}

% #1980 | Archive/2024/SP-K10-K11-LyVanHoang_TVN-023/11. C1-LuongGiac-lvhoang_TVN-023/dataC1K11/Bai-3-Dang-3-Xettinhchanlecuahamso.tex | mã gốc 1D1H4-4
\begin{bt}%[Bài giảng K10-11, Nguyễn Đắc Giáp]%[1D1H4-4]
 	[Cánh diều - Mức độ 2]
 	Xét tính chẵn lẻ của các hàm số sau
 	\begin{enumEX}{3}
 		\item $y=\sin x \cos x$;
 		\item $y=\tan x+\cot x$;
 		\item $y=\sin^2 x$.
 	\end{enumEX}
 	\loigiai{
 		\begin{enumerate}
 			\item Xét hàm số $y=f(x)=\sin x \cos x$ có tập xác định $\mathscr{D}=\mathbb{R}$.\\			
 			Do đó, nếu $x$ thuộc tập xác định $\mathscr{D}$ thì $-x$ cũng thuộc tập xác định $\mathscr{D}$ và
 			$$f(-x)=\sin (-x) \cdot \cos (-x)=-\sin x \cos x=-f(x).$$
 			Vậy $y=\sin x \cos x$ là hàm số lẻ.
 			\item  Xét hàm số $f(x)=\tan x+\cot x$ có tập xác định $\mathscr{D}=\mathbb{R} \setminus   \left\{k\pi ; \dfrac{\pi}{2}+k \pi \,\middle|\, k \in \mathbb{Z}\right\}$.\\
 			Do đó, nếu $x$ thuộc tập xác định $\mathscr{D}$ thì $-x$ cũng thuộc tập xác định $\mathscr{D}$ và
 			\begin{eqnarray*}
 				f(-x)&=& \tan (-x)+\cot (-x) \\
 				&=&  (-\tan x)+(-\cot x)\\
 				&=&  -(\tan x+\cot x)=-f(x).
 			\end{eqnarray*}
 			Do đó hàm số $y=\tan x+\cot x$ là hàm số lẻ.
 			\item Xét hàm số $f(x)=\sin^ 2x$ có tập xác định $\mathscr{D}=\mathbb{R}$.\\
 			Do đó, nếu $x$ thuộc tập xác định $\mathscr{D}$ thì $-x$ cũng thuộc tập xác định $\mathscr{D}$ và
 			$$f(-x)=\sin ^2(-x)=(-\sin x)^2=\sin ^2 x=f(x).$$
 			Vậy hàm số $y=\sin ^2 x$ là hàm số chẵn.
 		\end{enumerate}
 	}
 \end{bt}

% #1981 | Archive/2024/SP-K10-K11-LyVanHoang_TVN-023/11. C1-LuongGiac-lvhoang_TVN-023/dataC1K11/Bai-3-Dang-3-Xettinhchanlecuahamso.tex | mã gốc 1D1H4-4
\begin{bt}%[Bài giảng K10-11, Nguyễn Đắc Giáp]%[1D1H4-4]
 	[SBT Kết nối tri thức với cuộc sống - Mức độ 2]
 	Xét tính chẵn lẻ của các  hàm số sau:
 	\begin{enumEX}{2}
 		\item $y=x \sin 3x$;
 		\item  $y=\dfrac{1+\cos 3 x}{1-\cos 3x}$;
 		\item  $y=x^3 \cos 2x$;
 		\item  $y=\sin x-\cos x$.
 	\end{enumEX}
 	\loigiai{
 		\begin{enumerate}
 			\item Tập xác định của hàm số là $\mathscr{D}=\mathbb{R}$.\\
 			Nếu kí hiệu $f(x)=x \sin 3x$ thì với mọi $x \in \mathscr{D}$, ta có $-x \in \mathscr{D}$ và
 			$$	f(-x)=(-x) \sin (-3x)=x \sin 3 x=f(x).	$$
 			Vậy hàm số đã cho là hàm số chẵn.
 			\item  Tập xác định của hàm số là $\mathscr{D}=\mathbb{R}\setminus \left\{k \dfrac{2 \pi}{3} \,\middle|\, k \in \mathbb{Z}\right\}$.\\
 			Nếu kí hiệu $f(x)=\dfrac{1+\cos 3x}{1-\cos 3x}$ thì với mọi $x \in \mathscr{D}$, ta có $-x \in \mathscr{D}$ và
 			$$f(-x)=\dfrac{1+\cos (-3x)}{1-\cos (-3x)}=\dfrac{1+\cos 3 x}{1-\cos 3 x}=f(x).$$
 			Vậy hàm số đã cho là hàm số chẵn.
 			\item Tập xác định của hàm số là $\mathscr{D}=\mathbb{R}$.\\
 			Nếu kí hiệu $f(x)=x^3 \cos 2x$ thì với mọi $x \in \mathscr{D}$, ta có $-x \in \mathscr{D}$ và 
 			$$f(-x)=(-x)^3 \cos (-2x)=-x^3 \cos 2x=-f(x).$$
 			Vậy hàm số đã cho là hàm số lẻ.
 			\item Nếu kí hiệu $f(x)=\sin x-\cos x$
 			thì $f(-1)=\sin (-1)-\cos (-1)=-\sin 1-\cos 1 $ và $f(1)=\sin 1 -\cos 1$.\\
 			Vì $f(-1) \neq f(1)$ và $f(-1) \neq-f(1)$ nên hàm số đã cho không phải là hàm số chẵn, cũng không phải là hàm số lẻ.
 		\end{enumerate}
 	}
 \end{bt}

% #1982 | Archive/2024/SP-K10-K11-LyVanHoang_TVN-023/11. C1-LuongGiac-lvhoang_TVN-023/dataC1K11/Bai-3-Dang-3-Xettinhchanlecuahamso.tex | mã gốc 1D1V4-4
\begin{bt}%[Bài giảng K10-11, Nguyễn Đắc Giáp]%[1D1V4-4]
	[SBT Chân trời sáng tạo - Mức độ 3]
	Xét tính chẵn lẻ của các  hàm số sau:
	\begin{enumEX}{3}
		\item $y=\dfrac{\sin 3 x}{x}$;
		\item $y=-5 x^2+\cos \dfrac{x}{2}$;
		\item $y=x \sqrt{1+\cos 2 x}$;
		\item $y=\cot x-\dfrac{2}{\sin x}$;
		\item $y=|x|+\tan x$;
		\item $y=\tan \left(x+\dfrac{\pi}{4}\right)$.
	\end{enumEX}
	\loigiai{
		\begin{enumerate}
			\item Tập xác định của hàm số $y=\dfrac{\sin 3x}{x}$ là $\mathscr{D}=\mathbb{R} \setminus \{0\}$ thoả mãn điều kiện $-x \in \mathscr{D}$ với mọi $x \in \mathscr{D}$ và
			$$\dfrac{\sin [3(-x)]}{-x}=\dfrac{\sin (-3x)}{-x}=\dfrac{-\sin 3x}{-x}=\dfrac{\sin 3x}{x}.$$
			Vậy hàm số $y=\dfrac{\sin 3x}{x}$ là hàm số chẵn.
			\item Tập xác định của hàm số $y=-5x^2+\cos \dfrac{x}{2}$ là $\mathscr{D}=\mathbb{R}$ thoả mãn điều kiện $-x \in \mathscr{D}$ với mọi $x \in \mathscr{D}$ và
			$$-5(-x)^2+\cos \left(\dfrac{-x}{2}\right)=-5x^2 +\cos \dfrac{x}{2}.$$
			Vậy hàm số $y=-5x^2+\cos \dfrac{x}{2}$ là hàm số chẵn.
			\item Tập xác định của hàm số $y=x \sqrt{1+\cos 2x}$ là $\mathscr{D}=\mathbb{R}$ thoả mãn điều kiện $-x \in \mathscr{D}$ với mọi $x \in \mathscr{D}$ và
			$$(-x) \sqrt{1+\cos (-2x)}=-x \sqrt{1+\cos 2 x}.$$
			Vậy hàm số $y=x \sqrt{1+\cos 2 x}$ là hàm số lẻ.
			\item Tập xác định của hàm số $y=\cot x-\dfrac{2}{\sin x}$ là $\mathscr{D}=\mathbb{R} \setminus \{k \pi \mid k \in \mathbb{Z}\}$ thoả mãn điều kiện $-x \in \mathscr{D}$ với mọi $x \in \mathscr{D}$ và
			$$\cot (-x)-\dfrac{2}{\sin (-x)}=-\cot x+\dfrac{2}{\sin x}=-\left(\cot x-\dfrac{2}{\sin x}\right).$$
			Vậy hàm số $y=\cot x-\dfrac{2}{\sin x}$ là hàm số lẻ.
			\item Tập xác định của hàm số $y=|x|+\tan x$ là $\mathscr{D}=\mathbb{R} \setminus \left\{\dfrac{\pi}{2}+k \pi \,\middle|\, k \in \mathbb{Z}\right\}$ thoả mãn điều kiện $-x \in \mathscr{D}$ với mọi $x \in \mathscr{D}$.\\
			Đặt $f(x)=|x|+\tan x$. Xét hai giá trị $\dfrac{\pi}{4}$ và $-\dfrac{\pi}{4}$ thuộc $\mathscr{D}$, ta có
			$$
			f\left(\dfrac{\pi}{4}\right)=\left|\dfrac{\pi}{4}\right|+\tan \dfrac{\pi}{4}=\dfrac{\pi}{4}+1 \text { và } f\left(-\dfrac{\pi}{4}\right)=\left|-\dfrac{\pi}{4}\right|+\tan \left(-\dfrac{\pi}{4}\right)=\dfrac{\pi}{4}-1.
			$$		
			Do $f\left(-\dfrac{\pi}{4}\right) \neq f\left(\dfrac{\pi}{4}\right)$ và $f\left(-\dfrac{\pi}{4}\right) \neq-f\left(\dfrac{\pi}{4}\right)$ nên $y=|x|+\tan x$ không là hàm số chẵn cũng không là hàm số lẻ.
			\item Tập xác định của hàm số $y=\tan \left(x+\dfrac{\pi}{4}\right)$ là $\mathscr{D}=\mathbb{R} \setminus \left\{\dfrac{\pi}{4}+k \pi \,\middle|\, k \in \mathbb{Z}\right\}$ không thoả mãn điều kiện $-x \in \mathscr{D}$ với mọi $x \in \mathscr{D}$ vì $-\dfrac{\pi}{4} \in \mathscr{D}$ mà $\dfrac{\pi}{4} \notin \mathscr{D}$.\\
			Vậy hàm số $y=\tan \left(x+\dfrac{\pi}{4}\right)$ không là hàm số chẵn cũng không là hàm số lẻ.
		\end{enumerate}
	}
\end{bt}

% #1989 | Archive/2024/SP-K10-K11-LyVanHoang_TVN-023/11. C1-LuongGiac-lvhoang_TVN-023/dataC1K11/Bai1-Dang3.tex | mã gốc 1D1N1-4
\begin{vd}[MĐ1]%[1D1N1-4]
	Một đường tròn có bán kính $10$ cm. Tính độ dài các cung tròn bị chắn bởi các góc ở tâm có số đo là $\dfrac{2\pi}{3}$.
	\loigiai
	{
		Độ dài của cung tròn có số đo $\dfrac{2\pi}{3}$, bán kính $10$ cm là $\ell = 10 \cdot \dfrac{2\pi}{3} = \dfrac{20\pi}{3}$ cm.
	}
\end{vd}

% #1990 | Archive/2024/SP-K10-K11-LyVanHoang_TVN-023/11. C1-LuongGiac-lvhoang_TVN-023/dataC1K11/Bai1-Dang3.tex | mã gốc 1D1N1-4
\begin{vd}[MĐ1]%[1D1N1-4]
	Một đường tròn có bán kính $20$ cm. Tính độ dài các cung tròn bị chắn bởi các góc ở tâm có số đo là $45^\circ$.
	\loigiai
	{
		Độ dài của cung tròn có số đo $45^\circ = \dfrac{\pi}{4}$, bán kính $20$ cm là $\ell = 20 \cdot \dfrac{\pi}{4} = 5\pi$ cm.
	}
\end{vd}

% #1991 | Archive/2024/SP-K10-K11-LyVanHoang_TVN-023/11. C1-LuongGiac-lvhoang_TVN-023/dataC1K11/Bai1-Dang3.tex | mã gốc 1D1H1-4
\begin{vd}[MĐ2]%[1D1H1-4]
	Kim giờ của đồng hồ dài $7$ cm, kim phút dài $10$ cm. Tính quãng đường kim phút, kim giờ đi được trong $30$ phút.
	\loigiai
	{
		Trong $30$ phút kim giờ quét được một cung tròn có số đo bằng $\dfrac{\pi}{12}$ và kim phút quét được một cung tròn có số đo bằng $\pi$.\\
		Quãng đường kim giờ đi được trong $30$ phút là $7\cdot \dfrac{\pi}{12} = \dfrac{7\pi}{12} \approx 1{,}83$ cm.\\
		Quãng đường kim phút đi được trong $30$ phút là $10 \pi \approx 31{,}42$ cm.
	}
\end{vd}

% #1993 | Archive/2024/SP-K10-K11-LyVanHoang_TVN-023/11. C1-LuongGiac-lvhoang_TVN-023/dataC1K11/Bai1-Dang3.tex | mã gốc 1D1V1-4
\begin{vd}[MĐ3]%[1D1V1-4]
	Một vệ tinh được định vị tại vị trí $A$ trong không gian. Từ vị trí $A$, vệ tinh bắt đầu chuyển động quanh Trái Đất theo quỹ đạo là đường tròn với tâm là tâm $O$ của Trái Đất, bán kính $9~000$ km. Biết rằng vệ tinh chuyển động hết một vòng của quỹ đạo trong $2$ giờ. Vệ tinh chuyển động được quãng đường $200~000$ km sau bao nhiêu giờ (làm tròn đến kết quả hàng phần chục)?
	\shortans{3,5}
	\loigiai{
		Vệ tinh chuyển động được quãng đường $200~000$ km. Gọi $x$ là thời gian vệ tinh chuyển động. Khi đó
			$$200~000=x\cdot 2\pi \cdot 9~000\Leftrightarrow x\approx 3{,}5\text{ giờ}.$$
			
	}
\end{vd}

% #1994 | Archive/2024/SP-K10-K11-LyVanHoang_TVN-023/11. C1-LuongGiac-lvhoang_TVN-023/dataC1K11/Bai1-Dang3.tex | mã gốc 1D1H1-4
\begin{bt}%[1D1H1-4]
	Bánh xe có đường kính (tính cả lốp xe) là $55$ cm. Nếu xe chạy với vận tốc $40$ km/h thì trong một giây bánh xe quay được bao nhiêu vòng?
	\loigiai
	{
		Chu vi của bánh xe bằng $55\pi$ cm.\\
		Trong $1$ giây xe đi được $\dfrac{40\cdot 10^5}{3600} = \dfrac{10000}{9}$ cm.\\
		Như vây, trong $1$ giây bánh xe quay được $\dfrac{10000}{9 \cdot 55\pi} = \dfrac{2000}{99\pi}$ vòng.
	}
\end{bt}

% #1995 | Archive/2024/SP-K10-K11-LyVanHoang_TVN-023/11. C1-LuongGiac-lvhoang_TVN-023/dataC1K11/Bai1-Dang3.tex | mã gốc 1D1H1-4
\begin{bt}%[1D1H1-4]
	Kim phút và kim giờ của đồng hồ lớn Bưu điện Hà Nội theo thứ tự dài $1{,}75$ mét và $1{,}26$ mét. Hỏi trong $15$ phút, mũi kim phút và kim giờ vạch được cung tròn có độ dài bằng bao nhiêu mét?
	\loigiai
	{
		Trong $15$ phút kim giờ vạch được một cung tròn có số đo bằng $\dfrac{\pi}{24}$ và kim phút vạch được một cung tròn có số đo bằng $\dfrac{\pi}{2}$.\\
		Trong $15$ phút kim giờ vạch được cung tròn có độ dài bằng $1{,}26\cdot \dfrac{\pi}{24} = \dfrac{21\pi}{400} \approx 0{,}16$ m.\\
		Trong $15$ phút kim phút vạch được cung tròn có độ dài bằng $1{,}75 \cdot \dfrac{\pi}{2} = \dfrac{7\pi}{8} \approx 2{,}75$ m.
	}
\end{bt}

% #1996 | Archive/2024/SP-K10-K11-LyVanHoang_TVN-023/11. C1-LuongGiac-lvhoang_TVN-023/dataC1K11/Bai1-Dang3.tex | mã gốc 1D1H1-4
\begin{bt}%[1D1H1-4]
	Bánh xe của người đi xe đạp quay được $11$ vòng trong $5$ giây. Tính độ dài quãng đường mà người đi xe đã đi được trong $1$ phút, biết rằng đường kính của xe đạp là $680$ mm.
	\loigiai{
		Đổi $1$ phút $=60$ s.\\
		Trong $60$ giây, bánh xe quay được số vòng là $\dfrac{11}{5}\cdot60=132$ (vòng).\\
		Chu vi mỗi vòng xe là $680\pi$ (mm).\\
		Độ dài quãng đường người đó đi trong $1$ phút là $132\cdot680\pi=89760\pi$ (mm).
	}
\end{bt}

% #1997 | Archive/2024/SP-K10-K11-LyVanHoang_TVN-023/11. C1-LuongGiac-lvhoang_TVN-023/dataC1K11/Bai1-Dang3.tex | mã gốc 1D1H1-4
\begin{bt}%[1D1H1-4]
	\immini{Hải lí là một đơn vị chiều dài hàng hải, được tính bằng độ dài một cung chắn một góc $\alpha=\left(\dfrac{1}{60}\right)^{\circ}$ của đường kinh tuyến (Hình 17). Đổi số đo $\alpha$ sang radian và cho biết 1 hải lí bằng khoảng bao nhiêu kilômét, biết bán kính trung bình của Trái Đất là $6371\,$ km. Làm tròn kết quả đến hàng phần trăm.}
	{{\begin{tikzpicture}[scale=1.2, font=\footnotesize,line join=round, line cap=round, >=stealth];
				\draw[fill=cyan ] (0,0) circle (1.7);
				\draw (-2,0) -- (2,0) ;
				\draw (0,-2)node[right] {\scriptsize $\text{Cực Nam}$} -- (0,2) node[right] {\scriptsize $\text{Cực Bắc}$};
				\path
				(0,0) coordinate (O)
				(48:1.7) coordinate (B)
				(55:1.7) coordinate (C);
				\draw  (0:0)--(48:1.7) (0:0)--(55:1.7);
				\draw  (62:2) node[right,rotate=-43]{$\text{hải lí}$};
				\draw (0:1.7) arc (0:360:1.7);
				\draw (15:0.9) node{$\alpha=\left(\tfrac{1}{60}\right)^{\circ}$};
				\draw (0.5,0) node[below] {Đường xích đạo};
				\pic["$ $"{shift={(1pt,3pt)}},draw,angle radius=6mm,angle eccentricity=1.79] {angle = B--O--C};
				\path (270:2.8) node{Hình 17};
	\end{tikzpicture}}}
	\loigiai{
		$1$ hải lí = $\alpha\cdot R=\dfrac{1\cdot \pi }{60\cdot 180}\cdot 6371 \approx 1{,}85$ km.
		
	}
	
\end{bt}

% #1998 | Archive/2024/SP-K10-K11-LyVanHoang_TVN-023/11. C1-LuongGiac-lvhoang_TVN-023/dataC1K11/Bai1-Dang3.tex | mã gốc 1D1V1-4
\begin{bt}%[1D1V1-4]
	Một vệ tinh được định vị tại vị trí $A$ trong không gian. Từ vị trí $A$, vệ tinh bắt đầu chuyển động quanh Trái Đất theo quỹ đạo là đường tròn với tâm là tâm $O$ của Trái Đất, bán kính $9000 \mathrm{~km}$. Biết rằng vệ tinh chuyển động hết một vòng của quỹ đạo trong $2~$ giờ. Hãy tính quãng đường vệ tinh đã chuyển động được sau $30$ phút.
	\loigiai{
		Trong $2~$ giờ vệ tinh đi được quãng đường là $2\pi\cdot 9000=18000\pi $ km.\\
		Quãng đường vệ tinh chuyển động được trong $30$ phút $=\dfrac{1}{2}$ giờ là $\dfrac{1}{4}\cdot 18000\pi =4500\pi $ km.
	}
\end{bt}

% #2001 | Archive/2024/SP-K10-K11-LyVanHoang_TVN-023/11. C1-LuongGiac-lvhoang_TVN-023/dataC1K11/Bai1-Dang3.tex | mã gốc 1D1H1-4
\begin{ex}%[1D1H1-4]
	Từ một vị trí ban đầu trong không gian, vệ tinh $X$ chuyển động theo quỹ đạo là một đường tròn quanh Trái Đất và luôn cách tâm Trái Đất một khoảng bằng $9200$ km. Sau $2$ giờ thì vệ tinh $X$ hoàn thành hết một vòng di chuyển. Hãy cho biết tính đúng sai của mỗi khẳng định sau.
	\choiceTF[t]
	{\True Quãng đường vệ tinh $X$ chuyển động được sau $1$ giờ là $28902{,}65$ (km) (kết quả được làm tròn đến hàng phần trăm) (kết quả được làm tròn đến hàng phần trăm)}
	{\True Quãng đường vệ tinh $X$ chuyển động được sau $1{,}5$ giờ là $43353{,}98$ (km) (kết quả được làm tròn đến hàng phần trăm)}
	{Sau khoảng $5{,}3$ giờ thì $X$ di chuyển được quãng đường $240000$ km (kết quả được làm tròn đến hàng phần trăm)}
	{\True Giả sử vệ tinh di chuyển theo chiều dương của đường tròn, sau $4{,}5$ giờ thì vệ tinh vẽ nên một góc $\dfrac{9\pi}{2}$ rad}
	\loigiai{
		\begin{itemchoice}
			\itemch Một vòng di chuyển của $X$ chính là chu vi đường tròn:
			\[
			C=2\pi R=2\pi\cdot 9200=18400\pi\,\text{km}.
			\]
			Sau $1$ giờ, vệ tinh di chuyển nửa đường tròn với quãng đường là $\dfrac{1}{2}C=9200\pi\approx 28902{,}65$ (km).
			\itemch Sau $1{,}5$ giờ, vệ tinh di chuyển được $\dfrac{1{,}5\cdot 1}{2}$ đường tròn (hay $\dfrac{3}{4}$ đường tròn), quãng đường là $\dfrac{3}{4}C=\dfrac{3}{4}\cdot 18400\pi=13800\pi\approx 43353{,}98$ (km).
			\itemch Sau $5{,}3$ giờ, vệ tinh di chuyển được $\dfrac{5,3\cdot 1}{2}$ đường tròn (hay $\dfrac{53}{20}$ đường tròn), quãng đường là $\dfrac{53}{20}C=\dfrac{53}{20}\cdot 18400\pi=48760\pi\approx 153184{,}05$ (km).
			\itemch Sau $4{,}5$ giờ thì số vòng tròn mà vệ tinh $X$ di chuyển được là $\dfrac{4{,}5}{2}=\dfrac{9}{2}$ (vòng).\\
			Số đo góc lượng giác thu được là $\dfrac{9}{4}\cdot 2\pi=\dfrac{9\pi}{2}$ (rad).
		\end{itemchoice}
	}
\end{ex}

% #2002 | Archive/2024/SP-K10-K11-LyVanHoang_TVN-023/11. C1-LuongGiac-lvhoang_TVN-023/dataC1K11/Bai1-Dang3.tex | mã gốc 1D1H1-4
\begin{ex}%[1D1H1-4]
	\immini{Cho hình vẽ bên, biết đường tròn có bán kính $R=\dfrac{9}{\pi}$ cm. Tính độ dài của cung tròn $AM$.}{ \begin{tikzpicture}[declare function={r=1.5;},>=stealth]					
			\draw[->] (-1.35*r,0)--(1.5*r,0) node[shift={(-100:7pt)},font=\normalsize]{$ x $};
			\draw[->] (0,-1.35*r)--(0,1.5*r) node[shift={(190:7pt)},font=\normalsize]{$ y $};			
			\path (0,0) coordinate(O)
			(0:r) coordinate (A)
			(180:r) coordinate (B)	
			(120:r) coordinate (M);	
			\draw (O) circle (r);
			\draw[black!80] (A)--(O)--(M);		
			\path[->] pic["\scriptsize$\alpha=120^\circ$", angle eccentricity=1.2,draw,angle radius=12pt,right]{angle= A--O--M};		
			\foreach \t/\g in {O/-145,A/-37,M/90,B/150}{
				\draw[fill=white] (\t) circle (1pt) node[shift={(\g:7pt)},font=\scriptsize]{$ \t $};} 
			
	\end{tikzpicture} }
	\shortans{$12$}
	\loigiai{
		Ta có $120^\circ=\dfrac{4\pi}{3}$.\\
		Vậy độ dài cung tròn $AM$ là $\ell=\dfrac{9}{\pi}\cdot\dfrac{4\pi}{3}=12$ cm.
		}
\end{ex}

% #2003 | Archive/2024/SP-K10-K11-LyVanHoang_TVN-023/11. C1-LuongGiac-lvhoang_TVN-023/dataC1K11/Bai1-Dang3.tex | mã gốc 1D1V1-4
\begin{ex}%[1D1V1-4]
	\immini{Trên đường tròn lượng giác gốc $A$ có bán kính $R=\dfrac{24}{\pi}$, $M$ là điểm nằm trên đường tròn lượng giác thuộc góc phần tư thứ $I$ sao cho $\widehat{AOM}=\dfrac{\pi}{3}$. Gọi $N$ Là điểm đối xứng với $M$ qua trục $Oy$. Tính độ dài cung $AN$.
	}{		
		%-------------------------
		\begin{tikzpicture}[declare function={r=1.5;},>=stealth]					
			\draw[->] (-1.35*r,0)--(1.5*r,0) node[shift={(-100:7pt)},font=\normalsize]{$ x $};
			\draw[->] (0,-1.35*r)--(0,1.25*r) node[shift={(190:7pt)},font=\normalsize]{$ y $};
			\path (0,0) coordinate(O);
			\path (0:r) coordinate (A)
			(60:r) coordinate (M)
			(120:r) coordinate (N);
			\draw (O) circle (r);
			\draw(M)--(O)--(N)--cycle;
			\foreach \t/\g in {O/-132,A/-37,M/60,N/120}{
				\draw[fill=white] (\t) circle (1pt) node[shift={(\g:7pt)},font=\scriptsize]{$ \t $};}				
	\end{tikzpicture}}
	\shortans{$16$}
	\loigiai{
		\immini{Lấy thêm điểm $P$ như trong hình vẽ.\\	
			Vì $\widehat{AOM}=\dfrac{\pi}{3} \Rightarrow\widehat{MOP}=\dfrac{\pi}{6}.$ 	
			Lại có $N$ Là điểm đối xứng với $M$ qua trục $Oy\\ \Rightarrow\widehat{PON}=\dfrac{\pi}{6}.$\\
			$\Rightarrow \widehat{AON}=\widehat{AOM}+\widehat{MOP}+\widehat{PON}=\dfrac{\pi}{3}+\dfrac{\pi}{6}+\dfrac{\pi}{6}=\dfrac{2\pi}{3}\cdot$ \\
			Vậy độ dài cung $AN$ là $\ell=\dfrac{24}{\pi}\cdot \dfrac{2\pi}{3}=16$.
		}{
			%------------------
			\begin{tikzpicture}[declare function={r=1.5;},>=stealth]					
				\draw[->] (-1.35*r,0)--(1.5*r,0) node[shift={(-100:7pt)},font=\normalsize]{$ x $};
				\draw[->] (0,-1.35*r)--(0,1.25*r) node[shift={(190:7pt)},font=\normalsize]{$ y $};
				\path (0,0) coordinate(O)
				(0:r) coordinate (A)
				(60:r) coordinate (M)
				(90:r) coordinate (P)
				(120:r) coordinate (N);
				\draw (O) circle (r);
				\draw[green!80] (M)--(O)--(N)--cycle;		    		
				\foreach \t/\g in {O/-132,A/-37,M/60,N/120,P/55}{
					\draw[fill=white] (\t) circle (1pt) node[shift={(\g:7pt)},font=\scriptsize]{$ \t $};}
				
	\end{tikzpicture}}}
	%-----------------------------	
\end{ex}

% #2005 | Archive/2024/SP-K10-K11-LyVanHoang_TVN-023/11. C1-LuongGiac-lvhoang_TVN-023/dataC1K11/Bai1-Dang3.tex | mã gốc 1D1V1-4
\begin{ex}%[1D1V1-4]
	\immini{Trạm vũ trụ Quốc tế ISS (tên Tiếng Anh: International Space Station) nằm trong quỹ đạo tròn cách bề mặt Trái Đất khoảng $400$ km. Nếu trạm mặt đất theo dõi được trạm vũ trụ ISS khi nó nằm trong góc $45^\circ$ ở tâm của quỹ đạo tròn này phía trên ăng-ten theo dõi, thì trạm vũ trụ ISS đã di chuyển được bao nhiêu km trong khi nó đang được trạm mặt đất theo dõi? Giả sử rằng bán kính của Trái Đất là $6400$ km. Làm tròn kết quả đến hàng đơn vị.
	}
	{\begin{tikzpicture}[font=\footnotesize, line join=round, line cap=round, >=stealth,scale=0.8]
			\def\x{2}
			\def\y{3}
			\path 
			(0,0) coordinate (O)
			(0:\y) coordinate (A)	
			(70:\y) coordinate (B)
			(115:\y) coordinate (C)
			(70:2.5) node[right]{$400$ km}
			;
			\fill[gray!30] (O) circle(\x);
			\draw[dashed] (O) circle(\y);
			\draw[->] (O)--(A);
			\draw[->] (O)--(B);
			\draw[->] (O)--(C);
			\begin{scope}
				\clip (C)--(O)--(B);
				\draw[fill=yellow] (O) circle(.5);
				\path ($(O)+(90:0.8)$) node{$45^\circ$};
			\end{scope}
			\foreach \d in{O,A,B}
			\fill[black](\d) circle (1pt);
			%	($(\x)+(\g:5mm)$) node{$\x$};
	\end{tikzpicture}}
	\shortans{$5341$}
	\loigiai{
		Bán kính quỹ đạo của trạm vũ trụ quốc tế là $R=6400+400=6800$ km.\\
		Đổi $45^\circ =45\cdot \dfrac{\pi}{180}=\dfrac{\pi}{4}$ rad.\\
		Vậy trong khi được trạm mặt đất theo dõi, trạm ISS đã di chuyển một quãng đường có độ dài\\ $\ell=R\cdot \alpha =6800\cdot \dfrac{\pi}{4}\approx 5341$ km.
	}
\end{ex}

% #2006 | Archive/2024/SP-K10-K11-LyVanHoang_TVN-023/11. C1-LuongGiac-lvhoang_TVN-023/dataC1K11/Bai1-Dang3.tex | mã gốc 1D1V1-4
\begin{ex}%[1D1V1-4]%[1D1N1-5]%[1D1N1-5]
	Bánh xe của người đi xe đạp có bán kính $R=\dfrac{100}{\pi}$ cm quay được $12$ vòng trong $5$ giây. Giả sử tại thời điểm bắt đầu quay van xe nằm ở vị trí $V$ vuông góc với mặt đất. Hỏi trong một giây người đi xe đạp đi được quãng đường bao nhiêu (đơn vị tính bằng cm)?
	\begin{center}
		\begin{tikzpicture}[scale=0.8, font=\footnotesize, line join=round, line cap=round,>=stealth]
			\shade[shading=radial] (-1.5,-2) rectangle (1.5,-1.6);
			\shade[shading=radial] (4.5,-2) rectangle (7.5,-1.6);
			\def\Radius{1.5cm}
			\foreach \a in {0, 10, ..., 350} { 
				%\draw[red] (0, 0) -- (\a:\Radius);
				\draw[black] (0, 0) -- (\a+1.5:\Radius);
				%\draw[gray] (0, 0) -- (\a+3:\Radius);
				%\draw[gray!50] (0, 0) -- (\a+4.5:\Radius);
				\draw[blue] (0, 0) circle[radius=\Radius];
				\draw[line width=4pt] (0, 0) circle[radius=\Radius+0.15cm];
			}
			\foreach \h in {0, 60, ..., 300} {
				\draw[color=black,line width=0.5mm,rotate around={\h:(0,0)}] (0,0) edge[-,bend left=40] (0.2cm,0) edge[-,bend right=40] (0.2cm,0);
			}
			\draw[line width=1pt] (0,0) circle[radius=0.2cm];
			\draw[line width=4pt,rounded corners=1pt,rotate=0] (1.68,0.6)--(1.94,0.54);
			\draw[line width=0.3mm,yellow!60!black] (1.8,0) arc (0:150:1.8cm)--(0,0);
			\draw[line width=0.5mm,yellow!60!black] (1.8,0) arc (0:160:1.8cm);
			\foreach \b in {0, 10, ..., 350} { 
				% \draw[color=red,rotate around={\b:(6,0)}] (6, 0) -- (7.5,0);
				\draw[black,rotate around={\b+1.5:(6,0)}] (6, 0) -- (7.5,0);
				% \draw[gray,rotate around={\b+3:(6,0)}] (6, 0) -- (7.5,0);
				% \draw[gray!50,rotate around={\b+4.5:(6,0)}] (6, 0) -- (7.5,0);
				\draw[blue] (6, 0) circle[radius=\Radius];
				\draw[line width=4pt] (6, 0) circle[radius=\Radius+0.15cm];
			}
			\draw[black,line width=2pt] (6, 0)node[above right,fill=white,rounded corners]{$O$}--++ (-90:1.5)node[below,fill=white,rounded corners]{$V$};
			\draw[line width=0.3mm,yellow!60!black] (6,1.8) arc (90:190:1.8cm)--(6,0);
			\draw[line width=0.5mm,yellow!60!black] (6,1.8) arc (90:200:1.8cm);
			\draw[line width=1pt] (0,0.2cm)--(2.2,0.5cm) (0,-0.2cm)--(2.2,-0.5cm);
			\draw[line width=2pt] (2.2,0)--(2.6,-0.6) (2.2,0)--(1.8,0.6);
			\draw[line width=1mm,rounded corners=8pt,rotate around={25:(6,0)}] (6,0)--(6,1.3cm)--++(80:2.1cm);
			\draw[line width=1mm,rounded corners=8pt,rotate around={27:(6,0)}] (6,0)--(6,1.3cm)--++(80:2.1cm);
			\draw[line width=1mm,rounded corners=8pt,rotate around={26:(6,0)}] (6,0)--(6,1.3cm)--++(80:2.1cm)coordinate(a);
			\draw[line width=1.5mm] (a)--++(-80:0.15cm)--++(0:0.5cm)--++(180:0.05cm)coordinate(b);
			\draw[line width=1.1mm,rounded corners=4pt] (b)--++(-30:0.6cm)coordinate(c)--++(-160:0.4cm) (b)--++(150:0.6cm)coordinate(d)--++(200:0.4cm);
			\draw[line width=0.4mm,rounded corners=0pt] (c)--++(-95:0.1cm)--++(-160:0.25cm) (d)--++(-95:0.1cm)--++(200:0.25cm);
			\draw[line width=2mm,rotate around={35:(2.2,0)}] (2.2,0)--++(3.7,0)--++(75:0.4cm)--++(155:3.5)coordinate(e)--(2.2,0);
			\draw[line width=1mm] (0,0)--(2.2,0) (0,0)--(45:2.5cm) (0,0)--(46.5:2.5) (0,0)--(2:2.2);
			\draw[line width=1.2mm] (e)--++(106:0.9cm)--++(-74:0.1cm)coordinate(d);
			\foreach \c in {0, 72, ..., 288} {\draw[color=red,line width=1.8mm,rotate around={\c:(2.2,0)}] (2.2,0) edge[-,bend left=30] (2.69,0);}
			\draw[line width=2pt] (2.2,0) circle[radius=0.5cm];
			\fill (2.2,0) circle[radius=0.1cm];
			\draw[line width=4pt,rounded corners=1pt,rotate=0] (2.4,-0.58)--(2.68,-0.56);
			\draw[line width=2pt] (2.2,0)--(2.6,-0.6);
			\shade[top color=white,bottom color=black,rounded corners=2pt,rotate around={10:(d)}] (d)--++(220:2mm)--++(100:2mm)--++(180:3mm)--++(110:1.5mm)--++(160:2mm)--++(110:3.5mm)coordinate(i)--++(-15:10mm)--++(-5:5mm)--++(-20:3mm)coordinate(j)--++(-50:0.7mm)--++(180:4mm)--++(170:1mm)--++(-175:2mm)--(d);
			\shade[top color=gray!50,bottom color=black,rounded corners=2pt,rotate around={10:(d)}] (d)--++(220:2mm)--++(100:2mm)--++(180:3mm)--++(110:1.5mm)--++(160:2mm)--++(110:3.5mm)--++(-50:4.5mm)--++(-3:13mm)--++(-10:2mm)--++(180:4mm)--++(170:1mm)--++(-175:2mm)--(d);
			\filldraw[black] (0,0)circle(2pt) (6,0)circle(2pt);
			\draw[line width=0.3mm] (0,0)--++(105:2cm)--++(-2:0.7cm)coordinate(g)--+(-10:1.5cm) (g)--(0,0);
			\shade[top color=gray!10!white,bottom color=gray] (-0.8cm,1.9cm)--++(-3:1.4cm) arc (-90:90: 0.4mm)--++(174:1.4cm) arc (90:270: 0.65mm);
			\shade[top color=gray!20,bottom color=black] (-0.8cm,1.9cm)--++(-3:1.4cm) arc (-90:90: 0.4mm)--++(177:1.4cm) arc (90:270: 0.4mm);
			%\node[text=\fontsize{2pt}{0}\selectfont,color=white,rotate around={10:(3,2.72)}] at (3,2.72) {\tiny Nhóm Toán và \LaTeX};
		\end{tikzpicture}
	\end{center}
	\shortans{$480$}
	\loigiai{
		Trong một giây bánh xe quay được $\dfrac{12}{5}$ vòng.\\
		Suy ra trong một giây bánh xe quay được góc $\alpha=\dfrac{12}{5}\cdot 2 \pi\ (\mathrm{rad})$.\\
		Vậy trong một giây người đi xe đạp đi được quãng đường là $\ell=\dfrac{12}{5}\cdot 2 \pi\cdot \dfrac{100}{\pi}=480$ cm.}
\end{ex}

% #2007 | Archive/2024/SP-K10-K11-LyVanHoang_TVN-023/11. C1-LuongGiac-lvhoang_TVN-023/dataC1K11/Bai-3-Dang-2-Xetsudongbiennghichbienhslg.tex | mã gốc 1D1N4-3
\begin{vd}%[BG 11 - Nguyen Huynh]%[1D1N4-3]
	[Mức độ 1]
	Hãy nêu tất cả các hàm số trong các hàm số $ y = \sin x$, $y = \cos x$, $y = \tan x$, $y = \cot x $ thỏa mãn điều kiện đồng biến và nhận giá trị âm trong khoảng $ \left ( - \dfrac{\pi}{2};0 \right ) $.
	\loigiai{
		Ta thấy  hàm số $ y = \tan x  $ và $ y = \sin x $ thỏa mãn đề bài.
	}
\end{vd}

% #2008 | Archive/2024/SP-K10-K11-LyVanHoang_TVN-023/11. C1-LuongGiac-lvhoang_TVN-023/dataC1K11/Bai-3-Dang-2-Xetsudongbiennghichbienhslg.tex | mã gốc 1D1H4-3
\begin{vd}%[BG 11 - Nguyen Huynh]%[1D1H4-3]
	[Mức độ 1]
	Xét tính đồng biến, nghịch biến của hàm số $y=\sin 2x$ trên khoảng $\left(0;\dfrac{\pi}{4}\right)$.
	\loigiai{
		Ta có $x\in \left(0;\dfrac{\pi}{4}\right)\Rightarrow 2x\in \left(0;\dfrac{\pi}{2}\right)$ thuộc góc phần tư thứ I nên hàm số $y=\sin 2x$ đồng biến trên khoảng $\left(0;\dfrac{\pi}{4}\right)$.}
\end{vd}

% #2009 | Archive/2024/SP-K10-K11-LyVanHoang_TVN-023/11. C1-LuongGiac-lvhoang_TVN-023/dataC1K11/Bai-3-Dang-2-Xetsudongbiennghichbienhslg.tex | mã gốc 1D1H4-3
\begin{vd}%[BG 11 - Nguyen Huynh]%[1D1H4-3]
	[Mức độ 1]
	Xét tính đồng biến, nghịch biến của hàm số $y=\sin\left(x-\dfrac{\pi}{4}\right)$ trên khoảng $\left(0;\dfrac{\pi}{2}\right)$.
	\loigiai{
		Trên đường tròn lượng giác, khi $x\in \left(0;\dfrac{\pi}{2}\right)\Rightarrow x-\dfrac{\pi}{4} \in \left(-\dfrac{\pi}{4};\dfrac{\pi}{4}\right)\subset\left(-\dfrac{\pi}{2};\dfrac{\pi}{2}\right)$ nên hàm số $y=\sin\left(x-\dfrac{\pi}{4}\right)$ đồng biến trên khoảng $\left(0;\dfrac{\pi}{2}\right)$.
	}
\end{vd}

% #2012 | Archive/2024/SP-K10-K11-LyVanHoang_TVN-023/11. C1-LuongGiac-lvhoang_TVN-023/dataC1K11/Bai-3-Dang-2-Xetsudongbiennghichbienhslg.tex | mã gốc 1D1N4-3
\begin{bt}%[BG 11 - Nguyen Huynh]%[1D1N4-3]
	[Mức độ 1]
	Hàm số nào sau đây đồng biến trên khoảng $\left(\dfrac{\pi}{2};\pi\right)$?
	\begin{itemize}
		\item $y=\tan x$.
		\item $y=\cos x$.
		\item $y=\sin x$.
		\item $y=\cot x$.
	\end{itemize}
	\loigiai{
		Hàm số $y=\tan x$ đồng biến trên các khoảng xác định của nó nên nó đồng biến trên khoảng $\left(\dfrac{\pi}{2};\pi\right)$. }
\end{bt}

% #2013 | Archive/2024/SP-K10-K11-LyVanHoang_TVN-023/11. C1-LuongGiac-lvhoang_TVN-023/dataC1K11/Bai-3-Dang-2-Xetsudongbiennghichbienhslg.tex | mã gốc 1D1H4-3
\begin{bt}%[BG 11 - Nguyen Huynh]%[1D1H4-3]
	[Mức độ 2]
	Hàm số nào sau đây nghịch biến trên khoảng $\left(0; \dfrac{\pi}{2}\right)$?
	\begin{itemize}
		\item $y = \cos x$.
		\item $y = \tan x$.
		\item $y = \sin x$.
		\item $y = - \cot x$.
	\end{itemize}
	\loigiai{
		Các hàm số $y=\sin x$, $y=\tan x$, $y=-\cot x$ đồng biến trên khoảng $\left(0; \dfrac{\pi}{2}\right)$, hàm số $y=\cos x$ nghịch biến trên khoảng $\left(0; \dfrac{\pi}{2}\right)$.}
\end{bt}

% #2014 | Archive/2024/SP-K10-K11-LyVanHoang_TVN-023/11. C1-LuongGiac-lvhoang_TVN-023/dataC1K11/Bai-3-Dang-2-Xetsudongbiennghichbienhslg.tex | mã gốc 1D1H4-3
\begin{bt}%[BG 11 - Nguyen Huynh]%[1D1H4-3]
	[Mức độ 2]
	Trong các hàm số $y=\sin x$, $y=\cos x$, $y=\tan x$, hàm số nào đồng biến trên khoảng $(-\pi;0)$?
	\loigiai{
		Trên khoảng $(-\pi;0)$ hàm số $y=\cos x$ tăng từ $-1$ đến $1$.\\
		Vậy hàm số $y=\cos x$ đồng biến trên khoảng $(-\pi;0)$.
	}
\end{bt}

% #2015 | Archive/2024/SP-K10-K11-LyVanHoang_TVN-023/11. C1-LuongGiac-lvhoang_TVN-023/dataC1K11/Bai-3-Dang-2-Xetsudongbiennghichbienhslg.tex | mã gốc 1D1H4-3
\begin{bt}%[BG 11 - Nguyen Huynh]%[1D1H4-3]
	[Mức độ 2]
	Cho hàm số $y=\sin 2x$. Xét tính đơn điệu của hàm số trên khoảng $\left(0;\dfrac{\pi}{2}\right)$.
	\loigiai{
		Ta thấy, khi $x$ chạy từ $0$ đến $\dfrac{\pi}{2}$ thì $2x$ chạy từ $0$ đến $\pi$. Vậy
		\begin{itemize}
			\item Khi $x$ chạy từ $0$ đến $\dfrac{\pi}{4}$ thì $2x$ chạy từ $0$ đến $\dfrac{\pi}{2}$ nên $\sin 2x$ tăng từ $0$ đến $1$.
			\item Khi $x$ chạy từ $\dfrac{\pi}{4}$ đến $\dfrac{\pi}{2}$ thì $2x$ chạy từ $\dfrac{\pi}{2}$ đến $\pi$ nên $\sin 2x$ giảm từ $1$ về $0$.
		\end{itemize}
		Ta có bảng biến thiên như sau:
		\begin{center}
			\begin{tikzpicture}
				\tkzTabInit[nocadre=false,lgt=1,espcl=2.5]
				{$x$ /1,$y$ /2}
				{$0$,$\dfrac{\pi}{4}$,$\dfrac{\pi}{2}$}
				\tkzTabVar{-/$0$, +/$1$,-/$0$}
			\end{tikzpicture}
		\end{center}
	Vậy hàm số \(y=\sin 2x\) đồng biến trên khoảng $\left(0;\dfrac{\pi}{4}\right)$; nghịch biến trên khoảng $\left(\dfrac{\pi}{4}; \dfrac{\pi}{2}\right)$.
	}
\end{bt}

% #2024 | Archive/2024/SP-K10-K11-LyVanHoang_TVN-023/11. C1-LuongGiac-lvhoang_TVN-023/dataC1K11/Chuong1-Bai4-Dang5-Phuongtrinhduavecoban.tex | mã gốc 1D1H5-5
\begin{vd}%[Bài giảng 11, Trần Tony]%[1D1H5-5]
	Giải các phương trình sau
	\begin{listEX}[2]
		\item $2\sin x+\sqrt{2}=0$;
		\item $\sin2x-\cos x+2\sin x=1$;
		\item $3\sin ^2 x-5\sin x+2=0$;
		\item $\sqrt{3}\tan^2 x-2\tan x+\sqrt{3}=0$;
		\item $2\cos^2 2x-5\cos 2x+2=0$;
		\item $\sin^2\dfrac{x}{2}+\sin\dfrac{x}{2}-2=0$.
	\end{listEX}
	\loigiai{
		\begin{enumerate}[a)]
			\item $2\sin x+\sqrt{2}=0\Leftrightarrow\sin x=-\dfrac{\sqrt{2}}{2}\Leftrightarrow\sin x=\sin\left(-\dfrac{\pi}{4}\right)\Leftrightarrow\hoac{&x=-\dfrac{\pi}{4}+k2\pi\\&x=\dfrac{5\pi}{4}+k2\pi}\,(k\in\mathbb{Z})$;
			\item $\sin2x-\cos x+2\sin x=1\Leftrightarrow2\sin x\cos x-\cos x+2\sin x-1=0\Leftrightarrow(2\sin x-1)(\cos x+1)=0\\\Leftrightarrow\hoac{&\sin x=\dfrac{1}{2}\\&\cos x=-1}\Leftrightarrow\hoac{&\sin x=\sin\dfrac{\pi}{6}\\&x=(2k+1)\pi}\Leftrightarrow\hoac{&x=\dfrac{\pi}{6}+k2\pi\\&x=\dfrac{5\pi}{6}+k2\pi\\&x=(2k+1)\pi}\,(k\in\mathbb{Z})$;
			\item Đặt $t=\cos x$, $-1\le t\le 1$, phương trình đã cho trở thành $3t^2-5t+2=0$, ta được $t=1$ hoặc $t=\dfrac{2}{3}$.\\
			Với $t=1$ ta có $\cos x=1 \Leftrightarrow x= k2\pi,\, (k\in\mathbb{Z})$.\\
			Với $t=\dfrac{2}{3}$ ta có $\cos x=\dfrac{2}{3}=\cos\alpha\Leftrightarrow x=\pm \alpha +k2\pi,\, (k\in\mathbb{Z})$.\\
			Vậy tập nghiệm của phương trình đã cho là $S=\left\{ k2\pi, \pm \alpha+k2\pi, k\in\mathbb{Z} \right\}$.
			\item Đặt $t=\tan x$, phương trình đã cho trở thành $\sqrt{3} t^2-2t+\sqrt{3}=0$. Phương trình này vô nghiệm. Vậy phương trình đã cho vô nghiệm.
			\item $2\cos^2 2x-5\cos 2x+2=0\Leftrightarrow\hoac{&\cos2x=2\\&\cos2x=\dfrac{1}{2}}\Leftrightarrow\cos2x=\cos\dfrac{\pi}{3}\Leftrightarrow x=\pm\dfrac{\pi}{6}+k\pi,\, (k\in\mathbb{Z})$;
			\item $\sin^2\dfrac{x}{2}+\sin\dfrac{x}{2}-2=0\Leftrightarrow\hoac{&\sin\dfrac{x}{2}=1\\&\sin\dfrac{x}{2}=-2}\Leftrightarrow\dfrac{x}{2}=\dfrac{\pi}{2}+k2\pi\Leftrightarrow x=\pi+k4\pi,\, (k\in\mathbb{Z})$.
		\end{enumerate}
	}
\end{vd}

% #2025 | Archive/2024/SP-K10-K11-LyVanHoang_TVN-023/11. C1-LuongGiac-lvhoang_TVN-023/dataC1K11/Chuong1-Bai4-Dang5-Phuongtrinhduavecoban.tex | mã gốc 1D1H5-5
\begin{vd}%[Bài giảng 11, Trần Tony]%[1D1H5-5]
	Giải các phương trình sau
	\begin{listEX}[2]
		\item $\sqrt{3}\sin x+\cos x=2$.
		\item $3\sin x + \cos x = 2$.
		\item $\sqrt{3}\cos 2x+\sin 2x-\sqrt{3}=0$.
		\item $\sin x+3\cos x=5$.
	\end{listEX}
	\loigiai{
		\begin{enumerate}[a)]
			\item Phương trình tương đương $\dfrac{\sqrt{3}}{2}\sin x+\dfrac{1}{2}\cos x=1\Leftrightarrow \sin \dfrac{\pi}{3}\sin x +\cos \dfrac{\pi}{3}\cos x=1$\\
			$\Leftrightarrow \cos \left( x-\dfrac{\pi}{3}\right) =1\Leftrightarrow x=\dfrac{\pi}{3}+k2\pi$ $(k\in \mathbb{Z})$.
			
			\item  $3\sin x + \cos x = 2 \Leftrightarrow \dfrac{3}{\sqrt {10} }\sin x + \dfrac{1}{\sqrt {10}}\cos x = \dfrac{2}{\sqrt {10} }$.\hfill$(1)$\\
			Đặt $\dfrac{1}{\sqrt {10}} = \cos \alpha $ và
			$\dfrac{3}{\sqrt {10}} = \sin \alpha $,  thay vào (1), ta được
			{\allowdisplaybreaks
				\begin{align*}
					&\cos x\cos \alpha  + \sin x\sin \alpha  = \frac{2}{\sqrt {10} } \Leftrightarrow \cos (x - \alpha ) = \frac{2}{\sqrt {10}}=\cos\beta\\
					\Leftrightarrow& 
					\hoac{&x - \alpha  = \beta + k2\pi \\&
						x - \alpha  =  - \beta + k2\pi}
					\Leftrightarrow\hoac{&x = \alpha  + \beta + k2\pi \\&x = \alpha  - \beta + k2\pi.}
			\end{align*}}
			\item Ta có
			\begin{align*}
				\sqrt{3}\cos 2x+\sin 2x-\sqrt{3}=0&\Leftrightarrow \sqrt{3}\cos 2x+\sin 2x=\sqrt{3}.\\
				&\Leftrightarrow \dfrac{\sqrt{3}}{2}\cos 2x+\dfrac{1}{2}\sin 2x=\dfrac{\sqrt{3}}{2}\\
				&\Leftrightarrow \sin\left(2x+\dfrac{\pi}{3}\right)=\dfrac{\sqrt{3}}{2}\\
				&\Leftrightarrow \hoac{&2x+\dfrac{\pi}{3}=\dfrac{\pi}{3}+k2\pi\\ &2x+\dfrac{\pi}{3}=\dfrac{2\pi}{3}+k2\pi.}\Leftrightarrow \hoac{&x=k\pi\\ &x=\dfrac{\pi}{6}+k\pi}, (k\in\mathbb{Z}).
			\end{align*}
			\item Vì $1^2+3^2<5^2$ nên phương trình vô nghiệm.
		\end{enumerate}
	}
\end{vd}

% #2031 | Archive/2024/SP-K10-K11-LyVanHoang_TVN-023/11. C1-LuongGiac-lvhoang_TVN-023/dataC1K11/Chuong1-Bai4-Dang5-Phuongtrinhduavecoban.tex | mã gốc 1D1H5-5
\begin{bt}%[Bài giảng 11, Trần Tony]%[1D1H5-5]
	Giải các phương trình sau
	\begin{listEX}[2]
		\item $6\cos^2 x+5\sin x-7=0$.
		\item $\cos 2x+\cos x+1=0$.
		\item $3\sin^2 2x+7\cos 2x-3=0$.
		\item $7\tan x-4\cot x=12$.
	\end{listEX}
	\loigiai{
		\begin{enumerate}[a)]
			\item Đặt $t=\sin x$, $-1\le t\le 1$, phương trình đã cho trở thành $6(1-t^2)+5t-7=0$, ta được $t=\dfrac{1}{2}$ hoặc $t=\dfrac{1}{3}$.\\
			Với $t=\dfrac{1}{2}$, ta có $\sin x=\dfrac{1}{2} \Leftrightarrow \hoac{&x=\dfrac{\pi}{6}+k2\pi\\ &x=\dfrac{5\pi}{6}+k2\pi}\ (k\in\mathbb{Z})$.\\
			Với $t=\dfrac{1}{3}$, ta có $\sin x=\dfrac{1}{3}=\sin\alpha\Leftrightarrow \hoac{&x=\alpha + k2\pi\\ & x=\pi-\alpha + k2\pi}\ (k\in\mathbb{Z})$.
			
			\item Đặt $t=\cos x$, $-1\le t\le 1$, phương trình đã cho trở thành $2t^2-1+t+1=0$, ta được $t=0$ hoặc $t=-\dfrac{1}{2}$.\\
			Với $t=0$, ta có $\cos x=0\Leftrightarrow x=k2\pi$ ($k\in\mathbb{Z}$).\\
			Với $t=-\dfrac{1}{2}$, ta có $\cos x=-\dfrac{1}{2}\Leftrightarrow x=\pm \dfrac{2\pi}{3}+k2\pi$ ($k\in\mathbb{Z}$).
			\item Đặt $t=\cos 2x$, $-1\le t\le 1$. Phương trình đã cho trở thành
			\begin{eqnarray*}
				&&3\left(1-t^2\right) +7t -3=0\\
				&\Leftrightarrow &3t^2-7t=0\\
				&\Leftrightarrow &\hoac{&t=0\\ &t=\dfrac{7}{3}\quad(\text{loại}).}
			\end{eqnarray*}
			Ta có $\cos 2x=0$ hay $x=\dfrac{\pi}{4}+\dfrac{k\pi}{2}\ (k\in\mathbb{Z})$.\\
			Vậy tập nghiệm của phương trình đã cho là $S=\left\{\dfrac{\pi}{4}+\dfrac{k\pi}{2}, k\in\mathbb{Z}\right\}$.
			
			\item Đặt $t=\tan x$, $t\ne 0$. Phương trình đã cho trở thành
			\begin{eqnarray*}
				&&7t-\dfrac{4}{t}=12\\
				&\Leftrightarrow &7t^2-12t-4=0\\
				&\Leftrightarrow &\hoac{&t=2\\ &t=-\dfrac{2}{7}.}
			\end{eqnarray*}
			Với $t=2$ ta có $\tan x=2=\tan\alpha\Leftrightarrow x=\alpha +k\pi,\ (k\in\mathbb{Z})$.\\
			Với $t=-\dfrac{2}{7}$ ta có $\tan x=-\dfrac{2}{7}=\tan\beta\Leftrightarrow x=\beta+k\pi\ (k\in\mathbb{Z})$.\\
			Vậy tập nghiệm của phương trình đã cho là $S=\left\{\alpha +k\pi,\beta+k\pi, k\in\mathbb{Z} \right\}$.
		\end{enumerate}
	}
\end{bt}

% #2032 | Archive/2024/SP-K10-K11-LyVanHoang_TVN-023/11. C1-LuongGiac-lvhoang_TVN-023/dataC1K11/Chuong1-Bai4-Dang5-Phuongtrinhduavecoban.tex | mã gốc 1D1H5-5
\begin{bt}%[Bài giảng 11, Trần Tony]%[1D1H5-5]
	Giải các phương trình sau
	\begin{listEX}[2]
		\item $\dfrac{1 + \cos 4x - 2 \cos 2x}{\sin \left( x + \dfrac{\pi}{4} \right) \sin 2x} = 0$.
		\item $\dfrac{4\sin^22x+6\sin^2x-9-3\cos2x}{\cos x\left(\sin x-\dfrac{\sqrt{3}}{2}\right)}=0$.
	\end{listEX}
	\loigiai{
		\begin{enumerate}[a)]
			\item	Điều kiện xác định:\\
			$\sin \left( x + \dfrac{\pi}{4}\right) \sin 2x \ne 0 \Leftrightarrow \heva{&x \ne -\dfrac{\pi}{4}+k\pi \\& x \ne \dfrac{l\pi}{2}}$, với $k, l \in \mathbb{Z}$.\\
			Phương trình 
			\begin{eqnarray*}
				&&\dfrac{1 + \cos 4x - 2 \cos 2x}{\sin \left( x + \dfrac{\pi}{4}\right) \sin 2x} = 0
				\\&\Leftrightarrow& 2 \cos^2 2x - 2 \cos 2x = 0
				\\&\Leftrightarrow& \hoac{&\cos 2x = 0 \\& \cos 2x = 1 \text{ (loại)}}
				\\&\Leftrightarrow& x = \dfrac{\pi}{4} + \dfrac{k\pi}{2}	
			\end{eqnarray*}
			So với điều kiện xác định ở trên, ta có nghiệm của phương trình đã cho là $x = \dfrac{\pi}{4} + k\pi, k \in \mathbb{Z}$.
			\item Điều kiện:$\heva{&\cos x\ne0\\&\sin x\ne\dfrac{\sqrt{3}}{2}}\Leftrightarrow\heva{&x\ne\dfrac{\pi}{2}+k\pi\\&x\ne\dfrac{\pi}{3}+k\pi\\&x\ne\dfrac{2\pi}{3}+k\pi}$ \\
			Phương trình $\Leftrightarrow 4(1-\cos^22x)+3(1-\cos2x)-9-3\cos2x=0$\\
			$\Leftrightarrow4\cos^22x+6\cos2x+2=0\Leftrightarrow\hoac{&\cos2x=-1\\&\cos2x=-\dfrac{1}{2}}\Leftrightarrow\hoac{&x=\dfrac{\pi}{2}+k\pi\\&x=\pm\dfrac{\pi}{3}+k\pi}$ \\
			So điều kiện, nghiệm phương trình là $x=\dfrac{4\pi}{3}+k2\pi$ và $ x=-\dfrac{\pi}{3}+k2\pi$ với $k\in\mathbb{Z}$.
		\end{enumerate}
	}
\end{bt}

% #2033 | Archive/2024/SP-K10-K11-LyVanHoang_TVN-023/11. C1-LuongGiac-lvhoang_TVN-023/dataC1K11/Chuong1-Bai4-Dang5-Phuongtrinhduavecoban.tex | mã gốc 1D1H5-5
\begin{bt}%[Bài giảng 11, Trần Tony]%[1D1H5-5]
	Giải các phương trình sau
	\begin{listEX}[2]
		\item $\cos ^2x+\cos x-2=0$;
		\item $2\sin^2 x -5\sin x +2=0$;
		\item $3\tan^2 x-2\sqrt{3}\tan x+1=0$.
	\end{listEX}
	\loigiai{
		\begin{enumerate}[a)]
			\item Đặt $t=\cos x$, $-1\le t\le 1$. Phương trình trở thành $t^2+t-2$, ta được $t=1$ hoặc $t=-2$.\\Kết hợp điều kiện ta có $t=1$, do đó $\cos x=1$ hay $x=k2\pi\ (k\in\mathbb{Z})$.\\
			Vậy tập nghiệm của phương trình đã cho là $S=\left\{ k2\pi, k\in\mathbb{Z}  \right\}$.
			\item Đặt $t=\sin x$, $-1\le t\le 1$. Phương trình trở thành $2t^2-5t+2=0$, ta được $t=\dfrac{1}{2}$ hoặc $t=2$.\\Kết hợp điều kiện, ta có $t=\dfrac{1}{2}$, do đó $\sin x=\dfrac{1}{2}$ hay $x=\dfrac{\pi}{6}+k2\pi$ hoặc $x=\dfrac{5\pi}{6}+k2\pi$ ($k\in\mathbb{Z}$).\\
			Vậy tập nghiệm của phương trình đã cho là $S=\left\{\dfrac{\pi}{6}+k2\pi, \dfrac{5\pi}{6}+k2\pi, k\in\mathbb{Z}\right\}$.
			\item Đặt $t=\tan x$, phương trình đã cho trở thành $3t^2-2\sqrt{3}t+1=0$, ta được $t=\dfrac{1}{\sqrt{3}}$.\\Do đó $\tan x=\dfrac{1}{\sqrt{3}}$ hay $x=\dfrac{\pi}{6}+k\pi$ ($k\in\mathbb{Z}$).\\
			Vậy tập nghiệm của phương trình đã cho là $S=\left\{ \dfrac{\pi}{6}+k\pi, k\in\mathbb{Z}  \right\}$.
		\end{enumerate}
	}
\end{bt}

% #2039 | Archive/2024/SP-K10-K11-LyVanHoang_TVN-023/11. C1-LuongGiac-lvhoang_TVN-023/dataC1K11/Chuong1-Bai4-Dang5-Phuongtrinhduavecoban.tex | mã gốc 1D1H5-5
\begin{ex}%[Bài giảng 11, Trần Tony]%[1D1H5-5]
	Cho phương trình $\sin 4x+\sin 2x=\cos 4x+\cos 2x$
	\choiceTF[t]
	{Dùng công thứ biến đổi tổng thành tích, vế trái của phương trình đưa về dạng $\sin 3x\cos x$}
	{Dùng công thức biến đổi tổng thành tích, vế phải của phương trình đưa về dạng $\cos 3x\cos x$}
	{\True Nghiệm của phương trình đã cho là nghiệm của phương trình $\cos x=0$ và phương trình $\sin 3x=\cos 3x$}
	{\True Nghiệm của phương trình đã cho là $x=k2\pi$ và $x=\dfrac{\pi}{12}+k\dfrac{\pi}{3}$ ($k\in \mathbb{Z}$)}
	\loigiai{
		Dùng công thức biến đổi tổng thành tích, phương trình đưa về dạng
		\begin{eqnarray*}
			& &2\sin 3x\cos x=2\cos 3x\cos x\\
			&\Leftrightarrow&  \cos x(\sin 3x-\cos 3x)=0\\
			&\Leftrightarrow& \hoac{&\cos x=0\\&\sin 3x-\cos 3x=0}
		\end{eqnarray*}
		\begin{itemize}
			\item $\cos x=0\Leftrightarrow  x=\dfrac{\pi}{2}+k\pi$ ($k\in\mathbb{Z}$).
			\item $\sin 3x-\cos3x=0\Leftrightarrow \sin 3x=\cos 3x$.\\
			Nếu $\cos 3x=0$ thì phương trình đưa về dạng $\sin 3x=0$ (vô lí).\\
			Nếu $\cos 3x\ne 0$, phương trình đưa về dạng
			$$\tan 3x=1\Leftrightarrow 3x=\dfrac{\pi}{4}+k\pi\Leftrightarrow x=\dfrac{\pi}{12}+k\dfrac{\pi}{3}\, (k\in \mathbb{Z}).$$
			Cách khác: $\sin 3x-\cos3x=0\Leftrightarrow \sin 3x=\cos 3x=\sin\left(\dfrac{\pi}{2}-3x\right)\Leftrightarrow x=\dfrac{\pi}{12}+k\dfrac{\pi}{3}$.
		\end{itemize}
		Vậy nghiệm của phương trình đã cho là $x=\dfrac{\pi}{2}+k\pi$ và $x=\dfrac{\pi}{12}+k\dfrac{\pi}{3}$ ($k\in \mathbb{Z}$).
		\begin{itemchoice}
			\itemch Phát biểu đã cho sai.
			\itemch Phát biểu đã cho sai.
			\itemch Phát biểu đã cho đúng.
			\itemch Phát biểu đã cho đúng.
		\end{itemchoice}
	}
\end{ex}

% #2052 | Archive/2024/SP-K10-K11-LyVanHoang_TVN-023/11. C1-LuongGiac-lvhoang_TVN-023/dataC1K11/Bai-3-Dang-6-Dothihamsoluonggiac.tex | mã gốc 1D1H4-7
\begin{vd}%[1D1H4-7]%[BG K11, Nhật Thiện]
	[Cánh diều - Mức độ 2]
	Dùng đồ thị hàm số, tìm giá trị của $x$ trên đoạn $[-2 \pi ; 2 \pi]$ để
	\begin{enumEX}{2}
		\item Hàm số $y=\sin x$ nhận giá trị bằng $1$;	
		\item Hàm số $y=\cos x$ nhận giá trị bằng $-1$;	
	\end{enumEX}
	\loigiai{
		\begin{enumerate}
			\item Đồ thị hàm số $y=\sin x$ trên đoạn $[-2 \pi ; 2 \pi]$
			\begin{center}
				\begin{tikzpicture}[>=stealth,line join=round,line cap=round,font=\footnotesize,scale=0.8]
					\draw[->] (-7,0) -- (7,0) node[below] {\small $x$};
					\draw[->] (0,-1.4) -- (0,1.6) node[right] {\small $y$};
					\draw[samples=100,domain=-2*pi:2*pi] plot(\x,{sin((\x)*180/pi)});
					\draw (0,0) node [above left] {$O$};
					\draw (0,1) node [above left] {$1$};
					\draw (0,-1) node [above right] {$-1$};
					\draw (-6.28,0) node[below] {$-2\pi$};
					\draw (6.28,0) node[above] {$2\pi$};
					\draw[dashed] ({-3*pi/2},1)--({-3*pi/2},0) node[below] {$-\frac{3\pi}{2}$};
					\draw[dashed] (4.71,-1)--(4.71,0) node[above] {$\frac{3\pi}{2}$};
					\draw (-3.17,0) node[below left] {$-\pi$};
					\draw (3.14,0) node[below left] {$\pi$};
					\draw[dashed] (-pi/2,-1)--(-pi/2,0) node[above] {$-\frac{\pi}{2}$};
					\draw[dashed] (pi/2,1)--(pi/2,-0.2) node[below] {$\frac{\pi}{2}$};
					\draw[dashed] (-3*pi/2,1)--(pi/2,1);
					\draw[dashed] (-pi/2,-1)--(3*pi/2,-1);
				\end{tikzpicture}
			\end{center}
			Dựa vào đồ thị ta thấy trên đoạn $[-2 \pi ; 2 \pi]$ hàm số $y=\sin x$ nhận giá trị bằng $1$ khi $x\in\left\{ -\dfrac{3\pi}{2};\dfrac{\pi}{2}\right\} $.
			
			\item Đồ thị hàm số $y=\cos x$ trên đoạn $[-2 \pi ; 2 \pi]$
			\begin{center}
				\begin{tikzpicture}[>=stealth,line join=round,line cap=round,font=\footnotesize,scale=0.8]
					\draw[->] (-8.1,0) -- (8.5,0) node[below] {\small $x$};
					\draw[->] (0,-1.4) -- (0,1.6) node[right] {\small $y$};
					\draw[samples=100,domain=-2*pi:2*pi] plot(\x,{cos((\x)*180/pi)});
					\draw (0,0) node [above left] {$O$};
					\draw (0,1) node [above left] {$1$};
					\draw (0,-1) node [above right] {$-1$};
					\draw[dashed] (-pi,-1)--(-pi,0) node[above] {$-\pi$};
					\draw[dashed] (pi,-1)--(pi,0) node[above] {$\pi$};
					\draw[dashed] (-2*pi,1)--(-2*pi,0) node[below] {$-2\pi$};
					\draw[dashed] (2*pi,1)--(2*pi,0) node[below] {$2\pi$};
					\draw[dashed] (-1.57,0) node[below] {$-\frac{\pi}{2}$};
					\draw[dashed] (1.57,0) node[below] {$\frac{\pi}{2}$};
					\draw[dashed] (-2*pi,1)--(2*pi,1);
					\draw[dashed] (-pi,-1)--(pi,-1);
				\end{tikzpicture}
			\end{center}
			Dựa vào đồ thị ta thấy trên đoạn $[-2 \pi ; 2 \pi]$ hàm số $y=\cos x$ nhận giá trị bằng $-1$ khi $x=\pm\pi$. 
		\end{enumerate}
	}
\end{vd}

% #2053 | Archive/2024/SP-K10-K11-LyVanHoang_TVN-023/11. C1-LuongGiac-lvhoang_TVN-023/dataC1K11/Bai-3-Dang-6-Dothihamsoluonggiac.tex | mã gốc 1D1H4-7
\begin{vd}%[1D1H4-7]%[BG K11, Nhật Thiện]
	[Kết nối tri thức với cuộc sống - Mức độ 2]
	Sử dụng đồ thị ở hình dưới, hãy xác định các giá trị của $x$ trên đoạn $\left[-\dfrac{\pi}{2} ; 2 \pi\right]$ để hàm số $y=\cot x$ :
	\begin{center}
		\begin{tikzpicture}[>=stealth,scale=0.77,transform shape] 
			\path
			({-2*pi},0) coordinate (X1)
			({-1.5*pi},0) coordinate (X2)
			({-pi},0) coordinate (X3)
			({-0.5*pi},0) coordinate (X4)
			({0.5*pi},0) coordinate (X5)
			({pi},0) coordinate (X6)
			({1.5*pi},0) coordinate (X7)
			({2*pi},0) coordinate (X8)
			;
			\draw[->] (-7.5,0) -- (7.5,0) node[below] {\small $x$};
			\draw[->] (0,-3.5) -- (0,3.5) node[left] {\small $y$};
			\draw [dashed] ({-2*pi},3.5)--({-2*pi},-3.5) ({-pi},3.5)--({-pi},-3.5) ({pi},3.5)--({pi},-3.5) ({2*pi},3.5)--({2*pi},-3.5) ;
			\foreach \x/\g/\z in {X1/-150/-2\pi,X2/-130/-\tfrac{3\pi}{2},X3/-140/-\pi,X4/-100/-\tfrac{\pi}{2},X5/-100/\tfrac{\pi}{2},X6/-120/\pi,X7/-100/\tfrac{3\pi}{2}, X8/-120/2\pi}
			\fill[black] (\x) circle(1pt) +(\g:5mm) node {$\z$};
			\draw (0,0) node[below left]{$O$};
			\clip (-6.5,-3.5) rectangle (6.5,3.5) ;
			\draw[samples=100,domain={0.2}:{pi-0.2}] plot(\x,{cot((\x)*180/pi)});
			\draw[samples=100,domain={pi+0.2}:{2*pi-0.2}] plot(\x,{cot((\x)*180/pi)});
			\draw[samples=100,domain={-0.2}:{-pi+0.2}] plot(\x,{cot((\x)*180/pi)});
			\draw[samples=100,domain={-pi-0.2}:{-2*pi+0.2}] plot(\x,{cot((\x)*180/pi)});			
			\foreach \x in {-1.5*pi,-0.5*pi,0.5*pi,1.5*pi}
			{\draw[fill=black] (\x,0) circle (1pt);}
		\end{tikzpicture}
	\end{center}
	\begin{enumEX}{2}
		\item Nhận giá trị bằng $0$;
		\item Nhận giá trị âm.
	\end{enumEX}
	\loigiai{
		\begin{enumerate}
			\item Từ đồ thị ta suy ra trên đoạn $\left[-\dfrac{\pi}{2} ; 2 \pi\right]$ thì $y=0$ khi $x=-\dfrac{\pi}{2} ; x=\dfrac{\pi}{2} ; x=\dfrac{3 \pi}{2}$. 
			\item Hàm số nhận giá trị âm ứng với phần đồ thị nằm dưới trục hoành. \\
			Từ đồ thị ta suy ra trên đoạn $\left[-\dfrac{\pi}{2} ; 2 \pi\right]$ thì $y<0$ khi $x \in\left(-\dfrac{\pi}{2} ; 0\right) \cup\left(\dfrac{\pi}{2} ; \pi\right) \cup\left(\dfrac{3 \pi}{2} ; 2 \pi\right)$. 
		\end{enumerate}
	}
\end{vd}

% #2054 | Archive/2024/SP-K10-K11-LyVanHoang_TVN-023/11. C1-LuongGiac-lvhoang_TVN-023/dataC1K11/Bai-3-Dang-6-Dothihamsoluonggiac.tex | mã gốc 1D1H4-7
\begin{vd}%[1D1H4-7]%[BG K11, Nhật Thiện]
	[Chân trời sáng tạo - Mức độ 2]
	Dùng đồ thị hàm số, tìm giá trị của $x$ trên khoảng $\left(-\pi ; \dfrac{3 \pi}{2}\right)$ để
	\begin{enumEX}{2}
		\item Hàm số $y=\tan x$ nhận giá trị bằng $-1$;
		\item Hàm số $y=\tan x$ nhận giá trị bằng $0$;
		\item Hàm số $y=\cot x$ nhận giá trị bằng $1$;
		\item Hàm số $y=\cot x$ nhận giá trị bằng $0$.
	\end{enumEX}
	\loigiai{
		\begin{enumerate}
			\item Đồ thị hàm số $y=\tan x$ trên khoảng $\left(-\pi ; \dfrac{3 \pi}{2}\right)$
			\begin{center}
				\begin{tikzpicture}[>=stealth,line join=round,line cap=round,font=\footnotesize,scale=0.8]
					\draw[->] (-4,0) -- (5.5,0) node[below] {\small $x$};
					\draw[->] (0,-4) -- (0,5) node[right] {\small $y$};
					\draw (0,0) node[below right] {$O$};
					\draw (0,1) node[above right] {$1$};
					\draw (0,-1) node[below right] {$-1$};
					\draw[thick,samples=100,domain=-1.325:1.36] plot(\x,{tan((\x)*180/pi)});
					\draw (-3.17,0) node[below] {$-\pi$};
					\draw (3.14,0) node[below] {$\pi$};
					\draw (-1.57,0) node[below left=-0.1] {$-\frac{\pi}{2}$};
					\draw (1.57,-0.2) node[below left=-0.1] {$\frac{\pi}{2}$};
					\draw (4.71,0) node[below left=-0.1] {$\frac{3\pi}{2}$};
					\draw (2.36,0) node[above] {$\frac{3\pi}{4}$};
					\draw (-0.79,0) node[above] {$-\frac{\pi}{4}$};
					\draw[dashed] (1.57,-4)--(1.57,5);
					\draw[dashed] (-1.57,-4)--(-1.57,5);
					\draw[dashed] (4.71,-4)--(4.71,5);
					\draw[dashed] (2.36,0)--(2.36,-1);
					\draw[dashed] (-0.79,0)--(-0.79,-1);
					\draw[thick,samples=100,domain=1.82:4.5] plot(\x,{tan((\x)*180/pi)});
					\draw[thick,samples=100,domain=-3.15:-1.79] plot(\x,{tan((\x)*180/pi)});
					\draw[dashed] (-4,-1)--(5,-1);	
					\foreach \x in {-pi,pi}
					{\draw[fill=black] (\x,0) circle (1pt);}		
				\end{tikzpicture}
			\end{center}
			Dựa vào đồ thị ta thấy trên khoảng $\left(-\pi ; \dfrac{3 \pi}{2}\right)$ hàm số $y=\tan x$ nhận giá trị bằng $-1$ khi $x\in\left\lbrace -\dfrac{\pi}{4};\dfrac{3\pi}{4}\right\rbrace $. 
			\item Đồ thị hàm số $y=\tan x$ trên khoảng $\left(-\pi ; \dfrac{3 \pi}{2}\right)$
			\begin{center}
				\begin{tikzpicture}[>=stealth,line join=round,line cap=round,font=\footnotesize,scale=0.8]
					\draw[->] (-4,0) -- (5.5,0) node[below] {\small $x$};
					\draw[->] (0,-4) -- (0,5) node[right] {\small $y$};
					\draw (0,0) node[below right] {$O$};
					\draw (0,1) node[above right] {$1$};
					\draw (0,-1) node[below right] {$-1$};
					\draw[thick,samples=100,domain=-1.325:1.36] plot(\x,{tan((\x)*180/pi)});
					\draw (-3.17,0) node[below] {$-\pi$};
					\draw (3.14,0) node[below] {$\pi$};
					\draw (-1.57,0) node[below right=-0.1] {$-\frac{\pi}{2}$};
					\draw (1.57,-0.2) node[below right=-0.1] {$\frac{\pi}{2}$};
					\draw (4.71,0) node[below right=-0.1] {$\frac{3\pi}{2}$};
					\draw[dashed] (1.57,-4)--(1.57,5);
					\draw[dashed] (-1.57,-4)--(-1.57,5);
					\draw[dashed] (4.71,-4)--(4.71,5);
					\draw[thick,samples=100,domain=1.82:4.5] plot(\x,{tan((\x)*180/pi)});
					\draw[thick,samples=100,domain=-3.15:-1.79] plot(\x,{tan((\x)*180/pi)});			
				\end{tikzpicture}
			\end{center}
			Dựa vào đồ thị ta thấy trên khoảng $\left(-\pi ; \dfrac{3 \pi}{2}\right)$ hàm số $y=\tan x$ nhận giá trị bằng $0$ khi $x\in\left\lbrace 0;\pi\right\rbrace $.  
			\item Đồ thị hàm số $y=\cot x$ trên khoảng $\left(-\pi ; \dfrac{3 \pi}{2}\right)$.
			\begin{center}
				\begin{tikzpicture}[>=stealth,line join=round,line cap=round,font=\footnotesize,scale=0.8]
					\draw[->] (-4.5,0)--(5.5,0) node[below] {\small $x$};
					\draw[->] (0,-4)--(0,5) node[left] {\small $y$};
					\draw (0,0) node[below right] {$O$};
					\draw (0,1) node[above right] {$1$};
					\draw (0,-1) node[below right] {$-1$};
					\draw[thick,samples=100,domain=0.21:2.9] plot(\x,{cot((\x)*180/pi)});
					\draw[thick,samples=100,domain=0.21:2.9] plot(\x-pi,{cot((\x)*180/pi)});
					\draw[thick,samples=100,domain=0.21:1.57] plot(\x+pi,{cot((\x)*180/pi)});
					\draw (-3.17,0) node[below left=-0.1] {$-\pi$};
					\draw (3.14,0) node[below left=-0.1] {$\pi$};
					\draw (-1.57,0) node[above=-0.1] {$-\frac{\pi}{2}$};
					\draw (1.57,-0.2) node[below=-0.2] {$\frac{\pi}{2}$};
					\draw (4.71,0) node[below] {$\frac{3\pi}{2}$};
					\draw (-2.36,0) node[below] {$-\frac{3\pi}{4}$};
					\draw (0.79,0) node[below] {$\frac{\pi}{4}$};
					\draw (3.93,0) node[below] {$\frac{5\pi}{4}$};
					\draw[dashed] (3.14,-4)--(3.14,5);
					\draw[dashed] (-3.14,-4)--(-3.1,5);
					\draw[dashed] (-4,1)--(5,1);
					\draw[dashed] (0.79,0)--(0.79,1);
					\draw[dashed] (-2.36,0)--(-2.36,1);
					\draw[dashed] (3.93,0)--(3.93,1);
					\foreach \x in {1.5*pi,-0.5*pi,0.5*pi,1.5*pi}
					{\draw[fill=black] (\x,0) circle (1pt);}
				\end{tikzpicture}
			\end{center}
			Dựa vào đồ thị ta thấy trên khoảng $\left(-\pi ; \dfrac{3 \pi}{2}\right)$ hàm số $y=\cot x$ nhận giá trị bằng $1$ khi\\ $x\in\left\{ -\dfrac{3\pi}{4};\dfrac{\pi}{4};\dfrac{5\pi}{4}\right\} $.
			\item Đồ thị hàm số $y=\cot x$ trên khoảng $\left(-\pi ; \dfrac{3 \pi}{2}\right)$.
			\begin{center}
				\begin{tikzpicture}[>=stealth,line join=round,line cap=round,font=\footnotesize,scale=0.8]
					\draw[->] (-4.5,0)--(5.5,0) node[below] {\small $x$};
					\draw[->] (0,-4)--(0,5) node[left] {\small $y$};
					\draw (0,0) node[below right] {$O$};
					\draw (0,1) node[above right] {$1$};
					\draw (0,-1) node[below right] {$-1$};
					\draw[thick,samples=100,domain=0.21:2.9] plot(\x,{cot((\x)*180/pi)});
					\draw[thick,samples=100,domain=0.21:2.9] plot(\x-pi,{cot((\x)*180/pi)});
					\draw[thick,samples=100,domain=0.21:1.57] plot(\x+pi,{cot((\x)*180/pi)});
					\draw (-3.17,0) node[below left=-0.1] {$-\pi$};
					\draw (3.14,0) node[below left=-0.1] {$\pi$};
					\draw (-1.57,0) node[below left=-0.1] {$-\frac{\pi}{2}$};
					\draw (1.57,-0.2) node[below=-0.2] {$\frac{\pi}{2}$};
					\draw (4.71,0) node[below] {$\frac{3\pi}{2}$};
					\draw[dashed] (3.14,-4)--(3.14,5);
					\draw[dashed] (-3.14,-4)--(-3.1,5);
				\end{tikzpicture}
			\end{center}
			Dựa vào đồ thị ta thấy trên khoảng $\left(-\pi; \dfrac{3 \pi}{2}\right)$ hàm số $y=\cot x$ nhận giá trị bằng $1$ khi $x=\pm\dfrac{\pi}{2}$.
		\end{enumerate}
	}
\end{vd}

% #2056 | Archive/2024/SP-K10-K11-LyVanHoang_TVN-023/11. C1-LuongGiac-lvhoang_TVN-023/dataC1K11/Bai-3-Dang-6-Dothihamsoluonggiac.tex | mã gốc 1D1V4-7
\begin{vd}%[1D1V4-7]%[BG K11, Nhật Thiện]
	[Chân trời sáng tạo - Mức độ 3]
	Nhiệt độ ngoài trời $T$ (tính bằng ${ }^{\circ} \mathrm{C}$) vào thời điểm $t$ giờ trong một ngày ở một thành phố được tính bởi công thức $T=20+4 \sin \left(\dfrac{\pi}{12} t-\dfrac{5 \pi}{6}\right)$. Để bảo quản các tác phẩm nghệ thuật, hệ thống điều hoà nhiệt độ của một bào tàng sẽ được tự động bật khi nhiệt độ ngoài trời từ $22^{\circ} \mathrm{C}$ trở lên. Dựa vào đồ thị của hàm số $\sin$, hãy xác định khoảng thời gian $t$ trong ngày $(0 \leq t \leq 24)$ hệ thống điều hoà được bật.
	(Theo https://www.sciencedirect.com/science/article/abs/pii/0168192385900139).
	\loigiai
	{
		\immini{Ta có $T \geq 22$ khi và chỉ khi $20+4 \sin \left(\dfrac{\pi}{12} t-\dfrac{5 \pi}{6}\right) \geq 22$ hay $\sin \left(\dfrac{\pi}{12} t-\dfrac{5 \pi}{6}\right) \geq \dfrac{1}{2}$.\\ Vì $0 \leq t \leq 24$ nên $-\dfrac{5 \pi}{6} \leq \dfrac{\pi}{12} t-\dfrac{5 \pi}{6} \leq \dfrac{7 \pi}{6}$.\\
			Xét đồ thị hàm số $y=\sin x$ trên đoạn $\left[-\dfrac{5 \pi}{6} ; \dfrac{7 \pi}{6}\right]$.}
		{
			\begin{tikzpicture}[scale=1,>=stealth, font=\footnotesize, line join=round, line cap=round]
				\def\xmin{-2.8} \def\xmax{4.25} \def\ymin{-1.5} \def\ymax{1.8}
				\draw[->] (\xmin,0)--(\xmax,0) node [below]{$x$};
				\draw[->] (0,\ymin)--(0,\ymax) node [right]{$y$};
				\node at (0,0) [below left]{$O$};
				\clip (\xmin+0.1,\ymin+0.1) rectangle (\xmax-0.5,\ymax-0.1);
				\draw[smooth,samples=400,domain=\xmin:\xmax] plot(\x,{sin(\x r)});
				\draw[dashed] (\xmin,1)--(\xmax,1) (\xmin,-1)--(\xmax,-1);
				\foreach \x in {-2*pi,-1.5*pi,-pi,-0.5*pi,0}
				{\draw[fill=black] (\x,sin \x*180/pi) circle (1pt);
					\draw[dashed] (\x,sin \x*180/pi)--(\x,0);
					\draw[fill=black] (-\x,sin -\x*180/pi) circle (1pt);
					\draw[dashed] (-\x,sin -\x*180/pi)--(-\x,0);}
				\node at (0,1.3) [left]{$1$};
				\node at (0,-1.3) [left]{$-1$};
				\node at (0,0.6) [left]{$\frac{1}{2}$};
				\node at (-2*pi+0.15,0) [below]{$-2\pi$};
				\node at (-0.8*pi,0) [above]{$-\frac{5\pi}{6}$};
				\node at (0.16*pi,0) [below]{$\frac{\pi}{6}$};
				\node at (1.13*pi,0) [above]{$\frac{7\pi}{6}$};
				\node at (0.83*pi,0) [below]{$\frac{5\pi}{6}$};
				\draw[dashed] (0.53,0)--(0.53,0.5)--(0,0.5);
				\draw[dashed] (2.61,0)--(2.61,0.5)--(0,0.5);
				\draw[dashed] (3.66,0)--(3.66,-0.5);
				\draw[dashed] (-2.61,0)--(-2.61,-0.5);
		\end{tikzpicture}}
		Ta thấy $\sin \left(\dfrac{\pi}{12} t-\dfrac{5 \pi}{6}\right) \geq \dfrac{1}{2}$ khi và chỉ khi $\dfrac{\pi}{6} \leq \dfrac{\pi}{12} t-\dfrac{5 \pi}{6} \leq \dfrac{5 \pi}{6}$ hay $12 \leq t \leq 20$.\\
		Vậy hệ thống điều hoà được bật trong khoảng thời gian từ $12$ giờ đến $20$ giờ trong ngày.}
\end{vd}

% #2059 | Archive/2024/SP-K10-K11-LyVanHoang_TVN-023/11. C1-LuongGiac-lvhoang_TVN-023/dataC1K11/Bai-3-Dang-6-Dothihamsoluonggiac.tex | mã gốc 1D1H4-7
\begin{bt}%[1D1H4-7]%[BG K11, Nhật Thiện]
	[Mức độ 2]
	Cho hàm số $y=\sin x$ có đồ thị như hình vẽ dưới đây
	\begin{center}
	\begin{tikzpicture}[scale=0.8,>=stealth, font=\footnotesize, line join=round, line cap=round]
		\def\xmin{-10} \def\xmax{10.5} \def\ymin{-1.5} \def\ymax{1.8}
		\draw[->] (\xmin,0)--(\xmax,0) node [below]{$x$};
		\draw[->] (0,\ymin)--(0,\ymax) node [right]{$y$};
		\node at (0,0) [below right]{$O$};
		\clip (\xmin+0.1,\ymin+0.1) rectangle (\xmax-0.5,\ymax-0.1);
		\draw[smooth,samples=400,domain=\xmin:\xmax] plot(\x,{sin(\x r)});
		\draw[dashed] (\xmin,1)--(\xmax,1) (\xmin,-1)--(\xmax,-1);
		\foreach \x in {-3*pi,-2.5*pi,-2*pi,-1.5*pi,-pi,-0.5*pi,0}
		{\draw[fill=black] (\x,sin \x*180/pi) circle (1pt);
			\draw[dashed] (\x,sin \x*180/pi)--(\x,0);
			\draw[fill=black] (-\x,sin -\x*180/pi) circle (1pt);
			\draw[dashed] (-\x,sin -\x*180/pi)--(-\x,0);}
		\node at (0,1.3) [left]{$1$};
		\node at (0,-1.3) [left]{$-1$};
		\node at (-2*pi+0.15,0) [below]{$-2\pi$};
		\node at (-3*pi-0.15,0) [below]{$-3\pi$};
		\node at (-2.5*pi,0) [above]{$-\frac{5\pi}{2}$};
		\node at (-1.5*pi,0) [below]{$-\frac{3\pi}{2}$};
		\node at (-pi-0.15,0) [below]{$-\pi$};
		\node at (-0.5*pi,0) [above]{$-\frac{\pi}{2}$};
		\node at (0.5*pi,0) [below]{$\frac{\pi}{2}$};
		\node at (pi-0.1,0) [below]{$\pi$};
		\node at (1.5*pi,0) [above]{$\frac{3\pi}{2}$};
		\node at (2*pi+0.2,0) [below]{$2\pi$};
		\node at (2.5*pi+0.2,0) [below]{$\frac{5\pi}{2}$};
		\node at (3*pi+0.2,0) [below left]{$3\pi$};
		\node at (2.5*pi+1,0.8) [above]{$y=\sin x$};
	\end{tikzpicture}
	\end{center}
	Tìm những giá trị $x$ thuộc  $\left[-\dfrac{3\pi}{2};\dfrac{5\pi}{2}\right]$ để hàm số $y=\sin x$ nhận giá trị âm.
	\loigiai{
		Trên các khoảng $(\pi;0)$; $(\pi;2\pi)$ đồ thị hàm số nằm phía dưới trục hoành nên hàm số $y=\sin x$ nhận giá trị	âm.
	}
\end{bt}

% #2060 | Archive/2024/SP-K10-K11-LyVanHoang_TVN-023/11. C1-LuongGiac-lvhoang_TVN-023/dataC1K11/Bai-3-Dang-6-Dothihamsoluonggiac.tex | mã gốc 1D1H4-7
\begin{bt}%[1D1H4-7]%[BG K11, Nhật Thiện]
	[Cánh diều - Mức độ 2]
	Dùng đồ thị hàm số, tìm giá trị của $x$ trên đoạn $[-2 \pi ; 2 \pi]$ để
	\begin{enumEX}{2}
		\item Hàm số $y=\sin x$ nhận giá trị bằng $0$;
		\item Hàm số $y=\cos x$ nhận giá trị bằng $0$.
	\end{enumEX}
	\loigiai{
		\begin{enumerate}
			\item Đồ thị hàm số $y=\sin x$ trên đoạn $[-2 \pi ; 2 \pi]$
			\begin{center}
			\begin{tikzpicture}[>=stealth,line join=round,line cap=round,font=\footnotesize,scale=0.8]
				\draw[->] (-7,0) -- (7,0) node[below] {\small $x$};
				\draw[->] (0,-1.4) -- (0,1.6) node[right] {\small $y$};
				\draw[samples=100,domain=-2*pi:2*pi] plot(\x,{sin((\x)*180/pi)});
				\draw (0,0) node [above left] {$O$};
				\draw (0,1) node [above left] {$1$};
				\draw (0,-1) node [above right] {$-1$};
				\draw (-6.28,0) node[below] {$-2\pi$};
				\draw (6.28,0) node[above] {$2\pi$};
				\draw[dashed] ({-3*pi/2},1)--({-3*pi/2},0) node[below] {$-\frac{3\pi}{2}$};
				\draw[dashed] (4.71,-1)--(4.71,0) node[above] {$\frac{3\pi}{2}$};
				\draw (-3.17,0) node[below left] {$-\pi$};
				\draw (3.14,0) node[below left] {$\pi$};
				\draw[dashed] (-pi/2,-1)--(-pi/2,0) node[above] {$-\frac{\pi}{2}$};
				\draw[dashed] (pi/2,1)--(pi/2,-0.2) node[below] {$\frac{\pi}{2}$};
				\draw[dashed] (-3*pi/2,1)--(pi/2,1);
				\draw[dashed] (-pi/2,-1)--(3*pi/2,-1);
			\end{tikzpicture}
			\end{center}
			Dựa vào đồ thị ta thấy trên đoạn $[-2 \pi ; 2 \pi]$ hàm số $y=\sin x$ nhận giá trị bằng $0$ khi\\ $x\in\left\{ \pm2\pi;\pm\pi;0\right\} $.
			\item Đồ thị hàm số $y=\cos x$ trên đoạn $[-2 \pi ; 2 \pi]$
			\begin{center}
			\begin{tikzpicture}[>=stealth,line join=round,line cap=round,font=\footnotesize,scale=0.8]
				\draw[->] (-8.1,0) -- (8.5,0) node[below] {\small $x$};
				\draw[->] (0,-1.4) -- (0,1.6) node[right] {\small $y$};
				\draw[samples=100,domain=-2*pi:2*pi] plot(\x,{cos((\x)*180/pi)});
				\draw (0,0) node [above left] {$O$} (3*pi/2,0) node[shift={(-70:.35)}]{$\frac{3\pi}{2}$};
				\draw (0,1) node [above left] {$1$};
				\draw (0,-1) node [above right] {$-1$};
				\draw[dashed] (-pi,-1)--(-pi,0) node[above] {$-\pi$};
				\draw[dashed] (pi,-1)--(pi,0) node[above] {$\pi$};
				\draw[dashed] (-2*pi,1)--(-2*pi,0) node[below] {$-2\pi$};
				\draw[dashed] (2*pi,1)--(2*pi,0) node[below] {$2\pi$};
				\draw[dashed] (-1.57,0) node[below] {$-\frac{\pi}{2}$};
				\draw[dashed] (1.57,0) node[below] {$\frac{\pi}{2}$};
				\draw[dashed] (-2*pi,1)--(2*pi,1);
				\draw[dashed] (-pi,-1)--(pi,-1);
			\end{tikzpicture}
			\end{center}
			Dựa vào đồ thị ta thấy trên đoạn $[-2 \pi ; 2 \pi]$ hàm số $y=\cos x$ nhận giá trị bằng $0$ khi\\ $x\in\left\{ \pm\dfrac{3\pi}{2};\pm\dfrac{\pi}{2}\right\} $. 
		\end{enumerate}
	}
\end{bt}

% #2061 | Archive/2024/SP-K10-K11-LyVanHoang_TVN-023/11. C1-LuongGiac-lvhoang_TVN-023/dataC1K11/Bai-3-Dang-6-Dothihamsoluonggiac.tex | mã gốc 1D1H4-7
\begin{bt}%[1D1H4-7]%[BG K11, Nhật Thiện]
	[Mức độ 2]
	\immini{Cho đồ thị hàm số $y= 2\sin x+ 1$ như hình vẽ, dựa vào đồ thị hãy tìm
		\begin{enumerate}
			\item Giá trị lớn nhất và giá trị nhỏ nhất của hàm số trên đoạn $\left[- \pi; \pi\right]$.
			\item Các khoảng đồng biến và nghịch biến của hàm số trên đoạn $\left[- \pi; \pi\right]$.
	\end{enumerate}}
	{
		\begin{tikzpicture}[line join=round, line cap=round,>=stealth]
			%Giới hạn đồ thị và hàm số cần vẽ
			\def \xmin{-pi-1} 		\def \xmax{pi+1}
			\def \ymin{-2}		\def \ymax{4}
			\def \hamso{2*sin(\x r)+1}
			% Vẽ hệ trục
			\draw[->] (\xmin,0)--(\xmax,0) node[above right] {$x$}; 
			\draw[->] (0,\ymin)--(0,\ymax) node[below left] {$y$}; 
			\draw (0,0) node [below right] {$O$}; 
			% Vẽ các điểm trên 2 hệ trục  
			\foreach \x in {-pi,-pi/2,pi/2,pi} \draw[thin] (\x,1pt)--(\x,-1pt);
			\foreach \y in {-1,1,3} \draw[thin] (1pt,\y)--(-1pt,\y);
			\node at (-pi,0) [below]{$-\pi$};
			\node at (-pi/2,0)[above]{$-\dfrac{\pi}{2}$};
			\node at (pi,0) [below]{$\pi$};
			\node at (pi/2,0) [below]{$\dfrac{\pi}{2}$};
			\node at (0,1) [above left]{$1$};
			\node at (0,-1) [right]{$-1$};
			\node at (0,3)[left]{$3$};
			% Vẽ thêm mấy cái râu ria 
			\draw[dashed,thin](-pi/2,0)--(-pi/2,-1)--(0,-1); 
			\draw[dashed,thin](pi/2,0)--(pi/2,3)--(0,3); 
			\draw[dashed,thin] (-pi,0)--(-pi,1)--(0,1);
			\draw[dashed,thin] (pi,0)--(pi,1)--(0,1);
			% Vẽ hàm số
			\begin{scope}
				\clip (\xmin+0.1,\ymin+0.1) rectangle (\xmax-0.1,\ymax-0.1); 
				\draw[samples=350,domain=\xmin+0.1:\xmax-0.1,smooth,variable=\x] plot (\x,{\hamso}); 
			\end{scope} 
		\end{tikzpicture}
	}
	\loigiai{Dựa vào đồ thị ta có
		\begin{enumerate}
			\item
			Giá trị lớn nhất của hàm số đã cho trên đoạn $[-\pi; \pi]$ là $3$ khi $x= \dfrac{\pi}{2}$.\\
			Giá trị nhỏ nhất của hàm số đã cho trên đoạn $[-\pi; \pi]$ là $-1$ khi $x= -\dfrac{\pi}{2}$.
			\item  Trên đoạn $[-\pi;\pi]$, hàm số đã cho đồng biến trên khoảng $\left(\dfrac{-\pi}{2}; \dfrac{\pi}{2}\right)$.\\
			Trên đoạn $[-\pi;\pi]$, hàm số đã cho nghịch biến trên các khoảng $\left(-\pi; -\dfrac{\pi}{2}\right)$ và $\left(\dfrac{\pi}{2}; \pi\right).$
		\end{enumerate}
	}
\end{bt}

% #2063 | Archive/2024/SP-K10-K11-LyVanHoang_TVN-023/11. C1-LuongGiac-lvhoang_TVN-023/dataC1K11/Bai-3-Dang-6-Dothihamsoluonggiac.tex | mã gốc 1D1V4-7
\begin{bt}%[1D1V4-7]%[BG K11, Nhật Thiện]
	[Mức độ 3]
	Trong hình vẽ bên dưới, một chiếc máy bay $A$ bay ở độ cao $500$ m theo một đường thẳng đi ngang qua phía trên trạm quan sát $T$ ở mặt đất. Hình chiếu vuông góc của $A$ lên mặt đất là $H$, $\alpha$ là góc lượng giác $\left(Tx,TA\right)$ $\left(0<\alpha<\pi\right)$.
	\begin{center}
		\begin{tikzpicture}[>=stealth,x=1.0cm,y=1.0cm, scale=1]
			\draw[->](0,0) -- (14,0) node[below]{\footnotesize $x$ (m)};
			\draw[dashed,red](0,3) -- (14,3);
			\path
			(5,0) coordinate (T)
			(10,0) coordinate (H)
			(10,3) coordinate (A)
			;
			\draw[fill = black] (T) node [below]{\footnotesize$T$ (trạm quan sát)} circle (1pt) -- (H) node [below] {\footnotesize$H$} circle (1pt) -- (A) node [above]{\footnotesize$A$ (máy bay)} circle (1pt) -- (T);
			\draw (10,1.5) node [right]{\footnotesize $500$m};
			\draw pic[->,draw, angle radius=5mm]{angle=H--T--A} (T) node[shift={(10:.8)}]{$\alpha$};
			\path pic[draw,angle radius=6]{right angle=A--H--T};
		\end{tikzpicture}
	\end{center}
	\begin{enumerate}
		\item Biểu diễn tọa độ $x_H$ của điểm $H$ trên trục $Tx$ theo $\alpha$.
		\item Dựa vào đồ thị hàm số cô-tang, hãy cho biết với $\dfrac{\pi}{6} < \alpha < \dfrac{2\pi}{3}$ thì $x_H$ nằm trong khoảng nào? Làm tròn kết quả đến hàng phần mười.
	\end{enumerate}
	\loigiai{
		\begin{enumerate}
			\item Coi trạm quan sát $T$ là gốc tọa độ thì ta có
			\[x_H = \overline{TH} = AH\cdot\cot\alpha = 500\cot\alpha\;\left(\text{m}\right).\]
			\item \immini{
				Dựa vào đồ thị hàm số $\cot x$, ta thấy với $\dfrac{\pi}{6} < \alpha < \dfrac{2\pi}{3}$ thì $-\dfrac{\sqrt{3}}{3} < \cot \alpha < \sqrt{3}$.\\
				Do đó $-\dfrac{500\sqrt{3}}{3} < x_H < 500\sqrt{3}$. Làm tròn kết quả đến hàng phần mười ta được
				\[-288{,}7 < x_H< 866{,}0.\]
			}{
				\begin{tikzpicture}[font=\footnotesize,line join = round, line cap = round,>=stealth,scale=1]
					\draw[->] (-1,0) -- (4,0.)node[below]{\footnotesize $x$};
					\draw[->] (0,-3.1) -- (0,4.25)node[left]{\footnotesize $y$};
					\draw (0,0) node[above left] {\footnotesize $O$};
					\draw[dashed,smooth,samples=100,domain=0.25:2.8] plot(\x,{cot(((\x))*180/pi)});
					\draw[line width=1pt,smooth,red,samples=100,domain=pi/6:2*pi/3] plot(\x,{cot(((\x))*180/pi)});
					\draw[dashed] (pi,-3.1) -- (pi,0) node[below right]{$\pi$} -- (pi,4.2);
					\draw[dashed] (pi/6,0) node[below]{$\tfrac{\pi}{6}$} circle (1pt) --(pi/6,1.73)--(0,1.73) node[left]{$\sqrt{3}$};
					\draw[dashed] (2*pi/3,0) node[above]{$\tfrac{2\pi}{3}$} circle (1pt) --(2*pi/3,-0.577)--(0,-0.577) node[left]{$-\tfrac{\sqrt{3}}{3}$};
				\end{tikzpicture}
			}
		\end{enumerate}
	}
\end{bt}

% #2069 | Archive/2024/SP-K10-K11-LyVanHoang_TVN-023/11. C1-LuongGiac-lvhoang_TVN-023/dataC1K11/Bai-3-Dang-6-Dothihamsoluonggiac.tex | mã gốc 1D1V4-7
\begin{ex}[Mức độ 3]%[1D1V4-7]%[BG K11, Nhật Thiện]
	Trong Vật lý, phương trình tổng quát của một vật giao động điều hòa cho bởi công thức $x(t)=A\cos (\omega t+\varphi)$, trong đó $t$ là thời điểm (tính bằng giây), $x(t)$ (cm) là li độ của vật tại thời điểm $t$, $A$ là biên độ dao động (độ lệch lớn nhất so với vị trị cân bằng khi vật dao động, $A>0$), $\omega t+\varphi$ là pha của dao động tại thời điểm $t$ và $\varphi\in[-\pi;\pi]$ là pha ban đầu của dao động. Khi đó, chu kỳ $T$ của dao động là $T=\dfrac{2\pi}{\omega}$. Dao động của một con lắc đồng hồ theo thời gian $t$ (giây) được cho bởi đồ thị sau:
	\begin{center}
		\begin{tikzpicture}[line cap=round,line join=round,>=stealth]
			\draw [->] (-2,0)--(10.5,0) node [above right] {$t$};
			\draw [->] (0,-3)--(0,3.5) node [right] {$x$};
			\draw [samples=100, smooth,domain=0:10] plot (\x,{2*cos((\x*pi )r)});
			\foreach \x in {-2,-1,1,2,3,4,5,6,7,8,9,10} \draw (\x,-.045)--(\x,.045) node [below] {$\x$};
			\foreach \x in {-2,-1,1,2,3} \draw (-.045,\x)--(0.045,\x) node [left] {$\x$};
			\path (0,0) node [below left] {$0$};
			\draw[dashed] (0,2) -- (10,2) (0,-2) -- (10,-2);
		\end{tikzpicture}
	\end{center}
	Tính li độ của con lắc tại thời điểm $t=1{,}2$ giây (làm tròn kết quả đến hàng phần mười).  
	\shortans[]{$-1{,}6$}
	\loigiai{Dựa vào đồ thị ta có\\
		Biên độ dao động $A=2$ cm.\\
		Chu kì dao động $T=2$ giây.\\
		Suy ra $\omega=\dfrac{2\pi}{T}=\dfrac{2\pi}{2}=\pi$.\\
		Tại thời điểm ban đầu vật đi từ biên độ dương về vị trí cân bằng nên phương trình li độ có dạng $x=2\cos\left(\pi t\right)$ cm với $\varphi_0=0$.\\
		Tại thời điểm $t=1{,}2$ suy ra $x=2\cos \left(1{,}2\cdot\pi\right)=-1{,}6$.
	}
\end{ex}

% #2070 | Archive/2024/SP-K10-K11-LyVanHoang_TVN-023/11. C1-LuongGiac-lvhoang_TVN-023/dataC1K11/Bai-3-Dang-6-Dothihamsoluonggiac.tex | mã gốc 1D1V4-7
\begin{ex}%[1D1V4-7]%[BG K11, Nhật Thiện]
	[Mức độ 3]
	\immini{
		Khi đu quay hoạt động, vận tốc theo phương ngang của một cabin $M$ phụ thuộc vào góc lượng giác $\alpha = \left(Ox,OM\right)$ theo hàm số $v_x = 0{,}3\sin\alpha$ (m/s) (Hình vẽ bên). Dựa vào đồ thị của hàm số sin, biết trong vòng quay đầu tiên $\left(0\le \alpha\le 2\pi\right)$, góc $\alpha$ ở trong các khoảng $\left(0;\dfrac{\pi}{a}\right)$ và $\left(\dfrac{b\pi}{2};2\pi\right)$ thì $v_x$ tăng. Biết $a$, $b\in \mathbb{Z}$, tính $a+b$.
	}{\begin{tikzpicture}[scale=0.5,font=\small,line join=round,line cap=round,>=stealth]
		\tikzset{declare function={g=0;}}
		\coordinate (O) at (0,0);
		\coordinate[label=above right:$A$] (A) at (0:4.4);
		\coordinate[label=above right:$M$] (M) at (30+g:4.4);
		\pgfmathsetmacro{\gx}{60-g};
		\coordinate (v) at ($(M)+(120+g:3.3)$);
		\coordinate (vx) at ($(M)+(180:3.3*cos \gx)$);
		\draw[thick] (O) circle(4.4) (O) circle(3.53);
		\foreach \t in {0,1,...,11}{
			\draw[very thick,fill] (O)--(g+30*\t:4.4)circle(0.06);
			\begin{scope}[shift={(g+30*\t:4.4)}]
				\draw[fill=teal] (-0.27,0)--(0.27,0)--(0.36,-0.18)--(-0.36,-0.18)--cycle
				(0.32,-0.18)--(0.32,-0.54) (-0.32,-0.18)--(-0.32,-0.54)
				(-0.4,-0.54)--(0.4,-0.54)--(0.27,-1)--(-0.27,-1)--cycle;
		\end{scope}}
		\draw[fill=gray!70!yellow] (0.2,0)--(2.35,-5.53)--(3.65,-5.53)--(3.65,-6.02)--(1.95,-6.02)--(0,-0.85)--(-1.95,-6.02)--(-3.65,-6.02)--(-3.65,-5.53)--(-2.35,-5.53)--(-0.2,0)--cycle;
		\draw[fill=gray!10!yellow] (O) circle(0.34);
		\draw[fill] (-5,-5.85)--(0,-5.85) circle(0.08)--(6.4,-5.85);
		\draw[->,red] (-5,0)--(6.7,0) node [below]{\color{black}$x$};
		\draw[->,red] (M)--(v) node [left=-2pt]{\color{black}$\vec{v}$};
		\draw[->,red] (M)--(vx) node [below left=-5pt]{\color{black}$\vec{v}_{\!\!x}$};
		\draw[red] (O)--(M);
		\draw[dashed,red] (0,0)--(0,-5.85) (v)--(vx);
		\fill (0,0) circle(0.1);
		\path pic[draw=red,angle radius=5]{right angle=v--M--O} 
		(O) pic[draw,->,angle radius=16]{angle=A--O--M} node[shift={(14:23pt)}]{$\alpha$};
		\node at (O)[below]{$O$};
	\end{tikzpicture}	
	}
	\shortans[]{$5$}
	\loigiai{
		Ta có đồ thị của hàm số $v_x$ trên $\left[0;2\pi\right]$ như sau.
			\immini{
				Dựa vào đồ thị ta có trên các khoảng $\left(0;\dfrac{\pi}{2}\right)$ và $\left(\dfrac{3\pi}{2};2\pi\right)$ thì $v_x$ tăng.\\
				Suy ra $a=2$, $b=3$.\\
				Vậy $a+b=5$.
			}{
				\begin{tikzpicture}[>=stealth,x=1.0cm,y=1.0cm, scale=1]
					\draw[->](-1,0) -- (7,0) node[below]{\footnotesize $x$};
					\draw[->](0,-1.5) -- (0,1.85) node[right]{\footnotesize $v_x$};
					\draw[fill = black] (0,0) node [above left]{$O$} circle (1pt);
					\draw[fill = black] (pi,0) node [above right]{\footnotesize $\pi$} circle (1pt);
					\draw[fill = black] (2*pi,0) node [above]{\footnotesize $2\pi$} circle (1pt);
					\draw[fill = black] (pi/2,0) node [below]{\footnotesize $\tfrac{\pi}{2}$} circle (1pt);
					\draw[fill = black] (3*pi/2,0) node [above]{\footnotesize $\tfrac{3\pi}{2}$} circle (1pt);
					\draw[smooth, samples=100, domain=0:6.3]plot(\x,{sin((\x)*180/pi)});
					\draw[dashed] (pi/2,0)--(pi/2,1)--(0,1) node[left]{\footnotesize$0{,}3$} (3*pi/2,0)--(3*pi/2,-1)--(0,-1) node[left]{\footnotesize$-0{,3}$};
				\end{tikzpicture}
			}
	}
\end{ex}

% #2071 | Archive/2024/SP-K10-K11-LyVanHoang_TVN-023/11. C1-LuongGiac-lvhoang_TVN-023/dataC1K11/Bai2-Dang4-Congthucbiendoitongthanhtich.tex | mã gốc 1D1N3-4
\begin{vd}%[BG - 11 New - 4in1, Lê Đạt]%[1D1N3-4]
	Tính
	\begin{listEX}[2]
		\item $\sin \dfrac{11 \pi}{12}-\sin \dfrac{5 \pi}{12}$.
		\item $\cos 105^{\circ}+\cos 15^{\circ}$.	
	\end{listEX}	
	\loigiai{
		\begin{listEX}[1]
			\item $\sin \dfrac{11 \pi}{12}-\sin \dfrac{5 \pi}{12}=2 \cos \dfrac{\dfrac{11 \pi}{12}+\dfrac{5 \pi}{12}}{2} \sin \dfrac{\dfrac{11 \pi}{12}-\dfrac{5 \pi}{12}}{2}=2 \cos \dfrac{2 \pi}{3} \sin \dfrac{\pi}{4}=2 \cdot\left(-\dfrac{1}{2}\right) \cdot \dfrac{\sqrt{2}}{2}=-\dfrac{\sqrt{2}}{2}$. 
			\item $\cos 105^{\circ}+\cos 15^{\circ}=2 \cos \dfrac{105^{\circ}+15^{\circ}}{2} \cos \dfrac{105^{\circ}-15^{\circ}}{2}=2 \cos 60^{\circ} \cos 45^{\circ}=2 \cdot \dfrac{1}{2} \cdot \dfrac{\sqrt{2}}{2}=\dfrac{\sqrt{2}}{2}$.
		\end{listEX}	
	}
\end{vd}

% #2073 | Archive/2024/SP-K10-K11-LyVanHoang_TVN-023/11. C1-LuongGiac-lvhoang_TVN-023/dataC1K11/Bai2-Dang4-Congthucbiendoitongthanhtich.tex | mã gốc 1D1V3-4
\begin{vd}%[BG - 11 New - 4in1, Lê Đạt]%[1D1V3-4]
	Trong Vật lí, phương trình tổng quát của một vật dao động điều hoà cho bởi công thức $x(t)=A \cos (\omega t+\varphi)$, trong đó $t$ là thời điểm (tính bằng giây), $x(t)$ là li độ của vật tại thời điểm $t$, $A$ là biên độ dao động $(A>0)$ và $\varphi \in[-\pi ; \pi]$ là pha ban đầu của dao động. Xét hai dao động điều hoà có phương trình:
	$$
	\begin{aligned}
	& x_1(t)=2 \cos \left(\dfrac{\pi}{3} t+\dfrac{\pi}{6}\right)\,(\mathrm{cm}), \\
	& x_2(t)=2 \cos \left(\dfrac{\pi}{3} t-\dfrac{\pi}{3}\right)\,(\mathrm{cm}).
	\end{aligned}
	$$
	Gọi $x(t)=x_1(t)+x_2(t)$ là dao động tổng hợp. Xét tính đúng sai của các phát biểu sau.
	\choiceTF
	{\True Biên độ giao động của dao động $x_1(t)$ là $2$}
	{\True Pha giao động ban đầu của dao động $x_2(t)$ là $-\dfrac{\pi}{3}$}
	{\True Biên độ của dao động tổng hợp là $2\sqrt{2}$}
	{\True Pha dao động ban đầu của dao động tổng hợp là $-\dfrac{\pi}{12}$}
	\loigiai{
		\begin{itemchoice}
			\itemch Ta có $x_1(t)=2 \cos \left(\dfrac{\pi}{3} t+\dfrac{\pi}{6}\right)$ nên biên độ $A=2$.
			\itemch Ta có $x_2(t)=2 \cos \left(\dfrac{\pi}{3} t-\dfrac{\pi}{3}\right)$ nên ly độ ban đầu $\varphi=-\dfrac{\pi}{3}$.
			\itemch Dao động tổng hợp $x(t)$ có phương trình $x(t)=2 \cos \left(\dfrac{\pi}{3} t+\dfrac{\pi}{6}\right)+2 \cos \left(\dfrac{\pi}{3} t-\dfrac{\pi}{3}\right)$.\\
			Khi đó $x(t)=2 \left[\cos \left(\dfrac{\pi}{3} t+\dfrac{\pi}{6}\right)+ \cos \left(\dfrac{\pi}{3} t-\dfrac{\pi}{3}\right)\right]=2\cdot2\cos\left(\dfrac{\pi}{3}t-\dfrac{\pi}{12}\right)\cos\left(\dfrac{\pi}{4}\right)=2\sqrt{2}\cos\left(\dfrac{\pi}{3}t-\dfrac{\pi}{12}\right)$.\\
			Biên độ của dao động tổng hợp là $A=2\sqrt{2}$.
			\itemch Ta có $x(t)=2\sqrt{2}\cos\left(\dfrac{\pi}{3}t-\dfrac{\pi}{12}\right)$ nên dao động của dao động tổng hợp là $-\dfrac{\pi}{12}$.
		\end{itemchoice}
		
	}
\end{vd}

% #2074 | Archive/2024/SP-K10-K11-LyVanHoang_TVN-023/11. C1-LuongGiac-lvhoang_TVN-023/dataC1K11/Bai2-Dang4-Congthucbiendoitongthanhtich.tex | mã gốc 1D1N3-4
\begin{bt}%[BG - 11 New - 4in1, Lê Đạt]%[1D1N3-4]
	Tính $\cos \dfrac{7 \pi}{12} + \cos \dfrac{\pi}{12}$.
	\loigiai{
		Ta có $\cos \dfrac{7 \pi}{12} + \cos \dfrac{\pi}{12} = 2\cos \dfrac{\dfrac{7 \pi}{12} + \dfrac{\pi}{12}}{2} \cos \dfrac{\dfrac{7 \pi}{12} - \dfrac{\pi}{12}}{2} = 2 \cos \dfrac{\pi}{3} \cos \dfrac{\pi}{4} = 2 \cdot \dfrac{1}{2} \cdot \dfrac{\sqrt{2}}{2} = \dfrac{\sqrt{2}}{2}.$
	} 
\end{bt}

% #2075 | Archive/2024/SP-K10-K11-LyVanHoang_TVN-023/11. C1-LuongGiac-lvhoang_TVN-023/dataC1K11/Bai2-Dang4-Congthucbiendoitongthanhtich.tex | mã gốc 1D1V3-4
\begin{bt}%[BG - 11 New - 4in1, Lê Đạt]%[1D1V3-4]
	Rút gọn biểu thức: $A=\dfrac{\sin x+\sin 2 x+\sin 3 x}{\cos x+\cos 2 x+\cos 3 x}$.
	\loigiai{
		$A=\dfrac{(\sin x+\sin 3x)+\sin 2x}{(\cos x+\cos 3x)+\cos 2x}=\dfrac{2\sin2x\cos x+\sin 2x}{2\cos 2x\cos x+\cos 2x}=\dfrac{\sin 2x(2\cos x+1)}{\cos 2x(2\cos x+1)}=\tan 2x$.
	}
\end{bt}

% #2076 | Archive/2024/SP-K10-K11-LyVanHoang_TVN-023/11. C1-LuongGiac-lvhoang_TVN-023/dataC1K11/Bai2-Dang4-Congthucbiendoitongthanhtich.tex | mã gốc 1D1V3-4
\begin{bt}%[BG - 11 New - 4in1, Lê Đạt]%[1D1V3-4]
	Hiệu điện thế và cường độ dòng điện trong một thiết bị điện lần lượt được cho bởi các biểu thức sau
	$$
	\begin{aligned}
	& u=40 \sin (120 \pi t)+10 \sin (360 \pi t)(\mathrm{V}) \\
	& i=4 \sin (120 \pi t)+\sin (360 \pi t) \quad \text { (A). }
	\end{aligned}
	$$
	\begin{flushright}
		(Nguồn: Ron Larson, Intermediate Algebra, Cengage)
	\end{flushright}
	Biết rằng công suất tiêu thụ tức thời của thiết bị đó được tính theo công thức: $P=u \cdot i$ (W). Hãy viết biểu thức biểu thị công suất tiêu thụ tức thời ở dạng không có lũy thừa và tích của các biểu thức lượng giác.
	\loigiai{
		Ta có:
		$$
		\begin{aligned}
		P & =u \cdot i=[40 \sin (120 \pi t)+10 \sin (360 \pi t)] \cdot[4 \sin (120 \pi t)+\sin (360 \pi t)] \\
		& =160 \sin ^{2}(120 \pi t)+10 \sin ^{2}(360 \pi t)+80 \sin (120 \pi t) \sin (360 \pi t) \\
		& =80[1-\cos (240 \pi t)]+5[1-\cos (720 \pi t)]+40[\cos (360 \pi t-120 \pi t)-\cos (360 \pi t+120 \pi t)] \\
		& =85-80 \cos (240 \pi t)-5 \cos (720 \pi t)+40 \cos (240 \pi t)-40 \cos (480 \pi t) \\
		& =85-40 \cos (240 \pi t)-5 \cos (720 \pi t)-40 \cos (480 \pi t)(\mathrm{W}).
		\end{aligned}
		$$	
	}
\end{bt}

% #2085 | Archive/2024/SP-K10-K11-LyVanHoang_TVN-023/11. C1-LuongGiac-lvhoang_TVN-023/dataC1K11/Bai2-Dang2-Congthucnhandoi.tex | mã gốc 1D1H3-3
\begin{vd}%[Hoàng Thanh Phương, 24-25 BG K11]%[1D1H3-3]
	Cho $\sin a+\cos a=\dfrac{1}{2}$. Tính $\sin 2 a$, $\cos 4 a$.		
	\loigiai{
		\begin{listEX}[1]
			\item Do $\sin a+\cos a=\dfrac{1}{2}$ nên $(\sin a+\cos a)^{2}=\dfrac{1}{4} \Leftrightarrow \sin ^{2} a+\cos ^{2} a+2 \sin a \cos a=\dfrac{1}{4}$\\ hay $1+2 \sin a \cos a=\dfrac{1}{4}$. Suy ra $\sin 2 a=\dfrac{1}{4}-1=-\dfrac{3}{4}$.
			\item Áp dụng công thức nhân đôi, Ta có $\cos 4 a=\cos (2 \cdot 2 a)=1-2 \sin ^{2} 2 a=-\dfrac{1}{8}$.		
		\end{listEX}
	}
\end{vd}

% #2086 | Archive/2024/SP-K10-K11-LyVanHoang_TVN-023/11. C1-LuongGiac-lvhoang_TVN-023/dataC1K11/Bai2-Dang2-Congthucnhandoi.tex | mã gốc 1D1N3-3
\begin{vd}%[Hoàng Thanh Phương, 24-25 BG K11]%[1D1N3-3]
	Cho $\tan \dfrac{a}{2}=-2$. Tính $\tan a$.
	\loigiai{
		Áp dụng công thức nhân đôi, ta có
		$$\tan a=\dfrac{2 \tan\dfrac{a}{2}}{1-\tan ^2 \dfrac{a}{2}}=\dfrac{2 \cdot (-2)}{1-(-2)^2}=-\dfrac{4}{3}.$$
	}
\end{vd}

% #2087 | Archive/2024/SP-K10-K11-LyVanHoang_TVN-023/11. C1-LuongGiac-lvhoang_TVN-023/dataC1K11/Bai2-Dang2-Congthucnhandoi.tex | mã gốc 1D1N3-3
\begin{vd}%[Hoàng Thanh Phương, 24-25 BG K11]%[1D1N3-3]
	Biết $\cos \dfrac{\pi}{6}=\dfrac{\sqrt{3}}{2}$. Tính $\cos \dfrac{\pi}{12}$.
	\loigiai{
		Ta có $\cos ^{2} \dfrac{\pi}{12}=\dfrac{1+\cos \dfrac{\pi}{6}}{2}=\dfrac{2+\sqrt{3}}{4}$. Mà $\cos \dfrac{\pi}{12}>0$ nên $\cos \dfrac{\pi}{12}=\dfrac{\sqrt{2+\sqrt{3}}}{2}$.	
	}
\end{vd}

% #2088 | Archive/2024/SP-K10-K11-LyVanHoang_TVN-023/11. C1-LuongGiac-lvhoang_TVN-023/dataC1K11/Bai2-Dang2-Congthucnhandoi.tex | mã gốc 1D1N3-3
\begin{vd}%[Hoàng Thanh Phương, 24-25 BG K11]%[1D1N3-3]
	Không dùng máy tính, hãy Tính $\sin \dfrac{\pi}{8}$, $\cos \dfrac{\pi}{8}$.
	\loigiai{
		Ta có		
		\begin{itemize}
			\item $\sin ^{2} \dfrac{\pi}{8}=\dfrac{1-\cos \dfrac{\pi}{4}}{2}=\dfrac{2-\sqrt{2}}{4}$. Mà $\sin \dfrac{\pi}{8}>0$ nên $\sin \dfrac{\pi}{8}=\dfrac{\sqrt{2-\sqrt{2}}}{2}$.
			\item $\cos ^{2} \dfrac{\pi}{8}=\dfrac{1+\cos \dfrac{\pi}{4}}{2}=\dfrac{2+\sqrt{2}}{4}$. Mà $\cos \dfrac{\pi}{8}>0$ nên $\cos \dfrac{\pi}{8}=\dfrac{\sqrt{2+\sqrt{2}}}{2}$.					
		\end{itemize}
	}
\end{vd}

% #2091 | Archive/2024/SP-K10-K11-LyVanHoang_TVN-023/11. C1-LuongGiac-lvhoang_TVN-023/dataC1K11/Bai2-Dang2-Congthucnhandoi.tex | mã gốc 1D1H3-3
\begin{bt}%[Hoàng Thanh Phương, 24-25 BG K11]%[1D1H3-3]
	Cho $\sin a=\dfrac{2}{\sqrt{5}}$. Tính $\cos 2 a$, $\cos 4 a$.
	\loigiai{
		\begin{itemize}
			\item $\cos 2a=1-2\sin^2a=1-2\cdot \left(\dfrac{2}{\sqrt 5}\right)^2=-\dfrac{3}{5}$.
			\item $\cos 4a=2\cos^2(2a)-1=2\cdot \left(-\dfrac{3}{5} \right)^2-1=-\dfrac{7}{25}$.
		\end{itemize}
	}
\end{bt}

% #2092 | Archive/2024/SP-K10-K11-LyVanHoang_TVN-023/11. C1-LuongGiac-lvhoang_TVN-023/dataC1K11/Bai2-Dang2-Congthucnhandoi.tex | mã gốc 1D1H3-3
\begin{bt}%[Hoàng Thanh Phương, 24-25 BG K11]%[1D1H3-3]
	Cho $\sin a+\cos a=1$. Tính $\sin 2 a$. 
	\loigiai{
		Ta có $1=(\sin a+\cos a)^2=\sin^2a+2\sin a\cos a+\cos^2a=1+\sin 2a$.\\
		Suy ra: $\sin 2a=0$.
	}
\end{bt}

% #2093 | Archive/2024/SP-K10-K11-LyVanHoang_TVN-023/11. C1-LuongGiac-lvhoang_TVN-023/dataC1K11/Bai2-Dang2-Congthucnhandoi.tex | mã gốc 1D1H3-3
\begin{bt}%[Hoàng Thanh Phương, 24-25 BG K11]%[1D1H3-3]
	Cho $\cos 2 a=\dfrac{1}{3}$ với $\dfrac{\pi}{2}<a<\pi$. Tính $\sin a$, $\cos a$, $\tan a$.
	\loigiai{
		Do $\dfrac{\pi}{2}<a<\pi$ nên $\sin a>0,\, \cos a>0$ và $\tan a>0$.
		Ta có
		\begin{itemize}
			\item $\sin^2a=\dfrac{1-\cos 2a}{2}=\dfrac{1-\frac{1}{3}}{2}=\dfrac{1}{3}\Rightarrow \sin a=\dfrac{\sqrt{3}}{3}$.
			\item $\cos^2a=\dfrac{1+\cos 2a}{2}=\dfrac{1+\frac{1}{3}}{2}=\dfrac{2}{3}\Rightarrow \cos a=\dfrac{\sqrt{6}}{3}$.
			\item $\tan^2a=\dfrac{1-\cos 2a}{1+\cos 2a}=\dfrac{1-\frac{1}{3}}{1+\frac{1}{3}}=\dfrac{1}{2}\Rightarrow \tan a=\dfrac{\sqrt{2}}{2}$.
		\end{itemize}
	}
\end{bt}

% #2094 | Archive/2024/SP-K10-K11-LyVanHoang_TVN-023/11. C1-LuongGiac-lvhoang_TVN-023/dataC1K11/Bai2-Dang2-Congthucnhandoi.tex | mã gốc 1D1H3-3
\begin{bt}%[Hoàng Thanh Phương, 24-25 BG K11]%[1D1H3-3]
	Cho $\tan (a+b)=3$, $\tan (a-b)=2$. Tính $\tan 2 a$, $\tan 2 b$.
	\loigiai{
		\begin{itemize}
			\item $\tan 2a=\tan\left[(a+b)+(a-b)\right]=\dfrac{\tan(a+b)+\tan(a-b)}{1-\tan(a+b)\tan(a-b)}=\dfrac{3+2}{1-3\cdot 2}=-1$.
			\item $\tan 2b=\tan\left[(a+b)-(a-b)\right]=\dfrac{\tan(a+b)-\tan(a-b)}{1+\tan(a+b)\tan(a-b)}=\dfrac{3-2}{1+3\cdot 2}=\dfrac{1}{7}$.
		\end{itemize}
	}
\end{bt}

% #2095 | Archive/2024/SP-K10-K11-LyVanHoang_TVN-023/11. C1-LuongGiac-lvhoang_TVN-023/dataC1K11/Bai2-Dang2-Congthucnhandoi.tex | mã gốc 1D1H3-3
\begin{bt}%[Hoàng Thanh Phương, 24-25 BG K11]%[1D1H3-3]
	Tính các giá trị lượng giác của góc $2 \alpha$, biết
	\begin{listEX}[2] 
		\item[a)] $\sin \alpha = \dfrac{\sqrt{3}}{3}$ và $0 < \alpha < \dfrac{\pi}{2}$;
		\item[b)] $\sin \dfrac{\alpha}{2} = \dfrac{3}{4}$ và $\pi < \alpha < 2 \pi$.
	\end{listEX}
	\loigiai{
		\begin{itemize}
			\item[a)] Do $0 < \alpha < \dfrac{\pi}{2} \Rightarrow \cos \alpha > 0$.\\
			Ta có $\sin^2 \alpha + \cos^2 \alpha = 1 \Rightarrow \cos \alpha = \sqrt{1 - \sin^2 \alpha} = \sqrt{1 - \left(\dfrac{\sqrt{3}}{3}\right)^2} = \dfrac{\sqrt{6}}{3}$.\\
			Suy ra $\cos 2\alpha = 2\cos^2 \alpha - 1 = 2 \left(\dfrac{\sqrt{6}}{3}\right)^2 - 1 = \dfrac{1}{3}$;\\
			$\sin 2\alpha = 2\sin \alpha \cos \alpha = 2 \cdot \dfrac{\sqrt{3}}{3} \cdot \dfrac{\sqrt{6}}{3} = \dfrac{2\sqrt{2}}{3}$;\\
			$\tan 2\alpha  = \dfrac{\sin 2\alpha}{\cos 2\alpha} =2\sqrt{2}\Rightarrow \cot 2\alpha = \dfrac{1}{2\sqrt{2}} $.
			\item[b)] Ta có $\cos \alpha=\cos \left( 2\cdot\dfrac{\alpha}{2}\right) =1-2\sin^2 \dfrac{\alpha}{2}=1-2\cdot \left(\dfrac{3}{4}\right)^2=-\dfrac{1}{8}$.\\
			Ta có $\sin^2\alpha=1-\cos^2\alpha=1-\dfrac{1}{64}=\dfrac{63}{64}$.\\
			Vì $\pi<\alpha<2\pi$ nên $\sin \alpha<0$. Suy ra $\sin \alpha=-\sqrt{\dfrac{63}{64}}=-\dfrac{\sqrt{63}}{8}$.\\
			Do đó $\sin 2\alpha=2\sin\alpha \cos \alpha=2\left(-\dfrac{\sqrt{63}}{8}\right)\left(-\dfrac{1}{8}\right)=\dfrac{\sqrt{63}}{32}$;\\
			$\cos 2\alpha=2\cos^2\alpha-1=2\left(-\dfrac{1}{8}\right)^2-1=-\dfrac{31}{32}$.\\
			Suy ra $\tan 2\alpha=\dfrac{\sin2\alpha}{\cos 2\alpha}=-\dfrac{\sqrt{63}}{31}$ và $\cot 2\alpha=-\dfrac{31}{\sqrt{63}} $.
		\end{itemize}
	} 
\end{bt}

% #2096 | Archive/2024/SP-K10-K11-LyVanHoang_TVN-023/11. C1-LuongGiac-lvhoang_TVN-023/dataC1K11/Bai2-Dang2-Congthucnhandoi.tex | mã gốc 1D1H3-3
\begin{bt}%[Hoàng Thanh Phương, 24-25 BG K11]%[1D1H3-3]
	Tính các giá trị lượng giác của góc $\alpha$, biết
	\begin{listEX}[2] 
		\item[a)] $\cos 2 \alpha = \dfrac{2}{5}$ và $-\dfrac{\pi}{2} < \alpha < 0$;
		\item[b)] $\sin 2 \alpha = -\dfrac{4}{9}$ và $\dfrac{\pi}{2} < \alpha < \dfrac{3 \pi}{4}$.
	\end{listEX}
	\loigiai{
		\begin{itemize}
			\item[a)] Do $-\dfrac{\pi}{2} < \alpha < 0 \Rightarrow \heva{&\cos \alpha > 0\\&\sin \alpha < 0\\&\tan \alpha < 0\\&\cot \alpha < 0.}$\\
			Ta có $\cos 2 \alpha = \dfrac{2}{5} \Leftrightarrow 2\cos^2 \alpha - 1 = \dfrac{2}{5} \Leftrightarrow \cos^2 \alpha = \dfrac{7}{10} \Leftrightarrow \hoac{&\cos \alpha = \sqrt{\dfrac{7}{10}} \\&\cos \alpha = -\sqrt{\dfrac{7}{10}} \;\;\text{(loại)}.}$\\
			Do $\sin^2 \alpha + \cos^2 \alpha = 1 \Rightarrow \sin \alpha = - \sqrt{1 -\cos^2 \alpha} = - \sqrt{1 - \dfrac{7}{10}} = -\sqrt{\dfrac{3}{10}}$.\\
			Suy ra $\tan \alpha = \dfrac{\sin \alpha}{\cos \alpha} = - \sqrt{\dfrac{3}{7}} \Rightarrow \cot \alpha = - \sqrt{\dfrac{7}{3}}$.
			\item[b)] Do $\dfrac{\pi}{2} < \alpha < \dfrac{3 \pi}{4} \Rightarrow \heva{&\sin \alpha > 0\\&\cos \alpha < 0\\&\tan \alpha < 0\\&\cot \alpha < 0.}$\\
			Mặt khác $\dfrac{\pi}{2} < \alpha < \dfrac{3 \pi}{4} \Rightarrow \pi < 2\alpha < \dfrac{3 \pi}{2} \Rightarrow \cos 2\alpha < 0$.\\
			Ta có $\sin^2 2\alpha + \cos^2 2\alpha = 1 \Rightarrow \cos 2 \alpha = -\sqrt{1 - \sin^2 2\alpha} = -\sqrt{1 - \left(-\dfrac{4}{9}\right)^2} = -\dfrac{\sqrt{65}}{9}.$\\
			Ta có 
			$$\cos 2 \alpha = -\dfrac{\sqrt{65}}{9} \Leftrightarrow 2\cos^2 \alpha - 1 = -\dfrac{\sqrt{65}}{9} \Leftrightarrow \cos^2 \alpha = \dfrac{9 - \sqrt{65}}{18} \Leftrightarrow \hoac{&\cos \alpha = \sqrt{\dfrac{9 - \sqrt{65}}{18}} \;\;\text{(loại)}\\&\cos \alpha = -\sqrt{\dfrac{9 - \sqrt{65}}{18}}.}$$
			Do $\sin^2 \alpha + \cos^2 \alpha = 1 \Rightarrow \sin \alpha = \sqrt{1 -\cos^2 \alpha} = \sqrt{1 - \dfrac{9 - \sqrt{65}}{18}} = \sqrt{\dfrac{9 + \sqrt{65}}{18}}$.\\
			Suy ra $\tan \alpha = \dfrac{\sin \alpha}{\cos \alpha} = - \sqrt{\dfrac{9 + \sqrt{65}}{9 - \sqrt{65}}} \Rightarrow \cot \alpha = - \sqrt{\dfrac{9 - \sqrt{65}}{9 + \sqrt{65}}}$.
		\end{itemize}
	} 
\end{bt}

% #2097 | Archive/2024/SP-K10-K11-LyVanHoang_TVN-023/11. C1-LuongGiac-lvhoang_TVN-023/dataC1K11/Bai2-Dang2-Congthucnhandoi.tex | mã gốc 1D1H3-3
\begin{bt}%[Hoàng Thanh Phương, 24-25 BG K11]%[1D1H3-3]
	Tính $\cos \dfrac{\pi}{8}$ và $\tan \dfrac{\pi}{8}$.
	\loigiai{
		\begin{itemize}
			\item Ta có $\dfrac{\sqrt{2}}{2} = \cos \dfrac{\pi}{4} = \cos \left(2 \cdot \dfrac{\pi}{8}\right) = 2 \cos ^2 \dfrac{\pi}{8} - 1$. Suy ra $\cos ^2 \dfrac{\pi}{8} = \dfrac{2 + \sqrt{2}}{4}$.\\
			Vì $0 < \dfrac{\pi}{8} < \dfrac{\pi}{2}$ nên $\cos \dfrac{\pi}{8} > 0$.	Suy ra $\cos \dfrac{\pi}{8} = \dfrac{\sqrt{2 + \sqrt{2}}}{2}$.
			\item Ta có $1 = \tan \dfrac{\pi}{4} = \tan \left(2 \cdot \dfrac{\pi}{8}\right) = \dfrac{2\tan \dfrac{\pi}{8}}{1 - \tan^2 \dfrac{\pi}{8}}$.\\
			Suy ra $\tan^2 \dfrac{\pi}{8} + 2\tan \dfrac{\pi}{8} - 1= 0 \Leftrightarrow \hoac{&\tan \dfrac{\pi}{8} = -1 + \sqrt{2}\\&\tan \dfrac{\pi}{8} = -1 - \sqrt{2}.}$\\
			Vì $0 < \dfrac{\pi}{8} < \dfrac{\pi}{2}$ nên $\tan \dfrac{\pi}{8} > 0$.	Suy ra $\tan \dfrac{\pi}{8} = -1 + \sqrt{2}$.
		\end{itemize}
	} 
\end{bt}

% #2444 | Archive/SP_DA_SauLuoi/1-TL-THPT-NamKyKhoiNghia-GHKI-NH23-24.tex | mã gốc 1D1H2-2
\begin{bt}%[1D1H2-2][Dự án đề kiểm tra Toán 11 GHKI NH23-24- Bùi Thanh Cương]%[THPT Nam Kỳ Khởi Nghĩa]
	Cho $\sin x=-\dfrac{\sqrt{3}}{4}$ và $x \in\left(\pi ; \dfrac{3 \pi}{2}\right)$. Tính giá trị của $\cos x$, $\tan x$.
	\loigiai
	{Ta có $\sin^2x+\cos^2x=1$ nên $\cos^2x=1-\sin^2x=1-\dfrac{3}{16}=\dfrac{13}{16}\Rightarrow \cos x=\pm\dfrac{\sqrt{13}}{4}$.\\
		Do $x\in\left(\pi; \dfrac{3\pi}{2}\right)$ nên $\cos x<0$, ta chọn $\cos x=-\dfrac{\sqrt{13}}{4}$.\\
		Từ đó suy ra 
		$\tan x =\dfrac{\sin x }{\cos x}=\dfrac{-\dfrac{\sqrt{3}}{4}}{-\dfrac{\sqrt{13}}{4}}=\dfrac{\sqrt{39}}{13}$.\\
		Vậy giá trị  $\cos x= -\dfrac{\sqrt{13}}{4}$ và $\tan x=\dfrac{\sqrt{39}}{13}$.
	}
\end{bt}

% #2445 | Archive/SP_DA_SauLuoi/1-TL-THPT-NamKyKhoiNghia-GHKI-NH23-24.tex | mã gốc 1D1H3-3
\begin{bt}%[1D1H3-3][Dự án đề kiểm tra Toán 11 GHKI NH23-24- Bùi Thanh Cương]%[THPT Nam Kỳ Khởi Nghĩa]
	Cho $\tan x=\sqrt{2}$. Tính $\cos 2x$.
	\loigiai
	{Áp dụng công thức nhân đôi ta có $\cos 2x=\cos^2x-\sin^2x$. \hfill $(*)$\\
		Do $\tan x=\sqrt{2}$ nên $\cos x\neq 0$. Chia tử và mẫu biểu thức $(*)$ cho $\cos^2x$ ta được
		\[ \cos 2x=\dfrac{\dfrac{\cos^2x-\sin^2x}{\cos^2x}}{\dfrac{1}{\cos^2 x}}=\dfrac{1-\tan^2x}{1+\tan^2 x}=\dfrac{1-2}{1+2}=-\dfrac{1}{3}.\]  
		Vậy giá trị $\cos 2x=-\dfrac{1}{3}$.
	}
\end{bt}

% #2446 | Archive/SP_DA_SauLuoi/1-TL-THPT-NamKyKhoiNghia-GHKI-NH23-24.tex | mã gốc 1D1V3-5
\begin{bt}%[1D1V3-5]%[Dự án đề kiểm tra Toán 11 GHKI NH23-24- Bùi Thanh Cương]%[THPT Nam Kỳ Khởi Nghĩa]
	Rút gọn biểu thức $T=\dfrac{(\sin 4 x-\sin 2 x) \cdot \cos 5 x}{2 \cos ^2 4 x+\cos 2 x-1}$.
	\loigiai
	{Ta có
		\allowdisplaybreaks
		\begin{eqnarray*}
			T
			&=& \dfrac{(\sin 4 x-\sin 2 x) \cdot \cos 5 x}{2 \cos ^2 4 x+\cos 2 x-1} \\
			&=&\dfrac{(\sin 4 x-\sin 2 x) \cdot \cos 5 x}{(2\cos ^2 4 x-1)+\cos 2 x} \\
			&=& \dfrac{(\sin 4 x-\sin 2 x) \cdot \cos 5 x}{\cos 8x+\cos 2 x} \\
			&=&\dfrac{2\cos 3x\cdot \sin x\cdot \cos 5 x}{2\cos 5x\cdot \cos 3x} \\
			&=& \sin x.
		\end{eqnarray*}
		Vậy $T=\dfrac{(\sin 4 x-\sin 2 x) \cdot \cos 5 x}{2 \cos ^2 4 x+\cos 2 x-1}=\sin x$.
	}
\end{bt}

% #2459 | Archive/SP_DA_SauLuoi/1-TL-THPT-Marie-Curie-GHKI-NH23-24.tex | mã gốc 1D1V1-5
\begin{bt}%[1D1V1-5]%[Dự án đề kiểm tra Toán 11 GHKI NH23-24- Tư Đô Nguyên]%[THPT Marie-Curie]
\begin{listEX}
\item  Trên đường tròn lượng giác gốc $A$, biết góc lượng giác $(O A, O M)$ có số đo $410^{\circ}$. Hãy cho biết điểm $M$ nằm ở góc phần tư thứ mãy trên đường tròn lượng giác?
\item  Không dùng máy tính cầm tay, hãy tính giá trị $\sin \dfrac{13 \pi}{6}$.
\end{listEX}
\loigiai{
\begin{listEX}
	\item \immini{
		Ta có $410^\circ=360^\circ+50^\circ$. \\
		Vậy điểm $M$ biểu diễn góc lượng giác $410^\circ$ nằm trên phần đường tròn lượng giác thuộc góc phần tư thứ I sao cho $\widehat{AOM}=50^\circ$.%Điền câu hỏi
	}
	{
		\begin{tikzpicture}[line join=round, line cap=round,>=stealth,thick,scale=.6]
			\tikzset{every node/.style={scale=0.9}}
			\draw[->] (-2.4,0)--(2.4,0) node[below] {$x$};
			\draw[->] (0,-2.4)--(0,2.4) node[left] {$y$};
			\draw (0,0) node [below left] {$O$};
			\draw (0,0) circle (2);
			\path (2,0) coordinate (A);
	\foreach \x/\g in {A/55}
			\fill[black] 	(\x) circle (3pt)
			($(\g:5mm)+(\x)$) node {$\x$};
			\foreach \i in {0,...,1}
			{
				\draw plot[mark=*] coordinates {({410+\i*360/1}:2)}node[above right]{$M$};
			}
		\end{tikzpicture}%Chèn hình
	}
	\item $\sin \dfrac{13 \pi}{6}=\sin\left(2\pi+\dfrac{\pi}{6}\right)=\sin\dfrac{\pi}{6}=\dfrac{1}{2}$.
\end{listEX}

}
\end{bt}

% #2460 | Archive/SP_DA_SauLuoi/1-TL-THPT-Marie-Curie-GHKI-NH23-24.tex | mã gốc 1D1V2-4
\begin{bt}%[1D1V2-4]%[Dự án đề kiểm tra Toán 11 GHKI NH23-24- Tư Đô Nguyên]%[THPT Marie-Curie]
Cho góc $\alpha$ thỏa mãn $\cos \alpha=\dfrac{3}{5}$ và $0<\alpha<\dfrac{\pi}{2}$.
\begin{listEX}
\item  Tính giá tri $\sin \alpha$.
\item  Tính giá tri $\cos \left(\dfrac{\pi}{2}-2 \alpha\right)$.	
\end{listEX}
\loigiai{
\begin{listEX}
\item  
Vì $0<\alpha<\dfrac{\pi}{2}$ nên $\sin \alpha>0$. Khi đó:
$$
\sin ^2 \alpha=1-\cos ^2 \alpha=\dfrac{16}{25} \Leftrightarrow\hoac{
&	\sin \alpha=\dfrac{4}{5}\,(\text{Nhận}) \\
&	\sin \alpha=-\dfrac{4}{5}\,(\text{Loại}).
}
$$
\item Ta có
$$
\cos \left(\dfrac{\pi}{2}-2 \alpha\right)=\sin 2 \alpha=2 \sin \alpha \cos \alpha=\dfrac{24}{25}.
$$
\end{listEX}	
}
\end{bt}

% #2461 | Archive/SP_DA_SauLuoi/1-TL-THPT-Marie-Curie-GHKI-NH23-24.tex | mã gốc 1D1V2-3
\begin{bt}%[1D1V2-3]%[Dự án đề kiểm tra Toán 11 GHKI NH23-24- Tư Đô Nguyên]%[THPT Marie-Curie]
Chứng minh đẳng thức lượng giác sau trong điều kiện biểu thức có nghĩa
$$
\dfrac{1-\cos \alpha}{\sin ^2 \alpha}-\dfrac{1}{1+\cos \alpha}=0.
$$
\loigiai{
Ta có:
$$\dfrac{1-\cos \alpha}{\sin ^2 \alpha}-\dfrac{1}{1+\cos \alpha}=\dfrac{1-\cos ^2 \alpha-\sin ^2 \alpha}{\sin ^2 \alpha(1+\cos \alpha)}=\dfrac{1-\left(\cos ^2 \alpha+\sin ^2 \alpha\right)}{\sin ^2 \alpha(1+\cos \alpha)}=0.$$
}
\end{bt}

% #2462 | Archive/SP_DA_SauLuoi/1-TL-THPT-Marie-Curie-GHKI-NH23-24.tex | mã gốc 1D1B4-1
\begin{bt}%[1D1B4-1]%[Dự án Tex đề GHK-HK 2023-2024]%[Nhật Thiện]%[THPT Marie Curie]
    Tìm tập xác  định của hàm số $y=\dfrac{2003}{\sin x}+\cos (x-2004)$.
    \loigiai{
        Hàm số có nghĩa khi và chỉ khi $\sin x\ne 0\Leftrightarrow x\ne k\pi$, $k\in \mathbb{Z}$.\\
        Vậy tập xác định $\mathscr{D}=\mathbb{R}\setminus\left\{k\pi,k\in\mathbb{Z}\right\}$.
    }
\end{bt}

% #2463 | Archive/SP_DA_SauLuoi/1-TL-THPT-Marie-Curie-GHKI-NH23-24.tex | mã gốc 1D1B5-3
\begin{bt}%[1D1B5-3]%[Dự án Tex đề GHK-HK 2023-2024]%[Nhật Thiện]%[THPT Marie Curie]
    Giải các phương trình lượng giác sau
    \begin{enumEX}{2}
        \item $\cos  x=\dfrac{\sqrt{3}}{2}$;
        \item $\sin x+\cos \dfrac{\pi}{7}=0$.
    \end{enumEX}
    \loigiai{
        \begin{enumerate}
            \item Ta có  $\cos x=\dfrac{\sqrt{3}}{2}\Leftrightarrow \cos x=\cos \dfrac{\pi}{3}\Leftrightarrow \hoac{&x=\dfrac{\pi}{3}+k2\pi\\&x=-\dfrac{\pi}{3}+k2\pi}$, $k\in\mathbb{Z}$.
            \item Ta có
            \begin{eqnarray*}
                & &\sin x=-\cos \dfrac{\pi}{7}\Leftrightarrow \sin x=-\sin \left(\dfrac{5\pi}{14}\right)\Leftrightarrow \sin x=\sin \left(-\dfrac{5\pi}{14}\right)\\
                &\Leftrightarrow&  \hoac{&x=-\dfrac{5\pi}{14}+k2\pi\\&x=\pi-\dfrac{5\pi}{14}+k2\pi},\,k\in\mathbb{Z}\\
                &\Leftrightarrow& \hoac{&x=-\dfrac{5\pi}{14}+k2\pi\\&x=\dfrac{9\pi}{14}+k2\pi},\,k\in\mathbb{Z}.
            \end{eqnarray*}
        \end{enumerate}
    }
\end{bt}

% #2477 | Archive/SP_DA_SauLuoi/1-TL-THPT-NguyenThiMinhKhai-GHKI23-24.tex | mã gốc 1D1V5-5
\begin{bt}%[1D1V5-5]%[Dự án đề kiểm tra Toán 11 GHKI NH23-24- AHL]%[THPT - NguyenThiMinhKhai-HCM]
Giải các phương trình sau
\begin{multicols}{2}
\begin{enumerate}
\item $\sin 2 x=\sin \dfrac{\pi}{5}$.
\item $\cot \left(x+50^{\circ}\right)=\sqrt{3}$.
\item $\cos \left(3 x-\dfrac{\pi}{4}\right)=\cos 2 x$.
\item $\sin (x+6)=2 \cos ^2 5 x-1$.
\end{enumerate}
\end{multicols}
\loigiai{
\begin{enumerate}
\item  $\sin 2 x=\sin \dfrac{\pi}{5}$\\
$
\Leftrightarrow \hoac{ 
& 2 x = \dfrac { \pi } { 5 } + k 2 \pi  \\
& 2 x = \pi - \dfrac { \pi } { 5 } + k 2 \pi 
} \Leftrightarrow 
\hoac{
& x=\dfrac{\pi}{10}+k \pi \\
& x=\dfrac{2 \pi}{5}+k \pi
} (k \in \mathbb{Z})
$.
\item 
$  \cot \left(x+50^{\circ}\right)=\sqrt{3} \\
 \Leftrightarrow \cot \left(x+50^{\circ}\right)=\cot 30^{\circ} \\
 \Leftrightarrow x+50^{\circ}=30^{\circ}+k 180^{\circ} \\
  \Leftrightarrow x=-20^{\circ}+k 180^{\circ} \quad(k \in \mathbb{Z})
$.
\item $
 \cos \left(3 x-\dfrac{\pi}{4}\right)=\cos 2 x \\
 \Leftrightarrow \hoac{
& 3 x - \dfrac { \pi } { 4 } = 2 x + k 2 \pi  \\
& 3 x - \dfrac { \pi } { 4 } = - 2 x + k 2 \pi 
} \Leftrightarrow \hoac{
& x=\dfrac{\pi}{4}+k 2 \pi \\
& x=\dfrac{\pi}{20}+k \dfrac{2 \pi}{5}
} (k \in \mathbb{Z})
$.
\item  $\sin (x+6)=2 \cos ^2 5 x-1 \\
 \Leftrightarrow \sin (x+6)=\cos 10 x \\
 \Leftrightarrow \sin (x+6)=\sin \left(\dfrac{\pi}{2}-10 x\right) \\
 \Leftrightarrow \hoac{
& x + 6 = \dfrac { \pi } { 2 } - 1 0 x + k 2 \pi  \\
& x + 6 = \pi - \dfrac { \pi } { 2 } + 1 0 x + k 2 \pi 
} \Leftrightarrow \hoac{
& x=\dfrac{\pi}{22}-\dfrac{6}{11}+k \dfrac{2 \pi}{11} \\
& x=-\dfrac{\pi}{18}+\dfrac{2}{3}-k \dfrac{2 \pi}{9}
}(k \in \mathbb{Z})
$.
\end{enumerate}
}
\end{bt}

% #2480 | Archive/SP_DA_SauLuoi/1-TL-THPT-NguyenThiMinhKhai-GHKI23-24.tex | mã gốc 1D1V5-6
\begin{bt}%[1D1V5-6]%[Dự án đề kiểm tra Toán 11 GHKI NH23-24- Dương Phước Sang]%[THPT Nguyễn Thị Minh Khai - Tp HCM]
		Cho hai vật dao động điều hòa theo phương trình lần lượt là $x_1(t)=8\cos\left(4\pi t+\dfrac{\pi}{2}\right)$ $(\mathrm{cm})$ và $x_2(t)=-8\cos (4\pi t)$ $(\mathrm{cm})$, trong đó $x_1(t)$, $x_2(t)$ là li độ của hai vật tại thời điểm $t$ (giây). Khi hai vật dao động trong thời gian từ $0$ đến $10$ giây, hỏi chúng có cùng li độ mấy lần?
		\loigiai{
			Thời điểm hai vật có cùng li độ ta có $8\cos\left(4\pi t+\dfrac{\pi}{2}\right)=-8\cos (4\pi t)$
			$$\begin{aligned}
				&\Leftrightarrow 
				 \cos\left(4\pi t+\dfrac{\pi}{2}\right)=\cos (\pi-4\pi t)\\
				 & \Leftrightarrow \hoac{&4\pi t+\dfrac{\pi}{2}=\pi-4\pi t+k2\pi\\&4\pi t+\dfrac{\pi}{2}=-\pi+4\pi t+k2\pi}\\
				 & \Leftrightarrow \hoac{&8\pi t=\dfrac{\pi}{2}+k2\pi\\&0t=-\dfrac{3\pi}{2}+k2\pi\text{: vô nghiệm}}\\
				 & \Leftrightarrow 
				 t=\dfrac{1}{16}+\dfrac{k}{4} \ (k \in \mathbb{Z}).
			\end{aligned}$$
			Xét $t \in [0;10]$, ta có $0 \leq \dfrac{1}{16}+\dfrac{k}{4} \leq 10 \Leftrightarrow -\dfrac{1}{4} \leq t \leq \dfrac{159}{4} \Leftrightarrow k \in \{0;1;2;\ldots;39\}$ (do $k \in \mathbb{Z}$).\\
			Vậy trong $10$ giây đầu tiên, hai vật có cùng li độ $40$ lần.
		}
	\end{bt}

% #2490 | Archive/SP_DA_SauLuoi/1-TL-THPT-BuiThiXuan-GHKI-NH23-24.tex | mã gốc 1D1N2-2
\begin{bt}%[1D1N2-2]%[Dự án đề kiểm tra Toán 11 GHKI NH23-24- Đạt Lê]%[THPT BÙI THỊ XUÂN - Tp HCM]
	Cho $\tan \alpha=3$ với $0^{\circ}<\alpha<90^{\circ}$. Tính $\cot \alpha,$ $\cos \alpha,$ $\sin \alpha$.
	\loigiai{
		Ta có $ \cot \alpha  =\dfrac{1}{\tan \alpha }= \dfrac{1}{3}$.\\
		Ta lại có $\dfrac{1}{\cos^2\alpha}=1+\tan^2\alpha=1+3^2=10 \Rightarrow \cos^2\alpha=\dfrac{1}{10}\Rightarrow \cos\alpha=\pm\dfrac{\sqrt{10}}{10}$.\\
		Do $0^{\circ}<\alpha<90^{\circ}$ nên $\cos \alpha>0$, do đó $\cos\alpha=\dfrac{\sqrt{10}}{10}.$\\
		Từ đó suy ra $\sin \alpha = \tan \alpha \cdot \cos \alpha = 3\cdot \dfrac{\sqrt{10}}{10} = \dfrac{3\sqrt{10}}{10}. $
	}
\end{bt}

% #2491 | Archive/SP_DA_SauLuoi/1-TL-THPT-BuiThiXuan-GHKI-NH23-24.tex | mã gốc 1D1H3-2
\begin{bt}%[1D1H3-2]%[Dự án đề kiểm tra Toán 11 GHKI NH23-24- Đạt Lê]%[THPT BÙI THỊ XUÂN - Tp HCM]
	Chứng minh rằng $\dfrac{\cos 3 x}{\sin x}+\dfrac{\sin 3 x}{\cos x}=2 \cot 2 x$.
	\loigiai{
		\begin{align*}
			VT&=\dfrac{\cos 3 x}{\sin x}+\dfrac{\sin 3 x}{\cos x}\\
			&=\dfrac{\cos 3 x \cdot \cos x +\sin 3x \cdot \sin x}{\sin x\cdot \cos x}\\
			&=\dfrac{\cos (3x-x)}{\sin x\cdot\cos x}\\
			&=\dfrac{\cos 2x}{\sin x\cdot\cos x}\\
			&=\dfrac{2\cos 2x}{2\sin x\cdot\cos x}\\
			&=\dfrac{2\cos 2x}{\sin 2x}\\
			&=2 \cot 2x .
		\end{align*}
	}
\end{bt}

% #2492 | Archive/SP_DA_SauLuoi/1-TL-THPT-BuiThiXuan-GHKI-NH23-24.tex | mã gốc 1D1H5-3
\begin{bt}%[1D1H5-3]%[Dự án đề kiểm tra Toán 11 GHKI NH23-24- Võ Thị Thùy Trang]%[THPT BÙI THỊ XUÂN - Tp HCM]
	Giải phương trình $\cos \left(x+\dfrac{\pi}{3}\right)=\dfrac{1}{2}$.
	\loigiai{\allowdisplaybreaks
		\begin{eqnarray*}
			\cos \left(x+\dfrac{\pi}{3}\right)=\dfrac{1}{2}&\Leftrightarrow &\cos \left(x+\dfrac{\pi}{3}\right)=\cos \dfrac{\pi}{3}\\
			&\Leftrightarrow & \hoac{&x+\dfrac{\pi}{3}=\dfrac{\pi}{3}+k2\pi\\&x+\dfrac{\pi}{3}=-\dfrac{\pi}{3}+k2\pi} \\
			&\Leftrightarrow & \hoac{&x=k2\pi\\&x=-\dfrac{2\pi}{3}+k2\pi} (k\in \mathbb{Z}).
		\end{eqnarray*}
		Vậy nghiệm của phương trình là $x=k2\pi$ $(k\in \mathbb{Z})$ và $x=-\dfrac{2\pi}{3}+k2\pi$ $(k\in \mathbb{Z})$.
	}
\end{bt}

% #2496 | Archive/SP_DA_SauLuoi/1-TL-TN-THPT-NGUYENCHITHANH-GHKI-NH23-24.tex | mã gốc 1D1B2-2
\begin{bt}%[1D1B2-2]%[Dự án đề kiểm tra Toán 11 GHKI NH23-24- Đoàn Minh Tâm]%[THPT Nguyễn Chí Thanh- Tp HCM]
	Cho $\sin \alpha=\dfrac{12}{13}$ và $\dfrac{\pi}{2}<\alpha<\pi$. Hãy tính $\cos \alpha$, $\cos \left(\dfrac{\pi}{3}-\alpha\right)$.
	\loigiai{
		Do $\dfrac{\pi}{2}<\alpha<\pi$ nên $\heva{&\sin x>0\\&\cos x <0.}$\\
		Ta có $\sin^2\alpha+\cos^2\alpha=1 \Leftrightarrow \cos^2x =\dfrac{25}{169} \Leftrightarrow \hoac{&\cos x =-\dfrac{5}{13} &\text{(nhận)}\\&\cos x=\dfrac{5}{13} &\text{(loại)}.}$\\
		Ta có $\cos \left(\dfrac{\pi}{3}-\alpha\right)=\cos\dfrac{\pi}{3}\cos\alpha + \sin\dfrac{\pi}{3}\sin\alpha=\dfrac{-5+12\sqrt{3}}{26}$.
	}
\end{bt}

% #2497 | Archive/SP_DA_SauLuoi/1-TL-TN-THPT-NGUYENCHITHANH-GHKI-NH23-24.tex | mã gốc 1D1B4-1
\begin{bt}%[1D1B4-1]%[Dự án đề kiểm tra Toán 11 GHKI NH23-24- Đoàn Minh Tâm]%[THPT Nguyễn Chí Thanh- Tp HCM]
	Tìm tập xác định của hàm số $y=\cot \left(2 x-\dfrac{\pi}{3}\right)$.
	\loigiai{Điều kiện xác định $2 x-\dfrac{\pi}{3} \ne k\pi \Leftrightarrow x \ne \dfrac{\pi}{6}+k\dfrac{\pi}{2}$ ($k\in \mathbb{Z}$).\\
		Vậy tập xác định là $\mathscr{D}=\mathbb{R}\setminus\left\lbrace \dfrac{\pi}{6}+k\dfrac{\pi}{2} \mid k \in \mathbb{Z} \right\rbrace$.
	}
\end{bt}

% #2498 | Archive/SP_DA_SauLuoi/1-TL-TN-THPT-NGUYENCHITHANH-GHKI-NH23-24.tex | mã gốc 1D1Y1-2
\begin{bt}%[1D1Y1-2]%[Dự án đề kiểm tra Toán 11 GHKI NH23-24- Đoàn Minh Tâm]%[THPT Nguyễn Chí Thanh- Tp HCM]
	Giải phương trình $2 \cos \left(2 x-\dfrac{\pi}{3}\right)+1=0$.
	\loigiai{
		Ta có
		\begin{eqnarray*}
			&& 2 \cos \left(2 x-\dfrac{\pi}{3}\right)+1=0\\
			&\Leftrightarrow & \cos \left(2 x-\dfrac{\pi}{3}\right)=-\dfrac{1}{2}\\
			&\Leftrightarrow & \hoac{&2 x-\dfrac{\pi}{3}=\dfrac{2\pi}{3}+k2\pi\\&2 x-\dfrac{\pi}{3}=-\dfrac{2\pi}{3}+k2\pi}\\
			&\Leftrightarrow & \hoac{&x=\dfrac{\pi}{2}+k\pi\\&x=-\dfrac{\pi}{6}+k\pi} (k \in \mathbb{Z}).
		\end{eqnarray*}
		Vậy $S=\left\{\dfrac{\pi}{2}+k\pi;-\dfrac{\pi}{6}+k\pi\mid k \in \mathbb{Z}\right\}$.
	}
\end{bt}

% #2499 | Archive/SP_DA_SauLuoi/1-TL-TN-THPT-NGUYENCHITHANH-GHKI-NH23-24.tex | mã gốc 1D1B1-3
\begin{bt}%[1D1B1-3]%[Dự án đề kiểm tra Toán 11 GHKI NH23-24- Đoàn Minh Tâm]%[THPT Nguyễn Chí Thanh- Tp HCM]
	\immini{
		Viết công thức số đo tổng quát của các góc lượng giác $(O A, O M)$ và $(O A, O N)$ trong hình bên.
	}{
		\begin{tikzpicture}[scale=1, font=\footnotesize, line join=round, line cap=round, >=stealth]
			\draw[-stealth] (-2.5,0) -- (3,0)node[below] {$x$};
			\draw[-stealth] (0,-2.5) -- (0,2.5)node[right] {$y$};
			\draw[black] (0,0) circle (2);
			\draw[-stealth] (0:.8) arc (0:135:.8);
			\draw[-stealth] (0:1) arc (0:-120:1);
			\draw[fill=black] (0:0) node[below right]{$O$}--(135:2) circle(1pt) node[left]{$M$};
			\draw[fill=black] (0:0) node[below right]{$O$}--(-120:2) circle(1pt) node[left]{$N$};
			\path (-60:1.2) node[below=-2pt]{$-120^{\circ}$};
			\path (45:1.2) node[above]{$135^\circ$};
			\path (0:2) node[below right]{$A$};
		\end{tikzpicture}
	}
	\loigiai{
		\begin{itemize}
			\item $(O A, O M)=135^\circ + k 360^\circ$ ($k \in \mathbb{Z}$).
			\item $(O A, O N)=-120^\circ +k360^\circ$ ($k \in \mathbb{Z}$).
		\end{itemize}
	}
\end{bt}

% #2500 | Archive/SP_DA_SauLuoi/1-TL-TN-THPT-NGUYENCHITHANH-GHKI-NH23-24.tex | mã gốc 1D1K3-4
\begin{bt}%[1D1K3-4]%[Dự án đề kiểm tra Toán 11 GHKI NH23-24- Hiếu Mai]%[THPT Nguyễn Chí Thanh - Tp HCM]
	\begin{listEX}[1]
		\item Chứng minh biểu thức sau không phụ thuộc vào giá trị của $x$
		$$
		P=\sin \left(\dfrac{\pi}{2}-x\right)+\cos (\pi+x)-\tan \left(\dfrac{\pi}{2}-x\right) \tan (\pi-x).
		$$
		\item Chúng minh rằng: $\dfrac{\sin \alpha+\cos \alpha}{\sin ^3 \alpha}=\dfrac{1-\cot ^4 \alpha}{1-\cot \alpha}$ với mọi $\alpha \neq k \pi, \alpha \neq \dfrac{\pi}{4}+k \pi, k \in \mathbb{Z}$
	\end{listEX}
	\loigiai{
		\begin{listEX}[1]
			\item Ta thấy
			\begin{align*}
				P&= \cos x - \cos x - \cot x \cdot (-\tan x)
				\\
				&= \cot x \cdot \tan x
				\\
				&= 1.
			\end{align*}
			Nên $ P $ không phụ thuộc vào giá trị của $x$.
			\item Do
			\begin{align*}
				\dfrac{1-\cot ^4 \alpha}{1-\cot \alpha} 
				&= \dfrac{(1-\cot\alpha)(1+\cot\alpha)(1+\cot^2\alpha)}{1-\cot\alpha}
				\\
				&= \left(1+ \dfrac{\cos\alpha}{\sin\alpha}\right) \cdot \dfrac{1}{\sin^2\alpha}
				\\
				&= \dfrac{\sin \alpha+\cos \alpha}{\sin ^3 \alpha}
			\end{align*}
			Ta được điều phải chứng minh.
		\end{listEX}
	}
\end{bt}

% #2501 | Archive/SP_DA_SauLuoi/1-TL-TN-THPT-NGUYENCHITHANH-GHKI-NH23-24.tex | mã gốc 1D1T5-6
\begin{bt}%[1D1T5-6]%[Dự án đề kiểm tra Toán 11 GHKI NH23-24- Hiếu Mai]%[THPT Nguyễn Chí Thanh - Tp HCM]
	\immini
	{Một chiếc guồng nước có dạng hình tròn tâm $O$ bán kính $2{,}5$m, trên guồng nước có gắn một chiếc gàu múc nước; trục của guồng nước đặt tại $O$ cách mặt nước $2$m (hình bên). Biết rằng guồng nước quay đều theo chiều ngược chiều kim đồng hồ với tốc độ 1 vòng/phút. Giả sử ban đầu chiếc gàu múc nước ở vị trí $M$, sau $t$ phút $(t>0)$ guồng quay chiếc gàu múc nước đến vị trí điểm $A$. Gọi $h$ (mét) là khoảng cách tính từ điểm $A$ trên guồng nước đến mặt nước.}
	{
		\begin{tikzpicture}[
			>=stealth, line join=round, line cap=round, 
			scale=1, %rotate=90,
			declare function={a=2.5; p=60;}]
			% ve diem
			\path 
			(0,0) coordinate (o)
			(0:a) coordinate (M)
			(180:a) coordinate (E)
			(-90:2) coordinate (N)
			(60:a) coordinate (A)
			;
			\draw [name path=dtO]  (0,0) circle (a);
			
			\draw[fill=blue] (0,0) node[shift={(180:.5)}] {$ O $} circle (2pt) -- (M) circle (2pt) node[shift={(0:.5)}] {$ M $};
			\draw[fill=blue, <->] (0,0) -- (A) circle (2pt) node[shift={(60:.5)}] {$ A $} node[pos=.5, sloped, above] {$ 2,5 $m};
			\draw[dashed]
			(-2.5,-2) node[] (n1) {}-- (2.5,-2) node[] (n2) {}; 
			\draw[dashed, <->] (A) -- ($ (n1)!(A)!(n2) $) node[pos=.3, right] {$ h $};
			\draw[dashed, <->] (o) -- (N) node[pos=.5, left] {$ 2 $m};
			%	\node[anchor=south] at (90:{1.1*a}) {$ abc $};
			\node[] at (0,-3) {Mô phỏng guồng nước};
		\end{tikzpicture}
	}
	\begin{listEX}[1]
		\item Hãy lập công thức tính $h(m)$ theo thời gian $t$ (phút) tính từ
		khi gàu bắt đầu quay từ vị trí $ M $. (Quy ước nếu $h>0$ thì gàu múc ở trên mặt nước, nếu $h<0$ thì gàu múc ở dưới mặt nước).
		\item Tính các thời điểm gàu đạt độ cao lớn nhất so với mặt nước.
	\end{listEX}
	\loigiai{
		\begin{center}
			\begin{tikzpicture}[
				>=stealth, line join=round, line cap=round, 
				scale=1, %rotate=90,
				declare function={a=2.5; p=60;}]
				% ve diem
				\path 
				(0,0) coordinate (o)
				(0:a) coordinate (M)
				(180:a) coordinate (E)
				(-90:2) coordinate (N)
				(60:a) coordinate (A)
				(0:{a/2}) coordinate (K)
				;
				\draw [name path=dtO]  (0,0) circle (a);
				
				\draw[fill=blue] (0,0) node[shift={(180:.5)}] {$ O $} circle (2pt) -- (M) circle (2pt) node[shift={(0:.5)}] {$ M $};
				\draw[fill=blue, <->] (0,0) -- (A) circle (2pt) node[shift={(60:.5)}] {$ A $};
				\draw[dashed]
				(-2.5,-2) node[] (n1) {}-- (2.5,-2) node[] (n2) {}; 
				\draw[dashed, <->] (A) -- ($ (n1)!(A)!(n2) $);
				\draw[dashed, <->] (o) -- (N);
				\foreach \x/\g in {N/135, K/45}{
					\draw[fill=black] (\x) circle (1.5pt)+(\g:.4) node{$\x$};}
			\end{tikzpicture}
		\end{center}
		\begin{listEX}[1]
			\item Mỗi phút gàu quay được một vòng $ (2\pi \text{ rad}) $ nên góc quay $ \widehat{MOA} $ tại thời điểm $ t $ phút là
			$$
			\widehat{MOA} = (2\pi) \cdot t = 2\pi t.
			$$
			Chiều cao $ h $ của gàu so với mặt nước
			$$
			h = ON + AK
			= ON + OA \cdot \sin\widehat{MOA}
			= 2 + 2{,}5\sin(2\pi t).
			$$
			\item Gàu đạt độ cao lớn nhất so với mặt nước khi 
			$$
			h_\text{max} = 2 + 2{,}5 = 4{,}5.
			$$
			khi 
			$$
			\sin(2\pi t) = 1
			\Leftrightarrow
			2\pi t = \dfrac{\pi}{2} + k\cdot 2\pi
			\Leftrightarrow
			t = \dfrac{1}{4} + k
			\quad (k \in \mathbb{N}).
			$$
			Vậy các thời điểm gàu nước đạt độ cao lớn nhất so với mặt nước là $ t = \dfrac{1}{4} + k \quad (k \in \mathbb{N}) $.
		\end{listEX}
	}
\end{bt}

% #2518 | Archive/SP_DA_SauLuoi/1-TL-TN-THPT-VoThiSau-GHKI-NH23-24.tex | mã gốc 1D1H5-3
\begin{bt}%[1D1H5-3]%[Dự án đề kiểm tra Toán 11 GHKI NH23-24- Võ Thị Thùy Trang]%[THPT Võ Thị Sáu - Tp HCM]
	Giải phương trình sau $\cos 2 x+\cos \dfrac{3 \pi}{8}=0$
	\loigiai{
		\allowdisplaybreaks
		\begin{eqnarray*}
			\cos 2 x+\cos \dfrac{3 \pi}{8}=0 &\Leftrightarrow &\cos 2 x=\cos \dfrac{5 \pi}{8}\\ 
			&\Leftrightarrow &
			\hoac{&2x=\dfrac{5 \pi}{8}+k2\pi\\&2x=-\dfrac{5 \pi}{8}+k2\pi}\\
			&\Leftrightarrow & \hoac{&x=\dfrac{5 \pi}{16}+k\pi\\&x=-\dfrac{5 \pi}{16}+k\pi} (k\in \mathbb{Z}).
		\end{eqnarray*}
		Vậy nghiệm của phương trình là $x=\dfrac{5 \pi}{16}+k\pi$ $(k\in \mathbb{Z}) $ và $x=-\dfrac{5 \pi}{16}+k\pi$ $(k\in\mathbb{Z}) $.
	}
\end{bt}

% #2536 | Archive/SP_DA_SauLuoi/1-TL-TN-THPT-TANBINH-GHKI-NH22-23.tex | mã gốc 1D1H3-1
\begin{bt}%[1D1H3-1]%[Dự án đề kiểm tra Toán 11 GHKI NH23-24- Quang Vinh NT]%[THPT Tân Bình - Tp HCM]
	Cho $\sin x=\dfrac{3}{4}$ với $\dfrac{\pi}{2}<x<\pi$. Tính giá trị của $\cos x$ và $\cos \left(x+\dfrac{\pi}{6}\right)$
	\loigiai{
		Do $\dfrac{\pi}{2}<x<\pi$ nên $\cos x<0$.\\
		Suy ra $\cos x =-\sqrt{1-\sin^2 x}=-\sqrt{1-\left(\dfrac{3}{4}\right)^2}=-\dfrac{\sqrt{7}}{4}$.\\
		Ta có 
		\begin{eqnarray*}
		\cos \left(x+\dfrac{\pi}{6}\right)	
		&= & \cos x \cos \dfrac{\pi}{6}-\sin x \sin \dfrac{\pi}{6}\\
			&= & -\dfrac{\sqrt{7}}{4}\cdot \dfrac{\sqrt{3}}{2} - \dfrac{3}{4}\cdot \dfrac{1}{2}\\
				&= & -\dfrac{\sqrt{21}+3}{8}.
		\end{eqnarray*}
	}
\end{bt}

% #2538 | Archive/SP_DA_SauLuoi/1-TL-TN-THPT-TANBINH-GHKI-NH22-23.tex | mã gốc 1D1H5-5
\begin{bt}%[1D1H5-5]%[Dự án đề kiểm tra Toán 11 GHKI NH23-24- Quang Vinh NT]%[THPT Tân Bình - Tp HCM]
	Giải phương trình $\cos 2x +\cos \left(x-\dfrac{2\pi}{3}\right)=0$ và tìm các nghiệm thuộc $\left(-\pi;2\pi\right)$.
	\loigiai{
	Ta có
	\begin{eqnarray*}
		& & \cos 2x +\cos \left(x-\dfrac{2\pi}{3}\right)=0\\
		&\Leftrightarrow & \cos 2x =-\cos \left(x-\dfrac{2\pi}{3}\right)\\
		&\Leftrightarrow & \cos 2x =\cos \left(\pi+x-\dfrac{2\pi}{3}\right)\\		
		&\Leftrightarrow & \cos 2x =\cos \left(x+\dfrac{\pi}{3}\right)\\
		&\Leftrightarrow & \hoac{&2x = x+\dfrac{ \pi}{3}+k2\pi\\&2x = -x-\dfrac{ \pi}{3}+k2\pi}, ~k\in \mathbb{Z}\\
		&\Leftrightarrow & \hoac{&x = \dfrac{ \pi}{3}+k2\pi\\&3x = -\dfrac{ \pi}{3}+k2\pi}, ~k\in \mathbb{Z}\\
		&\Leftrightarrow & \hoac{&x = \dfrac{ \pi}{3}+k2\pi\\&x = -\dfrac{ \pi}{9}+\dfrac{k2\pi}{3}}, ~k\in \mathbb{Z}.
	\end{eqnarray*}
	Suy ra tập hợp các nghiệm thuộc khoảng $\left(-\pi;2\pi\right)$ là  $S=\left\{ \dfrac{\pi}{3};-\dfrac{\pi}{9};\dfrac{5\pi}{9};\dfrac{11\pi}{9}\right\}$.
	}
\end{bt}

% #2568 | Archive/SP_DA_SauLuoi/1-TL-TN-THPT-TruongTHD-GHKI-NH23-24.tex | mã gốc 1D1H3-3
\begin{bt}%%[1D1H3-3]%[Dự án đề kiểm tra Toán 11 GHKI NH23-24- Lê Minh Thiện Anh]%[THPT - Tp HCM]
	Cho $\sin x=\dfrac{1}{3}$ với $90^{\circ}<x<180^{\circ}$. Tính $\cos x$, $\sin 2x$ và $\cos \left(x-30^{\circ}\right)$.
	\loigiai{
		Ta có 
		\begin{eqnarray*}
		& & \sin^2x+\cos^2x=1\\
		&\Leftrightarrow & \left(\dfrac{1}{3}\right)^2+\cos^2x=1\\
		&\Leftrightarrow & \cos^2x=\dfrac{8}{9}\\
		&\Leftrightarrow & \cos x=\pm \dfrac{2\sqrt{2}}{3}.
		\end{eqnarray*}
		Vì $90^{\circ}<x<180^{\circ}$ nên $\cos x=-\dfrac{2\sqrt{2}}{3}$.\\
		Khi đó $\sin 2x=2\sin x\cos x=2\cdot \dfrac{1}{3}\cdot \left( -\dfrac{2\sqrt{2}}{3}\right)=-\dfrac{4\sqrt{2}}{9}$.\\
		$\cos \left(x-30^{\circ}\right)=\cos x\cos 30^\circ+\sin x\sin 30^\circ=-\dfrac{2\sqrt{2}}{3}\cdot\dfrac{\sqrt{3}}{2}+\dfrac{1}{3}\cdot \dfrac{1}{2}=\dfrac{1-2\sqrt{6}}{6}$.
	}
\end{bt}

% #2569 | Archive/SP_DA_SauLuoi/1-TL-TN-THPT-TruongTHD-GHKI-NH23-24.tex | mã gốc 1D1H4-2
\begin{bt}%[1D1H4-2]%[Dự án đề kiểm tra Toán 11 GHKI NH23-24- Lê Minh Thiện Anh]%[THPT - Tp HCM]
	Tìm tập xác định của hàm số $y=-3 \cot \left(x-\dfrac{\pi}{3}\right)+1$.
	\loigiai{
		Hàm số đã cho xác định khi 
		\begin{eqnarray*}
		& & \sin \left(x-\dfrac{\pi}{3}\right)\neq 0\\
		&\Leftrightarrow & x-\dfrac{\pi}{3}\neq k\pi\\
		&\Leftrightarrow & x\neq \dfrac{\pi}{3}+k\pi \,(k\in \mathbb{Z}).
		\end{eqnarray*}
		Vậy tập xác định của hàm số đã cho là $\mathscr{D}=\mathbb{R} \setminus \left\lbrace \dfrac{\pi}{3}+k\pi,k\in\mathbb{Z}\right\rbrace$.
	}
\end{bt}

% #2570 | Archive/SP_DA_SauLuoi/1-TL-TN-THPT-TruongTHD-GHKI-NH23-24.tex | mã gốc 1D1H5-5
\begin{bt}%[1D1H5-5]%[Dự án đề kiểm tra Toán 11 GHKI NH23-24- Lê Minh Thiện Anh]%[THPT - Tp HCM]
	Giải phương trình $\cos x-\sin \left(2x-\dfrac{\pi}{3}\right)=0$.
	\loigiai{
		Ta có
		\begin{eqnarray*}
		& & \cos x-\sin \left(2x-\dfrac{\pi}{3}\right)=0\\
		&\Leftrightarrow & \cos x=\sin \left(2x-\dfrac{\pi}{3}\right)\\
		&\Leftrightarrow & \cos x=\cos \left(\dfrac{5\pi}{6}-2x\right)\\
		&\Leftrightarrow & \hoac{& x= \dfrac{5\pi}{6}-2x+k2\pi \\& x= 2x-\dfrac{5\pi}{6}+k2\pi}\\
		&\Leftrightarrow & \hoac{& x=\dfrac{5\pi}{18}+\dfrac{k2\pi}{3} \\& x=\dfrac{5\pi}{6}+k2\pi} (k\in \mathbb{Z}).
		\end{eqnarray*}
		Vậy phương trình có tập nghiệm $S=\left\lbrace\dfrac{5\pi}{18}+\dfrac{k2\pi}{3},\dfrac{5\pi}{6}+k2\pi, k\in \mathbb{Z} \right\rbrace$.
	}
\end{bt}

% #2571 | Archive/SP_DA_SauLuoi/1-TL-TN-THPT-TruongTHD-GHKI-NH23-24.tex | mã gốc 1D1V4-7
\begin{bt}%[1D1V4-7]%[1D1V4-8]%[Dự án đề kiểm tra Toán 11 GHKI NH23-24- Thinh Ngo]%[THPT Trần Hưng Đạo]
	Nhiệt độ ngoài trời ở một thành phố vào các thời điểm khác nhau trong ngày có thể được mô phỏng bởi công thức $f(t)=29 + 8\cos\left[\dfrac{\pi}{4}\left( 4-\dfrac{t}{3}\right) \right]$, với $f$ tính bằng độ C và $t$ tính bằng giờ ($0\leq t \leq 24$).
	\begin{enumerate}
		\item Nhiệt độ cao nhất trong ngày là bao nhiêu độ C và vào lúc mấy giờ?
		\item Các công nhân công ty môi trường muốn cắt tỉa các cây trồng ở dải phân cách hai làn đường, họ bắt đầu làm việc từ $6$ giờ và kết thúc lúc $17$ giờ trong ngày. Dựa vào đồ thị hàm số côsin (xem hình bên dưới), hãy cho biết họ nên làm việc vào những thời điểm nào trong ngày để tránh nhiệt độ trên $36^\circ$C (kết quả làm tròn đến hàng phần mười).
		\begin{center}
			\begin{tikzpicture}[scale=0.8,>=stealth, font=\footnotesize, line join=round, line cap=round]
				\def\xmin{-10} \def\xmax{10.5} \def\ymin{-1.5} \def\ymax{1.8}
				\draw[->] (-7,0)--(6,0) node [below]{$x$};
				\draw[->] (0,\ymin)--(0,\ymax) node [right]{$y$};
				\node at (0,0) [below right]{$O$};
				\clip (\xmin+3.4,\ymin+0.1) rectangle (\xmax-5,\ymax-0.1);
				\draw[smooth,samples=400,domain=\xmin:\xmax] plot(\x,{cos(\x r)});
				\draw[dashed] (\xmin,1)--(0,1) (-3.14,-1)--(3.14,-1);		
				\foreach \x in {-2*pi,-1.5*pi,-pi,-0.5*pi,0}
				{\draw[fill=black] (\x,cos \x*180/pi) circle (1pt);
					\draw[dashed] (\x,cos \x*180/pi)--(\x,0);
					\draw[fill=black] (-\x,cos -\x*180/pi) circle (1pt);
					\draw[dashed] (-\x,cos \x*180/pi)--(-\x,0);}
				\node at (0,1.3) [left]{$1$};
				\node at (0,-1.3) [left]{$-1$};
				\node at (-2*pi+0.15,0) [below]{$-2\pi$};
				\node at (-1.5*pi,0) [below]{$-\frac{3\pi}{2}$};
				\node at (-pi-0.15,0) [above]{$-\pi$};
				\node at (-0.5*pi,0) [above]{$-\frac{\pi}{2}$};
				\node at (0.5*pi,0) [below]{$\frac{\pi}{2}$};
				\node at (pi-0.1,0) [above]{$\pi$};
				\node at (1.5*pi,0) [above]{$\frac{3\pi}{2}$};
			\end{tikzpicture}
		\end{center}
	\end{enumerate} 
	
	\loigiai{
	\begin{enumerate}
		\item Ta có
		\begin{eqnarray*} 
		&-1&\leq \cos\left[\dfrac{\pi}{4}\left( 4-\dfrac{t}{3}\right) \right]\leq 1\\ &\Leftrightarrow& - 8 \leq  8\cos\left[\dfrac{\pi}{4}\left( 4-\dfrac{t}{3}\right) \right]\leq 8\\ &\Leftrightarrow& 21 \leq 29 + 8\cos\left[\dfrac{\pi}{4}\left( 4-\dfrac{t}{3}\right) \right] \leq 37\\
		&\Leftrightarrow& 21 \leq f(t)\leq 37.
		\end{eqnarray*}
	$\text{Max} f(t) = 37$ khi $\cos\left[\dfrac{\pi}{4}\left( 4-\dfrac{t}{3}\right) \right] = 1\Leftrightarrow \dfrac{\pi}{4}\left( 4-\dfrac{t}{3}\right) = k2\pi \Leftrightarrow t = 3(4-8k)$.\\
	Do $0\leq t \leq 24$ cho nên $-\dfrac{1}{2}\leq k \leq \dfrac{1}{2}$ ta chọn $k = 0 \Rightarrow t = 12$.\\ 
	Vậy nhiệt độ cao nhất trong ngày là $37$ độ C vào lúc $12$ giờ.
	\item 
	\begin{center}
		\begin{tikzpicture}[scale=0.8,>=stealth, font=\footnotesize, line join=round, line cap=round]
			\def\xmin{-10} \def\xmax{10.5} \def\ymin{-1.5} \def\ymax{1.8}
			\draw[->] (-7,0)--(7,0) node [below]{$x$};
			\draw[->] (0,\ymin)--(0,\ymax) node [right]{$y$};
			\node at (0,0) [above right]{$O$};
			\clip (\xmin+3.4,\ymin+0.1) rectangle (\xmax-5,\ymax-0.1);
			\draw[smooth,samples=400,domain=\xmin:\xmax] plot(\x,{cos(\x r)});
			\draw[dashed] (\xmin,1)--(0,1) (-3.14,-1)--(3.14,-1);
			\draw
			(\xmin,7/8)--(\xmax,7/8);
			\foreach \x in {-2*pi,-1.5*pi,-pi,-0.5*pi,0}
			{\draw[fill=black] (\x,cos \x*180/pi) circle (1pt);
				\draw[dashed] (\x,cos \x*180/pi)--(\x,0);
				\draw[fill=black] (-\x,cos -\x*180/pi) circle (1pt);
				\draw[dashed] (-\x,cos \x*180/pi)--(-\x,0);}
			\node at (0,1.3) [left]{$1$};
			\node at (0,-1.3) [left]{$-1$};
			\node at (-2*pi+0.15,0) [below]{$-2\pi$};
			\node at (0.161*pi+0.15,0) [below  ]{$0.5$};
			\node at (-0.161*pi+0.15,0) [below left ]{$-0.5$};
			\node at (4.0,0.7) [above left ]{$y = \frac{7}{8}$};
			\node at (-1.5*pi,0) [below]{$-\frac{3\pi}{2}$};
			\node at (-pi-0.15,0) [above]{$-\pi$};
			\node at (-0.5*pi,0) [above]{$-\frac{\pi}{2}$};
			\node at (0.5*pi,0) [below]{$\frac{\pi}{2}$};
			\node at (pi-0.1,0) [above]{$\pi$};
			\node at (1.5*pi,0) [above]{$\frac{3\pi}{2}$};
			\draw[dashed] (0.5,0)--(0.5,7/8);
			\draw[dashed] (-0.5,0)--(-0.5,7/8);
		\end{tikzpicture}
	\end{center}
	Theo đề bài ta có
	\begin{eqnarray*} 
	 &f(t)& \leq 36 \\&\Leftrightarrow& 29 + 8\cos\left[\dfrac{\pi}{4}\left( 4-\dfrac{t}{3}\right) \right] \leq 36\\ &\Leftrightarrow & \cos\left[\dfrac{\pi}{4}\left( 4-\dfrac{t}{3}\right) \right] \leq \dfrac{7}{8}\\& \Leftrightarrow& -0{,}5  \leq \dfrac{\pi}{4}\left( 4-\dfrac{t}{3}\right) \leq 0{,}5\\& \Leftrightarrow& 10{,}1 \leq t \leq 13{,}9.
	 \end{eqnarray*}
	\end{enumerate}

	Như vậy để tránh nhiệt độ trên $36^\circ$C thì các công nhân làm việc từ $10{,1}$ giờ đến $13{,}9$ giờ.
	}
\end{bt}

% #2582 | Archive/SP_DA_SauLuoi/1-TL-TN-THPT-TranPhu-GHKI-NH23-24.tex | mã gốc 1D1V5-6
\begin{bt}%[Đề GHKI- THPT -Trần Phú- HCM]%[Lê Văn Toàn]%[1D1V5-6]
	Số giờ có ánh sáng mặt trời của một thành phố A trong ngày thứ $t$ của một năm không nhuận được cho bởi hàm số $d(t)=2 \sin \left[\dfrac{\pi(t-107)}{186}\right]+12$ với $t \in \mathbb{N}^*$ và $t \leq 365$. Khi đó ngày tháng nào trong năm thì thành phố A có 11 giờ ánh sáng?
	\loigiai{
		Thành phố A có đúng $11$ giờ có ánh sáng mặt trời thì $d(t)=11$.
		Khi đó
		$$\begin{aligned}
			& 11=2 \sin \left[\dfrac{\pi}{186}(t-107)\right]+12 \\
			& \Leftrightarrow \sin \left[\dfrac{\pi}{186}(t-107)\right]=-\dfrac{1}{2} \\
			& \Leftrightarrow \sin \left[\dfrac{\pi}{186}(t-107)\right]=\sin \left(-\dfrac{\pi}{6}\right) \\
			& \Leftrightarrow \hoac{&\dfrac{\pi}{186}(t-107)=-\dfrac{\pi}{6}+k 2 \pi\\ & \dfrac{\pi}{186}(t-107)=\pi + \dfrac{\pi}{6}+k 2 \pi} \\
			& \Leftrightarrow \hoac{&t=76+372k \\ &t = 262+372k} k\in \mathbb{Z}.
		\end{aligned}$$
		Mà $0 < t \leq 365$ nên
		$\hoac{&0 < 76+372k \leq 365\\& 0<262+372k \leq 365} \Leftrightarrow \hoac{&-0{,}2<k \leq 0{,}98\\&-0{,}7<k \leq 0{,}27.}$\\
		Mà $k \in \mathbb{Z}$ nên suy ra $k=0$ khi đó $t=76+372 \cdot 0=76$ và $t=262+372 \cdot 0=262$.\\
		Vậy Thành phố A có đúng $11$ giờ có ánh sáng mặt trời vào ngày thứ $76$ và $262$ trong năm.
	}
\end{bt}

% #2586 | Archive/SP_DA_SauLuoi/1-TL-THPT-Tenloman-GHKI-NH23-24.tex | mã gốc 1D1B1-5
\begin{bt}%[1D1B1-5]%ID%[Dự án đề kiểm tra Toán 11 GHKI NH23-24- DinhTri]%[THPT Tenloman - Tp HCM]
	\begin{enumerate}
		\item Cho $\cos a=\dfrac{1}{3}$. Tính giá trị của biểu thức $A=2-2 \cdot \sin ^2 a$.
		\item Cho $\sin a=\dfrac{4}{5}\left(\dfrac{\pi}{2}<a<\pi\right)$. Tính giá trị lượng giác: $\cos \left(\dfrac{\pi}{6}+a\right)$
		\item Tìm tập giá trị của hàm số $y=4 \cdot \sin 2 x+1$ trên tập hợp các số thực $\mathbb{R}$.
	\end{enumerate}
	\loigiai{
		\begin{enumerate}
			\item  $  \sin ^2 a=1-\left(\dfrac{1}{3}\right)^2=\dfrac{8}{9} \\
			 A=2 \cos ^2 a=2 \cdot\left(\dfrac{1}{3}\right)^2=\dfrac{2}{9}$
			\item $	 \cos a=\dfrac{-3}{5} \\
			 \cos \left(\dfrac{\pi}{6}+a\right)=\dfrac{-3}{5} \cdot \dfrac{\sqrt{3}}{2}-\dfrac{4}{5} \cdot \dfrac{1}{2}=-\dfrac{3 \sqrt{3}+4}{10}$
			\item $	 -1 \leq \sin 2 x \leq 1 \\
			 -3 \leq y \leq 5 \\
			 T=\left[-3;5\right]$
		\end{enumerate}
	}
\end{bt}

% #2587 | Archive/SP_DA_SauLuoi/1-TL-THPT-Tenloman-GHKI-NH23-24.tex | mã gốc 1D1K2-1
\begin{bt}%[1D1K2-1]%ID%[Dự án đề kiểm tra Toán 11 GHKI NH23-24- DinhTri]%[THPT Tenloman - Tp HCM]
	Giải các phương trình lượng giác sau:
	\begin{enumerate}
		\item $\cos 2 x=\cos \dfrac{\pi}{5}$
		\item $\sin x=\cos 2 x$
	\end{enumerate}
	\loigiai{
		
		\begin{enumerate}
			\item $\cos 2x=\cos \dfrac{\pi}{5}\Leftrightarrow \hoac{&2x=\dfrac{\pi}{5}+k2\pi \\&2x=-\dfrac{\pi}{5}+k2\pi}\Leftrightarrow \hoac{&x=\dfrac{\pi}{10}+k\pi,k\in \mathbb{Z} \\&x=-\dfrac{\pi}{10}+k\pi,k\in \mathbb{Z}}$.
			\item $\cos \left(\dfrac{\pi}{2}-x\right)=\cos 2x\Leftrightarrow \hoac{&-3x=\dfrac{-\pi}{2}+k2\pi \\&x=-\dfrac{\pi}{2}+k2\pi}\Leftrightarrow \hoac{&x=\dfrac{\pi}{6}-\dfrac{k2\pi}{3},k\in \mathbb{Z} \\&x=-\dfrac{\pi}{2}+k2\pi,k\in \mathbb{Z}}$.
		\end{enumerate}
	}
\end{bt}

% #2588 | Archive/SP_DA_SauLuoi/1-TL-THPT-Tenloman-GHKI-NH23-24.tex | mã gốc 1D1G1-5
\begin{bt}%[1D1G1-5]%ID%[Dự án đề kiểm tra Toán 11 GHKI NH23-24- DinhTri]%[THPT Tenloman - Tp HCM]
	Chứng minh đẳng thức sau: $\cos x\left(\dfrac{\sin x}{1+\cos x}+\dfrac{1+\cos x}{\sin x}\right)=2 \cot x$
	\loigiai{
		$VT=\cos x\left(\dfrac{\sin^2x+(1+\cos x)^2}{(1+\cos x)\sin}\right)=\dfrac{2\cos x(1+\cos x)}{\sin x(1+\cos x)}=2\cot x=VP$.	
	}
\end{bt}

% #2589 | Archive/SP_DA_SauLuoi/1-TL-THPT-Tenloman-GHKI-NH23-24.tex | mã gốc 1D1C4-8
\begin{bt}%[1D1C4-8]
\immini{
	\begin{enumerate}
	\item Nếu độ cao của mực nước thủy triều được tính theo công thức $$h(t)=4+2\cos (\dfrac{\pi t}{12}+\dfrac{3\pi}{2})$$ với t là thời gian trong ngày tính theo giờ và độ cao mực nước tính theo đơn vị m. Tính độ cao của mực nước vào lúc 6h sáng.\\
	\item Tuabin gió chuyển đổi động lực di chuyển của gió thành năng lượng điện. Hình bên mô phỏng ba cánh quạt của tuabin gió, các điểm M, N và P là vị trí ban đầu của ba cánh quạt. Giả sử điểm O cách mặt đất 130m và mỗi cánh của tuabin gió dài 70m. Khi điểm N quay được 503 vòng thì khoảng cách từ M đến mặt đất là bao nhiêu biết quạt quay theo chiều kim đồng hồ.(Kết quả làm tròn đến chữ số thập phân thứ hai sau dấu phẩy).
\end{enumerate}}{\begin{tikzpicture}[scale=1, font=\footnotesize, line join=round, line cap=round,>=stealth]  
\def\R{3}
\path
(0,0) coordinate (O)
(20:\R) coordinate (M)
(120:\R) coordinate (N)
(-110:\R) coordinate (P)
(-\R-.2,0) coordinate (E)
(-\R-.2,-2*\R) coordinate (F)
;
\draw[dashed,thick,green] (O) circle (\R);
\draw (N)--(O)--(M) (O)--(P)(E)--(F);

\foreach \x/\g in {O/-45,M/45,N/120,P/-130}\draw[fill=black] (\x) circle (.05) +(\g:.3) node{$\x$};
\draw[fill=black] (E) circle (.05) ;
\draw[fill=black] (F) circle (.05) ;
\node[right] at (-\R-.2,-\R) {$130$ m};
\end{tikzpicture}}
\loigiai{\begin{enumerate}
		\item $\underline{\underline{h}(6)}=6 \mathrm{~m}$
		\item Ta có góc $\widehat{A O M}=15^{\circ}$.Quay 503 vòng thì các điểm vẫn quay lại vị trí ban đầu do đó khoảng cách từ khoảng cách từ $\mathrm{M}$ đến mặt đất là:
		$$
		130+70 \sin \left(15^{\circ}\right)=148.12 m
		$$
	\end{enumerate}
}
\end{bt}

% #2591 | Archive/SP_DA_SauLuoi/1-TL-TN-THPT-PHAMPHUTHU-GHKI-NH23-24.tex | mã gốc 1D1K3-3
\begin{bt}%[1D1K3-3]%[Dự án đề kiểm tra Toán 11 GHKI NH23-24- AHL]%[THPT Pham Phu Thu - Tp HCM]
	Cho $\sin \alpha=-\dfrac{3}{5}$ và $\pi<\alpha<\dfrac{3 \pi}{2}$. Tính
	\begin{multicols}{3}
		\begin{enumerate}
			\item $\cos \alpha$;
			\item $\cos 2 \alpha$;
			\item $\sin \left(\alpha-\dfrac{\pi}{3}\right)$.
		\end{enumerate}
	\end{multicols}
	\loigiai{
		\begin{enumerate}
			\item Do $\pi<\alpha<\dfrac{3\pi}{2} \Rightarrow \cos \alpha<0$.\\
			Ta có $\sin ^2 \alpha+\cos ^2 \alpha=1 \Rightarrow \cos \alpha= - \sqrt{1-\sin ^2 \alpha}=-\sqrt{1-\left(-\dfrac{3}{5}\right)^2}=-\dfrac{4}{5}$.
			\item  Ta có $\cos 2 \alpha=2 \cos ^2 \alpha-1=2\left(-\dfrac{4}{5}\right)^2-1=\dfrac{7}{25}$.
			\item Ta có
			\begin{eqnarray*}
				\sin \left(\alpha-\dfrac{\pi}{3}\right) & = &\sin \alpha \cos \dfrac{\pi}{3}- \cos \alpha \sin \dfrac{\pi}{3} \\
				& = & -\dfrac{3}{5} \cdot \dfrac{1}{2} + \dfrac{4}{5} \cdot \dfrac{\sqrt{3}}{2} \\
				& = & \dfrac{-3+4\sqrt{3}}{10}.
			\end{eqnarray*}
		\end{enumerate}
	}
\end{bt}

% #2592 | Archive/SP_DA_SauLuoi/1-TL-TN-THPT-PHAMPHUTHU-GHKI-NH23-24.tex | mã gốc 1D1K3-4
\begin{bt}%[1D1K3-4]%[Dự án đề kiểm tra Toán 11 GHKI NH23-24- AHL]%[THPT Pham Phu Thu - Tp HCM]
	Với $x$ là giá trị sao cho hai vế đều có nghĩa, chứng minh đẳng thức sau
	$$
	\dfrac{1-(\sin 2 x-\cos 2 x)^2}{\sin 10 x-\sin 2 x}=\dfrac{1}{2 \cos 6 x}.
	$$
	\loigiai{
		Ta có
		\begin{eqnarray*}
			VT &=&\dfrac{1-\left( \sin^2{2x}+\cos^2{2x}-2\sin2x \cos2x \right)}{2\cos6x \sin4x}\\
			&=&\dfrac{1-\left( 1-\sin4x \right)}{2 \cos6x \sin4x}\\
			&=&\dfrac{1}{2 \cos6x}\\
			&=&VP.
		\end{eqnarray*}	
	}
\end{bt}

% #2593 | Archive/SP_DA_SauLuoi/1-TL-TN-THPT-PHAMPHUTHU-GHKI-NH23-24.tex | mã gốc 1D1B4-1
\begin{bt}%[1D1B4-1]%[Dự án đề kiểm tra Toán 11 GHKI NH23-24- AHL]%[THPT Pham Phu Thu - Tp HCM]
	Tìm tập xác định của các hàm số sau
	\begin{multicols}{2}
		\begin{enumerate}
			\item $y=\sin x+\tan \left(x+\dfrac{\pi}{3}\right)$;
			\item $y=\dfrac{\cos x}{1-\sin 4 x}$.
		\end{enumerate}
	\end{multicols}
	\loigiai{
		\begin{enumerate}
			\item Hàm số xác định khi
			$$x+\dfrac{\pi}{3} \ne \dfrac{\pi}{2} + k\pi  \Leftrightarrow x \ne \dfrac{\pi}{6} +k\pi \quad \left( k \in \mathbb{Z} \right). $$
			Vậy tập xác định của hàm số là $\mathscr{D}= \mathbb{R} \setminus \left\{ \dfrac{\pi}{6} +k\pi \mid k \in \mathbb{Z}  \right\}$.
			\item  Hàm số xác định khi
			$$ 1- \sin4x \ne 0 \Leftrightarrow 4x \ne \dfrac{\pi}{2} +k2\pi \Leftrightarrow x \ne \dfrac{\pi}{8} + \dfrac{k\pi}{2} \quad  \left( k \in \mathbb{Z} \right). $$
			Vậy tập xác định của hàm số là $\mathscr{D}= \mathbb{R} \setminus \left\{ \dfrac{\pi}{8} +\dfrac{k\pi}{2} \mid k \in \mathbb{Z}   \right\}$.
		\end{enumerate}
	}
\end{bt}

% #2594 | Archive/SP_DA_SauLuoi/1-TL-TN-THPT-PHAMPHUTHU-GHKI-NH23-24.tex | mã gốc 1D1B4-5
\begin{bt}%%[1D1B4-5]%[Dự án đề kiểm tra Toán 11 GHKI NH23-24- An Do]%[THPT - Tp HCM]
	Tìm tập giá trị của hàm số $y = 2\cos^2 x + 5$ trên $\mathbb{R}$.
	\loigiai{
		Ta có $0\leq \cos^2 x \leq 1 \Leftrightarrow 0\leq 2\cos^2x \leq 2 \Leftrightarrow 5\leq 2\cos^2x + 5 \leq 6\Leftrightarrow 5\leq y \leq 7$.
		\\Vậy tập giá trị của hàm số là $\left[5; 7\right]$.
	}
\end{bt}

% #2595 | Archive/SP_DA_SauLuoi/1-TL-TN-THPT-PHAMPHUTHU-GHKI-NH23-24.tex | mã gốc 1D1B5-3
\begin{bt}%%[1D1B5-3]%[Dự án đề kiểm tra Toán 11 GHKI NH23-24- An Do]%[THPT - Tp HCM]
	Giải các phương trình sau:
	\begin{listEX}[2]
		\item $\cot 3x - \dfrac{\sqrt{3}}{3} = 0$;
		\item $\sin \left(2x - \dfrac{\pi}{5}\right)  -\cos x =0$.
	\end{listEX}
	\loigiai{
		\begin{listEX}[2]
			\item \begin{eqnarray*}
				&&\,\cot 3x - \dfrac{\sqrt{3}}{3} = 0\\
				&\Leftrightarrow&\cot 3x = \dfrac{\sqrt{3}}{3}\\
				&\Leftrightarrow&\cot 3x = \cot \dfrac{\pi}{3}\\
				&\Leftrightarrow&3x = \dfrac{\pi}{3} + k\pi \qquad (k \in \mathbb{Z})\\
				&\Leftrightarrow&x = \dfrac{\pi}{9} + \dfrac{k\pi}{3}.\qquad (k \in \mathbb{Z})
			\end{eqnarray*}
			\item \begin{eqnarray*}
				&&\sin \left(2x - \dfrac{\pi}{5}\right)  -\cos x =0\\
				&\Leftrightarrow&\sin \left(2x - \dfrac{\pi}{5}\right) = \cos x\\
				&\Leftrightarrow& \sin \left(2x - \dfrac{\pi}{5}\right) = \sin \left(\dfrac{\pi }{2} - x\right)\\
				&\Leftrightarrow&\hoac{&2x - \dfrac{\pi}{5} = \dfrac{\pi}{2}-x + k2\pi\\&2x - \dfrac{\pi}{5} = \pi - \dfrac{\pi}{2}+x + k2\pi} \quad(k\in \mathbb{Z})\\
				&\Leftrightarrow&\hoac{&x = \dfrac{7\pi}{30} + \dfrac{k2\pi}{3}\\&x = \dfrac{7\pi}{10} + k2\pi.}\quad(k\in \mathbb{Z})
			\end{eqnarray*}
		\end{listEX}
	}
\end{bt}

% #2596 | Archive/SP_DA_SauLuoi/1-TL-TN-THPT-PHAMPHUTHU-GHKI-NH23-24.tex | mã gốc 1D1K5-6
\begin{bt}%[1D1K5-6][Dự án đề kiểm tra Toán 11 GHKI NH23-24- An Do]%[THPT - Tp HCM]
	Huyết áp là áp lực cần thiết tác động lên thành của động mạch để đưa máu từ tim đên nuôi dưỡng các mô trong cơ thể. Huyết áp được tạo ra do lực co bóp của cơ tim và sức cản của thành động mạch. Mỗi lần tim đập, huyết áp của chúng ta tăng rồi giảm giữa các nhịp. Huyết áp tối đa và huyết áp tối thiểu được gọi tương ứng là huyết áp tâm thu và tâm trương. Chi số huyết áp của chúng ta được viết là huyết áp tâm thu/huyêt áp tâm trương. Chi số huyết áp $120/80$ là bình thường. Giả sử huyết áp của một người nào đó được mô hình hóa bởi hàm số
	$$
	P(t)=100+20 \sin \left(\dfrac{5 \pi}{2} t\right),
	$$
	Trong đó $P\ (t)$ là huyết áp tính theo đơn vị mmHg (milimét thủy ngân) và thời gian $t$ tính theo giây.
	\begin{enumerate}
		\item  Hãy xác định các thời điểm huyết áp của người đó là $80 \mathrm{mmHg}$
		\item  Trong khoảng từ 0 đến 2 giây, hãy xác định số lần huyết áp của người đó là $80 \mathrm{mmHg}$
	\end{enumerate}
	\loigiai{
		\begin{enumerate}
			\item Giải phương trình \begin{eqnarray*}
				&&100+20 \sin \left(\dfrac{5 \pi}{2} t\right)  =  80\\
				&\Leftrightarrow& \sin \left(\dfrac{5 \pi}{2} t\right) =  -1\\
				&\Leftrightarrow& \dfrac{5\pi}{2}t = \dfrac{-\pi}{2} + k2\pi \qquad (k \in \mathbb{Z})\\
				&\Leftrightarrow& t =-\dfrac{1}{5} + \dfrac{4k}{5}. \qquad (k \in \mathbb{Z})\\
			\end{eqnarray*}\\
			Vì $t \geq 0$ nên $t  = -\dfrac{1}{5} + \dfrac{4k}{5}$ với $k \in \mathbb{N^*}$.
			\\Vậy các thời điểm huyết áp người đó là $80$mmHg được biểu diễn bởi công thức $$t  = -\dfrac{1}{5} + \dfrac{4k}{5} \text{~với~} k \in \mathbb{N^*}.$$
			\item Theo yêu cầu đề bài ta cần tìm $k$ thỏa \begin{eqnarray*}
				&&0\leq -\dfrac{1}{5} + \dfrac{4k}{5} \leq 2 \qquad\text{với } k\in \mathbb{N^*}\\
				&\Leftrightarrow& \dfrac{1}{5} \leq \dfrac{4k}{5} \leq \dfrac{11}{5} \qquad\text{với } k\in \mathbb{N^*}\\
				&\Leftrightarrow&\dfrac{1}{4} \leq k \leq \dfrac{11}{4} \qquad\text{với } k\in \mathbb{N^*}\\
				&\Leftrightarrow& k \in \left\{1;2\right\}.
			\end{eqnarray*}
			Vậy có hai lần huyết áp của người đó là $80$mmHg trong khoảng thời gian từ $0$ đến $2$ giây.
		\end{enumerate}
	}
\end{bt}

% #2625 | Archive/SP_DA_SauLuoi/1-TL-THPT-TranKhaiNguyen-GHKI-NH23-24.tex | mã gốc 1D1N1-6
\begin{bt}[1 điểm]%[1D1N1-6]%[Dự án đề kiểm tra Toán 11 GHKI NH23-24 - Don Lee]%[THPT Trần Khai Nguyên - Tp HCM]
	Xác định điểm $M$ và $N$ trên đường tròn lượng giác lần lượt là điểm biểu diễn của các góc lượng giác có số đo bằng $-135^\circ$, $\dfrac{9\pi}{4}$.
	\loigiai{
		\begin{center}
			\begin{tikzpicture}[line join = round, line cap = round, >=stealth, font=\footnotesize, scale=0.8]
				\tikzset{label style/.style={font=\footnotesize}}
				\path (0,0) coordinate (O)
				(3,0) coordinate (A)
				(0:0)++(-135:3) coordinate (M)
				(0:0)++(45:3) coordinate (N)
				;
				\draw[->] (-4,0) -- (4,0) node[above,blue]{$x$};
				\draw[->] (0,-4) -- (0.,4) node[left,blue]{$y$};
				\draw (O) circle (3cm);
				\draw[red,-stealth,smooth,samples=100] plot[domain =0:(9/4)*pi]({.6*(1.12)^(\x) *cos(\x r)},{.6*(1.12)^(\x)*sin(\x r)});
				\draw[blue,-stealth,smooth,samples=100] plot[domain =0:(-3/4)*pi]({1.5*(1.12)^(\x) *cos(\x r)},{1.5*(1.12)^(\x)*sin(\x r)});
				\draw[red] (O)--(N);
				\draw[blue] (O)--(M);
				\foreach \p/\r in {A/-45,M/-100,N/40,O/-40}
				\fill (\p) circle (1pt) node[shift={(\r:3mm)},blue]{$\p$};
				\draw (25:2)node{$\dfrac{9\pi}{4}$} (-50:2)node{$-135^\circ$};
			\end{tikzpicture}
		\end{center}
	}
\end{bt}

% #2626 | Archive/SP_DA_SauLuoi/1-TL-THPT-TranKhaiNguyen-GHKI-NH23-24.tex | mã gốc 1D1H5-3
\begin{bt}[1 điểm]%[1D1H5-3]%[Dự án đề kiểm tra Toán 11 GHKI NH23-24 - Don Lee]%[THPT Trần Khai Nguyên - Tp HCM]
	Giải phương trình $\sin\left(x+\dfrac{\pi}{6}\right)=\dfrac{1}{2}$.
	\loigiai{
		Ta có \allowdisplaybreaks
		\begin{eqnarray*}
			&&\sin\left(x+\dfrac{\pi}{6}\right)=\dfrac{1}{2}\\
			&\Leftrightarrow&\hoac{& x+\dfrac{\pi}{6}=\dfrac{\pi}{6}+k2\pi \\ & x+\dfrac{\pi}{6}=\pi-\dfrac{\pi}{6}+l2\pi}\\
			&\Leftrightarrow&\hoac{& x=k2\pi \\ & x=\dfrac{2\pi}{3}+l2\pi}\;\left(k,l\in\mathbb{Z}\right).
		\end{eqnarray*}
		Vậy phương trình có nghiệm là $\hoac{& x=k2\pi \\ & x=\dfrac{2\pi}{3}+l2\pi}\;\left(k,l\in\mathbb{Z}\right)$.
	}
\end{bt}

% #2627 | Archive/SP_DA_SauLuoi/1-TL-THPT-TranKhaiNguyen-GHKI-NH23-24.tex | mã gốc 1D1T5-6
\begin{bt}[1,5 điểm]%[1D1T5-6]%[Dự án đề kiểm tra Toán 11 GHKI NH23-24 - Don Lee]%[THPT Trần Khai Nguyên - Tp HCM]
	Chiều cao $h$ mét của một cabin trên vòng quay vào thời điểm $t$ giây sau khi bắt đầu chuyển động được cho bởi công thức $h(t)=25+20\sin\left(\dfrac{\pi}{30}t+\dfrac{\pi}{3}\right)$.
	\begin{enumerate}[a)]
		\item Cabin đạt độ cao bao nhiêu khi bắt đầu chuyển động và sau khi chuyển động được $30$ giây.
		\item Tìm những thời điểm trong $2,5$ phút đầu tiên mà cabin đạt độ cao $45$ m.
	\end{enumerate}
	\loigiai{
		\begin{enumerate}[a)]
			\item Ta có
			\begin{itemize}
				\item $h(0)=25+20\sin\left(\dfrac{\pi}{3}\right)=25+10\sqrt{3}$ m.
				\item $h(30)=25+20\sin\left(\dfrac{\pi}{30}\cdot 30+\dfrac{\pi}{3}\right)=25-10\sqrt{3}$ m.
			\end{itemize}
			\item Xét phương trình\\
			\allowdisplaybreaks
			$\begin{aligned}[t]
				&\quad 25+20\sin\left(\dfrac{\pi}{30}t+\dfrac{\pi}{3}\right)=45 \Leftrightarrow 20\sin\left(\dfrac{\pi}{30}t+\dfrac{\pi}{3}\right)=20\\
				&\Leftrightarrow \sin\left(\dfrac{\pi}{30}t+\dfrac{\pi}{3}\right)=1 \Leftrightarrow \dfrac{\pi}{30}t+\dfrac{\pi}{3}=\dfrac{\pi}{2}+k2\pi\\
				&\Leftrightarrow t=5+60k, \quad (k\in\mathbb{Z}).
			\end{aligned}$\\
			Có $0\le 5+60k\le 150 \Leftrightarrow -\dfrac{5}{60}\le k\le \dfrac{145}{60}$.\\
			Mà $k\in\mathbb{Z}$ suy ra $k\in \{1;2\}$.\\
			Vậy trong $2,5$ phút đầu tiên cabin đạt độ cao $45$ m tại $t_1=65$ giây và $t_2=125$ giây.
		\end{enumerate}
	}
\end{bt}

% #2632 | Archive/SP_DA_SauLuoi/1-TL-THPT-BaDiem-GHKI-NH23-24.tex | mã gốc 1D1V4-2
\begin{bt}%[1D1V4-2]%[Dự án đề kiểm tra Toán 11 GHKI NH23-24- Tổng Nguyễn]%[THPT Bà Điểm - Tp HCM]
	Tìm tập xác định của hàm số $y=\dfrac{\cot 2x}{2\sin x-1}$.
	\loigiai{Hàm số $y=\dfrac{\cot 2x}{2\sin x-1}$
		xác định khi 
		\[\heva{&2\sin x -1 \neq 0\\&\sin 2x \neq 0}\Leftrightarrow \heva{& x \neq \dfrac{\pi}{6}+k2\pi\\& x\neq \dfrac{5\pi}{6}+k2\pi\\& 2x \neq k\pi}\Leftrightarrow \heva{& x \neq \dfrac{\pi}{6}+k2\pi\\& x\neq \dfrac{5\pi}{6}+k2\pi\\& x \neq \dfrac{k\pi}{2}}, \  \ k \in \mathbb{Z}.\]
		Vậy tập xác định của hàm số đã cho là $\mathscr{D}=\mathbb{R} \setminus \left\{\dfrac{\pi}{6}+k2\pi, \ \dfrac{5\pi}{6}+k2\pi, \  \dfrac{k\pi}{2}\  \text{với } \ k \in \mathbb{Z}\right\}$.
		
	}
\end{bt}

% #2633 | Archive/SP_DA_SauLuoi/1-TL-THPT-BaDiem-GHKI-NH23-24.tex | mã gốc 1D1V5-5
\begin{bt}%[1D1V5-5]%[Dự án đề kiểm tra Toán 11 GHKI NH23-24- Tổng Nguyễn]%[THPT Bà Điểm - Tp HCM]
	Giải các phương trình sau
	\begin{listEX}[3]
		\item $\cos \left(x+\dfrac{5\pi}{6}\right)=-1$; 
		\item $\sin \left(2x-\dfrac{\pi}{6}\right)-\cos x =0$; 
		\item $\cos^2 3x+\cos^2 5x =1$.
	\end{listEX}
	\loigiai{ 
		\begin{listEX}[1]
			\item $\cos \left(x+\dfrac{5\pi}{6}\right)=-1$.\\
			Ta có 
			\[\cos \left(x+\dfrac{5\pi}{6}\right)=-1\Leftrightarrow x+\dfrac{5\pi}{6} =\pi +k2\pi\Leftrightarrow x =\dfrac{\pi}{6}+k2\pi, \ k \in \mathbb{Z}.\]
			Vậy phương trình đã cho  có nghiệm  $x =\dfrac{\pi}{6}+k2\pi, \ k \in \mathbb{Z}$.
			\item $\sin \left(2x-\dfrac{\pi}{6}\right)-\cos x =0$.\\
			Ta có 
			\begin{eqnarray*}
				\sin \left(2x-\dfrac{\pi}{6}\right)-\cos x =0
				&\Leftrightarrow&\sin \left(2x-\dfrac{\pi}{6}\right)=\cos x\\
				&\Leftrightarrow& \sin \left(2x-\dfrac{\pi}{6}\right)=\sin\left(\dfrac{\pi}{2}-x\right)\\
				&\Leftrightarrow& \hoac{&2x-\dfrac{\pi}{6}=\dfrac{\pi}{2}-x+k2\pi\\&2x-\dfrac{\pi}{6}=\pi-\dfrac{\pi}{2}+x+k2\pi}\\
				&\Leftrightarrow &\hoac{&x=\dfrac{2\pi}{9}+\dfrac{k2\pi}{3}\\&x=\dfrac{2\pi}{3}+k2\pi}, \ k \in \mathbb{Z}.
			\end{eqnarray*}
			Vậy phương trình đã cho có nghiệm $x=\dfrac{2\pi}{9}+\dfrac{k2\pi}{3}, \ k \in \mathbb{Z} $ và $x=\dfrac{2\pi}{3}+k2\pi, \ k \in \mathbb{Z}$.
			\item $\cos^2 3x+\cos^2 5x =1$.\\
			Ta có 
			\begin{eqnarray*}
				\cos^2 3x+\cos^2 5x =1 &\Leftrightarrow& \dfrac{1+\cos 6x}{2}+\dfrac{1+\cos 10x}{2}=1\\
				&\Leftrightarrow& \cos 6x +\cos 10=0\\
				&\Leftrightarrow& \cos 6x = \cos \left(\pi -10x\right)\\
				&\Leftrightarrow& \hoac{&6x=\pi-10x+k2\pi\\&6x=-\pi+10x+k2\pi}\\
				&\Leftrightarrow&\hoac{&x=\dfrac{\pi}{16}+\dfrac{k\pi}{8}\\&x=\dfrac{\pi}{4}-\dfrac{k\pi}{2}}, \ k \in \mathbb{Z}.
			\end{eqnarray*}
			Vậy phương trình đã cho có nghiệm $x=\dfrac{\pi}{16}+\dfrac{k\pi}{8},\ k \in \mathbb{Z}$ và $ x=\dfrac{\pi}{4}-\dfrac{k\pi}{2}, \ k \in \mathbb{Z}$.
		\end{listEX}
	}
\end{bt}

% #2636 | Archive/SP_KhoLatex/ID6_old/1-CK1-OT-De-35.tex | mã gốc 1D1V5-5
\begin{bt}%[1D1V5-5]
 Cho phương trình $4\sin \left(x+\dfrac{\pi}{3} \right)\cos \left(x-\dfrac{\pi}{6} \right)=a^2+\sqrt{3} \sin 2x-\cos 2x \  (1)$. Tìm tất cả các giá trị nguyên của tham số $a$ để phương trình $(1)$ có nghiệm.
 \loigiai{
Ta có
\allowdisplaybreaks
\begin{eqnarray*}
&&4\sin \left(x+\dfrac{\pi}{3} \right)\cos \left(x-\dfrac{\pi}{6} \right)=a^2+\sqrt{3} \sin 2x-\cos 2x\\
&\Leftrightarrow& 2\left[1+\sin \left(2x+\dfrac{\pi}{6} \right)\right]=a^2+\sqrt{3} \sin 2x-\cos 2x\\
&\Leftrightarrow& 2+\sqrt{3} \sin 2x+\cos 2x=a^2+\sqrt{3} \sin 2x-\cos 2x\\
&\Leftrightarrow& \cos 2x=\dfrac{a^2-2}{2}.
\end{eqnarray*}
  Do $-1\le \cos 2x\le 1$ nên phương trình có nghiệm khi 
  \allowdisplaybreaks
  \begin{eqnarray*}
  &&-1\le \dfrac{a^2-2}{2} \le 1\\
  &\Leftrightarrow& 0\le a^2 \le 4\\ 
   &\Leftrightarrow& -2\le a\le 2.
  \end{eqnarray*}
  Do $a \in \mathbb{Z}$ nên $a\in\{-2;-1;0;1;2\}$.\\
  Vậy có $5$ số nguyên thỏa mãn yêu cầu.
  }
 \end{bt}

% #2704 | Archive/SP_KhoLatex/ID6_old/1-CK1-OT-De-34.tex | mã gốc 1D1V1-6
\begin{bt}%[1D1V1-6]
Một đồng hồ treo tường có kim giờ dài $10,25$ cm, kim phút dài $13,25$ cm. Trong 30 phút kim giờ vạch nên cung tròn có độ dài bao nhiêu?
\loigiai{
Trong 30 phút, kim giờ quay được $\dfrac{1}{2}$ vòng tròn.\\
Khi đó, kim giờ vạch nên cung tròn có độ dài là $\dfrac{1}{2}\cdot 2\pi \cdot 10{,}25=10{,}25\pi\approx 32{,}2$ (cm).}
\end{bt}

% #2732 | Archive/SP_KhoLatex/ID6_old/1-CK1-OT-De-33.tex | mã gốc 1D1H1-4
\begin{bt}%[1D1H1-4]
	Bánh xe máy có đường kính kể cả lốp xe là $55  $ cm. Nếu xe chạy với vận tốc $40 $ km/h thì trong một giây bánh xe quay được bao nhiêu vòng?
	\loigiai{
		Ta có $C=2 \pi R=d \pi=55 \pi$ cm.\\
		Đổi $40 $ km/h $=\dfrac{10000}{9}$ (cm/s).\\
		Vậy trong một giây bánh xe quay được $\dfrac{\dfrac{10000}{9}}{55 \pi} \approx 6{,}43$ (vòng).
	}
\end{bt}

% #2802 | Archive/SP_KhoLatex/ID6_old/1-CK1-OT-De-32.tex | mã gốc 1D1V6-8
\begin{bt}%[Dự án EX-CK1 10 và 11 mới - Thành Đức Trung]%[1D1V6-8]
	Số giờ có ánh sáng mặt trời của một thành phố A trong ngày thứ $t$ của năm $2023$ được cho bởi một hàm số $y = 4 \sin\left[\dfrac{\pi}{178}(t-60)\right]+10$ với $t \in \mathbb{Z}$ và $0 < t \leq 365$. Vào ngày nào trong năm thì thành phố A có nhiều giờ có ánh sáng mặt trời nhất?
	\loigiai{
		Vì $\sin \left[\dfrac{\pi}{178}(t-60)\right]\leq 1\Rightarrow y=4\sin\left[\dfrac{\pi}{178}(t-60)\right]+10\leq 14$.\\
		Ngày có ánh sáng mặt trời nhiều nhất
		$$\begin{aligned}
			& \Leftrightarrow y=14 \\
			& \Leftrightarrow \sin\left[\dfrac{\pi}{178}(t-60)\right]=1 \\
			& \Leftrightarrow \dfrac{\pi}{178}(t-60)=\dfrac{\pi}{2}+k2\pi \\
			& \Leftrightarrow t=149+356k.
		\end{aligned}$$
		Do $0<t\leq 365\Rightarrow 0<149+356k\leq 365\Leftrightarrow -\dfrac{149}{356}<k\leq \dfrac{54}{89}$, mà $k\in \mathbb{Z}$ nên $k=0$.\\
		Với $k=0\Rightarrow t=149$ rơi vào ngày $29$ tháng $5$ (vì ta đã biết tháng $1$ và $3$ có $31$ ngày, tháng $4$ có $30$ ngày, riêng đối với năm $2023$ thì không phải năm nhuận nên tháng $2$ có $28$ ngày hoặc dựa vào dữ kiện $0<t\leq 365$ thì ta biết năm này tháng $2$ chỉ có $28$ ngày).
	}
\end{bt}

% #2827 | Archive/SP_KhoLatex/ID6_old/1-CK1-OT-De-18.tex | mã gốc 1D1V6-6
\begin{bt}%[1D1V6-6]
Giải phương trình $\dfrac{1}{\cos x}+\dfrac{\sqrt{3}}{\sin x}=2 \sin \left(x+\dfrac{\pi}{3}\right)$.
	\loigiai{
		Điều kiện $\heva{&\sin x \neq 0 \\ &\cos x \neq 0} \Leftrightarrow \sin 2 x \neq 0 \Leftrightarrow 2 x \neq k \pi \Leftrightarrow x \neq \dfrac{k \pi}{2} \quad(k \in \mathbb{Z})$.\\
		Theo đề bài ta có
		\begin{eqnarray*}
			&& \dfrac{1}{\cos x}+\dfrac{\sqrt{3}}{\sin x}=2 \sin \left(x+\dfrac{\pi}{3}\right) \\
			& \Leftrightarrow& \dfrac{\sin x+\sqrt{3} \cos x}{\sin x \cos x}=2 \sin \left(x+\dfrac{\pi}{3}\right) \\
			&\Leftrightarrow& \dfrac{\dfrac{1}{2} \sin x+\dfrac{\sqrt{3}}{2} \cos x}{\sin x \cos x}=\sin \left(x+\dfrac{\pi}{3}\right) \\
			& \Leftrightarrow& \dfrac{\sin \left(x+\dfrac{\pi}{3}\right)}{\sin x \cos x}=\sin \left(x+\dfrac{\pi}{3}\right) \\
			 &\Leftrightarrow& \sin \left(x+\dfrac{\pi}{3}\right)\left(\dfrac{1}{\sin x \cos x}-1\right)=0\\
	&	\Leftrightarrow& \sin \left(x+\dfrac{\pi}{3}\right)\left(\dfrac{2}{\sin 2 x}-1\right)=0 \\
		& \Leftrightarrow&\hoac{&
			\sin \left(x+\dfrac{\pi}{3}\right)=0 \\
			&\sin 2 x=2	} 
		\end{eqnarray*}
	Với $\sin \left(x+\dfrac{\pi}{3}\right)=0 \Leftrightarrow x+\dfrac{\pi}{3}=k \pi \Leftrightarrow x=-\dfrac{\pi}{3}+k \pi \quad(k \in \mathbb{Z}).$\\		
		Với $\sin 2 x=2$ phương trình vô nghiệm.\\
		Kết hợp với điều kiện, vậy nghiệm của phương trình là $x=-\dfrac{\pi}{3}+k \pi \quad(k \in \mathbb{Z})$.
	}
\end{bt}

% #2896 | Archive/SP_KhoLatex/ID6_old/1-CK1-OT-De-31.tex | mã gốc 1D1K5-6
\begin{ex}%[1D1K5-6]
	Số giờ có ánh sáng mặt trời của một thành phố $A$ trong ngày thứ $t$ của năm $2017$ được cho bởi một hàm số $y=4\sin \left[ \dfrac{\pi }{178}\left( t-60 \right) \right]+10$ với $t\in \mathbb{Z}$ và $0<t\le 365$. Vào ngày nào trong năm thì thành phố $A$ có nhiều giờ có ánh sáng mặt trời nhất?
	\loigiai{
		Vì $\sin \left[ \dfrac{\pi }{178}\left( t-60 \right) \right]\le 1\Rightarrow y=4\sin \left[ \dfrac{\pi }{178}\left( t-60 \right) \right]+10\le 14.$\\
		Ngày có ánh sáng mặt trời nhiều nhất\\ 
		$$\Leftrightarrow y=14\Leftrightarrow \sin \left[ \dfrac{\pi }{178}\left( t-60 \right) \right]=1
		\Leftrightarrow \dfrac{\pi }{178}\left( t-60 \right)=\dfrac{\pi }{2}+k2\pi \Leftrightarrow t=149+356k.$$
		Do $0<t\le 365\Rightarrow 0<149+356k\le 365\Leftrightarrow -\dfrac{149}{356}<k\le \dfrac{54}{89}.\,\text{Suy ra}\,{k\in \mathbb{Z}}$ nên $k=0$.\\
		Với $k=0\Rightarrow t=149$ rơi vào ngày $29$ tháng $5$ (vì ta đã biết tháng $1$ và $3$ có $31$ ngày, tháng $4$ có $30$ ngày, riêng đối với năm $2017$ thì không phải năm nhuận nên tháng $2$ có $28$ ngày hoặc dựa vào dữ kiện $0<t\le 365$ thì ta biết năm này tháng $2$ chỉ có $28$ ngày).}
\end{ex}

% #2898 | Archive/SP_KhoLatex/ID6_old/1-CK1-OT-De-31.tex | mã gốc 1D1K1-3
\begin{ex}%[1D1K1-3]
	Một công ty xây dựng khảo sát khách hàng xem họ có nhu cầu mua nhà ở mức giá nào. Kết quả khảo sát được ghi lại ở bảng sau\\
	\begin{center}
		\begin{tabular}{|c|c|c|c|c|c|}
			\hline \begin{tabular}{c} 
				Mức giá \\
				(triệu đồng $\left./ \mathrm{m}^2\right)$
			\end{tabular} & {$[10 ; 14)$} & {$[14 ; 18)$} & {$[18 ; 22)$} & {$[22 ; 26)$} & {$[26 ; 30)$} \\
			\hline Số khách hàng & 54 & 78 & 120 & 45 & 12 \\
			\hline
		\end{tabular}
	\end{center}
	Dựa vào bảng số liệu trên\\
	Công ty nên xây nhà ở mức giá nào để nhiều người có nhu cầu mua nhất?
	\loigiai{
		Nhóm chứa mốt của mẫu số liệu trên là nhóm $\left[ 18;22 \right)$.\\
		Do đó ${{u}_{m}}=18,{{n}_{m-1}}=78,{{n}_{m}}=120,{{n}_{m+1}}=45,{{u}_{m+1}}-{{u}_{m}}=22-18=4$.\\
		Mốt của mẫu số liệu ghép nhóm là\\
		$${{M}_{o}}=18+\dfrac{120-78}{\left( 120-78 \right)+\left( 120-45 \right)}\cdot 4=\dfrac{758}{39}\approx 19{,}4.$$
		Dựa vào kết quả trên ta có thể dự đoán rằng nếu công ty xây nhà ở mức giá 19,4 triệu đồng/$m^2$ thì sẽ có nhiều người có nhu cầu mua nhất.}
\end{ex}

% #2967 | Archive/SP_KhoLatex/ID6_old/1D1Y3-5.tex | mã gốc 1D1Y3-5
\begin{bt}%[Đề thi Học Kì 1, 2016 - 2017 trường THPT Nguyễn Thị Minh Khai, Hồ Chí Minh]%[Phạm Doãn Lê Bình, 11-HK1-18]%[1D1Y3-5]
	Giải phương trình $\cos^2 x - \sin^2 x = \sin 3x + \cos 4x$.
\loigiai{
	Phương trình ban đầu tương đương với 
		\begin{align*}
		& \cos 2x = \sin 3x + \cos 4x \Leftrightarrow \sin 3x + \cos 4x - \cos 2x = 0 \\
		\Leftrightarrow & \sin 3x - 2\sin 3x \sin x = 0 \Leftrightarrow \sin 3x (1-2\sin x)=0\\
		\Leftrightarrow & \hoac{& \sin 3x = 0 \\ & \sin x = \dfrac{1}{2}} \Leftrightarrow \hoac{& x = \dfrac{k\pi}{3} \\ & x = \dfrac{\pi}{6}+k2\pi \\ & x = \dfrac{5\pi}{6}+k2\pi}, (k\in\mathbb{Z}).
		\end{align*}
Vậy tập nghiệm của hệ phương trình là $S=\{ \dfrac{k\pi}{3}; \dfrac{\pi}{6}+k2\pi; \dfrac{5\pi}{6}+k2\pi|k\in\mathbb{Z}\}$.
}
\end{bt}

% #3036 | Archive/SP_KhoLatex/ID6_old/1-GK1-KNTT-De16-NH23-24.tex | mã gốc 1D1B2-2
\begin{bt}%[1D1B2-2]
	Cho $\sin \alpha=\dfrac{3}{5}$ với $0<\alpha<\dfrac{\pi}{2}$. Tính $\cos \alpha, \tan \alpha$.
	\loigiai{
	Ta có \begin{eqnarray*}
			& & \sin ^2 \alpha+\cos ^2 \alpha=1\\ 
			& \Leftrightarrow & \cos ^2 \alpha=1-\sin ^2 \alpha\\
			& \Leftrightarrow &  \cos ^2 \alpha=1-\left(\dfrac{3}{5}\right)^2\\
			& \Leftrightarrow & \cos ^2 \alpha=\dfrac{16}{25}\\ 
			& \Leftrightarrow & \cos \alpha= \pm \dfrac{4}{5}.
		\end{eqnarray*}	
	Vì $0<\alpha<\dfrac{\pi}{2} \Rightarrow \cos \alpha>0 \Rightarrow \cos \alpha=\dfrac{4}{5}$. Ta có $\tan \alpha=\dfrac{\sin \alpha}{\cos \alpha}=\dfrac{\dfrac{3}{5}}{\dfrac{4}{5}}=\dfrac{3}{4}$.
	}
\end{bt}

% #3038 | Archive/SP_KhoLatex/ID6_old/1-GK1-KNTT-De16-NH23-24.tex | mã gốc 1D1G5-5
\begin{bt}%[1D1G5-5]
Giải phương trình: $\cot x+\sin x\left(1+\tan x \tan \dfrac{x}{2}\right)=4$.
	\loigiai{
	Điều kiện: $
	\heva{&\sin x \neq 0\\&\cos x \neq 0\\&\cos \dfrac{x}{2} \neq 0}
	\Leftrightarrow
	\heva{&\sin x \neq 0\\&\cos x \neq 0}
    \Leftrightarrow \sin 2 x \neq 0$.
    	\begin{eqnarray*}
    	& & \cot x+\sin x\left(1+\tan x \tan \dfrac{x}{2}\right)=4\\ 
    	& \Leftrightarrow & \dfrac{\cos x}{\sin x}+\sin x\left(1+\dfrac{\sin x}{\cos x} \cdot \dfrac{\sin \dfrac{x}{2}}{\cos \dfrac{x}{2}}\right)=4\\
    	& \Leftrightarrow &  \dfrac{\cos x}{\sin x}+\sin x\left(\dfrac{\cos x \cos \dfrac{x}{2}+\sin x \cdot \sin \dfrac{x}{2}}{\cos x \cdot \cos \dfrac{x}{2}}\right)=4\\
    	& \Leftrightarrow & \dfrac{\cos x}{\sin x}+\sin x \cdot \dfrac{\cos \dfrac{x}{2}}{\cos x \cdot \cos \dfrac{x}{2}}=4\\ 
    	& \Leftrightarrow & \dfrac{\cos x}{\sin x}+\dfrac{\sin x}{\cos x}=4\\
    	& \Leftrightarrow & \dfrac{\cos ^2 x+\sin ^2 x}{\sin x \cdot \cos x}=4\\
        & \Leftrightarrow & 4 \sin x \cdot \cos x=1\\
        & \Leftrightarrow & 2 \sin 2 x=1 \Leftrightarrow \sin 2 x=\dfrac{1}{2}\\
        & \Leftrightarrow & \hoac{&x=\dfrac{\pi}{12}+k \pi\\&x=\dfrac{5 \pi}{12}+k \pi}, k \in \mathbb{Z}.
    \end{eqnarray*}	
So với điều kiện, ta nhận 
$\hoac{&x=\dfrac{\pi}{12}+k \pi \\&x=\dfrac{5 \pi}{12}+k \pi}, k \in \mathbb{Z}$.
	}
\end{bt}

% #3088 | Archive/SP_KhoLatex/ID6_old/11-OT-PTLuongGiacCoBan.tex | mã gốc 1D1Y2-1
\begin{bt}%[ĐGNL-ĐGTD-2022, Tran Tony]%[1D1Y2-1]
	Tìm tất cả các giá trị của số thực $a$ để phương trình $\cos x = a^2$ có nghiệm.
	\loigiai{
		Phương trình có nghiệm khi $a^2\le 1\Leftrightarrow -1\le a\le 1$.
	}
\end{bt}

% #3089 | Archive/SP_KhoLatex/ID6_old/11-OT-PTLuongGiacCoBan.tex | mã gốc 1D1B2-1
\begin{bt}%[ĐGNL-ĐGTD-2022, Tran Tony]%[1D1B2-1]
	Tìm tất cả các giá trị của số thực $m$ để phương trình $m\sin x = m^2-m$ có nghiệm.
	\loigiai{
	\begin{itemize}
		\item Với $m=0$, phương trình trở thành $0=0$, phương trình có tập nghiệm $S=\mathbb{R}$.
		\item Với $m\neq 0$, phương trình tương đương với $\sin x=m-1$.\\
		Phương trình có nghiệm $\Leftrightarrow \heva{&m\neq 0\\&-1\leq m-1\leq 1}\Leftrightarrow 0< m\leq 2$.
	\end{itemize}
		Từ hai trường hợp trên suy ra các giá trị của $m$ để phương trình có nghiệm là $0\leq m\leq 2$.
	}
\end{bt}

% #3090 | Archive/SP_KhoLatex/ID6_old/11-OT-PTLuongGiacCoBan.tex | mã gốc 1D1B2-1
\begin{bt}%[ĐGNL-ĐGTD-2022, Tran Tony]%[1D1B2-1]
	Giải phương trình $\cos \left(2x+\dfrac{\pi}{3}\right)=\cos \left(x+\dfrac{\pi}{4}\right)$.
	\loigiai{
		Ta có 	
		$$\cos \left(2x+\dfrac{\pi}{3}\right)=\cos \left(x+\dfrac{\pi}{4}\right)\Leftrightarrow \hoac{&2x+\dfrac{\pi}{3}=x+\dfrac{\pi}{4}+k2\pi\\& 2x+\dfrac{\pi}{3}=-x-\dfrac{\pi}{4}+k2\pi}\Leftrightarrow \hoac{&x=-\dfrac{\pi}{12}+k2\pi\\& x=-\dfrac{7\pi}{36}+\dfrac{k2\pi}{3}}\, (k\in \mathbb{Z}).$$
	}
\end{bt}

% #3091 | Archive/SP_KhoLatex/ID6_old/11-OT-PTLuongGiacCoBan.tex | mã gốc 1D1B2-1
\begin{bt}%[ĐGNL-ĐGTD-2022, Tran Tony]%[1D1B2-1]
	Giải phương trình $\cos \left(\dfrac{x}{2}+20^\circ \right)=-\dfrac{\sqrt{3}}{2}$.
	\loigiai{
		Ta có 
		$$\cos \left(\dfrac{x}{2}+20^\circ \right)=-\dfrac{\sqrt{3}}{2}\Leftrightarrow \hoac{&\dfrac{x}{2}+20^\circ=150^\circ+k360^\circ\\& \dfrac{x}{2}+20^\circ=-150^\circ+k360^\circ}\Leftrightarrow \hoac{&x=260^\circ+k720^\circ\\& x=-340^\circ+k720^\circ}\, (k\in \mathbb{Z}).$$	
	}
\end{bt}

% #3092 | Archive/SP_KhoLatex/ID6_old/11-OT-PTLuongGiacCoBan.tex | mã gốc 1D1B2-1
\begin{bt}%[ĐGNL-ĐGTD-2022, Tran Tony]%[1D1B2-1]
	Giải phương trình $\tan 2x=1$.
	\loigiai{
		Ta có
		$$\tan 2x=1 \Leftrightarrow  x=\dfrac{\pi}{8}+\dfrac{k\pi}{2}, k\in \mathbb{Z}.$$
	}
\end{bt}

% #3093 | Archive/SP_KhoLatex/ID6_old/11-OT-PTLuongGiacCoBan.tex | mã gốc 1D1B2-1
\begin{bt}%[ĐGNL-ĐGTD-2022, Tran Tony]%[1D1B2-1]
	Giải phương trình $\cos 3x = \dfrac{1}{3}$.
	\loigiai{Ta có  $\cos 3x = \dfrac{1}{3}\Leftrightarrow \hoac{&3x=\arccos\dfrac{1}{3}+k2\pi\\&3x=-\arccos\dfrac{1}{3}+k2\pi}\Leftrightarrow \hoac{&x=\dfrac{1}{3}\arccos\dfrac{1}{3}+\dfrac{k2\pi}{3}\\&x=\dfrac{1}{3}\arccos\dfrac{1}{3}+\dfrac{k2\pi}{3}}, (k\in \mathbb{Z}).$
	}	
\end{bt}

% #3094 | Archive/SP_KhoLatex/ID6_old/11-OT-PTLuongGiacCoBan.tex | mã gốc 1D1B2-1
\begin{bt}%[ĐGNL-ĐGTD-2022, Tran Tony]%[1D1B2-1]
	Giải phương trình $\sin\left(x-\dfrac{\pi}{4}\right).\cos\left(x-\dfrac{\pi}{6}\right) = 0$.
	\loigiai{
		Ta có  
		\allowdisplaybreaks
		\begin{eqnarray*}
			&&\sin \left(x-\dfrac{\pi}{4}\right)\cdot\cos \left(x-\dfrac{\pi}{6}\right) = 0\\
			&\Leftrightarrow&\hoac{&\sin \left(x-\dfrac{\pi}{4}\right)=0\\&\cos \left(x-\dfrac{\pi}{6}\right) = 0}\\
			&\Leftrightarrow&\hoac{&x-\dfrac{\pi}{4}=k\pi\\&x-\dfrac{\pi}{6}=\dfrac{\pi}{2}+k\pi}\\
			&\Leftrightarrow&\hoac{&x=\dfrac{\pi}{4}+k\pi\\&x=\dfrac{2\pi}{3}+k\pi}(k\in \mathbb{Z}).
		\end{eqnarray*}
		}
\end{bt}

% #3095 | Archive/SP_KhoLatex/ID6_old/11-OT-PTLuongGiacCoBan.tex | mã gốc 1D1B2-1
\begin{bt}%[ĐGNL-ĐGTD-2022, Tran Tony]%[1D1B2-1]
	Giải phương trình $\sin x+\cos \left( x-\dfrac{\pi}{2}\right)=2$.
	\loigiai{
		$\sin x+\cos \left( x-\dfrac{\pi}{2}\right)=2\Leftrightarrow 2\sin x=2\Leftrightarrow \sin x=1\Leftrightarrow x=\dfrac{\pi}{2}+k2\pi,k\in\mathbb{Z}$.
	}
\end{bt}

% #3096 | Archive/SP_KhoLatex/ID6_old/11-OT-PTLuongGiacCoBan.tex | mã gốc 1D1B2-1
\begin{bt}%[ĐGNL-ĐGTD-2022, Tran Tony]%[1D1B2-1]
	Giải phương trình $\tan 3x\cdot\cot 2x=1$.
	\loigiai{
		Điều kiện: $\heva{&\cos 3x\neq 0 \\ &\sin 2x\neq 0}\Leftrightarrow\heva{&x\neq \dfrac{\pi }{6}+k\dfrac{\pi }{3} \\ &x\neq k\dfrac{\pi }{2}}, k\in \mathbb{Z}$.\\
		Phương trình $\Leftrightarrow \tan 3x=\dfrac{1}{\cot 2x}=\tan 2x$ $\Leftrightarrow 3x=2x+k\pi $
		$\Leftrightarrow x=k\pi$, $\ k\in \mathbb{Z}$ (nghiệm này bị loại).\\
		Vậy phương trình đã cho vô nghiệm.
	}
\end{bt}

% #3097 | Archive/SP_KhoLatex/ID6_old/11-OT-PTLuongGiacCoBan.tex | mã gốc 1D1K2-1
\begin{bt}%[ĐGNL-ĐGTD-2022, Tran Tony]%[1D1K2-1]
	Giải phương trình $\sin 4x \cos x =\sin5x \cos2x$.
	\loigiai{Phương trình đã cho tương đương với  $$ \dfrac{1}{2} \left(\sin5x + \sin3x\right) =  \dfrac{1}{2} \left(\sin7x + \sin3x\right) \Leftrightarrow \sin5x = \sin7x \Leftrightarrow \hoac{&x= k\pi\\& x=\dfrac{\pi}{12}+\dfrac{k\pi}{6}}\, (k\in \mathbb{Z}).$$
	}
\end{bt}

% #3098 | Archive/SP_KhoLatex/ID6_old/11-OT-PTLuongGiacCoBan.tex | mã gốc 1D1B2-1
\begin{bt}%[ĐGNL-ĐGTD-2022, Tran Tony]%[1D1B2-1]
	Giải phương trình $\cos x=\dfrac{\sqrt{3}}{2}$ trong khoảng $\left(-\dfrac{\pi}{2},0\right)$.
	\loigiai{
		Phương trình đã cho tương đương
		\[x=\pm \dfrac{\pi}{6}+k2\pi\ (k\in\mathbb{Z}). \]
		Do đó nghiệm của phương trình trên khoảng $\left(-\dfrac{\pi}{2};0\right)$ là $-\dfrac{\pi}{6}$. 
	}
\end{bt}

% #3099 | Archive/SP_KhoLatex/ID6_old/11-OT-PTLuongGiacCoBan.tex | mã gốc 1D1B2-1
\begin{bt}%[ĐGNL-ĐGTD-2022, Tran Tony]%[1D1B2-1]
	Tìm nghiệm dương nhỏ nhất của phương trình $\sin\left(x-\dfrac{\pi}{3}\right)=1$.
	\loigiai{
		Phương trình đã cho tương đương
		\[x=\dfrac{5\pi}{6}+k2\pi\ (k\in \mathbb{Z}).\]
		Suy ra nghiệm cần tìm là $\dfrac{5\pi}{6}$. 
	}
\end{bt}

% #3100 | Archive/SP_KhoLatex/ID6_old/11-OT-PTLuongGiacCoBan.tex | mã gốc 1D1B2-1
\begin{bt}%[ĐGNL-ĐGTD-2022, Tran Tony]%[1D1B2-1]
	Tính tổng các nghiệm của phương trình $\sin\left(2x+\dfrac{\pi}{4}\right)=-1$ thuộc đoạn $\left[0;2\pi\right]$.
	\loigiai{
		Phương trình đã cho tương đương
		\[2x+\dfrac{\pi}{4}=-\dfrac{\pi}{2}+k2\pi\Leftrightarrow x=-\dfrac{3\pi}{8}+k\pi\ (k\in \mathbb{Z}). \]
		Các nghiệm cần tìm là $\dfrac{5\pi}{8}$ và $\dfrac{13\pi}{8}$. Tổng của chúng là $\dfrac{9\pi}{4}$.
	}
\end{bt}

% #3101 | Archive/SP_KhoLatex/ID6_old/11-OT-PTLuongGiacCoBan.tex | mã gốc 1D1K2-1
\begin{bt}%[ĐGNL-ĐGTD-2022, Tran Tony]%[1D1K2-1]
	Tìm nghiệm thuộc khoảng $(0;2\pi)$ của phương trình $\sin x + 2\sin2x + \sin3x=0$.
	\loigiai{
		Phương trình đã cho tương đương với 
		$$ 2\sin 2x\cos x+2\sin 2x=0\Leftrightarrow \sin2x(\cos{x}+1)=0\Leftrightarrow\hoac{&\sin 2x=0\\& \cos x=-1}\Leftrightarrow\hoac{&x=k\dfrac{\pi}{2}\\&x=\pi+k2\pi}\, (k\in\mathbb{Z}).$$
		Do $x\in (0; 2\pi)\Rightarrow x\in \left\{\dfrac{\pi}{2}, \pi, \dfrac{3\pi}{2}\right\}$.
	}
\end{bt}

% #3102 | Archive/SP_KhoLatex/ID6_old/11-OT-PTLuongGiacCoBan.tex | mã gốc 1D1K2-1
\begin{bt}%[ĐGNL-ĐGTD-2022, Tran Tony]%[1D1K2-1]
	Tìm số nghiệm trong đoạn $[-4\pi;5\pi]$ của phương trình $\tan 2x + \tan x = 0$.
	\loigiai{
		Điều kiện: $\heva{ 2x \neq \dfrac{\pi}{2}+k\pi \\
			x\neq \dfrac{\pi}{2}+ k\pi}
		\Leftrightarrow \heva{x\neq \dfrac{\pi}{4}+\dfrac{k\pi}{2} \\ 
			x \neq \dfrac{\pi}{2}+ k\pi}$, $k\in \mathbb{Z}$. \\
		Khi đó,
		\allowdisplaybreaks
		\begin{eqnarray*}
			&&\tan 2x + \tan x= 0\\
			&\Leftrightarrow&\tan 2x = - \tan x\\
			&\Leftrightarrow&\tan 2x = \tan (-x)\\
			&\Leftrightarrow&2x = -x + k\pi\\
			&\Leftrightarrow&3x = k\pi \Leftrightarrow x = \dfrac{k\pi}{3}, k\in \mathbb{Z} \text{ (thỏa điều kiện)}.
		\end{eqnarray*}
		Mà $x \in [-4\pi;5\pi]$ nên $-4\pi \leq \dfrac{k\pi}{3} \leq 5\pi \Leftrightarrow -12 \leq k \leq 15$. \\ 
		Vậy, số nghiệm là $28$.
	}
\end{bt}

% #3104 | Archive/SP_KhoLatex/ID6_old/1.3.1.tex | mã gốc 1D1B3-1
\begin{ex}%[1D1B3-1]%::Cau 1::
	
	Không dùng máy tính cầm tay, tính các giá trị lượng giác của các góc
	\begin{multicols}{2}
		\begin{enumerate}[a)]
		\item  $\dfrac{7\pi }{12}$;
		\item $-195^{\circ}  $.
		\end{enumerate}
	\end{multicols}
	\loigiai{
		\begin{enumerate}[a)]
		\item Ta có
		\begin{itemize}
			\item $\sin \dfrac{7\pi }{12}=\sin \left( \dfrac{\pi }{3}+\dfrac{\pi }{4} \right)=\sin \dfrac{\pi }{3}\cdot \cos \dfrac{\pi }{4}+\cos \dfrac{\pi }{3}\sin \dfrac{\pi }{4}=\dfrac{\sqrt{6}+\sqrt{2}}{4}$.\\
			\item $\cos \dfrac{7\pi }{12}=\cos \left( \dfrac{\pi }{3}+\dfrac{\pi }{4} \right)=\cos \dfrac{\pi }{3}\cdot\cos \dfrac{\pi }{4}-\sin \dfrac{\pi }{3}\sin \dfrac{\pi }{4}=\dfrac{\sqrt{2}-\sqrt{6}}{4}$.\\
			\item $\tan \dfrac{7\pi }{12}=\dfrac{\sqrt{2}+\sqrt{6}}{\sqrt{2}-\sqrt{6}}=-2-\sqrt{3}$.\\
		\end{itemize}
		\item  Ta có
		\begin{itemize}
			\item $\sin \left( -195^{\circ}  \right)=-\sin \left( 180^{\circ}  +15^{\circ}  \right)=\sin 15^{\circ}  =\sin \left( 45^{\circ}  -30^{\circ}  \right)=\dfrac{\sqrt{6}-\sqrt{2}}{4}$.\\
			\item $\cos \left( -195^{\circ}  \right)=\cos 195^{\circ}  =\cos \left( 180^{\circ}  +15^{\circ}  \right)=-\cos 15^{\circ}  =-\cos \left( 45^{\circ}  -30^{\circ}  \right)=\dfrac{\sqrt{6}+\sqrt{2}}{4}$.\\
			\item $\tan \left( -195^{\circ}  \right)=-2+\sqrt{3}.$
		\end{itemize}
	\end{enumerate}
	}
\end{ex}

% #3105 | Archive/SP_KhoLatex/ID6_old/1.3.1.tex | mã gốc 1D1B3-1
\begin{ex}%[1D1B3-1]%::Cau 2::
	
	Tính $\sin \left( \alpha +\dfrac{\pi }{3} \right),\cos \left( \dfrac{\pi }{6}-\alpha \right)$ biết $\cos \alpha =\dfrac{12}{13}$ và $\dfrac{3\pi }{2}<a<2\pi $.
	\loigiai{
		Ta có 
		\begin{itemize}
			\item $\dfrac{3\pi }{2}<a<2\pi \Rightarrow \sin \alpha <0 \Rightarrow \sin \alpha =-\sqrt{1-{\cos ^2}\alpha }=-\dfrac{5}{13}$.\\
			\item $\sin \left( \alpha +\dfrac{\pi }{3} \right)=\sin \alpha \cdot \cos \dfrac{\pi }{3}+\cos \alpha \cdot\sin \dfrac{\pi }{3}=\dfrac{-5+12\sqrt{3}}{26}$.\\
			\item $\cos \left( \dfrac{\pi }{6}+\alpha \right)=\cos \dfrac{\pi }{6}\cdot\cos \alpha -\sin \dfrac{\pi }{6}\cdot \sin \alpha =\dfrac{5+12\sqrt{3}}{26}$.
		\end{itemize}
	}
\end{ex}

% #3106 | Archive/SP_KhoLatex/ID6_old/1.3.1.tex | mã gốc 1D1B3-2
\begin{ex}%[1D1B3-2]%::Cau 3::
	
	Tính giá trị lượng giác của góc $2\alpha $ biết rằng
	\begin{multicols}{2}
		\begin{enumerate}[a)]
			\item  $\sin \alpha =\dfrac{5}{13},\,\,\dfrac{\pi }{2}<\alpha <\pi $;
			\item $\sin \dfrac{\alpha }{2}=\dfrac{1}{4}$ và $\pi <\alpha <2\pi $.
		\end{enumerate}
	\end{multicols}
	\loigiai{
		\begin{enumerate}[a)]
			\item Ta có
		Do $\dfrac{\pi }{2}<\alpha <\pi \Rightarrow \cos \alpha <0\Rightarrow \cos \alpha =-\sqrt{1-{\sin ^2}\alpha }=\dfrac{-12}{13}$.
		\begin{itemize}
			\item $\sin 2\alpha =2\sin \alpha \cdot \cos \alpha =\dfrac{-120}{169}$.\\
			\item $\cos 2\alpha =2{\cos ^2}\alpha -1=\dfrac{119}{169}$.\\
			\item $\tan 2\alpha =\dfrac{\sin 2\alpha }{\cos 2\alpha }=\dfrac{-120}{119}$.\\
			\item $\cot 2\alpha =\dfrac{-119}{120}$.
		\end{itemize}
		\item Ta có
		Do $\pi <\alpha <2\pi \Rightarrow \dfrac{\pi }{2}<\dfrac{\alpha }{2}<\pi \Rightarrow \cos \dfrac{\alpha }{2}<0$.\\
		$\Rightarrow \cos \dfrac{\alpha }{2}=-\sqrt{1-\dfrac{1}{16}}=\dfrac{-\sqrt{15}}{4}$\\
		$\Rightarrow \sin \alpha =2\sin \dfrac{\alpha }{2}\cdot \cos \dfrac{\alpha }{2}=\dfrac{-\sqrt{15}}{8}$.\\
		$\cos \alpha =2{\cos ^2}\dfrac{\alpha }{2}-1=\dfrac{7}{8}$.\\
		$\Rightarrow \cos 2\alpha =2{\cos ^2}\alpha -1=\dfrac{17}{32}$.\\
		$\sin 2\alpha =\dfrac{-7\sqrt{15}}{32}$.\\
		$\tan 2\alpha =\dfrac{-7\sqrt{15}}{17}$, $\cot 2\alpha =\dfrac{-17\sqrt{15}}{105}$.
		\end{enumerate}	
		}
\end{ex}

% #3107 | Archive/SP_KhoLatex/ID6_old/1.3.1.tex | mã gốc 1D1B3-4
\begin{ex}%[1D1B3-4]%::Cau 4::
	
	Rút gọn biểu thức sau
	\begin{multicols}{2}
		\begin{enumerate}[a)]
			\item  $\sqrt{2}\cos \left( \alpha -\dfrac{\pi }{4} \right)-\sin \alpha $;
			\item ${\sin ^4}\alpha +{\cos ^4}\alpha -\dfrac{1}{4}\cos 4\alpha $.
		\end{enumerate}
	\end{multicols}
	\loigiai{
		\begin{enumerate}[a)]
			\item $\sqrt{2}\cos \left( \alpha -\dfrac{\pi }{4} \right)-\sin \alpha =\cos \alpha +\sin \alpha -\sin \alpha =\cos \alpha $.\\
			\item 
			\begin{eqnarray*}
				{\sin ^4}\alpha +{\cos ^4}\alpha -\dfrac{1}{4}\cos 4\alpha &=&{{\left( {\sin ^2}\alpha +{\cos ^2}\alpha \right)}^2}-2{\sin ^2}\alpha \cdot {\cos ^2}\alpha -\dfrac{1}{4}\cos 4\alpha \\
				&=& 1-\dfrac{1}{2}{\sin ^2}2\alpha -\dfrac{1}{4}\cos 4\alpha\\
				&=& 1-\dfrac{1}{4}\left( 1-\cos 4\alpha \right)-\dfrac{1}{4}\cos 4\alpha\\
				&=&\dfrac{3}{4}+\dfrac{1}{4}\cos 4\alpha -\dfrac{1}{4}\cos 4\alpha =\dfrac{3}{4}.
			\end{eqnarray*}
		\end{enumerate}
	}
\end{ex}

% #3108 | Archive/SP_KhoLatex/ID6_old/1.3.1.tex | mã gốc 1D1B3-4
\begin{ex}%[1D1B3-4]%::Cau 5::
	
	Tính giá trị lượng giác của góc $\alpha $, biết
	\begin{multicols}{2}
		\begin{enumerate}[a)]
		\item  $\cos 2\alpha =\dfrac{-3}{5}$ và $-\pi <\alpha <\dfrac{-\pi }{2}$;
		\item $\sin 2\alpha =\dfrac{-4}{5},\,\dfrac{\pi }{2}<\alpha <\dfrac{3\pi }{4}$.	
		\end{enumerate}
	\end{multicols}
	\loigiai{
		\begin{enumerate}[a)]
			\item 
			$-\pi <\alpha <\dfrac{-\pi }{2}\Rightarrow \sin \alpha <0,\cos \alpha <0$\\
			$\Rightarrow $ $\cos 2\alpha =1-2{\sin ^2}\alpha =\dfrac{-3}{5}\Rightarrow {\sin ^2}\alpha =\dfrac{4}{5}\Rightarrow \sin \alpha =\dfrac{-2\sqrt{5}}{5}$.\\
			$\cos \alpha =\dfrac{-\sqrt{5}}{5}$.\\
			$\tan \alpha =\dfrac{1}{2}$, $\cot \alpha =2$.
			\item 
			$\dfrac{\pi }{2}<\alpha <\dfrac{3\pi }{4}\Rightarrow \,\pi <2\alpha <\dfrac{3\pi }{2}\Rightarrow \cos 2\alpha <0$.\\
			$\cos 2\alpha =\dfrac{-3}{5}\Rightarrow 2{\cos ^2\alpha}-1=\dfrac{-3}{5}\Rightarrow \cos \alpha =\dfrac{-\sqrt{5}}{5}$.\\
			$\sin 2\alpha =2\sin \alpha .\cos \alpha =\dfrac{-4}{5}\Rightarrow \sin \alpha =\dfrac{2\sqrt{5}}{5}$.
	\end{enumerate}}
\end{ex}

% #3109 | Archive/SP_KhoLatex/ID6_old/1.3.1.tex | mã gốc 1D1B3-4
\begin{ex}%[1D1B3-4]%::Cau 6::
	
	Chứng minh rằng, trong tam giác $ABC$, ta có $\cos \dfrac{A}{2}=\sin \dfrac{B}{2}\cos \dfrac{C}{2}+\cos \dfrac{B}{2}\sin \dfrac{C}{2}$.
	\loigiai{
		Ta có \\
		$\sin \left( \dfrac{B+C}{2} \right)=\sin \left( \dfrac{\pi }{2}-\dfrac{A}{2} \right)=\cos \dfrac{A}{2}\Rightarrow \cos \dfrac{A}{2}=\sin \dfrac{B}{2}\cos \dfrac{C}{2}+\cos \dfrac{B}{2}\sin \dfrac{C}{2}$.}
\end{ex}

% #3110 | Archive/SP_KhoLatex/ID6_old/1.3.1.tex | mã gốc 1D1B3-4
\begin{ex}%[1D1B3-4]%::Cau 7::
	\immini{
		Cho tam giác $ABC$ vuông tại $A$ có $AB=4,AC=1$. Trên tia đối của tia $CA$ lấy điểm $D$ sao cho $\widehat{CBD}=45^\circ $. Tính tan $\widehat{ABD}$ và độ dài cạnh $AD$.}
	{\begin{tikzpicture}[line join = round, line cap = round, scale=.6] 
			\coordinate (A) at (0,0);
			\coordinate (B) at (0,6);	
			\coordinate (C) at (1.5,0);
			\coordinate (D) at (8,0); 
			\draw (A)--(B)--(D)--cycle (B)--(C); 
			\foreach \x/\g in {A/180,B/90,C/-90,D/0}\fill[black](\x) circle (2pt)($(\x)+(\g:3.5mm)$) node{$\x$};
		\end{tikzpicture} 
	}
	\loigiai{
		Ta có $\tan \widehat{ABC}=\dfrac{AC}{AB}=\dfrac{1}{4}$.\\
		$\tan \widehat{ABD}=\tan \left( \widehat{ABC}+\widehat{CBD} \right)=\dfrac{\tan \widehat{ABC}+\tan \widehat{CBD}}{1-\tan \widehat{ABC}\cdot\tan \widehat{CBD}}=\dfrac{\dfrac{1}{4}+\tan 45^\circ}{1-\dfrac{1}{4}\cdot\tan 45^\circ}=\dfrac{5}{3}$.\\
		$\tan \widehat{ABD}=\dfrac{AD}{AB}\Rightarrow AD=AB\cdot\tan \widehat{ABD}=4\cdot\dfrac{5}{3}=\dfrac{20}{3}$.}
\end{ex}

% #3212 | Archive/SP_KhoLatex/ID6_old/DeThamKhao-02.tex | mã gốc 1D1B2-1
\begin{ex}%[Dự án ĐGNL-ĐGTD-2022 GĐ1]%[Lê Minh Thiện Anh]%[1D1B2-1]
	Phương trình $2\sin 2x+\sqrt 2=0$ có bao nhiêu 
	nghiệm thuộc khoảng $(0 ;\pi)$?\\
	\textit{Kết quả:} \tikz\draw (0,0) rectangle (3,0.6);
	\loigiai
	{	
		$$2\sin 2x+\sqrt 2=0 \Leftrightarrow \sin 2x=-\dfrac{\sqrt{2}}{2} \Leftrightarrow \hoac{& x=\dfrac{7\pi}{8} +k\pi , k\in\mathbb{Z} \\ & x=\dfrac{5\pi}{8}+k\pi , k\in\mathbb{Z}.}$$
		Với $x=\dfrac{7\pi}{8} +k\pi$ ta có $0\le\dfrac{7\pi}{8} +k\pi \le \pi \Leftrightarrow -\dfrac{7}{8} \le k \le \dfrac{1}{8}$ suy ra có $1$ giá trị $k$ thỏa mãn.\\
		Với $x=\dfrac{5\pi}{8} +k\pi$ ta có $0\le\dfrac{5\pi}{8} +k\pi \le \pi \Leftrightarrow -\dfrac{5}{8} \le k \le \dfrac{3}{8}$ suy ra có $1$ giá trị $k$ thỏa mãn.\\ 
		Vậy có phương trình đã cho có $2$ nghiệm thuộc khoảng $(0;\pi)$.
	}
\end{ex}

% #3215 | Archive/SP_KhoLatex/ID6_old/DeThamKhao-02.tex | mã gốc 1D1B2-1
\begin{ex}%[Dự án ĐGNL-ĐGTD-2022 GĐ1]%[Lê Minh Thiện Anh]%[1D1B2-1]
	Phương trình $\sqrt{3}\cos x+\sin x=-2$ có bao nhiêu nghiệm thuộc đoạn $[0;4035\pi]$?\\
	\textit{Kết quả:} \tikz\draw (0,0) rectangle (3,0.6);
	\loigiai{
		Ta có
		\begin{eqnarray*}
			\sqrt{3}\cos x+\sin x=-2&\Leftrightarrow& \dfrac{\sqrt{3}}{2}\cos x+\dfrac{1}{2}\sin x=-1\\
			&\Leftrightarrow& \sin\left(x+\dfrac{\pi}{3}\right)=-1\\
			&\Leftrightarrow& x+\dfrac{\pi}{3}=-\dfrac{\pi}{2}+k2\pi\\
			&\Leftrightarrow& x=-\dfrac{5\pi}{6}+k2\pi, k\in \mathbb{Z}.
		\end{eqnarray*}
		Trên mỗi nửa khoảng $\left[k2\pi;(k+1)2\pi\right)$, phương trình đã cho có đúng $1$	nghiệm. Trên $[4034\pi; 4035\pi]$ phương trình vô nghiệm.\\
		Do đó phương trình đã cho có $\dfrac{4034}{2}=2017$ nghiệm.
	}
\end{ex}

% #3271 | Archive/SP_KhoLatex/ID6_old/GK1.2.tex | mã gốc 1D1K5-6
\begin{bt}%[1D1K5-6]
	(1,0 điểm) Hằng ngày mực nước của con kênh lên xuống theo thủy triều. Độ sâu $h$ (mét) của mực nước trong kênh được tính tại thời điểm $t$ (giờ) trong một ngày ($t>0$) bởi công thức $h=4\sin (\dfrac{\pi t}{8}+\dfrac{5\pi }{8})+16$. Mực nước của kênh cao nhất khi $t$ bằng bao nhiêu?
	\loigiai{
		Mực nước của kênh cao nhất khi $h$ lớn nhất
		$$ \sin (\dfrac{\pi t}{8}+\dfrac{5\pi }{8})=1\Leftrightarrow \dfrac{\pi }{8}t+\dfrac{5\pi }{8}=\dfrac{\pi }{2}+k2\pi \Leftrightarrow t=-1+16k, k\in\mathbb{Z}.$$
		Với $0<t\le 24$ suy ra $0<-1+16k\le 24\Leftrightarrow \dfrac{1}{16}<k\le \dfrac{25}{16}$.\\
		Do $k\in \mathbb{Z}$ nên $k=1$ thỏa mãn. Khi $k=1$ thì $t=15h$.}
\end{bt}

% #3274 | Archive/SP_KhoLatex/ID6_old/1.3.2.tex | mã gốc 1D1B3-1
\begin{ex}%[1D1B3-1]%::Cau 10::
	
	Chứng minh rằng trong mọi tam giác $ABC$, ta có $\cos A=\sin B\sin C-\cos B\cos C$.
	\loigiai{
		Ta có \[\cos A=\cos \left( \pi -\left( B+C \right) \right)=-\cos \left( B+C \right)=-\left( \cos B\cos C-\sin B\sin C \right)=\sin B\sin C-\cos B\cos C.\]
		}
\end{ex}

% #3276 | Archive/SP_KhoLatex/ID6_old/1.3.2.tex | mã gốc 1D1K3-1
\begin{ex}%[1D1K3-1]%::Cau 12::
	\immini{
	Cho tam giác $ABC$ vuông tại $B$ có $AB=6; BC=8$. Trên tia đối của tia $CB$ lấy điểm $D$ sao cho $CD=8$. Tính $\tan \widehat{CAD}$ ?}
	{\begin{tikzpicture}[line join = round, line cap = round, scale=.6] 
			\coordinate (B) at (0,0);
			\coordinate (A) at (0,4);	
			\coordinate (C) at (3,0);
			\coordinate (D) at (6,0); 
			\draw (A)--(B)--(D)--cycle (A)--(C); 
			\foreach \x/\g in {A/90,B/-90,C/-90,D/-90}\fill[black](\x) circle (2pt)($(\x)+(\g:3.5mm)$) node{$\x$};
	\end{tikzpicture}}
	\loigiai{
		Trong tam giác $ABC$ ta có $\tan \widehat{BAC}=\dfrac{BC}{AB}=\dfrac{8}{6}=\dfrac{4}{3}$.\\
		Trong tam giác $ABD$ ta có $\tan \widehat{BAD}=\dfrac{BD}{AB}=\dfrac{8+8}{6}=\dfrac{8}{3}$.\\
		Do đó \[\tan \widehat{CAD}=\tan \left( \widehat{BAD}-\widehat{BAC} \right)=\dfrac{\tan \widehat{BAD}-\tan \widehat{BAC}}{1+\tan \widehat{BAD}\cdot\tan \widehat{BAC}}=\dfrac{\dfrac{8}{3}-\dfrac{4}{3}}{1+\dfrac{8}{3}\cdot\dfrac{4}{3}}=\dfrac{12}{41}.\]	
		}
\end{ex}

% #3277 | Archive/SP_KhoLatex/ID6_old/1.3.2.tex | mã gốc 1D1K3-1
\begin{ex}%[1D1K3-1]%::Cau 13::
	\immini{
		Cho tam giác $AHO$ vuông tại $H$ có $AH=14; HO=15$. Trên cạnh $AH$ lấy điểm $B$ sao cho $HB=12$. Tính $\tan \widehat{AOB}$ ?}
	{\begin{tikzpicture}[line join = round, line cap = round, scale=.6] 
			\coordinate (H) at (0,0);
			\coordinate (A) at (0,4);	
			\coordinate (B) at (0,3);
			\coordinate (O) at (6,0); 
			\draw (H)--(A)--(O)--cycle (O)--(B); 
			\foreach \x/\g in {H/-90,A/90,O/-90,B/180}\fill[black](\x) circle (2pt)($(\x)+(\g:3.5mm)$) node{$\x$};
	\end{tikzpicture}}
	\loigiai{
		Xét tam giác $AOH$ vuông tại $H$, ta có $\tan \widehat{AOH}=\dfrac{AH}{HO}=\dfrac{14}{15}$.\\
		Xét tam giác $BOH$ vuông tại $H$, ta có $\tan \widehat{BOH}=\dfrac{BH}{HO}=\dfrac{12}{15}=\dfrac{4}{5}$.\\
		Do đó 
		\[\tan \widehat{AOB}=\tan \left( \widehat{AOH}-\widehat{BOH} \right)=\dfrac{\tan \widehat{AOH}-\tan \widehat{BOH}}{1+\tan \widehat{AOH}\cdot\tan \widehat{BOH}}=\dfrac{\dfrac{14}{15}-\dfrac{4}{5}}{1+\dfrac{14}{15}\cdot\dfrac{4}{5}}=\dfrac{\dfrac{2}{15}}{\dfrac{131}{75}}=\dfrac{10}{131}.\]
		Vậy $\tan\widehat{AOB}=\dfrac{10}{131}$.	
		}
\end{ex}

% #3278 | Archive/SP_KhoLatex/ID6_old/1.3.2.tex | mã gốc 1D1K3-1
\begin{ex}%[1D1K3-1]%::Cau 14::
	
	Trong Vật lí, phương trình tổng quát của một vật dao động điều hòa cho bởi công thức\\ $x\left( t \right)=A\cos \left( \omega t+\varphi \right)$, trong đó $t$ là thời điểm (tính bằng giây), $x\left( t \right)$ là li độ của vật tại thời điểm $t,A$ là biên độ dao động $\left( A>0 \right)$ và $\varphi \in \left[ -\pi ;\pi \right]$ là pha ban đầu của dao động. Xét hai dao động điều hòa có phương trình ${x_1}\left( t \right)=5.\cos \left( \dfrac{\pi }{4}t+\dfrac{\pi }{3} \right)$ và ${x_2}\left( t \right)=5.\cos \left( \dfrac{\pi }{4}t-\dfrac{\pi }{6} \right)$.Tìm dao động tổng hợp $x\left( t \right)=x_1\left( t \right)+x_2\left( t \right)$ và sử dụng công thức biến đổi tổng thành tích để tìm biên độ và pha ban đầu của dao động tổng hợp này
	\loigiai{
		Dao động tổng hợp $x\left( t \right)=x_1\left( \text{t} \right)+x_2\left( t \right)$.\\
		Suy ra $x\left( \text{t} \right)=5\cos \left( \dfrac{\pi }{4}t+\dfrac{\pi }{3} \right)+5\cos \left( \dfrac{\pi }{4}t-\dfrac{\pi }{6} \right)$ (cm).\\
		Ta có 
		\begin{eqnarray*}
			5\cos \left( \dfrac{\pi }{4}t+\dfrac{\pi }{3} \right)+5\cos \left( \dfrac{\pi }{4}t-\dfrac{\pi }{6} \right)
			&=&5\left[ \cos \left( \dfrac{\pi }{4}t+\dfrac{\pi }{3} \right)+\cos \left( \dfrac{\pi }{4}t-\dfrac{\pi }{6} \right) \right]\\
			&=&5\cdot2\cos \dfrac{\left( \dfrac{\pi }{4}t+\dfrac{\pi }{3} \right)+\left( \dfrac{\pi }{4}t-\dfrac{\pi }{6} \right)}{2}\cos \dfrac{\left( \dfrac{\pi }{4}t+\dfrac{\pi }{3} \right)-\left( \dfrac{\pi }{4}t-\dfrac{\pi }{6} \right)}{2}\\
			&=& 10\cos \left( \dfrac{\pi }{4}t+\dfrac{\pi }{12} \right)\cos \dfrac{\pi }{4}\\
			&=& 10\cos \left( \dfrac{\pi }{4}t+\dfrac{\pi }{12} \right)\cdot \dfrac{\sqrt{2}}{2}\\
			&=& 5\sqrt{2}\cos \left( \dfrac{\pi }{4}t+\dfrac{\pi }{12} \right).
		\end{eqnarray*}
		Vậy dao động tổng hợp có phương trình là $x\left( t \right)=2\sqrt{2}\cos \left( \dfrac{\pi }{6}t+\dfrac{\pi }{12} \right)$ với biên độ $A=5\sqrt{2}$ và pha ban đầu là $\varphi =\dfrac{\pi }{12}$.\\
	}
\end{ex}

% #3279 | Archive/SP_KhoLatex/ID6_old/DeThamKhao-03.tex | mã gốc 1D1K2-1
\begin{ex}%[Dự án ĐGNL-ĐGTD-2022 GĐ1]%[Duong Xuan Loi]%[1D1K2-1]
	Phương trình $\cos 2x +\cos x=0$ có bao nhiêu nghiệm thuộc khoảng $\left(-\pi; \pi\right)$?
	
	\textit{Kết quả:} \tikz\draw (0,0) rectangle (3,0.6);
	\loigiai{
		Ta có 
		\begin{eqnarray*}
			\cos 2x +\cos x=0 &\Leftrightarrow& \cos 2x =\cos (\pi-x)\\
			&\Leftrightarrow& 
			\hoac{&2x=\pi -x+k 2\pi\\&2x=-\pi +x+k2\pi} \\
			&\Leftrightarrow& \hoac{&x=\dfrac{\pi}{3}+\dfrac{k2\pi}{3}\\&x=-\pi+k2\pi}
		\end{eqnarray*}
		Với $ x=\dfrac{\pi}{3}+\dfrac{k2\pi}{3}$ ta có 
		\[-\pi <x<\pi \Leftrightarrow -\pi <\dfrac{\pi}{3}+\dfrac{k2\pi}{3}<\pi\Leftrightarrow -2<k<1\]
		Mà $ k \in\mathbb{Z}$ nên $k=-1;k=0$ hay $x= -\dfrac{\pi}{3}$, $x=\dfrac {\pi}{3}$.\\
		Với $x= -\pi +k2 \pi $ ta có 
		\[-\pi <x<\pi \Leftrightarrow -\pi <-\pi+k 2\pi <\pi\Leftrightarrow 0<k<1\]
		Mà $ k \in \mathbb{Z}$ nên không tồn tại $k$.\\
		Vậy phương trình đã cho  có $2$ nghiệm thuộc $\left(-\pi ; \pi\right)$.
	}
\end{ex}

% #3294 | Archive/SP_KhoLatex/ID6_old/DeThamKhao-07.tex | mã gốc 1D1K3-3
\begin{ex}%[Dự án ĐGNL-ĐGTD-2022 GĐ1]%[Dao-V- Thuy]%[1D1K3-3]
	Số nghiệm của phương trình $\sin 2x=2\cos x\left(\sqrt{3}\cos x+2 \right)$ trên đoạn $[0; 4\pi]$ là\\
	\textit{Kết quả:} \tikz\draw (0,0) rectangle (3,0.6);
	\loigiai{
		Ta có$\colon$ 
\begin{eqnarray*}
			& & \sin 2x=2\cos x\left(\sqrt{3}\cos x+2 \right)\\
			& \Leftrightarrow & 2\sin x\cos x-2\cos x\left(\sqrt{3}\cos x+2\right)=0\\
			& \Leftrightarrow & 2\cos x\left(\sin x-\sqrt{3}\cos x-2\right)=0\\
			& \Leftrightarrow & \hoac{& \cos x=0	\quad &(1)\\& \sin x-\sqrt{3}\cos x-2=0. \quad &(2)}
\end{eqnarray*}
		\noindent
		Phương trình $(1) \Leftrightarrow x=\dfrac{\pi}{2}+k\pi \quad (k \in \mathbb{Z})$.\\
		Phương trình $(2) \Leftrightarrow \sin x-\sqrt{3}\cos x=2 \Leftrightarrow \sin \left(x-\dfrac{\pi}{3}\right)=1 \Leftrightarrow x=\dfrac{5\pi}{6}+k2\pi 
		\quad (k \in \mathbb{Z})$.\\
		Vậy phương trình đã cho có các nghiệm là $x=\dfrac{\pi}{2}+k\pi$, $x=\dfrac{5\pi}{6}+k2\pi$ $\quad (k \in \mathbb{Z})$.
\begin{itemize}
			\item Với $x=\dfrac{\pi}{2}+k\pi \; (k \in \mathbb{Z})$ ta có $0 \le x \le 4\pi \Leftrightarrow 0 \le \dfrac{\pi}{2}+k\pi \le 4\pi \Leftrightarrow -\dfrac{1}{2} \le k \le \dfrac{7}{2}$.\\[6pt]
			Vì $k \in \mathbb{Z}$ nên $k \in \{0; 1; 2; 3\} \Rightarrow$ có 4 nghiệm của phương trình $(1)$ là $x\in \left\{\dfrac{\pi}{2}; \dfrac{3\pi}{2}; \dfrac{5\pi}{2}; \dfrac{7\pi}{2}\right\}$.
			\item Với $x=\dfrac{5\pi}{6}+k2\pi\; (k \in \mathbb{Z})$ ta có $0 \le x \le 4\pi \Leftrightarrow 0 \le \dfrac{5\pi}{6}+k2\pi \le 4\pi \Leftrightarrow -\dfrac{5}{12} \le k \le \dfrac{19}{12}$.\\[6pt]
			Vì $k \in \mathbb{Z}$ nên $k \in \{0; 1\} \Rightarrow$ có 2 nghiệm của phương trình $(2)$ là $x\in \left\{\dfrac{5\pi}{6}; \dfrac{17\pi}{6}\right\}$.
\end{itemize}
		Vậy phương trình đã cho có 6 nghiệm thuộc khoảng $[0; 4\pi]$.
	}
\end{ex}

% #3330 | Archive/SP_KhoLatex/ID6_old/GK1.6.tex | mã gốc 1D1B5-3
\begin{bt}%[1D1B5-3]
	Giải phương trình $2\sin x-1=0$.
	\loigiai{
	Ta có:
	\begin{eqnarray*}
		&& 2\sin x-1=0\Leftrightarrow 2\sin x=1 \\
		& \Leftrightarrow& \sin x=\dfrac{1}{2}=\sin \dfrac{\pi}{6}\\
		&\Leftrightarrow& \hoac{& x=\dfrac{\pi}{6}+k2\pi \\
		& x=\pi-\dfrac{\pi}{6}+k2\pi} \\
		&\Leftrightarrow& \hoac{& x=\dfrac{\pi}{6}+k2\pi \\ & x=\dfrac{5\pi}{6}+k2\pi}(k\in \mathbb{Z}).
	\end{eqnarray*}
	}
\end{bt}

% #3332 | Archive/SP_KhoLatex/ID6_old/GK1.6.tex | mã gốc 1D1G3-5
\begin{bt}%[1D1G3-5]
	Cho tam giác $ABC,AB=c,AC=b,BC=a$. Chứng minh rằng nếu $\dfrac{{b^3}+{c^3}-{a^3}}{b+c-a}={a^2}$
	và $\cos (A+C)+3\cos B=1$ thì $\triangle ABC$ đều.
	\loigiai{
		Ta có:
		\begin{eqnarray*}
			&&\dfrac{{b^3}+{c^3}-{a^3}}{b+c-a}={a^2}\\
			&\Leftrightarrow& {b^3}+{c^3}-{a^3}={a^2}( b+c-a )\\
			&\Leftrightarrow& {b^3}+{c^3}={a^2}( b+c )\\
			&\Leftrightarrow& ( b+c )( {b^2}-bc+{c^2} )={a^2}( b+c )\\
			&\Leftrightarrow& {b^2}+{c^2}-{a^2}-bc=0\, (\text{do}\,b+c\ne 0 )\\
			&\Leftrightarrow& 2bc\cdot \cos A-bc=0\\
			&\Leftrightarrow& \cos A=\dfrac{1}{2}\\
			&\Leftrightarrow& \widehat{A}={60^\circ}.
		\end{eqnarray*}
		Ta có $A+B+C={180^\circ}\Rightarrow \cos B=\cos [ {180^\circ}-( A+C ) ]=-\cos( A+C )$.\\
		Suy ra $ \cos( A+C )+3\cos B=1\Leftrightarrow -\cos B+3\cos B=1\Leftrightarrow \cos B=\dfrac{1}{2}\Leftrightarrow \widehat{B}={60^\circ}$.\\
		Do đó $\widehat{C}={60^\circ}$.\\
		Vậy tam giác$\triangle ABC$ đều.
	}
\end{bt}

% #3341 | Archive/SP_KhoLatex/ID6_old/1-GK1-KNTT-De14-NH23-24.tex | mã gốc 1D1B2-1
\begin{bt}%[1D1B2-1]% ::Cau 31::
	Giải các phương trình sau:
	\begin{enumEX}{2}
		\item $\sin \left( 4x+\dfrac{\pi}{4} \right)=\sin \left( x-\dfrac{\pi}{5} \right)$; 
		\item $\cot \left( 2x-40^{\circ} \right)=-\sqrt{3}$.
	\end{enumEX}
	\loigiai{
		\begin{enumEX}{1}
		\item Ta có 
		\allowdisplaybreaks
		\begin{eqnarray*}
			&& \sin \left( 4x+\dfrac{\pi}{4} \right)=\sin \left( x-\dfrac{\pi}{5} \right)\\
			&\Leftrightarrow& \hoac{
				& 4x+\dfrac{\pi}{4}=x-\dfrac{\pi}{5}+k2\pi\\
				& 4x+\dfrac{\pi}{4}=\pi-x+\dfrac{\pi}{5}+k2\pi}\\
			&\Leftrightarrow& \hoac{
				& 3x=-\dfrac{9\pi}{20}+k2\pi\\
				& 5x=\dfrac{19\pi}{20}+k2\pi}\\
			&\Leftrightarrow& \hoac{
				& x=-\dfrac{3\pi}{20}+k\dfrac{2\pi}{3} \\
				& x=\dfrac{19\pi}{100}+k\dfrac{2\pi}{5}},\,\left( k\in \mathbb{Z} \right).
		\end{eqnarray*}
		\item Ta có $\cot \left( 2x-40^{\circ} \right)=-\sqrt{3}\Leftrightarrow 2x-40^{\circ} =-30^{\circ} +k180^{\circ} \Leftrightarrow x=5^{\circ} +k90^{\circ} ,\,\left( k\in \mathbb{Z} \right)$.
	\end{enumEX}
	}
\end{bt}

% #3342 | Archive/SP_KhoLatex/ID6_old/1-GK1-KNTT-De14-NH23-24.tex | mã gốc 1D1K3-8
\begin{bt}%[1D1K3-8]% ::Cau 32::
	Số giờ có ánh sáng của một thành phố $X$ ở vĩ độ $40^{\circ} $ Bắc trong ngày thứ $t$ của một năm không nhuận được cho bởi hàm số $d\left( t \right)=3\sin \left[ \dfrac{\pi}{182}\left( t-80 \right) \right]+12$, $t\in \mathbb{Z}$ và $0<t\le 365$. Vào ngày nào trong năm thì thành phố $X$ có nhiều giờ ánh sáng nhất?
	\loigiai{
		Ta có $d\left( t \right)=3\sin \left[ \dfrac{\pi}{182}\left( t-80 \right) \right]+12\le 3+12=15$.\\
		Dấu bằng xảy ra khi $$\sin \left[ \dfrac{\pi}{182}\left( t-80 \right) \right]=1\Leftrightarrow \dfrac{\pi}{182}\left( t-80 \right)=\dfrac{\pi}{2}+k2\pi\,\,\left( k\in \mathbb{Z} \right) \Leftrightarrow t=171+364k.$$
		Mặt khác $t\in \left( 0;365 \right]$ nên $0<171+364k\le 365\Leftrightarrow -\dfrac{171}{364}<k\le \dfrac{194}{364}$.\\
		Mà $k\in \mathbb{Z}$ nên $k=0$.\\
		Vậy $t=171$.}
\end{bt}

% #3343 | Archive/SP_KhoLatex/ID6_old/1-GK1-KNTT-De14-NH23-24.tex | mã gốc 1D1K3-8
\begin{bt}%[1D1K3-8]% ::Cau 33::
	Hàng ngày mực nước của một con kênh lên xuống theo thủy triều. Độ sâu $h$ (mét) của mực nước trong kênh tính theo thời gian $t$ (giờ) $\left( 0\le t\le 24 \right)$ được mô tả bởi công thức $h=A\cos \left( \dfrac{\pi t}{6}+1 \right)+B$, với $A, B$ là các số thực dương cho trước. Biết độ sâu của mực nước lớn nhất là $15$ mét khi thủy triều lên cao và khi thủy triều xuống thấp thì độ sâu của mực nước thấp nhất là $9$ mét. Tính thời điểm độ sâu của mực nước là $13{,}5$ mét (tính chính xác đến $\dfrac{1}{100}$ giờ).
	\loigiai{
		Với mọi $0\le t\le 24$, ta có
		\allowdisplaybreaks
		\begin{eqnarray*}
		&& -1\le \cos \left( \dfrac{\pi t}{6}+1 \right)\le 1 \\
		&\Leftrightarrow& -A+B\le A\cos \left( \dfrac{\pi t}{6}+1 \right)+B\le A+B.
		\end{eqnarray*}
		Độ sâu của mực nước lớn nhất bằng $A+B$ khi $\cos \left( \dfrac{\pi t}{6}+1 \right)=1$ và thấp nhất bằng $-A+B$ khi $\cos \left( \dfrac{\pi t}{6}+1 \right)=-1$.\\
		Ta có hệ $\heva{
			& A+B=15 \\
			& -A+B=9}
		\Leftrightarrow \heva{
			& B=12 \\
			& A=3.}$\\
		Ta được $h=3\cos \left( \dfrac{\pi t}{6}+1 \right)+12$.\\
		Theo đề, ta tìm thời điểm mà độ sâu 
		\allowdisplaybreaks
		\begin{eqnarray*}
		&& h=13{,}5\Leftrightarrow 3\cos \left( \dfrac{\pi t}{6}+1 \right)+12=13{,}5\Leftrightarrow \cos \left( \dfrac{\pi t}{6}+1 \right)=\dfrac{1}{2}\\
		&\Leftrightarrow& \hoac{
			& \dfrac{\pi t}{6}+1=\dfrac{\pi}{3}+k2\pi\\
			& \dfrac{\pi t}{6}+1=-\dfrac{\pi}{3}+k2\pi},\left( k\in \mathbb{Z} \right)
		\Leftrightarrow \heva{
			& t=\left( -1+\dfrac{\pi}{3} \right)\cdot \dfrac{6}{\pi}+12k \\
			& t=\left( -1-\dfrac{\pi}{3} \right)\cdot \dfrac{6}{\pi}+12k},\left( k\in \mathbb{Z} \right).
		\end{eqnarray*}
		Do $0\le t\le 24; k\in \mathbb{Z}$ nên $t=0{,}09$ (giờ); $t=12{,}09$ (giờ); $t=8{,}09$ (giờ); $t=20{,}09$ (giờ).}
\end{bt}

% #3446 | Archive/SP_KhoLatex/ID6_old/GK1.5.tex | mã gốc 1D1K3-4
\begin{bt}%[1D1K3-4]
Cho $\triangle ABC$ có $ M$ là trung điểm cạnh $ BC$ và $\dfrac{\sin B+\sin C}{\cos B+\cos C}=\sin A$. Chứng minh rằng $ BC=2AM$.
\loigiai{
	Ta có $A+B+C=180^\circ  \Rightarrow \dfrac{A+B+C}{2}=90^\circ \Rightarrow \dfrac{B+C}{2}=90^\circ -\dfrac{A}{2}$.\\
	Khi đó
	\allowdisplaybreaks \vspace*{-0.7cm}
	\begin{eqnarray*} 
		& & \dfrac{\sin B+\sin C}{\cos B+\cos C}=\sin A \\
		& \Leftrightarrow & \dfrac{2\sin\dfrac{B+C}{2}\cos\dfrac{B-C}{2}}{2\cos\dfrac{B+C}{2}\cos\dfrac{B-C}{2}}=\sin A\\
		& \Leftrightarrow & \dfrac{\sin\dfrac{B+C}{2}}{\cos\dfrac{B+C}{2}}=2\sin\dfrac{A}{2}\cos\dfrac{A}{2}\\
		& \Leftrightarrow &\dfrac{\cos\dfrac{A}{2}}{\sin\dfrac{A}{2}}=2\sin\dfrac{A}{2}\cos\dfrac{A}{2}\\
		& \Leftrightarrow & \dfrac{1}{\sin\dfrac{A}{2}}=2\sin\dfrac{A}{2}\\
		& \Leftrightarrow & 1-2\sin^2\dfrac{A}{2}=0\\ 
		& \Leftrightarrow & \cos A=0 \\
		& \Leftrightarrow & A=90^\circ.
	\end{eqnarray*}
	Suy ra $\triangle ABC$ vuông tại $ A$.\\
	Do đó trung tuyến $ AM$ ứng với cạnh huyền $ BC$ nên $BC=2AM$.
	}
\end{bt}

% #3452 | Archive/SP_KhoLatex/ID6_old/GK1.4.tex | mã gốc 1D1K3-5
\begin{ex}%[1D1K3-5]
	Chứng minh trong mọi tam giác ABC ta đều có $\cos A+\cos B+\cos C\le \dfrac{3}{2}$.
	\loigiai{
		Ta có $\cos A+\cos B+\cos C=2\cos \dfrac{A+B}{2} \cos \dfrac{A-B}{2}+\cos C$.\\		
		Vì $\dfrac{A+B}{2}=\dfrac{\pi}{2}-\dfrac{C}{2}$ nên $\cos \dfrac{A+B}{2}=\sin \dfrac{C}{2}$.\\		
		Mặt khác $\cos C=1-2\sin^2 \dfrac{C}{2}$ do đó		
		\begin{eqnarray*}
			\cos A+\cos B+\cos C
			& =& 2\sin \dfrac{C}{2} \cos \dfrac{A-B}{2}+1-2\sin^2 \dfrac{C}{2}\\
			&= & -2\left( \sin^2 \dfrac{C}{2}-\sin \dfrac{C}{2} \cos \dfrac{A-B}{2}-\dfrac{1}{2} \right)\\\\
			&= & -2\left( \sin^2 \dfrac{C}{2}-2\sin \dfrac{C}{2}\cdot\dfrac{1}{2} \cos \dfrac{A-B}{2}+\dfrac{1}{4} \cos^2 \dfrac{A-B}{2} \right)+1+\dfrac{1}{2} \cos^2 \dfrac{A-B}{2} \\\\
			&= &-2(\sin \dfrac{C}{2}+\dfrac{1}{2} \cos \dfrac{A-B}{2} )^2+1+\dfrac{1}{2} \cos^2 \dfrac{A-B}{2}. 
		\end{eqnarray*}
		Vì $\left|\cos \dfrac{A-B}{2} \right|\le 1\Rightarrow \cos^2 \dfrac{A-B}{2} \le 1$ nên 	
		$\cos A+\cos B+\cos C\le 1+\dfrac{1}{2}=\dfrac{3}{2}$.\\
		Vậy yêu cầu bài toán đã được chứng minh.
		
	}
	\end{ex}

% #3509 | Archive/SP_KhoLatex/ID6_old/1-CK1-OT-De-14.tex | mã gốc 1D1H5-3
\begin{bt}%[1D1H5-3] 
	Giải phương trình $\sin\left(2x+\dfrac{\pi}{3}\right)=\dfrac{\sqrt{3}}{2}$.
	\loigiai{
		Ta có
		\begin{eqnarray*}
			&&\sin\left(2x+\dfrac{\pi}{3}\right)=\dfrac{\sqrt{3}}{2}\Leftrightarrow\sin\left(2x+\dfrac{\pi}{3}\right)=\sin\dfrac{\pi}{3} \\
			&\Leftrightarrow&\hoac{&2x+\dfrac{\pi}{3}=\dfrac{\pi}{3}+k2\pi\\&2x+\dfrac{\pi}{3}=\pi-\dfrac{\pi}{3}+k2\pi}\Leftrightarrow\hoac{&2x=k2\pi\\&2x=\dfrac{\pi}{3}+k2\pi}\\
			&\Leftrightarrow&\hoac{&x=k\pi\\&x=\dfrac{\pi}{6}+k\pi}(k\in\mathbb{Z}).
		\end{eqnarray*}
	 }
\end{bt}

% #3545 | Archive/SP_KhoLatex/ID6_old/1K1-3.tex | mã gốc 1D1K1-5
\begin{vd}%[Dự án 1 - TLDH - TeamTeXHoa - Lê Minh Cường]%[1D1K1-5]
	Tìm giá trị lớn nhất, giá trị nhỏ nhất của hàm số
	\begin{enumerate}
		\item $y=\cos^2 x+2\sin x+2$;
		\item $y=\sin^4 x-2\cos^2 x+1$
	\end{enumerate}
	\loigiai{
		\begin{enumerate}
			\item Ta có
			\begin{eqnarray*}
				y&=&\cos^2 x+2\sin x+2\\
				&=& \left(1-\sin^2 x\right)+2\sin x+2\\
				&=& -\sin^2 x+2\sin x+3\\
				&=& -\left(\sin x-1\right)^2+4.
			\end{eqnarray*}
			Vì
			\begin{eqnarray*}
				-1\le \sin x\le 1&\Rightarrow&-2\le \sin x-1\le 0\\
				&\Rightarrow& 4\ge \left(\sin x-1\right)^2\ge 0\\
				&\Rightarrow&-4\le -\left(\sin x-1\right)^2\le 0\\
				&\Rightarrow&0\le -\left(\sin x-1\right)^2+4\le 4.
			\end{eqnarray*}
			Hay $0\le y\le 4$.\\
			Do đó
			\begin{itemize}
				\item $\max y=4$ khi $\sin x=1\Leftrightarrow x=\dfrac{\pi}{2}+k2\pi$, $k\in\mathbb{Z}$.
				\item $\min y=0$ khi $\sin x=-1\Leftrightarrow x=-\dfrac{\pi}{2}+k2\pi$, $k\in\mathbb{Z}$.
			\end{itemize}
			\begin{luuy}
				\immini{
					Nếu đặt $t=\sin x$ thì $t\in [-1;1]$.\\
					Ta có $(P)\colon y=f(t)=-t^2+2t+3$ xác định với mọi $t\in[-1;1]$, $(P)$ có hoành độ đỉnh $t=1$ và trên đoạn $[-1;1]$ hàm số đồng biến nên hàm số đạt giá trị nhỏ nhất tại $t=-1$ hay $\sin x=-1$ và giá trị lớn nhất tại $t=1$ hay $\sin x=1$.
				}{
					\begin{tikzpicture}[line join=round, line cap=round,>=stealth,thick,scale=.7]
						\tikzset{label style/.style={font=\footnotesize}}
						\def \xmin{-2}
						\def \xmax{2}
						\def \ymin{-1}
						\def \ymax{5}
						\def \hamso{-(\x)^2+2*(\x)+3}
						\draw[->] (\xmin,0)--(\xmax,0) node[below left] {$x$};
						\draw[->] (0,\ymin)--(0,\ymax) node[below left] {$y$};
						\draw (0,0) node [below right] {$O$};
						\begin{scope}
							\clip (\xmin+0.01,\ymin+0.01) rectangle (\xmax-0.01,\ymax-0.01);
							\draw[samples=350,domain=\xmin+0.01:\xmax-0.01,smooth,variable=\x] plot (\x,{\hamso});
							\foreach \y in {1,2,3}
							\fill (0,\y) circle (1.5pt) node[right]{$\y$};
							\foreach \x in {-1,1}
							\fill (\x,0) circle (1.5pt) node[below]{$\x$};
						\end{scope}
					\end{tikzpicture}
				}
			\end{luuy}
			\item Ta có
			\begin{eqnarray*}
				y&=&\sin^4 x-2\cos^2 x+1\\
				&=& \left(1-\cos^2 x\right)^2-2\cos^2 x+1\\
				&=& \cos^4 x-4\cos^2 x+2\\
				&=& \left(\cos^2 x-2\right)^2-2.
			\end{eqnarray*}
			Vì
			\begin{eqnarray*}
				0\le \cos^2 x\le 1&\Rightarrow&-2\le \cos^2 x-2\le -1\\
				&\Rightarrow& 4\ge \left(\cos^2 x-2\right)^2\ge 1\\
				&\Rightarrow&2\ge \left(\cos^2 x-2\right)^2-2\ge -1.
			\end{eqnarray*}
			Hay $-1\le y\le 2$.\\
			Do đó
			\begin{itemize}
				\item $\max y=2$ khi $\cos^2 x=0\Leftrightarrow \cos x=0\Leftrightarrow x=\dfrac{\pi}{2}+k\pi$, $k\in\mathbb{Z}$.
				\item $\min y=-1$ khi $\cos^2 x=1\Leftrightarrow \sin x=0\Leftrightarrow x=k\pi$, $k\in\mathbb{Z}$.
			\end{itemize}
			\begin{luuy}
				\immini{
					Nếu đặt $t=\cos^2 x$ thì $t\in [0;1]$.\\
					Ta có $(P)\colon y=f(t)=t^2-4t+2$ xác định với mọi $t\in[0;1]$, $(P)$ có hoành độ đỉnh $t=2$ và trên đoạn $[0;1]$ hàm số nghịch biến nên hàm số đạt giá trị nhỏ nhất tại $t=1$ và giá trị lớn nhất tại $t=0$.
				}{
					\begin{tikzpicture}[line join=round, line cap=round,>=stealth,thick,scale=.7]
						\tikzset{label style/.style={font=\footnotesize}}
						\def \xmin{-1}
						\def \xmax{2}
						\def \ymin{-3}
						\def \ymax{2}
						\def \hamso{(\x)^2-4*(\x)+2}
						\draw[->] (\xmin,0)--(\xmax,0) node[below left] {$x$};
						\draw[->] (0,\ymin)--(0,\ymax) node[below left] {$y$};
						\draw (0,0) node [above left] {$O$};
						\begin{scope}
							\clip (\xmin+0.01,\ymin+0.01) rectangle (\xmax-0.01,\ymax-0.01);
							\draw[samples=350,domain=\xmin+0.01:\xmax-0.01,smooth,variable=\x] plot (\x,{\hamso});
							\foreach \y in {1,-1,-2}
							\fill (0,\y) circle (1.5pt) node[left]{$\y$};
							\foreach \x in {1}
							\fill (\x,0) circle (1.5pt) node[below]{$\x$};
						\end{scope}
					\end{tikzpicture}
				}
			\end{luuy}
		\end{enumerate}
	}
\end{vd}

% #3546 | Archive/SP_KhoLatex/ID6_old/1K1-3.tex | mã gốc 1D1K1-5
\begin{vd}%[Dự án 1 - TLDH - TeamTeXHoa - Lê Minh Cường]%[1D1K1-5]
	Tìm giá trị lớn nhất và giá trị nhỏ nhất của hàm số $y=\dfrac{2\sin x-\cos x+1}{\sin x+\cos x-2}$.
	\loigiai{
		Ta có $\sin x+\cos x-2=\sqrt{2}\sin\left(x+\dfrac{\pi}{4}\right)-2$.\\
		Với mọi $x\in\mathbb{R}$, ta có
		\begin{eqnarray*}
			-1\le \sin\left(x+\dfrac{\pi}{4}\right)\le 1&\Rightarrow&-\sqrt{2}\le \sqrt{2}\sin\left(x+\dfrac{\pi}{4}\right)\le \sqrt{2}\\
			&\Rightarrow& \sqrt{2}\sin\left(x+\dfrac{\pi}{4}\right)-2\le \sqrt{2}-2<0\\
			&\Rightarrow& \sin x+\cos x-2=\sqrt{2}\sin\left(x+\dfrac{\pi}{4}\right)-2\ne 0.
		\end{eqnarray*}
		Do đó $\mathscr{D}=\mathbb{R}$.\\
		Biến đổi 
		\begin{eqnarray*}
			y=\dfrac{2\sin x-\cos x+1}{\sin x+\cos x-2}&\Leftrightarrow&y\sin x+y\cos x-2y=2\sin x-\cos x+1\\
			&\Leftrightarrow& (y-2)\sin x+(y+1)\cos x=2y+1\quad (*).
		\end{eqnarray*}
		$(*)$ có nghiệm thực khi và chỉ khi
		\begin{eqnarray*}
			(y+2)^2+(y+1)^2\ge (2y+1)^2&\Leftrightarrow&2y^2+6y-4\le 0\\
			&\Leftrightarrow& \dfrac{-3-\sqrt{17}}{2}\le y\le \dfrac{-3+\sqrt{17}}{2}.
		\end{eqnarray*}
		Vậy
		\begin{itemize}
			\item $\max y=\dfrac{-3+\sqrt{17}}{2}$.
			\item $\min y=\dfrac{-3-\sqrt{17}}{2}$.
		\end{itemize}
	}
\end{vd}

% #3547 | Archive/SP_KhoLatex/ID6_old/1K1-3.tex | mã gốc 1D1K1-4
\begin{vd}%[Dự án 1 - TLDH - TeamTeXHoa - Lê Minh Cường]%[1D1K1-4]
	Chứng minh rằng các hàm số sau là những hàm số tuần hoàn với chu kỳ cơ sở $T_0$
	\begin{enumerate}
		\item $f(x)=\sin x$, $T_0=2\pi$.
		\item $f(x)=\tan 2x$, $T_0=\dfrac{\pi}{2}$.
	\end{enumerate}
	\loigiai{
		\begin{enumerate}
			\item Ta có $f(x+2\pi)=f(x)$, $\forall x\in\mathbb{R}$.\\
			Giả sử có số thực dương $T<2\pi$ thoả $$f(x+T)=f(x)\Leftrightarrow \sin (x+T)=\sin x,\, \forall x\in\mathbb{R}.\quad (*)$$
			Cho $x=\dfrac{\pi}{2}$, ta có
			\begin{itemize}
				\item VT (*)$=\sin\left(\dfrac{\pi}{2}+T\right)=\cos T<1$.
				\item VP (*)$=\sin\left(\dfrac{\pi}{2}\right)=1$.
			\end{itemize}
			Do đó (*) không xảy ra với mọi số thực $x$.\\
			Vậy hàm số đã cho tuần hoàn với chu kỳ $T_0=2\pi$.
			\item Ta có $f\left(x+\dfrac{\pi}{2}\right)=f(x)$, $\forall x\in\mathscr{D}$.\\
			Giả sử có số thực dương $T<\dfrac{\pi}{2}$ thoả $$f(x+T)=f(x)\Leftrightarrow \tan (2x+2T)=\tan 2x,\, \forall x\in\mathscr{D}. \quad (**)$$
			Cho $x=0$, ta có
			\begin{itemize}
				\item VT (**)$=\tan 2T\ne 0$.
				\item VP (**)$=0$.
			\end{itemize}
			Do đó (**) không xảy ra với mọi $x\in\mathscr{D}$.\\
			Vậy hàm số đã cho tuần hoàn với chu kỳ $T_0=\dfrac{\pi}{2}$.
		\end{enumerate}
	}
\end{vd}

% #3634 | Archive/SP_KhoLatex/ID6_old/1.2.2.tex | mã gốc 1D1B2-4
\begin{ex}%[1D1B2-4]
	Chứng minh các đẳng thức lượng giác sau
	\begin{multicols}{2}
	\begin{enumerate}[a)]
		\item ${\sin ^4}a-{\cos ^4}a=2{\sin ^2}a-1$
		\item $\dfrac{\sin a}{\cot a}-\dfrac{\tan a}{\sin a}=-\cos a$
	\end{enumerate}
	\end{multicols}
	\loigiai{
		\begin{enumerate}[a)]
		\item $\sin ^4 a-\cos ^4 a=2\sin ^2a-1$.
		\allowdisplaybreaks
		\begin{eqnarray*}
		\sin ^4 a-\cos ^4 a &=& \left(\sin ^2 a+\cos ^2 a \right)\left(\sin ^2 a-\cos ^2 a \right)=\sin ^2 a-\cos ^2 a\\
			&=& \sin ^2 a-\left( 1-\sin ^2 a \right)=2\sin ^2a-1.
		\end{eqnarray*}
		\item  $\dfrac{\sin a}{\cot a}-\dfrac{\tan a}{\sin a}=-\cos a$\\
		$\dfrac{\sin a}{\cot a}-\dfrac{\tan a}{\sin a}=\dfrac{{\sin ^2}a}{\cos a}-\dfrac{\sin a}{\sin a\cos a}=\dfrac{{\sin ^3}a-\sin a}{\sin a\cos a}=\dfrac{{\sin ^2}a-1}{\cos a}=-\cos a$~(đpcm).
		\end{enumerate}
		}
\end{ex}

% #3635 | Archive/SP_KhoLatex/ID6_old/1.2.2.tex | mã gốc 1D1B2-4
\begin{ex}%[1D1B2-4]
	 Chứng minh các đẳng thức lượng giác sau
	 	\begin{multicols}{2}
	 	\begin{enumerate}[a)]
	\item $\dfrac{\sin ^2 a-\cos ^2 a}{\tan a-\cot a}=\sin a\cos a$
    \item $\dfrac{\tan a-\cot a}{\tan a+\cot a}=1-2\cos ^2 a$
		\end{enumerate}
\end{multicols}
	\loigiai{
		\begin{enumerate}[a)]
		\item  $\dfrac{\sin ^2 a-\cos ^2a}{\tan a-\cot a}=\sin a\cos a$\\
		$\dfrac{\sin ^2 a-\cos ^2 a}{\tan a-\cot a}=\dfrac{\sin ^2 a-\cos ^2 a}{\dfrac{\sin a}{\cos a}-\dfrac{\cos a}{\sin a}}=\dfrac{{\sin ^2}a-{\cos ^2}a}{\dfrac{\sin ^2 a-\cos ^2 a}{\cos a\cdot \sin a}}=\sin a\cdot \cos a$~(đpcm).
		\item  $\dfrac{\tan a-\cot a}{\tan a+\cot a}=1-2{\cos ^2}a$
		\begin{eqnarray*}
				\dfrac{\tan a-\cot a}{\tan a+\cot a}&=&\dfrac{\dfrac{\sin a}{\cos a}-\dfrac{\cos a}{\sin a}}{\dfrac{\sin a}{\cos a}+\dfrac{\cos a}{\sin a}}=\dfrac{\sin ^2 a-\cos ^2 a}{\sin ^2 a+\cos ^2 a}=\sin ^2 a-\cos ^2 a\\
			&=& 1-\cos ^2 a-\cos ^2 a=1-2\cos ^2 a.
		\end{eqnarray*}
	\end{enumerate}
		}
\end{ex}

% #3636 | Archive/SP_KhoLatex/ID6_old/1.2.2.tex | mã gốc 1D1B2-4
\begin{ex}%[1D1B2-4]
	Rút gọn các biểu thức sau
	\begin{multicols}{2}
		\begin{enumerate}[a)]
	\item $\dfrac{1+\cot \alpha }{1+\tan \alpha }-\cot \alpha $	\\
	\item  $\sin \left( \dfrac{\pi }{2}-\alpha \right)+\cos \left( \pi -\alpha \right)$\\
	\item  $\sin \left( \alpha +7\pi \right)+\cos \left( \alpha -\dfrac{\pi }{2} \right)+\tan \left( 3\pi +\alpha \right).\cot \left( 2\pi -\alpha \right)$
		\end{enumerate}
\end{multicols}
	\loigiai{
        \begin{enumerate}[a)]
		\item $\dfrac{1+\cot \alpha }{1+\tan \alpha }-\cot \alpha =\dfrac{1+\cot \alpha -\cot \alpha \left( 1+\tan \alpha \right)}{1+\tan \alpha }=\dfrac{1+\cot \alpha -\cot \alpha -\cot \alpha \tan \alpha }{1+\tan \alpha }=0$\\
		\item  $\sin \left( \dfrac{\pi }{2}-\alpha \right)+\cos \left( \pi -\alpha \right)=\cos \alpha -\cos \alpha =0$.\\
		\item  $\sin \left( \alpha +7\pi \right)+\cos \left( \alpha -\dfrac{\pi }{2} \right)+\tan \left( 3\pi +\alpha \right).\cot \left( 2\pi -\alpha \right)
		=-\sin \alpha +\sin \alpha +\tan \alpha \left( -\cot \alpha \right)=-1$.
	\end{enumerate}
		}
\end{ex}

% #3638 | Archive/SP_KhoLatex/ID6_old/1.2.2.tex | mã gốc 1D1K2-5
\begin{ex}%[1D1K2-5]
	\immini{
		Thanh $OM$ quay ngược chiều kim đồng hồ quanh trục $O$ của nó trên mặt phẳng thẳng đứng và in bóng vuông góc với mặt đất (tham khảo hình bên). Vị trí ban đầu của thanh là $OA$. Hỏi độ dài bóng $O'M'$ của $OM$ khi thanh quay được $5\dfrac{1}{6}$ vòng là bao nhiêu, biết rằng độ dài thanh $OM$ là $25$ cm? Kết quả làm tròn đến hàng phần mười.
	}
{\vspace{-0.5cm}
	\begin{tikzpicture}[scale=0.9, font=\footnotesize,line join=round, line cap=round, >=stealth,x=2cm,y=2cm]
		\def\x{1}
		\def\a{30}
		\pgfmathsetmacro{\y}{\x*sin(\a)}
		\pgfmathsetmacro{\z}{\x*cos(\a)}
		\fill (0,0) node [below right]{$O$} circle(1pt);
		\draw[dashed](1,0) arc (0:360:2cm);
		\fill (1,0) node[above right]{$A$}circle(1pt);
		\fill (210:1) node[left]{$M$} circle(1pt);
		\coordinate (P) at (-1.2,-1.2);
		\coordinate (Q) at (1.2,-1.2);
		\draw (1,0)--(0,0)--(210:1);
		\draw[->] (.2,0) arc (0:210:.4cm);
		\node at (-.2,.2){$\alpha$};
		\coordinate (M') at ($(P)!(210:1)!(Q)$);
		\coordinate (O') at ($(P)!(0,0)!(Q)$);
		\foreach \i in {M',O'}{\fill[](\i) circle (1pt);}
		\draw (P)--(Q) (210:1)--(M') (0,0)--(O');
		\foreach \i in {M',O'}{\draw(\i) node[scale=0.9,below]{$\i$};}
		\draw[->] (0,1.5)--(0,1.1);\draw[->](-0.2,1.5)--(-0.2,1.1);
		\draw[->](0.2,1.5)--(0.2,1.1);
		\node at (-.5,-1.3){bóng};
		\node at (.6,1.3){ánh sáng};	
	\end{tikzpicture}	}
\loigiai{
	Chọn hệ trục như hình vẽ.
	\begin{center}
			\begin{tikzpicture}[scale=0.9, font=\footnotesize,line join=round, line cap=round, >=stealth,x=2cm,y=2cm]
			\def\x{1}
			\def\a{30}
			\pgfmathsetmacro{\y}{\x*sin(\a)}
			\pgfmathsetmacro{\z}{\x*cos(\a)}
			\fill (0,0) node [below right]{$O$} circle(1pt);
			\draw[->] (-1.1,0)--(1.1,0)node[below]{$x$};
			\draw[->] (0,-1.1)--(0,1.1)node[left]{$y$};
			\draw[dashed](1,0) arc (0:360:2cm);
			\fill (1,0) node[above right]{$A$}circle(1pt);
			\fill (210:1) node[left]{$M$} circle(1pt);
			\coordinate (P) at (-1.2,-1.2);
			\coordinate (Q) at (1.2,-1.2);
			\draw (1,0)--(0,0)--(210:1);
			\draw[->] (.2,0) arc (0:210:.4cm);
			\draw[dashed] (210:1)--(-0.867,0)node[above right]{$H$}circle(1pt);
			\node at (-.2,.2){$\alpha$};
			\coordinate (M') at ($(P)!(210:1)!(Q)$);
			\coordinate (O') at ($(P)!(0,0)!(Q)$);
			\foreach \i in {M',O'}{\fill[](\i) circle (1pt);}
			\draw (P)--(Q) (210:1)--(M') (0,0)--(O');
			\foreach \i in {M',O'}{\draw(\i) node[scale=0.9,below]{$\i$};}
			\draw[->] (0,1.5)--(0,1.1);\draw[->](-0.2,1.5)--(-0.2,1.1);
			\draw[->](0.2,1.5)--(0.2,1.1);
			\node at (-.5,-1.3){bóng};
			\node at (.6,1.3){ánh sáng};	
		\end{tikzpicture}
	\end{center}
	Gọi $H$ là hình chiếu vuông góc của $M$ lên trục $Ox$.\\
	Thanh $OM$ quay được một góc lượng giác $\alpha =5\dfrac{1}{6}\cdot \left(360^\circ \right)=1860^\circ$.\\
	Khi đó điểm $M$ biểu diễn cho góc lượng giác $\alpha $ là $M\left( OM\cdot \cos \alpha ;OM\cdot \sin \alpha \right)$.\\
	$\Rightarrow OH=OM\cdot \cos \left(1860^\circ\right)=12{,}5~ \text{cm}\Rightarrow O'M'=12{,}5$ cm.}	
\end{ex}

% #3648 | Archive/SP_KhoLatex/ID6_old/1-CK1-OT-De-29.tex | mã gốc 1D1H4-6
\begin{bt}%Câu 36%[1D1H4-6]
	(0.5 điểm) Số giờ có ánh sáng mặt trời của thành phố $A$ ở vĩ độ $40^\circ $ Bắc trong ngày thứ $t$ của một năm không nhuận được cho bởi hàm số\\
	$d(t)=3\sin\left[\dfrac{\pi}{182}\cdot \left(t-80\right)\right]+12$ với $t\in\mathbb{Z}$ và $0\le t\le 365$.\\
	Hỏi thành phố $A$ có đúng $12$ giờ có ánh sáng mặt trời vào ngày nào trong năm?
	\loigiai{
		Ta có
		\begin{eqnarray*}
			&&
			d(t)=12\\
			&\Leftrightarrow &3\sin\left[\dfrac{\pi}{182}\cdot \left(t-80\right)\right]+12=12\\ 
			&\Leftrightarrow&\sin\left[\dfrac{\pi}{182}\cdot\left(t-80\right)\right]=0\\ 
			&\Leftrightarrow&\dfrac{\pi}{182}\cdot\left(t-80\right)=k\pi\\ 
			&\Leftrightarrow &t-80=182k\\ 
			&\Leftrightarrow &t=80+182k \,,\,\,\,k\in\mathbb{Z}. 
		\end{eqnarray*}
		Mà 
		\begin{eqnarray*}
			&&
			0\le t\le 365\\
			&\Leftrightarrow &0\le 80+182k\le 365\\ 
			&\Leftrightarrow&-\dfrac{40}{91}\le k\le\dfrac{285}{182}\,\,,\,\,\,k\in\mathbb{Z}\\ 
			&\Rightarrow &k\in\left\{ 0;1\right\}. 
		\end{eqnarray*}
		Vậy thành phố $A$ có đúng $12$ giờ có ánh sáng mặt trời vào ngày thứ 80 (ứng với $k=0$) và ngày thứ 262 (ứng với $k=1$).}
\end{bt}

% #3692 | Archive/SP_KhoLatex/ID6_old/12-OT-PTLuongGiacThuongGap.tex | mã gốc 1D1B3-1
\begin{bt}%[1D1B3-1]
	Giải phương trình: $4\cos^2x-4\sin x-1=0$. 
	\loigiai{\allowdisplaybreaks
	Ta có	
		\begin{eqnarray*}
			& 4\cos^2x-4\sin x-1=0 & \Leftrightarrow 4(1-\sin^2x)-4\sin x-1=0\\
			&  & \Leftrightarrow 4-4\sin^2x-4\sin x-1=0\\
			&  & \Leftrightarrow 4\sin^2{x}+4\sin{x}-3=0.
		\end{eqnarray*}
		Đặt $t=\sin x$ $(-1\le t\le 1)$. Khi đó, phương trình trở thành
		\begin{displaymath}
			4t^2+4t-3=0 \Leftrightarrow (2t-1)(2t+3)=0 \Leftrightarrow \hoac{&t=\dfrac{1}{2}\\&t=\dfrac{-3}{2}.}
		\end{displaymath}
		Vì $-1\le t\le 1$ nên $t=\sin x=\dfrac{1}{2}\Leftrightarrow \hoac{&x=\dfrac{\pi}{6}+k2\pi\\&x=\dfrac{5\pi}{6}+k2\pi}, (k\in\mathbb{Z})$.
	}
\end{bt}

% #3693 | Archive/SP_KhoLatex/ID6_old/12-OT-PTLuongGiacThuongGap.tex | mã gốc 1D1B3-1
\begin{bt}%[1D1B3-1]
	Giải phương trình $\cos 2x-3\cos x+2=0$. 
	\loigiai{
	Ta có	\allowdisplaybreaks
		\begin{eqnarray*}
			& \cos 2x-3\cos x+2=0 & \Leftrightarrow 2\cos^2x-1-3\cos x+2=0\\
			& & \Leftrightarrow 2\cos^2x-3\cos x+1=0.
		\end{eqnarray*}
		Đặt $t=\cos x$ $(-1\le t\le 1)$. Khi đó, phương trình trở thành
		\begin{displaymath}
			2t^2-3t+1=0 \Leftrightarrow (2t-1)(t-1)=0 \Leftrightarrow \hoac{&t=\dfrac{1}{2}\\&t=1.}
		\end{displaymath}
		Vì $-1\le t\le 1$ nên $\hoac{&t=\cos x=\dfrac{1}{2}\\&t=\cos x=1} \Leftrightarrow \hoac{&x=k2\pi\\&x=\dfrac{-\pi}{3}+k2\pi\\&x=\dfrac{\pi}{3}+k2\pi}, (k\in\mathbb{Z})$.
	}
\end{bt}

% #3694 | Archive/SP_KhoLatex/ID6_old/12-OT-PTLuongGiacThuongGap.tex | mã gốc 1D1K3-1
\begin{bt}%[1D1K3-1]
	Giải phương trình $4\sin^4x+5\cos^2x-4=0$. 
	\loigiai{
	Ta có	\allowdisplaybreaks
		\begin{eqnarray*}
			& 4\sin^4x+5\cos^2x-4=0 & \Leftrightarrow 4\sin^4x+5(1-\sin^2x)-4=0\\
			& & \Leftrightarrow 4\sin^4x-5\sin^2x+1=0.
		\end{eqnarray*}
		Đặt $t=\sin^2x$ $(0\le t\le 1)$. Khi đó, phương trình trở thành
		\begin{displaymath}
			4t^2-5t+1=0 \Leftrightarrow  (4t-1)(t-1)=0 \Leftrightarrow \hoac{&t=\dfrac{1}{4}\\&t=1.}
		\end{displaymath}
		Vì $0\le t\le 1$ nên $\hoac{&t=\sin^2x=\dfrac{1}{4}\\&t=\sin^2x=1}\Leftrightarrow\hoac{&t=\sin x=\pm\dfrac{1}{2}\\&t=\sin x=\pm1}\Leftrightarrow\hoac{&x=\dfrac{-\pi}{2}+k\pi\\&x=\dfrac{-\pi}{6}+k\pi\\&x=\dfrac{\pi}{6}+k\pi}, (k\in\mathbb{Z})$.}
\end{bt}

% #3695 | Archive/SP_KhoLatex/ID6_old/12-OT-PTLuongGiacThuongGap.tex | mã gốc 1D1B3-3
\begin{bt}%[1D1B3-3]
	Giải phương trình $\sin x-\sqrt{3}\cos x=-\sqrt{2}$.
	\loigiai{
		Ta có
		\begin{eqnarray*}
		&&	\sin x-\sqrt{3}\cos x=-\sqrt{2}\\
		&\Leftrightarrow & \dfrac{1}{2}\sin x-\dfrac{\sqrt{3}}{2}\cos x=-\dfrac{\sqrt{2}}{2}\\
			&\Leftrightarrow & \sin x\cos \dfrac{\pi}{3}-\cos x\sin \dfrac{\pi}{3}=-\dfrac{\sqrt{2}}{2}\\
			&\Leftrightarrow &\sin \left(x-\dfrac{\pi}{3}\right) =-\dfrac{\sqrt{2}}{2}\\
			&\Leftrightarrow &\hoac{&x-\dfrac{\pi}{3}=-\dfrac{\pi}{4}+k2\pi\\&x-\dfrac{\pi}{3}=\pi-\left( -\dfrac{\pi}{4}\right)+k2\pi},  \left(k\in\mathbb{Z}\right)\\
			&\Leftrightarrow &\hoac{&x=\dfrac{\pi}{12}+k2\pi\\&x=\dfrac{19\pi}{12}+k2\pi}, \left(k\in\mathbb{Z}\right).
		\end{eqnarray*}
		Vậy phương trình có các nghiệm $x=\dfrac{\pi}{12}+k2\pi$; $x=\dfrac{19\pi}{12}+k2\pi$, với $k\in\mathbb{Z}$.
	}
\end{bt}

% #3696 | Archive/SP_KhoLatex/ID6_old/12-OT-PTLuongGiacThuongGap.tex | mã gốc 1D1B3-3
\begin{bt}%[1D1B3-3]
	Giải phương trình $\sin 3x - \sqrt{3}\cos 3x-1=0$.
	\loigiai{
	Ta có
		\begin{eqnarray*}
			&&\sin 3x - \sqrt{3}\cos 3x-1=0 \\
			&\Leftrightarrow&\dfrac{1}{2}\sin 3x - \dfrac{\sqrt{3}}{2}\cos3x = \dfrac{1}{2}\\ &\Leftrightarrow& \cos\dfrac{\pi}{3}\sin 3x - \sin\dfrac{\pi}{3}\cos3x = \dfrac{1}{2}\\
			&\Leftrightarrow& \sin \left(3x-\dfrac{\pi}{3}\right)=\sin \dfrac{\pi}{6}\\
			&\Leftrightarrow& \hoac{&3x-\dfrac{\pi}{3}=\dfrac{\pi}{6}+k2\pi\\&3x-\dfrac{\pi}{3}=\dfrac{5\pi}{6}+k2\pi} , \left(k\in\mathbb{Z}\right)\\
			&\Leftrightarrow& \hoac{&x=\dfrac{\pi}{6}+k\dfrac{2\pi}{3}\\&x=\dfrac{7\pi}{18}+k\dfrac{2\pi}{3}} , \left(k\in\mathbb{Z}\right).
		\end{eqnarray*}
	}
\end{bt}

% #3697 | Archive/SP_KhoLatex/ID6_old/12-OT-PTLuongGiacThuongGap.tex | mã gốc 1D1B3-3
\begin{bt}%[1D1B3-3]
	Giải phương trình $\left(\sin\dfrac{x}{2}+\cos\dfrac{x}{2}\right)^2+\sqrt{3}\cos x=2$.
	\loigiai{
		{\allowdisplaybreaks
			\begin{eqnarray*}
				& & \left(\sin\dfrac{x}{2}+\cos\dfrac{x}{2}\right)^2+\sqrt{3}\cos x=2\\
				&\Leftrightarrow&
				\sin x+\sqrt{3}\cos x=1\\
				&\Leftrightarrow&
				\dfrac{1}{2}\sin x+\dfrac{\sqrt{3}}{2}\cos x=\dfrac{1}{2}\\
				&\Leftrightarrow&
				\sin\left(x+\dfrac{\pi}{3}\right)=\dfrac{1}{2}
				\\ &\Leftrightarrow&
				\hoac{&x+\dfrac{\pi}{3}=\dfrac{\pi}{6}+k2\pi \\&x+\dfrac{\pi}{3}=\dfrac{5\pi}{6}+k2\pi}\\
				&\Leftrightarrow&
				\hoac{&x=-\dfrac{\pi}{6}+k2\pi\\&x=\dfrac{\pi}{2}+k2\pi}, \left(k\in\mathbb{Z}\right).
			\end{eqnarray*}
		}
	}
\end{bt}

% #3698 | Archive/SP_KhoLatex/ID6_old/12-OT-PTLuongGiacThuongGap.tex | mã gốc 1D1K3-6
\begin{bt}%[1D1K3-6]
	Giải phương trình $\dfrac{\sqrt{3}\sin x-\cos x-1}{\sin 2x}=0$.
	\loigiai{
		Điều kiện $\sin 2x\ne 0\Leftrightarrow x\ne \dfrac{k\pi}{2}$ với $k\in\mathbb{Z}$.\\
		Khi đó phương trình đã cho trở thành
		\allowdisplaybreaks
		\begin{eqnarray*}
			\sqrt{3}\sin x-\cos x-1=0&\Leftrightarrow&\dfrac{\sqrt{3}}{2}\sin x-\dfrac{1}{2}\cos x=\dfrac{1}{2}\\
			&\Leftrightarrow&\sin\left(x-\dfrac{\pi}{6}\right)=\dfrac{1}{2}\\
			&\Leftrightarrow&\hoac{&x-\dfrac{\pi}{6}=\dfrac{\pi}{6}+k2\pi\\&x-\dfrac{\pi}{6}=\dfrac{5\pi}{6}+k2\pi}\\
			&\Leftrightarrow&\hoac{&x=\dfrac{\pi}{3}+k2\pi\\&x=\pi+k2\pi}, \left(k\in\mathbb{Z}\right).
		\end{eqnarray*}
		So sánh với điều kiện xác định của phương trình ta có nghiệm của phương trình là $x=\dfrac{\pi}{3}+k2\pi$ với $k\in\mathbb{Z}$.
	}
\end{bt}

% #3699 | Archive/SP_KhoLatex/ID6_old/12-OT-PTLuongGiacThuongGap.tex | mã gốc 1D1B3-2
\begin{bt}%[1D1B3-2]
	Giải phương trình lượng giác $4\sin^2x-6\sin x\cos x-\cos 2x=5$.
	\loigiai{
		Ta có $ 4\sin^2x-6\sin x\cos x-\cos 2x=5\Leftrightarrow 2\sin^2 x-3\sin x\cos x-\cos^2x=2\ \ (1)  $.\\
		Do $2=2\left(\sin^2{x}+\cos^2{x}\right)$ nên phương trình $(1)$ là phương trình đẳng cấp bậc hai với $\sin{x}$ và $\cos{x}$.
		\begin{itemize}
			\item Trường hợp $ \cos x=0\Leftrightarrow x=\dfrac{\pi}{2}+k\pi $, thay vào phương trình $ (1) $, ta được $ 2=2  \text{ (luôn đúng)}$, suy ra $ x=\dfrac{\pi}{2}+k\pi,\, k\in\mathbb{Z}$ là nghiệm của phương trình $ (1) $.
			\item Trường hợp $ \cos x\ne 0 $, chia hai vế phương trình $ (1) $ cho $ \cos^2x $, ta được $$ 2\tan^2x-3\tan x-1=2(1+\tan^2x)\Leftrightarrow \tan x=-1\Leftrightarrow x=-\dfrac{\pi}{4}+k\pi,\, k\in\mathbb{Z}.  $$
		\end{itemize}
		Vậy các nghiệm của phương trình $ (1) $ là $ x=\dfrac{\pi}{2}+k\pi $, $ x=-\dfrac{\pi}{4}+k\pi$ với $k\in\mathbb{Z}$.
	}
\end{bt}

% #3700 | Archive/SP_KhoLatex/ID6_old/12-OT-PTLuongGiacThuongGap.tex | mã gốc 1D1K3-4
\begin{bt}%[1D1K3-4]
	Giải phương trình $ \sin 2x+(2-\sqrt{2})(\sin x+\cos x)+1-2\sqrt{2}=0 $.
	\loigiai{
	%	$ \sin 2x+(2-\sqrt{2})(\sin x+\cos x)+1-2\sqrt{2}=0 $. \tagEX{1}
	\noindent	Đặt $ t=\sin x+\cos x $	($ |t|\leq \sqrt{2} $), suy ra $ t^{2}=1+2\sin x\cos x\Rightarrow \sin 2x=t^{2}-1 $.\\
	Phương trình đã cho trở thành
		% Thay vào phương trình (1) ta được
		\[
		t^{2}-1+(2-\sqrt{2})t+1-2\sqrt{2}=0\Leftrightarrow t^{2}+(2-\sqrt{2})t-2\sqrt{2}=0\Leftrightarrow\hoac{& t=\sqrt{2}\text{ (thỏa mãn)}\\ &t=-2\text{ (loại)}.}
		\]
		Với $ t=\sqrt{2} $, ta có
		$$ \sin x+\cos x=\sqrt{2}\Leftrightarrow \dfrac{1}{\sqrt{2}}\sin x+\dfrac{1}{\sqrt{2}}\cos x=1\Leftrightarrow \cos \left(x-\dfrac{\pi}{4}\right)=1\Leftrightarrow x=\dfrac{\pi}{4}+k2\pi, \left(k\in\mathbb{Z}\right).$$
	}
\end{bt}

% #3701 | Archive/SP_KhoLatex/ID6_old/12-OT-PTLuongGiacThuongGap.tex | mã gốc 1D1B3-4
\begin{bt}%[1D1B3-4]
	Giải phương trình $ 2(\tan^{2}x+\cot^{2}x)-(4-\sqrt{2})(\tan x+\cot x)+4+2\sqrt{2}=0 $.
	\loigiai
	{
		Điều kiện xác định $ \heva{& \sin x\neq 0\\ &\cos x\neq 0}\Leftrightarrow x\neq k\dfrac{\pi}{2} $.\\
		Đặt $ t=\tan x+\cot x $	($ |t|\geq 2 $), suy ra $ t^{2}=\tan^{2}x+\cot^{2}x+2\Rightarrow \tan^{2}x+\cot^{2}x=t^{2}-2 $. Thay vào (1) ta được
		\[
		2(t^{2}-2)-(4-\sqrt{2})t+4+2\sqrt{2}=0\Leftrightarrow 2t^{2}-(4-\sqrt{2})t+2\sqrt{2}=0\text{ (vô nghiệm)}.
		\]
		Vậy phương trình vô nghiệm.
	}
\end{bt}

% #3702 | Archive/SP_KhoLatex/ID6_old/12-OT-PTLuongGiacThuongGap.tex | mã gốc 1D1B3-5
\begin{bt}%[1D1B3-5]
	Giải phương trình $1 + \cos x + \cos 2x + \cos 3x = 0$.
	\loigiai{
		Ta có
		\begin{eqnarray*}
			1 + \cos x + \cos 2x + \cos 3x = 0 & \Leftrightarrow &(\cos 2x + 1) + (\cos x + \cos 3x) = 0\\
			&\Leftrightarrow & 2 \cos ^2x + 2\cos 2x \cos x = 0\\
			& \Leftrightarrow & \cos x (\cos x + \cos 2x) = 0\\
			&\Leftrightarrow & \hoac{&\cos x = 0\\&\cos x + \cos 2x = 0}\\
			&\Leftrightarrow & \hoac{&\cos x = 0\\& \cos 2x = \cos(\pi-x) }\\
			& \Leftrightarrow & \hoac{&x = \dfrac{\pi}{2} + k \pi\\ &x = -\pi + k2\pi\\ &x = \dfrac{\pi}{3} + k\dfrac{2\pi}{3}}, \left(k\in\mathbb{Z}\right).
	\end{eqnarray*}}
\end{bt}

% #3703 | Archive/SP_KhoLatex/ID6_old/12-OT-PTLuongGiacThuongGap.tex | mã gốc 1D1B3-5
\begin{bt}%[1D1B3-5]
	Giải phương trình $\sin 2x\cos x=\sin 7x\cos 4x$.
	\loigiai
	{
		Ta có
		\allowdisplaybreaks
		\begin{eqnarray*}
			&& \sin 2x\cos x=\sin 7x\cos 4x \\
			&\Leftrightarrow & \dfrac{1}{2}\left(\sin 3x + \sin x\right) = \dfrac{1}{2}\left(\sin 11x + \sin 3x\right)\\
			&\Leftrightarrow & \sin x = \sin 11x \\
			&\Leftrightarrow & \left[\begin{aligned}&11x=x+k2\pi	\\&11x=\pi-x+k2\pi\end{aligned}\right. , \left(k\in\mathbb{Z}\right)\\
			&\Leftrightarrow & \left[\begin{aligned}&x=\dfrac{k\pi}{5}	\\&x=\dfrac{\pi}{12}+\dfrac{k\pi}{6}\end{aligned}\right. , \left(k\in\mathbb{Z}\right).
		\end{eqnarray*}
		Vậy phương trình đã cho có các nghiệm $x=\dfrac{k\pi}{5}$, $x=\dfrac{\pi}{12}+\dfrac{k\pi}{6}$ với $k\in\mathbb{Z}$.
	}
\end{bt}

% #3704 | Archive/SP_KhoLatex/ID6_old/12-OT-PTLuongGiacThuongGap.tex | mã gốc 1D1B3-5
\begin{bt}%[1D1B3-5]
	Giải phương trình $2\cos 3x\cos x-\cos 4x-3\cos x=-2$.
	\loigiai
	{
		Ta có
		\allowdisplaybreaks
		\begin{eqnarray*}
			&& 2\cos 3x\cos x-\cos 4x-3\cos x=-2 \\
			&\Leftrightarrow & \cos 4x + \cos 2x - \cos 4x -3\cos x + 2=0 \\
			&\Leftrightarrow & 2\cos^2 x-3\cos x+1=0 \\
			&\Leftrightarrow & \left[\begin{aligned}&\cos x=1 \\&\cos x=\dfrac{1}{2}\end{aligned}\right. \\
			&\Leftrightarrow & \left[\begin{aligned}&x=k2\pi \\&x=\pm\dfrac{\pi}{3}+k2\pi\end{aligned}\right. , \left(k\in\mathbb{Z}\right).
		\end{eqnarray*}
		Vậy phương trình đã cho có các nghiệm $x=k2\pi$, $x=\pm\dfrac{\pi}{3}+k2\pi$ với $k\in\mathbb{Z}$.
	}
\end{bt}

% #3705 | Archive/SP_KhoLatex/ID6_old/12-OT-PTLuongGiacThuongGap.tex | mã gốc 1D1B3-6
\begin{bt}%[1D1B3-6]
	Giải phương trình $\dfrac{\cos x+ \sin 2x}{\cos 3x}+1=0$.
	\loigiai{
	\begin{itemize} 
 \item Điều kiện $ \cos 3x \ne 0 \Leftrightarrow x\ne \dfrac{\pi}{6}+\dfrac{k \pi }{3},\,\,(k\in \mathbb{Z})$.
 \item Phương trình đã cho tương đương với
 \begin{eqnarray*} 
 	& & \cos x-\cos 3x +\sin 2x=0\\ 
 	&\Leftrightarrow & 2\sin 2x \sin x +\sin 2x=0\\ 
 	&\Leftrightarrow & \hoac{&\sin 2x =0\\&\sin x=-\dfrac{1}{2}} \Leftrightarrow \hoac{&x=\dfrac{k\pi}{2}\\&x=-\dfrac{\pi}{6}+k2\pi \\& x=\dfrac{7\pi}{6}+k2\pi}, \left(k\in\mathbb{Z}\right). 
 \end{eqnarray*} 
 Vì $\cos 3x=4\cos^3 x-3\cos x$ nên đối chiếu điều kiện ta phải loại đi những giá trị mà $\cos x=0$.\\
 Cuối cùng ta được $\hoac{&x=k\pi \\&x=-\dfrac{\pi}{6}+k2\pi \\& x=\dfrac{7\pi}{6}+k2\pi},\,\,(k\in \mathbb{Z})$.
 \end{itemize}
	} 
\end{bt}

% #3747 | Archive/SP_KhoLatex/ID6_old/1.2.1.tex | mã gốc 1D1B2-4
\begin{ex}%[1D1B2-4]
Chứng minh các đẳng thức lượng giác sau
\begin{multicols}{2}
\begin{enumerate}[a)]
\item  $\sin^4x+\cos^4x=1-2\sin^2x\cdot\cos^2x$	
\item $\dfrac{\sin x+\cos x}{\sin^3x}=\cot^3x+\cot^2x+\cot x+1$.
\end{enumerate}
\end{multicols}
\loigiai{
	\begin{enumerate}[a)]
\item Ta có $\sin^4x+\cos^4x=\left(\sin^2x\right)^2+\left(\cos^2x\right)^2
=\left(\sin^2x+\cos^2x\right)^2-2\sin^2x.\cos^2x=1-2\sin^2x\cdot \cos^2x$.
\item Ta có $\dfrac{\sin x+\cos x}{\sin^3x}=\dfrac{1}{\sin^2x}+\dfrac{\cos x}{\sin^3x}$.\\
Mà $\cot^2x+1=\dfrac{1}{\sin^2x}$ và $\cot x=\dfrac{\cos x}{\sin x}$ nên
$\dfrac{1}{\sin^2x}+\dfrac{\cos x}{\sin^3x}=\cot^2x+1+\cot x\left(\cot^2x+1\right)=\cot^3x+\cot^2x+\cot x+1$.
\end{enumerate}
}
\end{ex}

% #3748 | Archive/SP_KhoLatex/ID6_old/1.2.1.tex | mã gốc 1D1B2-4
\begin{ex}%[1D1B2-4]
Chứng minh các đẳng thức lượng giác sau
\begin{multicols}{2}
\begin{enumerate}[a)] 
\item $\sin^4x+2\cos^2x=1+\cos^4x$
\item $\dfrac{\cot^2x-\cot^2y}{\cot^2x\cdot \cot^2y}=\dfrac{\cos^2x-\cos^2y}{\cos^2x\cdot \cos^2y}$.
\end{enumerate}
\end{multicols}
\loigiai{
	\begin{enumerate}[a)] 
\item Đẳng thức tương đương với $\sin^4x=1-2\cos^2x+\left(\cos^2x\right)^2$
$\Leftrightarrow{\sin^4}x=\left(1-\cos^2x\right)^2\quad(*)$\\
Mà $\sin^2x+\cos^2x=1\Leftrightarrow{\sin^2}x=1-\cos^2x$\\
Do đó (*) $\Leftrightarrow \cos^4 x=\left(\cos^2x\right)^2$
\item  Ta có 
\begin{eqnarray*}
	\dfrac{\cot^2x-\cot^2y}{\cot^2x\cdot \cot^2y}&=&\dfrac{1}{\cot^2y}-\dfrac{1}{\cot^2x}=\tan^2y-\tan^2x
	=\left(\dfrac{1}{\cos^2y}-1\right)-\left(\dfrac{1}{\cos^2x}-1\right)\\
	&=&\dfrac{1}{\cos^2y}-\dfrac{1}{\cos^2x}=\dfrac{\cos^2x-\cos^2y}{\cos^2x\cdot\cos^2y}.
\end{eqnarray*}
\end{enumerate}
}
\end{ex}

% #3749 | Archive/SP_KhoLatex/ID6_old/1.2.1.tex | mã gốc 1D1B2-4
\begin{ex}%[1D1B2-4]
Rút gọn các biểu thức sau
\begin{enumerate}[a)]
\item  $A=\dfrac{\sin \alpha +\cos \alpha}{\cos \alpha \left(1+\tan \alpha\right)}$	
\item  $B=\sin \left(\dfrac{5\pi}{2}-\alpha\right)+\cos \left(3\pi +\alpha\right)$
\item  $C=\cos \left(5\pi -\alpha\right)-\sin \left(\dfrac{3\pi}{2}+\alpha\right)+\tan \left(\dfrac{3\pi}{2}-\alpha\right)+\cot \left(3\pi -\alpha\right)$.
\end{enumerate}
\loigiai{
	\begin{enumerate}[a)]
\item  $A=\dfrac{\sin \alpha +\cos \alpha}{\cos \alpha \left(1+\dfrac{\sin \alpha}{\text{cos}\alpha}\right)}=\dfrac{\sin \alpha +\cos \alpha}{\cos \alpha \left(\dfrac{\text{cos}\alpha +\sin \alpha}{\text{cos}\alpha}\right)}=1$.
\item $B=\sin \left(2\pi +\dfrac{\pi}{2}-\alpha\right)+\cos \left(2\pi +\pi +\alpha\right)=\sin \left(\dfrac{\pi}{2}-\alpha\right)+\cos \left(\pi +\alpha\right)=\cos \alpha -\cos \alpha =0$.
\item Ta có 
\begin{itemize}
	\item $\cos (5\pi -\alpha)=\cos \left(\pi -\alpha +2\cdot 2\pi\right)
	=\cos \left(\pi -\alpha\right)=-\cos \alpha$.
	\item $\sin \left(\dfrac{3\pi}{2}+\alpha\right)=\sin \left(\pi +\dfrac{\pi}{2}+\alpha\right)=-\sin \left(\dfrac{\pi}{2}+\alpha\right)=-\cos \alpha$.
	\item $\tan \left(\dfrac{3\pi}{2}-\alpha\right)=\tan \left(\pi +\dfrac{\pi}{2}-\alpha\right)=\tan \left(\dfrac{\pi}{2}-\alpha\right)=\cot \alpha$.
	\item $\cot (3\pi -\alpha)=\cot \left(-\alpha\right)=-\cot \alpha $.
\end{itemize}
Suy ra $C=-\cos \alpha -\left(-\cos \alpha\right)+\cot \alpha +\left(-\cot \alpha\right)=0$.
\end{enumerate}}
\end{ex}

% #3750 | Archive/SP_KhoLatex/ID6_old/1.2.1.tex | mã gốc 1D1K2-5
\begin{ex}%[1D1K2-5]
Thanh $OM$ quay ngược chiều kim đồng hồ quanh trục $O$ của nó trên một mặt phẳng thẳng đứng và in bóng vuông góc xuống mặt đất như Hình vẽ. Vị trí ban đầu của thanh là $OA$. Hỏi độ dài bóng $O'M'$ của $OM$ khi thanh quay được $2\dfrac{1}{3}$ vòng là bao nhiêu, biết độ dài thanh $OM$ là $20 ~\mathrm{cm}$? Kết quả làm tròn đến hàng phần mười.
\begin{center}
\begin{tikzpicture}[scale=0.9, font=\footnotesize,line join=round, line cap=round, >=stealth,x=2cm,y=2cm]
	\def\x{1}
	\def\a{30}
	\pgfmathsetmacro{\y}{\x*sin(\a)}
	\pgfmathsetmacro{\z}{\x*cos(\a)}
	\fill (0,0) node [below right]{$O$} circle(1pt);
	\draw[dashed](1,0) arc (0:360:2cm);
	\fill (1,0) node[above right]{$A$}circle(1pt);
	\fill (210:1) node[left]{$M$} circle(1pt);
	\coordinate (P) at (-1.2,-1.2);
	\coordinate (Q) at (1.2,-1.2);
	\draw (1,0)--(0,0)--(210:1);
	\draw[->] (.2,0) arc (0:210:.4cm);
	\node at (-.2,.2){$\alpha$};
	\coordinate (M') at ($(P)!(210:1)!(Q)$);
	\coordinate (O') at ($(P)!(0,0)!(Q)$);
	\foreach \i in {M',O'}{\fill[](\i) circle (1pt);}
	\draw (P)--(Q) (210:1)--(M') (0,0)--(O');
	\foreach \i in {M',O'}{\draw(\i) node[scale=0.9,below]{$\i$};}
	\draw[->] (0,1.5)--(0,1.1);\draw[->](-0.2,1.5)--(-0.2,1.1);
	\draw[->](0.2,1.5)--(0.2,1.1);
	\node at (-.5,-1.3){bóng};
	\node at (.6,1.3){ánh sáng};	
\end{tikzpicture}
\end{center}
\loigiai{
Đặt hệ trục tọa độ như hình vẽ
\begin{center}
		\begin{tikzpicture}[scale=0.9, font=\footnotesize,line join=round, line cap=round, >=stealth,x=2cm,y=2cm]
		\def\x{1}
		\def\a{30}
		\pgfmathsetmacro{\y}{\x*sin(\a)}
		\pgfmathsetmacro{\z}{\x*cos(\a)}
		\fill (0,0) node [below right]{$O$} circle(1pt);
		\draw[->] (-1.1,0)--(1.1,0)node[below]{$x$};
		\draw[->] (0,-1.1)--(0,1.1)node[left]{$y$};
		\draw[dashed](1,0) arc (0:360:2cm);
		\fill (1,0) node[above right]{$A$}circle(1pt);
		\fill (210:1) node[left]{$M$} circle(1pt);
		\coordinate (P) at (-1.2,-1.2);
		\coordinate (Q) at (1.2,-1.2);
		\draw (1,0)--(0,0)--(210:1);
		\draw[->] (.2,0) arc (0:210:.4cm);
		\draw[dashed] (210:1)--(-0.867,0)node[above right]{$H$}circle(1pt);
		\node at (-.2,.2){$\alpha$};
		\coordinate (M') at ($(P)!(210:1)!(Q)$);
		\coordinate (O') at ($(P)!(0,0)!(Q)$);
		\foreach \i in {M',O'}{\fill[](\i) circle (1pt);}
		\draw (P)--(Q) (210:1)--(M') (0,0)--(O');
		\foreach \i in {M',O'}{\draw(\i) node[scale=0.9,below]{$\i$};}
		\draw[->] (0,1.5)--(0,1.1);\draw[->](-0.2,1.5)--(-0.2,1.1);
		\draw[->](0.2,1.5)--(0.2,1.1);
		\node at (-.5,-1.3){bóng};
		\node at (.6,1.3){ánh sáng};	
	\end{tikzpicture}
\end{center}
Điểm $M$ là điểm biểu diễn góc lượng giác $\alpha $. Kẻ $MH$ vuông góc với $Ox$.\\
Ta có $\alpha =2\dfrac{1}{3}\cdot 360^\circ =840^\circ $. Khi đó $M\left(20\cos 840^\circ;20\sin 840^\circ\right)$\\
Suy ra $OH=\left| 20\cos 840^\circ\right|=10\left(\mathrm{cm}\right)$.\\ Vậy độ dài bóng $O'M'$ của $OM$ khi thanh quay được $2\dfrac{1}{3}$ vòng là $10 ~\mathrm{cm}$.
}
\end{ex}

% #3773 | Archive/SP_KhoLatex/ID6_old/1K1-ot.tex | mã gốc 1D1B2-2
\begin{bt}%[Dự án 2 - TLDH - TeamTeXHoa - Quang Ngô]%[1D1B2-2]
    Cho góc $\alpha $ thoả mãn $\dfrac{\pi}{2}<\alpha <\pi$, $\cos\alpha =-\dfrac{1}{\sqrt{3}}$. Tính giá trị của các biểu thức sau:
    \begin{listEX}[2]
        \item $ \sin\left(\alpha+\dfrac{\pi}{6}\right)$;
        \item $ \cos\left(\alpha+\dfrac{\pi}{6}\right)$;
        \item $ \sin\left(\alpha-\dfrac{\pi}{3}\right)$; 
        \item $ \cos\left(\alpha-\dfrac{\pi}{6}\right)$.
    \end{listEX}
    \loigiai{
        Vì $\dfrac{\pi}{2}<\alpha <\pi $ nên $ \sin\alpha >0$. Mặt khác từ $ \sin^2\alpha+ \cos^2\alpha =1$.\\ 
        Suy ra  $ \sin\alpha =\sqrt{1- \cos^2\alpha}=\sqrt{1-{\left(-\dfrac{1}{\sqrt{3}}\right)}^2}=\dfrac{\sqrt{6}}{3}.$
        \begin{enumerate}
            \item $ \sin\left(\alpha+\dfrac{\pi}{6}\right)= \sin\alpha \cos\dfrac{\pi}{6}+ \cos\alpha \sin\dfrac{\pi}{6}$
            $=\dfrac{\sqrt{6}}{3}\cdot \dfrac{\sqrt{3}}{2}+\left(-\dfrac{1}{\sqrt{3}}\right)\cdot \dfrac{1}{2}=\dfrac{3\sqrt{2}-\sqrt{3}}{6}. $
            \item $ \cos\left(\alpha+\dfrac{\pi}{6}\right)= \cos\alpha \cos\dfrac{\pi}{6}- \sin\alpha \sin\dfrac{\pi}{6}$
            $=-\dfrac{1}{\sqrt{3}}\cdot \dfrac{\sqrt{3}}{2}-\dfrac{\sqrt{6}}{3}\cdot \dfrac{1}{2}=\dfrac{-3-\sqrt{6}}{6}. $
            \item $ \sin\left(\alpha-\dfrac{\pi}{3}\right)= \sin\alpha \cos\dfrac{\pi}{3}- \cos\alpha \sin\dfrac{\pi}{3}$
            $=\dfrac{\sqrt{6}}{3}\cdot \dfrac{1}{2}-\left(-\dfrac{1}{\sqrt{3}}\right)\cdot \dfrac{\sqrt{3}}{2}=\dfrac{\sqrt{6}+3}{6}.$
            \item $ \cos\left(\alpha-\dfrac{\pi}{6}\right)= \cos\alpha \cos\dfrac{\pi}{6}+ \sin\alpha \sin\dfrac{\pi}{6}$
            $=-\dfrac{1}{\sqrt{3}}\cdot \dfrac{\sqrt{3}}{2}+\dfrac{\sqrt{6}}{3}\cdot \dfrac{1}{2}=\dfrac{-3+\sqrt{6}}{6}.$
        \end{enumerate}
    }
\end{bt}

% #3774 | Archive/SP_KhoLatex/ID6_old/1K1-ot.tex | mã gốc 1D1B3-2
\begin{bt}%[Dự án 2 - TLDH - TeamTeXHoa - Quang Ngô]%[1D1B3-2]
    Cho góc bất kì $\alpha $. Chứng minh các đẳng thức sau:
    \begin{listEX}[2]
        \item $(\sin\alpha+ \cos\alpha)^2=1+ \sin2\alpha $. 
        \item $ \cos^4\alpha- \sin^4\alpha = \cos2\alpha $.
    \end{listEX}
    \loigiai{
        \begin{enumerate}
            \item Áp dụng hệ thức lượng giác cơ bản: $ \sin^2\alpha+ \cos^2\alpha =1$ và công thức nhân đôi: $ \sin2\alpha =2 \sin\alpha \cos\alpha $.\\
            Ta có: $VT=(\sin a+ \cos a)^2= \sin^2a+ \cos^2a+2 \sin a \cos a=1+ \sin2a=VP$ (đpcm).
            \item Áp dụng hệ thức Iượng giác cơ bản: $ \sin^2a+ \cos^2\alpha =1$ và công thức nhân đôi: $ \cos2\alpha = \cos^2\alpha- \sin^2\alpha $.\\
            Ta có: $VT= \cos^4a-\sin^4a=\left(\cos^2a\right)^2-\left(\sin^2a\right)^2$
            $=\left(\cos^2a+ \sin^2a\right)\left(\cos^2a- \sin^2a\right)=1\cdot \cos2a= \cos2a=VP$. (đpcm)
        \end{enumerate}
    }
\end{bt}

% #3775 | Archive/SP_KhoLatex/ID6_old/1K1-ot.tex | mã gốc 1D1B4-5
\begin{bt}%[Dự án 2 - TLDH - TeamTeXHoa - Quang Ngô]%[1D1B4-5]
    Tìm tập giá trị của các hàm sau:
    \begin{listEX}[2]
        \item $y=2 \cos\left(2x-\dfrac{\pi}{3}\right)-1$. 
        \item $y= \sin x+ \cos x$.
    \end{listEX}
    \loigiai{
        \begin{enumerate}
            \item Ta có\\
            \begin{eqnarray*}
                && -1\le \cos\left(2x-\dfrac{\pi}{3}\right)\le 1  ,\;\forall x\in \mathbb{R}. \\ 
                & \Leftrightarrow&-2\le 2 \cos\left(2x-\dfrac{\pi}{3}\right)\le 2  ,\;\forall x\in \mathbb{R}. \\ 
                & \Leftrightarrow&-2-1\le 2 \cos\left(2x-\dfrac{\pi}{3}\right)-1\le 2-1  ,\;\forall  x\in \mathbb{R}\\ 
                & \Leftrightarrow&-3\le 2 \cos\left(2x-\dfrac{\pi}{3}\right)-1\le 1  ,\;\forall  x\in \mathbb{R} \\ 
                & \Leftrightarrow&-3\le y\le 1  ,\;\forall  x\in \mathbb{R}.
            \end{eqnarray*}
            Vậy tập giá trị của hàm số $y=2\cos\left(2x-\dfrac{\pi}{3}\right)-1$ là $\left[ -3;1\right]$.
            \item Ta có
            \begin{eqnarray*}
                 \sin x+\cos x &=&\sqrt{2}\left(\dfrac{1}{\sqrt{2}} \sin x+\dfrac{1}{\sqrt{2}} \cos x\right) \\ 
                 &=&\sqrt{2}\left(\cos\dfrac{\pi}{4} \sin x+ \sin\dfrac{\pi}{4} \cos x\right) \\ 
                 &=&\sqrt{2}\left(\sin x \cos\dfrac{\pi}{4}+ \cos x \sin\dfrac{\pi}{4}\right) \\ 
                 &=&\sqrt{2} \sin\left(x+\dfrac{\pi}{4}\right).
            \end{eqnarray*}
            Khi đó ta có hàm số $ y=\sqrt{2} \sin\left(x+\dfrac{\pi}{4}\right)$.\\
            Lại có: $-1\le \sin\left(x+\dfrac{\pi}{4}\right)\le 1$ với mọi $ x\in \mathbb{R}$.\\
            $\Leftrightarrow-\sqrt{2}\le \sqrt{2} \sin\left(x+\dfrac{\pi}{4}\right)\le \sqrt{2}$ với mọi $x\in \mathbb{R}$.\\
            $\Leftrightarrow-\sqrt{2}\le y\le \sqrt{2}$ với mọi $x\in \mathbb{R}$.\\
            Vậy tập giá trị của hàm số $ y= \sin x+ \cos x$ là $\left[-\sqrt{2};\sqrt{2}\right]$.
        \end{enumerate}
    }
\end{bt}

% #3776 | Archive/SP_KhoLatex/ID6_old/1K1-ot.tex | mã gốc 1D1B5-5
\begin{bt}%[Dự án 2 - TLDH - TeamTeXHoa - Quang Ngô]%[1D1B5-5]
    Giải các phương trình sau:
    \begin{enumerate}
        \item $ \cos\left(3x-\dfrac{\pi}{4}\right)=-\dfrac{\sqrt{2}}{2}$. 
        \item $2 \sin^2x-1+ \cos3x=0$
        \item $ \tan\left(2x+\dfrac{\pi}{5}\right)= \tan\left(x-\dfrac{\pi}{6}\right)$
    \end{enumerate}
\loigiai{
    \begin{enumerate}
        \item $\cos\left(3x-\dfrac{\pi}{4}\right)=-\dfrac{\sqrt{2}}{2} \\
        \Leftrightarrow\cos\left(3x-\dfrac{\pi}{4}\right)=\cos\dfrac{3\pi}{4} \\
        \Leftrightarrow \hoac{
            &3x-\dfrac{\pi}{4}=\dfrac{3\pi}{4}+k2\pi \\
            &3x-\dfrac{\pi}{4}=-\dfrac{3\pi}{4}+k2\pi}\left(k\in \mathbb{Z}\right)\\
        \Leftrightarrow \hoac{
            &3x=\pi+k2\pi \\
            &3x=-\dfrac{\pi}{2}+k2\pi}\left(k\in \mathbb{Z}\right)\\
        \Leftrightarrow \hoac{
            &x=\dfrac{\pi}{3}+k\dfrac{2\pi}{3} \\
            &x=-\dfrac{\pi}{6}+k\dfrac{2\pi}{3}}\left(k\in \mathbb{Z}\right)$\\
        Vậy phương trình đã cho có các nghiệm là $x=\dfrac{\pi}{3}+k\dfrac{2\pi}{3},k\in \mathbb{Z}$ và $x=-\dfrac{\pi}{6}+k\dfrac{2\pi}{3},k\in \mathbb{Z}$.
        \item $2\sin^2x-1+ \cos3x=0$\\
        $\Leftrightarrow-\left(1-2 \sin^2x\right)+ \cos3 x=0$\\
        $\Leftrightarrow- \cos2 x+ \cos3 x=0$\\
        $\Leftrightarrow \cos3 x= \cos2 x$\\
        $\Leftrightarrow \hoac{
            &3x=2x+k2\pi \\
            &3x=-2x+k2\pi} \left(k\in \mathbb{Z}\right)$\\
        $\Leftrightarrow \hoac{
            &x=k2\pi \\
            &5x=k2\pi}\left(k\in \mathbb{Z}\right)$\\
        $\Leftrightarrow \hoac{
            &x=k2\pi \\
            &x=k\dfrac{2\pi}{5}}\left(k\in \mathbb{Z}\right)$
        \item $ \tan\left(2x+\dfrac{\pi}{5}\right)= \tan\left(x-\dfrac{\pi}{6}\right)$\\
        $\Leftrightarrow 2x+\dfrac{\pi}{5}=x-\dfrac{\pi}{6}+k\pi,k\in \mathbb{Z}$\\
        $\Leftrightarrow x=-\dfrac{11\pi}{30}+k\pi,k\in \mathbb{Z}$\\
        Vậy phương trình đã cho có các nghiệm là  $$x=-\dfrac{11\pi}{30}+k\pi,k\in \mathbb{Z}.$$
    \end{enumerate}}
\end{bt}

% #3777 | Archive/SP_KhoLatex/ID6_old/1K1-ot.tex | mã gốc 1D1K2-5
\begin{bt}%[Dự án 2 - TLDH - TeamTeXHoa - Quang Ngô]%[1D1K2-5]
    Huyết áp là áp lực cần thiết tác động lên thành của động mạch để đưa máu từ tim đến nuôi dưỡng các mô trong cơ thể. Huyết áp được tạo ra do lực co bóp của cơ tim và sức cản của thành động mạch. Mỗi lần tim đập, huyết áp của chúng ta tăng rồi giảm giữa các nhịp. Huyết áp tối đa và huyết áp tối thiểu được gọi tương ứng là huyết áp tâm thu và tâm trương. Chỉ số huyết áp của chúng ta được viết là huyết áp tâm thu/huyết áp tâm trương. Chỉ số huyết áp 120/80 là bình thường. Giả sử huyết áp của một người nào đó được mô hình hoá bởi hàm số
    trong đó $p(t)$. là huyết áp tính theo đơn vị mmHg. (milimét thuỷ ngân) và thời gian $t$ tính theo phút.
    \begin{enumerate}
        \item Tìm chu kì của hàm số $p(t)$.
        \item Tìm số nhịp tim mỗi phút.
        \item Tìm chỉ số huyết áp. So sánh huyết áp của người này với huyết áp bình thường.
    \end{enumerate}
    \loigiai{
        \begin{enumerate}
            \item Chu kì của hàm số $ p\left(t\right)$ là $T=\dfrac{2\pi}{160\pi}=\dfrac{1}{80}$.
            \item Thời gian giữa hai lần tim đập là $T=\dfrac{1}{80}$ (phút)\\
            Số nhịp tim mỗi phút là $1:\dfrac{1}{80}=80$ nhịp.
            \item Ta có: $-1\le \sin\left(160\pi t\right)\le 1$ với mọi $t\in \mathbb{R}$\\
            $\Leftrightarrow-25\le 25 \sin\left(160\pi t\right)\le 25$ với mọi $ t\in \mathbb{R}$\\
            $\Leftrightarrow 115+\left(-25\right)\le 115+25 \sin\left(160\pi t\right)\le 115+25$ với mọi $ t\in \mathbb{R}$\\
            $\Leftrightarrow 90\le p (t)\le 140$ với mọi $ t\in \mathbb{R}$.\\
            Do đó, chỉ số huyết áp của người này là 140/90 và chỉ số huyết áp của người này cao hơn mức bình thường.
    \end{enumerate}}
\end{bt}

% #3778 | Archive/SP_KhoLatex/ID6_old/1K1-ot.tex | mã gốc 1D1K5-6
\begin{bt}%[Dự án 2 - TLDH - TeamTeXHoa - Quang Ngô]%[1D1K5-6]
    Khi một tia sáng truyền từ ông khí vào mặt nước thì một phần tia sáng bị phản xạ trên bề mặt, phần còn lại bị khúc xạ như trong Hình. Góc tới $i$ liên hệ với góc khúc xạ $r$ bởi Định luật khúc xạ ánh $\dfrac{\sin i}{\sin r}=\dfrac{n_2}{n_1}. $
    \begin{center}
        \begin{tikzpicture}[line join = round, line cap = round, >=stealth, font=\footnotesize, scale=1]
            
            \path 
            (-1,0) coordinate (A)
            (5,0) coordinate (B)
            (2,0) coordinate (C)
            (-1,2) coordinate (D)
            (5,2) coordinate (E)
            (2,3) coordinate (F)
            (2,-3) coordinate (G)
            (4,-3) coordinate (H)
            ;
            \fill[color=blue!20] (A)--(B)--(5,-3.2)--(-1,-3.2)--(A);
            \draw
            (A)--(B)
            ($(C)!.7!(H)$)--(H)node[right]{$R$}
            ($(D)!.3!(C)$)--(C)node[below left]{$I$}
            ($(C)!.7!(E)$)--(E)node[above]{$S'$}
            pic[draw,black,angle radius=0.3cm] {angle = F--C--D}
            (C)++(120:.5)node{$i$}
            pic[draw,black,angle radius=0.4cm] {angle = E--C--F}
            (C)++(60:.6)node{$i'$}
            pic[draw,black,angle radius=0.3cm] {angle = G--C--H}
            pic[draw,black,angle radius=0.4cm] {angle = G--C--H}
            (C)++(-70:.6)node{$r$}
            ;
            \draw[->]
            (C)--($(C)!.7!(H)$);
            \draw[->]
            (D)node[above]{$S$}--($(D)!.3!(C)$);
            \draw[->]           
            (C)--($(C)!.7!(E)$);
            \draw[dashed] (F)node[left]{$N$}--(G)node[left]{$N'$};
            \draw 
            (A)++(60:.5) circle (2mm)node{$1$}
            (A)++(-60:.5) circle (2mm)node{$2$}
            ;
        \end{tikzpicture}
    \end{center}
    Ở đây, $n_1$ và $n_2$ tương ứng là chiết suất của môi trường 1 (không khí) và môi trường 2 (nước). Cho biết góc tới $i=50^\circ$, hãy tính góc khúc xạ, biết rằng chiết suất của không khí bằng 1 còn chiết suất của nước là $1,33$.
    \loigiai{
        Theo bài ra ta có: $ i=50^\circ$, $n_1=1$, $n_2=1,33$, thay vào $\dfrac{\sin i}{\sin r}=\dfrac{n_2}{n_1}$ ta được:\\
        $\dfrac{\sin{50}^\circ}{\sin r}=\dfrac{1,33}{1}$ (đk $\sin r\neq 0$)\\
        $\Rightarrow \sin r=\dfrac{\sin50^\circ}{1,33}$$\Leftrightarrow \sin r\approx 0,57597$ (thỏa mãn điều kiện)\\
        $\Leftrightarrow \sin r\approx \sin\left(35^\circ10^\prime\right)$\\
        $\Leftrightarrow \hoac{
            &r\approx 35^\circ10^\prime+k360^\circ \\
            &r\approx 180^\circ-35^\circ10^\prime+k360^\circ 
        }\left(k\in \mathbb{Z}\right) \\
        \Leftrightarrow \hoac{
            &r\approx 35^\circ10^\prime+k360^\circ \\
            &r\approx 144^\circ50^\prime+k360^\circ 
        }\left(k\in \mathbb{Z}\right)$ \\
        Mà $0^\circ<r<90^\circ$ nên $r\approx 35^\circ10^\prime$.\\
        Vậy góc khúc xạ $r\approx 35^\circ10^\prime.$
    }
\end{bt}

% #3871 | Archive/SP_KhoLatex/SP_Đề Theo Bài-Khối 11/1-B3.tex | mã gốc 1D1C1-1
\begin{ex}%[1D1C1-1]
	Một chất điểm dao động điều hòa xung quanh điểm  $O$ theo phương nằm ngang. Biết rằng vị trí của chất điểm đó so với $O$  tại thời điểm $t$ (giây) được xác định theo phương trình $x=3\cos \left(\dfrac{2t}{3}+\dfrac{\pi}{3}\right)$ (m). Trong $5$ (giây) đầu tiên, thời điểm nào thì chất điểm nằm xa $O$ nhất? (kết quả làm tròn đến hàng phần trăm)
	\shortans{$3,14$}
	\loigiai{Ta mô tả di chuyển của chất điểm $M$ như hình dưới với điểm cách xa $O$ nhất là $A,B$.
		\begin{center}
			\begin{tikzpicture}
				\draw[line width = 1pt] (-7.5,0)--(0,0);
				\draw (-7.5,0.2) node[above]{$A$};
				\draw(-7.5,0) circle (0.05);
				\draw (-3.5,0.2) node[above]{$O$};
				\draw(-3.5,0) circle (0.05);
				\draw (-1.5,0.2) node[above]{$M$};
				\draw(-1.5,0) circle (0.05);
				\draw (0,0.2) node[above]{$B$};
				\draw(0,0) circle (0.05);
			\end{tikzpicture}
		\end{center}
		Do $-3 \leq 3 \cos \left(\dfrac{2 t}{3}+\dfrac{\pi}{3}\right) \leq 3, \forall t$.\\
		Dấu bằng xảy $\mathrm{ra} \Leftrightarrow \cos \left(\dfrac{2 t}{3}+\dfrac{\pi}{3}\right)= \pm 1 \Leftrightarrow \sin \left(\dfrac{2 t}{3}+\dfrac{\pi}{3}\right)=0 \Leftrightarrow \dfrac{2 t}{3}+\dfrac{\pi}{3}=k \pi, k \in \mathbb{Z}$ $\Leftrightarrow t=-\dfrac{\pi}{2}+\dfrac{k 3 \pi}{2}(k \in \mathbb{Z})$.\\
		Với $0 \leq t \leq 5 \Rightarrow 0 \leq-\dfrac{\pi}{2}+\dfrac{k 3 \pi}{2} \leq 5 \Leftrightarrow \dfrac{1}{3} \leq k<1,4$ và $k \in \mathbb{Z}$ nên $k=1$ suy ra $t=\pi$.\\
		Vậy trong $5$ (giây) đầu tiên, thời điềm $t=\pi\approx 3,14$ ( giây) thì chất điểm nằm xa $O$ nhất.}
\end{ex}

% #3873 | Archive/SP_KhoLatex/SP_Đề Theo Bài-Khối 11/1-B3.tex | mã gốc 1D1V1-1
\begin{ex}%[1D1V1-1]
	Hằng ngày, Mặt trời chiếu sáng, bóng của một tòa chung cư cao $40$m in trên mặt đất, độ dài bóng của tòa nhà này được tính bằng công thức $S(t)=40\cot\left |\dfrac{\pi}{12}t \right | $, ở đó $S$ được tính bằng mét, còn $t$ là số giờ tính từ $6$ giờ sáng đến $18$ giờ chiều. Hãy cho biết có bao nhiêu giá trị của $t$ thỏa điều kiện bóng của tòa nhà bằng chiều cao của tòa nhà?
	\shortans{$2$}
	\loigiai{Độ dài bóng của tòa nhà bằng chiều cao tòa nhà khi: \\
		$S(t)=40\Leftrightarrow 40\cot\left |\dfrac{\pi}{12}t \right |=40 \Leftrightarrow \cot \dfrac{\pi}{12}t=\pm 1$
		$\Leftrightarrow\dfrac{\pi}{12}t=\pm \dfrac{\pi}{4} +k \pi \Leftrightarrow t=\pm 3+12k \quad (k\in \mathbb{Z})$.\\
		\text { Vì } $0 \leq t \leq 12 \text { nên }\left[\begin{array}{l}
			0 \leq 3+12 k \leq 12 \\
			0 \leq-3+12 k \leq 12
		\end{array}\right.$ \\
		Vì $ \Leftrightarrow\left[\begin{array}{l}
			-0,25 \leq k \leq 0,75 \\
			0,25 \leq k \leq 1,25
		\end{array}(k \in \mathbb{Z})\right. \\
		\Rightarrow k \in\{1 ; 2\} \Rightarrow t=3 \text { hoặc } t=9$.}
\end{ex}

% #3874 | Archive/SP_KhoLatex/SP_Đề Theo Bài-Khối 11/1-B3.tex | mã gốc 1D1V1-1
\begin{ex}%[1D1V1-1]
	Một con lắc lò xo dao động điều hòa quanh vị trí cân bằng theo phương trình $y=25 \sin 4 \pi t$, trong đó $y$ được tính bằng centimét, còn $t$ được tính bằng giây. Hãy cho biết tần số giao động của con lắc lò xo, tức là số lần dao động trong một giây?
	\shortans{$2$}
	\loigiai{
		Hàm số $y=25 \sin 4 \pi t$ tuần hoàn với chu kỳ $T=\dfrac{2\pi}{4\pi}=\dfrac{1}{2}$ (giây).\\
		Vì chu kỳ dao động của con lắc là $T=\dfrac{1}{2}$ (giây) nên trong một giây con lắc dao động $2$ lần.}
\end{ex}

% #3895 | Archive/SP_KhoLatex/SP_Đề Theo Bài-Khối 11/1-B4.tex | mã gốc 1D1H5-5
\begin{ex}%[1D1H5-5]
	Nghiệm của phương trình $\cos2x-\sin x=0$ là
	\haicot
	{$x=\dfrac{\pi}{6}+\dfrac{k2\pi}{3}; \; \; x=\dfrac{\pi}{2}+k2\pi \; \; \left(k\in \mathbb{Z}\right)$}
	{$x=\dfrac{\pi}{6}+k2\pi; \; \; x=-\dfrac{\pi}{2}+k2\pi \; \; \left(k\in \mathbb{Z}\right)$}
	{\True $x=\dfrac{\pi}{6}+\dfrac{k2\pi}{3}; \; \; x=-\dfrac{\pi}{2}+k2\pi \; \left(k\in \mathbb{Z}\right)$}
	{$x=\dfrac{\pi}{3}+\dfrac{k2\pi}{3}; \; \; x=-\dfrac{\pi}{2}+k2\pi \; \left(k\in \mathbb{Z}\right)$}
	\loigiai{
		Ta có
		\begin{eqnarray*}
			& &\cos2x-\sin x=0\Leftrightarrow \cos 2x=\sin x\Leftrightarrow \cos2x=\cos\left(\dfrac{\pi}{2}-x\right)\\
			&\Leftrightarrow& \hoac{&{2x=\dfrac{\pi}{2}-x+k2\pi} \\&{2x=-\dfrac{\pi}{2}+x+k2\pi}&}\Leftrightarrow \hoac{&{x=\dfrac{\pi}{6}+\dfrac{k2\pi}{3}} \\&{x=-\dfrac{\pi}{2}+k2\pi}&} \left(k\in \mathbb{Z}\right).
		\end{eqnarray*}
	}
\end{ex}

% #3897 | Archive/SP_KhoLatex/SP_Đề Theo Bài-Khối 11/1-B4.tex | mã gốc 1D1C5-6
\begin{ex}%[1D1C5-6]
	Số giờ nắng gắt của tỉnh Quảng Nam ở vĩ độ $40^\circ$ bắc trong ngày thứ $n$ của một năm không nhuận được cho bởi hàm số: $f(n)=3\sin \left[\dfrac{{\it \pi}}{182} \left(n-80\right)\right]+12$, $n\in \mathbb{Z}$ và $0< n\le 365$. Vào ngày thứ bao nhiêu trong năm đó thì tỉnh Quảng Nam chịu nhiều giờ nắng gắt nhất?
	\shortans{$171$}
	\loigiai{
		Số giờ nắng gắt trong ngày thứ $n$ được tính bởi công thức \[f(n)=3\sin \left[\dfrac{{\it \pi}}{182} \left(n-80\right)\right]+12.\]
		Vậy Tỉnh Quảng Nam chịu nhiều giờ nắng gắt nhất nghĩa là $f(n)$ đạt giá trị lớn nhất.\\
		Ta có $-1\le \sin \left[\dfrac{\pi}{182} \left(n-80\right)\right]\le 1, \forall n$ $\Rightarrow f(n)\le 15$.\\
		Suy ra 
		\begin{eqnarray*}
			\max f(n)=15&\Leftrightarrow& \sin \left[\dfrac{\pi}{182} \left(n-80\right)\right]=1\\
			&\Leftrightarrow& \dfrac{\pi}{182} \left(n-80\right)=\dfrac{\pi}{2}+k2\pi \\
			&\Leftrightarrow& n=171+k\cdot 364,\forall k\in {\bf \mathbb{Z}}.
		\end{eqnarray*}
		Mà $0< n\le 365$ nên $n=171$.
	}
\end{ex}

% #3898 | Archive/SP_KhoLatex/SP_Đề Theo Bài-Khối 11/1-B4.tex | mã gốc 1D1C5-6
\begin{ex}%[1D1C5-6]
	\immini{Giả sử vật giao động điều hòa xung quanh vị trí cân bằng theo phương trình $x=2\cos \left(5t-\dfrac{\pi}{6} \right)$. Ở đây, thời gian $t$ tính bằng giây và quãng đường $x$ tính bằng cm. Hãy cho biết trong khoảng thời gian từ 0đến 6giây, vật đi qua vị trí cân bằng bao nhiêu lần? (hình ảnh minh họa vị trí cân bằng)}
	{\definecolor{darkpastelgreen}{rgb}{0.01, 0.75, 0.24}
		\definecolor{almond}{rgb}{0.94, 0.87, 0.8}
		\definecolor{arsenic}{rgb}{0.23, 0.27, 0.29}
		\definecolor{eggplant}{rgb}{0.38, 0.25, 0.32}
		\definecolor{crimson}{rgb}{0.86, 0.08, 0.24}
		\begin{tikzpicture}[line join=round, line cap=round,scale=.5,transform shape]
			\clip (-2.5,-2.5) rectangle (4.5,4.2);
			\tikzset{q_h/.pic={
					% Tre
					\draw[fill=darkpastelgreen] (-.05,4)--(.9,-1.8)
					..controls +(30:.05) and +(-30:.05) ..(1,-1.75)--(.05,4);
					\draw[fill=darkpastelgreen] (-.1,4)--(.6,-1.4)
					..controls +(30:.05) and +(-30:.05) ..(.5,-1.4)--(-.2,4)
					;
					%----------------------
					\def\M{ 
						(-.2,1.76)
						..controls +(-80:.1) and +(80:.1) ..(-.18,1.6)
						..controls +(-40:.1) and +(30:.06) ..(-.18,1.5)--(-.17,1.46)--(-.22,1.44)--(-.19,1.43)--(-.2,1.43)
						..controls +(-80:.1) and +(-12:.05) ..(-.3,1.35)--(-.35,1.28)--(-.55,1.35)
						..controls +(60:.1) and +(-90:.05) ..(-.45,1.5)
						..controls +(100:.1) and +(140:.05) ..(-.38,1.55)
						..controls +(70:.1) and +(-140:.05) ..cycle
						;
					}
					
					\draw[black]\M;
					\fill[almond] \M;
					\draw (-.2,1.76)
					..controls +(-140:.1) and +(50:.05) ..(-.3,1.6)
					(-.2,1.76)
					..controls +(-120:.1) and +(40:.05) ..(-.25,1.6)
					(-.28,1.578)
					..controls +(-20:.02) and +(-150:.02) ..(-.22,1.58)
					(-.25,1.568)--(-.23,1.567);
					%------------------Áo
					\def\A3{ 
						(-.1,.9)
						..controls +(-30:.05) and +(110:0) ..(.18,1)
						--(.19,.95)
						..controls +(-140:.05) and +(30:.05) ..(-.1,.7)--cycle
						;
					}
					\fill[arsenic] \A3;
					\draw[black]\A3;
					\def\A2{ 
						(-.35,1.28)
						..controls +(-30:.05) and +(110:.05) ..(-.25,1.18)
						..controls +(-70:.05) and +(90:.02) ..(-.05,.9)
						..controls +(-90:.02) and +(50:.05) ..(-.05,.4)
						..controls +(-130:.02) and +(50:.05) ..(.6,-.8)
						..controls +(-130:.02) and +(50:.05) ..(.9,-1.8)
						..controls +(-130:.2) and +(30:.3) ..(-.5,-2.3)
						..controls +(130:.2) and +(-90:.1) ..(-.2,-1.5)
						..controls +(130:.2) and +(-140:.1) ..cycle
						;
					}
					\fill[arsenic] \A2;
					\draw[black]\A2;
					\def\A{ 
						(-.55,1.35)
						..controls +(-150:.2) and +(120:.1) ..(-.45,0)
						..controls +(-60:.1) and +(50:.1) ..(-.3,-1.6)
						..controls +(-10:.1) and +(-120:.1) ..(-.15,-1.5)
						..controls +(60:.1) and +(-120:.1) ..(0,-1.45)
						..controls +(120:.3) and +(-50:.1) ..(-.35,.77)
						..controls +(-50:.1) and +(130:.1) ..(-.15,.48)
						..controls +(-50:.1) and +(-130:.1) ..(.43,.93)
						..controls +(-130:.05) and +(-40:.05) ..(.3,1.05)
						..controls +(-130:.05) and +(40:.05) ..(-.05,.75)
						..controls +(-140:.05) and +(-40:.05) ..(-.3,1.15)
						..controls +(100:.05) and +(-40:.05) ..(-.35,1.28)--cycle
						;
					}
					\fill[eggplant] \A;
					\draw[black]\A;
					%-------------------------------
					%nón
					\def\N{ 
						(-.14,1.95)
						..controls +(-80:.1) and +(80:.1) ..(-.17,1.7)
						..controls +(110:.02) and +(-40:.02) ..(-.2,1.76)
						..controls +(-110:.1) and +(50:.4) ..(-.48,1.45)--(-.45,1.35)
						..controls +(-150:.05) and +(20:.05) ..(-.7,1.32)
						..controls +(50:.05) and +(-120:.05) ..(-.65,1.45)
						..controls +(160:.02) and +(-50:.02) ..(-.7,1.48)
						..controls +(160:.02) and +(-50:.02) ..(-.8,1.6)--(-.75,1.63)
						..controls +(-20:.06) and +(130:.04) ..(-.6,1.47)
						..controls +(130:.2) and +(-90:.1) ..(-.7,1.78)
						..controls +(90:.1) and +(170:.1) ..(-.4,1.9)
						..controls +(20:.1) and +(160:.1) ..cycle
						;
					}
					\fill[arsenic] \N;
					\draw[black]\N;
					\draw (-.14,1.95)
					..controls +(-100:.05) and +(40:.06) ..(-.2,1.76)
					(-.5,1.55)
					..controls +(30:.1) and +(-100:.06) ..(-.3,1.76)
					(-.6,1.7)
					..controls +(-55:.1) and +(120:.06) ..(-.55,1.5)
					;
					
					\def\T1{ 
						(.3,1.03)
						..controls +(60:.05) and +(-120:.05) ..(.35,1.15)
						..controls +(60:.05) and +(-140:.05) ..(.45,1.26)
						..controls +(30:.05) and +(60:.04) ..(.5,1.22)
						..controls +(30:.03) and +(50:.03) ..(.54,1.18)
						..controls +(30:.03) and +(50:.03) ..(.56,1.12)
						..controls +(30:.03) and +(50:.03) ..(.58,1.07)
						..controls +(160:.03) and +(40:.03) ..(.49,1)
						..controls +(-140:.03) and +(40:.03) ..(.41,.92)--cycle
						;
					}
					\fill[almond] \T1;
					\draw[black]\T1;
					
					\fill[crimson](-.27,.4)
					..controls +(-80:.6) and +(-80:.1) ..(-.5,-1.9)--(-.45,-1.85)
					..controls +(90:.1) and +(-80:1) ..(-.28,.48)--cycle;
					
					\fill[crimson](-.27,.4)
					..controls +(-80:.6) and +(-80:.1) ..(-.2,-1.9)--(-.15,-1.85)
					..controls +(90:.1) and +(-80:1) ..(-.28,.48)--cycle;
					
			}}
			\tikzset{no/.pic={
					\draw[crimson] (1.8,-1)
					..controls +(150:.4) and +(-120:.5) ..(1.7,-1)
					..controls +(30:.4) and +(-40:.3) ..(1.75,-1)
					..controls +(-30:.2) and +(40:.2) ..(1.7,-1.4)
					(1.75,-1)
					..controls +(-50:.1) and +(20:.1) ..(1.7,-1.2)
					
					;
			}}
			\draw[thick] (4,-1.25) arc(-60:-120:6cm);
			\draw[dashed,thick] (4,-1.25)--(1.3,4)--(-2,-1.25) (1.3,4)--(1.3,-2);
			\node at (1.5,.5) [rotate=90]{\tiny Vị trí cân bằng};
			\draw[thick,<->] (.35,-2)--(1.3,-2);
			\node at (.85,-2) [above]{\tiny $x$};
			\fill (4,-1.25) circle (0.05);
			\fill (-2,-1.25) circle (0.05);
			\fill (1.3,4) circle (0.05);
			\path
			(0,0)pic[scale=1,rotate=-20]{q_h}(-.88,1.48)pic[scale=1, scale=.7,rotate=-20]{no};
			
			
	\end{tikzpicture}}
	\shortans{$9$}
	\loigiai{
		Khi vật ở vị trí cân bằng thì $x=0$
		\begin{eqnarray*}
			&\Leftrightarrow& 3\cos \left(5t-\dfrac{\pi}{6} \right)=0\Leftrightarrow \cos\left(5t-\dfrac{\pi}{6} \right)=0\\
			&\Leftrightarrow& 5t-\dfrac{\pi}{6}=\dfrac{\pi}{2}+k\pi, k\in {\bf \mathbb{Z}}\Leftrightarrow t=\dfrac{2\pi}{15}+\dfrac{k\pi}{5}, k\in {\bf \mathbb{Z}}.
		\end{eqnarray*}
		Mà theo giả thiết 
		\begin{eqnarray*}
			0< t < 6&\Leftrightarrow& 0< \dfrac{2\pi}{15}+\dfrac{k\pi}{5} < 6\\
			&\Leftrightarrow&-\dfrac{2}{3} < k < \dfrac{30}{\pi}-\dfrac{2}{3}\\
			k\in {\bf \mathbb{Z}}&\Rightarrow& k\in \left\{0;1;2;3;4;5;6;7;8\right\}.
		\end{eqnarray*}
	}
\end{ex}

% #4354 | Archive/SP-4Phan-Dot3/Data/CD_1_OTHK2_De06.tex | mã gốc 1D1V2-3
\begin{ex}%[Làm đề GHK, HK, 4 phần - 2025]%[Huỳnh Xuân Tín]%[1D1V2-3]
	Hộp $I$ chứa $4$ viên bi trắng, $5$ viên bi đỏ và $6$ viên bi xanh. Hộp $II$ có $7$ viên bi trắng, $6$ viên bi đỏ và $5$ viên bi xanh. Lấy ngẫu nhiên từ mỗi hộp $1$ viên bi. Xác suất để hai viên bi được lấy ra có cùng màu là $\dfrac{a}{b}$ với $\dfrac{a}{b}$ là phân số tối giản và $a, b\in\mathbb{Z}$. Tính $T = a + b$.
	\shortans{$179$}
	\loigiai{
		Gọi $A$ là biến cố: \lq\lq Hai viên bi được lấy ra có cùng màu\rq\rq.\\
		Xét các biến cố sau
		\begin{itemize}
			\item $A_1\colon$ \lq\lq Hai viên bi lấy ra màu trắng\rq\rq. Ta được $\mathrm{P}(A_1)=\dfrac{4}{15}\cdot \dfrac{7}{18}$.
			\item  $A_2\colon$ \lq\lq Hai viên bi lấy ra màu đỏ\rq\rq. Ta được $\mathrm{P}(A_2)=\dfrac{5}{15}\cdot \dfrac{6}{18}$.
			\item $A_3\colon$ \lq\lq Hai viên bi lấy ra màu xanh\rq\rq. Ta được $\mathrm{P}(A_3)=\dfrac{6}{15}\cdot \dfrac{5}{18}$.
		\end{itemize}
		Lúc đó $A=A_1\cup A_2\cup A_3$ và $A_1$, $A_2$, $A_3$ là các biến cố xung khắc nên
		\begin{eqnarray*}
			\mathrm{P}(A) = \mathrm{P}(A_1) + \mathrm{P}(A_2) + \mathrm{P}(A_3) = \dfrac{44}{135}.
		\end{eqnarray*}
		Do đó $\heva{&a = 44\\&b = 135.}$\\
		Vậy $T = a + b = 44 + 135 = 179$.
	}
\end{ex}

% #4453 | Archive/2025-Sangtac/Khoi10/Data/0-CTST/CTST_0-OTHK1_De05.tex | mã gốc 1D1V5-3
\begin{ex}%[Câu 4]%[1D1V5-3]
	Mẫu số liệu dưới đây thống kê thời gian chờ xe buýt (đơn vị: phút) của $10$ học sinh ở cùng một bến.
	$$1;4;5;6;6;8;10;11;12;25$$
	\choiceTF
	{\True Số trung bình cộng của mẫu số liệu là $\overline{x}=8{,}8$ (phút)}
	{Khoảng tứ phân vị của mẫu số liệu là $\Delta_Q=5$ (phút)}
	{Độ lệch chuẩn của mẫu số liệu là $s \approx 5{,}27$ (phút)}
	{\True $25$ là giá trị bất thường của mẫu số liệu}
	\loigiai{
		\begin{itemchoice}
			\itemch Đúng. Vì $\overline{x}=\dfrac{1+4+5+6+6+8+10+11+12+25}{10}=8{,}8$ (phút).
			\itemch Đúng. Trung vị của mẫu số liệu là $\dfrac{6+8}{2}=7$ (phút).\\
			Trung vị của dãy $1;4;5;6;6$ là $5$. Trung vị của dãy $8;10;11;12;25$ là $11$.\\
			Vậy tứ phân vị của mẫu số liệu là $Q_1=5$ (phút), $Q_2=7$ (phút), $Q_3=11$ (phút).\\
			Khoảng biến thiên của mẫu số liệu là $R=25-1=24$ (phút).\\
			Khoảng tứ phân vị của mẫu số liệu là $\Delta_Q=11-5=6$ (phút).
			\itemch Sai. Ta có $(1-8{,}8)^2+(4-8{,}8)^2+(5-8{,}8)^2+(6-8{,}8)^2\cdot2+(8-8{,8})^2+(10-8{,}8)^2+(11-8{,}8)^2+(12-8{,}8)^2+(25-8{,}8)^2=393{,}6$.\\
			Suy ra phương sai của mẫu số liệu là $s^2=\dfrac{393{,6}}{10}=39{,}36$.\\
			Độ lệch chuẩn của mẫu số liệu là $s=\sqrt{39{,}36}\approx6{,}27$ (phút).
			\itemch Đúng. Ta có $Q_1-\dfrac{3}{2}\Delta_Q=5-\dfrac{3}{2}\cdot6=-4$, $Q_3+\dfrac{3}{2}\cdot\Delta_Q=11+\dfrac{3}{2}=20$. Vì $25>20$ nên $25$ là giá trị bất thường của mẫu số liệu.
		\end{itemchoice}
	}
\end{ex}

% #4454 | Archive/2025-Sangtac/Khoi10/Data/0-CTST/CTST_0-OTHK1_De05.tex | mã gốc 1D1H2-3
\begin{ex}%[Câu 1]%[1D1H2-3]
	Một người nông dân thả $1\,000$ con cá giống vào hồ nuôi vừa mới đào. Biết rằng sau mỗi năm thì số lượng cá trong hồ tăng thêm $x$ lần số lượng cá ban đầu và $x$ không đổi. Bằng cách thay đổi kĩ thuật nuôi và thức ăn cho cá. Hỏi sau hai năm để số cá trong hồ là $9000$ con thì tốc độ tăng trưởng cá trong hồ là bao nhiêu? Biết tốc độ tăng mỗi năm là không đổi
	\shortans[0]{$5$}
	\loigiai{
		Sau một năm số lượng các trong hồ là $1\,000+1\,000x=1\,000(1+x)$ (con).\\
		Sau hai năm số lượng cá trong hồ là $1\,000(1+x)+1\,000(1+x)x=1\,000(1+x)^2$ (con).\\
		Điều kiện $x>0$. Để số lượng cá trong hồ sau hai năm là $9000$ thì ta có:\\
		$\Rightarrow 1\,000(1+x)^2=9\,000 \Leftrightarrow (1+x)^2=9 \Leftrightarrow \hoac{&x=5\text{ (nhận)}\\&x=-7\text{ (loại)}.}$\\
		Vậy tốc độ tăng thêm số lượng cá trong hồ sau mỗi năm là $5$ lần số lượng cá ban đầu.
	}
\end{ex}

% #4459 | Archive/2025-Sangtac/Khoi10/Data/0-CTST/CTST_0-OTHK1_De05.tex | mã gốc 1D1C4-6
\begin{ex}%[Câu 6]%[1D1C4-6]
	Bảng số liệu sau thống kê nhiệt độ tại Thành phố Hồ Chí Minh trong một lần đo vào một ngày của năm 2021 
	\begin{center}
		\begin{tabular}{|c|c|c|c|c|c|c|c|c|}
	\hline
	Giờ đo			& $1$h & $4$h & $7$h & $10$h & $13$h & $16$h & $19$h & $22$h \\
	\hline
	Nhiệt độ (độ C)	& $27$ & $26$ & $28$ & $32$ & $34$ & $35$ & $30$ & $28$ \\
	\hline
	\end{tabular}
	\end{center}
	Tìm độ lệch chuẩn của mẫu số liệu đã cho (làm tròn kết quả đến hàng phần trăm).
	\shortans[0]{$3{,}12$}
	\loigiai{
		Số trung bình là $\overline{x}=\dfrac{27+26+\ldots+30+28}{8}=30$.\\
		Phương sai $s^2=\dfrac{1}{8}\left[(x_1-\overline{x})^2+(x_2-\overline{x})^2+\ldots+(x_8-\overline{x})^2\right]$\\$s^2=\dfrac{1}{8}\left[(27-30)^2+(26-30)^2+\ldots+(28-30)^2\right]=9{,}75$.\\
		Độ lệch chuẩn $s=\sqrt{s^2}\approx3{,}12$.
	}
\end{ex}

% #4697 | Archive/2025-Sangtac/Khoi11/Data/1-CD/CD_1-OTHK1_De05.tex | mã gốc 1D1K5-6
\begin{ex}%[1D1K5-6]
	Chiều cao $h$ (m) của một cabin trên vòng quay vào thời điểm $t$ giây sau khi bắt đầu chuyển động được cho bởi công thức $h(t)=30 + 20\sin \left(\dfrac{\pi}{25}t + \dfrac{\pi}{3}\right)$. Sau bao nhiêu giây thì cabin đạt độ cao $40$ m lần đầu tiên?
	\shortans{$12{,}5$}
	\begin{center}
			\begin{tikzpicture}[scale=0.4,font=\small,line join=round,line cap=round,>=stealth]
			\tikzset{declare function={g=0;}}
			\coordinate (O) at (0,0);
		%	\coordinate[label=above right:$A$] (A) at (0:4.4);
		%	\coordinate[label=above right:$M$] (M) at (30+g:4.4);
			\pgfmathsetmacro{\gx}{60-g};
			\coordinate (v) at ($(M)+(120+g:3.3)$);
			\coordinate (vx) at ($(M)+(180:3.3*cos \gx)$);
			\draw[thick] (O) circle(4.4) (O) circle(3.53);
			\foreach \t in {0,1,...,11}{
				\draw[very thick,fill] (O)--(g+30*\t:4.4)circle(0.06);
				\begin{scope}[shift={(g+30*\t:4.4)}]
					\draw[fill=teal] (-0.27,0)--(0.27,0)--(0.36,-0.18)--(-0.36,-0.18)--cycle
					(0.32,-0.18)--(0.32,-0.54) (-0.32,-0.18)--(-0.32,-0.54)
					(-0.4,-0.54)--(0.4,-0.54)--(0.27,-1)--(-0.27,-1)--cycle;
			\end{scope}}
			\draw[fill=gray!70!yellow] (0.2,0)--(2.35,-5.53)--(3.65,-5.53)--(3.65,-6.02)--(1.95,-6.02)--(0,-0.85)--(-1.95,-6.02)--(-3.65,-6.02)--(-3.65,-5.53)--(-2.35,-5.53)--(-0.2,0)--cycle;
			\draw[fill=gray!10!yellow] (O) circle(0.34);
			\draw[fill] (-5,-5.85)--(0,-5.85) circle(0.08)--(6.4,-5.85);
		%	\draw[->,red] (-5,0)--(6.7,0) node %[below]{\color{black}$x$};
%			\draw[->,red] (M)--(v) node [left=-2pt]{\color{black}$\vec{v}$};
%			\draw[->,red] (M)--(vx) node [below left=-5pt]{\color{black}$\vec{v}_{\!\!x}$};
%			\draw[red] (O)--(M);
%			\draw[dashed,red] (0,0)--(0,-5.85) (v)--(vx);
%			\fill (0,0) circle(0.1);
%			\path pic[draw=red,angle radius=5]{right angle=v--M--O} 
%			(O) pic[draw,->,angle radius=16]{angle=A--O--M} node[shift={(14:23pt)}]{$\alpha$};
		%	\node at (O)[below]{$O$};
		\end{tikzpicture}
	\end{center}
	\loigiai{
		Để cabin đạt độ cao $40$ m thì 
		\begin{eqnarray*}
			& & h(t)=40 \Leftrightarrow 20\sin \left(\dfrac{\pi}{25}t + \dfrac{\pi}{3}\right) =10 \Leftrightarrow  \sin \left(\dfrac{\pi}{25}t + \dfrac{\pi}{3}\right) = \dfrac{1}{2}\\
			&\Leftrightarrow & \hoac{&  \dfrac{\pi}{25}t + \dfrac{\pi}{3} = \dfrac{\pi}{6} + k2\pi \\& \dfrac{\pi}{25}t + \dfrac{\pi}{3} = \dfrac{5\pi}{6} + k2\pi} \Leftrightarrow \hoac{& t = -\dfrac{25}{6}+50k \\ & t = \dfrac{25}{2} + 50k}, k \in \mathbb{Z}.
		\end{eqnarray*}
		Dễ thấy, giá trị $t>0$ nhỏ nhất thoả mãn là $t = 12{,}5$ (giây).
	}
\end{ex}

% #4707 | Archive/2025-Sangtac/Khoi11/Data/1-CD/CD_1-OTHK1_De10.tex | mã gốc 1D1V4-6
\begin{ex}%[1D1V4-6]
	Số giờ có ánh sáng mặt trời của một thành phố tại vĩ độ $40^\circ$ Bắc trong ngày thứ $t$ của năm được cho bởi hàm số $d(t) = 4 \sin\left[\dfrac{\pi}{182}(t - 88)\right] + 10$ (giờ), với $t$ là số ngày trong năm từ $1$ đến $365$. Ngày đầu tiên tạo $10$ giờ ánh sáng mặt trời là ngày.
	\shortans[0]{$88$}
	\loigiai{
		Với $t\in[1;365]$ ta có 
		\begin{eqnarray*}
			& & \sin\left[\dfrac{\pi}{182}(t - 88)\right] + 10 = 10\\
			&\Leftrightarrow & \sin\left[\dfrac{\pi}{182}(t - 88)\right] = 0\\
			&\Leftrightarrow & \dfrac{\pi}{182}(t - 88)=k\pi\\
			&\Leftrightarrow & t=88+182k\\
			&\Leftrightarrow & t \in \{88, 270\}.
		\end{eqnarray*}
		Số ngày tạo $10$ giờ ánh sáng mặt trời ngày đầu tiên của năm là $88$.}
\end{ex}

% #4744 | Archive/2025-Sangtac/Khoi11/Data/1-CD/CD_1-OTGK1_De08.tex | mã gốc 1D1V1-6
\begin{ex}%[1D1V1-6]
	Một bánh xe có đường kính kể cả lốp xe là $55$ cm. Nếu xe chạy với tốc độ $50$ km$/$h thì trong một giây bánh xe quay được bao nhiêu vòng? (Kết quả được làm tròn đến hàng phần trăm).
	\shortans[0]{$8{,}04$}
	\loigiai{
	Tốc độ xe là $50$ km$/$h $=\dfrac{50\cdot 100000}{3600}$ cm$/$s $=\dfrac{12500}{9}$ cm$/$s.\\
	Mỗi vòng bánh xe có chiều dài $2\pi R=2\pi \cdot \dfrac{55}{2}=55\pi$ (cm).\\
	Vậy mỗi giây thì bánh xe lăn được số vòng là $\dfrac{12500}{9}:(55\pi)\approx 8{,}04$ (vòng).
}
\end{ex}

% #4828 | Archive/2025-Sangtac/Khoi11/Data/1-KNTT/KNTT_1-OTGK1_De02.tex | mã gốc 1D1V4-8
\begin{ex}%[1D1V4-8]
	Hằng ngày, mực nước của một con kênh lên xuống theo thuỷ triều. Độ sâu $h$ (m) của mực nước trong kênh tính theo thời gian $t$ (giờ) trong một ngày $(0 \leq t<24)$ cho bởi công thức $h=\cos \left(\dfrac{\pi t}{6}+2 \pi\right)+5$. Hỏi trong ngày vị trí nước xuống mức thấp nhất là bao nhiêu mét?
	
	\shortans[]{4}
	\loigiai{
		Với mọi $0 \leq t<24$ ta có
		$$-1\leq \cos \left(\dfrac{\pi t}{6}+2 \pi\right)\leq 1 \Leftrightarrow 4\leq \cos \left(\dfrac{\pi t}{6}+2 \pi\right)+5\leq 6.$$
			Mực nước thấp nhất của kênh $h=4$ (m) khi
		\begin{align*}
			\cos \left(\dfrac{\pi t}{6}+2 \pi\right)=-1 \Leftrightarrow \dfrac{\pi t}{6}+2\pi=\pi+k2\pi \Leftrightarrow t=-6+12k, k\in \mathbb{Z}.
		\end{align*} 
		Do $0\le t< 24$ nên $k=1$, $t=6$ và  $k=2$, $t=18$.\\
		Vậy lúc $6$ và $18$ giờ mực nước kênh là thấp nhất và $h=4$ (m).
	}
\end{ex}

% #4869 | Archive/2025-Sangtac/Khoi11/Data/1-CTST/CTST_1-OTGK1_De08.tex | mã gốc 1D1V5-6
\begin{ex}%[1D1V5-6]
Một quả đạn pháo được bắn ra khỏi nòng pháo với vận tốc ban đầu $ v_0=500$ m/s hợp với phương ngang một góc $ \alpha $. Trong Vật lí, ta biết rằng, nếu bỏ qua sức cản của không khí và coi quả đạn pháo được bắn ra từ mặt đất thì quỹ đạo của quả đạn tuân theo phương trình $ y=\dfrac{-g}{2v_0^2\cos^2\alpha}\cdot x^2+x\tan \alpha $, trong đó $g = 10$ m/s$ ^2 $ là gia tốc trọng trường. Góc bắn $ \alpha $ để quả đạn bay xa nhất là bao nhiêu độ?
\shortans[0]{45}
\loigiai{
Ta có \begin{eqnarray*}
y=0&\Leftrightarrow& \dfrac{-g}{2v_0^2\cos^2\alpha}\cdot x^2+x\tan \alpha=0\\
&\Leftrightarrow & x\left(\dfrac{-g}{2v_0^2\cos^2\alpha}\cdot x+\tan \alpha \right)=0 \\
&\Leftrightarrow & \hoac{&x=0\\&x=\dfrac{v_0^2\sin2\alpha}{g}= 25~000\sin2\alpha.}
\end{eqnarray*}
Vậy tầm bay xa mà quả đạn đạt tới là $25~000\sin2\alpha$ m.\\
Vậy đạn bay xa nhất khi $ \sin2\alpha=1\Leftrightarrow 2\alpha=\dfrac{\pi}{2}+k2\pi\Leftrightarrow \alpha =\dfrac{\pi}{4}+k\pi $.\\
Vì $ 0\le\alpha\le \dfrac{\pi}{2} $ nên $ \alpha=\dfrac{\pi}{4} $.\\
Vậy góc bắn bằng $ 45^\circ $ thì đạn bay xa nhất.
}
\end{ex}

% #4898 | Archive/2025-Sangtac/Khoi11/Data/1-CTST/CTST_1-OTHK1_De02.tex | mã gốc 1D1V5-3
\begin{ex}%[1D1V5-3]
Hằng ngày mực nước của con kênh lên xuống theo thủy triều. Độ sâu $h$ (mét) của mực nước trong kênh được tính tại thời điểm $t$ (giờ) trong một ngày ($t > 0$) bởi công thức $h=4\sin (\dfrac{\pi t}{8}+\dfrac{5\pi}{8})+16$. Mực nước của kênh cao nhất khi $t$ bằng bao nhiêu?
\shortans[6]{$15$}
\loigiai{
Mực nước của kênh cao nhất khi $h$ lớn nhất
$$ \sin (\dfrac{\pi t}{8}+\dfrac{5\pi}{8})=1\Leftrightarrow \dfrac{\pi}{8}t+\dfrac{5\pi}{8}=\dfrac{\pi}{2}+k2\pi \Leftrightarrow t=-1+16k, k\in\mathbb{Z}.$$
Với $0 < t\le 24$ suy ra $0 <-1+16k\le 24\Leftrightarrow \dfrac{1}{16} < k\le \dfrac{25}{16}$.\\
Do $k\in \mathbb{Z}$ nên $k=1$ thỏa mãn. Khi $k=1$ thì $t=15$ h.
}
\end{ex}

% #4908 | Archive/2025-Sangtac/Khoi11/Data/1-CTST/CTST_1-OTHK1_De03.tex | mã gốc 1D1V4-8
\begin{ex}%[1D1V4-8]
	Hằng ngày, mực nước của một con kênh lên xuống theo thuỷ triều. Độ sâu $h$ (m) của mực nước trong kênh tính theo thời gian $t$ (giờ) trong một ngày $(0 \leq t<24)$ cho bởi công thức $h=\cos \left(\dfrac{\pi t}{6}+2 \pi\right)+5$. Hỏi trong ngày vị trí nước xuống mức thấp nhất trễ nhất là mấy giờ?
		\shortans[0]{$21$}
	\loigiai{
		Với mọi $0 \leq t<24$ ta có
		$$-1\leq \cos \left(\dfrac{\pi t}{6}+2 \pi\right)\leq 1 \Leftrightarrow 4\leq \cos \left(\dfrac{\pi t}{6}+2 \pi\right)+5\leq 6.$$
		Mực nước thấp nhất của kênh $h=4$ (m) khi
		\begin{align*}
			\cos \left(\dfrac{\pi t}{6}+2 \pi\right)=-1 \Leftrightarrow \dfrac{\pi t}{6}+2\pi=-\dfrac{\pi }{2}+k2\pi \Leftrightarrow t=-15+12k, k\in \mathbb{Z}.
		\end{align*} 
		Do $0\le t\le 24$ nên $k=2$, $t=9$ và  $k=3$, $t=21$.\\
		Vậy lúc $9$ và $21$ giờ mực nước kênh là thấp nhất và $h=4$ (m). Vậy mực nước thấp trễ nhất lúc $21$ giờ.
	}
\end{ex}

% #4918 | Archive/2025-Sangtac/Khoi11/Data/1-CTST/CTST_1-OTHK1_De01.tex | mã gốc 1D1V5-3
\begin{ex}%[1D1V5-3]
Giả sử một vật dao động điều hòa xung quanh vị trí cân bằng theo phương trình $x=2\cos\left(5t-\dfrac{\pi}{6} \right)$. Ở đây, thời gian $t$ tính bằng giây và quãng đường $x$ tính bằng centimét. Hãy cho biết trong khoảng thời gian từ $0$ đến $6$ giây, vật đi qua vị trí cân bằng bao nhiêu lần?
\shortans[6]{$9$}
\loigiai{
Vật đi qua vị trí cân bằng thì $x=0$
\begin{eqnarray*}
&\Leftrightarrow&2\cos\left(5t-\dfrac{\pi}{6} \right)=0\\
&\Rightarrow&5t-\dfrac{\pi}{6}=\dfrac{\pi}{2}+k\pi\\
&\Rightarrow&t=\dfrac{2\pi}{15}+k\dfrac{\pi}{5},\,k\in \mathbb{Z}.	
\end{eqnarray*}	
Trong khoảng thời gian từ $0$ đến $6$ giây 
\begin{eqnarray*}
&\Rightarrow&0\le t\le 6\\
&\Leftrightarrow&0\le\dfrac{2\pi}{15}+k\dfrac{\pi}{5}\le 6\\
&\Leftrightarrow &-\dfrac{2}{3}\le k\le 8{,}88\\
&\Rightarrow&k\in \left\lbrace 0;1;\ldots; 8 \right\rbrace.
\end{eqnarray*}	
Vậy trong khoảng thời gian từ $0$ đến $6$ giây, vật đi qua vị trí cân bằng $9$ lần.
}
\end{ex}

% #4979 | Archive/2025-Sangtac/Khoi11/Data/1-CTST/CTST_1-OTGK1_De03.tex | mã gốc 1D1N6-8
\begin{ex}%[1D1N6-8] 
		Hằng ngày mực nước tại một cảng biển lên xuống theo thuỷ triều. Độ sâu $h$ m của mực nước theo thời gian $t$ (giờ) trong một ngày được cho bởi công thức
		\centerline{$h=11+2\sin\left(\dfrac{\pi}{12}t\right) ~\text{với} ~0\leq t\leq 24$.}
		Tìm thời điểm mà mực nước tại cảng là cao nhất.
		\shortans[3]{$6$}
		\loigiai{
			Ta có
			\begin{eqnarray*}
				&&-1\le\sin\left(\dfrac{\pi}{12}t\right)\le 1\\
				&\Leftrightarrow & 11-2\le 11+2\sin\left(\dfrac{\pi}{12}t\right)\le 11+2\\
				&\Leftrightarrow& 9\le h\le 13
			\end{eqnarray*}
			Vậy mực nước tại cảng cao nhất bằng $13$ m khi
			\begin{eqnarray*}
				&&\sin\left(\dfrac{\pi}{12}t\right)=1\\
				&\Leftrightarrow & \left(\dfrac{\pi}{12}t\right)=\dfrac{\pi}{2}+2k\pi\\
				&\Leftrightarrow & t=6+24k ~\left(k\in\mathbb{Z}\right).
			\end{eqnarray*}
			Mà $0\le t\le 24$ nên $t=6$.\\ Vậy thời điểm mà mực nước tại cảng cao nhất là lúc $6$ (giờ).
		}
	\end{ex}

% #5004 | Archive/2025-Sangtac/Khoi11/Data/1-CTST/CTST_1-OTGK1_De05.tex | mã gốc 1D1V4-8
\begin{ex}%[Câu 1]%[1D1V4-8]
	Khi đu quay hoạt động, vận tốc theo phương ngang của một cabin $M$ phụ thuộc vào góc lượng giác $\alpha = (Ox, OM)$ theo hàm số $v_x = 0{,}25\sin\alpha$ (m/s). Vận tốc lớn nhất của cabin là
	\begin{center}
		\begin{tikzpicture}[scale=0.7,font=\small,line join=round,line cap=round,>=stealth]
			\tikzset{declare function={g=0;}}
			\coordinate (O) at (0,0);
			%	\coordinate[label=above right:$A$] (A) at (0:4.4);
			%	\coordinate[label=above right:$M$] (M) at (30+g:4.4);
			\pgfmathsetmacro{\gx}{60-g};
			\coordinate (v) at ($(M)+(120+g:3.3)$);
			\coordinate (vx) at ($(M)+(180:3.3*cos \gx)$);
			\draw[thick] (O) circle(4.4) (O) circle(3.53);
			\foreach \t in {0,1,...,11}{
				\draw[very thick,fill] (O)--(g+30*\t:4.4)circle(0.06);
				\begin{scope}[shift={(g+30*\t:4.4)}]
					\draw[fill=teal] (-0.27,0)--(0.27,0)--(0.36,-0.18)--(-0.36,-0.18)--cycle
					(0.32,-0.18)--(0.32,-0.54) (-0.32,-0.18)--(-0.32,-0.54)
					(-0.4,-0.54)--(0.4,-0.54)--(0.27,-1)--(-0.27,-1)--cycle;
			\end{scope}}
			\draw[fill=gray!70!yellow] (0.2,0)--(2.35,-5.53)--(3.65,-5.53)--(3.65,-6.02)--(1.95,-6.02)--(0,-0.85)--(-1.95,-6.02)--(-3.65,-6.02)--(-3.65,-5.53)--(-2.35,-5.53)--(-0.2,0)--cycle;
			\draw[fill=gray!10!yellow] (O) circle(0.34);
			\draw[fill] (-5,-5.85)--(0,-5.85) circle(0.08)--(6.4,-5.85);
			\draw[->] (-5,0)--(6.7,0) node [below]{\color{black}$x$};
			\path 
			(3.82,2.22) coordinate (M)
			($(M)-(1.8,0)$) coordinate (N)
			($(N)+(0,2.5)$) coordinate (P)
			;
			\fill[black] (M) circle(1pt) node[above right]{$M$};
			\fill[black] (N) circle(0pt) node[below left]{$\overrightarrow{v}_x$};
			\fill[black] (P) circle(0pt) node[above]{$\overrightarrow{v}$};
			\fill[black] (O) circle(2pt);
			\fill[black] (O)+(-90:0.3) circle(0pt) node[below]{$O$};
			
			\draw[->] (M)--(N);
			\draw[->] (M)--(P);
			\draw[dashed] (N)--(P);
			\draw (O)--(M);
			\draw[->] (O)+(0:0.8) arc (0:30:0.8) node[below right]{$\alpha$};
			\draw pic[draw, angle radius = 4mm] {right angle=O--M--P};
		\end{tikzpicture}
	\end{center}
	\shortans{$0{,}25$}
	\loigiai{
		Với mọi  Ta có  $-1\le \sin\alpha \le 1$, $\forall \alpha\in \mathbb{R}$.\\
		Do đó $-0{,}25 \le 0{,}25\sin\alpha \le 0{,}25 \Leftrightarrow -0{,}25 \le v(x) \le 0{,}25$.\\
		Vậy vận tốc lớn nhất của cabin là $0{,}25$.
	}
\end{ex}

% #5014 | Archive/2025-Sangtac/Khoi11/Data/1-CTST/CTST_1-OTGK1_De04.tex | mã gốc 1D1V4-8
\begin{ex}%[1D1V4-8]
	\immini{Khi đu quay hoạt động, vận tốc theo phương ngang của một cabin $M$ phụ thuộc vào góc lượng giác $\alpha = \left(Ox,OM\right)$ theo hàm số $v_x = 0{,}5\sin\alpha$ (m/s). Giá trị lớn nhất của $ v_x $ là}
	{
		\begin{tikzpicture}[scale=0.6,font=\small,line join=round,line cap=round,>=stealth]
			\tikzset{declare function={g=0;}}
			\coordinate (O) at (0,0);
			%	\coordinate[label=above right:$A$] (A) at (0:4.4);
			%	\coordinate[label=above right:$M$] (M) at (30+g:4.4);
			\pgfmathsetmacro{\gx}{60-g};
			\coordinate (v) at ($(M)+(120+g:3.3)$);
			\coordinate (vx) at ($(M)+(180:3.3*cos \gx)$);
			\draw[thick] (O) circle(4.4) (O) circle(3.53);
			\foreach \t in {0,1,...,11}{
				\draw[very thick,fill] (O)--(g+30*\t:4.4)circle(0.06);
				\begin{scope}[shift={(g+30*\t:4.4)}]
					\draw[fill=teal] (-0.27,0)--(0.27,0)--(0.36,-0.18)--(-0.36,-0.18)--cycle
					(0.32,-0.18)--(0.32,-0.54) (-0.32,-0.18)--(-0.32,-0.54)
					(-0.4,-0.54)--(0.4,-0.54)--(0.27,-1)--(-0.27,-1)--cycle;
			\end{scope}}
			\draw[fill=gray!70!yellow] (0.2,0)--(2.35,-5.53)--(3.65,-5.53)--(3.65,-6.02)--(1.95,-6.02)--(0,-0.85)--(-1.95,-6.02)--(-3.65,-6.02)--(-3.65,-5.53)--(-2.35,-5.53)--(-0.2,0)--cycle;
			\draw[fill=gray!10!yellow] (O) circle(0.34);
			\draw[fill] (-5,-5.85)--(0,-5.85) circle(0.08)--(6.4,-5.85);
			\draw[->] (-5,0)--(6.7,0) node [below]{\color{black}$x$};
			\path 
			(3.82,2.22) coordinate (M)
			($(M)-(1.8,0)$) coordinate (N)
			($(N)+(0,2.5)$) coordinate (P)
			;
			\fill[black] (M) circle(1pt) node[above right]{$M$};
			\fill[black] (N) circle(0pt) node[below left]{$\overrightarrow{v}_x$};
			\fill[black] (P) circle(0pt) node[above]{$\overrightarrow{v}$};
			\fill[black] (O) circle(2pt);
			\fill[black] (O)+(-90:0.3) circle(0pt) node[below]{$O$};
			\draw[->] (M)--(N);
			\draw[->] (M)--(P);
			\draw[dashed] (N)--(P);
			\draw (O)--(M);
			\draw[->] (O)+(0:0.8) arc (0:30:0.8) node[below right]{$\alpha$};
			\draw pic[draw, angle radius = 4mm] {right angle=O--M--P};
		\end{tikzpicture}
	}
	\shortans{$ 0{,}5 $}
	\loigiai{ Với mọi $ \alpha $ ta có\\
		$ -1\le \sin\alpha \le 1 \Rightarrow -0{,}5\le 0{,}5\sin\alpha \le 0{,}5 \Rightarrow -0{,}5 \le v_x \le 0{,}5$.\\
		Vậy giá trị lớn nhất của $v_x$ là $0{,}5$ (m/s).
	}
\end{ex}

% #5015 | Archive/2025-Sangtac/Khoi11/Data/1-CTST/CTST_1-OTGK1_De04.tex | mã gốc 1D1V4-8
\begin{ex}%[Câu 19]%[1D1V4-8]
	Cho vận tốc $v$ (cm/s) của một con lắc đơn theo thời gian $t$ (giây) được cho bởi công thức $v=-4\sin\left(1{,}5t+\dfrac{\pi}{3}\right)$ với $0\le t \le 2$. Xác định thời điểm vận tốc con lắc bằng $2$ cm/s. (Làm tròn kết quả đến hàng chục).
	\shortans[0]{$1{,}7$}
	\loigiai{
		Ta có
		\begin{eqnarray*}
			& & -4\sin\left(1{,}5t+\dfrac{\pi}{3}\right)=2\\
			&\Leftrightarrow & \sin\left(1{,}5t+\dfrac{\pi}{3}\right)=\dfrac{-1}{2}\\
			&\Leftrightarrow &	\hoac{&1{,}5t+\dfrac{\pi}{3}=\dfrac{-\pi}{6}+k2\pi \\ & 1{,}5t+\dfrac{\pi}{3}=\dfrac{7\pi}{6}+k2\pi}\\
			&\Leftrightarrow & \hoac{& t=\dfrac{-\pi}{3}+\dfrac{k4\pi}{3} \\ & t= \dfrac{5\pi}{9} + \dfrac{k4\pi}{3}} (k\in\mathbb{Z}).
		\end{eqnarray*}
		TH1: $ t=\dfrac{-\pi}{3}+\dfrac{k4\pi}{3}$, $0\le t \le 2 $
		$ \Rightarrow 0\le \dfrac{-\pi}{3}+\dfrac{k4\pi}{3} \le 2 \Rightarrow \dfrac{1}{4}\le k \le \dfrac{6+\pi}{4\pi} \approx 0{,}73$.\\
		Mà $ k\in \mathbb{Z} $ suy ra $ k\in \varnothing$.\\
		TH2: $ t= \dfrac{5\pi}{9} + \dfrac{k4\pi}{3}, 0\le t\le 2 $
		$ \Rightarrow 0\le \dfrac{5\pi}{9} + \dfrac{k4\pi}{3} \le 2 \Rightarrow \dfrac{-5}{3} \le t \le \dfrac{18-5\pi}{12\pi}\approx 0{,}06 $.\\
		Mà $ k\in \mathbb{Z} $ suy ra $ k=0 \Rightarrow t=\dfrac{5\pi}{9}\approx 1{,}7 $ (giây).\\
	}
\end{ex}

% #5022 | Archive/2025-Sangtac/Khoi12/Data/1-CD/CD-1-OTGK1-De02.tex | mã gốc 1D1H1-4
\begin{ex}%[1D1H1-4]
	Đường kính của một bánh xe đạp là $65$ (cm). Để người đi xe đạp đi được quãng đường $2041$ (m) thì mỗi bánh xe phải lăn bao nhiêu vòng? (lấy $\pi =3{,}14$).
	\shortans{$1000$}
	\loigiai{
		Đổi $2041$ (m) $=204100$ (cm).\\
		Chu vi bánh xe là $\ell =\alpha \cdot R=2\pi R=65\cdot 3{,}14=204{,}1$ (cm).\\
		Số vòng bánh xe cần lăn để đi được quãng đường dài $2041$ (m) là $204100 : 204{,}1=1000$ (vòng). }
\end{ex}

% #5023 | Archive/2025-Sangtac/Khoi12/Data/1-CD/CD-1-OTGK1-De02.tex | mã gốc 1D1H4-8
\begin{ex}%[1D1H4-8]
	\immini[thm]{Đồ thị ở hình bên mô tả sự thay đổi chiều cao của mực nước tại cảng trong vòng $24$ giờ tính từ lúc nửa đêm. Biết chiều cao của mực nước $h$ (m) theo thời gian $t$ (h) ($0\leq t\leq 24$) được cho bởi công thức $h=19+5\cos\left(\dfrac{\pi}{12}t+ \dfrac{\pi}{4}\right)$. Tìm thời điểm đầu tiên trong ngày khi chiều cao của mực nước là $14$ m.}
	{\begin{tikzpicture}[>=stealth,line join=round,line cap=round,font=\footnotesize,scale=0.55]
			\begin{scope}[scale=0.3]
				\draw [->](-1,0)--(0,0)node[above right]{$O$}--(27,0)node[below]{$t(h)$};
				\draw [->](0,-1)--(0,20)node[left]{$h(m)$};
				\draw [dashed](12,0)node[below]{$12$}|-(0,10)node[left]{$12$} (24,0)node[below]{$24$}|-(0,16)node[left]{$17$};
				\draw[smooth,samples=400,domain=0:24] plot(\x,{13+3*cos(pi/12*\x r)});
			\end{scope}
	\end{tikzpicture}}
	\shortans[0]{$9$}
	\loigiai{ 
		Do chiều cao của mực nước là $14$ m nên
		\allowdisplaybreaks 
		\begin{eqnarray*} 
			19+5\cos\left(\dfrac{\pi}{12}t+\dfrac{\pi}{4}\right)=14\\
			&\Leftrightarrow& \cos\left(\dfrac{\pi}{12}t+\dfrac{\pi}{4}\right)=-1\\ 
			&\Leftrightarrow& t=9+24k.
		\end{eqnarray*}
	Do $0\leq t\leq 24$ suy ra $t=9$.\\
		Vậy thời điểm $9$ giờ mà mực nước tại cảng là $14$ m. 
	}
\end{ex}

% #5025 | Archive/2025-Sangtac/Khoi12/Data/1-CD/CD-1-OTGK1-De02.tex | mã gốc 1D1C5-6
\begin{ex}%[1D1C5-6]
	\immini[thm]{Giả sử vật giao động điều hòa xung quanh vị trí cân bằng theo phương trình $x=2\cos \left(5t-\dfrac{\pi}{6} \right)$. Ở đây, thời gian $t$ tính bằng giây và quãng đường $x$ tính bằng cm. Hãy cho biết trong khoảng thời gian từ 0đến 6giây, vật đi qua vị trí cân bằng bao nhiêu lần? (hình ảnh minh họa vị trí cân bằng)}
	{\definecolor{darkpastelgreen}{rgb}{0.01, 0.75, 0.24}
		\definecolor{almond}{rgb}{0.94, 0.87, 0.8}
		\definecolor{arsenic}{rgb}{0.23, 0.27, 0.29}
		\definecolor{eggplant}{rgb}{0.38, 0.25, 0.32}
		\definecolor{crimson}{rgb}{0.86, 0.08, 0.24}
		\begin{tikzpicture}[line join=round, line cap=round,scale=.7,transform shape]
			\clip (-2.5,-2.5) rectangle (4.5,4.2);
			\tikzset{q_h/.pic={
					% Tre
					\draw[fill=darkpastelgreen] (-.05,4)--(.9,-1.8)
					..controls +(30:.05) and +(-30:.05) ..(1,-1.75)--(.05,4);
					\draw[fill=darkpastelgreen] (-.1,4)--(.6,-1.4)
					..controls +(30:.05) and +(-30:.05) ..(.5,-1.4)--(-.2,4)
					;
					%----------------------
					\def\M{ 
						(-.2,1.76)
						..controls +(-80:.1) and +(80:.1) ..(-.18,1.6)
						..controls +(-40:.1) and +(30:.06) ..(-.18,1.5)--(-.17,1.46)--(-.22,1.44)--(-.19,1.43)--(-.2,1.43)
						..controls +(-80:.1) and +(-12:.05) ..(-.3,1.35)--(-.35,1.28)--(-.55,1.35)
						..controls +(60:.1) and +(-90:.05) ..(-.45,1.5)
						..controls +(100:.1) and +(140:.05) ..(-.38,1.55)
						..controls +(70:.1) and +(-140:.05) ..cycle
						;
					}
					
					\draw[black]\M;
					\fill[almond] \M;
					\draw (-.2,1.76)
					..controls +(-140:.1) and +(50:.05) ..(-.3,1.6)
					(-.2,1.76)
					..controls +(-120:.1) and +(40:.05) ..(-.25,1.6)
					(-.28,1.578)
					..controls +(-20:.02) and +(-150:.02) ..(-.22,1.58)
					(-.25,1.568)--(-.23,1.567);
					%------------------Áo
					\def\A3{ 
						(-.1,.9)
						..controls +(-30:.05) and +(110:0) ..(.18,1)
						--(.19,.95)
						..controls +(-140:.05) and +(30:.05) ..(-.1,.7)--cycle
						;
					}
					\fill[arsenic] \A3;
					\draw[black]\A3;
					\def\A2{ 
						(-.35,1.28)
						..controls +(-30:.05) and +(110:.05) ..(-.25,1.18)
						..controls +(-70:.05) and +(90:.02) ..(-.05,.9)
						..controls +(-90:.02) and +(50:.05) ..(-.05,.4)
						..controls +(-130:.02) and +(50:.05) ..(.6,-.8)
						..controls +(-130:.02) and +(50:.05) ..(.9,-1.8)
						..controls +(-130:.2) and +(30:.3) ..(-.5,-2.3)
						..controls +(130:.2) and +(-90:.1) ..(-.2,-1.5)
						..controls +(130:.2) and +(-140:.1) ..cycle
						;
					}
					\fill[arsenic] \A2;
					\draw[black]\A2;
					\def\A{ 
						(-.55,1.35)
						..controls +(-150:.2) and +(120:.1) ..(-.45,0)
						..controls +(-60:.1) and +(50:.1) ..(-.3,-1.6)
						..controls +(-10:.1) and +(-120:.1) ..(-.15,-1.5)
						..controls +(60:.1) and +(-120:.1) ..(0,-1.45)
						..controls +(120:.3) and +(-50:.1) ..(-.35,.77)
						..controls +(-50:.1) and +(130:.1) ..(-.15,.48)
						..controls +(-50:.1) and +(-130:.1) ..(.43,.93)
						..controls +(-130:.05) and +(-40:.05) ..(.3,1.05)
						..controls +(-130:.05) and +(40:.05) ..(-.05,.75)
						..controls +(-140:.05) and +(-40:.05) ..(-.3,1.15)
						..controls +(100:.05) and +(-40:.05) ..(-.35,1.28)--cycle
						;
					}
					\fill[eggplant] \A;
					\draw[black]\A;
					%-------------------------------
					%nón
					\def\N{ 
						(-.14,1.95)
						..controls +(-80:.1) and +(80:.1) ..(-.17,1.7)
						..controls +(110:.02) and +(-40:.02) ..(-.2,1.76)
						..controls +(-110:.1) and +(50:.4) ..(-.48,1.45)--(-.45,1.35)
						..controls +(-150:.05) and +(20:.05) ..(-.7,1.32)
						..controls +(50:.05) and +(-120:.05) ..(-.65,1.45)
						..controls +(160:.02) and +(-50:.02) ..(-.7,1.48)
						..controls +(160:.02) and +(-50:.02) ..(-.8,1.6)--(-.75,1.63)
						..controls +(-20:.06) and +(130:.04) ..(-.6,1.47)
						..controls +(130:.2) and +(-90:.1) ..(-.7,1.78)
						..controls +(90:.1) and +(170:.1) ..(-.4,1.9)
						..controls +(20:.1) and +(160:.1) ..cycle
						;
					}
					\fill[arsenic] \N;
					\draw[black]\N;
					\draw (-.14,1.95)
					..controls +(-100:.05) and +(40:.06) ..(-.2,1.76)
					(-.5,1.55)
					..controls +(30:.1) and +(-100:.06) ..(-.3,1.76)
					(-.6,1.7)
					..controls +(-55:.1) and +(120:.06) ..(-.55,1.5)
					;
					
					\def\T1{ 
						(.3,1.03)
						..controls +(60:.05) and +(-120:.05) ..(.35,1.15)
						..controls +(60:.05) and +(-140:.05) ..(.45,1.26)
						..controls +(30:.05) and +(60:.04) ..(.5,1.22)
						..controls +(30:.03) and +(50:.03) ..(.54,1.18)
						..controls +(30:.03) and +(50:.03) ..(.56,1.12)
						..controls +(30:.03) and +(50:.03) ..(.58,1.07)
						..controls +(160:.03) and +(40:.03) ..(.49,1)
						..controls +(-140:.03) and +(40:.03) ..(.41,.92)--cycle
						;
					}
					\fill[almond] \T1;
					\draw[black]\T1;
					
					\fill[crimson](-.27,.4)
					..controls +(-80:.6) and +(-80:.1) ..(-.5,-1.9)--(-.45,-1.85)
					..controls +(90:.1) and +(-80:1) ..(-.28,.48)--cycle;
					
					\fill[crimson](-.27,.4)
					..controls +(-80:.6) and +(-80:.1) ..(-.2,-1.9)--(-.15,-1.85)
					..controls +(90:.1) and +(-80:1) ..(-.28,.48)--cycle;
					
			}}
			\tikzset{no/.pic={
					\draw[crimson] (1.8,-1)
					..controls +(150:.4) and +(-120:.5) ..(1.7,-1)
					..controls +(30:.4) and +(-40:.3) ..(1.75,-1)
					..controls +(-30:.2) and +(40:.2) ..(1.7,-1.4)
					(1.75,-1)
					..controls +(-50:.1) and +(20:.1) ..(1.7,-1.2)
					
					;
			}}
			\draw[thick] (4,-1.25) arc(-60:-120:6cm);
			\draw[dashed,thick] (4,-1.25)--(1.3,4)--(-2,-1.25) (1.3,4)--(1.3,-2);
			\node at (1.5,.5) [rotate=90]{\tiny Vị trí cân bằng};
			\draw[thick,<->] (.35,-2)--(1.3,-2);
			\node at (.85,-2) [above]{\tiny $x$};
			\fill (4,-1.25) circle (0.05);
			\fill (-2,-1.25) circle (0.05);
			\fill (1.3,4) circle (0.05);
			\path
			(0,0)pic[scale=1,rotate=-20]{q_h}(-.88,1.48)pic[scale=1, scale=.7,rotate=-20]{no};
			
			
	\end{tikzpicture}}
	\shortans{$9$}
	\loigiai{
		Khi vật ở vị trí cân bằng thì $x=0$
		\begin{eqnarray*}
			&\Leftrightarrow& 3\cos \left(5t-\dfrac{\pi}{6} \right)=0\Leftrightarrow \cos\left(5t-\dfrac{\pi}{6} \right)=0\\
			&\Leftrightarrow& 5t-\dfrac{\pi}{6}=\dfrac{\pi}{2}+k\pi, k\in {\bf \mathbb{Z}}\Leftrightarrow t=\dfrac{2\pi}{15}+\dfrac{k\pi}{5}, k\in {\bf \mathbb{Z}}.
		\end{eqnarray*}
		Mà theo giả thiết 
		\begin{eqnarray*}
			0< t < 6&\Leftrightarrow& 0< \dfrac{2\pi}{15}+\dfrac{k\pi}{5} < 6\\
			&\Leftrightarrow&-\dfrac{2}{3} < k < \dfrac{30}{\pi}-\dfrac{2}{3}\\
			k\in {\bf \mathbb{Z}}&\Rightarrow& k\in \left\{0;1;2;3;4;5;6;7;8\right\}.
		\end{eqnarray*}
	}
\end{ex}

% #5031 | Archive/2025-Sangtac/Khoi12/Data/1-CD/CD-1-OTGK1-De01.tex | mã gốc 1D1V1-6
\begin{ex}%[1D1V1-6]
	Bánh xe của một xe máy có đường kính $100$ cm quay được $6390$ vòng trong $3$ phút. Hỏi vận tốc $v$ (km/h) của người đi xe máy đó là bao nhiêu? (làm tròn đến hàng đơn vị).
	\shortans{$40$}
	\loigiai{
		Bán kính của bánh xe là $100: 2=50$ (cm).\\
		Trong $1$ phút bánh xe quay được $6390\colon 3=2130$ (vòng).\\
		Trong $1$ giờ bánh xe quay được $2130\cdot 60=127800$ (vòng).\\
		Trong $1$ giờ bánh xe đi được quãng đường là $127800\cdot 2\pi \cdot 50\approx 40149554$ (cm).\\
		Vậy vận tốc của người đi xe máy là $40149554: 1=40149554$ (cm/h) $\approx 40$ (km/h). }
\end{ex}

% #5032 | Archive/2025-Sangtac/Khoi12/Data/1-CD/CD-1-OTGK1-De01.tex | mã gốc 1D1H4-8
\begin{ex}%[1D1H4-8]
	Hằng ngày mực nước tại một cảng biển lên xuống theo thuỷ triều. Độ sâu $h$ (mét) của mực nước theo thời gian $t$ (giờ) trong một ngày được cho bởi công thức $h=16 + 9\sin\left( \dfrac{t}{10}\pi \right)$ với $0\leq t\leq 24$. Thời điểm đầu tiên $t_0$ (giờ) mà mực nước tại cảng cao nhất là $h_0$ (mét). Khi đó $t_0 + h_0$ bằng
	\shortans[0]{$30$}
	\loigiai{ 
		Ta có $-1\leq \sin\left( \dfrac{t}{10}\pi \right)\leq 1 \Leftrightarrow 7\leq h \leq 25.$ \\
		Vậy mực nước tại cảng cao nhất bằng $25$ m khi $$\sin\left( \dfrac{t}{10}\pi \right)= 1\Leftrightarrow \left( \dfrac{t}{10}\pi \right) = \dfrac{\pi}{2}+k2\pi\Leftrightarrow t= 5+20k.$$ 
		Mà $0\leq t\leq 24$ nên thời điểm đầu tiên mà mực nước tại cảng cao nhất là 5 (giờ).\\ 
		Khi đó $t_0 + h_0=30$.  
}
\end{ex}

% #5042 | Archive/2025-Sangtac/Khoi12/Data/1-KNTT/KNTT_1-OTGK1_De08.tex | mã gốc 1D1V2-5
\begin{ex}%[1D1V2-5]
	\immini{Một chiếc đu quay quay quanh tâm $O$. Biết tâm $O$ cách mặt đất $12$ m. Một người ngồi trong cabin ở vị trí thấp nhất trên đu quay cách mặt đất $3$m. Đu quay đưa người đó đến vị trí $B$ (như hình vẽ), biết góc lượng giác tạo bởi tia đầu $OA$ và tia cuối $OB$ bằng $60^\circ$.
		Hỏi lúc này, người đó ở vị trí cách mặt đất bao nhiêu mét? Kết quả làm tròn đến một chữ số thập phân. }{\begin{tikzpicture}[scale=0.7,>=stealth]
			\def\r{3}
			\def\a{\r/10}
			\def\n{12}%Số nan đu quay
			\def\m{12}%Số lồng đu quay
			\pgfmathsetmacro{\g}{20}
			\pgfmathsetmacro{\b}{2*\a/3}
			\pgfmathsetmacro{\xo}{\a *cos(\g)}
			\pgfmathsetmacro{\yo}{\b *sin(\g)}
			\def\longdq(#1:#2){
				\begin{scope}[shift={(#1:#2)}]
					\draw (0,-\b) ellipse ({\b/2} and {\b});
					\fill[blue!60] (\xo,\yo-\b) arc (\g:-180-\g:{\a} and {\b})--cycle;
					\draw (0,-\b) ellipse({\a} and {\b});
				\end{scope}
			}
			\foreach \i in {1,...,\n}{
				\pgfmathsetmacro{\gm}{360*\i/\n}
				\pgfmathsetmacro{\gh}{(360*\i+180)/\n}
				\pgfmathsetmacro{\gb}{360*(\i+1)/\n}
				\draw (\gm:\r-\a)--(\gh:\r)--(\gb:\r-\a);
				\draw (\gm:\r-\a)--(\gm:\a);
			}
			\draw[double] (0:0) circle (\r);
			\draw (0:0) circle (\r-\a);
			\foreach \i in {1,...,\m}{\longdq(360*\i/\m:\r)}
			\draw[top color =gray!10, bottom color=gray!30] (-30:\a) arc(-30:0:\a) --(-70:1.5*\r)--++(\a/2,0)--([turn]-70:\r/2) --([turn]-110:3*\a)--(-30:\a);
			\draw[top color =gray!10, bottom color=gray!30, xscale=-1] (-30:\a) arc(-30:0:\a) --(-70:1.5*\r)--++(\a/2,0)--([turn]-70:\r/2) --([turn]-110:3*\a)--(-30:\a);
			\fill[ball color=orange] (0:0) circle (\a);
			\draw [red, line width=2](0,-4.5)--(0,3.05);
			\draw (0,0) node[above left]{$O$};
			\coordinate (O) at (0,0);
			\coordinate[scale=1.5,label=right:$A$] (A) at ($(O) + (0:3)$);\fill (A)circle(3pt);
			
			\coordinate [scale=1.5,label=above right:$B$](M) at ($(O)!1!60:(A)$);\fill (M)circle(3pt);
			\coordinate (v) at ($(M) + (120:3)$);
			\coordinate (Vx) at ($(M) + (-180:1.5)$);
			\draw[double] (0:0) circle (\a);
			\draw[red, line width=2] (0,0)--(A) (0,0)--(M);
			\draw [shift={(0,0)},->,line width=0.3pt] (0:.9) arc (0:60:.9);		
			\draw (.75,.5) node[right]{\scriptsize $60^\circ$};	
	\end{tikzpicture}}
	\shortans{19{,}8}
	\loigiai{
		Theo đề bài, ta có $OI=12$ m; $OB=OH=9$ m; $OK=BB'$.\\
		Xét tam giác $OBB'$, $\sin 60^\circ =\dfrac{BB'}{OB}\Rightarrow BB'=OB\sin 60^\circ =9\cdot \dfrac{\sqrt{3}}{2}$.\\
		Suy ra, ở vị trí $B$, người này cách mặt đất một khoảng bằng
		$$KI=OK+OI=12+9\cdot \dfrac{\sqrt{3}}{2}\approx 19{,}8 \text{ m}.$$}
\end{ex}

% #5043 | Archive/2025-Sangtac/Khoi12/Data/1-KNTT/KNTT_1-OTGK1_De08.tex | mã gốc 1D1V4-6
\begin{ex}%[1D1V4-6]
	Trong môn cầu lông, khi phát cầu, người chơi cần đánh cầu qua khỏi lưới sang phía sân đối phương và không được để cho cầu rơi ngoài biên.
	Trong mặt phẳng toạ độ $O x y$, chọn điểm có tọa độ $(0 ; y_0)$ là điểm xuất phát thì phương trình quỹ đạo của cầu lông khi rời khỏi mặt vợt là $$y=\dfrac{-g \cdot x^2}{2 \cdot v_0^2 \cdot \cos ^2 \alpha}+\tan (\alpha) \cdot x+y_0;$$ trong đó:
	\begin{itemize}
		\item $g$ là gia tốc trọng truờng (thường được chọn là $9{,}8 \mathrm{~m} / \mathrm{s}^{2}$);
		\item $\alpha$ là góc phát cầu (so với phương ngang của mặt đất);
		\item $v_0$ là vận tốc ban đầu của cầu;
		\item $y_0$ là khoảng cách từ vị trí phát cầu đến mặt đất.
	\end{itemize}
	Đây là một hàm số bậc hai nên quỹ đạo chuyển động của cầu lông là một parabol.
	Một người chơi cầu lông đang đứng khoảng cách từ vị trí người này đến vị trí cầu rơi chạm đất (tầm bay xa) là $6{,}68 \mathrm{~m}$. Quan sát hình bên dưới, hỏi người chơi đã phát cầu góc lớn nhất khoảng bao nhiêu độ so với mặt đất? (biết cầu rời mặt vợt ở độ cao $0{,}7 \mathrm{~m}$ so với mặt đất và vận tốc xuất phát của cầu là $8 \mathrm{~m} / \mathrm{s}$, bỏ qua sức cản của gió và xem quỹ đạo của cầu luôn nằm trong mặt phẳng phẳng đứng).
	\begin{center}
		\begin{tikzpicture}[scale=1, line join=round, line cap=round, >=stealth]
			\tikzset{every node/.style={scale=1}}
			\def\xmin{-1}\def\xmax{12}\def\ymin{-1}\def\ymax{5}
			\draw[->] (\xmin-0.2,0)--(\xmax+0.2,0) node[below]{$x$};
			\draw[->] (0,\ymin-0.2)--(0,\ymax+0.2) node[right]{$y$};
			\draw (0,0) node[below left]{$O$};
			\draw (3,3.84) node[above]{Quả cầu};\fill (3,3.84)circle(5pt);
			\foreach \x in {10}\draw (\x,0.1)--(\x,-0.1) node[below]{Điểm chạm đất};\fill (10,0)circle(3pt);
			\foreach \y in {2}\draw (0.1,\y)--(-0.1,\y) node[left]{$0{,}7~\mathrm{m}$};
			\clip (\xmin,\ymin) rectangle (\xmax,\ymax);
			\draw[thick,smooth,samples=200,domain=0:10] plot(\x,{0.010316090813929726*(\x)^(3.0)-0.25038214657326735*(\x)^(2.0)+1.2722123843397009*(\x)+2.0});
			\draw [->](0,2)--(1.5,3);
			\draw [dashed](0,2)--(2,2);
			\draw (.5,2.25) node[right]{$\alpha$};
			\draw (1.5,1) node[right]{$y=\dfrac{-g \cdot x^2}{2 \cdot v_0^2 \cdot \cos ^2 \alpha}+\tan (\alpha) \cdot x+y_0$};
		\end{tikzpicture}
	\end{center}
	\shortans{54}
	\loigiai{
		Với $g=9{,}8 \mathrm{~m} / \mathrm{s}^2$, vận tốc ban đầu $v_0=8 \mathrm{~m} / \mathrm{s}$, phương trình quỹ đạo của cầu:
		$$y=\dfrac{-g \cdot x^{2}}{2 \cdot v_{0}^{2} \cdot \cos ^{2} \alpha}+\tan (\alpha) \cdot x+y_{0}.$$
		Khoảng cách từ vị trí người này đến vị trí cầu rơi chạm đất (tầm bay xa) là $6{,}68 \mathrm{~m}$; nghĩa là $x=6{,}68 \mathrm{~m}$. Ta có
		\allowdisplaybreaks
		\begin{eqnarray*}
			&&\dfrac{-9{,}8 \cdot(6{,}68)^2}{128 \cdot \cos ^2 \alpha}+\tan (\alpha) \cdot(6{,}68)+0,7=0 \\
			&\Leftrightarrow& \dfrac{-9{,}8 \cdot(6{,}68)^2}{128}\left(1+\tan ^2 \alpha\right)+\tan (\alpha) \cdot(6{,}68)+0,7=0 \\
			&\Leftrightarrow& \hoac{&\tan \alpha \approx 1{,}3 7 8  \\&
				\tan \alpha \approx 0{,}5 7 6} \Leftrightarrow \hoac{&
				\alpha \approx 54{,}04^\circ \\&
				\alpha \approx 29{,}97^\circ.}
		\end{eqnarray*}
		Vậy người chơi đã phát cầu một góc gần $54^\circ$ hoặc gần $30^\circ$ so với mặt đất.}
\end{ex}

% #5044 | Archive/2025-Sangtac/Khoi12/Data/1-KNTT/KNTT_1-OTGK1_De08.tex | mã gốc 1D1V5-6
\begin{ex}%[1D1V5-6]
	Số giờ có ánh nắng mặt trời của một thành phố A trong ngày thứ $t$ (ở đây $t$ là số ngày tính từ ngày $1$ tháng Giêng) của một năm không nhuận được mô hình hóa bởi hàm số
	$$L(t)=12+2{,}83\sin\left[\dfrac{2\pi}{365}(t-80)\right]\text{ với } t\in \mathbb{Z} \text{ và } 0<t\le 365.$$
	Vào hai  ngày $a$ và $b$ trong năm, thành phố A có khoảng $10$ giờ ánh sáng mặt trời, tính $a+b$.
	
	\shortans{342}
	\loigiai 
	{Ngày thành phố A có $10$ giờ ánh sáng mặt trời nhất ứng với 
		\allowdisplaybreaks
		\begin{eqnarray*}
			12+2{,}83\sin\left[\dfrac{2\pi}{365}(t-80)\right]=10&\Leftrightarrow&\sin\left[\dfrac{2\pi}{365}(t-80)\right]=-\dfrac{200}{283}\\
			&\Leftrightarrow&\hoac{&\dfrac{2\pi}{365}(t-80)\approx -0{,}78+k2\pi\\&\dfrac{2\pi}{365}(t-80)\approx 3{,}93+k2\pi}\\
			&\Leftrightarrow&\hoac{&t\approx 34{,}69+365k\\&t\approx 308{,}30+365k}\,(k\in\mathbb{Z}).
		\end{eqnarray*}
		Vì $0<t\le 365$ nên $k=0$, suy ra $t\approx 34{,}69$ hoặc $t\approx 308{,}30$.\\
		Như vậy, vào ngày thứ $34$ của năm, tức là khoảng ngày $3$ tháng $2$ và ngày thứ $308$ của năm, tức là ngày $4$ tháng $11$ thì thành phố A sẽ có $10$ giờ ánh sáng mặt trời.	
	}
\end{ex}

% #5048 | Archive/2025-Sangtac/Khoi12/Data/1-KNTT/KNTT_1-OTGK1_De01.tex | mã gốc 1D1V1-6
\begin{ex}%[1D1V1-6]
	\immini{Thanh $OM$ quay ngược chiều kim đồng hồ quanh trục $O$ của nó trên một mặt phẳng thẳng đứng và in bóng vuông góc xuống mặt đất như Hình vẽ. Vị trí ban đầu của thanh là $OA$. Hỏi độ dài bóng $O'M'$ của $OM$ khi thanh quay được $2\dfrac{1}{3}$ vòng là bao nhiêu cm, biết độ dài thanh $OM$ là $20 ~\mathrm{cm}$?  }{\begin{tikzpicture}[scale=0.9, font=\footnotesize,line join=round, line cap=round, >=stealth,x=2cm,y=2cm]
			\def\x{1}
			\def\a{30}
			\pgfmathsetmacro{\y}{\x*sin(\a)}
			\pgfmathsetmacro{\z}{\x*cos(\a)}
			\fill (0,0) node [below right]{$O$} circle(1pt);
			\draw[dashed](1,0) arc (0:360:2cm);
			\fill (1,0) node[above right]{$A$}circle(1pt);
			\fill (210:1) node[left]{$M$} circle(1pt);
			\coordinate (P) at (-1.2,-1.2);
			\coordinate (Q) at (1.2,-1.2);
			\draw (1,0)--(0,0)--(210:1);
			\draw[->] (.2,0) arc (0:210:.4cm);
			\node at (-.2,.2){$\alpha$};
			\coordinate (M') at ($(P)!(210:1)!(Q)$);
			\coordinate (O') at ($(P)!(0,0)!(Q)$);
			\foreach \i in {M',O'}{\fill[](\i) circle (1pt);}
			\draw (P)--(Q) (210:1)--(M') (0,0)--(O');
			\foreach \i in {M',O'}{\draw(\i) node[scale=0.9,below]{$\i$};}
			\draw[->] (0,1.5)--(0,1.1);\draw[->](-0.2,1.5)--(-0.2,1.1);
			\draw[->](0.2,1.5)--(0.2,1.1);
			\node at (-.5,-1.3){bóng};
			\node at (.6,1.3){ánh sáng};	
		\end{tikzpicture}
	}
	\shortans{10}
	\loigiai{
		Đặt hệ trục tọa độ như hình vẽ
		\begin{center}
			\begin{tikzpicture}[scale=0.9, font=\footnotesize,line join=round, line cap=round, >=stealth,x=2cm,y=2cm]
				\def\x{1}
				\def\a{30}
				\pgfmathsetmacro{\y}{\x*sin(\a)}
				\pgfmathsetmacro{\z}{\x*cos(\a)}
				\fill (0,0) node [below right]{$O$} circle(1pt);
				\draw[->] (-1.1,0)--(1.1,0)node[below]{$x$};
				\draw[->] (0,-1.1)--(0,1.1)node[left]{$y$};
				\draw[dashed](1,0) arc (0:360:2cm);
				\fill (1,0) node[above right]{$A$}circle(1pt);
				\fill (210:1) node[left]{$M$} circle(1pt);
				\coordinate (P) at (-1.2,-1.2);
				\coordinate (Q) at (1.2,-1.2);
				\draw (1,0)--(0,0)--(210:1);
				\draw[->] (.2,0) arc (0:210:.4cm);
				\draw[dashed] (210:1)--(-0.867,0)node[above right]{$H$}circle(1pt);
				\node at (-.2,.2){$\alpha$};
				\coordinate (M') at ($(P)!(210:1)!(Q)$);
				\coordinate (O') at ($(P)!(0,0)!(Q)$);
				\foreach \i in {M',O'}{\fill[](\i) circle (1pt);}
				\draw (P)--(Q) (210:1)--(M') (0,0)--(O');
				\foreach \i in {M',O'}{\draw(\i) node[scale=0.9,below]{$\i$};}
				\draw[->] (0,1.5)--(0,1.1);\draw[->](-0.2,1.5)--(-0.2,1.1);
				\draw[->](0.2,1.5)--(0.2,1.1);
				\node at (-.5,-1.3){bóng};
				\node at (.6,1.3){ánh sáng};	
			\end{tikzpicture}
		\end{center}
		Điểm $M$ là điểm biểu diễn góc lượng giác $\alpha $. Kẻ $MH$ vuông góc với $Ox$.\\
		Ta có $\alpha =2\dfrac{1}{3}\cdot 360^\circ =840^\circ $. Khi đó $M\left(20\cos 840^\circ;20\sin 840^\circ\right)$\\
		Suy ra $OH=\left| 20\cos 840^\circ\right|=10\left(\mathrm{cm}\right)$.\\ Vậy độ dài bóng $O'M'$ của $OM$ khi thanh quay được $2\dfrac{1}{3}$ vòng là $10 ~\mathrm{cm}$.
	}
\end{ex}

% #5049 | Archive/2025-Sangtac/Khoi12/Data/1-KNTT/KNTT_1-OTGK1_De01.tex | mã gốc 1D1V2-5
\begin{ex}%[1D1V2-5]
	Một dòng điện xoay chiều có cường độ dòng điện $i$ (ampe) tại thời điểm $t$ (giây) được tính bởi công thức $i=2\cos\left(\dfrac{131\pi}{12}t\right)$. Tính giá trị của cường độ dòng điện tại thời điểm  $t=1$ (giây). \textit{(Làm tròn đến hàng phần chục)} 
	
	\shortans{-1{,}9}
	\loigiai{ Ta có $i=2\cos\left(\dfrac{131\pi}{12}t\right)$ nên tại thời điểm  $t=1$ (giây) thì
		\allowdisplaybreaks
		\begin{eqnarray*}
			i&=&2\cos\left(\dfrac{131\pi}{12}t\right)=2\cos\left(11\pi-\dfrac{1}{12}\pi\right)=2.\left[\cos11\pi.\cos\dfrac{1}{12}\pi+\sin11\pi.\sin\dfrac{1}{12}\pi\right]\\
			&=&2\cdot \left[-1\cdot\dfrac{\sqrt{6}+\sqrt{2}}{4}+0\cdot\dfrac{\sqrt{6}+\sqrt{2}}{4}\right]=-\left(\dfrac{\sqrt{6}+\sqrt{2}}{2}\right)\approx-1{,}93 \text{ (ampe)}.
		\end{eqnarray*}
	}	
\end{ex}

% #5050 | Archive/2025-Sangtac/Khoi12/Data/1-KNTT/KNTT_1-OTGK1_De01.tex | mã gốc 1D1V4-8
\begin{ex}%[1D1V4-8]
	\immini{
		Một chiếc guồng nước có dạng hình tròn bán kính $2{,}5$ m; trục của nó đặt cách mặt nước $2$ m (hình bên). Khi guồng quay đều, khoảng cách $h$ (m) tính từ một chiếc gầu tại điểm $A$ trên guồng đến mặt nước là $h=|y|$ trong đó $$y=2+2{,}5\sin 2\pi\left(x-\dfrac{1}{4}\right)$$ với $x$ là thời gian quay của guồng ($x\ge 0$), tính bằng phút; ta quy ước rằng $y>0$ khi gầu ở trên mặt nước và $y<0$ khi gầu ở dưới mặt nước. Chiếc gầu cách mặt nước $2$ mét lần đầu tiên tại thời điểm mấy giây? 
	}
	{
		\begin{tikzpicture}[>=stealth,line join=round,line cap=round,font=\footnotesize,scale=1]
			\def\r{2.5}
			\coordinate[label=above:$O$] (O) at (0,0);
			\draw[name path=(C)] (O) circle (\r);
			\coordinate (M) at (0,-2);
			\path 
			(O)++(45:\r)coordinate(A)node[above right]{$A$}
			(-3,-2)coordinate(B)
			(3,-2)coordinate(C)
			($(B)!(A)!(C)$)coordinate(H)
			;
			\draw[dashed] (A)--(H)node[midway,right]{$h$ m}
			(B)--(C)
			;
			\draw[<->](O)--(A)node[midway,left]{$2{,}5$ m};
			\draw[<->](O)--(M)node[midway,right]{$2$ m};
			\draw[->] (A) arc (45:60:2.5);
		\end{tikzpicture}
	}
	\shortans{15}
	\loigiai 
	{Gầu cách mặt nước $2$ m nên 
		\allowdisplaybreaks
		\begin{eqnarray*}
			2+2{,}5\sin 2\pi\left(x-\dfrac{1}{2}\right)=2&\Leftrightarrow&\sin 2\pi\left(x-\dfrac{1}{2}\right)=0\\
			&\Leftrightarrow&2\pi\left(x-\dfrac{1}{4}\right)=k\pi\\
			&\Leftrightarrow&x=\dfrac{1}{4}+\dfrac{k}{2},\,(k\in\mathbb{Z}).
		\end{eqnarray*}
		Vậy chiếc gầu cách mặt nước $2$ m lần đầu tiên tại thời điểm $x=\dfrac{1}{4}$ phút $=15$ giây.
	}
\end{ex}

% #5064 | Archive/2025-Sangtac/Khoi12/Data/1-KNTT/KNTT_1-OTGK1_De05.tex | mã gốc 1D1V3-7
\begin{ex}%[1D1V3-7]
	Trong Vật lí, phương trình tổng quát của một vật dao động điều hoà cho bởi công thức $x(t)=A\cos(\omega t+\varphi)$, $t$ trong đó   là thời điểm (tính bằng giây), $x(t)$  là li độ của vật tại thời điểm $t$, $A$ là biên độ dao động $(A>0)$ và $\varphi \in \left[-\pi; \pi \right]$ là pha ban đầu của dao động. Xét hai dao động điều hoà có phương trình
	\allowdisplaybreaks
	\begin{eqnarray*}
		x_{1}(t)&=&2\cos \left(\dfrac{\pi}{2}t+\dfrac{\pi}{5} \right)\text{ (cm)},\\
		x_{2}(t)&=&2\cos \left(\dfrac{\pi}{2}t -\dfrac{3\pi}{10}\right)\text{ (cm)}.
	\end{eqnarray*}
	Xét dao động tổng hợp $x(t)=x_{1}(t) + x_{2}(t)$, giả sử pha ban đầu của dao động tổng hợp này là $\dfrac{\pi}{a}$. Tính $a$.

	\shortans{-20}
	\loigiai{	Dao động tổng hợp là \allowdisplaybreaks
		\begin{eqnarray*}
			x(t)&=&x_{1}(t)+x_{2}(t)= 2\cos \left(\dfrac{\pi}{2}t+\dfrac{\pi}{5} \right)+2\cos \left(\dfrac{\pi}{2}t-\dfrac{3\pi}{10}\right)\\
			\Rightarrow x(t)&=&2\left[\cos\left(\dfrac{\pi}{2}t+\dfrac{\pi}{5}\right)+\cos\left(\dfrac{\pi}{2}t-\dfrac{3\pi}{10}\right)\right]\\
			\Rightarrow x(t)&=&4\cos\left[\dfrac{\left(\dfrac{\pi}{2}t+\dfrac{\pi}{5}\right)+\left(\dfrac{\pi}{2}t-\dfrac{3\pi}{10}\right)}{2}\right].\cos\left[\dfrac{\left(\dfrac{\pi}{2}t+\dfrac{\pi}{5}\right)-\left(\dfrac{\pi}{2}t-\dfrac{3\pi}{10}\right)}{2}\right].\\
			\Rightarrow x(t)&=&4\cdot \cos\left(\dfrac{\pi}{2}t-\dfrac{\pi}{20}\right)\cdot \cos\dfrac{\pi}{4}\Rightarrow x(t)=2\sqrt{2}.\cos\left(\dfrac{\pi}{2}t-\dfrac{\pi}{20}\right).
		\end{eqnarray*}
		Vậy biên độ của dao động tổng hợp là $A=2\sqrt{2}$ và pha ban đầu của dao động tổng hợp là $\varphi=-\dfrac{\pi}{20}$. Do đó $a=-20$.
	}	
\end{ex}

% #5065 | Archive/2025-Sangtac/Khoi12/Data/1-KNTT/KNTT_1-OTGK1_De05.tex | mã gốc 1D1V4-8
\begin{ex}%[1D1V4-8]
	Số giờ có ánh nắng mặt trời của một thành phố A trong ngày thứ $t$ (ở đây $t$ là số ngày tính từ ngày $1$ tháng Giêng) của một năm không nhuận được mô hình hóa bởi hàm số
	$$L(t)=12+2{,}83\sin\left[\dfrac{2\pi}{365}(t-80)\right]\text{ với } t\in \mathbb{Z} \text{ và } 0<t\le 365.$$
	Vào ngày thứ mấy  trong năm thì thành phố A có nhiều giờ ánh sáng mặt trời nhất?
	
	\shortans{171}
	\loigiai 
	{
		Ngày thành phố A có nhiều giờ ánh sáng mặt trời nhất ứng với 
		$$\sin\left[\dfrac{2\pi}{365}(t-80)\right]=1\Leftrightarrow\dfrac{2\pi}{365}(t-80)=\dfrac{\pi}{2}+k2\pi\Leftrightarrow t=171{,}25+365k,\,(k\in\mathbb{Z}).$$
		Vì $0<t\le 365$ nên $k=0$, suy ra $t=171{,}25$.\\
		Như vậy, vào ngày thứ $171$ của năm, tức là khoảng ngày $20$ tháng $6$ thì thành phố A sẽ có nhiều giờ ánh sáng mặt trời nhất.
	}
\end{ex}

% #5070 | Archive/2025-Sangtac/Khoi12/Data/1-CTST/CTST-1-OTGK1-De07.tex | mã gốc 1D1V4-8
\begin{ex}%[1D1V4-8]
	\immini[thm]{
		Hội Lim (tỉnh Bắc Ninh) vào mùa xuân thường có trò chơi đánh đu. Khi người chơi đu nhún đều, cây đu sẽ đưa người chơi đu dao động qua lại vị trí cân bằng (\textit{tham khảo hình vẽ bên}). Nghiên cứu trò chơi này, người ta thấy khoảng cách $h$ (m) từ người chơi đu đến vị trí cân bằng được biểu diễn qua thời gian $t$ (s) (với $t\ge 0$) bởi hệ thức $h=|d|$ với $d=3\cos\left[\dfrac{\pi}{3}(2t-1)\right]$, trong đó quy ước $d>0$ khi vị trí cân bằng ở  phía sau lưng người chơi đu và $d<0$ trong trường hợp ngược lại (\textit{Nguồn: R.Larson and B.Edwards, Calculus 10e Cengage}). Tìm thời điểm đầu tiên mà khoảng cách $h$ lớn nhất. (Viết kết quả dưới dạng số thập phân).
		\shortans{$0{,}5$}
	}
	{\begin{tikzpicture}[scale=1.2, font=\footnotesize,line join=round, line cap=round, >=stealth]
			\def\r{3}
			\path (0,0) coordinate (O)
			(-55:\r) coordinate (A)
			(-125:\r) coordinate (B)
			(-90:\r+.2) coordinate (H)
			(-110:\r) coordinate (M)node[below, scale=.6]{Người đu}
			($(O)!(M)!(H)$) coordinate (K)
			;
			\draw (A) arc (-55:-125:\r) (O)--(M);
			\draw[<->,dashed,thin] (M)--(K)node[midway, above, scale=.6]{$h$};
			\draw[dashed] (A)--(O)--(B) (H)--(O)node[midway, below, sloped,scale=.6]{Vị trí cân bằng};
			\foreach \i in {A,B,M,O}{\draw[fill=black](\i) circle (1pt);}
	\end{tikzpicture}}
	\loigiai{
		Do $-1 \leq \cos \left[\dfrac{\pi}{3}(2 t-1)\right] \leq 1$ nên $-3 \leq 3 \cos \left[\dfrac{\pi}{3}(2 t-1)\right] \leq 3$ hay $-3 \leq d \leq 3$.\\
		Do đó, $0 \leq|d| \leq 3$. Vậy $h$ lớn nhất bằng 3 khi $|d|=3$ hay
		\begin{eqnarray*}
			\cos\left[\dfrac{\pi}{3}(2 t-1)\right]&=& \pm 1\\
			&\Leftrightarrow& \sin \left[\dfrac{\pi}{3}(2 t-1)\right]=0 \\
			& \Leftrightarrow& \dfrac{\pi}{3}(2 t-1)=k \pi\\
			&\Leftrightarrow& t=\dfrac{1+3 k}{2} \text { với } k \in \mathbb{N}.
		\end{eqnarray*}
		Thời điểm đầu tiên mà khoång cách $h$ lớn nhất là $t=0{,}5$ s (ứng với $k=0$).
	}
\end{ex}

% #5071 | Archive/2025-Sangtac/Khoi12/Data/1-CTST/CTST-1-OTGK1-De07.tex | mã gốc 1D1V4-8
\begin{ex}%[1D1V4-8]
Một vệ tinh bay quanh Trái Đất theo một quỹ đạo hình Elip (như hình vẽ). Độ cao $h$ (tính bằng kilômet) của vệ tinh so với bề mặt Trái Đất được xác định bởi công thức $h=550+450 \cdot \cos \dfrac{\pi}{50} t$. Trong đó $t$ là thời gian tính bằng phút kể từ lúc vệ tinh bay vào quỹ đạo. Người ta cần thực hiện một thí nghiệm khoa học khi vệ tinh cách mặt đất $250$ km. Trong khoảng $60$ phút đầu tiên kể từ lúc vệ tinh bay vào quỹ đạo, hãy tìm thời điểm để có thể thực hiện thí nghiệm đó? (Làm tròn đến hàng phần trăm)
\shortans{$36{,}61$}
\begin{center}
\begin{tikzpicture}[scale=0.7, line join=round, line cap=round, >=stealth]
\draw (0,0) ellipse ({8.25} and {6.75});
\draw (-2.25,0) circle (5);
\fill (5,5.3)circle(9pt);
\draw (-11,0) node{$h=100$};
\draw [->](-9.5,0)--(-8.25,0);
\draw [->](-6,0)--(-7.25,0);
\draw [<->](4.8,5.05)--(2.1,2.5);
\draw ($(4.8,5.05)! .5 ! (2.1,2.5)$) node[below right]{$h$};
\draw [<->](2.75,0)--(8.25,0);
\draw ($(2.75,0)! .5 ! (8.25,0)$) node[below]{$1\,000$};
\draw (5,5.3) node[above right]{Vệ tinh};
\draw (-2.25,0) node{Trái đất};
		\end{tikzpicture}
\end{center}
\loigiai{
Ta có phương trình 
\allowdisplaybreaks
\begin{eqnarray*}
&&550+450 \cdot \cos \dfrac{\pi}{50} t=250 \Leftrightarrow \cos \dfrac{\pi}{50} t=-\dfrac{2}{3}\Leftrightarrow \cos \dfrac{\pi}{50} t=\cos(2{,}3) \\
&\Leftrightarrow&\hoac{&\dfrac {\pi}{50}t \approx 2{,}3 + k2\pi\\&\dfrac{\pi}{50}t \approx -2{,}3+k2\pi}\Leftrightarrow \hoac{&t\approx 36{,}61+k100\\&
t\approx-36{,}61+k100}, k\in \mathbb{Z}
\end{eqnarray*}
Vậy trong khoảng $60$ phút đầu tiên kể từ lúc vệ tinh bay vào quỹ đạo, tại thời điểm $t \approx 36{,}61$ (phút) thì ta có thể thực hiện thí nghiệm đó.}
\end{ex}

% #5091 | Archive/2025-Sangtac/SP-Main/Data/1-KNTT/KNTT_1-OTGK1_De02.tex | mã gốc 1D1V4-8
\begin{ex}%[1D1V4-8]
	Hằng ngày, mực nước của một con kênh lên xuống theo thuỷ triều. Độ sâu $h$ (m) của mực nước trong kênh tính theo thời gian $t$ (giờ) trong một ngày $(0 \leq t<24)$ cho bởi công thức $h=\cos \left(\dfrac{\pi t}{6}+2 \pi\right)+5$. Hỏi trong ngày vị trí nước xuống mức thấp nhất là bao nhiêu mét?
	\shortans[0]{$4$}
	\loigiai{
		Với mọi $0 \leq t<24$ ta có
		$$-1\leq \cos \left(\dfrac{\pi t}{6}+2 \pi\right)\leq 1 \Leftrightarrow 4\leq \cos \left(\dfrac{\pi t}{6}+2 \pi\right)+5\leq 6.$$
			Mực nước thấp nhất của kênh $h=4$ (m) khi
		\begin{align*}
			\cos \left(\dfrac{\pi t}{6}+2 \pi\right)=-1 \Leftrightarrow \dfrac{\pi t}{6}+2\pi=\pi+k2\pi \Leftrightarrow t=-6+12k, k\in \mathbb{Z}.
		\end{align*} 
		Do $0\le t< 24$ nên $k=1$, $t=6$ và  $k=2$, $t=18$.\\
		Vậy lúc $6$ và $18$ giờ mực nước kênh là thấp nhất và $h=4$ (m).
	}
\end{ex}

% #5092 | Archive/2025-Sangtac/SP-Main/Data/1-KNTT/KNTT_1-OTGK1_De02.tex | mã gốc 1D1V5-6
\begin{ex}%[1D1V5-6]
	Cho hai vật dao động điều hòa theo phương trình lần lượt là $x_1(t)=8\cos\left(4\pi t+\dfrac{\pi}{2}\right)$ $(\mathrm{cm})$ và $x_2(t)=-8\cos (4\pi t)$ $(\mathrm{cm})$, trong đó $x_1(t)$, $x_2(t)$ là li độ của hai vật tại thời điểm $t$ (giây). Khi hai vật dao động trong thời gian từ $0$ đến $10$ giây, hỏi chúng có cùng li độ mấy lần?
	\shortans[0]{$40$}
	\loigiai{
			Thời điểm hai vật có cùng li độ, ta có $8\cos\left(4\pi t+\dfrac{\pi}{2}\right)=-8\cos (4\pi t)$
			$$\begin{aligned}
				&\Leftrightarrow 
				 \cos\left(4\pi t+\dfrac{\pi}{2}\right)=\cos (\pi-4\pi t)\\
				 & \Leftrightarrow \hoac{&4\pi t+\dfrac{\pi}{2}=\pi-4\pi t+k2\pi\\&4\pi t+\dfrac{\pi}{2}=-\pi+4\pi t+k2\pi}\\
				 & \Leftrightarrow \hoac{&8\pi t=\dfrac{\pi}{2}+k2\pi\\&0t=-\dfrac{3\pi}{2}+k2\pi\text{: vô nghiệm}}\\
				 & \Leftrightarrow 
				 t=\dfrac{1}{16}+\dfrac{k}{4} \ (k \in \mathbb{Z}).
			\end{aligned}$$
			Xét $t \in [0;10]$, ta có $0 \leq \dfrac{1}{16}+\dfrac{k}{4} \leq 10 \Leftrightarrow -\dfrac{1}{4} \leq k \leq \dfrac{159}{4} \Leftrightarrow k \in \{0;1;2;\ldots;39\}$ (do $k \in \mathbb{Z}$).\\
			Vậy trong $10$ giây đầu tiên, hai vật có cùng li độ $40$ lần.
		}
\end{ex}

% #5099 | Archive/2025-Sangtac/SP-Main/Data/1-KNTT/KNTT_1-OTGK1_De06.tex | mã gốc 1D1V2-5
\begin{ex}%[2025-K11-SangtacdeGKI]%[ChuongtrinhKNTT, Trần Lê Tuấn Anh]%[1D1V2-5]
	Trong thực tế, người ta có sử dụng lượng giác trong việc sử dụng GPS để kiểm tra vị trí của vệ tinh từ tọa độ trên Trái Đất. Cho $ f(x) = 26\,750 \cos x\cos(\pi-x)$ (km); $g(x)=26\,750\cos x\sin(\pi-x)$ (km); $h(x)=26\,750\sin x$ (km) ($x$ tính theo đơn vị radian) lần lượt là vị trí tọa độ của vệ tinh xác định trong không gian thay đổi theo tọa độ $x$ trên Trái Đất. Khi vị trí trên Trái Đất ở điểm $x=30^\circ$ thì vị trí của vệ tinh sẽ được biểu thị dưới dạng tọa độ là $(a;b;c)$, với $a=f(x),b=g(x),c=h(x)$. Khi đó giá trị của biểu thức $P=\dfrac{a-b+c}{100}$ là bao nhiêu (làm tròn kết quả đến hàng đơn vị)?
	\shortans[0]{$-183$}
	\loigiai{
		Ta có $30^\circ=\dfrac{\pi}{6}$.\\

		$a=f\left(\dfrac{\pi}{6}\right)= 26\,750 \cos\left(\dfrac{\pi}{6}\right)\cos\left(\pi-\dfrac{\pi}{6}\right)=-20\,062{,}5$.\\[2mm]

		$b=g\left(\dfrac{\pi}{6}\right)=26\,750\cos\left(\dfrac{\pi}{6}\right)\sin\left(\pi-\dfrac{\pi}{6}\right)=\dfrac{13\,375\sqrt{3}}{2}$.\\[2mm]

		$c=h\left(\dfrac{\pi}{6}\right)=26\,750\sin\left(\dfrac{\pi}{6}\right)=13\,375$.\\[2mm]

		$P=\dfrac{a-b+c}{100}=\dfrac{-20\,062{,}5-\dfrac{13\,375\sqrt{3}}{2}+13\,375}{100}\approx -183$.
	}
\end{ex}

% #5100 | Archive/2025-Sangtac/SP-Main/Data/1-KNTT/KNTT_1-OTGK1_De06.tex | mã gốc 1D1V5-6
\begin{ex}%[2025-K11-SangtacdeGKI]%[ChuongtrinhKNTT, Trần Lê Tuấn Anh]%[1D1V5-6]
	Một con lắc đơn dao động điều hòa xung quanh vị trí cân bằng theo phương trình $x=\sqrt{2}\cos\left(2\pi t+\dfrac{\pi}{3}\right)$ (cm) với $t$ đơn vị là giây. Kể từ lúc bắt đầu chuyển động, sau $8$ giây, vật đã đi qua vị trí cân bằng mấy lần?
	\shortans[0]{$16$}
	\loigiai{
		Vị trí cân bằng của vật dao động điều hòa là khi $x=0$, ta có
		\begin{eqnarray*}
			\sqrt{2}\cos\left(2\pi t+\dfrac{\pi}{3}\right)&=&0\\
			\cos\left(2\pi t+\dfrac{\pi}{3}\right)&=&0\\
			2\pi t+\dfrac{\pi}{3}&=&\dfrac{\pi}{2}+k\pi, k \in \mathbb{Z}\\
			t&=&\dfrac{1}{12}+\dfrac{k}{2}, k \in \mathbb{Z}
		\end{eqnarray*}
		Trong khoảng thời gian từ $0$ đến $8$ giây, tức là $0\leq t\leq 8$ hay $0\leq \dfrac{1}{12}+\dfrac{k}{2}\leq 8$\\
		$\Rightarrow \dfrac{-1}{6}\leq k\leq \dfrac{95}{6}$.\\
		Vì $k\in \mathbb{Z}$ nên $k \in \{0;1;2;3;4;5;6;7;8;9;10;11;12;13;14;15\}$. \\
		Vậy trong khoảng thời gian từ $0$ đến $8$  giây, vật đi qua vị trí cân bằng $16$ lần.
	}
\end{ex}

% #5107 | Archive/2025-Sangtac/SP-Main/Data/1-KNTT/KNTT_1-OTGK1_De04.tex | mã gốc 1D1H4-8
\begin{ex}%[1D1H4-8]%[Dự án soạn đề kiểm tra Toán 11 GHKI NH24-25 - Nguyen Huynh]%[KNTT_1-OTGK1-De04]
	Hằng ngày mực nước con kênh lên xuống theo thủy triều. Độ sâu $ h$ (mét) của mực nước trong kênh được tính tại thời điểm $t$ giờ ($0\le t\le24$) trong một ngày bởi công thức $ h=3\cos\left(\dfrac{\pi t}{8}+\dfrac{\pi}{4}\right)+12$. Biết mực nước của kênh cao nhất khi $t=t_0$, tính giá trị của $t_0$.
	\shortans[0]{14}
	\loigiai{
		Mực nước cao nhất của kênh đạt được khi
		\begin{align*}
			\cos \left(\frac{\pi t_0}{8}+\frac{\pi}{4}\right)=1 \Leftrightarrow \frac{\pi t_0}{8}+\frac{\pi}{4}= k2\pi, k\in\mathbb{Z} \Leftrightarrow t_0=16k -2, k\in \mathbb{Z}.
		\end{align*} 
		Do $0\le t_0\le 24$ nên $k=1$ và $t_0=14$.
	}
\end{ex}
